\documentclass[landscape]{article}

% Test sheet for the Group 4 derived letters (hooks, tails, curls, topbars, descenders), 2026-11-11.
% Grids: .venv/bin/python run_print_combo_matrix.py (from the repo root, after run_patch_font.py).
% PDF: cd tex && xelatex derived-group4.tex

\usepackage[table]{xcolor}
\usepackage{fontspec}
\usepackage{geometry}
\usepackage{array}

\geometry{a4paper, left=1.5cm, right=1.5cm, top=1.5cm, bottom=1.5cm}
\setlength{\parindent}{0pt}
\addtolength{\parskip}{2pt}

% always use `Path = ../fonts/`, always include the extension
\setmainfont{Libertinus Serif}[
	Path = ../fonts/,
	UprightFont    = LibertinusSerif-Regular.otf,
	ItalicFont     = LibertinusSerif-Italic.otf,
	BoldFont       = LibertinusSerif-Semibold.otf,
	BoldItalicFont = LibertinusSerif-SemiboldItalic.otf
]

\newfontfamily\LibertinusSerifPatch[
	Path = ../fonts/,
	UprightFont    = LibertinusSerif-Regular-patch.otf,
	ItalicFont     = LibertinusSerif-Italic-patch.otf,
	BoldFont       = LibertinusSerif-Semibold-patch.otf,
	BoldItalicFont = LibertinusSerif-SemiboldItalic-patch.otf
]{Libertinus Serif Patch}

\newcommand{\gs}{\hspace{0.5mm}}
\newcommand{\cgj}{\char"034F }

\definecolor{supportedanchor}{RGB}{0,0,0}
\definecolor{supportedprec}{RGB}{0,0,180}
\definecolor{supportedfall}{RGB}{180,0,0}
\definecolor{unsupportedanchor}{RGB}{180,180,180}
\definecolor{unsupportedprec}{RGB}{150,190,255}
\definecolor{unsupportedfall}{RGB}{255,180,200}
\definecolor{missingglyph}{RGB}{200,200,200}
\newcommand{\SUPPORTEDANCH}[1]{{\color{supportedanchor}#1}}
\newcommand{\SUPPORTEDPREC}[1]{{\color{supportedprec}#1}}
\newcommand{\SUPPORTEDFALL}[1]{{\color{supportedfall}#1}}
\newcommand{\UNSUPPORTEDANCH}[1]{{\color{unsupportedanchor}#1}}
\newcommand{\UNSUPPORTEDPREC}[1]{{\color{unsupportedprec}#1}}
\newcommand{\UNSUPPORTEDFALL}[1]{{\color{unsupportedfall}#1}}
\newcommand{\MISSINGGLYPH}[1]{{\color{missingglyph}#1}}

\begin{document}

\section*{Group 4 derived letters: test sheet}

Letters with a hook, tail, curl, topbar or descender, added to stage 8 on 2026-11-11 (\verb|data/derived_letters.py|), each after its base letter.
Each grid row is one IPA mark on every letter; black is positioned by an anchor, red is the fallback (no anchor), blue a precomposed glyph, gray a missing glyph.
The original fonts come first, then the patched fonts.

\begin{tabular}{>{\raggedright}p{0.17\linewidth}>{\raggedright}p{0.28\linewidth}p{0.47\linewidth}}
letters & rule & why \\
\hline
Ƥ ƃ ƌ Ƴ ƴ ƙ ƈ ẚ ꭒ & inherit both (fit) & the hook, topbar or ring stays clear of the marks; italic and semibold italic redrawings keep stage 6's candidate \\
Ƭ ƭ & stem shift (\verb|stem_axis|) & the hook at the end of T's bar spoils the fit above; ƭ's hook replaces the top of t's stem \\
ƥ ʠ ƭ above & \verb|derived_y|: $X$ from p, q, t, $Y$ at the ascender reference & the hook rises from the stem to the ascender, under the mark (it would move $Y$ by $+102$ to $+256$) \\
ƫ Ʈ ȡ ȴ ȵ ȶ Ɽ Ⱨ ⱨ Ⱪ ⱪ Ɱ ỿ & \verb|bottom_tail| & stems matched above the tail; below, the candidate's $Y$ where the tail is under the mark, else the base letter's (Ⱨ) \\
Ⱬ ⱬ & not derived & Z has no stems: stage 6's candidate \\
\end{tabular}

\newpage
\subsection*{Regular, Group 4, above, Above}

% grid. columns are bases, rows are marks.
\SUPPORTEDANCH{\char"0048\char"0300} \SUPPORTEDFALL{\char"2C67\char"0300} \SUPPORTEDANCH{\char"004B\char"0300} \SUPPORTEDANCH{\char"2C69\char"0300} \SUPPORTEDANCH{\char"004D\char"0300} \SUPPORTEDANCH{\char"2C6E\char"0300} \SUPPORTEDANCH{\char"0050\char"0300} \SUPPORTEDANCH{\char"01A4\char"0300} \SUPPORTEDANCH{\char"0052\char"0300} \SUPPORTEDFALL{\char"2C64\char"0300} \SUPPORTEDANCH{\char"0054\char"0300} \SUPPORTEDFALL{\char"01AC\char"0300} \SUPPORTEDFALL{\char"01AE\char"0300} \SUPPORTEDPREC{\char"0059\char"0300} \SUPPORTEDFALL{\char"01B3\char"0300} \SUPPORTEDPREC{\char"0061\char"0300} \SUPPORTEDFALL{\char"1E9A\char"0300} \SUPPORTEDFALL{\char"0062\char"0300} \SUPPORTEDFALL{\char"0183\char"0300} \SUPPORTEDANCH{\char"0063\char"0300} \SUPPORTEDFALL{\char"0188\char"0300} \SUPPORTEDANCH{\char"0064\char"0300} \SUPPORTEDANCH{\char"018C\char"0300} \SUPPORTEDFALL{\char"0221\char"0300} \SUPPORTEDANCH{\char"0068\char"0300} \SUPPORTEDANCH{\char"2C68\char"0300} \SUPPORTEDANCH{\char"006B\char"0300} \SUPPORTEDFALL{\char"0199\char"0300} \SUPPORTEDANCH{\char"2C6A\char"0300} \SUPPORTEDANCH{\char"006C\char"0300} \SUPPORTEDFALL{\char"0234\char"0300} \SUPPORTEDPREC{\char"006E\char"0300} \SUPPORTEDANCH{\char"0235\char"0300} \SUPPORTEDANCH{\char"0070\char"0300} \SUPPORTEDFALL{\char"01A5\char"0300} \SUPPORTEDANCH{\char"0071\char"0300} \SUPPORTEDANCH{\char"02A0\char"0300} \SUPPORTEDFALL{\char"0074\char"0300} \SUPPORTEDFALL{\char"01AB\char"0300} \SUPPORTEDFALL{\char"01AD\char"0300} \SUPPORTEDFALL{\char"0236\char"0300} \SUPPORTEDPREC{\char"0075\char"0300} \SUPPORTEDANCH{\char"AB52\char"0300} \SUPPORTEDPREC{\char"0079\char"0300} \SUPPORTEDANCH{\char"01B4\char"0300} \SUPPORTEDANCH{\char"1EFF\char"0300}

\SUPPORTEDANCH{\char"0048\char"0301} \SUPPORTEDFALL{\char"2C67\char"0301} \SUPPORTEDPREC{\char"004B\char"0301} \SUPPORTEDANCH{\char"2C69\char"0301} \SUPPORTEDPREC{\char"004D\char"0301} \SUPPORTEDANCH{\char"2C6E\char"0301} \SUPPORTEDPREC{\char"0050\char"0301} \SUPPORTEDANCH{\char"01A4\char"0301} \SUPPORTEDPREC{\char"0052\char"0301} \SUPPORTEDFALL{\char"2C64\char"0301} \SUPPORTEDANCH{\char"0054\char"0301} \SUPPORTEDFALL{\char"01AC\char"0301} \SUPPORTEDFALL{\char"01AE\char"0301} \SUPPORTEDPREC{\char"0059\char"0301} \SUPPORTEDFALL{\char"01B3\char"0301} \SUPPORTEDPREC{\char"0061\char"0301} \SUPPORTEDFALL{\char"1E9A\char"0301} \SUPPORTEDFALL{\char"0062\char"0301} \SUPPORTEDFALL{\char"0183\char"0301} \SUPPORTEDPREC{\char"0063\char"0301} \SUPPORTEDFALL{\char"0188\char"0301} \SUPPORTEDANCH{\char"0064\char"0301} \SUPPORTEDANCH{\char"018C\char"0301} \SUPPORTEDFALL{\char"0221\char"0301} \SUPPORTEDANCH{\char"0068\char"0301} \SUPPORTEDANCH{\char"2C68\char"0301} \SUPPORTEDPREC{\char"006B\char"0301} \SUPPORTEDFALL{\char"0199\char"0301} \SUPPORTEDANCH{\char"2C6A\char"0301} \SUPPORTEDPREC{\char"006C\char"0301} \SUPPORTEDFALL{\char"0234\char"0301} \SUPPORTEDPREC{\char"006E\char"0301} \SUPPORTEDANCH{\char"0235\char"0301} \SUPPORTEDPREC{\char"0070\char"0301} \SUPPORTEDFALL{\char"01A5\char"0301} \SUPPORTEDANCH{\char"0071\char"0301} \SUPPORTEDANCH{\char"02A0\char"0301} \SUPPORTEDFALL{\char"0074\char"0301} \SUPPORTEDFALL{\char"01AB\char"0301} \SUPPORTEDFALL{\char"01AD\char"0301} \SUPPORTEDFALL{\char"0236\char"0301} \SUPPORTEDPREC{\char"0075\char"0301} \SUPPORTEDANCH{\char"AB52\char"0301} \SUPPORTEDPREC{\char"0079\char"0301} \SUPPORTEDANCH{\char"01B4\char"0301} \SUPPORTEDANCH{\char"1EFF\char"0301}

\SUPPORTEDPREC{\char"0048\char"0302} \SUPPORTEDFALL{\char"2C67\char"0302} \SUPPORTEDANCH{\char"004B\char"0302} \SUPPORTEDANCH{\char"2C69\char"0302} \SUPPORTEDANCH{\char"004D\char"0302} \SUPPORTEDANCH{\char"2C6E\char"0302} \SUPPORTEDANCH{\char"0050\char"0302} \SUPPORTEDANCH{\char"01A4\char"0302} \SUPPORTEDANCH{\char"0052\char"0302} \SUPPORTEDFALL{\char"2C64\char"0302} \SUPPORTEDANCH{\char"0054\char"0302} \SUPPORTEDFALL{\char"01AC\char"0302} \SUPPORTEDFALL{\char"01AE\char"0302} \SUPPORTEDPREC{\char"0059\char"0302} \SUPPORTEDFALL{\char"01B3\char"0302} \SUPPORTEDPREC{\char"0061\char"0302} \SUPPORTEDFALL{\char"1E9A\char"0302} \SUPPORTEDFALL{\char"0062\char"0302} \SUPPORTEDFALL{\char"0183\char"0302} \SUPPORTEDPREC{\char"0063\char"0302} \SUPPORTEDFALL{\char"0188\char"0302} \SUPPORTEDANCH{\char"0064\char"0302} \SUPPORTEDANCH{\char"018C\char"0302} \SUPPORTEDFALL{\char"0221\char"0302} \SUPPORTEDPREC{\char"0068\char"0302} \SUPPORTEDANCH{\char"2C68\char"0302} \SUPPORTEDANCH{\char"006B\char"0302} \SUPPORTEDFALL{\char"0199\char"0302} \SUPPORTEDANCH{\char"2C6A\char"0302} \SUPPORTEDANCH{\char"006C\char"0302} \SUPPORTEDFALL{\char"0234\char"0302} \SUPPORTEDANCH{\char"006E\char"0302} \SUPPORTEDANCH{\char"0235\char"0302} \SUPPORTEDANCH{\char"0070\char"0302} \SUPPORTEDFALL{\char"01A5\char"0302} \SUPPORTEDANCH{\char"0071\char"0302} \SUPPORTEDANCH{\char"02A0\char"0302} \SUPPORTEDFALL{\char"0074\char"0302} \SUPPORTEDFALL{\char"01AB\char"0302} \SUPPORTEDFALL{\char"01AD\char"0302} \SUPPORTEDFALL{\char"0236\char"0302} \SUPPORTEDPREC{\char"0075\char"0302} \SUPPORTEDANCH{\char"AB52\char"0302} \SUPPORTEDPREC{\char"0079\char"0302} \SUPPORTEDANCH{\char"01B4\char"0302} \SUPPORTEDANCH{\char"1EFF\char"0302}

\SUPPORTEDANCH{\char"0048\char"0303} \SUPPORTEDFALL{\char"2C67\char"0303} \SUPPORTEDANCH{\char"004B\char"0303} \SUPPORTEDANCH{\char"2C69\char"0303} \SUPPORTEDANCH{\char"004D\char"0303} \SUPPORTEDANCH{\char"2C6E\char"0303} \SUPPORTEDANCH{\char"0050\char"0303} \SUPPORTEDANCH{\char"01A4\char"0303} \SUPPORTEDANCH{\char"0052\char"0303} \SUPPORTEDFALL{\char"2C64\char"0303} \SUPPORTEDANCH{\char"0054\char"0303} \SUPPORTEDFALL{\char"01AC\char"0303} \SUPPORTEDFALL{\char"01AE\char"0303} \SUPPORTEDPREC{\char"0059\char"0303} \SUPPORTEDFALL{\char"01B3\char"0303} \SUPPORTEDPREC{\char"0061\char"0303} \SUPPORTEDFALL{\char"1E9A\char"0303} \SUPPORTEDFALL{\char"0062\char"0303} \SUPPORTEDFALL{\char"0183\char"0303} \SUPPORTEDANCH{\char"0063\char"0303} \SUPPORTEDFALL{\char"0188\char"0303} \SUPPORTEDANCH{\char"0064\char"0303} \SUPPORTEDANCH{\char"018C\char"0303} \SUPPORTEDFALL{\char"0221\char"0303} \SUPPORTEDANCH{\char"0068\char"0303} \SUPPORTEDANCH{\char"2C68\char"0303} \SUPPORTEDANCH{\char"006B\char"0303} \SUPPORTEDFALL{\char"0199\char"0303} \SUPPORTEDANCH{\char"2C6A\char"0303} \SUPPORTEDANCH{\char"006C\char"0303} \SUPPORTEDFALL{\char"0234\char"0303} \SUPPORTEDPREC{\char"006E\char"0303} \SUPPORTEDANCH{\char"0235\char"0303} \SUPPORTEDANCH{\char"0070\char"0303} \SUPPORTEDFALL{\char"01A5\char"0303} \SUPPORTEDANCH{\char"0071\char"0303} \SUPPORTEDANCH{\char"02A0\char"0303} \SUPPORTEDFALL{\char"0074\char"0303} \SUPPORTEDFALL{\char"01AB\char"0303} \SUPPORTEDFALL{\char"01AD\char"0303} \SUPPORTEDFALL{\char"0236\char"0303} \SUPPORTEDPREC{\char"0075\char"0303} \SUPPORTEDANCH{\char"AB52\char"0303} \SUPPORTEDPREC{\char"0079\char"0303} \SUPPORTEDANCH{\char"01B4\char"0303} \SUPPORTEDANCH{\char"1EFF\char"0303}

\SUPPORTEDANCH{\char"0048\char"0304} \SUPPORTEDFALL{\char"2C67\char"0304} \SUPPORTEDANCH{\char"004B\char"0304} \SUPPORTEDANCH{\char"2C69\char"0304} \SUPPORTEDANCH{\char"004D\char"0304} \SUPPORTEDANCH{\char"2C6E\char"0304} \SUPPORTEDANCH{\char"0050\char"0304} \SUPPORTEDANCH{\char"01A4\char"0304} \SUPPORTEDANCH{\char"0052\char"0304} \SUPPORTEDFALL{\char"2C64\char"0304} \SUPPORTEDANCH{\char"0054\char"0304} \SUPPORTEDFALL{\char"01AC\char"0304} \SUPPORTEDFALL{\char"01AE\char"0304} \SUPPORTEDPREC{\char"0059\char"0304} \SUPPORTEDFALL{\char"01B3\char"0304} \SUPPORTEDPREC{\char"0061\char"0304} \SUPPORTEDFALL{\char"1E9A\char"0304} \SUPPORTEDFALL{\char"0062\char"0304} \SUPPORTEDFALL{\char"0183\char"0304} \SUPPORTEDANCH{\char"0063\char"0304} \SUPPORTEDFALL{\char"0188\char"0304} \SUPPORTEDANCH{\char"0064\char"0304} \SUPPORTEDANCH{\char"018C\char"0304} \SUPPORTEDFALL{\char"0221\char"0304} \SUPPORTEDANCH{\char"0068\char"0304} \SUPPORTEDANCH{\char"2C68\char"0304} \SUPPORTEDANCH{\char"006B\char"0304} \SUPPORTEDFALL{\char"0199\char"0304} \SUPPORTEDANCH{\char"2C6A\char"0304} \SUPPORTEDANCH{\char"006C\char"0304} \SUPPORTEDFALL{\char"0234\char"0304} \SUPPORTEDANCH{\char"006E\char"0304} \SUPPORTEDANCH{\char"0235\char"0304} \SUPPORTEDANCH{\char"0070\char"0304} \SUPPORTEDFALL{\char"01A5\char"0304} \SUPPORTEDANCH{\char"0071\char"0304} \SUPPORTEDANCH{\char"02A0\char"0304} \SUPPORTEDFALL{\char"0074\char"0304} \SUPPORTEDFALL{\char"01AB\char"0304} \SUPPORTEDFALL{\char"01AD\char"0304} \SUPPORTEDFALL{\char"0236\char"0304} \SUPPORTEDPREC{\char"0075\char"0304} \SUPPORTEDANCH{\char"AB52\char"0304} \SUPPORTEDPREC{\char"0079\char"0304} \SUPPORTEDANCH{\char"01B4\char"0304} \SUPPORTEDANCH{\char"1EFF\char"0304}

\SUPPORTEDANCH{\char"0048\char"0306} \SUPPORTEDFALL{\char"2C67\char"0306} \SUPPORTEDANCH{\char"004B\char"0306} \SUPPORTEDANCH{\char"2C69\char"0306} \SUPPORTEDANCH{\char"004D\char"0306} \SUPPORTEDANCH{\char"2C6E\char"0306} \SUPPORTEDANCH{\char"0050\char"0306} \SUPPORTEDANCH{\char"01A4\char"0306} \SUPPORTEDANCH{\char"0052\char"0306} \SUPPORTEDFALL{\char"2C64\char"0306} \SUPPORTEDANCH{\char"0054\char"0306} \SUPPORTEDFALL{\char"01AC\char"0306} \SUPPORTEDFALL{\char"01AE\char"0306} \SUPPORTEDANCH{\char"0059\char"0306} \SUPPORTEDFALL{\char"01B3\char"0306} \SUPPORTEDPREC{\char"0061\char"0306} \SUPPORTEDFALL{\char"1E9A\char"0306} \SUPPORTEDFALL{\char"0062\char"0306} \SUPPORTEDFALL{\char"0183\char"0306} \SUPPORTEDANCH{\char"0063\char"0306} \SUPPORTEDFALL{\char"0188\char"0306} \SUPPORTEDANCH{\char"0064\char"0306} \SUPPORTEDANCH{\char"018C\char"0306} \SUPPORTEDFALL{\char"0221\char"0306} \SUPPORTEDANCH{\char"0068\char"0306} \SUPPORTEDANCH{\char"2C68\char"0306} \SUPPORTEDANCH{\char"006B\char"0306} \SUPPORTEDFALL{\char"0199\char"0306} \SUPPORTEDANCH{\char"2C6A\char"0306} \SUPPORTEDANCH{\char"006C\char"0306} \SUPPORTEDFALL{\char"0234\char"0306} \SUPPORTEDANCH{\char"006E\char"0306} \SUPPORTEDANCH{\char"0235\char"0306} \SUPPORTEDANCH{\char"0070\char"0306} \SUPPORTEDFALL{\char"01A5\char"0306} \SUPPORTEDANCH{\char"0071\char"0306} \SUPPORTEDANCH{\char"02A0\char"0306} \SUPPORTEDFALL{\char"0074\char"0306} \SUPPORTEDFALL{\char"01AB\char"0306} \SUPPORTEDFALL{\char"01AD\char"0306} \SUPPORTEDFALL{\char"0236\char"0306} \SUPPORTEDPREC{\char"0075\char"0306} \SUPPORTEDANCH{\char"AB52\char"0306} \SUPPORTEDANCH{\char"0079\char"0306} \SUPPORTEDANCH{\char"01B4\char"0306} \SUPPORTEDANCH{\char"1EFF\char"0306}

\SUPPORTEDPREC{\char"0048\char"0307} \SUPPORTEDFALL{\char"2C67\char"0307} \SUPPORTEDANCH{\char"004B\char"0307} \SUPPORTEDANCH{\char"2C69\char"0307} \SUPPORTEDPREC{\char"004D\char"0307} \SUPPORTEDANCH{\char"2C6E\char"0307} \SUPPORTEDPREC{\char"0050\char"0307} \SUPPORTEDANCH{\char"01A4\char"0307} \SUPPORTEDPREC{\char"0052\char"0307} \SUPPORTEDFALL{\char"2C64\char"0307} \SUPPORTEDPREC{\char"0054\char"0307} \SUPPORTEDFALL{\char"01AC\char"0307} \SUPPORTEDFALL{\char"01AE\char"0307} \SUPPORTEDPREC{\char"0059\char"0307} \SUPPORTEDFALL{\char"01B3\char"0307} \SUPPORTEDPREC{\char"0061\char"0307} \SUPPORTEDFALL{\char"1E9A\char"0307} \SUPPORTEDPREC{\char"0062\char"0307} \SUPPORTEDFALL{\char"0183\char"0307} \SUPPORTEDPREC{\char"0063\char"0307} \SUPPORTEDFALL{\char"0188\char"0307} \SUPPORTEDPREC{\char"0064\char"0307} \SUPPORTEDANCH{\char"018C\char"0307} \SUPPORTEDFALL{\char"0221\char"0307} \SUPPORTEDPREC{\char"0068\char"0307} \SUPPORTEDANCH{\char"2C68\char"0307} \SUPPORTEDANCH{\char"006B\char"0307} \SUPPORTEDFALL{\char"0199\char"0307} \SUPPORTEDANCH{\char"2C6A\char"0307} \SUPPORTEDANCH{\char"006C\char"0307} \SUPPORTEDFALL{\char"0234\char"0307} \SUPPORTEDPREC{\char"006E\char"0307} \SUPPORTEDANCH{\char"0235\char"0307} \SUPPORTEDPREC{\char"0070\char"0307} \SUPPORTEDFALL{\char"01A5\char"0307} \SUPPORTEDANCH{\char"0071\char"0307} \SUPPORTEDANCH{\char"02A0\char"0307} \SUPPORTEDPREC{\char"0074\char"0307} \SUPPORTEDFALL{\char"01AB\char"0307} \SUPPORTEDFALL{\char"01AD\char"0307} \SUPPORTEDFALL{\char"0236\char"0307} \SUPPORTEDANCH{\char"0075\char"0307} \SUPPORTEDANCH{\char"AB52\char"0307} \SUPPORTEDPREC{\char"0079\char"0307} \SUPPORTEDANCH{\char"01B4\char"0307} \SUPPORTEDANCH{\char"1EFF\char"0307}

\SUPPORTEDPREC{\char"0048\char"0308} \SUPPORTEDFALL{\char"2C67\char"0308} \SUPPORTEDANCH{\char"004B\char"0308} \SUPPORTEDANCH{\char"2C69\char"0308} \SUPPORTEDANCH{\char"004D\char"0308} \SUPPORTEDANCH{\char"2C6E\char"0308} \SUPPORTEDANCH{\char"0050\char"0308} \SUPPORTEDANCH{\char"01A4\char"0308} \SUPPORTEDANCH{\char"0052\char"0308} \SUPPORTEDFALL{\char"2C64\char"0308} \SUPPORTEDANCH{\char"0054\char"0308} \SUPPORTEDFALL{\char"01AC\char"0308} \SUPPORTEDFALL{\char"01AE\char"0308} \SUPPORTEDPREC{\char"0059\char"0308} \SUPPORTEDFALL{\char"01B3\char"0308} \SUPPORTEDPREC{\char"0061\char"0308} \SUPPORTEDFALL{\char"1E9A\char"0308} \SUPPORTEDFALL{\char"0062\char"0308} \SUPPORTEDFALL{\char"0183\char"0308} \SUPPORTEDANCH{\char"0063\char"0308} \SUPPORTEDFALL{\char"0188\char"0308} \SUPPORTEDANCH{\char"0064\char"0308} \SUPPORTEDANCH{\char"018C\char"0308} \SUPPORTEDFALL{\char"0221\char"0308} \SUPPORTEDPREC{\char"0068\char"0308} \SUPPORTEDANCH{\char"2C68\char"0308} \SUPPORTEDANCH{\char"006B\char"0308} \SUPPORTEDFALL{\char"0199\char"0308} \SUPPORTEDANCH{\char"2C6A\char"0308} \SUPPORTEDANCH{\char"006C\char"0308} \SUPPORTEDFALL{\char"0234\char"0308} \SUPPORTEDANCH{\char"006E\char"0308} \SUPPORTEDANCH{\char"0235\char"0308} \SUPPORTEDANCH{\char"0070\char"0308} \SUPPORTEDFALL{\char"01A5\char"0308} \SUPPORTEDANCH{\char"0071\char"0308} \SUPPORTEDANCH{\char"02A0\char"0308} \SUPPORTEDPREC{\char"0074\char"0308} \SUPPORTEDFALL{\char"01AB\char"0308} \SUPPORTEDFALL{\char"01AD\char"0308} \SUPPORTEDFALL{\char"0236\char"0308} \SUPPORTEDPREC{\char"0075\char"0308} \SUPPORTEDANCH{\char"AB52\char"0308} \SUPPORTEDPREC{\char"0079\char"0308} \SUPPORTEDANCH{\char"01B4\char"0308} \SUPPORTEDANCH{\char"1EFF\char"0308}

\SUPPORTEDANCH{\char"0048\char"0309} \SUPPORTEDFALL{\char"2C67\char"0309} \SUPPORTEDANCH{\char"004B\char"0309} \SUPPORTEDANCH{\char"2C69\char"0309} \SUPPORTEDANCH{\char"004D\char"0309} \SUPPORTEDANCH{\char"2C6E\char"0309} \SUPPORTEDANCH{\char"0050\char"0309} \SUPPORTEDANCH{\char"01A4\char"0309} \SUPPORTEDANCH{\char"0052\char"0309} \SUPPORTEDFALL{\char"2C64\char"0309} \SUPPORTEDANCH{\char"0054\char"0309} \SUPPORTEDFALL{\char"01AC\char"0309} \SUPPORTEDFALL{\char"01AE\char"0309} \SUPPORTEDPREC{\char"0059\char"0309} \SUPPORTEDFALL{\char"01B3\char"0309} \SUPPORTEDPREC{\char"0061\char"0309} \SUPPORTEDFALL{\char"1E9A\char"0309} \SUPPORTEDFALL{\char"0062\char"0309} \SUPPORTEDFALL{\char"0183\char"0309} \SUPPORTEDANCH{\char"0063\char"0309} \SUPPORTEDFALL{\char"0188\char"0309} \SUPPORTEDANCH{\char"0064\char"0309} \SUPPORTEDANCH{\char"018C\char"0309} \SUPPORTEDFALL{\char"0221\char"0309} \SUPPORTEDANCH{\char"0068\char"0309} \SUPPORTEDANCH{\char"2C68\char"0309} \SUPPORTEDANCH{\char"006B\char"0309} \SUPPORTEDFALL{\char"0199\char"0309} \SUPPORTEDANCH{\char"2C6A\char"0309} \SUPPORTEDANCH{\char"006C\char"0309} \SUPPORTEDFALL{\char"0234\char"0309} \SUPPORTEDANCH{\char"006E\char"0309} \SUPPORTEDANCH{\char"0235\char"0309} \SUPPORTEDANCH{\char"0070\char"0309} \SUPPORTEDFALL{\char"01A5\char"0309} \SUPPORTEDANCH{\char"0071\char"0309} \SUPPORTEDANCH{\char"02A0\char"0309} \SUPPORTEDFALL{\char"0074\char"0309} \SUPPORTEDFALL{\char"01AB\char"0309} \SUPPORTEDFALL{\char"01AD\char"0309} \SUPPORTEDFALL{\char"0236\char"0309} \SUPPORTEDPREC{\char"0075\char"0309} \SUPPORTEDANCH{\char"AB52\char"0309} \SUPPORTEDPREC{\char"0079\char"0309} \SUPPORTEDANCH{\char"01B4\char"0309} \SUPPORTEDANCH{\char"1EFF\char"0309}

\SUPPORTEDANCH{\char"0048\char"030A} \SUPPORTEDFALL{\char"2C67\char"030A} \SUPPORTEDANCH{\char"004B\char"030A} \SUPPORTEDANCH{\char"2C69\char"030A} \SUPPORTEDANCH{\char"004D\char"030A} \SUPPORTEDANCH{\char"2C6E\char"030A} \SUPPORTEDANCH{\char"0050\char"030A} \SUPPORTEDANCH{\char"01A4\char"030A} \SUPPORTEDANCH{\char"0052\char"030A} \SUPPORTEDFALL{\char"2C64\char"030A} \SUPPORTEDANCH{\char"0054\char"030A} \SUPPORTEDFALL{\char"01AC\char"030A} \SUPPORTEDFALL{\char"01AE\char"030A} \SUPPORTEDANCH{\char"0059\char"030A} \SUPPORTEDFALL{\char"01B3\char"030A} \SUPPORTEDPREC{\char"0061\char"030A} \SUPPORTEDFALL{\char"1E9A\char"030A} \SUPPORTEDFALL{\char"0062\char"030A} \SUPPORTEDFALL{\char"0183\char"030A} \SUPPORTEDANCH{\char"0063\char"030A} \SUPPORTEDFALL{\char"0188\char"030A} \SUPPORTEDANCH{\char"0064\char"030A} \SUPPORTEDANCH{\char"018C\char"030A} \SUPPORTEDFALL{\char"0221\char"030A} \SUPPORTEDANCH{\char"0068\char"030A} \SUPPORTEDANCH{\char"2C68\char"030A} \SUPPORTEDANCH{\char"006B\char"030A} \SUPPORTEDFALL{\char"0199\char"030A} \SUPPORTEDANCH{\char"2C6A\char"030A} \SUPPORTEDANCH{\char"006C\char"030A} \SUPPORTEDFALL{\char"0234\char"030A} \SUPPORTEDANCH{\char"006E\char"030A} \SUPPORTEDANCH{\char"0235\char"030A} \SUPPORTEDANCH{\char"0070\char"030A} \SUPPORTEDFALL{\char"01A5\char"030A} \SUPPORTEDANCH{\char"0071\char"030A} \SUPPORTEDANCH{\char"02A0\char"030A} \SUPPORTEDFALL{\char"0074\char"030A} \SUPPORTEDFALL{\char"01AB\char"030A} \SUPPORTEDFALL{\char"01AD\char"030A} \SUPPORTEDFALL{\char"0236\char"030A} \SUPPORTEDPREC{\char"0075\char"030A} \SUPPORTEDANCH{\char"AB52\char"030A} \SUPPORTEDPREC{\char"0079\char"030A} \SUPPORTEDANCH{\char"01B4\char"030A} \SUPPORTEDANCH{\char"1EFF\char"030A}

\SUPPORTEDANCH{\char"0048\char"030B} \SUPPORTEDFALL{\char"2C67\char"030B} \SUPPORTEDANCH{\char"004B\char"030B} \SUPPORTEDANCH{\char"2C69\char"030B} \SUPPORTEDANCH{\char"004D\char"030B} \SUPPORTEDANCH{\char"2C6E\char"030B} \SUPPORTEDANCH{\char"0050\char"030B} \SUPPORTEDANCH{\char"01A4\char"030B} \SUPPORTEDANCH{\char"0052\char"030B} \SUPPORTEDFALL{\char"2C64\char"030B} \SUPPORTEDANCH{\char"0054\char"030B} \SUPPORTEDFALL{\char"01AC\char"030B} \SUPPORTEDFALL{\char"01AE\char"030B} \SUPPORTEDANCH{\char"0059\char"030B} \SUPPORTEDFALL{\char"01B3\char"030B} \SUPPORTEDANCH{\char"0061\char"030B} \SUPPORTEDFALL{\char"1E9A\char"030B} \SUPPORTEDFALL{\char"0062\char"030B} \SUPPORTEDFALL{\char"0183\char"030B} \SUPPORTEDANCH{\char"0063\char"030B} \SUPPORTEDFALL{\char"0188\char"030B} \SUPPORTEDANCH{\char"0064\char"030B} \SUPPORTEDANCH{\char"018C\char"030B} \SUPPORTEDFALL{\char"0221\char"030B} \SUPPORTEDANCH{\char"0068\char"030B} \SUPPORTEDANCH{\char"2C68\char"030B} \SUPPORTEDANCH{\char"006B\char"030B} \SUPPORTEDFALL{\char"0199\char"030B} \SUPPORTEDANCH{\char"2C6A\char"030B} \SUPPORTEDANCH{\char"006C\char"030B} \SUPPORTEDFALL{\char"0234\char"030B} \SUPPORTEDANCH{\char"006E\char"030B} \SUPPORTEDANCH{\char"0235\char"030B} \SUPPORTEDANCH{\char"0070\char"030B} \SUPPORTEDFALL{\char"01A5\char"030B} \SUPPORTEDANCH{\char"0071\char"030B} \SUPPORTEDANCH{\char"02A0\char"030B} \SUPPORTEDFALL{\char"0074\char"030B} \SUPPORTEDFALL{\char"01AB\char"030B} \SUPPORTEDFALL{\char"01AD\char"030B} \SUPPORTEDFALL{\char"0236\char"030B} \SUPPORTEDPREC{\char"0075\char"030B} \SUPPORTEDANCH{\char"AB52\char"030B} \SUPPORTEDANCH{\char"0079\char"030B} \SUPPORTEDANCH{\char"01B4\char"030B} \SUPPORTEDANCH{\char"1EFF\char"030B}

\SUPPORTEDPREC{\char"0048\char"030C} \SUPPORTEDFALL{\char"2C67\char"030C} \SUPPORTEDPREC{\char"004B\char"030C} \SUPPORTEDANCH{\char"2C69\char"030C} \SUPPORTEDANCH{\char"004D\char"030C} \SUPPORTEDANCH{\char"2C6E\char"030C} \SUPPORTEDANCH{\char"0050\char"030C} \SUPPORTEDANCH{\char"01A4\char"030C} \SUPPORTEDPREC{\char"0052\char"030C} \SUPPORTEDFALL{\char"2C64\char"030C} \SUPPORTEDPREC{\char"0054\char"030C} \SUPPORTEDFALL{\char"01AC\char"030C} \SUPPORTEDFALL{\char"01AE\char"030C} \SUPPORTEDANCH{\char"0059\char"030C} \SUPPORTEDFALL{\char"01B3\char"030C} \SUPPORTEDPREC{\char"0061\char"030C} \SUPPORTEDFALL{\char"1E9A\char"030C} \SUPPORTEDFALL{\char"0062\char"030C} \SUPPORTEDFALL{\char"0183\char"030C} \SUPPORTEDPREC{\char"0063\char"030C} \SUPPORTEDFALL{\char"0188\char"030C} \SUPPORTEDPREC{\char"0064\char"030C} \SUPPORTEDANCH{\char"018C\char"030C} \SUPPORTEDFALL{\char"0221\char"030C} \SUPPORTEDPREC{\char"0068\char"030C} \SUPPORTEDANCH{\char"2C68\char"030C} \SUPPORTEDPREC{\char"006B\char"030C} \SUPPORTEDFALL{\char"0199\char"030C} \SUPPORTEDANCH{\char"2C6A\char"030C} \SUPPORTEDPREC{\char"006C\char"030C} \SUPPORTEDFALL{\char"0234\char"030C} \SUPPORTEDPREC{\char"006E\char"030C} \SUPPORTEDANCH{\char"0235\char"030C} \SUPPORTEDANCH{\char"0070\char"030C} \SUPPORTEDFALL{\char"01A5\char"030C} \SUPPORTEDANCH{\char"0071\char"030C} \SUPPORTEDANCH{\char"02A0\char"030C} \SUPPORTEDPREC{\char"0074\char"030C} \SUPPORTEDFALL{\char"01AB\char"030C} \SUPPORTEDFALL{\char"01AD\char"030C} \SUPPORTEDFALL{\char"0236\char"030C} \SUPPORTEDPREC{\char"0075\char"030C} \SUPPORTEDANCH{\char"AB52\char"030C} \SUPPORTEDANCH{\char"0079\char"030C} \SUPPORTEDANCH{\char"01B4\char"030C} \SUPPORTEDANCH{\char"1EFF\char"030C}

\SUPPORTEDANCH{\char"0048\char"030D} \SUPPORTEDFALL{\char"2C67\char"030D} \SUPPORTEDANCH{\char"004B\char"030D} \SUPPORTEDANCH{\char"2C69\char"030D} \SUPPORTEDANCH{\char"004D\char"030D} \SUPPORTEDANCH{\char"2C6E\char"030D} \SUPPORTEDANCH{\char"0050\char"030D} \SUPPORTEDANCH{\char"01A4\char"030D} \SUPPORTEDANCH{\char"0052\char"030D} \SUPPORTEDFALL{\char"2C64\char"030D} \SUPPORTEDANCH{\char"0054\char"030D} \SUPPORTEDFALL{\char"01AC\char"030D} \SUPPORTEDFALL{\char"01AE\char"030D} \SUPPORTEDANCH{\char"0059\char"030D} \SUPPORTEDFALL{\char"01B3\char"030D} \SUPPORTEDANCH{\char"0061\char"030D} \SUPPORTEDFALL{\char"1E9A\char"030D} \SUPPORTEDFALL{\char"0062\char"030D} \SUPPORTEDFALL{\char"0183\char"030D} \SUPPORTEDANCH{\char"0063\char"030D} \SUPPORTEDFALL{\char"0188\char"030D} \SUPPORTEDANCH{\char"0064\char"030D} \SUPPORTEDANCH{\char"018C\char"030D} \SUPPORTEDFALL{\char"0221\char"030D} \SUPPORTEDANCH{\char"0068\char"030D} \SUPPORTEDANCH{\char"2C68\char"030D} \SUPPORTEDANCH{\char"006B\char"030D} \SUPPORTEDFALL{\char"0199\char"030D} \SUPPORTEDANCH{\char"2C6A\char"030D} \SUPPORTEDANCH{\char"006C\char"030D} \SUPPORTEDFALL{\char"0234\char"030D} \SUPPORTEDANCH{\char"006E\char"030D} \SUPPORTEDANCH{\char"0235\char"030D} \SUPPORTEDANCH{\char"0070\char"030D} \SUPPORTEDFALL{\char"01A5\char"030D} \SUPPORTEDANCH{\char"0071\char"030D} \SUPPORTEDANCH{\char"02A0\char"030D} \SUPPORTEDFALL{\char"0074\char"030D} \SUPPORTEDFALL{\char"01AB\char"030D} \SUPPORTEDFALL{\char"01AD\char"030D} \SUPPORTEDFALL{\char"0236\char"030D} \SUPPORTEDANCH{\char"0075\char"030D} \SUPPORTEDANCH{\char"AB52\char"030D} \SUPPORTEDANCH{\char"0079\char"030D} \SUPPORTEDANCH{\char"01B4\char"030D} \SUPPORTEDANCH{\char"1EFF\char"030D}

\SUPPORTEDANCH{\char"0048\char"030F} \SUPPORTEDFALL{\char"2C67\char"030F} \SUPPORTEDANCH{\char"004B\char"030F} \SUPPORTEDANCH{\char"2C69\char"030F} \SUPPORTEDANCH{\char"004D\char"030F} \SUPPORTEDANCH{\char"2C6E\char"030F} \SUPPORTEDANCH{\char"0050\char"030F} \SUPPORTEDANCH{\char"01A4\char"030F} \SUPPORTEDPREC{\char"0052\char"030F} \SUPPORTEDFALL{\char"2C64\char"030F} \SUPPORTEDANCH{\char"0054\char"030F} \SUPPORTEDFALL{\char"01AC\char"030F} \SUPPORTEDFALL{\char"01AE\char"030F} \SUPPORTEDANCH{\char"0059\char"030F} \SUPPORTEDFALL{\char"01B3\char"030F} \SUPPORTEDPREC{\char"0061\char"030F} \SUPPORTEDFALL{\char"1E9A\char"030F} \SUPPORTEDFALL{\char"0062\char"030F} \SUPPORTEDFALL{\char"0183\char"030F} \SUPPORTEDANCH{\char"0063\char"030F} \SUPPORTEDFALL{\char"0188\char"030F} \SUPPORTEDANCH{\char"0064\char"030F} \SUPPORTEDANCH{\char"018C\char"030F} \SUPPORTEDFALL{\char"0221\char"030F} \SUPPORTEDANCH{\char"0068\char"030F} \SUPPORTEDANCH{\char"2C68\char"030F} \SUPPORTEDANCH{\char"006B\char"030F} \SUPPORTEDFALL{\char"0199\char"030F} \SUPPORTEDANCH{\char"2C6A\char"030F} \SUPPORTEDANCH{\char"006C\char"030F} \SUPPORTEDFALL{\char"0234\char"030F} \SUPPORTEDANCH{\char"006E\char"030F} \SUPPORTEDANCH{\char"0235\char"030F} \SUPPORTEDANCH{\char"0070\char"030F} \SUPPORTEDFALL{\char"01A5\char"030F} \SUPPORTEDANCH{\char"0071\char"030F} \SUPPORTEDANCH{\char"02A0\char"030F} \SUPPORTEDFALL{\char"0074\char"030F} \SUPPORTEDFALL{\char"01AB\char"030F} \SUPPORTEDFALL{\char"01AD\char"030F} \SUPPORTEDFALL{\char"0236\char"030F} \SUPPORTEDPREC{\char"0075\char"030F} \SUPPORTEDANCH{\char"AB52\char"030F} \SUPPORTEDANCH{\char"0079\char"030F} \SUPPORTEDANCH{\char"01B4\char"030F} \SUPPORTEDANCH{\char"1EFF\char"030F}

\SUPPORTEDANCH{\char"0048\char"0310} \SUPPORTEDFALL{\char"2C67\char"0310} \SUPPORTEDANCH{\char"004B\char"0310} \SUPPORTEDANCH{\char"2C69\char"0310} \SUPPORTEDANCH{\char"004D\char"0310} \SUPPORTEDANCH{\char"2C6E\char"0310} \SUPPORTEDANCH{\char"0050\char"0310} \SUPPORTEDANCH{\char"01A4\char"0310} \SUPPORTEDANCH{\char"0052\char"0310} \SUPPORTEDFALL{\char"2C64\char"0310} \SUPPORTEDANCH{\char"0054\char"0310} \SUPPORTEDFALL{\char"01AC\char"0310} \SUPPORTEDFALL{\char"01AE\char"0310} \SUPPORTEDANCH{\char"0059\char"0310} \SUPPORTEDFALL{\char"01B3\char"0310} \SUPPORTEDANCH{\char"0061\char"0310} \SUPPORTEDFALL{\char"1E9A\char"0310} \SUPPORTEDFALL{\char"0062\char"0310} \SUPPORTEDFALL{\char"0183\char"0310} \SUPPORTEDANCH{\char"0063\char"0310} \SUPPORTEDFALL{\char"0188\char"0310} \SUPPORTEDANCH{\char"0064\char"0310} \SUPPORTEDANCH{\char"018C\char"0310} \SUPPORTEDFALL{\char"0221\char"0310} \SUPPORTEDANCH{\char"0068\char"0310} \SUPPORTEDANCH{\char"2C68\char"0310} \SUPPORTEDANCH{\char"006B\char"0310} \SUPPORTEDFALL{\char"0199\char"0310} \SUPPORTEDANCH{\char"2C6A\char"0310} \SUPPORTEDANCH{\char"006C\char"0310} \SUPPORTEDFALL{\char"0234\char"0310} \SUPPORTEDANCH{\char"006E\char"0310} \SUPPORTEDANCH{\char"0235\char"0310} \SUPPORTEDANCH{\char"0070\char"0310} \SUPPORTEDFALL{\char"01A5\char"0310} \SUPPORTEDANCH{\char"0071\char"0310} \SUPPORTEDANCH{\char"02A0\char"0310} \SUPPORTEDFALL{\char"0074\char"0310} \SUPPORTEDFALL{\char"01AB\char"0310} \SUPPORTEDFALL{\char"01AD\char"0310} \SUPPORTEDFALL{\char"0236\char"0310} \SUPPORTEDANCH{\char"0075\char"0310} \SUPPORTEDANCH{\char"AB52\char"0310} \SUPPORTEDANCH{\char"0079\char"0310} \SUPPORTEDANCH{\char"01B4\char"0310} \SUPPORTEDANCH{\char"1EFF\char"0310}

\SUPPORTEDANCH{\char"0048\char"0311} \SUPPORTEDFALL{\char"2C67\char"0311} \SUPPORTEDANCH{\char"004B\char"0311} \SUPPORTEDANCH{\char"2C69\char"0311} \SUPPORTEDANCH{\char"004D\char"0311} \SUPPORTEDANCH{\char"2C6E\char"0311} \SUPPORTEDANCH{\char"0050\char"0311} \SUPPORTEDANCH{\char"01A4\char"0311} \SUPPORTEDPREC{\char"0052\char"0311} \SUPPORTEDFALL{\char"2C64\char"0311} \SUPPORTEDANCH{\char"0054\char"0311} \SUPPORTEDFALL{\char"01AC\char"0311} \SUPPORTEDFALL{\char"01AE\char"0311} \SUPPORTEDANCH{\char"0059\char"0311} \SUPPORTEDFALL{\char"01B3\char"0311} \SUPPORTEDPREC{\char"0061\char"0311} \SUPPORTEDFALL{\char"1E9A\char"0311} \SUPPORTEDFALL{\char"0062\char"0311} \SUPPORTEDFALL{\char"0183\char"0311} \SUPPORTEDANCH{\char"0063\char"0311} \SUPPORTEDFALL{\char"0188\char"0311} \SUPPORTEDANCH{\char"0064\char"0311} \SUPPORTEDANCH{\char"018C\char"0311} \SUPPORTEDFALL{\char"0221\char"0311} \SUPPORTEDANCH{\char"0068\char"0311} \SUPPORTEDANCH{\char"2C68\char"0311} \SUPPORTEDANCH{\char"006B\char"0311} \SUPPORTEDFALL{\char"0199\char"0311} \SUPPORTEDANCH{\char"2C6A\char"0311} \SUPPORTEDANCH{\char"006C\char"0311} \SUPPORTEDFALL{\char"0234\char"0311} \SUPPORTEDANCH{\char"006E\char"0311} \SUPPORTEDANCH{\char"0235\char"0311} \SUPPORTEDANCH{\char"0070\char"0311} \SUPPORTEDFALL{\char"01A5\char"0311} \SUPPORTEDANCH{\char"0071\char"0311} \SUPPORTEDANCH{\char"02A0\char"0311} \SUPPORTEDFALL{\char"0074\char"0311} \SUPPORTEDFALL{\char"01AB\char"0311} \SUPPORTEDFALL{\char"01AD\char"0311} \SUPPORTEDFALL{\char"0236\char"0311} \SUPPORTEDPREC{\char"0075\char"0311} \SUPPORTEDANCH{\char"AB52\char"0311} \SUPPORTEDANCH{\char"0079\char"0311} \SUPPORTEDANCH{\char"01B4\char"0311} \SUPPORTEDANCH{\char"1EFF\char"0311}

\SUPPORTEDANCH{\char"0048\char"0312} \SUPPORTEDFALL{\char"2C67\char"0312} \SUPPORTEDANCH{\char"004B\char"0312} \SUPPORTEDANCH{\char"2C69\char"0312} \SUPPORTEDANCH{\char"004D\char"0312} \SUPPORTEDANCH{\char"2C6E\char"0312} \SUPPORTEDANCH{\char"0050\char"0312} \SUPPORTEDANCH{\char"01A4\char"0312} \SUPPORTEDANCH{\char"0052\char"0312} \SUPPORTEDFALL{\char"2C64\char"0312} \SUPPORTEDANCH{\char"0054\char"0312} \SUPPORTEDFALL{\char"01AC\char"0312} \SUPPORTEDFALL{\char"01AE\char"0312} \SUPPORTEDANCH{\char"0059\char"0312} \SUPPORTEDFALL{\char"01B3\char"0312} \SUPPORTEDANCH{\char"0061\char"0312} \SUPPORTEDFALL{\char"1E9A\char"0312} \SUPPORTEDFALL{\char"0062\char"0312} \SUPPORTEDFALL{\char"0183\char"0312} \SUPPORTEDANCH{\char"0063\char"0312} \SUPPORTEDFALL{\char"0188\char"0312} \SUPPORTEDANCH{\char"0064\char"0312} \SUPPORTEDANCH{\char"018C\char"0312} \SUPPORTEDFALL{\char"0221\char"0312} \SUPPORTEDANCH{\char"0068\char"0312} \SUPPORTEDANCH{\char"2C68\char"0312} \SUPPORTEDANCH{\char"006B\char"0312} \SUPPORTEDFALL{\char"0199\char"0312} \SUPPORTEDANCH{\char"2C6A\char"0312} \SUPPORTEDANCH{\char"006C\char"0312} \SUPPORTEDFALL{\char"0234\char"0312} \SUPPORTEDANCH{\char"006E\char"0312} \SUPPORTEDANCH{\char"0235\char"0312} \SUPPORTEDANCH{\char"0070\char"0312} \SUPPORTEDFALL{\char"01A5\char"0312} \SUPPORTEDANCH{\char"0071\char"0312} \SUPPORTEDANCH{\char"02A0\char"0312} \SUPPORTEDFALL{\char"0074\char"0312} \SUPPORTEDFALL{\char"01AB\char"0312} \SUPPORTEDFALL{\char"01AD\char"0312} \SUPPORTEDFALL{\char"0236\char"0312} \SUPPORTEDANCH{\char"0075\char"0312} \SUPPORTEDANCH{\char"AB52\char"0312} \SUPPORTEDANCH{\char"0079\char"0312} \SUPPORTEDANCH{\char"01B4\char"0312} \SUPPORTEDANCH{\char"1EFF\char"0312}

\SUPPORTEDANCH{\char"0048\char"033D} \SUPPORTEDFALL{\char"2C67\char"033D} \SUPPORTEDANCH{\char"004B\char"033D} \SUPPORTEDANCH{\char"2C69\char"033D} \SUPPORTEDANCH{\char"004D\char"033D} \SUPPORTEDANCH{\char"2C6E\char"033D} \SUPPORTEDANCH{\char"0050\char"033D} \SUPPORTEDANCH{\char"01A4\char"033D} \SUPPORTEDANCH{\char"0052\char"033D} \SUPPORTEDFALL{\char"2C64\char"033D} \SUPPORTEDANCH{\char"0054\char"033D} \SUPPORTEDFALL{\char"01AC\char"033D} \SUPPORTEDFALL{\char"01AE\char"033D} \SUPPORTEDANCH{\char"0059\char"033D} \SUPPORTEDFALL{\char"01B3\char"033D} \SUPPORTEDANCH{\char"0061\char"033D} \SUPPORTEDFALL{\char"1E9A\char"033D} \SUPPORTEDFALL{\char"0062\char"033D} \SUPPORTEDFALL{\char"0183\char"033D} \SUPPORTEDANCH{\char"0063\char"033D} \SUPPORTEDFALL{\char"0188\char"033D} \SUPPORTEDANCH{\char"0064\char"033D} \SUPPORTEDANCH{\char"018C\char"033D} \SUPPORTEDFALL{\char"0221\char"033D} \SUPPORTEDANCH{\char"0068\char"033D} \SUPPORTEDANCH{\char"2C68\char"033D} \SUPPORTEDANCH{\char"006B\char"033D} \SUPPORTEDFALL{\char"0199\char"033D} \SUPPORTEDANCH{\char"2C6A\char"033D} \SUPPORTEDANCH{\char"006C\char"033D} \SUPPORTEDFALL{\char"0234\char"033D} \SUPPORTEDANCH{\char"006E\char"033D} \SUPPORTEDANCH{\char"0235\char"033D} \SUPPORTEDANCH{\char"0070\char"033D} \SUPPORTEDFALL{\char"01A5\char"033D} \SUPPORTEDANCH{\char"0071\char"033D} \SUPPORTEDANCH{\char"02A0\char"033D} \SUPPORTEDFALL{\char"0074\char"033D} \SUPPORTEDFALL{\char"01AB\char"033D} \SUPPORTEDFALL{\char"01AD\char"033D} \SUPPORTEDFALL{\char"0236\char"033D} \SUPPORTEDANCH{\char"0075\char"033D} \SUPPORTEDANCH{\char"AB52\char"033D} \SUPPORTEDANCH{\char"0079\char"033D} \SUPPORTEDANCH{\char"01B4\char"033D} \SUPPORTEDANCH{\char"1EFF\char"033D}

\SUPPORTEDANCH{\char"0048\char"033E} \SUPPORTEDFALL{\char"2C67\char"033E} \SUPPORTEDANCH{\char"004B\char"033E} \SUPPORTEDANCH{\char"2C69\char"033E} \SUPPORTEDANCH{\char"004D\char"033E} \SUPPORTEDANCH{\char"2C6E\char"033E} \SUPPORTEDANCH{\char"0050\char"033E} \SUPPORTEDANCH{\char"01A4\char"033E} \SUPPORTEDANCH{\char"0052\char"033E} \SUPPORTEDFALL{\char"2C64\char"033E} \SUPPORTEDANCH{\char"0054\char"033E} \SUPPORTEDFALL{\char"01AC\char"033E} \SUPPORTEDFALL{\char"01AE\char"033E} \SUPPORTEDANCH{\char"0059\char"033E} \SUPPORTEDFALL{\char"01B3\char"033E} \SUPPORTEDANCH{\char"0061\char"033E} \SUPPORTEDFALL{\char"1E9A\char"033E} \SUPPORTEDFALL{\char"0062\char"033E} \SUPPORTEDFALL{\char"0183\char"033E} \SUPPORTEDANCH{\char"0063\char"033E} \SUPPORTEDFALL{\char"0188\char"033E} \SUPPORTEDANCH{\char"0064\char"033E} \SUPPORTEDANCH{\char"018C\char"033E} \SUPPORTEDFALL{\char"0221\char"033E} \SUPPORTEDANCH{\char"0068\char"033E} \SUPPORTEDANCH{\char"2C68\char"033E} \SUPPORTEDANCH{\char"006B\char"033E} \SUPPORTEDFALL{\char"0199\char"033E} \SUPPORTEDANCH{\char"2C6A\char"033E} \SUPPORTEDANCH{\char"006C\char"033E} \SUPPORTEDFALL{\char"0234\char"033E} \SUPPORTEDANCH{\char"006E\char"033E} \SUPPORTEDANCH{\char"0235\char"033E} \SUPPORTEDANCH{\char"0070\char"033E} \SUPPORTEDFALL{\char"01A5\char"033E} \SUPPORTEDANCH{\char"0071\char"033E} \SUPPORTEDANCH{\char"02A0\char"033E} \SUPPORTEDFALL{\char"0074\char"033E} \SUPPORTEDFALL{\char"01AB\char"033E} \SUPPORTEDFALL{\char"01AD\char"033E} \SUPPORTEDFALL{\char"0236\char"033E} \SUPPORTEDANCH{\char"0075\char"033E} \SUPPORTEDANCH{\char"AB52\char"033E} \SUPPORTEDANCH{\char"0079\char"033E} \SUPPORTEDANCH{\char"01B4\char"033E} \SUPPORTEDANCH{\char"1EFF\char"033E}

\SUPPORTEDANCH{\char"0048\char"033F} \SUPPORTEDFALL{\char"2C67\char"033F} \SUPPORTEDANCH{\char"004B\char"033F} \SUPPORTEDANCH{\char"2C69\char"033F} \SUPPORTEDANCH{\char"004D\char"033F} \SUPPORTEDANCH{\char"2C6E\char"033F} \SUPPORTEDANCH{\char"0050\char"033F} \SUPPORTEDANCH{\char"01A4\char"033F} \SUPPORTEDANCH{\char"0052\char"033F} \SUPPORTEDFALL{\char"2C64\char"033F} \SUPPORTEDANCH{\char"0054\char"033F} \SUPPORTEDFALL{\char"01AC\char"033F} \SUPPORTEDFALL{\char"01AE\char"033F} \SUPPORTEDANCH{\char"0059\char"033F} \SUPPORTEDFALL{\char"01B3\char"033F} \SUPPORTEDANCH{\char"0061\char"033F} \SUPPORTEDFALL{\char"1E9A\char"033F} \SUPPORTEDFALL{\char"0062\char"033F} \SUPPORTEDFALL{\char"0183\char"033F} \SUPPORTEDANCH{\char"0063\char"033F} \SUPPORTEDFALL{\char"0188\char"033F} \SUPPORTEDANCH{\char"0064\char"033F} \SUPPORTEDANCH{\char"018C\char"033F} \SUPPORTEDFALL{\char"0221\char"033F} \SUPPORTEDANCH{\char"0068\char"033F} \SUPPORTEDANCH{\char"2C68\char"033F} \SUPPORTEDANCH{\char"006B\char"033F} \SUPPORTEDFALL{\char"0199\char"033F} \SUPPORTEDANCH{\char"2C6A\char"033F} \SUPPORTEDANCH{\char"006C\char"033F} \SUPPORTEDFALL{\char"0234\char"033F} \SUPPORTEDANCH{\char"006E\char"033F} \SUPPORTEDANCH{\char"0235\char"033F} \SUPPORTEDANCH{\char"0070\char"033F} \SUPPORTEDFALL{\char"01A5\char"033F} \SUPPORTEDANCH{\char"0071\char"033F} \SUPPORTEDANCH{\char"02A0\char"033F} \SUPPORTEDFALL{\char"0074\char"033F} \SUPPORTEDFALL{\char"01AB\char"033F} \SUPPORTEDFALL{\char"01AD\char"033F} \SUPPORTEDFALL{\char"0236\char"033F} \SUPPORTEDANCH{\char"0075\char"033F} \SUPPORTEDANCH{\char"AB52\char"033F} \SUPPORTEDANCH{\char"0079\char"033F} \SUPPORTEDANCH{\char"01B4\char"033F} \SUPPORTEDANCH{\char"1EFF\char"033F}

\SUPPORTEDANCH{\char"0048\char"030E} \SUPPORTEDFALL{\char"2C67\char"030E} \SUPPORTEDANCH{\char"004B\char"030E} \SUPPORTEDANCH{\char"2C69\char"030E} \SUPPORTEDANCH{\char"004D\char"030E} \SUPPORTEDANCH{\char"2C6E\char"030E} \SUPPORTEDANCH{\char"0050\char"030E} \SUPPORTEDANCH{\char"01A4\char"030E} \SUPPORTEDANCH{\char"0052\char"030E} \SUPPORTEDFALL{\char"2C64\char"030E} \SUPPORTEDANCH{\char"0054\char"030E} \SUPPORTEDFALL{\char"01AC\char"030E} \SUPPORTEDFALL{\char"01AE\char"030E} \SUPPORTEDANCH{\char"0059\char"030E} \SUPPORTEDFALL{\char"01B3\char"030E} \SUPPORTEDANCH{\char"0061\char"030E} \SUPPORTEDFALL{\char"1E9A\char"030E} \SUPPORTEDFALL{\char"0062\char"030E} \SUPPORTEDFALL{\char"0183\char"030E} \SUPPORTEDANCH{\char"0063\char"030E} \SUPPORTEDFALL{\char"0188\char"030E} \SUPPORTEDANCH{\char"0064\char"030E} \SUPPORTEDANCH{\char"018C\char"030E} \SUPPORTEDFALL{\char"0221\char"030E} \SUPPORTEDANCH{\char"0068\char"030E} \SUPPORTEDANCH{\char"2C68\char"030E} \SUPPORTEDANCH{\char"006B\char"030E} \SUPPORTEDFALL{\char"0199\char"030E} \SUPPORTEDANCH{\char"2C6A\char"030E} \SUPPORTEDANCH{\char"006C\char"030E} \SUPPORTEDFALL{\char"0234\char"030E} \SUPPORTEDANCH{\char"006E\char"030E} \SUPPORTEDANCH{\char"0235\char"030E} \SUPPORTEDANCH{\char"0070\char"030E} \SUPPORTEDFALL{\char"01A5\char"030E} \SUPPORTEDANCH{\char"0071\char"030E} \SUPPORTEDANCH{\char"02A0\char"030E} \SUPPORTEDFALL{\char"0074\char"030E} \SUPPORTEDFALL{\char"01AB\char"030E} \SUPPORTEDFALL{\char"01AD\char"030E} \SUPPORTEDFALL{\char"0236\char"030E} \SUPPORTEDANCH{\char"0075\char"030E} \SUPPORTEDANCH{\char"AB52\char"030E} \SUPPORTEDANCH{\char"0079\char"030E} \SUPPORTEDANCH{\char"01B4\char"030E} \SUPPORTEDANCH{\char"1EFF\char"030E}

\SUPPORTEDANCH{\char"0048\char"0346} \SUPPORTEDFALL{\char"2C67\char"0346} \SUPPORTEDANCH{\char"004B\char"0346} \SUPPORTEDANCH{\char"2C69\char"0346} \SUPPORTEDANCH{\char"004D\char"0346} \SUPPORTEDANCH{\char"2C6E\char"0346} \SUPPORTEDANCH{\char"0050\char"0346} \SUPPORTEDANCH{\char"01A4\char"0346} \SUPPORTEDANCH{\char"0052\char"0346} \SUPPORTEDFALL{\char"2C64\char"0346} \SUPPORTEDANCH{\char"0054\char"0346} \SUPPORTEDFALL{\char"01AC\char"0346} \SUPPORTEDFALL{\char"01AE\char"0346} \SUPPORTEDANCH{\char"0059\char"0346} \SUPPORTEDFALL{\char"01B3\char"0346} \SUPPORTEDANCH{\char"0061\char"0346} \SUPPORTEDFALL{\char"1E9A\char"0346} \SUPPORTEDFALL{\char"0062\char"0346} \SUPPORTEDFALL{\char"0183\char"0346} \SUPPORTEDANCH{\char"0063\char"0346} \SUPPORTEDFALL{\char"0188\char"0346} \SUPPORTEDANCH{\char"0064\char"0346} \SUPPORTEDANCH{\char"018C\char"0346} \SUPPORTEDFALL{\char"0221\char"0346} \SUPPORTEDANCH{\char"0068\char"0346} \SUPPORTEDANCH{\char"2C68\char"0346} \SUPPORTEDANCH{\char"006B\char"0346} \SUPPORTEDFALL{\char"0199\char"0346} \SUPPORTEDANCH{\char"2C6A\char"0346} \SUPPORTEDANCH{\char"006C\char"0346} \SUPPORTEDFALL{\char"0234\char"0346} \SUPPORTEDANCH{\char"006E\char"0346} \SUPPORTEDANCH{\char"0235\char"0346} \SUPPORTEDANCH{\char"0070\char"0346} \SUPPORTEDFALL{\char"01A5\char"0346} \SUPPORTEDANCH{\char"0071\char"0346} \SUPPORTEDANCH{\char"02A0\char"0346} \SUPPORTEDFALL{\char"0074\char"0346} \SUPPORTEDFALL{\char"01AB\char"0346} \SUPPORTEDFALL{\char"01AD\char"0346} \SUPPORTEDFALL{\char"0236\char"0346} \SUPPORTEDANCH{\char"0075\char"0346} \SUPPORTEDANCH{\char"AB52\char"0346} \SUPPORTEDANCH{\char"0079\char"0346} \SUPPORTEDANCH{\char"01B4\char"0346} \SUPPORTEDANCH{\char"1EFF\char"0346}

\SUPPORTEDANCH{\char"0048\char"034A} \SUPPORTEDFALL{\char"2C67\char"034A} \SUPPORTEDANCH{\char"004B\char"034A} \SUPPORTEDANCH{\char"2C69\char"034A} \SUPPORTEDANCH{\char"004D\char"034A} \SUPPORTEDANCH{\char"2C6E\char"034A} \SUPPORTEDANCH{\char"0050\char"034A} \SUPPORTEDANCH{\char"01A4\char"034A} \SUPPORTEDANCH{\char"0052\char"034A} \SUPPORTEDFALL{\char"2C64\char"034A} \SUPPORTEDANCH{\char"0054\char"034A} \SUPPORTEDFALL{\char"01AC\char"034A} \SUPPORTEDFALL{\char"01AE\char"034A} \SUPPORTEDANCH{\char"0059\char"034A} \SUPPORTEDFALL{\char"01B3\char"034A} \SUPPORTEDANCH{\char"0061\char"034A} \SUPPORTEDFALL{\char"1E9A\char"034A} \SUPPORTEDFALL{\char"0062\char"034A} \SUPPORTEDFALL{\char"0183\char"034A} \SUPPORTEDANCH{\char"0063\char"034A} \SUPPORTEDFALL{\char"0188\char"034A} \SUPPORTEDANCH{\char"0064\char"034A} \SUPPORTEDANCH{\char"018C\char"034A} \SUPPORTEDFALL{\char"0221\char"034A} \SUPPORTEDANCH{\char"0068\char"034A} \SUPPORTEDANCH{\char"2C68\char"034A} \SUPPORTEDANCH{\char"006B\char"034A} \SUPPORTEDFALL{\char"0199\char"034A} \SUPPORTEDANCH{\char"2C6A\char"034A} \SUPPORTEDANCH{\char"006C\char"034A} \SUPPORTEDFALL{\char"0234\char"034A} \SUPPORTEDANCH{\char"006E\char"034A} \SUPPORTEDANCH{\char"0235\char"034A} \SUPPORTEDANCH{\char"0070\char"034A} \SUPPORTEDFALL{\char"01A5\char"034A} \SUPPORTEDANCH{\char"0071\char"034A} \SUPPORTEDANCH{\char"02A0\char"034A} \SUPPORTEDFALL{\char"0074\char"034A} \SUPPORTEDFALL{\char"01AB\char"034A} \SUPPORTEDFALL{\char"01AD\char"034A} \SUPPORTEDFALL{\char"0236\char"034A} \SUPPORTEDANCH{\char"0075\char"034A} \SUPPORTEDANCH{\char"AB52\char"034A} \SUPPORTEDANCH{\char"0079\char"034A} \SUPPORTEDANCH{\char"01B4\char"034A} \SUPPORTEDANCH{\char"1EFF\char"034A}

\SUPPORTEDANCH{\char"0048\char"034B} \SUPPORTEDFALL{\char"2C67\char"034B} \SUPPORTEDANCH{\char"004B\char"034B} \SUPPORTEDANCH{\char"2C69\char"034B} \SUPPORTEDANCH{\char"004D\char"034B} \SUPPORTEDANCH{\char"2C6E\char"034B} \SUPPORTEDANCH{\char"0050\char"034B} \SUPPORTEDANCH{\char"01A4\char"034B} \SUPPORTEDANCH{\char"0052\char"034B} \SUPPORTEDFALL{\char"2C64\char"034B} \SUPPORTEDANCH{\char"0054\char"034B} \SUPPORTEDFALL{\char"01AC\char"034B} \SUPPORTEDFALL{\char"01AE\char"034B} \SUPPORTEDANCH{\char"0059\char"034B} \SUPPORTEDFALL{\char"01B3\char"034B} \SUPPORTEDANCH{\char"0061\char"034B} \SUPPORTEDFALL{\char"1E9A\char"034B} \SUPPORTEDFALL{\char"0062\char"034B} \SUPPORTEDFALL{\char"0183\char"034B} \SUPPORTEDANCH{\char"0063\char"034B} \SUPPORTEDFALL{\char"0188\char"034B} \SUPPORTEDANCH{\char"0064\char"034B} \SUPPORTEDANCH{\char"018C\char"034B} \SUPPORTEDFALL{\char"0221\char"034B} \SUPPORTEDANCH{\char"0068\char"034B} \SUPPORTEDANCH{\char"2C68\char"034B} \SUPPORTEDANCH{\char"006B\char"034B} \SUPPORTEDFALL{\char"0199\char"034B} \SUPPORTEDANCH{\char"2C6A\char"034B} \SUPPORTEDANCH{\char"006C\char"034B} \SUPPORTEDFALL{\char"0234\char"034B} \SUPPORTEDANCH{\char"006E\char"034B} \SUPPORTEDANCH{\char"0235\char"034B} \SUPPORTEDANCH{\char"0070\char"034B} \SUPPORTEDFALL{\char"01A5\char"034B} \SUPPORTEDANCH{\char"0071\char"034B} \SUPPORTEDANCH{\char"02A0\char"034B} \SUPPORTEDFALL{\char"0074\char"034B} \SUPPORTEDFALL{\char"01AB\char"034B} \SUPPORTEDFALL{\char"01AD\char"034B} \SUPPORTEDFALL{\char"0236\char"034B} \SUPPORTEDANCH{\char"0075\char"034B} \SUPPORTEDANCH{\char"AB52\char"034B} \SUPPORTEDANCH{\char"0079\char"034B} \SUPPORTEDANCH{\char"01B4\char"034B} \SUPPORTEDANCH{\char"1EFF\char"034B}

\SUPPORTEDANCH{\char"0048\char"034C} \SUPPORTEDFALL{\char"2C67\char"034C} \SUPPORTEDANCH{\char"004B\char"034C} \SUPPORTEDANCH{\char"2C69\char"034C} \SUPPORTEDANCH{\char"004D\char"034C} \SUPPORTEDANCH{\char"2C6E\char"034C} \SUPPORTEDANCH{\char"0050\char"034C} \SUPPORTEDANCH{\char"01A4\char"034C} \SUPPORTEDANCH{\char"0052\char"034C} \SUPPORTEDFALL{\char"2C64\char"034C} \SUPPORTEDANCH{\char"0054\char"034C} \SUPPORTEDFALL{\char"01AC\char"034C} \SUPPORTEDFALL{\char"01AE\char"034C} \SUPPORTEDANCH{\char"0059\char"034C} \SUPPORTEDFALL{\char"01B3\char"034C} \SUPPORTEDANCH{\char"0061\char"034C} \SUPPORTEDFALL{\char"1E9A\char"034C} \SUPPORTEDFALL{\char"0062\char"034C} \SUPPORTEDFALL{\char"0183\char"034C} \SUPPORTEDANCH{\char"0063\char"034C} \SUPPORTEDFALL{\char"0188\char"034C} \SUPPORTEDANCH{\char"0064\char"034C} \SUPPORTEDANCH{\char"018C\char"034C} \SUPPORTEDFALL{\char"0221\char"034C} \SUPPORTEDANCH{\char"0068\char"034C} \SUPPORTEDANCH{\char"2C68\char"034C} \SUPPORTEDANCH{\char"006B\char"034C} \SUPPORTEDFALL{\char"0199\char"034C} \SUPPORTEDANCH{\char"2C6A\char"034C} \SUPPORTEDANCH{\char"006C\char"034C} \SUPPORTEDFALL{\char"0234\char"034C} \SUPPORTEDANCH{\char"006E\char"034C} \SUPPORTEDANCH{\char"0235\char"034C} \SUPPORTEDANCH{\char"0070\char"034C} \SUPPORTEDFALL{\char"01A5\char"034C} \SUPPORTEDANCH{\char"0071\char"034C} \SUPPORTEDANCH{\char"02A0\char"034C} \SUPPORTEDFALL{\char"0074\char"034C} \SUPPORTEDFALL{\char"01AB\char"034C} \SUPPORTEDFALL{\char"01AD\char"034C} \SUPPORTEDFALL{\char"0236\char"034C} \SUPPORTEDANCH{\char"0075\char"034C} \SUPPORTEDANCH{\char"AB52\char"034C} \SUPPORTEDANCH{\char"0079\char"034C} \SUPPORTEDANCH{\char"01B4\char"034C} \SUPPORTEDANCH{\char"1EFF\char"034C}

\SUPPORTEDANCH{\char"0048\char"0350} \SUPPORTEDFALL{\char"2C67\char"0350} \SUPPORTEDANCH{\char"004B\char"0350} \SUPPORTEDANCH{\char"2C69\char"0350} \SUPPORTEDANCH{\char"004D\char"0350} \SUPPORTEDANCH{\char"2C6E\char"0350} \SUPPORTEDANCH{\char"0050\char"0350} \SUPPORTEDANCH{\char"01A4\char"0350} \SUPPORTEDANCH{\char"0052\char"0350} \SUPPORTEDFALL{\char"2C64\char"0350} \SUPPORTEDANCH{\char"0054\char"0350} \SUPPORTEDFALL{\char"01AC\char"0350} \SUPPORTEDFALL{\char"01AE\char"0350} \SUPPORTEDANCH{\char"0059\char"0350} \SUPPORTEDFALL{\char"01B3\char"0350} \SUPPORTEDANCH{\char"0061\char"0350} \SUPPORTEDFALL{\char"1E9A\char"0350} \SUPPORTEDFALL{\char"0062\char"0350} \SUPPORTEDFALL{\char"0183\char"0350} \SUPPORTEDANCH{\char"0063\char"0350} \SUPPORTEDFALL{\char"0188\char"0350} \SUPPORTEDANCH{\char"0064\char"0350} \SUPPORTEDANCH{\char"018C\char"0350} \SUPPORTEDFALL{\char"0221\char"0350} \SUPPORTEDANCH{\char"0068\char"0350} \SUPPORTEDANCH{\char"2C68\char"0350} \SUPPORTEDANCH{\char"006B\char"0350} \SUPPORTEDFALL{\char"0199\char"0350} \SUPPORTEDANCH{\char"2C6A\char"0350} \SUPPORTEDANCH{\char"006C\char"0350} \SUPPORTEDFALL{\char"0234\char"0350} \SUPPORTEDANCH{\char"006E\char"0350} \SUPPORTEDANCH{\char"0235\char"0350} \SUPPORTEDANCH{\char"0070\char"0350} \SUPPORTEDFALL{\char"01A5\char"0350} \SUPPORTEDANCH{\char"0071\char"0350} \SUPPORTEDANCH{\char"02A0\char"0350} \SUPPORTEDFALL{\char"0074\char"0350} \SUPPORTEDFALL{\char"01AB\char"0350} \SUPPORTEDFALL{\char"01AD\char"0350} \SUPPORTEDFALL{\char"0236\char"0350} \SUPPORTEDANCH{\char"0075\char"0350} \SUPPORTEDANCH{\char"AB52\char"0350} \SUPPORTEDANCH{\char"0079\char"0350} \SUPPORTEDANCH{\char"01B4\char"0350} \SUPPORTEDANCH{\char"1EFF\char"0350}

\SUPPORTEDANCH{\char"0048\char"0351} \SUPPORTEDFALL{\char"2C67\char"0351} \SUPPORTEDANCH{\char"004B\char"0351} \SUPPORTEDANCH{\char"2C69\char"0351} \SUPPORTEDANCH{\char"004D\char"0351} \SUPPORTEDANCH{\char"2C6E\char"0351} \SUPPORTEDANCH{\char"0050\char"0351} \SUPPORTEDANCH{\char"01A4\char"0351} \SUPPORTEDANCH{\char"0052\char"0351} \SUPPORTEDFALL{\char"2C64\char"0351} \SUPPORTEDANCH{\char"0054\char"0351} \SUPPORTEDFALL{\char"01AC\char"0351} \SUPPORTEDFALL{\char"01AE\char"0351} \SUPPORTEDANCH{\char"0059\char"0351} \SUPPORTEDFALL{\char"01B3\char"0351} \SUPPORTEDANCH{\char"0061\char"0351} \SUPPORTEDFALL{\char"1E9A\char"0351} \SUPPORTEDFALL{\char"0062\char"0351} \SUPPORTEDFALL{\char"0183\char"0351} \SUPPORTEDANCH{\char"0063\char"0351} \SUPPORTEDFALL{\char"0188\char"0351} \SUPPORTEDANCH{\char"0064\char"0351} \SUPPORTEDANCH{\char"018C\char"0351} \SUPPORTEDFALL{\char"0221\char"0351} \SUPPORTEDANCH{\char"0068\char"0351} \SUPPORTEDANCH{\char"2C68\char"0351} \SUPPORTEDANCH{\char"006B\char"0351} \SUPPORTEDFALL{\char"0199\char"0351} \SUPPORTEDANCH{\char"2C6A\char"0351} \SUPPORTEDANCH{\char"006C\char"0351} \SUPPORTEDFALL{\char"0234\char"0351} \SUPPORTEDANCH{\char"006E\char"0351} \SUPPORTEDANCH{\char"0235\char"0351} \SUPPORTEDANCH{\char"0070\char"0351} \SUPPORTEDFALL{\char"01A5\char"0351} \SUPPORTEDANCH{\char"0071\char"0351} \SUPPORTEDANCH{\char"02A0\char"0351} \SUPPORTEDFALL{\char"0074\char"0351} \SUPPORTEDFALL{\char"01AB\char"0351} \SUPPORTEDFALL{\char"01AD\char"0351} \SUPPORTEDFALL{\char"0236\char"0351} \SUPPORTEDANCH{\char"0075\char"0351} \SUPPORTEDANCH{\char"AB52\char"0351} \SUPPORTEDANCH{\char"0079\char"0351} \SUPPORTEDANCH{\char"01B4\char"0351} \SUPPORTEDANCH{\char"1EFF\char"0351}

\SUPPORTEDANCH{\char"0048\char"0352} \SUPPORTEDFALL{\char"2C67\char"0352} \SUPPORTEDANCH{\char"004B\char"0352} \SUPPORTEDANCH{\char"2C69\char"0352} \SUPPORTEDANCH{\char"004D\char"0352} \SUPPORTEDANCH{\char"2C6E\char"0352} \SUPPORTEDANCH{\char"0050\char"0352} \SUPPORTEDANCH{\char"01A4\char"0352} \SUPPORTEDANCH{\char"0052\char"0352} \SUPPORTEDFALL{\char"2C64\char"0352} \SUPPORTEDANCH{\char"0054\char"0352} \SUPPORTEDFALL{\char"01AC\char"0352} \SUPPORTEDFALL{\char"01AE\char"0352} \SUPPORTEDANCH{\char"0059\char"0352} \SUPPORTEDFALL{\char"01B3\char"0352} \SUPPORTEDANCH{\char"0061\char"0352} \SUPPORTEDFALL{\char"1E9A\char"0352} \SUPPORTEDFALL{\char"0062\char"0352} \SUPPORTEDFALL{\char"0183\char"0352} \SUPPORTEDANCH{\char"0063\char"0352} \SUPPORTEDFALL{\char"0188\char"0352} \SUPPORTEDANCH{\char"0064\char"0352} \SUPPORTEDANCH{\char"018C\char"0352} \SUPPORTEDFALL{\char"0221\char"0352} \SUPPORTEDANCH{\char"0068\char"0352} \SUPPORTEDANCH{\char"2C68\char"0352} \SUPPORTEDANCH{\char"006B\char"0352} \SUPPORTEDFALL{\char"0199\char"0352} \SUPPORTEDANCH{\char"2C6A\char"0352} \SUPPORTEDANCH{\char"006C\char"0352} \SUPPORTEDFALL{\char"0234\char"0352} \SUPPORTEDANCH{\char"006E\char"0352} \SUPPORTEDANCH{\char"0235\char"0352} \SUPPORTEDANCH{\char"0070\char"0352} \SUPPORTEDFALL{\char"01A5\char"0352} \SUPPORTEDANCH{\char"0071\char"0352} \SUPPORTEDANCH{\char"02A0\char"0352} \SUPPORTEDFALL{\char"0074\char"0352} \SUPPORTEDFALL{\char"01AB\char"0352} \SUPPORTEDFALL{\char"01AD\char"0352} \SUPPORTEDFALL{\char"0236\char"0352} \SUPPORTEDANCH{\char"0075\char"0352} \SUPPORTEDANCH{\char"AB52\char"0352} \SUPPORTEDANCH{\char"0079\char"0352} \SUPPORTEDANCH{\char"01B4\char"0352} \SUPPORTEDANCH{\char"1EFF\char"0352}

\SUPPORTEDANCH{\char"0048\char"0357} \SUPPORTEDFALL{\char"2C67\char"0357} \SUPPORTEDANCH{\char"004B\char"0357} \SUPPORTEDANCH{\char"2C69\char"0357} \SUPPORTEDANCH{\char"004D\char"0357} \SUPPORTEDANCH{\char"2C6E\char"0357} \SUPPORTEDANCH{\char"0050\char"0357} \SUPPORTEDANCH{\char"01A4\char"0357} \SUPPORTEDANCH{\char"0052\char"0357} \SUPPORTEDFALL{\char"2C64\char"0357} \SUPPORTEDANCH{\char"0054\char"0357} \SUPPORTEDFALL{\char"01AC\char"0357} \SUPPORTEDFALL{\char"01AE\char"0357} \SUPPORTEDANCH{\char"0059\char"0357} \SUPPORTEDFALL{\char"01B3\char"0357} \SUPPORTEDANCH{\char"0061\char"0357} \SUPPORTEDFALL{\char"1E9A\char"0357} \SUPPORTEDFALL{\char"0062\char"0357} \SUPPORTEDFALL{\char"0183\char"0357} \SUPPORTEDANCH{\char"0063\char"0357} \SUPPORTEDFALL{\char"0188\char"0357} \SUPPORTEDANCH{\char"0064\char"0357} \SUPPORTEDANCH{\char"018C\char"0357} \SUPPORTEDFALL{\char"0221\char"0357} \SUPPORTEDANCH{\char"0068\char"0357} \SUPPORTEDANCH{\char"2C68\char"0357} \SUPPORTEDANCH{\char"006B\char"0357} \SUPPORTEDFALL{\char"0199\char"0357} \SUPPORTEDANCH{\char"2C6A\char"0357} \SUPPORTEDANCH{\char"006C\char"0357} \SUPPORTEDFALL{\char"0234\char"0357} \SUPPORTEDANCH{\char"006E\char"0357} \SUPPORTEDANCH{\char"0235\char"0357} \SUPPORTEDANCH{\char"0070\char"0357} \SUPPORTEDFALL{\char"01A5\char"0357} \SUPPORTEDANCH{\char"0071\char"0357} \SUPPORTEDANCH{\char"02A0\char"0357} \SUPPORTEDFALL{\char"0074\char"0357} \SUPPORTEDFALL{\char"01AB\char"0357} \SUPPORTEDFALL{\char"01AD\char"0357} \SUPPORTEDFALL{\char"0236\char"0357} \SUPPORTEDANCH{\char"0075\char"0357} \SUPPORTEDANCH{\char"AB52\char"0357} \SUPPORTEDANCH{\char"0079\char"0357} \SUPPORTEDANCH{\char"01B4\char"0357} \SUPPORTEDANCH{\char"1EFF\char"0357}

\SUPPORTEDANCH{\char"0048\char"035B} \SUPPORTEDFALL{\char"2C67\char"035B} \SUPPORTEDANCH{\char"004B\char"035B} \SUPPORTEDANCH{\char"2C69\char"035B} \SUPPORTEDANCH{\char"004D\char"035B} \SUPPORTEDANCH{\char"2C6E\char"035B} \SUPPORTEDANCH{\char"0050\char"035B} \SUPPORTEDANCH{\char"01A4\char"035B} \SUPPORTEDANCH{\char"0052\char"035B} \SUPPORTEDFALL{\char"2C64\char"035B} \SUPPORTEDANCH{\char"0054\char"035B} \SUPPORTEDFALL{\char"01AC\char"035B} \SUPPORTEDFALL{\char"01AE\char"035B} \SUPPORTEDANCH{\char"0059\char"035B} \SUPPORTEDFALL{\char"01B3\char"035B} \SUPPORTEDANCH{\char"0061\char"035B} \SUPPORTEDFALL{\char"1E9A\char"035B} \SUPPORTEDFALL{\char"0062\char"035B} \SUPPORTEDFALL{\char"0183\char"035B} \SUPPORTEDANCH{\char"0063\char"035B} \SUPPORTEDFALL{\char"0188\char"035B} \SUPPORTEDANCH{\char"0064\char"035B} \SUPPORTEDANCH{\char"018C\char"035B} \SUPPORTEDFALL{\char"0221\char"035B} \SUPPORTEDANCH{\char"0068\char"035B} \SUPPORTEDANCH{\char"2C68\char"035B} \SUPPORTEDANCH{\char"006B\char"035B} \SUPPORTEDFALL{\char"0199\char"035B} \SUPPORTEDANCH{\char"2C6A\char"035B} \SUPPORTEDANCH{\char"006C\char"035B} \SUPPORTEDFALL{\char"0234\char"035B} \SUPPORTEDANCH{\char"006E\char"035B} \SUPPORTEDANCH{\char"0235\char"035B} \SUPPORTEDANCH{\char"0070\char"035B} \SUPPORTEDFALL{\char"01A5\char"035B} \SUPPORTEDANCH{\char"0071\char"035B} \SUPPORTEDANCH{\char"02A0\char"035B} \SUPPORTEDFALL{\char"0074\char"035B} \SUPPORTEDFALL{\char"01AB\char"035B} \SUPPORTEDFALL{\char"01AD\char"035B} \SUPPORTEDFALL{\char"0236\char"035B} \SUPPORTEDANCH{\char"0075\char"035B} \SUPPORTEDANCH{\char"AB52\char"035B} \SUPPORTEDANCH{\char"0079\char"035B} \SUPPORTEDANCH{\char"01B4\char"035B} \SUPPORTEDANCH{\char"1EFF\char"035B}


\newpage
\subsection*{Regular patched, Group 4, above, Above}

{\LibertinusSerifPatch
% grid. columns are bases, rows are marks.
\SUPPORTEDANCH{\char"0048\char"0300} \SUPPORTEDANCH{\char"2C67\char"0300} \SUPPORTEDANCH{\char"004B\char"0300} \SUPPORTEDANCH{\char"2C69\char"0300} \SUPPORTEDANCH{\char"004D\char"0300} \SUPPORTEDANCH{\char"2C6E\char"0300} \SUPPORTEDANCH{\char"0050\char"0300} \SUPPORTEDANCH{\char"01A4\char"0300} \SUPPORTEDANCH{\char"0052\char"0300} \SUPPORTEDANCH{\char"2C64\char"0300} \SUPPORTEDANCH{\char"0054\char"0300} \SUPPORTEDANCH{\char"01AC\char"0300} \SUPPORTEDANCH{\char"01AE\char"0300} \SUPPORTEDPREC{\char"0059\char"0300} \SUPPORTEDANCH{\char"01B3\char"0300} \SUPPORTEDPREC{\char"0061\char"0300} \SUPPORTEDANCH{\char"1E9A\char"0300} \SUPPORTEDANCH{\char"0062\char"0300} \SUPPORTEDANCH{\char"0183\char"0300} \SUPPORTEDANCH{\char"0063\char"0300} \SUPPORTEDANCH{\char"0188\char"0300} \SUPPORTEDANCH{\char"0064\char"0300} \SUPPORTEDANCH{\char"018C\char"0300} \SUPPORTEDANCH{\char"0221\char"0300} \SUPPORTEDANCH{\char"0068\char"0300} \SUPPORTEDANCH{\char"2C68\char"0300} \SUPPORTEDANCH{\char"006B\char"0300} \SUPPORTEDANCH{\char"0199\char"0300} \SUPPORTEDANCH{\char"2C6A\char"0300} \SUPPORTEDANCH{\char"006C\char"0300} \SUPPORTEDANCH{\char"0234\char"0300} \SUPPORTEDPREC{\char"006E\char"0300} \SUPPORTEDANCH{\char"0235\char"0300} \SUPPORTEDANCH{\char"0070\char"0300} \SUPPORTEDANCH{\char"01A5\char"0300} \SUPPORTEDANCH{\char"0071\char"0300} \SUPPORTEDANCH{\char"02A0\char"0300} \SUPPORTEDANCH{\char"0074\char"0300} \SUPPORTEDANCH{\char"01AB\char"0300} \SUPPORTEDANCH{\char"01AD\char"0300} \SUPPORTEDANCH{\char"0236\char"0300} \SUPPORTEDPREC{\char"0075\char"0300} \SUPPORTEDANCH{\char"AB52\char"0300} \SUPPORTEDPREC{\char"0079\char"0300} \SUPPORTEDANCH{\char"01B4\char"0300} \SUPPORTEDANCH{\char"1EFF\char"0300}

\SUPPORTEDANCH{\char"0048\char"0301} \SUPPORTEDANCH{\char"2C67\char"0301} \SUPPORTEDPREC{\char"004B\char"0301} \SUPPORTEDANCH{\char"2C69\char"0301} \SUPPORTEDPREC{\char"004D\char"0301} \SUPPORTEDANCH{\char"2C6E\char"0301} \SUPPORTEDPREC{\char"0050\char"0301} \SUPPORTEDANCH{\char"01A4\char"0301} \SUPPORTEDPREC{\char"0052\char"0301} \SUPPORTEDANCH{\char"2C64\char"0301} \SUPPORTEDANCH{\char"0054\char"0301} \SUPPORTEDANCH{\char"01AC\char"0301} \SUPPORTEDANCH{\char"01AE\char"0301} \SUPPORTEDPREC{\char"0059\char"0301} \SUPPORTEDANCH{\char"01B3\char"0301} \SUPPORTEDPREC{\char"0061\char"0301} \SUPPORTEDANCH{\char"1E9A\char"0301} \SUPPORTEDANCH{\char"0062\char"0301} \SUPPORTEDANCH{\char"0183\char"0301} \SUPPORTEDPREC{\char"0063\char"0301} \SUPPORTEDANCH{\char"0188\char"0301} \SUPPORTEDANCH{\char"0064\char"0301} \SUPPORTEDANCH{\char"018C\char"0301} \SUPPORTEDANCH{\char"0221\char"0301} \SUPPORTEDANCH{\char"0068\char"0301} \SUPPORTEDANCH{\char"2C68\char"0301} \SUPPORTEDPREC{\char"006B\char"0301} \SUPPORTEDANCH{\char"0199\char"0301} \SUPPORTEDANCH{\char"2C6A\char"0301} \SUPPORTEDPREC{\char"006C\char"0301} \SUPPORTEDANCH{\char"0234\char"0301} \SUPPORTEDPREC{\char"006E\char"0301} \SUPPORTEDANCH{\char"0235\char"0301} \SUPPORTEDPREC{\char"0070\char"0301} \SUPPORTEDANCH{\char"01A5\char"0301} \SUPPORTEDANCH{\char"0071\char"0301} \SUPPORTEDANCH{\char"02A0\char"0301} \SUPPORTEDANCH{\char"0074\char"0301} \SUPPORTEDANCH{\char"01AB\char"0301} \SUPPORTEDANCH{\char"01AD\char"0301} \SUPPORTEDANCH{\char"0236\char"0301} \SUPPORTEDPREC{\char"0075\char"0301} \SUPPORTEDANCH{\char"AB52\char"0301} \SUPPORTEDPREC{\char"0079\char"0301} \SUPPORTEDANCH{\char"01B4\char"0301} \SUPPORTEDANCH{\char"1EFF\char"0301}

\SUPPORTEDPREC{\char"0048\char"0302} \SUPPORTEDANCH{\char"2C67\char"0302} \SUPPORTEDANCH{\char"004B\char"0302} \SUPPORTEDANCH{\char"2C69\char"0302} \SUPPORTEDANCH{\char"004D\char"0302} \SUPPORTEDANCH{\char"2C6E\char"0302} \SUPPORTEDANCH{\char"0050\char"0302} \SUPPORTEDANCH{\char"01A4\char"0302} \SUPPORTEDANCH{\char"0052\char"0302} \SUPPORTEDANCH{\char"2C64\char"0302} \SUPPORTEDANCH{\char"0054\char"0302} \SUPPORTEDANCH{\char"01AC\char"0302} \SUPPORTEDANCH{\char"01AE\char"0302} \SUPPORTEDPREC{\char"0059\char"0302} \SUPPORTEDANCH{\char"01B3\char"0302} \SUPPORTEDPREC{\char"0061\char"0302} \SUPPORTEDANCH{\char"1E9A\char"0302} \SUPPORTEDANCH{\char"0062\char"0302} \SUPPORTEDANCH{\char"0183\char"0302} \SUPPORTEDPREC{\char"0063\char"0302} \SUPPORTEDANCH{\char"0188\char"0302} \SUPPORTEDANCH{\char"0064\char"0302} \SUPPORTEDANCH{\char"018C\char"0302} \SUPPORTEDANCH{\char"0221\char"0302} \SUPPORTEDPREC{\char"0068\char"0302} \SUPPORTEDANCH{\char"2C68\char"0302} \SUPPORTEDANCH{\char"006B\char"0302} \SUPPORTEDANCH{\char"0199\char"0302} \SUPPORTEDANCH{\char"2C6A\char"0302} \SUPPORTEDANCH{\char"006C\char"0302} \SUPPORTEDANCH{\char"0234\char"0302} \SUPPORTEDANCH{\char"006E\char"0302} \SUPPORTEDANCH{\char"0235\char"0302} \SUPPORTEDANCH{\char"0070\char"0302} \SUPPORTEDANCH{\char"01A5\char"0302} \SUPPORTEDANCH{\char"0071\char"0302} \SUPPORTEDANCH{\char"02A0\char"0302} \SUPPORTEDANCH{\char"0074\char"0302} \SUPPORTEDANCH{\char"01AB\char"0302} \SUPPORTEDANCH{\char"01AD\char"0302} \SUPPORTEDANCH{\char"0236\char"0302} \SUPPORTEDPREC{\char"0075\char"0302} \SUPPORTEDANCH{\char"AB52\char"0302} \SUPPORTEDPREC{\char"0079\char"0302} \SUPPORTEDANCH{\char"01B4\char"0302} \SUPPORTEDANCH{\char"1EFF\char"0302}

\SUPPORTEDANCH{\char"0048\char"0303} \SUPPORTEDANCH{\char"2C67\char"0303} \SUPPORTEDANCH{\char"004B\char"0303} \SUPPORTEDANCH{\char"2C69\char"0303} \SUPPORTEDANCH{\char"004D\char"0303} \SUPPORTEDANCH{\char"2C6E\char"0303} \SUPPORTEDANCH{\char"0050\char"0303} \SUPPORTEDANCH{\char"01A4\char"0303} \SUPPORTEDANCH{\char"0052\char"0303} \SUPPORTEDANCH{\char"2C64\char"0303} \SUPPORTEDANCH{\char"0054\char"0303} \SUPPORTEDANCH{\char"01AC\char"0303} \SUPPORTEDANCH{\char"01AE\char"0303} \SUPPORTEDPREC{\char"0059\char"0303} \SUPPORTEDANCH{\char"01B3\char"0303} \SUPPORTEDPREC{\char"0061\char"0303} \SUPPORTEDANCH{\char"1E9A\char"0303} \SUPPORTEDANCH{\char"0062\char"0303} \SUPPORTEDANCH{\char"0183\char"0303} \SUPPORTEDANCH{\char"0063\char"0303} \SUPPORTEDANCH{\char"0188\char"0303} \SUPPORTEDANCH{\char"0064\char"0303} \SUPPORTEDANCH{\char"018C\char"0303} \SUPPORTEDANCH{\char"0221\char"0303} \SUPPORTEDANCH{\char"0068\char"0303} \SUPPORTEDANCH{\char"2C68\char"0303} \SUPPORTEDANCH{\char"006B\char"0303} \SUPPORTEDANCH{\char"0199\char"0303} \SUPPORTEDANCH{\char"2C6A\char"0303} \SUPPORTEDANCH{\char"006C\char"0303} \SUPPORTEDANCH{\char"0234\char"0303} \SUPPORTEDPREC{\char"006E\char"0303} \SUPPORTEDANCH{\char"0235\char"0303} \SUPPORTEDANCH{\char"0070\char"0303} \SUPPORTEDANCH{\char"01A5\char"0303} \SUPPORTEDANCH{\char"0071\char"0303} \SUPPORTEDANCH{\char"02A0\char"0303} \SUPPORTEDANCH{\char"0074\char"0303} \SUPPORTEDANCH{\char"01AB\char"0303} \SUPPORTEDANCH{\char"01AD\char"0303} \SUPPORTEDANCH{\char"0236\char"0303} \SUPPORTEDPREC{\char"0075\char"0303} \SUPPORTEDANCH{\char"AB52\char"0303} \SUPPORTEDPREC{\char"0079\char"0303} \SUPPORTEDANCH{\char"01B4\char"0303} \SUPPORTEDANCH{\char"1EFF\char"0303}

\SUPPORTEDANCH{\char"0048\char"0304} \SUPPORTEDANCH{\char"2C67\char"0304} \SUPPORTEDANCH{\char"004B\char"0304} \SUPPORTEDANCH{\char"2C69\char"0304} \SUPPORTEDANCH{\char"004D\char"0304} \SUPPORTEDANCH{\char"2C6E\char"0304} \SUPPORTEDANCH{\char"0050\char"0304} \SUPPORTEDANCH{\char"01A4\char"0304} \SUPPORTEDANCH{\char"0052\char"0304} \SUPPORTEDANCH{\char"2C64\char"0304} \SUPPORTEDANCH{\char"0054\char"0304} \SUPPORTEDANCH{\char"01AC\char"0304} \SUPPORTEDANCH{\char"01AE\char"0304} \SUPPORTEDPREC{\char"0059\char"0304} \SUPPORTEDANCH{\char"01B3\char"0304} \SUPPORTEDPREC{\char"0061\char"0304} \SUPPORTEDANCH{\char"1E9A\char"0304} \SUPPORTEDANCH{\char"0062\char"0304} \SUPPORTEDANCH{\char"0183\char"0304} \SUPPORTEDANCH{\char"0063\char"0304} \SUPPORTEDANCH{\char"0188\char"0304} \SUPPORTEDANCH{\char"0064\char"0304} \SUPPORTEDANCH{\char"018C\char"0304} \SUPPORTEDANCH{\char"0221\char"0304} \SUPPORTEDANCH{\char"0068\char"0304} \SUPPORTEDANCH{\char"2C68\char"0304} \SUPPORTEDANCH{\char"006B\char"0304} \SUPPORTEDANCH{\char"0199\char"0304} \SUPPORTEDANCH{\char"2C6A\char"0304} \SUPPORTEDANCH{\char"006C\char"0304} \SUPPORTEDANCH{\char"0234\char"0304} \SUPPORTEDANCH{\char"006E\char"0304} \SUPPORTEDANCH{\char"0235\char"0304} \SUPPORTEDANCH{\char"0070\char"0304} \SUPPORTEDANCH{\char"01A5\char"0304} \SUPPORTEDANCH{\char"0071\char"0304} \SUPPORTEDANCH{\char"02A0\char"0304} \SUPPORTEDANCH{\char"0074\char"0304} \SUPPORTEDANCH{\char"01AB\char"0304} \SUPPORTEDANCH{\char"01AD\char"0304} \SUPPORTEDANCH{\char"0236\char"0304} \SUPPORTEDPREC{\char"0075\char"0304} \SUPPORTEDANCH{\char"AB52\char"0304} \SUPPORTEDPREC{\char"0079\char"0304} \SUPPORTEDANCH{\char"01B4\char"0304} \SUPPORTEDANCH{\char"1EFF\char"0304}

\SUPPORTEDANCH{\char"0048\char"0306} \SUPPORTEDANCH{\char"2C67\char"0306} \SUPPORTEDANCH{\char"004B\char"0306} \SUPPORTEDANCH{\char"2C69\char"0306} \SUPPORTEDANCH{\char"004D\char"0306} \SUPPORTEDANCH{\char"2C6E\char"0306} \SUPPORTEDANCH{\char"0050\char"0306} \SUPPORTEDANCH{\char"01A4\char"0306} \SUPPORTEDANCH{\char"0052\char"0306} \SUPPORTEDANCH{\char"2C64\char"0306} \SUPPORTEDANCH{\char"0054\char"0306} \SUPPORTEDANCH{\char"01AC\char"0306} \SUPPORTEDANCH{\char"01AE\char"0306} \SUPPORTEDANCH{\char"0059\char"0306} \SUPPORTEDANCH{\char"01B3\char"0306} \SUPPORTEDPREC{\char"0061\char"0306} \SUPPORTEDANCH{\char"1E9A\char"0306} \SUPPORTEDANCH{\char"0062\char"0306} \SUPPORTEDANCH{\char"0183\char"0306} \SUPPORTEDANCH{\char"0063\char"0306} \SUPPORTEDANCH{\char"0188\char"0306} \SUPPORTEDANCH{\char"0064\char"0306} \SUPPORTEDANCH{\char"018C\char"0306} \SUPPORTEDANCH{\char"0221\char"0306} \SUPPORTEDANCH{\char"0068\char"0306} \SUPPORTEDANCH{\char"2C68\char"0306} \SUPPORTEDANCH{\char"006B\char"0306} \SUPPORTEDANCH{\char"0199\char"0306} \SUPPORTEDANCH{\char"2C6A\char"0306} \SUPPORTEDANCH{\char"006C\char"0306} \SUPPORTEDANCH{\char"0234\char"0306} \SUPPORTEDANCH{\char"006E\char"0306} \SUPPORTEDANCH{\char"0235\char"0306} \SUPPORTEDANCH{\char"0070\char"0306} \SUPPORTEDANCH{\char"01A5\char"0306} \SUPPORTEDANCH{\char"0071\char"0306} \SUPPORTEDANCH{\char"02A0\char"0306} \SUPPORTEDANCH{\char"0074\char"0306} \SUPPORTEDANCH{\char"01AB\char"0306} \SUPPORTEDANCH{\char"01AD\char"0306} \SUPPORTEDANCH{\char"0236\char"0306} \SUPPORTEDPREC{\char"0075\char"0306} \SUPPORTEDANCH{\char"AB52\char"0306} \SUPPORTEDANCH{\char"0079\char"0306} \SUPPORTEDANCH{\char"01B4\char"0306} \SUPPORTEDANCH{\char"1EFF\char"0306}

\SUPPORTEDPREC{\char"0048\char"0307} \SUPPORTEDANCH{\char"2C67\char"0307} \SUPPORTEDANCH{\char"004B\char"0307} \SUPPORTEDANCH{\char"2C69\char"0307} \SUPPORTEDPREC{\char"004D\char"0307} \SUPPORTEDANCH{\char"2C6E\char"0307} \SUPPORTEDPREC{\char"0050\char"0307} \SUPPORTEDANCH{\char"01A4\char"0307} \SUPPORTEDPREC{\char"0052\char"0307} \SUPPORTEDANCH{\char"2C64\char"0307} \SUPPORTEDPREC{\char"0054\char"0307} \SUPPORTEDANCH{\char"01AC\char"0307} \SUPPORTEDANCH{\char"01AE\char"0307} \SUPPORTEDPREC{\char"0059\char"0307} \SUPPORTEDANCH{\char"01B3\char"0307} \SUPPORTEDPREC{\char"0061\char"0307} \SUPPORTEDANCH{\char"1E9A\char"0307} \SUPPORTEDPREC{\char"0062\char"0307} \SUPPORTEDANCH{\char"0183\char"0307} \SUPPORTEDPREC{\char"0063\char"0307} \SUPPORTEDANCH{\char"0188\char"0307} \SUPPORTEDPREC{\char"0064\char"0307} \SUPPORTEDANCH{\char"018C\char"0307} \SUPPORTEDANCH{\char"0221\char"0307} \SUPPORTEDPREC{\char"0068\char"0307} \SUPPORTEDANCH{\char"2C68\char"0307} \SUPPORTEDANCH{\char"006B\char"0307} \SUPPORTEDANCH{\char"0199\char"0307} \SUPPORTEDANCH{\char"2C6A\char"0307} \SUPPORTEDANCH{\char"006C\char"0307} \SUPPORTEDANCH{\char"0234\char"0307} \SUPPORTEDPREC{\char"006E\char"0307} \SUPPORTEDANCH{\char"0235\char"0307} \SUPPORTEDPREC{\char"0070\char"0307} \SUPPORTEDANCH{\char"01A5\char"0307} \SUPPORTEDANCH{\char"0071\char"0307} \SUPPORTEDANCH{\char"02A0\char"0307} \SUPPORTEDPREC{\char"0074\char"0307} \SUPPORTEDANCH{\char"01AB\char"0307} \SUPPORTEDANCH{\char"01AD\char"0307} \SUPPORTEDANCH{\char"0236\char"0307} \SUPPORTEDANCH{\char"0075\char"0307} \SUPPORTEDANCH{\char"AB52\char"0307} \SUPPORTEDPREC{\char"0079\char"0307} \SUPPORTEDANCH{\char"01B4\char"0307} \SUPPORTEDANCH{\char"1EFF\char"0307}

\SUPPORTEDPREC{\char"0048\char"0308} \SUPPORTEDANCH{\char"2C67\char"0308} \SUPPORTEDANCH{\char"004B\char"0308} \SUPPORTEDANCH{\char"2C69\char"0308} \SUPPORTEDANCH{\char"004D\char"0308} \SUPPORTEDANCH{\char"2C6E\char"0308} \SUPPORTEDANCH{\char"0050\char"0308} \SUPPORTEDANCH{\char"01A4\char"0308} \SUPPORTEDANCH{\char"0052\char"0308} \SUPPORTEDANCH{\char"2C64\char"0308} \SUPPORTEDANCH{\char"0054\char"0308} \SUPPORTEDANCH{\char"01AC\char"0308} \SUPPORTEDANCH{\char"01AE\char"0308} \SUPPORTEDPREC{\char"0059\char"0308} \SUPPORTEDANCH{\char"01B3\char"0308} \SUPPORTEDPREC{\char"0061\char"0308} \SUPPORTEDANCH{\char"1E9A\char"0308} \SUPPORTEDANCH{\char"0062\char"0308} \SUPPORTEDANCH{\char"0183\char"0308} \SUPPORTEDANCH{\char"0063\char"0308} \SUPPORTEDANCH{\char"0188\char"0308} \SUPPORTEDANCH{\char"0064\char"0308} \SUPPORTEDANCH{\char"018C\char"0308} \SUPPORTEDANCH{\char"0221\char"0308} \SUPPORTEDPREC{\char"0068\char"0308} \SUPPORTEDANCH{\char"2C68\char"0308} \SUPPORTEDANCH{\char"006B\char"0308} \SUPPORTEDANCH{\char"0199\char"0308} \SUPPORTEDANCH{\char"2C6A\char"0308} \SUPPORTEDANCH{\char"006C\char"0308} \SUPPORTEDANCH{\char"0234\char"0308} \SUPPORTEDANCH{\char"006E\char"0308} \SUPPORTEDANCH{\char"0235\char"0308} \SUPPORTEDANCH{\char"0070\char"0308} \SUPPORTEDANCH{\char"01A5\char"0308} \SUPPORTEDANCH{\char"0071\char"0308} \SUPPORTEDANCH{\char"02A0\char"0308} \SUPPORTEDPREC{\char"0074\char"0308} \SUPPORTEDANCH{\char"01AB\char"0308} \SUPPORTEDANCH{\char"01AD\char"0308} \SUPPORTEDANCH{\char"0236\char"0308} \SUPPORTEDPREC{\char"0075\char"0308} \SUPPORTEDANCH{\char"AB52\char"0308} \SUPPORTEDPREC{\char"0079\char"0308} \SUPPORTEDANCH{\char"01B4\char"0308} \SUPPORTEDANCH{\char"1EFF\char"0308}

\SUPPORTEDANCH{\char"0048\char"0309} \SUPPORTEDANCH{\char"2C67\char"0309} \SUPPORTEDANCH{\char"004B\char"0309} \SUPPORTEDANCH{\char"2C69\char"0309} \SUPPORTEDANCH{\char"004D\char"0309} \SUPPORTEDANCH{\char"2C6E\char"0309} \SUPPORTEDANCH{\char"0050\char"0309} \SUPPORTEDANCH{\char"01A4\char"0309} \SUPPORTEDANCH{\char"0052\char"0309} \SUPPORTEDANCH{\char"2C64\char"0309} \SUPPORTEDANCH{\char"0054\char"0309} \SUPPORTEDANCH{\char"01AC\char"0309} \SUPPORTEDANCH{\char"01AE\char"0309} \SUPPORTEDPREC{\char"0059\char"0309} \SUPPORTEDANCH{\char"01B3\char"0309} \SUPPORTEDPREC{\char"0061\char"0309} \SUPPORTEDANCH{\char"1E9A\char"0309} \SUPPORTEDANCH{\char"0062\char"0309} \SUPPORTEDANCH{\char"0183\char"0309} \SUPPORTEDANCH{\char"0063\char"0309} \SUPPORTEDANCH{\char"0188\char"0309} \SUPPORTEDANCH{\char"0064\char"0309} \SUPPORTEDANCH{\char"018C\char"0309} \SUPPORTEDANCH{\char"0221\char"0309} \SUPPORTEDANCH{\char"0068\char"0309} \SUPPORTEDANCH{\char"2C68\char"0309} \SUPPORTEDANCH{\char"006B\char"0309} \SUPPORTEDANCH{\char"0199\char"0309} \SUPPORTEDANCH{\char"2C6A\char"0309} \SUPPORTEDANCH{\char"006C\char"0309} \SUPPORTEDANCH{\char"0234\char"0309} \SUPPORTEDANCH{\char"006E\char"0309} \SUPPORTEDANCH{\char"0235\char"0309} \SUPPORTEDANCH{\char"0070\char"0309} \SUPPORTEDANCH{\char"01A5\char"0309} \SUPPORTEDANCH{\char"0071\char"0309} \SUPPORTEDANCH{\char"02A0\char"0309} \SUPPORTEDANCH{\char"0074\char"0309} \SUPPORTEDANCH{\char"01AB\char"0309} \SUPPORTEDANCH{\char"01AD\char"0309} \SUPPORTEDANCH{\char"0236\char"0309} \SUPPORTEDPREC{\char"0075\char"0309} \SUPPORTEDANCH{\char"AB52\char"0309} \SUPPORTEDPREC{\char"0079\char"0309} \SUPPORTEDANCH{\char"01B4\char"0309} \SUPPORTEDANCH{\char"1EFF\char"0309}

\SUPPORTEDANCH{\char"0048\char"030A} \SUPPORTEDANCH{\char"2C67\char"030A} \SUPPORTEDANCH{\char"004B\char"030A} \SUPPORTEDANCH{\char"2C69\char"030A} \SUPPORTEDANCH{\char"004D\char"030A} \SUPPORTEDANCH{\char"2C6E\char"030A} \SUPPORTEDANCH{\char"0050\char"030A} \SUPPORTEDANCH{\char"01A4\char"030A} \SUPPORTEDANCH{\char"0052\char"030A} \SUPPORTEDANCH{\char"2C64\char"030A} \SUPPORTEDANCH{\char"0054\char"030A} \SUPPORTEDANCH{\char"01AC\char"030A} \SUPPORTEDANCH{\char"01AE\char"030A} \SUPPORTEDANCH{\char"0059\char"030A} \SUPPORTEDANCH{\char"01B3\char"030A} \SUPPORTEDPREC{\char"0061\char"030A} \SUPPORTEDANCH{\char"1E9A\char"030A} \SUPPORTEDANCH{\char"0062\char"030A} \SUPPORTEDANCH{\char"0183\char"030A} \SUPPORTEDANCH{\char"0063\char"030A} \SUPPORTEDANCH{\char"0188\char"030A} \SUPPORTEDANCH{\char"0064\char"030A} \SUPPORTEDANCH{\char"018C\char"030A} \SUPPORTEDANCH{\char"0221\char"030A} \SUPPORTEDANCH{\char"0068\char"030A} \SUPPORTEDANCH{\char"2C68\char"030A} \SUPPORTEDANCH{\char"006B\char"030A} \SUPPORTEDANCH{\char"0199\char"030A} \SUPPORTEDANCH{\char"2C6A\char"030A} \SUPPORTEDANCH{\char"006C\char"030A} \SUPPORTEDANCH{\char"0234\char"030A} \SUPPORTEDANCH{\char"006E\char"030A} \SUPPORTEDANCH{\char"0235\char"030A} \SUPPORTEDANCH{\char"0070\char"030A} \SUPPORTEDANCH{\char"01A5\char"030A} \SUPPORTEDANCH{\char"0071\char"030A} \SUPPORTEDANCH{\char"02A0\char"030A} \SUPPORTEDANCH{\char"0074\char"030A} \SUPPORTEDANCH{\char"01AB\char"030A} \SUPPORTEDANCH{\char"01AD\char"030A} \SUPPORTEDANCH{\char"0236\char"030A} \SUPPORTEDPREC{\char"0075\char"030A} \SUPPORTEDANCH{\char"AB52\char"030A} \SUPPORTEDPREC{\char"0079\char"030A} \SUPPORTEDANCH{\char"01B4\char"030A} \SUPPORTEDANCH{\char"1EFF\char"030A}

\SUPPORTEDANCH{\char"0048\char"030B} \SUPPORTEDANCH{\char"2C67\char"030B} \SUPPORTEDANCH{\char"004B\char"030B} \SUPPORTEDANCH{\char"2C69\char"030B} \SUPPORTEDANCH{\char"004D\char"030B} \SUPPORTEDANCH{\char"2C6E\char"030B} \SUPPORTEDANCH{\char"0050\char"030B} \SUPPORTEDANCH{\char"01A4\char"030B} \SUPPORTEDANCH{\char"0052\char"030B} \SUPPORTEDANCH{\char"2C64\char"030B} \SUPPORTEDANCH{\char"0054\char"030B} \SUPPORTEDANCH{\char"01AC\char"030B} \SUPPORTEDANCH{\char"01AE\char"030B} \SUPPORTEDANCH{\char"0059\char"030B} \SUPPORTEDANCH{\char"01B3\char"030B} \SUPPORTEDANCH{\char"0061\char"030B} \SUPPORTEDANCH{\char"1E9A\char"030B} \SUPPORTEDANCH{\char"0062\char"030B} \SUPPORTEDANCH{\char"0183\char"030B} \SUPPORTEDANCH{\char"0063\char"030B} \SUPPORTEDANCH{\char"0188\char"030B} \SUPPORTEDANCH{\char"0064\char"030B} \SUPPORTEDANCH{\char"018C\char"030B} \SUPPORTEDANCH{\char"0221\char"030B} \SUPPORTEDANCH{\char"0068\char"030B} \SUPPORTEDANCH{\char"2C68\char"030B} \SUPPORTEDANCH{\char"006B\char"030B} \SUPPORTEDANCH{\char"0199\char"030B} \SUPPORTEDANCH{\char"2C6A\char"030B} \SUPPORTEDANCH{\char"006C\char"030B} \SUPPORTEDANCH{\char"0234\char"030B} \SUPPORTEDANCH{\char"006E\char"030B} \SUPPORTEDANCH{\char"0235\char"030B} \SUPPORTEDANCH{\char"0070\char"030B} \SUPPORTEDANCH{\char"01A5\char"030B} \SUPPORTEDANCH{\char"0071\char"030B} \SUPPORTEDANCH{\char"02A0\char"030B} \SUPPORTEDANCH{\char"0074\char"030B} \SUPPORTEDANCH{\char"01AB\char"030B} \SUPPORTEDANCH{\char"01AD\char"030B} \SUPPORTEDANCH{\char"0236\char"030B} \SUPPORTEDPREC{\char"0075\char"030B} \SUPPORTEDANCH{\char"AB52\char"030B} \SUPPORTEDANCH{\char"0079\char"030B} \SUPPORTEDANCH{\char"01B4\char"030B} \SUPPORTEDANCH{\char"1EFF\char"030B}

\SUPPORTEDPREC{\char"0048\char"030C} \SUPPORTEDANCH{\char"2C67\char"030C} \SUPPORTEDPREC{\char"004B\char"030C} \SUPPORTEDANCH{\char"2C69\char"030C} \SUPPORTEDANCH{\char"004D\char"030C} \SUPPORTEDANCH{\char"2C6E\char"030C} \SUPPORTEDANCH{\char"0050\char"030C} \SUPPORTEDANCH{\char"01A4\char"030C} \SUPPORTEDPREC{\char"0052\char"030C} \SUPPORTEDANCH{\char"2C64\char"030C} \SUPPORTEDPREC{\char"0054\char"030C} \SUPPORTEDANCH{\char"01AC\char"030C} \SUPPORTEDANCH{\char"01AE\char"030C} \SUPPORTEDANCH{\char"0059\char"030C} \SUPPORTEDANCH{\char"01B3\char"030C} \SUPPORTEDPREC{\char"0061\char"030C} \SUPPORTEDANCH{\char"1E9A\char"030C} \SUPPORTEDANCH{\char"0062\char"030C} \SUPPORTEDANCH{\char"0183\char"030C} \SUPPORTEDPREC{\char"0063\char"030C} \SUPPORTEDANCH{\char"0188\char"030C} \SUPPORTEDPREC{\char"0064\char"030C} \SUPPORTEDANCH{\char"018C\char"030C} \SUPPORTEDANCH{\char"0221\char"030C} \SUPPORTEDPREC{\char"0068\char"030C} \SUPPORTEDANCH{\char"2C68\char"030C} \SUPPORTEDPREC{\char"006B\char"030C} \SUPPORTEDANCH{\char"0199\char"030C} \SUPPORTEDANCH{\char"2C6A\char"030C} \SUPPORTEDPREC{\char"006C\char"030C} \SUPPORTEDANCH{\char"0234\char"030C} \SUPPORTEDPREC{\char"006E\char"030C} \SUPPORTEDANCH{\char"0235\char"030C} \SUPPORTEDANCH{\char"0070\char"030C} \SUPPORTEDANCH{\char"01A5\char"030C} \SUPPORTEDANCH{\char"0071\char"030C} \SUPPORTEDANCH{\char"02A0\char"030C} \SUPPORTEDPREC{\char"0074\char"030C} \SUPPORTEDANCH{\char"01AB\char"030C} \SUPPORTEDANCH{\char"01AD\char"030C} \SUPPORTEDANCH{\char"0236\char"030C} \SUPPORTEDPREC{\char"0075\char"030C} \SUPPORTEDANCH{\char"AB52\char"030C} \SUPPORTEDANCH{\char"0079\char"030C} \SUPPORTEDANCH{\char"01B4\char"030C} \SUPPORTEDANCH{\char"1EFF\char"030C}

\SUPPORTEDANCH{\char"0048\char"030D} \SUPPORTEDANCH{\char"2C67\char"030D} \SUPPORTEDANCH{\char"004B\char"030D} \SUPPORTEDANCH{\char"2C69\char"030D} \SUPPORTEDANCH{\char"004D\char"030D} \SUPPORTEDANCH{\char"2C6E\char"030D} \SUPPORTEDANCH{\char"0050\char"030D} \SUPPORTEDANCH{\char"01A4\char"030D} \SUPPORTEDANCH{\char"0052\char"030D} \SUPPORTEDANCH{\char"2C64\char"030D} \SUPPORTEDANCH{\char"0054\char"030D} \SUPPORTEDANCH{\char"01AC\char"030D} \SUPPORTEDANCH{\char"01AE\char"030D} \SUPPORTEDANCH{\char"0059\char"030D} \SUPPORTEDANCH{\char"01B3\char"030D} \SUPPORTEDANCH{\char"0061\char"030D} \SUPPORTEDANCH{\char"1E9A\char"030D} \SUPPORTEDANCH{\char"0062\char"030D} \SUPPORTEDANCH{\char"0183\char"030D} \SUPPORTEDANCH{\char"0063\char"030D} \SUPPORTEDANCH{\char"0188\char"030D} \SUPPORTEDANCH{\char"0064\char"030D} \SUPPORTEDANCH{\char"018C\char"030D} \SUPPORTEDANCH{\char"0221\char"030D} \SUPPORTEDANCH{\char"0068\char"030D} \SUPPORTEDANCH{\char"2C68\char"030D} \SUPPORTEDANCH{\char"006B\char"030D} \SUPPORTEDANCH{\char"0199\char"030D} \SUPPORTEDANCH{\char"2C6A\char"030D} \SUPPORTEDANCH{\char"006C\char"030D} \SUPPORTEDANCH{\char"0234\char"030D} \SUPPORTEDANCH{\char"006E\char"030D} \SUPPORTEDANCH{\char"0235\char"030D} \SUPPORTEDANCH{\char"0070\char"030D} \SUPPORTEDANCH{\char"01A5\char"030D} \SUPPORTEDANCH{\char"0071\char"030D} \SUPPORTEDANCH{\char"02A0\char"030D} \SUPPORTEDANCH{\char"0074\char"030D} \SUPPORTEDANCH{\char"01AB\char"030D} \SUPPORTEDANCH{\char"01AD\char"030D} \SUPPORTEDANCH{\char"0236\char"030D} \SUPPORTEDANCH{\char"0075\char"030D} \SUPPORTEDANCH{\char"AB52\char"030D} \SUPPORTEDANCH{\char"0079\char"030D} \SUPPORTEDANCH{\char"01B4\char"030D} \SUPPORTEDANCH{\char"1EFF\char"030D}

\SUPPORTEDANCH{\char"0048\char"030F} \SUPPORTEDANCH{\char"2C67\char"030F} \SUPPORTEDANCH{\char"004B\char"030F} \SUPPORTEDANCH{\char"2C69\char"030F} \SUPPORTEDANCH{\char"004D\char"030F} \SUPPORTEDANCH{\char"2C6E\char"030F} \SUPPORTEDANCH{\char"0050\char"030F} \SUPPORTEDANCH{\char"01A4\char"030F} \SUPPORTEDPREC{\char"0052\char"030F} \SUPPORTEDANCH{\char"2C64\char"030F} \SUPPORTEDANCH{\char"0054\char"030F} \SUPPORTEDANCH{\char"01AC\char"030F} \SUPPORTEDANCH{\char"01AE\char"030F} \SUPPORTEDANCH{\char"0059\char"030F} \SUPPORTEDANCH{\char"01B3\char"030F} \SUPPORTEDPREC{\char"0061\char"030F} \SUPPORTEDANCH{\char"1E9A\char"030F} \SUPPORTEDANCH{\char"0062\char"030F} \SUPPORTEDANCH{\char"0183\char"030F} \SUPPORTEDANCH{\char"0063\char"030F} \SUPPORTEDANCH{\char"0188\char"030F} \SUPPORTEDANCH{\char"0064\char"030F} \SUPPORTEDANCH{\char"018C\char"030F} \SUPPORTEDANCH{\char"0221\char"030F} \SUPPORTEDANCH{\char"0068\char"030F} \SUPPORTEDANCH{\char"2C68\char"030F} \SUPPORTEDANCH{\char"006B\char"030F} \SUPPORTEDANCH{\char"0199\char"030F} \SUPPORTEDANCH{\char"2C6A\char"030F} \SUPPORTEDANCH{\char"006C\char"030F} \SUPPORTEDANCH{\char"0234\char"030F} \SUPPORTEDANCH{\char"006E\char"030F} \SUPPORTEDANCH{\char"0235\char"030F} \SUPPORTEDANCH{\char"0070\char"030F} \SUPPORTEDANCH{\char"01A5\char"030F} \SUPPORTEDANCH{\char"0071\char"030F} \SUPPORTEDANCH{\char"02A0\char"030F} \SUPPORTEDANCH{\char"0074\char"030F} \SUPPORTEDANCH{\char"01AB\char"030F} \SUPPORTEDANCH{\char"01AD\char"030F} \SUPPORTEDANCH{\char"0236\char"030F} \SUPPORTEDPREC{\char"0075\char"030F} \SUPPORTEDANCH{\char"AB52\char"030F} \SUPPORTEDANCH{\char"0079\char"030F} \SUPPORTEDANCH{\char"01B4\char"030F} \SUPPORTEDANCH{\char"1EFF\char"030F}

\SUPPORTEDANCH{\char"0048\char"0310} \SUPPORTEDANCH{\char"2C67\char"0310} \SUPPORTEDANCH{\char"004B\char"0310} \SUPPORTEDANCH{\char"2C69\char"0310} \SUPPORTEDANCH{\char"004D\char"0310} \SUPPORTEDANCH{\char"2C6E\char"0310} \SUPPORTEDANCH{\char"0050\char"0310} \SUPPORTEDANCH{\char"01A4\char"0310} \SUPPORTEDANCH{\char"0052\char"0310} \SUPPORTEDANCH{\char"2C64\char"0310} \SUPPORTEDANCH{\char"0054\char"0310} \SUPPORTEDANCH{\char"01AC\char"0310} \SUPPORTEDANCH{\char"01AE\char"0310} \SUPPORTEDANCH{\char"0059\char"0310} \SUPPORTEDANCH{\char"01B3\char"0310} \SUPPORTEDANCH{\char"0061\char"0310} \SUPPORTEDANCH{\char"1E9A\char"0310} \SUPPORTEDANCH{\char"0062\char"0310} \SUPPORTEDANCH{\char"0183\char"0310} \SUPPORTEDANCH{\char"0063\char"0310} \SUPPORTEDANCH{\char"0188\char"0310} \SUPPORTEDANCH{\char"0064\char"0310} \SUPPORTEDANCH{\char"018C\char"0310} \SUPPORTEDANCH{\char"0221\char"0310} \SUPPORTEDANCH{\char"0068\char"0310} \SUPPORTEDANCH{\char"2C68\char"0310} \SUPPORTEDANCH{\char"006B\char"0310} \SUPPORTEDANCH{\char"0199\char"0310} \SUPPORTEDANCH{\char"2C6A\char"0310} \SUPPORTEDANCH{\char"006C\char"0310} \SUPPORTEDANCH{\char"0234\char"0310} \SUPPORTEDANCH{\char"006E\char"0310} \SUPPORTEDANCH{\char"0235\char"0310} \SUPPORTEDANCH{\char"0070\char"0310} \SUPPORTEDANCH{\char"01A5\char"0310} \SUPPORTEDANCH{\char"0071\char"0310} \SUPPORTEDANCH{\char"02A0\char"0310} \SUPPORTEDANCH{\char"0074\char"0310} \SUPPORTEDANCH{\char"01AB\char"0310} \SUPPORTEDANCH{\char"01AD\char"0310} \SUPPORTEDANCH{\char"0236\char"0310} \SUPPORTEDANCH{\char"0075\char"0310} \SUPPORTEDANCH{\char"AB52\char"0310} \SUPPORTEDANCH{\char"0079\char"0310} \SUPPORTEDANCH{\char"01B4\char"0310} \SUPPORTEDANCH{\char"1EFF\char"0310}

\SUPPORTEDANCH{\char"0048\char"0311} \SUPPORTEDANCH{\char"2C67\char"0311} \SUPPORTEDANCH{\char"004B\char"0311} \SUPPORTEDANCH{\char"2C69\char"0311} \SUPPORTEDANCH{\char"004D\char"0311} \SUPPORTEDANCH{\char"2C6E\char"0311} \SUPPORTEDANCH{\char"0050\char"0311} \SUPPORTEDANCH{\char"01A4\char"0311} \SUPPORTEDPREC{\char"0052\char"0311} \SUPPORTEDANCH{\char"2C64\char"0311} \SUPPORTEDANCH{\char"0054\char"0311} \SUPPORTEDANCH{\char"01AC\char"0311} \SUPPORTEDANCH{\char"01AE\char"0311} \SUPPORTEDANCH{\char"0059\char"0311} \SUPPORTEDANCH{\char"01B3\char"0311} \SUPPORTEDPREC{\char"0061\char"0311} \SUPPORTEDANCH{\char"1E9A\char"0311} \SUPPORTEDANCH{\char"0062\char"0311} \SUPPORTEDANCH{\char"0183\char"0311} \SUPPORTEDANCH{\char"0063\char"0311} \SUPPORTEDANCH{\char"0188\char"0311} \SUPPORTEDANCH{\char"0064\char"0311} \SUPPORTEDANCH{\char"018C\char"0311} \SUPPORTEDANCH{\char"0221\char"0311} \SUPPORTEDANCH{\char"0068\char"0311} \SUPPORTEDANCH{\char"2C68\char"0311} \SUPPORTEDANCH{\char"006B\char"0311} \SUPPORTEDANCH{\char"0199\char"0311} \SUPPORTEDANCH{\char"2C6A\char"0311} \SUPPORTEDANCH{\char"006C\char"0311} \SUPPORTEDANCH{\char"0234\char"0311} \SUPPORTEDANCH{\char"006E\char"0311} \SUPPORTEDANCH{\char"0235\char"0311} \SUPPORTEDANCH{\char"0070\char"0311} \SUPPORTEDANCH{\char"01A5\char"0311} \SUPPORTEDANCH{\char"0071\char"0311} \SUPPORTEDANCH{\char"02A0\char"0311} \SUPPORTEDANCH{\char"0074\char"0311} \SUPPORTEDANCH{\char"01AB\char"0311} \SUPPORTEDANCH{\char"01AD\char"0311} \SUPPORTEDANCH{\char"0236\char"0311} \SUPPORTEDPREC{\char"0075\char"0311} \SUPPORTEDANCH{\char"AB52\char"0311} \SUPPORTEDANCH{\char"0079\char"0311} \SUPPORTEDANCH{\char"01B4\char"0311} \SUPPORTEDANCH{\char"1EFF\char"0311}

\SUPPORTEDANCH{\char"0048\char"0312} \SUPPORTEDANCH{\char"2C67\char"0312} \SUPPORTEDANCH{\char"004B\char"0312} \SUPPORTEDANCH{\char"2C69\char"0312} \SUPPORTEDANCH{\char"004D\char"0312} \SUPPORTEDANCH{\char"2C6E\char"0312} \SUPPORTEDANCH{\char"0050\char"0312} \SUPPORTEDANCH{\char"01A4\char"0312} \SUPPORTEDANCH{\char"0052\char"0312} \SUPPORTEDANCH{\char"2C64\char"0312} \SUPPORTEDANCH{\char"0054\char"0312} \SUPPORTEDANCH{\char"01AC\char"0312} \SUPPORTEDANCH{\char"01AE\char"0312} \SUPPORTEDANCH{\char"0059\char"0312} \SUPPORTEDANCH{\char"01B3\char"0312} \SUPPORTEDANCH{\char"0061\char"0312} \SUPPORTEDANCH{\char"1E9A\char"0312} \SUPPORTEDANCH{\char"0062\char"0312} \SUPPORTEDANCH{\char"0183\char"0312} \SUPPORTEDANCH{\char"0063\char"0312} \SUPPORTEDANCH{\char"0188\char"0312} \SUPPORTEDANCH{\char"0064\char"0312} \SUPPORTEDANCH{\char"018C\char"0312} \SUPPORTEDANCH{\char"0221\char"0312} \SUPPORTEDANCH{\char"0068\char"0312} \SUPPORTEDANCH{\char"2C68\char"0312} \SUPPORTEDANCH{\char"006B\char"0312} \SUPPORTEDANCH{\char"0199\char"0312} \SUPPORTEDANCH{\char"2C6A\char"0312} \SUPPORTEDANCH{\char"006C\char"0312} \SUPPORTEDANCH{\char"0234\char"0312} \SUPPORTEDANCH{\char"006E\char"0312} \SUPPORTEDANCH{\char"0235\char"0312} \SUPPORTEDANCH{\char"0070\char"0312} \SUPPORTEDANCH{\char"01A5\char"0312} \SUPPORTEDANCH{\char"0071\char"0312} \SUPPORTEDANCH{\char"02A0\char"0312} \SUPPORTEDANCH{\char"0074\char"0312} \SUPPORTEDANCH{\char"01AB\char"0312} \SUPPORTEDANCH{\char"01AD\char"0312} \SUPPORTEDANCH{\char"0236\char"0312} \SUPPORTEDANCH{\char"0075\char"0312} \SUPPORTEDANCH{\char"AB52\char"0312} \SUPPORTEDANCH{\char"0079\char"0312} \SUPPORTEDANCH{\char"01B4\char"0312} \SUPPORTEDANCH{\char"1EFF\char"0312}

\SUPPORTEDANCH{\char"0048\char"033D} \SUPPORTEDANCH{\char"2C67\char"033D} \SUPPORTEDANCH{\char"004B\char"033D} \SUPPORTEDANCH{\char"2C69\char"033D} \SUPPORTEDANCH{\char"004D\char"033D} \SUPPORTEDANCH{\char"2C6E\char"033D} \SUPPORTEDANCH{\char"0050\char"033D} \SUPPORTEDANCH{\char"01A4\char"033D} \SUPPORTEDANCH{\char"0052\char"033D} \SUPPORTEDANCH{\char"2C64\char"033D} \SUPPORTEDANCH{\char"0054\char"033D} \SUPPORTEDANCH{\char"01AC\char"033D} \SUPPORTEDANCH{\char"01AE\char"033D} \SUPPORTEDANCH{\char"0059\char"033D} \SUPPORTEDANCH{\char"01B3\char"033D} \SUPPORTEDANCH{\char"0061\char"033D} \SUPPORTEDANCH{\char"1E9A\char"033D} \SUPPORTEDANCH{\char"0062\char"033D} \SUPPORTEDANCH{\char"0183\char"033D} \SUPPORTEDANCH{\char"0063\char"033D} \SUPPORTEDANCH{\char"0188\char"033D} \SUPPORTEDANCH{\char"0064\char"033D} \SUPPORTEDANCH{\char"018C\char"033D} \SUPPORTEDANCH{\char"0221\char"033D} \SUPPORTEDANCH{\char"0068\char"033D} \SUPPORTEDANCH{\char"2C68\char"033D} \SUPPORTEDANCH{\char"006B\char"033D} \SUPPORTEDANCH{\char"0199\char"033D} \SUPPORTEDANCH{\char"2C6A\char"033D} \SUPPORTEDANCH{\char"006C\char"033D} \SUPPORTEDANCH{\char"0234\char"033D} \SUPPORTEDANCH{\char"006E\char"033D} \SUPPORTEDANCH{\char"0235\char"033D} \SUPPORTEDANCH{\char"0070\char"033D} \SUPPORTEDANCH{\char"01A5\char"033D} \SUPPORTEDANCH{\char"0071\char"033D} \SUPPORTEDANCH{\char"02A0\char"033D} \SUPPORTEDANCH{\char"0074\char"033D} \SUPPORTEDANCH{\char"01AB\char"033D} \SUPPORTEDANCH{\char"01AD\char"033D} \SUPPORTEDANCH{\char"0236\char"033D} \SUPPORTEDANCH{\char"0075\char"033D} \SUPPORTEDANCH{\char"AB52\char"033D} \SUPPORTEDANCH{\char"0079\char"033D} \SUPPORTEDANCH{\char"01B4\char"033D} \SUPPORTEDANCH{\char"1EFF\char"033D}

\SUPPORTEDANCH{\char"0048\char"033E} \SUPPORTEDANCH{\char"2C67\char"033E} \SUPPORTEDANCH{\char"004B\char"033E} \SUPPORTEDANCH{\char"2C69\char"033E} \SUPPORTEDANCH{\char"004D\char"033E} \SUPPORTEDANCH{\char"2C6E\char"033E} \SUPPORTEDANCH{\char"0050\char"033E} \SUPPORTEDANCH{\char"01A4\char"033E} \SUPPORTEDANCH{\char"0052\char"033E} \SUPPORTEDANCH{\char"2C64\char"033E} \SUPPORTEDANCH{\char"0054\char"033E} \SUPPORTEDANCH{\char"01AC\char"033E} \SUPPORTEDANCH{\char"01AE\char"033E} \SUPPORTEDANCH{\char"0059\char"033E} \SUPPORTEDANCH{\char"01B3\char"033E} \SUPPORTEDANCH{\char"0061\char"033E} \SUPPORTEDANCH{\char"1E9A\char"033E} \SUPPORTEDANCH{\char"0062\char"033E} \SUPPORTEDANCH{\char"0183\char"033E} \SUPPORTEDANCH{\char"0063\char"033E} \SUPPORTEDANCH{\char"0188\char"033E} \SUPPORTEDANCH{\char"0064\char"033E} \SUPPORTEDANCH{\char"018C\char"033E} \SUPPORTEDANCH{\char"0221\char"033E} \SUPPORTEDANCH{\char"0068\char"033E} \SUPPORTEDANCH{\char"2C68\char"033E} \SUPPORTEDANCH{\char"006B\char"033E} \SUPPORTEDANCH{\char"0199\char"033E} \SUPPORTEDANCH{\char"2C6A\char"033E} \SUPPORTEDANCH{\char"006C\char"033E} \SUPPORTEDANCH{\char"0234\char"033E} \SUPPORTEDANCH{\char"006E\char"033E} \SUPPORTEDANCH{\char"0235\char"033E} \SUPPORTEDANCH{\char"0070\char"033E} \SUPPORTEDANCH{\char"01A5\char"033E} \SUPPORTEDANCH{\char"0071\char"033E} \SUPPORTEDANCH{\char"02A0\char"033E} \SUPPORTEDANCH{\char"0074\char"033E} \SUPPORTEDANCH{\char"01AB\char"033E} \SUPPORTEDANCH{\char"01AD\char"033E} \SUPPORTEDANCH{\char"0236\char"033E} \SUPPORTEDANCH{\char"0075\char"033E} \SUPPORTEDANCH{\char"AB52\char"033E} \SUPPORTEDANCH{\char"0079\char"033E} \SUPPORTEDANCH{\char"01B4\char"033E} \SUPPORTEDANCH{\char"1EFF\char"033E}

\SUPPORTEDANCH{\char"0048\char"033F} \SUPPORTEDANCH{\char"2C67\char"033F} \SUPPORTEDANCH{\char"004B\char"033F} \SUPPORTEDANCH{\char"2C69\char"033F} \SUPPORTEDANCH{\char"004D\char"033F} \SUPPORTEDANCH{\char"2C6E\char"033F} \SUPPORTEDANCH{\char"0050\char"033F} \SUPPORTEDANCH{\char"01A4\char"033F} \SUPPORTEDANCH{\char"0052\char"033F} \SUPPORTEDANCH{\char"2C64\char"033F} \SUPPORTEDANCH{\char"0054\char"033F} \SUPPORTEDANCH{\char"01AC\char"033F} \SUPPORTEDANCH{\char"01AE\char"033F} \SUPPORTEDANCH{\char"0059\char"033F} \SUPPORTEDANCH{\char"01B3\char"033F} \SUPPORTEDANCH{\char"0061\char"033F} \SUPPORTEDANCH{\char"1E9A\char"033F} \SUPPORTEDANCH{\char"0062\char"033F} \SUPPORTEDANCH{\char"0183\char"033F} \SUPPORTEDANCH{\char"0063\char"033F} \SUPPORTEDANCH{\char"0188\char"033F} \SUPPORTEDANCH{\char"0064\char"033F} \SUPPORTEDANCH{\char"018C\char"033F} \SUPPORTEDANCH{\char"0221\char"033F} \SUPPORTEDANCH{\char"0068\char"033F} \SUPPORTEDANCH{\char"2C68\char"033F} \SUPPORTEDANCH{\char"006B\char"033F} \SUPPORTEDANCH{\char"0199\char"033F} \SUPPORTEDANCH{\char"2C6A\char"033F} \SUPPORTEDANCH{\char"006C\char"033F} \SUPPORTEDANCH{\char"0234\char"033F} \SUPPORTEDANCH{\char"006E\char"033F} \SUPPORTEDANCH{\char"0235\char"033F} \SUPPORTEDANCH{\char"0070\char"033F} \SUPPORTEDANCH{\char"01A5\char"033F} \SUPPORTEDANCH{\char"0071\char"033F} \SUPPORTEDANCH{\char"02A0\char"033F} \SUPPORTEDANCH{\char"0074\char"033F} \SUPPORTEDANCH{\char"01AB\char"033F} \SUPPORTEDANCH{\char"01AD\char"033F} \SUPPORTEDANCH{\char"0236\char"033F} \SUPPORTEDANCH{\char"0075\char"033F} \SUPPORTEDANCH{\char"AB52\char"033F} \SUPPORTEDANCH{\char"0079\char"033F} \SUPPORTEDANCH{\char"01B4\char"033F} \SUPPORTEDANCH{\char"1EFF\char"033F}

\SUPPORTEDANCH{\char"0048\char"030E} \SUPPORTEDANCH{\char"2C67\char"030E} \SUPPORTEDANCH{\char"004B\char"030E} \SUPPORTEDANCH{\char"2C69\char"030E} \SUPPORTEDANCH{\char"004D\char"030E} \SUPPORTEDANCH{\char"2C6E\char"030E} \SUPPORTEDANCH{\char"0050\char"030E} \SUPPORTEDANCH{\char"01A4\char"030E} \SUPPORTEDANCH{\char"0052\char"030E} \SUPPORTEDANCH{\char"2C64\char"030E} \SUPPORTEDANCH{\char"0054\char"030E} \SUPPORTEDANCH{\char"01AC\char"030E} \SUPPORTEDANCH{\char"01AE\char"030E} \SUPPORTEDANCH{\char"0059\char"030E} \SUPPORTEDANCH{\char"01B3\char"030E} \SUPPORTEDANCH{\char"0061\char"030E} \SUPPORTEDANCH{\char"1E9A\char"030E} \SUPPORTEDANCH{\char"0062\char"030E} \SUPPORTEDANCH{\char"0183\char"030E} \SUPPORTEDANCH{\char"0063\char"030E} \SUPPORTEDANCH{\char"0188\char"030E} \SUPPORTEDANCH{\char"0064\char"030E} \SUPPORTEDANCH{\char"018C\char"030E} \SUPPORTEDANCH{\char"0221\char"030E} \SUPPORTEDANCH{\char"0068\char"030E} \SUPPORTEDANCH{\char"2C68\char"030E} \SUPPORTEDANCH{\char"006B\char"030E} \SUPPORTEDANCH{\char"0199\char"030E} \SUPPORTEDANCH{\char"2C6A\char"030E} \SUPPORTEDANCH{\char"006C\char"030E} \SUPPORTEDANCH{\char"0234\char"030E} \SUPPORTEDANCH{\char"006E\char"030E} \SUPPORTEDANCH{\char"0235\char"030E} \SUPPORTEDANCH{\char"0070\char"030E} \SUPPORTEDANCH{\char"01A5\char"030E} \SUPPORTEDANCH{\char"0071\char"030E} \SUPPORTEDANCH{\char"02A0\char"030E} \SUPPORTEDANCH{\char"0074\char"030E} \SUPPORTEDANCH{\char"01AB\char"030E} \SUPPORTEDANCH{\char"01AD\char"030E} \SUPPORTEDANCH{\char"0236\char"030E} \SUPPORTEDANCH{\char"0075\char"030E} \SUPPORTEDANCH{\char"AB52\char"030E} \SUPPORTEDANCH{\char"0079\char"030E} \SUPPORTEDANCH{\char"01B4\char"030E} \SUPPORTEDANCH{\char"1EFF\char"030E}

\SUPPORTEDANCH{\char"0048\char"0346} \SUPPORTEDANCH{\char"2C67\char"0346} \SUPPORTEDANCH{\char"004B\char"0346} \SUPPORTEDANCH{\char"2C69\char"0346} \SUPPORTEDANCH{\char"004D\char"0346} \SUPPORTEDANCH{\char"2C6E\char"0346} \SUPPORTEDANCH{\char"0050\char"0346} \SUPPORTEDANCH{\char"01A4\char"0346} \SUPPORTEDANCH{\char"0052\char"0346} \SUPPORTEDANCH{\char"2C64\char"0346} \SUPPORTEDANCH{\char"0054\char"0346} \SUPPORTEDANCH{\char"01AC\char"0346} \SUPPORTEDANCH{\char"01AE\char"0346} \SUPPORTEDANCH{\char"0059\char"0346} \SUPPORTEDANCH{\char"01B3\char"0346} \SUPPORTEDANCH{\char"0061\char"0346} \SUPPORTEDANCH{\char"1E9A\char"0346} \SUPPORTEDANCH{\char"0062\char"0346} \SUPPORTEDANCH{\char"0183\char"0346} \SUPPORTEDANCH{\char"0063\char"0346} \SUPPORTEDANCH{\char"0188\char"0346} \SUPPORTEDANCH{\char"0064\char"0346} \SUPPORTEDANCH{\char"018C\char"0346} \SUPPORTEDANCH{\char"0221\char"0346} \SUPPORTEDANCH{\char"0068\char"0346} \SUPPORTEDANCH{\char"2C68\char"0346} \SUPPORTEDANCH{\char"006B\char"0346} \SUPPORTEDANCH{\char"0199\char"0346} \SUPPORTEDANCH{\char"2C6A\char"0346} \SUPPORTEDANCH{\char"006C\char"0346} \SUPPORTEDANCH{\char"0234\char"0346} \SUPPORTEDANCH{\char"006E\char"0346} \SUPPORTEDANCH{\char"0235\char"0346} \SUPPORTEDANCH{\char"0070\char"0346} \SUPPORTEDANCH{\char"01A5\char"0346} \SUPPORTEDANCH{\char"0071\char"0346} \SUPPORTEDANCH{\char"02A0\char"0346} \SUPPORTEDANCH{\char"0074\char"0346} \SUPPORTEDANCH{\char"01AB\char"0346} \SUPPORTEDANCH{\char"01AD\char"0346} \SUPPORTEDANCH{\char"0236\char"0346} \SUPPORTEDANCH{\char"0075\char"0346} \SUPPORTEDANCH{\char"AB52\char"0346} \SUPPORTEDANCH{\char"0079\char"0346} \SUPPORTEDANCH{\char"01B4\char"0346} \SUPPORTEDANCH{\char"1EFF\char"0346}

\SUPPORTEDANCH{\char"0048\char"034A} \SUPPORTEDANCH{\char"2C67\char"034A} \SUPPORTEDANCH{\char"004B\char"034A} \SUPPORTEDANCH{\char"2C69\char"034A} \SUPPORTEDANCH{\char"004D\char"034A} \SUPPORTEDANCH{\char"2C6E\char"034A} \SUPPORTEDANCH{\char"0050\char"034A} \SUPPORTEDANCH{\char"01A4\char"034A} \SUPPORTEDANCH{\char"0052\char"034A} \SUPPORTEDANCH{\char"2C64\char"034A} \SUPPORTEDANCH{\char"0054\char"034A} \SUPPORTEDANCH{\char"01AC\char"034A} \SUPPORTEDANCH{\char"01AE\char"034A} \SUPPORTEDANCH{\char"0059\char"034A} \SUPPORTEDANCH{\char"01B3\char"034A} \SUPPORTEDANCH{\char"0061\char"034A} \SUPPORTEDANCH{\char"1E9A\char"034A} \SUPPORTEDANCH{\char"0062\char"034A} \SUPPORTEDANCH{\char"0183\char"034A} \SUPPORTEDANCH{\char"0063\char"034A} \SUPPORTEDANCH{\char"0188\char"034A} \SUPPORTEDANCH{\char"0064\char"034A} \SUPPORTEDANCH{\char"018C\char"034A} \SUPPORTEDANCH{\char"0221\char"034A} \SUPPORTEDANCH{\char"0068\char"034A} \SUPPORTEDANCH{\char"2C68\char"034A} \SUPPORTEDANCH{\char"006B\char"034A} \SUPPORTEDANCH{\char"0199\char"034A} \SUPPORTEDANCH{\char"2C6A\char"034A} \SUPPORTEDANCH{\char"006C\char"034A} \SUPPORTEDANCH{\char"0234\char"034A} \SUPPORTEDANCH{\char"006E\char"034A} \SUPPORTEDANCH{\char"0235\char"034A} \SUPPORTEDANCH{\char"0070\char"034A} \SUPPORTEDANCH{\char"01A5\char"034A} \SUPPORTEDANCH{\char"0071\char"034A} \SUPPORTEDANCH{\char"02A0\char"034A} \SUPPORTEDANCH{\char"0074\char"034A} \SUPPORTEDANCH{\char"01AB\char"034A} \SUPPORTEDANCH{\char"01AD\char"034A} \SUPPORTEDANCH{\char"0236\char"034A} \SUPPORTEDANCH{\char"0075\char"034A} \SUPPORTEDANCH{\char"AB52\char"034A} \SUPPORTEDANCH{\char"0079\char"034A} \SUPPORTEDANCH{\char"01B4\char"034A} \SUPPORTEDANCH{\char"1EFF\char"034A}

\SUPPORTEDANCH{\char"0048\char"034B} \SUPPORTEDANCH{\char"2C67\char"034B} \SUPPORTEDANCH{\char"004B\char"034B} \SUPPORTEDANCH{\char"2C69\char"034B} \SUPPORTEDANCH{\char"004D\char"034B} \SUPPORTEDANCH{\char"2C6E\char"034B} \SUPPORTEDANCH{\char"0050\char"034B} \SUPPORTEDANCH{\char"01A4\char"034B} \SUPPORTEDANCH{\char"0052\char"034B} \SUPPORTEDANCH{\char"2C64\char"034B} \SUPPORTEDANCH{\char"0054\char"034B} \SUPPORTEDANCH{\char"01AC\char"034B} \SUPPORTEDANCH{\char"01AE\char"034B} \SUPPORTEDANCH{\char"0059\char"034B} \SUPPORTEDANCH{\char"01B3\char"034B} \SUPPORTEDANCH{\char"0061\char"034B} \SUPPORTEDANCH{\char"1E9A\char"034B} \SUPPORTEDANCH{\char"0062\char"034B} \SUPPORTEDANCH{\char"0183\char"034B} \SUPPORTEDANCH{\char"0063\char"034B} \SUPPORTEDANCH{\char"0188\char"034B} \SUPPORTEDANCH{\char"0064\char"034B} \SUPPORTEDANCH{\char"018C\char"034B} \SUPPORTEDANCH{\char"0221\char"034B} \SUPPORTEDANCH{\char"0068\char"034B} \SUPPORTEDANCH{\char"2C68\char"034B} \SUPPORTEDANCH{\char"006B\char"034B} \SUPPORTEDANCH{\char"0199\char"034B} \SUPPORTEDANCH{\char"2C6A\char"034B} \SUPPORTEDANCH{\char"006C\char"034B} \SUPPORTEDANCH{\char"0234\char"034B} \SUPPORTEDANCH{\char"006E\char"034B} \SUPPORTEDANCH{\char"0235\char"034B} \SUPPORTEDANCH{\char"0070\char"034B} \SUPPORTEDANCH{\char"01A5\char"034B} \SUPPORTEDANCH{\char"0071\char"034B} \SUPPORTEDANCH{\char"02A0\char"034B} \SUPPORTEDANCH{\char"0074\char"034B} \SUPPORTEDANCH{\char"01AB\char"034B} \SUPPORTEDANCH{\char"01AD\char"034B} \SUPPORTEDANCH{\char"0236\char"034B} \SUPPORTEDANCH{\char"0075\char"034B} \SUPPORTEDANCH{\char"AB52\char"034B} \SUPPORTEDANCH{\char"0079\char"034B} \SUPPORTEDANCH{\char"01B4\char"034B} \SUPPORTEDANCH{\char"1EFF\char"034B}

\SUPPORTEDANCH{\char"0048\char"034C} \SUPPORTEDANCH{\char"2C67\char"034C} \SUPPORTEDANCH{\char"004B\char"034C} \SUPPORTEDANCH{\char"2C69\char"034C} \SUPPORTEDANCH{\char"004D\char"034C} \SUPPORTEDANCH{\char"2C6E\char"034C} \SUPPORTEDANCH{\char"0050\char"034C} \SUPPORTEDANCH{\char"01A4\char"034C} \SUPPORTEDANCH{\char"0052\char"034C} \SUPPORTEDANCH{\char"2C64\char"034C} \SUPPORTEDANCH{\char"0054\char"034C} \SUPPORTEDANCH{\char"01AC\char"034C} \SUPPORTEDANCH{\char"01AE\char"034C} \SUPPORTEDANCH{\char"0059\char"034C} \SUPPORTEDANCH{\char"01B3\char"034C} \SUPPORTEDANCH{\char"0061\char"034C} \SUPPORTEDANCH{\char"1E9A\char"034C} \SUPPORTEDANCH{\char"0062\char"034C} \SUPPORTEDANCH{\char"0183\char"034C} \SUPPORTEDANCH{\char"0063\char"034C} \SUPPORTEDANCH{\char"0188\char"034C} \SUPPORTEDANCH{\char"0064\char"034C} \SUPPORTEDANCH{\char"018C\char"034C} \SUPPORTEDANCH{\char"0221\char"034C} \SUPPORTEDANCH{\char"0068\char"034C} \SUPPORTEDANCH{\char"2C68\char"034C} \SUPPORTEDANCH{\char"006B\char"034C} \SUPPORTEDANCH{\char"0199\char"034C} \SUPPORTEDANCH{\char"2C6A\char"034C} \SUPPORTEDANCH{\char"006C\char"034C} \SUPPORTEDANCH{\char"0234\char"034C} \SUPPORTEDANCH{\char"006E\char"034C} \SUPPORTEDANCH{\char"0235\char"034C} \SUPPORTEDANCH{\char"0070\char"034C} \SUPPORTEDANCH{\char"01A5\char"034C} \SUPPORTEDANCH{\char"0071\char"034C} \SUPPORTEDANCH{\char"02A0\char"034C} \SUPPORTEDANCH{\char"0074\char"034C} \SUPPORTEDANCH{\char"01AB\char"034C} \SUPPORTEDANCH{\char"01AD\char"034C} \SUPPORTEDANCH{\char"0236\char"034C} \SUPPORTEDANCH{\char"0075\char"034C} \SUPPORTEDANCH{\char"AB52\char"034C} \SUPPORTEDANCH{\char"0079\char"034C} \SUPPORTEDANCH{\char"01B4\char"034C} \SUPPORTEDANCH{\char"1EFF\char"034C}

\SUPPORTEDANCH{\char"0048\char"0350} \SUPPORTEDANCH{\char"2C67\char"0350} \SUPPORTEDANCH{\char"004B\char"0350} \SUPPORTEDANCH{\char"2C69\char"0350} \SUPPORTEDANCH{\char"004D\char"0350} \SUPPORTEDANCH{\char"2C6E\char"0350} \SUPPORTEDANCH{\char"0050\char"0350} \SUPPORTEDANCH{\char"01A4\char"0350} \SUPPORTEDANCH{\char"0052\char"0350} \SUPPORTEDANCH{\char"2C64\char"0350} \SUPPORTEDANCH{\char"0054\char"0350} \SUPPORTEDANCH{\char"01AC\char"0350} \SUPPORTEDANCH{\char"01AE\char"0350} \SUPPORTEDANCH{\char"0059\char"0350} \SUPPORTEDANCH{\char"01B3\char"0350} \SUPPORTEDANCH{\char"0061\char"0350} \SUPPORTEDANCH{\char"1E9A\char"0350} \SUPPORTEDANCH{\char"0062\char"0350} \SUPPORTEDANCH{\char"0183\char"0350} \SUPPORTEDANCH{\char"0063\char"0350} \SUPPORTEDANCH{\char"0188\char"0350} \SUPPORTEDANCH{\char"0064\char"0350} \SUPPORTEDANCH{\char"018C\char"0350} \SUPPORTEDANCH{\char"0221\char"0350} \SUPPORTEDANCH{\char"0068\char"0350} \SUPPORTEDANCH{\char"2C68\char"0350} \SUPPORTEDANCH{\char"006B\char"0350} \SUPPORTEDANCH{\char"0199\char"0350} \SUPPORTEDANCH{\char"2C6A\char"0350} \SUPPORTEDANCH{\char"006C\char"0350} \SUPPORTEDANCH{\char"0234\char"0350} \SUPPORTEDANCH{\char"006E\char"0350} \SUPPORTEDANCH{\char"0235\char"0350} \SUPPORTEDANCH{\char"0070\char"0350} \SUPPORTEDANCH{\char"01A5\char"0350} \SUPPORTEDANCH{\char"0071\char"0350} \SUPPORTEDANCH{\char"02A0\char"0350} \SUPPORTEDANCH{\char"0074\char"0350} \SUPPORTEDANCH{\char"01AB\char"0350} \SUPPORTEDANCH{\char"01AD\char"0350} \SUPPORTEDANCH{\char"0236\char"0350} \SUPPORTEDANCH{\char"0075\char"0350} \SUPPORTEDANCH{\char"AB52\char"0350} \SUPPORTEDANCH{\char"0079\char"0350} \SUPPORTEDANCH{\char"01B4\char"0350} \SUPPORTEDANCH{\char"1EFF\char"0350}

\SUPPORTEDANCH{\char"0048\char"0351} \SUPPORTEDANCH{\char"2C67\char"0351} \SUPPORTEDANCH{\char"004B\char"0351} \SUPPORTEDANCH{\char"2C69\char"0351} \SUPPORTEDANCH{\char"004D\char"0351} \SUPPORTEDANCH{\char"2C6E\char"0351} \SUPPORTEDANCH{\char"0050\char"0351} \SUPPORTEDANCH{\char"01A4\char"0351} \SUPPORTEDANCH{\char"0052\char"0351} \SUPPORTEDANCH{\char"2C64\char"0351} \SUPPORTEDANCH{\char"0054\char"0351} \SUPPORTEDANCH{\char"01AC\char"0351} \SUPPORTEDANCH{\char"01AE\char"0351} \SUPPORTEDANCH{\char"0059\char"0351} \SUPPORTEDANCH{\char"01B3\char"0351} \SUPPORTEDANCH{\char"0061\char"0351} \SUPPORTEDANCH{\char"1E9A\char"0351} \SUPPORTEDANCH{\char"0062\char"0351} \SUPPORTEDANCH{\char"0183\char"0351} \SUPPORTEDANCH{\char"0063\char"0351} \SUPPORTEDANCH{\char"0188\char"0351} \SUPPORTEDANCH{\char"0064\char"0351} \SUPPORTEDANCH{\char"018C\char"0351} \SUPPORTEDANCH{\char"0221\char"0351} \SUPPORTEDANCH{\char"0068\char"0351} \SUPPORTEDANCH{\char"2C68\char"0351} \SUPPORTEDANCH{\char"006B\char"0351} \SUPPORTEDANCH{\char"0199\char"0351} \SUPPORTEDANCH{\char"2C6A\char"0351} \SUPPORTEDANCH{\char"006C\char"0351} \SUPPORTEDANCH{\char"0234\char"0351} \SUPPORTEDANCH{\char"006E\char"0351} \SUPPORTEDANCH{\char"0235\char"0351} \SUPPORTEDANCH{\char"0070\char"0351} \SUPPORTEDANCH{\char"01A5\char"0351} \SUPPORTEDANCH{\char"0071\char"0351} \SUPPORTEDANCH{\char"02A0\char"0351} \SUPPORTEDANCH{\char"0074\char"0351} \SUPPORTEDANCH{\char"01AB\char"0351} \SUPPORTEDANCH{\char"01AD\char"0351} \SUPPORTEDANCH{\char"0236\char"0351} \SUPPORTEDANCH{\char"0075\char"0351} \SUPPORTEDANCH{\char"AB52\char"0351} \SUPPORTEDANCH{\char"0079\char"0351} \SUPPORTEDANCH{\char"01B4\char"0351} \SUPPORTEDANCH{\char"1EFF\char"0351}

\SUPPORTEDANCH{\char"0048\char"0352} \SUPPORTEDANCH{\char"2C67\char"0352} \SUPPORTEDANCH{\char"004B\char"0352} \SUPPORTEDANCH{\char"2C69\char"0352} \SUPPORTEDANCH{\char"004D\char"0352} \SUPPORTEDANCH{\char"2C6E\char"0352} \SUPPORTEDANCH{\char"0050\char"0352} \SUPPORTEDANCH{\char"01A4\char"0352} \SUPPORTEDANCH{\char"0052\char"0352} \SUPPORTEDANCH{\char"2C64\char"0352} \SUPPORTEDANCH{\char"0054\char"0352} \SUPPORTEDANCH{\char"01AC\char"0352} \SUPPORTEDANCH{\char"01AE\char"0352} \SUPPORTEDANCH{\char"0059\char"0352} \SUPPORTEDANCH{\char"01B3\char"0352} \SUPPORTEDANCH{\char"0061\char"0352} \SUPPORTEDANCH{\char"1E9A\char"0352} \SUPPORTEDANCH{\char"0062\char"0352} \SUPPORTEDANCH{\char"0183\char"0352} \SUPPORTEDANCH{\char"0063\char"0352} \SUPPORTEDANCH{\char"0188\char"0352} \SUPPORTEDANCH{\char"0064\char"0352} \SUPPORTEDANCH{\char"018C\char"0352} \SUPPORTEDANCH{\char"0221\char"0352} \SUPPORTEDANCH{\char"0068\char"0352} \SUPPORTEDANCH{\char"2C68\char"0352} \SUPPORTEDANCH{\char"006B\char"0352} \SUPPORTEDANCH{\char"0199\char"0352} \SUPPORTEDANCH{\char"2C6A\char"0352} \SUPPORTEDANCH{\char"006C\char"0352} \SUPPORTEDANCH{\char"0234\char"0352} \SUPPORTEDANCH{\char"006E\char"0352} \SUPPORTEDANCH{\char"0235\char"0352} \SUPPORTEDANCH{\char"0070\char"0352} \SUPPORTEDANCH{\char"01A5\char"0352} \SUPPORTEDANCH{\char"0071\char"0352} \SUPPORTEDANCH{\char"02A0\char"0352} \SUPPORTEDANCH{\char"0074\char"0352} \SUPPORTEDANCH{\char"01AB\char"0352} \SUPPORTEDANCH{\char"01AD\char"0352} \SUPPORTEDANCH{\char"0236\char"0352} \SUPPORTEDANCH{\char"0075\char"0352} \SUPPORTEDANCH{\char"AB52\char"0352} \SUPPORTEDANCH{\char"0079\char"0352} \SUPPORTEDANCH{\char"01B4\char"0352} \SUPPORTEDANCH{\char"1EFF\char"0352}

\SUPPORTEDANCH{\char"0048\char"0357} \SUPPORTEDANCH{\char"2C67\char"0357} \SUPPORTEDANCH{\char"004B\char"0357} \SUPPORTEDANCH{\char"2C69\char"0357} \SUPPORTEDANCH{\char"004D\char"0357} \SUPPORTEDANCH{\char"2C6E\char"0357} \SUPPORTEDANCH{\char"0050\char"0357} \SUPPORTEDANCH{\char"01A4\char"0357} \SUPPORTEDANCH{\char"0052\char"0357} \SUPPORTEDANCH{\char"2C64\char"0357} \SUPPORTEDANCH{\char"0054\char"0357} \SUPPORTEDANCH{\char"01AC\char"0357} \SUPPORTEDANCH{\char"01AE\char"0357} \SUPPORTEDANCH{\char"0059\char"0357} \SUPPORTEDANCH{\char"01B3\char"0357} \SUPPORTEDANCH{\char"0061\char"0357} \SUPPORTEDANCH{\char"1E9A\char"0357} \SUPPORTEDANCH{\char"0062\char"0357} \SUPPORTEDANCH{\char"0183\char"0357} \SUPPORTEDANCH{\char"0063\char"0357} \SUPPORTEDANCH{\char"0188\char"0357} \SUPPORTEDANCH{\char"0064\char"0357} \SUPPORTEDANCH{\char"018C\char"0357} \SUPPORTEDANCH{\char"0221\char"0357} \SUPPORTEDANCH{\char"0068\char"0357} \SUPPORTEDANCH{\char"2C68\char"0357} \SUPPORTEDANCH{\char"006B\char"0357} \SUPPORTEDANCH{\char"0199\char"0357} \SUPPORTEDANCH{\char"2C6A\char"0357} \SUPPORTEDANCH{\char"006C\char"0357} \SUPPORTEDANCH{\char"0234\char"0357} \SUPPORTEDANCH{\char"006E\char"0357} \SUPPORTEDANCH{\char"0235\char"0357} \SUPPORTEDANCH{\char"0070\char"0357} \SUPPORTEDANCH{\char"01A5\char"0357} \SUPPORTEDANCH{\char"0071\char"0357} \SUPPORTEDANCH{\char"02A0\char"0357} \SUPPORTEDANCH{\char"0074\char"0357} \SUPPORTEDANCH{\char"01AB\char"0357} \SUPPORTEDANCH{\char"01AD\char"0357} \SUPPORTEDANCH{\char"0236\char"0357} \SUPPORTEDANCH{\char"0075\char"0357} \SUPPORTEDANCH{\char"AB52\char"0357} \SUPPORTEDANCH{\char"0079\char"0357} \SUPPORTEDANCH{\char"01B4\char"0357} \SUPPORTEDANCH{\char"1EFF\char"0357}

\SUPPORTEDANCH{\char"0048\char"035B} \SUPPORTEDANCH{\char"2C67\char"035B} \SUPPORTEDANCH{\char"004B\char"035B} \SUPPORTEDANCH{\char"2C69\char"035B} \SUPPORTEDANCH{\char"004D\char"035B} \SUPPORTEDANCH{\char"2C6E\char"035B} \SUPPORTEDANCH{\char"0050\char"035B} \SUPPORTEDANCH{\char"01A4\char"035B} \SUPPORTEDANCH{\char"0052\char"035B} \SUPPORTEDANCH{\char"2C64\char"035B} \SUPPORTEDANCH{\char"0054\char"035B} \SUPPORTEDANCH{\char"01AC\char"035B} \SUPPORTEDANCH{\char"01AE\char"035B} \SUPPORTEDANCH{\char"0059\char"035B} \SUPPORTEDANCH{\char"01B3\char"035B} \SUPPORTEDANCH{\char"0061\char"035B} \SUPPORTEDANCH{\char"1E9A\char"035B} \SUPPORTEDANCH{\char"0062\char"035B} \SUPPORTEDANCH{\char"0183\char"035B} \SUPPORTEDANCH{\char"0063\char"035B} \SUPPORTEDANCH{\char"0188\char"035B} \SUPPORTEDANCH{\char"0064\char"035B} \SUPPORTEDANCH{\char"018C\char"035B} \SUPPORTEDANCH{\char"0221\char"035B} \SUPPORTEDANCH{\char"0068\char"035B} \SUPPORTEDANCH{\char"2C68\char"035B} \SUPPORTEDANCH{\char"006B\char"035B} \SUPPORTEDANCH{\char"0199\char"035B} \SUPPORTEDANCH{\char"2C6A\char"035B} \SUPPORTEDANCH{\char"006C\char"035B} \SUPPORTEDANCH{\char"0234\char"035B} \SUPPORTEDANCH{\char"006E\char"035B} \SUPPORTEDANCH{\char"0235\char"035B} \SUPPORTEDANCH{\char"0070\char"035B} \SUPPORTEDANCH{\char"01A5\char"035B} \SUPPORTEDANCH{\char"0071\char"035B} \SUPPORTEDANCH{\char"02A0\char"035B} \SUPPORTEDANCH{\char"0074\char"035B} \SUPPORTEDANCH{\char"01AB\char"035B} \SUPPORTEDANCH{\char"01AD\char"035B} \SUPPORTEDANCH{\char"0236\char"035B} \SUPPORTEDANCH{\char"0075\char"035B} \SUPPORTEDANCH{\char"AB52\char"035B} \SUPPORTEDANCH{\char"0079\char"035B} \SUPPORTEDANCH{\char"01B4\char"035B} \SUPPORTEDANCH{\char"1EFF\char"035B}

}

\newpage
\subsection*{Italic, Group 4, above, Above}

{\itshape
% grid. columns are bases, rows are marks.
\SUPPORTEDANCH{\char"0048\char"0300} \SUPPORTEDFALL{\char"2C67\char"0300} \SUPPORTEDANCH{\char"004B\char"0300} \SUPPORTEDANCH{\char"2C69\char"0300} \SUPPORTEDANCH{\char"004D\char"0300} \MISSINGGLYPH{\char"2C6E\char"0300} \SUPPORTEDANCH{\char"0050\char"0300} \SUPPORTEDFALL{\char"01A4\char"0300} \SUPPORTEDANCH{\char"0052\char"0300} \SUPPORTEDFALL{\char"2C64\char"0300} \SUPPORTEDANCH{\char"0054\char"0300} \SUPPORTEDFALL{\char"01AC\char"0300} \SUPPORTEDFALL{\char"01AE\char"0300} \SUPPORTEDPREC{\char"0059\char"0300} \SUPPORTEDFALL{\char"01B3\char"0300} \SUPPORTEDPREC{\char"0061\char"0300} \SUPPORTEDFALL{\char"1E9A\char"0300} \SUPPORTEDFALL{\char"0062\char"0300} \SUPPORTEDFALL{\char"0183\char"0300} \SUPPORTEDANCH{\char"0063\char"0300} \SUPPORTEDFALL{\char"0188\char"0300} \SUPPORTEDANCH{\char"0064\char"0300} \SUPPORTEDANCH{\char"018C\char"0300} \SUPPORTEDFALL{\char"0221\char"0300} \SUPPORTEDANCH{\char"0068\char"0300} \SUPPORTEDANCH{\char"2C68\char"0300} \SUPPORTEDANCH{\char"006B\char"0300} \SUPPORTEDFALL{\char"0199\char"0300} \SUPPORTEDANCH{\char"2C6A\char"0300} \SUPPORTEDANCH{\char"006C\char"0300} \SUPPORTEDFALL{\char"0234\char"0300} \SUPPORTEDPREC{\char"006E\char"0300} \SUPPORTEDANCH{\char"0235\char"0300} \SUPPORTEDANCH{\char"0070\char"0300} \SUPPORTEDFALL{\char"01A5\char"0300} \SUPPORTEDANCH{\char"0071\char"0300} \SUPPORTEDANCH{\char"02A0\char"0300} \SUPPORTEDFALL{\char"0074\char"0300} \SUPPORTEDFALL{\char"01AB\char"0300} \SUPPORTEDFALL{\char"01AD\char"0300} \SUPPORTEDFALL{\char"0236\char"0300} \SUPPORTEDPREC{\char"0075\char"0300} \MISSINGGLYPH{\char"AB52\char"0300} \SUPPORTEDPREC{\char"0079\char"0300} \SUPPORTEDANCH{\char"01B4\char"0300} \MISSINGGLYPH{\char"1EFF\char"0300}

\SUPPORTEDANCH{\char"0048\char"0301} \SUPPORTEDFALL{\char"2C67\char"0301} \SUPPORTEDPREC{\char"004B\char"0301} \SUPPORTEDANCH{\char"2C69\char"0301} \SUPPORTEDPREC{\char"004D\char"0301} \MISSINGGLYPH{\char"2C6E\char"0301} \SUPPORTEDPREC{\char"0050\char"0301} \SUPPORTEDFALL{\char"01A4\char"0301} \SUPPORTEDPREC{\char"0052\char"0301} \SUPPORTEDFALL{\char"2C64\char"0301} \SUPPORTEDANCH{\char"0054\char"0301} \SUPPORTEDFALL{\char"01AC\char"0301} \SUPPORTEDFALL{\char"01AE\char"0301} \SUPPORTEDPREC{\char"0059\char"0301} \SUPPORTEDFALL{\char"01B3\char"0301} \SUPPORTEDPREC{\char"0061\char"0301} \SUPPORTEDFALL{\char"1E9A\char"0301} \SUPPORTEDFALL{\char"0062\char"0301} \SUPPORTEDFALL{\char"0183\char"0301} \SUPPORTEDPREC{\char"0063\char"0301} \SUPPORTEDFALL{\char"0188\char"0301} \SUPPORTEDANCH{\char"0064\char"0301} \SUPPORTEDANCH{\char"018C\char"0301} \SUPPORTEDFALL{\char"0221\char"0301} \SUPPORTEDANCH{\char"0068\char"0301} \SUPPORTEDANCH{\char"2C68\char"0301} \SUPPORTEDPREC{\char"006B\char"0301} \SUPPORTEDFALL{\char"0199\char"0301} \SUPPORTEDANCH{\char"2C6A\char"0301} \SUPPORTEDPREC{\char"006C\char"0301} \SUPPORTEDFALL{\char"0234\char"0301} \SUPPORTEDPREC{\char"006E\char"0301} \SUPPORTEDANCH{\char"0235\char"0301} \SUPPORTEDPREC{\char"0070\char"0301} \SUPPORTEDFALL{\char"01A5\char"0301} \SUPPORTEDANCH{\char"0071\char"0301} \SUPPORTEDANCH{\char"02A0\char"0301} \SUPPORTEDFALL{\char"0074\char"0301} \SUPPORTEDFALL{\char"01AB\char"0301} \SUPPORTEDFALL{\char"01AD\char"0301} \SUPPORTEDFALL{\char"0236\char"0301} \SUPPORTEDPREC{\char"0075\char"0301} \MISSINGGLYPH{\char"AB52\char"0301} \SUPPORTEDPREC{\char"0079\char"0301} \SUPPORTEDANCH{\char"01B4\char"0301} \MISSINGGLYPH{\char"1EFF\char"0301}

\SUPPORTEDPREC{\char"0048\char"0302} \SUPPORTEDFALL{\char"2C67\char"0302} \SUPPORTEDANCH{\char"004B\char"0302} \SUPPORTEDANCH{\char"2C69\char"0302} \SUPPORTEDANCH{\char"004D\char"0302} \MISSINGGLYPH{\char"2C6E\char"0302} \SUPPORTEDANCH{\char"0050\char"0302} \SUPPORTEDFALL{\char"01A4\char"0302} \SUPPORTEDANCH{\char"0052\char"0302} \SUPPORTEDFALL{\char"2C64\char"0302} \SUPPORTEDANCH{\char"0054\char"0302} \SUPPORTEDFALL{\char"01AC\char"0302} \SUPPORTEDFALL{\char"01AE\char"0302} \SUPPORTEDPREC{\char"0059\char"0302} \SUPPORTEDFALL{\char"01B3\char"0302} \SUPPORTEDPREC{\char"0061\char"0302} \SUPPORTEDFALL{\char"1E9A\char"0302} \SUPPORTEDFALL{\char"0062\char"0302} \SUPPORTEDFALL{\char"0183\char"0302} \SUPPORTEDPREC{\char"0063\char"0302} \SUPPORTEDFALL{\char"0188\char"0302} \SUPPORTEDANCH{\char"0064\char"0302} \SUPPORTEDANCH{\char"018C\char"0302} \SUPPORTEDFALL{\char"0221\char"0302} \SUPPORTEDPREC{\char"0068\char"0302} \SUPPORTEDANCH{\char"2C68\char"0302} \SUPPORTEDANCH{\char"006B\char"0302} \SUPPORTEDFALL{\char"0199\char"0302} \SUPPORTEDANCH{\char"2C6A\char"0302} \SUPPORTEDANCH{\char"006C\char"0302} \SUPPORTEDFALL{\char"0234\char"0302} \SUPPORTEDANCH{\char"006E\char"0302} \SUPPORTEDANCH{\char"0235\char"0302} \SUPPORTEDANCH{\char"0070\char"0302} \SUPPORTEDFALL{\char"01A5\char"0302} \SUPPORTEDANCH{\char"0071\char"0302} \SUPPORTEDANCH{\char"02A0\char"0302} \SUPPORTEDFALL{\char"0074\char"0302} \SUPPORTEDFALL{\char"01AB\char"0302} \SUPPORTEDFALL{\char"01AD\char"0302} \SUPPORTEDFALL{\char"0236\char"0302} \SUPPORTEDPREC{\char"0075\char"0302} \MISSINGGLYPH{\char"AB52\char"0302} \SUPPORTEDPREC{\char"0079\char"0302} \SUPPORTEDANCH{\char"01B4\char"0302} \MISSINGGLYPH{\char"1EFF\char"0302}

\SUPPORTEDANCH{\char"0048\char"0303} \SUPPORTEDFALL{\char"2C67\char"0303} \SUPPORTEDANCH{\char"004B\char"0303} \SUPPORTEDANCH{\char"2C69\char"0303} \SUPPORTEDANCH{\char"004D\char"0303} \MISSINGGLYPH{\char"2C6E\char"0303} \SUPPORTEDANCH{\char"0050\char"0303} \SUPPORTEDFALL{\char"01A4\char"0303} \SUPPORTEDANCH{\char"0052\char"0303} \SUPPORTEDFALL{\char"2C64\char"0303} \SUPPORTEDANCH{\char"0054\char"0303} \SUPPORTEDFALL{\char"01AC\char"0303} \SUPPORTEDFALL{\char"01AE\char"0303} \SUPPORTEDPREC{\char"0059\char"0303} \SUPPORTEDFALL{\char"01B3\char"0303} \SUPPORTEDPREC{\char"0061\char"0303} \SUPPORTEDFALL{\char"1E9A\char"0303} \SUPPORTEDFALL{\char"0062\char"0303} \SUPPORTEDFALL{\char"0183\char"0303} \SUPPORTEDANCH{\char"0063\char"0303} \SUPPORTEDFALL{\char"0188\char"0303} \SUPPORTEDANCH{\char"0064\char"0303} \SUPPORTEDANCH{\char"018C\char"0303} \SUPPORTEDFALL{\char"0221\char"0303} \SUPPORTEDANCH{\char"0068\char"0303} \SUPPORTEDANCH{\char"2C68\char"0303} \SUPPORTEDANCH{\char"006B\char"0303} \SUPPORTEDFALL{\char"0199\char"0303} \SUPPORTEDANCH{\char"2C6A\char"0303} \SUPPORTEDANCH{\char"006C\char"0303} \SUPPORTEDFALL{\char"0234\char"0303} \SUPPORTEDPREC{\char"006E\char"0303} \SUPPORTEDANCH{\char"0235\char"0303} \SUPPORTEDANCH{\char"0070\char"0303} \SUPPORTEDFALL{\char"01A5\char"0303} \SUPPORTEDANCH{\char"0071\char"0303} \SUPPORTEDANCH{\char"02A0\char"0303} \SUPPORTEDFALL{\char"0074\char"0303} \SUPPORTEDFALL{\char"01AB\char"0303} \SUPPORTEDFALL{\char"01AD\char"0303} \SUPPORTEDFALL{\char"0236\char"0303} \SUPPORTEDPREC{\char"0075\char"0303} \MISSINGGLYPH{\char"AB52\char"0303} \SUPPORTEDPREC{\char"0079\char"0303} \SUPPORTEDANCH{\char"01B4\char"0303} \MISSINGGLYPH{\char"1EFF\char"0303}

\SUPPORTEDANCH{\char"0048\char"0304} \SUPPORTEDFALL{\char"2C67\char"0304} \SUPPORTEDANCH{\char"004B\char"0304} \SUPPORTEDANCH{\char"2C69\char"0304} \SUPPORTEDANCH{\char"004D\char"0304} \MISSINGGLYPH{\char"2C6E\char"0304} \SUPPORTEDANCH{\char"0050\char"0304} \SUPPORTEDFALL{\char"01A4\char"0304} \SUPPORTEDANCH{\char"0052\char"0304} \SUPPORTEDFALL{\char"2C64\char"0304} \SUPPORTEDANCH{\char"0054\char"0304} \SUPPORTEDFALL{\char"01AC\char"0304} \SUPPORTEDFALL{\char"01AE\char"0304} \SUPPORTEDPREC{\char"0059\char"0304} \SUPPORTEDFALL{\char"01B3\char"0304} \SUPPORTEDPREC{\char"0061\char"0304} \SUPPORTEDFALL{\char"1E9A\char"0304} \SUPPORTEDFALL{\char"0062\char"0304} \SUPPORTEDFALL{\char"0183\char"0304} \SUPPORTEDANCH{\char"0063\char"0304} \SUPPORTEDFALL{\char"0188\char"0304} \SUPPORTEDANCH{\char"0064\char"0304} \SUPPORTEDANCH{\char"018C\char"0304} \SUPPORTEDFALL{\char"0221\char"0304} \SUPPORTEDANCH{\char"0068\char"0304} \SUPPORTEDANCH{\char"2C68\char"0304} \SUPPORTEDANCH{\char"006B\char"0304} \SUPPORTEDFALL{\char"0199\char"0304} \SUPPORTEDANCH{\char"2C6A\char"0304} \SUPPORTEDANCH{\char"006C\char"0304} \SUPPORTEDFALL{\char"0234\char"0304} \SUPPORTEDANCH{\char"006E\char"0304} \SUPPORTEDANCH{\char"0235\char"0304} \SUPPORTEDANCH{\char"0070\char"0304} \SUPPORTEDFALL{\char"01A5\char"0304} \SUPPORTEDANCH{\char"0071\char"0304} \SUPPORTEDANCH{\char"02A0\char"0304} \SUPPORTEDFALL{\char"0074\char"0304} \SUPPORTEDFALL{\char"01AB\char"0304} \SUPPORTEDFALL{\char"01AD\char"0304} \SUPPORTEDFALL{\char"0236\char"0304} \SUPPORTEDPREC{\char"0075\char"0304} \MISSINGGLYPH{\char"AB52\char"0304} \SUPPORTEDPREC{\char"0079\char"0304} \SUPPORTEDANCH{\char"01B4\char"0304} \MISSINGGLYPH{\char"1EFF\char"0304}

\SUPPORTEDANCH{\char"0048\char"0306} \SUPPORTEDFALL{\char"2C67\char"0306} \SUPPORTEDANCH{\char"004B\char"0306} \SUPPORTEDANCH{\char"2C69\char"0306} \SUPPORTEDANCH{\char"004D\char"0306} \MISSINGGLYPH{\char"2C6E\char"0306} \SUPPORTEDANCH{\char"0050\char"0306} \SUPPORTEDFALL{\char"01A4\char"0306} \SUPPORTEDANCH{\char"0052\char"0306} \SUPPORTEDFALL{\char"2C64\char"0306} \SUPPORTEDANCH{\char"0054\char"0306} \SUPPORTEDFALL{\char"01AC\char"0306} \SUPPORTEDFALL{\char"01AE\char"0306} \SUPPORTEDANCH{\char"0059\char"0306} \SUPPORTEDFALL{\char"01B3\char"0306} \SUPPORTEDPREC{\char"0061\char"0306} \SUPPORTEDFALL{\char"1E9A\char"0306} \SUPPORTEDFALL{\char"0062\char"0306} \SUPPORTEDFALL{\char"0183\char"0306} \SUPPORTEDANCH{\char"0063\char"0306} \SUPPORTEDFALL{\char"0188\char"0306} \SUPPORTEDANCH{\char"0064\char"0306} \SUPPORTEDANCH{\char"018C\char"0306} \SUPPORTEDFALL{\char"0221\char"0306} \SUPPORTEDANCH{\char"0068\char"0306} \SUPPORTEDANCH{\char"2C68\char"0306} \SUPPORTEDANCH{\char"006B\char"0306} \SUPPORTEDFALL{\char"0199\char"0306} \SUPPORTEDANCH{\char"2C6A\char"0306} \SUPPORTEDANCH{\char"006C\char"0306} \SUPPORTEDFALL{\char"0234\char"0306} \SUPPORTEDANCH{\char"006E\char"0306} \SUPPORTEDANCH{\char"0235\char"0306} \SUPPORTEDANCH{\char"0070\char"0306} \SUPPORTEDFALL{\char"01A5\char"0306} \SUPPORTEDANCH{\char"0071\char"0306} \SUPPORTEDANCH{\char"02A0\char"0306} \SUPPORTEDFALL{\char"0074\char"0306} \SUPPORTEDFALL{\char"01AB\char"0306} \SUPPORTEDFALL{\char"01AD\char"0306} \SUPPORTEDFALL{\char"0236\char"0306} \SUPPORTEDPREC{\char"0075\char"0306} \MISSINGGLYPH{\char"AB52\char"0306} \SUPPORTEDANCH{\char"0079\char"0306} \SUPPORTEDANCH{\char"01B4\char"0306} \MISSINGGLYPH{\char"1EFF\char"0306}

\SUPPORTEDPREC{\char"0048\char"0307} \SUPPORTEDFALL{\char"2C67\char"0307} \SUPPORTEDANCH{\char"004B\char"0307} \SUPPORTEDANCH{\char"2C69\char"0307} \SUPPORTEDPREC{\char"004D\char"0307} \MISSINGGLYPH{\char"2C6E\char"0307} \SUPPORTEDPREC{\char"0050\char"0307} \SUPPORTEDFALL{\char"01A4\char"0307} \SUPPORTEDPREC{\char"0052\char"0307} \SUPPORTEDFALL{\char"2C64\char"0307} \SUPPORTEDPREC{\char"0054\char"0307} \SUPPORTEDFALL{\char"01AC\char"0307} \SUPPORTEDFALL{\char"01AE\char"0307} \SUPPORTEDPREC{\char"0059\char"0307} \SUPPORTEDFALL{\char"01B3\char"0307} \SUPPORTEDPREC{\char"0061\char"0307} \SUPPORTEDFALL{\char"1E9A\char"0307} \SUPPORTEDPREC{\char"0062\char"0307} \SUPPORTEDFALL{\char"0183\char"0307} \SUPPORTEDPREC{\char"0063\char"0307} \SUPPORTEDFALL{\char"0188\char"0307} \SUPPORTEDPREC{\char"0064\char"0307} \SUPPORTEDANCH{\char"018C\char"0307} \SUPPORTEDFALL{\char"0221\char"0307} \SUPPORTEDPREC{\char"0068\char"0307} \SUPPORTEDANCH{\char"2C68\char"0307} \SUPPORTEDANCH{\char"006B\char"0307} \SUPPORTEDFALL{\char"0199\char"0307} \SUPPORTEDANCH{\char"2C6A\char"0307} \SUPPORTEDANCH{\char"006C\char"0307} \SUPPORTEDFALL{\char"0234\char"0307} \SUPPORTEDPREC{\char"006E\char"0307} \SUPPORTEDANCH{\char"0235\char"0307} \SUPPORTEDPREC{\char"0070\char"0307} \SUPPORTEDFALL{\char"01A5\char"0307} \SUPPORTEDANCH{\char"0071\char"0307} \SUPPORTEDANCH{\char"02A0\char"0307} \SUPPORTEDPREC{\char"0074\char"0307} \SUPPORTEDFALL{\char"01AB\char"0307} \SUPPORTEDFALL{\char"01AD\char"0307} \SUPPORTEDFALL{\char"0236\char"0307} \SUPPORTEDANCH{\char"0075\char"0307} \MISSINGGLYPH{\char"AB52\char"0307} \SUPPORTEDPREC{\char"0079\char"0307} \SUPPORTEDANCH{\char"01B4\char"0307} \MISSINGGLYPH{\char"1EFF\char"0307}

\SUPPORTEDPREC{\char"0048\char"0308} \SUPPORTEDFALL{\char"2C67\char"0308} \SUPPORTEDANCH{\char"004B\char"0308} \SUPPORTEDANCH{\char"2C69\char"0308} \SUPPORTEDANCH{\char"004D\char"0308} \MISSINGGLYPH{\char"2C6E\char"0308} \SUPPORTEDANCH{\char"0050\char"0308} \SUPPORTEDFALL{\char"01A4\char"0308} \SUPPORTEDANCH{\char"0052\char"0308} \SUPPORTEDFALL{\char"2C64\char"0308} \SUPPORTEDANCH{\char"0054\char"0308} \SUPPORTEDFALL{\char"01AC\char"0308} \SUPPORTEDFALL{\char"01AE\char"0308} \SUPPORTEDPREC{\char"0059\char"0308} \SUPPORTEDFALL{\char"01B3\char"0308} \SUPPORTEDPREC{\char"0061\char"0308} \SUPPORTEDFALL{\char"1E9A\char"0308} \SUPPORTEDFALL{\char"0062\char"0308} \SUPPORTEDFALL{\char"0183\char"0308} \SUPPORTEDANCH{\char"0063\char"0308} \SUPPORTEDFALL{\char"0188\char"0308} \SUPPORTEDANCH{\char"0064\char"0308} \SUPPORTEDANCH{\char"018C\char"0308} \SUPPORTEDFALL{\char"0221\char"0308} \SUPPORTEDPREC{\char"0068\char"0308} \SUPPORTEDANCH{\char"2C68\char"0308} \SUPPORTEDANCH{\char"006B\char"0308} \SUPPORTEDFALL{\char"0199\char"0308} \SUPPORTEDANCH{\char"2C6A\char"0308} \SUPPORTEDANCH{\char"006C\char"0308} \SUPPORTEDFALL{\char"0234\char"0308} \SUPPORTEDANCH{\char"006E\char"0308} \SUPPORTEDANCH{\char"0235\char"0308} \SUPPORTEDANCH{\char"0070\char"0308} \SUPPORTEDFALL{\char"01A5\char"0308} \SUPPORTEDANCH{\char"0071\char"0308} \SUPPORTEDANCH{\char"02A0\char"0308} \SUPPORTEDPREC{\char"0074\char"0308} \SUPPORTEDFALL{\char"01AB\char"0308} \SUPPORTEDFALL{\char"01AD\char"0308} \SUPPORTEDFALL{\char"0236\char"0308} \SUPPORTEDPREC{\char"0075\char"0308} \MISSINGGLYPH{\char"AB52\char"0308} \SUPPORTEDPREC{\char"0079\char"0308} \SUPPORTEDANCH{\char"01B4\char"0308} \MISSINGGLYPH{\char"1EFF\char"0308}

\SUPPORTEDANCH{\char"0048\char"0309} \SUPPORTEDFALL{\char"2C67\char"0309} \SUPPORTEDANCH{\char"004B\char"0309} \SUPPORTEDANCH{\char"2C69\char"0309} \SUPPORTEDANCH{\char"004D\char"0309} \MISSINGGLYPH{\char"2C6E\char"0309} \SUPPORTEDANCH{\char"0050\char"0309} \SUPPORTEDFALL{\char"01A4\char"0309} \SUPPORTEDANCH{\char"0052\char"0309} \SUPPORTEDFALL{\char"2C64\char"0309} \SUPPORTEDANCH{\char"0054\char"0309} \SUPPORTEDFALL{\char"01AC\char"0309} \SUPPORTEDFALL{\char"01AE\char"0309} \SUPPORTEDPREC{\char"0059\char"0309} \SUPPORTEDFALL{\char"01B3\char"0309} \SUPPORTEDPREC{\char"0061\char"0309} \SUPPORTEDFALL{\char"1E9A\char"0309} \SUPPORTEDFALL{\char"0062\char"0309} \SUPPORTEDFALL{\char"0183\char"0309} \SUPPORTEDANCH{\char"0063\char"0309} \SUPPORTEDFALL{\char"0188\char"0309} \SUPPORTEDANCH{\char"0064\char"0309} \SUPPORTEDANCH{\char"018C\char"0309} \SUPPORTEDFALL{\char"0221\char"0309} \SUPPORTEDANCH{\char"0068\char"0309} \SUPPORTEDANCH{\char"2C68\char"0309} \SUPPORTEDANCH{\char"006B\char"0309} \SUPPORTEDFALL{\char"0199\char"0309} \SUPPORTEDANCH{\char"2C6A\char"0309} \SUPPORTEDANCH{\char"006C\char"0309} \SUPPORTEDFALL{\char"0234\char"0309} \SUPPORTEDANCH{\char"006E\char"0309} \SUPPORTEDANCH{\char"0235\char"0309} \SUPPORTEDANCH{\char"0070\char"0309} \SUPPORTEDFALL{\char"01A5\char"0309} \SUPPORTEDANCH{\char"0071\char"0309} \SUPPORTEDANCH{\char"02A0\char"0309} \SUPPORTEDFALL{\char"0074\char"0309} \SUPPORTEDFALL{\char"01AB\char"0309} \SUPPORTEDFALL{\char"01AD\char"0309} \SUPPORTEDFALL{\char"0236\char"0309} \SUPPORTEDPREC{\char"0075\char"0309} \MISSINGGLYPH{\char"AB52\char"0309} \SUPPORTEDPREC{\char"0079\char"0309} \SUPPORTEDANCH{\char"01B4\char"0309} \MISSINGGLYPH{\char"1EFF\char"0309}

\SUPPORTEDANCH{\char"0048\char"030A} \SUPPORTEDFALL{\char"2C67\char"030A} \SUPPORTEDANCH{\char"004B\char"030A} \SUPPORTEDANCH{\char"2C69\char"030A} \SUPPORTEDANCH{\char"004D\char"030A} \MISSINGGLYPH{\char"2C6E\char"030A} \SUPPORTEDANCH{\char"0050\char"030A} \SUPPORTEDFALL{\char"01A4\char"030A} \SUPPORTEDANCH{\char"0052\char"030A} \SUPPORTEDFALL{\char"2C64\char"030A} \SUPPORTEDANCH{\char"0054\char"030A} \SUPPORTEDFALL{\char"01AC\char"030A} \SUPPORTEDFALL{\char"01AE\char"030A} \SUPPORTEDANCH{\char"0059\char"030A} \SUPPORTEDFALL{\char"01B3\char"030A} \SUPPORTEDPREC{\char"0061\char"030A} \SUPPORTEDFALL{\char"1E9A\char"030A} \SUPPORTEDFALL{\char"0062\char"030A} \SUPPORTEDFALL{\char"0183\char"030A} \SUPPORTEDANCH{\char"0063\char"030A} \SUPPORTEDFALL{\char"0188\char"030A} \SUPPORTEDANCH{\char"0064\char"030A} \SUPPORTEDANCH{\char"018C\char"030A} \SUPPORTEDFALL{\char"0221\char"030A} \SUPPORTEDANCH{\char"0068\char"030A} \SUPPORTEDANCH{\char"2C68\char"030A} \SUPPORTEDANCH{\char"006B\char"030A} \SUPPORTEDFALL{\char"0199\char"030A} \SUPPORTEDANCH{\char"2C6A\char"030A} \SUPPORTEDANCH{\char"006C\char"030A} \SUPPORTEDFALL{\char"0234\char"030A} \SUPPORTEDANCH{\char"006E\char"030A} \SUPPORTEDANCH{\char"0235\char"030A} \SUPPORTEDANCH{\char"0070\char"030A} \SUPPORTEDFALL{\char"01A5\char"030A} \SUPPORTEDANCH{\char"0071\char"030A} \SUPPORTEDANCH{\char"02A0\char"030A} \SUPPORTEDFALL{\char"0074\char"030A} \SUPPORTEDFALL{\char"01AB\char"030A} \SUPPORTEDFALL{\char"01AD\char"030A} \SUPPORTEDFALL{\char"0236\char"030A} \SUPPORTEDPREC{\char"0075\char"030A} \MISSINGGLYPH{\char"AB52\char"030A} \SUPPORTEDPREC{\char"0079\char"030A} \SUPPORTEDANCH{\char"01B4\char"030A} \MISSINGGLYPH{\char"1EFF\char"030A}

\SUPPORTEDANCH{\char"0048\char"030B} \SUPPORTEDFALL{\char"2C67\char"030B} \SUPPORTEDANCH{\char"004B\char"030B} \SUPPORTEDANCH{\char"2C69\char"030B} \SUPPORTEDANCH{\char"004D\char"030B} \MISSINGGLYPH{\char"2C6E\char"030B} \SUPPORTEDANCH{\char"0050\char"030B} \SUPPORTEDFALL{\char"01A4\char"030B} \SUPPORTEDANCH{\char"0052\char"030B} \SUPPORTEDFALL{\char"2C64\char"030B} \SUPPORTEDANCH{\char"0054\char"030B} \SUPPORTEDFALL{\char"01AC\char"030B} \SUPPORTEDFALL{\char"01AE\char"030B} \SUPPORTEDANCH{\char"0059\char"030B} \SUPPORTEDFALL{\char"01B3\char"030B} \SUPPORTEDANCH{\char"0061\char"030B} \SUPPORTEDFALL{\char"1E9A\char"030B} \SUPPORTEDFALL{\char"0062\char"030B} \SUPPORTEDFALL{\char"0183\char"030B} \SUPPORTEDANCH{\char"0063\char"030B} \SUPPORTEDFALL{\char"0188\char"030B} \SUPPORTEDANCH{\char"0064\char"030B} \SUPPORTEDANCH{\char"018C\char"030B} \SUPPORTEDFALL{\char"0221\char"030B} \SUPPORTEDANCH{\char"0068\char"030B} \SUPPORTEDANCH{\char"2C68\char"030B} \SUPPORTEDANCH{\char"006B\char"030B} \SUPPORTEDFALL{\char"0199\char"030B} \SUPPORTEDANCH{\char"2C6A\char"030B} \SUPPORTEDANCH{\char"006C\char"030B} \SUPPORTEDFALL{\char"0234\char"030B} \SUPPORTEDANCH{\char"006E\char"030B} \SUPPORTEDANCH{\char"0235\char"030B} \SUPPORTEDANCH{\char"0070\char"030B} \SUPPORTEDFALL{\char"01A5\char"030B} \SUPPORTEDANCH{\char"0071\char"030B} \SUPPORTEDANCH{\char"02A0\char"030B} \SUPPORTEDFALL{\char"0074\char"030B} \SUPPORTEDFALL{\char"01AB\char"030B} \SUPPORTEDFALL{\char"01AD\char"030B} \SUPPORTEDFALL{\char"0236\char"030B} \SUPPORTEDPREC{\char"0075\char"030B} \MISSINGGLYPH{\char"AB52\char"030B} \SUPPORTEDANCH{\char"0079\char"030B} \SUPPORTEDANCH{\char"01B4\char"030B} \MISSINGGLYPH{\char"1EFF\char"030B}

\SUPPORTEDPREC{\char"0048\char"030C} \SUPPORTEDFALL{\char"2C67\char"030C} \SUPPORTEDPREC{\char"004B\char"030C} \SUPPORTEDANCH{\char"2C69\char"030C} \SUPPORTEDANCH{\char"004D\char"030C} \MISSINGGLYPH{\char"2C6E\char"030C} \SUPPORTEDANCH{\char"0050\char"030C} \SUPPORTEDFALL{\char"01A4\char"030C} \SUPPORTEDPREC{\char"0052\char"030C} \SUPPORTEDFALL{\char"2C64\char"030C} \SUPPORTEDPREC{\char"0054\char"030C} \SUPPORTEDFALL{\char"01AC\char"030C} \SUPPORTEDFALL{\char"01AE\char"030C} \SUPPORTEDANCH{\char"0059\char"030C} \SUPPORTEDFALL{\char"01B3\char"030C} \SUPPORTEDPREC{\char"0061\char"030C} \SUPPORTEDFALL{\char"1E9A\char"030C} \SUPPORTEDFALL{\char"0062\char"030C} \SUPPORTEDFALL{\char"0183\char"030C} \SUPPORTEDPREC{\char"0063\char"030C} \SUPPORTEDFALL{\char"0188\char"030C} \SUPPORTEDPREC{\char"0064\char"030C} \SUPPORTEDANCH{\char"018C\char"030C} \SUPPORTEDFALL{\char"0221\char"030C} \SUPPORTEDPREC{\char"0068\char"030C} \SUPPORTEDANCH{\char"2C68\char"030C} \SUPPORTEDPREC{\char"006B\char"030C} \SUPPORTEDFALL{\char"0199\char"030C} \SUPPORTEDANCH{\char"2C6A\char"030C} \SUPPORTEDPREC{\char"006C\char"030C} \SUPPORTEDFALL{\char"0234\char"030C} \SUPPORTEDPREC{\char"006E\char"030C} \SUPPORTEDANCH{\char"0235\char"030C} \SUPPORTEDANCH{\char"0070\char"030C} \SUPPORTEDFALL{\char"01A5\char"030C} \SUPPORTEDANCH{\char"0071\char"030C} \SUPPORTEDANCH{\char"02A0\char"030C} \SUPPORTEDPREC{\char"0074\char"030C} \SUPPORTEDFALL{\char"01AB\char"030C} \SUPPORTEDFALL{\char"01AD\char"030C} \SUPPORTEDFALL{\char"0236\char"030C} \SUPPORTEDPREC{\char"0075\char"030C} \MISSINGGLYPH{\char"AB52\char"030C} \SUPPORTEDANCH{\char"0079\char"030C} \SUPPORTEDANCH{\char"01B4\char"030C} \MISSINGGLYPH{\char"1EFF\char"030C}

\SUPPORTEDANCH{\char"0048\char"030D} \SUPPORTEDFALL{\char"2C67\char"030D} \SUPPORTEDANCH{\char"004B\char"030D} \SUPPORTEDANCH{\char"2C69\char"030D} \SUPPORTEDANCH{\char"004D\char"030D} \MISSINGGLYPH{\char"2C6E\char"030D} \SUPPORTEDANCH{\char"0050\char"030D} \SUPPORTEDFALL{\char"01A4\char"030D} \SUPPORTEDANCH{\char"0052\char"030D} \SUPPORTEDFALL{\char"2C64\char"030D} \SUPPORTEDANCH{\char"0054\char"030D} \SUPPORTEDFALL{\char"01AC\char"030D} \SUPPORTEDFALL{\char"01AE\char"030D} \SUPPORTEDANCH{\char"0059\char"030D} \SUPPORTEDFALL{\char"01B3\char"030D} \SUPPORTEDANCH{\char"0061\char"030D} \SUPPORTEDFALL{\char"1E9A\char"030D} \SUPPORTEDFALL{\char"0062\char"030D} \SUPPORTEDFALL{\char"0183\char"030D} \SUPPORTEDANCH{\char"0063\char"030D} \SUPPORTEDFALL{\char"0188\char"030D} \SUPPORTEDANCH{\char"0064\char"030D} \SUPPORTEDANCH{\char"018C\char"030D} \SUPPORTEDFALL{\char"0221\char"030D} \SUPPORTEDANCH{\char"0068\char"030D} \SUPPORTEDANCH{\char"2C68\char"030D} \SUPPORTEDANCH{\char"006B\char"030D} \SUPPORTEDFALL{\char"0199\char"030D} \SUPPORTEDANCH{\char"2C6A\char"030D} \SUPPORTEDANCH{\char"006C\char"030D} \SUPPORTEDFALL{\char"0234\char"030D} \SUPPORTEDANCH{\char"006E\char"030D} \SUPPORTEDANCH{\char"0235\char"030D} \SUPPORTEDANCH{\char"0070\char"030D} \SUPPORTEDFALL{\char"01A5\char"030D} \SUPPORTEDANCH{\char"0071\char"030D} \SUPPORTEDANCH{\char"02A0\char"030D} \SUPPORTEDFALL{\char"0074\char"030D} \SUPPORTEDFALL{\char"01AB\char"030D} \SUPPORTEDFALL{\char"01AD\char"030D} \SUPPORTEDFALL{\char"0236\char"030D} \SUPPORTEDANCH{\char"0075\char"030D} \MISSINGGLYPH{\char"AB52\char"030D} \SUPPORTEDANCH{\char"0079\char"030D} \SUPPORTEDANCH{\char"01B4\char"030D} \MISSINGGLYPH{\char"1EFF\char"030D}

\SUPPORTEDANCH{\char"0048\char"030F} \SUPPORTEDFALL{\char"2C67\char"030F} \SUPPORTEDANCH{\char"004B\char"030F} \SUPPORTEDANCH{\char"2C69\char"030F} \SUPPORTEDANCH{\char"004D\char"030F} \MISSINGGLYPH{\char"2C6E\char"030F} \SUPPORTEDANCH{\char"0050\char"030F} \SUPPORTEDFALL{\char"01A4\char"030F} \SUPPORTEDPREC{\char"0052\char"030F} \SUPPORTEDFALL{\char"2C64\char"030F} \SUPPORTEDANCH{\char"0054\char"030F} \SUPPORTEDFALL{\char"01AC\char"030F} \SUPPORTEDFALL{\char"01AE\char"030F} \SUPPORTEDANCH{\char"0059\char"030F} \SUPPORTEDFALL{\char"01B3\char"030F} \SUPPORTEDPREC{\char"0061\char"030F} \SUPPORTEDFALL{\char"1E9A\char"030F} \SUPPORTEDFALL{\char"0062\char"030F} \SUPPORTEDFALL{\char"0183\char"030F} \SUPPORTEDANCH{\char"0063\char"030F} \SUPPORTEDFALL{\char"0188\char"030F} \SUPPORTEDANCH{\char"0064\char"030F} \SUPPORTEDANCH{\char"018C\char"030F} \SUPPORTEDFALL{\char"0221\char"030F} \SUPPORTEDANCH{\char"0068\char"030F} \SUPPORTEDANCH{\char"2C68\char"030F} \SUPPORTEDANCH{\char"006B\char"030F} \SUPPORTEDFALL{\char"0199\char"030F} \SUPPORTEDANCH{\char"2C6A\char"030F} \SUPPORTEDANCH{\char"006C\char"030F} \SUPPORTEDFALL{\char"0234\char"030F} \SUPPORTEDANCH{\char"006E\char"030F} \SUPPORTEDANCH{\char"0235\char"030F} \SUPPORTEDANCH{\char"0070\char"030F} \SUPPORTEDFALL{\char"01A5\char"030F} \SUPPORTEDANCH{\char"0071\char"030F} \SUPPORTEDANCH{\char"02A0\char"030F} \SUPPORTEDFALL{\char"0074\char"030F} \SUPPORTEDFALL{\char"01AB\char"030F} \SUPPORTEDFALL{\char"01AD\char"030F} \SUPPORTEDFALL{\char"0236\char"030F} \SUPPORTEDPREC{\char"0075\char"030F} \MISSINGGLYPH{\char"AB52\char"030F} \SUPPORTEDANCH{\char"0079\char"030F} \SUPPORTEDANCH{\char"01B4\char"030F} \MISSINGGLYPH{\char"1EFF\char"030F}

\SUPPORTEDANCH{\char"0048\char"0310} \SUPPORTEDFALL{\char"2C67\char"0310} \SUPPORTEDANCH{\char"004B\char"0310} \SUPPORTEDANCH{\char"2C69\char"0310} \SUPPORTEDANCH{\char"004D\char"0310} \MISSINGGLYPH{\char"2C6E\char"0310} \SUPPORTEDANCH{\char"0050\char"0310} \SUPPORTEDFALL{\char"01A4\char"0310} \SUPPORTEDANCH{\char"0052\char"0310} \SUPPORTEDFALL{\char"2C64\char"0310} \SUPPORTEDANCH{\char"0054\char"0310} \SUPPORTEDFALL{\char"01AC\char"0310} \SUPPORTEDFALL{\char"01AE\char"0310} \SUPPORTEDANCH{\char"0059\char"0310} \SUPPORTEDFALL{\char"01B3\char"0310} \SUPPORTEDANCH{\char"0061\char"0310} \SUPPORTEDFALL{\char"1E9A\char"0310} \SUPPORTEDFALL{\char"0062\char"0310} \SUPPORTEDFALL{\char"0183\char"0310} \SUPPORTEDANCH{\char"0063\char"0310} \SUPPORTEDFALL{\char"0188\char"0310} \SUPPORTEDANCH{\char"0064\char"0310} \SUPPORTEDANCH{\char"018C\char"0310} \SUPPORTEDFALL{\char"0221\char"0310} \SUPPORTEDANCH{\char"0068\char"0310} \SUPPORTEDANCH{\char"2C68\char"0310} \SUPPORTEDANCH{\char"006B\char"0310} \SUPPORTEDFALL{\char"0199\char"0310} \SUPPORTEDANCH{\char"2C6A\char"0310} \SUPPORTEDANCH{\char"006C\char"0310} \SUPPORTEDFALL{\char"0234\char"0310} \SUPPORTEDANCH{\char"006E\char"0310} \SUPPORTEDANCH{\char"0235\char"0310} \SUPPORTEDANCH{\char"0070\char"0310} \SUPPORTEDFALL{\char"01A5\char"0310} \SUPPORTEDANCH{\char"0071\char"0310} \SUPPORTEDANCH{\char"02A0\char"0310} \SUPPORTEDFALL{\char"0074\char"0310} \SUPPORTEDFALL{\char"01AB\char"0310} \SUPPORTEDFALL{\char"01AD\char"0310} \SUPPORTEDFALL{\char"0236\char"0310} \SUPPORTEDANCH{\char"0075\char"0310} \MISSINGGLYPH{\char"AB52\char"0310} \SUPPORTEDANCH{\char"0079\char"0310} \SUPPORTEDANCH{\char"01B4\char"0310} \MISSINGGLYPH{\char"1EFF\char"0310}

\SUPPORTEDANCH{\char"0048\char"0311} \SUPPORTEDFALL{\char"2C67\char"0311} \SUPPORTEDANCH{\char"004B\char"0311} \SUPPORTEDANCH{\char"2C69\char"0311} \SUPPORTEDANCH{\char"004D\char"0311} \MISSINGGLYPH{\char"2C6E\char"0311} \SUPPORTEDANCH{\char"0050\char"0311} \SUPPORTEDFALL{\char"01A4\char"0311} \SUPPORTEDPREC{\char"0052\char"0311} \SUPPORTEDFALL{\char"2C64\char"0311} \SUPPORTEDANCH{\char"0054\char"0311} \SUPPORTEDFALL{\char"01AC\char"0311} \SUPPORTEDFALL{\char"01AE\char"0311} \SUPPORTEDANCH{\char"0059\char"0311} \SUPPORTEDFALL{\char"01B3\char"0311} \SUPPORTEDPREC{\char"0061\char"0311} \SUPPORTEDFALL{\char"1E9A\char"0311} \SUPPORTEDFALL{\char"0062\char"0311} \SUPPORTEDFALL{\char"0183\char"0311} \SUPPORTEDANCH{\char"0063\char"0311} \SUPPORTEDFALL{\char"0188\char"0311} \SUPPORTEDANCH{\char"0064\char"0311} \SUPPORTEDANCH{\char"018C\char"0311} \SUPPORTEDFALL{\char"0221\char"0311} \SUPPORTEDANCH{\char"0068\char"0311} \SUPPORTEDANCH{\char"2C68\char"0311} \SUPPORTEDANCH{\char"006B\char"0311} \SUPPORTEDFALL{\char"0199\char"0311} \SUPPORTEDANCH{\char"2C6A\char"0311} \SUPPORTEDANCH{\char"006C\char"0311} \SUPPORTEDFALL{\char"0234\char"0311} \SUPPORTEDANCH{\char"006E\char"0311} \SUPPORTEDANCH{\char"0235\char"0311} \SUPPORTEDANCH{\char"0070\char"0311} \SUPPORTEDFALL{\char"01A5\char"0311} \SUPPORTEDANCH{\char"0071\char"0311} \SUPPORTEDANCH{\char"02A0\char"0311} \SUPPORTEDFALL{\char"0074\char"0311} \SUPPORTEDFALL{\char"01AB\char"0311} \SUPPORTEDFALL{\char"01AD\char"0311} \SUPPORTEDFALL{\char"0236\char"0311} \SUPPORTEDPREC{\char"0075\char"0311} \MISSINGGLYPH{\char"AB52\char"0311} \SUPPORTEDANCH{\char"0079\char"0311} \SUPPORTEDANCH{\char"01B4\char"0311} \MISSINGGLYPH{\char"1EFF\char"0311}

\SUPPORTEDANCH{\char"0048\char"0312} \SUPPORTEDFALL{\char"2C67\char"0312} \SUPPORTEDANCH{\char"004B\char"0312} \SUPPORTEDANCH{\char"2C69\char"0312} \SUPPORTEDANCH{\char"004D\char"0312} \MISSINGGLYPH{\char"2C6E\char"0312} \SUPPORTEDANCH{\char"0050\char"0312} \SUPPORTEDFALL{\char"01A4\char"0312} \SUPPORTEDANCH{\char"0052\char"0312} \SUPPORTEDFALL{\char"2C64\char"0312} \SUPPORTEDANCH{\char"0054\char"0312} \SUPPORTEDFALL{\char"01AC\char"0312} \SUPPORTEDFALL{\char"01AE\char"0312} \SUPPORTEDANCH{\char"0059\char"0312} \SUPPORTEDFALL{\char"01B3\char"0312} \SUPPORTEDANCH{\char"0061\char"0312} \SUPPORTEDFALL{\char"1E9A\char"0312} \SUPPORTEDFALL{\char"0062\char"0312} \SUPPORTEDFALL{\char"0183\char"0312} \SUPPORTEDANCH{\char"0063\char"0312} \SUPPORTEDFALL{\char"0188\char"0312} \SUPPORTEDANCH{\char"0064\char"0312} \SUPPORTEDANCH{\char"018C\char"0312} \SUPPORTEDFALL{\char"0221\char"0312} \SUPPORTEDANCH{\char"0068\char"0312} \SUPPORTEDANCH{\char"2C68\char"0312} \SUPPORTEDANCH{\char"006B\char"0312} \SUPPORTEDFALL{\char"0199\char"0312} \SUPPORTEDANCH{\char"2C6A\char"0312} \SUPPORTEDANCH{\char"006C\char"0312} \SUPPORTEDFALL{\char"0234\char"0312} \SUPPORTEDANCH{\char"006E\char"0312} \SUPPORTEDANCH{\char"0235\char"0312} \SUPPORTEDANCH{\char"0070\char"0312} \SUPPORTEDFALL{\char"01A5\char"0312} \SUPPORTEDANCH{\char"0071\char"0312} \SUPPORTEDANCH{\char"02A0\char"0312} \SUPPORTEDFALL{\char"0074\char"0312} \SUPPORTEDFALL{\char"01AB\char"0312} \SUPPORTEDFALL{\char"01AD\char"0312} \SUPPORTEDFALL{\char"0236\char"0312} \SUPPORTEDANCH{\char"0075\char"0312} \MISSINGGLYPH{\char"AB52\char"0312} \SUPPORTEDANCH{\char"0079\char"0312} \SUPPORTEDANCH{\char"01B4\char"0312} \MISSINGGLYPH{\char"1EFF\char"0312}

\SUPPORTEDANCH{\char"0048\char"033D} \SUPPORTEDFALL{\char"2C67\char"033D} \SUPPORTEDANCH{\char"004B\char"033D} \SUPPORTEDANCH{\char"2C69\char"033D} \SUPPORTEDANCH{\char"004D\char"033D} \MISSINGGLYPH{\char"2C6E\char"033D} \SUPPORTEDANCH{\char"0050\char"033D} \SUPPORTEDFALL{\char"01A4\char"033D} \SUPPORTEDANCH{\char"0052\char"033D} \SUPPORTEDFALL{\char"2C64\char"033D} \SUPPORTEDANCH{\char"0054\char"033D} \SUPPORTEDFALL{\char"01AC\char"033D} \SUPPORTEDFALL{\char"01AE\char"033D} \SUPPORTEDANCH{\char"0059\char"033D} \SUPPORTEDFALL{\char"01B3\char"033D} \SUPPORTEDANCH{\char"0061\char"033D} \SUPPORTEDFALL{\char"1E9A\char"033D} \SUPPORTEDFALL{\char"0062\char"033D} \SUPPORTEDFALL{\char"0183\char"033D} \SUPPORTEDANCH{\char"0063\char"033D} \SUPPORTEDFALL{\char"0188\char"033D} \SUPPORTEDANCH{\char"0064\char"033D} \SUPPORTEDANCH{\char"018C\char"033D} \SUPPORTEDFALL{\char"0221\char"033D} \SUPPORTEDANCH{\char"0068\char"033D} \SUPPORTEDANCH{\char"2C68\char"033D} \SUPPORTEDANCH{\char"006B\char"033D} \SUPPORTEDFALL{\char"0199\char"033D} \SUPPORTEDANCH{\char"2C6A\char"033D} \SUPPORTEDANCH{\char"006C\char"033D} \SUPPORTEDFALL{\char"0234\char"033D} \SUPPORTEDANCH{\char"006E\char"033D} \SUPPORTEDANCH{\char"0235\char"033D} \SUPPORTEDANCH{\char"0070\char"033D} \SUPPORTEDFALL{\char"01A5\char"033D} \SUPPORTEDANCH{\char"0071\char"033D} \SUPPORTEDANCH{\char"02A0\char"033D} \SUPPORTEDFALL{\char"0074\char"033D} \SUPPORTEDFALL{\char"01AB\char"033D} \SUPPORTEDFALL{\char"01AD\char"033D} \SUPPORTEDFALL{\char"0236\char"033D} \SUPPORTEDANCH{\char"0075\char"033D} \MISSINGGLYPH{\char"AB52\char"033D} \SUPPORTEDANCH{\char"0079\char"033D} \SUPPORTEDANCH{\char"01B4\char"033D} \MISSINGGLYPH{\char"1EFF\char"033D}

\SUPPORTEDANCH{\char"0048\char"033E} \SUPPORTEDFALL{\char"2C67\char"033E} \SUPPORTEDANCH{\char"004B\char"033E} \SUPPORTEDANCH{\char"2C69\char"033E} \SUPPORTEDANCH{\char"004D\char"033E} \MISSINGGLYPH{\char"2C6E\char"033E} \SUPPORTEDANCH{\char"0050\char"033E} \SUPPORTEDFALL{\char"01A4\char"033E} \SUPPORTEDANCH{\char"0052\char"033E} \SUPPORTEDFALL{\char"2C64\char"033E} \SUPPORTEDANCH{\char"0054\char"033E} \SUPPORTEDFALL{\char"01AC\char"033E} \SUPPORTEDFALL{\char"01AE\char"033E} \SUPPORTEDANCH{\char"0059\char"033E} \SUPPORTEDFALL{\char"01B3\char"033E} \SUPPORTEDANCH{\char"0061\char"033E} \SUPPORTEDFALL{\char"1E9A\char"033E} \SUPPORTEDFALL{\char"0062\char"033E} \SUPPORTEDFALL{\char"0183\char"033E} \SUPPORTEDANCH{\char"0063\char"033E} \SUPPORTEDFALL{\char"0188\char"033E} \SUPPORTEDANCH{\char"0064\char"033E} \SUPPORTEDANCH{\char"018C\char"033E} \SUPPORTEDFALL{\char"0221\char"033E} \SUPPORTEDANCH{\char"0068\char"033E} \SUPPORTEDANCH{\char"2C68\char"033E} \SUPPORTEDANCH{\char"006B\char"033E} \SUPPORTEDFALL{\char"0199\char"033E} \SUPPORTEDANCH{\char"2C6A\char"033E} \SUPPORTEDANCH{\char"006C\char"033E} \SUPPORTEDFALL{\char"0234\char"033E} \SUPPORTEDANCH{\char"006E\char"033E} \SUPPORTEDANCH{\char"0235\char"033E} \SUPPORTEDANCH{\char"0070\char"033E} \SUPPORTEDFALL{\char"01A5\char"033E} \SUPPORTEDANCH{\char"0071\char"033E} \SUPPORTEDANCH{\char"02A0\char"033E} \SUPPORTEDFALL{\char"0074\char"033E} \SUPPORTEDFALL{\char"01AB\char"033E} \SUPPORTEDFALL{\char"01AD\char"033E} \SUPPORTEDFALL{\char"0236\char"033E} \SUPPORTEDANCH{\char"0075\char"033E} \MISSINGGLYPH{\char"AB52\char"033E} \SUPPORTEDANCH{\char"0079\char"033E} \SUPPORTEDANCH{\char"01B4\char"033E} \MISSINGGLYPH{\char"1EFF\char"033E}

\SUPPORTEDANCH{\char"0048\char"033F} \SUPPORTEDFALL{\char"2C67\char"033F} \SUPPORTEDANCH{\char"004B\char"033F} \SUPPORTEDANCH{\char"2C69\char"033F} \SUPPORTEDANCH{\char"004D\char"033F} \MISSINGGLYPH{\char"2C6E\char"033F} \SUPPORTEDANCH{\char"0050\char"033F} \SUPPORTEDFALL{\char"01A4\char"033F} \SUPPORTEDANCH{\char"0052\char"033F} \SUPPORTEDFALL{\char"2C64\char"033F} \SUPPORTEDANCH{\char"0054\char"033F} \SUPPORTEDFALL{\char"01AC\char"033F} \SUPPORTEDFALL{\char"01AE\char"033F} \SUPPORTEDANCH{\char"0059\char"033F} \SUPPORTEDFALL{\char"01B3\char"033F} \SUPPORTEDANCH{\char"0061\char"033F} \SUPPORTEDFALL{\char"1E9A\char"033F} \SUPPORTEDFALL{\char"0062\char"033F} \SUPPORTEDFALL{\char"0183\char"033F} \SUPPORTEDANCH{\char"0063\char"033F} \SUPPORTEDFALL{\char"0188\char"033F} \SUPPORTEDANCH{\char"0064\char"033F} \SUPPORTEDANCH{\char"018C\char"033F} \SUPPORTEDFALL{\char"0221\char"033F} \SUPPORTEDANCH{\char"0068\char"033F} \SUPPORTEDANCH{\char"2C68\char"033F} \SUPPORTEDANCH{\char"006B\char"033F} \SUPPORTEDFALL{\char"0199\char"033F} \SUPPORTEDANCH{\char"2C6A\char"033F} \SUPPORTEDANCH{\char"006C\char"033F} \SUPPORTEDFALL{\char"0234\char"033F} \SUPPORTEDANCH{\char"006E\char"033F} \SUPPORTEDANCH{\char"0235\char"033F} \SUPPORTEDANCH{\char"0070\char"033F} \SUPPORTEDFALL{\char"01A5\char"033F} \SUPPORTEDANCH{\char"0071\char"033F} \SUPPORTEDANCH{\char"02A0\char"033F} \SUPPORTEDFALL{\char"0074\char"033F} \SUPPORTEDFALL{\char"01AB\char"033F} \SUPPORTEDFALL{\char"01AD\char"033F} \SUPPORTEDFALL{\char"0236\char"033F} \SUPPORTEDANCH{\char"0075\char"033F} \MISSINGGLYPH{\char"AB52\char"033F} \SUPPORTEDANCH{\char"0079\char"033F} \SUPPORTEDANCH{\char"01B4\char"033F} \MISSINGGLYPH{\char"1EFF\char"033F}

\SUPPORTEDANCH{\char"0048\char"030E} \SUPPORTEDFALL{\char"2C67\char"030E} \SUPPORTEDANCH{\char"004B\char"030E} \SUPPORTEDANCH{\char"2C69\char"030E} \SUPPORTEDANCH{\char"004D\char"030E} \MISSINGGLYPH{\char"2C6E\char"030E} \SUPPORTEDANCH{\char"0050\char"030E} \SUPPORTEDFALL{\char"01A4\char"030E} \SUPPORTEDANCH{\char"0052\char"030E} \SUPPORTEDFALL{\char"2C64\char"030E} \SUPPORTEDANCH{\char"0054\char"030E} \SUPPORTEDFALL{\char"01AC\char"030E} \SUPPORTEDFALL{\char"01AE\char"030E} \SUPPORTEDANCH{\char"0059\char"030E} \SUPPORTEDFALL{\char"01B3\char"030E} \SUPPORTEDANCH{\char"0061\char"030E} \SUPPORTEDFALL{\char"1E9A\char"030E} \SUPPORTEDFALL{\char"0062\char"030E} \SUPPORTEDFALL{\char"0183\char"030E} \SUPPORTEDANCH{\char"0063\char"030E} \SUPPORTEDFALL{\char"0188\char"030E} \SUPPORTEDANCH{\char"0064\char"030E} \SUPPORTEDANCH{\char"018C\char"030E} \SUPPORTEDFALL{\char"0221\char"030E} \SUPPORTEDANCH{\char"0068\char"030E} \SUPPORTEDANCH{\char"2C68\char"030E} \SUPPORTEDANCH{\char"006B\char"030E} \SUPPORTEDFALL{\char"0199\char"030E} \SUPPORTEDANCH{\char"2C6A\char"030E} \SUPPORTEDANCH{\char"006C\char"030E} \SUPPORTEDFALL{\char"0234\char"030E} \SUPPORTEDANCH{\char"006E\char"030E} \SUPPORTEDANCH{\char"0235\char"030E} \SUPPORTEDANCH{\char"0070\char"030E} \SUPPORTEDFALL{\char"01A5\char"030E} \SUPPORTEDANCH{\char"0071\char"030E} \SUPPORTEDANCH{\char"02A0\char"030E} \SUPPORTEDFALL{\char"0074\char"030E} \SUPPORTEDFALL{\char"01AB\char"030E} \SUPPORTEDFALL{\char"01AD\char"030E} \SUPPORTEDFALL{\char"0236\char"030E} \SUPPORTEDANCH{\char"0075\char"030E} \MISSINGGLYPH{\char"AB52\char"030E} \SUPPORTEDANCH{\char"0079\char"030E} \SUPPORTEDANCH{\char"01B4\char"030E} \MISSINGGLYPH{\char"1EFF\char"030E}

\SUPPORTEDANCH{\char"0048\char"0346} \SUPPORTEDFALL{\char"2C67\char"0346} \SUPPORTEDANCH{\char"004B\char"0346} \SUPPORTEDANCH{\char"2C69\char"0346} \SUPPORTEDANCH{\char"004D\char"0346} \MISSINGGLYPH{\char"2C6E\char"0346} \SUPPORTEDANCH{\char"0050\char"0346} \SUPPORTEDFALL{\char"01A4\char"0346} \SUPPORTEDANCH{\char"0052\char"0346} \SUPPORTEDFALL{\char"2C64\char"0346} \SUPPORTEDANCH{\char"0054\char"0346} \SUPPORTEDFALL{\char"01AC\char"0346} \SUPPORTEDFALL{\char"01AE\char"0346} \SUPPORTEDANCH{\char"0059\char"0346} \SUPPORTEDFALL{\char"01B3\char"0346} \SUPPORTEDANCH{\char"0061\char"0346} \SUPPORTEDFALL{\char"1E9A\char"0346} \SUPPORTEDFALL{\char"0062\char"0346} \SUPPORTEDFALL{\char"0183\char"0346} \SUPPORTEDANCH{\char"0063\char"0346} \SUPPORTEDFALL{\char"0188\char"0346} \SUPPORTEDANCH{\char"0064\char"0346} \SUPPORTEDANCH{\char"018C\char"0346} \SUPPORTEDFALL{\char"0221\char"0346} \SUPPORTEDANCH{\char"0068\char"0346} \SUPPORTEDANCH{\char"2C68\char"0346} \SUPPORTEDANCH{\char"006B\char"0346} \SUPPORTEDFALL{\char"0199\char"0346} \SUPPORTEDANCH{\char"2C6A\char"0346} \SUPPORTEDANCH{\char"006C\char"0346} \SUPPORTEDFALL{\char"0234\char"0346} \SUPPORTEDANCH{\char"006E\char"0346} \SUPPORTEDANCH{\char"0235\char"0346} \SUPPORTEDANCH{\char"0070\char"0346} \SUPPORTEDFALL{\char"01A5\char"0346} \SUPPORTEDANCH{\char"0071\char"0346} \SUPPORTEDANCH{\char"02A0\char"0346} \SUPPORTEDFALL{\char"0074\char"0346} \SUPPORTEDFALL{\char"01AB\char"0346} \SUPPORTEDFALL{\char"01AD\char"0346} \SUPPORTEDFALL{\char"0236\char"0346} \SUPPORTEDANCH{\char"0075\char"0346} \MISSINGGLYPH{\char"AB52\char"0346} \SUPPORTEDANCH{\char"0079\char"0346} \SUPPORTEDANCH{\char"01B4\char"0346} \MISSINGGLYPH{\char"1EFF\char"0346}

\SUPPORTEDANCH{\char"0048\char"034A} \SUPPORTEDFALL{\char"2C67\char"034A} \SUPPORTEDANCH{\char"004B\char"034A} \SUPPORTEDANCH{\char"2C69\char"034A} \SUPPORTEDANCH{\char"004D\char"034A} \MISSINGGLYPH{\char"2C6E\char"034A} \SUPPORTEDANCH{\char"0050\char"034A} \SUPPORTEDFALL{\char"01A4\char"034A} \SUPPORTEDANCH{\char"0052\char"034A} \SUPPORTEDFALL{\char"2C64\char"034A} \SUPPORTEDANCH{\char"0054\char"034A} \SUPPORTEDFALL{\char"01AC\char"034A} \SUPPORTEDFALL{\char"01AE\char"034A} \SUPPORTEDANCH{\char"0059\char"034A} \SUPPORTEDFALL{\char"01B3\char"034A} \SUPPORTEDANCH{\char"0061\char"034A} \SUPPORTEDFALL{\char"1E9A\char"034A} \SUPPORTEDFALL{\char"0062\char"034A} \SUPPORTEDFALL{\char"0183\char"034A} \SUPPORTEDANCH{\char"0063\char"034A} \SUPPORTEDFALL{\char"0188\char"034A} \SUPPORTEDANCH{\char"0064\char"034A} \SUPPORTEDANCH{\char"018C\char"034A} \SUPPORTEDFALL{\char"0221\char"034A} \SUPPORTEDANCH{\char"0068\char"034A} \SUPPORTEDANCH{\char"2C68\char"034A} \SUPPORTEDANCH{\char"006B\char"034A} \SUPPORTEDFALL{\char"0199\char"034A} \SUPPORTEDANCH{\char"2C6A\char"034A} \SUPPORTEDANCH{\char"006C\char"034A} \SUPPORTEDFALL{\char"0234\char"034A} \SUPPORTEDANCH{\char"006E\char"034A} \SUPPORTEDANCH{\char"0235\char"034A} \SUPPORTEDANCH{\char"0070\char"034A} \SUPPORTEDFALL{\char"01A5\char"034A} \SUPPORTEDANCH{\char"0071\char"034A} \SUPPORTEDANCH{\char"02A0\char"034A} \SUPPORTEDFALL{\char"0074\char"034A} \SUPPORTEDFALL{\char"01AB\char"034A} \SUPPORTEDFALL{\char"01AD\char"034A} \SUPPORTEDFALL{\char"0236\char"034A} \SUPPORTEDANCH{\char"0075\char"034A} \MISSINGGLYPH{\char"AB52\char"034A} \SUPPORTEDANCH{\char"0079\char"034A} \SUPPORTEDANCH{\char"01B4\char"034A} \MISSINGGLYPH{\char"1EFF\char"034A}

\SUPPORTEDANCH{\char"0048\char"034B} \SUPPORTEDFALL{\char"2C67\char"034B} \SUPPORTEDANCH{\char"004B\char"034B} \SUPPORTEDANCH{\char"2C69\char"034B} \SUPPORTEDANCH{\char"004D\char"034B} \MISSINGGLYPH{\char"2C6E\char"034B} \SUPPORTEDANCH{\char"0050\char"034B} \SUPPORTEDFALL{\char"01A4\char"034B} \SUPPORTEDANCH{\char"0052\char"034B} \SUPPORTEDFALL{\char"2C64\char"034B} \SUPPORTEDANCH{\char"0054\char"034B} \SUPPORTEDFALL{\char"01AC\char"034B} \SUPPORTEDFALL{\char"01AE\char"034B} \SUPPORTEDANCH{\char"0059\char"034B} \SUPPORTEDFALL{\char"01B3\char"034B} \SUPPORTEDANCH{\char"0061\char"034B} \SUPPORTEDFALL{\char"1E9A\char"034B} \SUPPORTEDFALL{\char"0062\char"034B} \SUPPORTEDFALL{\char"0183\char"034B} \SUPPORTEDANCH{\char"0063\char"034B} \SUPPORTEDFALL{\char"0188\char"034B} \SUPPORTEDANCH{\char"0064\char"034B} \SUPPORTEDANCH{\char"018C\char"034B} \SUPPORTEDFALL{\char"0221\char"034B} \SUPPORTEDANCH{\char"0068\char"034B} \SUPPORTEDANCH{\char"2C68\char"034B} \SUPPORTEDANCH{\char"006B\char"034B} \SUPPORTEDFALL{\char"0199\char"034B} \SUPPORTEDANCH{\char"2C6A\char"034B} \SUPPORTEDANCH{\char"006C\char"034B} \SUPPORTEDFALL{\char"0234\char"034B} \SUPPORTEDANCH{\char"006E\char"034B} \SUPPORTEDANCH{\char"0235\char"034B} \SUPPORTEDANCH{\char"0070\char"034B} \SUPPORTEDFALL{\char"01A5\char"034B} \SUPPORTEDANCH{\char"0071\char"034B} \SUPPORTEDANCH{\char"02A0\char"034B} \SUPPORTEDFALL{\char"0074\char"034B} \SUPPORTEDFALL{\char"01AB\char"034B} \SUPPORTEDFALL{\char"01AD\char"034B} \SUPPORTEDFALL{\char"0236\char"034B} \SUPPORTEDANCH{\char"0075\char"034B} \MISSINGGLYPH{\char"AB52\char"034B} \SUPPORTEDANCH{\char"0079\char"034B} \SUPPORTEDANCH{\char"01B4\char"034B} \MISSINGGLYPH{\char"1EFF\char"034B}

\SUPPORTEDANCH{\char"0048\char"034C} \SUPPORTEDFALL{\char"2C67\char"034C} \SUPPORTEDANCH{\char"004B\char"034C} \SUPPORTEDANCH{\char"2C69\char"034C} \SUPPORTEDANCH{\char"004D\char"034C} \MISSINGGLYPH{\char"2C6E\char"034C} \SUPPORTEDANCH{\char"0050\char"034C} \SUPPORTEDFALL{\char"01A4\char"034C} \SUPPORTEDANCH{\char"0052\char"034C} \SUPPORTEDFALL{\char"2C64\char"034C} \SUPPORTEDANCH{\char"0054\char"034C} \SUPPORTEDFALL{\char"01AC\char"034C} \SUPPORTEDFALL{\char"01AE\char"034C} \SUPPORTEDANCH{\char"0059\char"034C} \SUPPORTEDFALL{\char"01B3\char"034C} \SUPPORTEDANCH{\char"0061\char"034C} \SUPPORTEDFALL{\char"1E9A\char"034C} \SUPPORTEDFALL{\char"0062\char"034C} \SUPPORTEDFALL{\char"0183\char"034C} \SUPPORTEDANCH{\char"0063\char"034C} \SUPPORTEDFALL{\char"0188\char"034C} \SUPPORTEDANCH{\char"0064\char"034C} \SUPPORTEDANCH{\char"018C\char"034C} \SUPPORTEDFALL{\char"0221\char"034C} \SUPPORTEDANCH{\char"0068\char"034C} \SUPPORTEDANCH{\char"2C68\char"034C} \SUPPORTEDANCH{\char"006B\char"034C} \SUPPORTEDFALL{\char"0199\char"034C} \SUPPORTEDANCH{\char"2C6A\char"034C} \SUPPORTEDANCH{\char"006C\char"034C} \SUPPORTEDFALL{\char"0234\char"034C} \SUPPORTEDANCH{\char"006E\char"034C} \SUPPORTEDANCH{\char"0235\char"034C} \SUPPORTEDANCH{\char"0070\char"034C} \SUPPORTEDFALL{\char"01A5\char"034C} \SUPPORTEDANCH{\char"0071\char"034C} \SUPPORTEDANCH{\char"02A0\char"034C} \SUPPORTEDFALL{\char"0074\char"034C} \SUPPORTEDFALL{\char"01AB\char"034C} \SUPPORTEDFALL{\char"01AD\char"034C} \SUPPORTEDFALL{\char"0236\char"034C} \SUPPORTEDANCH{\char"0075\char"034C} \MISSINGGLYPH{\char"AB52\char"034C} \SUPPORTEDANCH{\char"0079\char"034C} \SUPPORTEDANCH{\char"01B4\char"034C} \MISSINGGLYPH{\char"1EFF\char"034C}

\SUPPORTEDANCH{\char"0048\char"0350} \SUPPORTEDFALL{\char"2C67\char"0350} \SUPPORTEDANCH{\char"004B\char"0350} \SUPPORTEDANCH{\char"2C69\char"0350} \SUPPORTEDANCH{\char"004D\char"0350} \MISSINGGLYPH{\char"2C6E\char"0350} \SUPPORTEDANCH{\char"0050\char"0350} \SUPPORTEDFALL{\char"01A4\char"0350} \SUPPORTEDANCH{\char"0052\char"0350} \SUPPORTEDFALL{\char"2C64\char"0350} \SUPPORTEDANCH{\char"0054\char"0350} \SUPPORTEDFALL{\char"01AC\char"0350} \SUPPORTEDFALL{\char"01AE\char"0350} \SUPPORTEDANCH{\char"0059\char"0350} \SUPPORTEDFALL{\char"01B3\char"0350} \SUPPORTEDANCH{\char"0061\char"0350} \SUPPORTEDFALL{\char"1E9A\char"0350} \SUPPORTEDFALL{\char"0062\char"0350} \SUPPORTEDFALL{\char"0183\char"0350} \SUPPORTEDANCH{\char"0063\char"0350} \SUPPORTEDFALL{\char"0188\char"0350} \SUPPORTEDANCH{\char"0064\char"0350} \SUPPORTEDANCH{\char"018C\char"0350} \SUPPORTEDFALL{\char"0221\char"0350} \SUPPORTEDANCH{\char"0068\char"0350} \SUPPORTEDANCH{\char"2C68\char"0350} \SUPPORTEDANCH{\char"006B\char"0350} \SUPPORTEDFALL{\char"0199\char"0350} \SUPPORTEDANCH{\char"2C6A\char"0350} \SUPPORTEDANCH{\char"006C\char"0350} \SUPPORTEDFALL{\char"0234\char"0350} \SUPPORTEDANCH{\char"006E\char"0350} \SUPPORTEDANCH{\char"0235\char"0350} \SUPPORTEDANCH{\char"0070\char"0350} \SUPPORTEDFALL{\char"01A5\char"0350} \SUPPORTEDANCH{\char"0071\char"0350} \SUPPORTEDANCH{\char"02A0\char"0350} \SUPPORTEDFALL{\char"0074\char"0350} \SUPPORTEDFALL{\char"01AB\char"0350} \SUPPORTEDFALL{\char"01AD\char"0350} \SUPPORTEDFALL{\char"0236\char"0350} \SUPPORTEDANCH{\char"0075\char"0350} \MISSINGGLYPH{\char"AB52\char"0350} \SUPPORTEDANCH{\char"0079\char"0350} \SUPPORTEDANCH{\char"01B4\char"0350} \MISSINGGLYPH{\char"1EFF\char"0350}

\SUPPORTEDANCH{\char"0048\char"0351} \SUPPORTEDFALL{\char"2C67\char"0351} \SUPPORTEDANCH{\char"004B\char"0351} \SUPPORTEDANCH{\char"2C69\char"0351} \SUPPORTEDANCH{\char"004D\char"0351} \MISSINGGLYPH{\char"2C6E\char"0351} \SUPPORTEDANCH{\char"0050\char"0351} \SUPPORTEDFALL{\char"01A4\char"0351} \SUPPORTEDANCH{\char"0052\char"0351} \SUPPORTEDFALL{\char"2C64\char"0351} \SUPPORTEDANCH{\char"0054\char"0351} \SUPPORTEDFALL{\char"01AC\char"0351} \SUPPORTEDFALL{\char"01AE\char"0351} \SUPPORTEDANCH{\char"0059\char"0351} \SUPPORTEDFALL{\char"01B3\char"0351} \SUPPORTEDANCH{\char"0061\char"0351} \SUPPORTEDFALL{\char"1E9A\char"0351} \SUPPORTEDFALL{\char"0062\char"0351} \SUPPORTEDFALL{\char"0183\char"0351} \SUPPORTEDANCH{\char"0063\char"0351} \SUPPORTEDFALL{\char"0188\char"0351} \SUPPORTEDANCH{\char"0064\char"0351} \SUPPORTEDANCH{\char"018C\char"0351} \SUPPORTEDFALL{\char"0221\char"0351} \SUPPORTEDANCH{\char"0068\char"0351} \SUPPORTEDANCH{\char"2C68\char"0351} \SUPPORTEDANCH{\char"006B\char"0351} \SUPPORTEDFALL{\char"0199\char"0351} \SUPPORTEDANCH{\char"2C6A\char"0351} \SUPPORTEDANCH{\char"006C\char"0351} \SUPPORTEDFALL{\char"0234\char"0351} \SUPPORTEDANCH{\char"006E\char"0351} \SUPPORTEDANCH{\char"0235\char"0351} \SUPPORTEDANCH{\char"0070\char"0351} \SUPPORTEDFALL{\char"01A5\char"0351} \SUPPORTEDANCH{\char"0071\char"0351} \SUPPORTEDANCH{\char"02A0\char"0351} \SUPPORTEDFALL{\char"0074\char"0351} \SUPPORTEDFALL{\char"01AB\char"0351} \SUPPORTEDFALL{\char"01AD\char"0351} \SUPPORTEDFALL{\char"0236\char"0351} \SUPPORTEDANCH{\char"0075\char"0351} \MISSINGGLYPH{\char"AB52\char"0351} \SUPPORTEDANCH{\char"0079\char"0351} \SUPPORTEDANCH{\char"01B4\char"0351} \MISSINGGLYPH{\char"1EFF\char"0351}

\SUPPORTEDANCH{\char"0048\char"0352} \SUPPORTEDFALL{\char"2C67\char"0352} \SUPPORTEDANCH{\char"004B\char"0352} \SUPPORTEDANCH{\char"2C69\char"0352} \SUPPORTEDANCH{\char"004D\char"0352} \MISSINGGLYPH{\char"2C6E\char"0352} \SUPPORTEDANCH{\char"0050\char"0352} \SUPPORTEDFALL{\char"01A4\char"0352} \SUPPORTEDANCH{\char"0052\char"0352} \SUPPORTEDFALL{\char"2C64\char"0352} \SUPPORTEDANCH{\char"0054\char"0352} \SUPPORTEDFALL{\char"01AC\char"0352} \SUPPORTEDFALL{\char"01AE\char"0352} \SUPPORTEDANCH{\char"0059\char"0352} \SUPPORTEDFALL{\char"01B3\char"0352} \SUPPORTEDANCH{\char"0061\char"0352} \SUPPORTEDFALL{\char"1E9A\char"0352} \SUPPORTEDFALL{\char"0062\char"0352} \SUPPORTEDFALL{\char"0183\char"0352} \SUPPORTEDANCH{\char"0063\char"0352} \SUPPORTEDFALL{\char"0188\char"0352} \SUPPORTEDANCH{\char"0064\char"0352} \SUPPORTEDANCH{\char"018C\char"0352} \SUPPORTEDFALL{\char"0221\char"0352} \SUPPORTEDANCH{\char"0068\char"0352} \SUPPORTEDANCH{\char"2C68\char"0352} \SUPPORTEDANCH{\char"006B\char"0352} \SUPPORTEDFALL{\char"0199\char"0352} \SUPPORTEDANCH{\char"2C6A\char"0352} \SUPPORTEDANCH{\char"006C\char"0352} \SUPPORTEDFALL{\char"0234\char"0352} \SUPPORTEDANCH{\char"006E\char"0352} \SUPPORTEDANCH{\char"0235\char"0352} \SUPPORTEDANCH{\char"0070\char"0352} \SUPPORTEDFALL{\char"01A5\char"0352} \SUPPORTEDANCH{\char"0071\char"0352} \SUPPORTEDANCH{\char"02A0\char"0352} \SUPPORTEDFALL{\char"0074\char"0352} \SUPPORTEDFALL{\char"01AB\char"0352} \SUPPORTEDFALL{\char"01AD\char"0352} \SUPPORTEDFALL{\char"0236\char"0352} \SUPPORTEDANCH{\char"0075\char"0352} \MISSINGGLYPH{\char"AB52\char"0352} \SUPPORTEDANCH{\char"0079\char"0352} \SUPPORTEDANCH{\char"01B4\char"0352} \MISSINGGLYPH{\char"1EFF\char"0352}

\SUPPORTEDANCH{\char"0048\char"0357} \SUPPORTEDFALL{\char"2C67\char"0357} \SUPPORTEDANCH{\char"004B\char"0357} \SUPPORTEDANCH{\char"2C69\char"0357} \SUPPORTEDANCH{\char"004D\char"0357} \MISSINGGLYPH{\char"2C6E\char"0357} \SUPPORTEDANCH{\char"0050\char"0357} \SUPPORTEDFALL{\char"01A4\char"0357} \SUPPORTEDANCH{\char"0052\char"0357} \SUPPORTEDFALL{\char"2C64\char"0357} \SUPPORTEDANCH{\char"0054\char"0357} \SUPPORTEDFALL{\char"01AC\char"0357} \SUPPORTEDFALL{\char"01AE\char"0357} \SUPPORTEDANCH{\char"0059\char"0357} \SUPPORTEDFALL{\char"01B3\char"0357} \SUPPORTEDANCH{\char"0061\char"0357} \SUPPORTEDFALL{\char"1E9A\char"0357} \SUPPORTEDFALL{\char"0062\char"0357} \SUPPORTEDFALL{\char"0183\char"0357} \SUPPORTEDANCH{\char"0063\char"0357} \SUPPORTEDFALL{\char"0188\char"0357} \SUPPORTEDANCH{\char"0064\char"0357} \SUPPORTEDANCH{\char"018C\char"0357} \SUPPORTEDFALL{\char"0221\char"0357} \SUPPORTEDANCH{\char"0068\char"0357} \SUPPORTEDANCH{\char"2C68\char"0357} \SUPPORTEDANCH{\char"006B\char"0357} \SUPPORTEDFALL{\char"0199\char"0357} \SUPPORTEDANCH{\char"2C6A\char"0357} \SUPPORTEDANCH{\char"006C\char"0357} \SUPPORTEDFALL{\char"0234\char"0357} \SUPPORTEDANCH{\char"006E\char"0357} \SUPPORTEDANCH{\char"0235\char"0357} \SUPPORTEDANCH{\char"0070\char"0357} \SUPPORTEDFALL{\char"01A5\char"0357} \SUPPORTEDANCH{\char"0071\char"0357} \SUPPORTEDANCH{\char"02A0\char"0357} \SUPPORTEDFALL{\char"0074\char"0357} \SUPPORTEDFALL{\char"01AB\char"0357} \SUPPORTEDFALL{\char"01AD\char"0357} \SUPPORTEDFALL{\char"0236\char"0357} \SUPPORTEDANCH{\char"0075\char"0357} \MISSINGGLYPH{\char"AB52\char"0357} \SUPPORTEDANCH{\char"0079\char"0357} \SUPPORTEDANCH{\char"01B4\char"0357} \MISSINGGLYPH{\char"1EFF\char"0357}

\SUPPORTEDANCH{\char"0048\char"035B} \SUPPORTEDFALL{\char"2C67\char"035B} \SUPPORTEDANCH{\char"004B\char"035B} \SUPPORTEDANCH{\char"2C69\char"035B} \SUPPORTEDANCH{\char"004D\char"035B} \MISSINGGLYPH{\char"2C6E\char"035B} \SUPPORTEDANCH{\char"0050\char"035B} \SUPPORTEDFALL{\char"01A4\char"035B} \SUPPORTEDANCH{\char"0052\char"035B} \SUPPORTEDFALL{\char"2C64\char"035B} \SUPPORTEDANCH{\char"0054\char"035B} \SUPPORTEDFALL{\char"01AC\char"035B} \SUPPORTEDFALL{\char"01AE\char"035B} \SUPPORTEDANCH{\char"0059\char"035B} \SUPPORTEDFALL{\char"01B3\char"035B} \SUPPORTEDANCH{\char"0061\char"035B} \SUPPORTEDFALL{\char"1E9A\char"035B} \SUPPORTEDFALL{\char"0062\char"035B} \SUPPORTEDFALL{\char"0183\char"035B} \SUPPORTEDANCH{\char"0063\char"035B} \SUPPORTEDFALL{\char"0188\char"035B} \SUPPORTEDANCH{\char"0064\char"035B} \SUPPORTEDANCH{\char"018C\char"035B} \SUPPORTEDFALL{\char"0221\char"035B} \SUPPORTEDANCH{\char"0068\char"035B} \SUPPORTEDANCH{\char"2C68\char"035B} \SUPPORTEDANCH{\char"006B\char"035B} \SUPPORTEDFALL{\char"0199\char"035B} \SUPPORTEDANCH{\char"2C6A\char"035B} \SUPPORTEDANCH{\char"006C\char"035B} \SUPPORTEDFALL{\char"0234\char"035B} \SUPPORTEDANCH{\char"006E\char"035B} \SUPPORTEDANCH{\char"0235\char"035B} \SUPPORTEDANCH{\char"0070\char"035B} \SUPPORTEDFALL{\char"01A5\char"035B} \SUPPORTEDANCH{\char"0071\char"035B} \SUPPORTEDANCH{\char"02A0\char"035B} \SUPPORTEDFALL{\char"0074\char"035B} \SUPPORTEDFALL{\char"01AB\char"035B} \SUPPORTEDFALL{\char"01AD\char"035B} \SUPPORTEDFALL{\char"0236\char"035B} \SUPPORTEDANCH{\char"0075\char"035B} \MISSINGGLYPH{\char"AB52\char"035B} \SUPPORTEDANCH{\char"0079\char"035B} \SUPPORTEDANCH{\char"01B4\char"035B} \MISSINGGLYPH{\char"1EFF\char"035B}

}

\newpage
\subsection*{Italic patched, Group 4, above, Above}

{\LibertinusSerifPatch\itshape
% grid. columns are bases, rows are marks.
\SUPPORTEDANCH{\char"0048\char"0300} \SUPPORTEDANCH{\char"2C67\char"0300} \SUPPORTEDANCH{\char"004B\char"0300} \SUPPORTEDANCH{\char"2C69\char"0300} \SUPPORTEDANCH{\char"004D\char"0300} \MISSINGGLYPH{\char"2C6E\char"0300} \SUPPORTEDANCH{\char"0050\char"0300} \SUPPORTEDANCH{\char"01A4\char"0300} \SUPPORTEDANCH{\char"0052\char"0300} \SUPPORTEDANCH{\char"2C64\char"0300} \SUPPORTEDANCH{\char"0054\char"0300} \SUPPORTEDANCH{\char"01AC\char"0300} \SUPPORTEDANCH{\char"01AE\char"0300} \SUPPORTEDPREC{\char"0059\char"0300} \SUPPORTEDANCH{\char"01B3\char"0300} \SUPPORTEDPREC{\char"0061\char"0300} \SUPPORTEDANCH{\char"1E9A\char"0300} \SUPPORTEDANCH{\char"0062\char"0300} \SUPPORTEDANCH{\char"0183\char"0300} \SUPPORTEDANCH{\char"0063\char"0300} \SUPPORTEDANCH{\char"0188\char"0300} \SUPPORTEDANCH{\char"0064\char"0300} \SUPPORTEDANCH{\char"018C\char"0300} \SUPPORTEDANCH{\char"0221\char"0300} \SUPPORTEDANCH{\char"0068\char"0300} \SUPPORTEDANCH{\char"2C68\char"0300} \SUPPORTEDANCH{\char"006B\char"0300} \SUPPORTEDANCH{\char"0199\char"0300} \SUPPORTEDANCH{\char"2C6A\char"0300} \SUPPORTEDANCH{\char"006C\char"0300} \SUPPORTEDANCH{\char"0234\char"0300} \SUPPORTEDPREC{\char"006E\char"0300} \SUPPORTEDANCH{\char"0235\char"0300} \SUPPORTEDANCH{\char"0070\char"0300} \SUPPORTEDANCH{\char"01A5\char"0300} \SUPPORTEDANCH{\char"0071\char"0300} \SUPPORTEDANCH{\char"02A0\char"0300} \SUPPORTEDANCH{\char"0074\char"0300} \SUPPORTEDANCH{\char"01AB\char"0300} \SUPPORTEDANCH{\char"01AD\char"0300} \SUPPORTEDANCH{\char"0236\char"0300} \SUPPORTEDPREC{\char"0075\char"0300} \MISSINGGLYPH{\char"AB52\char"0300} \SUPPORTEDPREC{\char"0079\char"0300} \SUPPORTEDANCH{\char"01B4\char"0300} \MISSINGGLYPH{\char"1EFF\char"0300}

\SUPPORTEDANCH{\char"0048\char"0301} \SUPPORTEDANCH{\char"2C67\char"0301} \SUPPORTEDPREC{\char"004B\char"0301} \SUPPORTEDANCH{\char"2C69\char"0301} \SUPPORTEDPREC{\char"004D\char"0301} \MISSINGGLYPH{\char"2C6E\char"0301} \SUPPORTEDPREC{\char"0050\char"0301} \SUPPORTEDANCH{\char"01A4\char"0301} \SUPPORTEDPREC{\char"0052\char"0301} \SUPPORTEDANCH{\char"2C64\char"0301} \SUPPORTEDANCH{\char"0054\char"0301} \SUPPORTEDANCH{\char"01AC\char"0301} \SUPPORTEDANCH{\char"01AE\char"0301} \SUPPORTEDPREC{\char"0059\char"0301} \SUPPORTEDANCH{\char"01B3\char"0301} \SUPPORTEDPREC{\char"0061\char"0301} \SUPPORTEDANCH{\char"1E9A\char"0301} \SUPPORTEDANCH{\char"0062\char"0301} \SUPPORTEDANCH{\char"0183\char"0301} \SUPPORTEDPREC{\char"0063\char"0301} \SUPPORTEDANCH{\char"0188\char"0301} \SUPPORTEDANCH{\char"0064\char"0301} \SUPPORTEDANCH{\char"018C\char"0301} \SUPPORTEDANCH{\char"0221\char"0301} \SUPPORTEDANCH{\char"0068\char"0301} \SUPPORTEDANCH{\char"2C68\char"0301} \SUPPORTEDPREC{\char"006B\char"0301} \SUPPORTEDANCH{\char"0199\char"0301} \SUPPORTEDANCH{\char"2C6A\char"0301} \SUPPORTEDPREC{\char"006C\char"0301} \SUPPORTEDANCH{\char"0234\char"0301} \SUPPORTEDPREC{\char"006E\char"0301} \SUPPORTEDANCH{\char"0235\char"0301} \SUPPORTEDPREC{\char"0070\char"0301} \SUPPORTEDANCH{\char"01A5\char"0301} \SUPPORTEDANCH{\char"0071\char"0301} \SUPPORTEDANCH{\char"02A0\char"0301} \SUPPORTEDANCH{\char"0074\char"0301} \SUPPORTEDANCH{\char"01AB\char"0301} \SUPPORTEDANCH{\char"01AD\char"0301} \SUPPORTEDANCH{\char"0236\char"0301} \SUPPORTEDPREC{\char"0075\char"0301} \MISSINGGLYPH{\char"AB52\char"0301} \SUPPORTEDPREC{\char"0079\char"0301} \SUPPORTEDANCH{\char"01B4\char"0301} \MISSINGGLYPH{\char"1EFF\char"0301}

\SUPPORTEDPREC{\char"0048\char"0302} \SUPPORTEDANCH{\char"2C67\char"0302} \SUPPORTEDANCH{\char"004B\char"0302} \SUPPORTEDANCH{\char"2C69\char"0302} \SUPPORTEDANCH{\char"004D\char"0302} \MISSINGGLYPH{\char"2C6E\char"0302} \SUPPORTEDANCH{\char"0050\char"0302} \SUPPORTEDANCH{\char"01A4\char"0302} \SUPPORTEDANCH{\char"0052\char"0302} \SUPPORTEDANCH{\char"2C64\char"0302} \SUPPORTEDANCH{\char"0054\char"0302} \SUPPORTEDANCH{\char"01AC\char"0302} \SUPPORTEDANCH{\char"01AE\char"0302} \SUPPORTEDPREC{\char"0059\char"0302} \SUPPORTEDANCH{\char"01B3\char"0302} \SUPPORTEDPREC{\char"0061\char"0302} \SUPPORTEDANCH{\char"1E9A\char"0302} \SUPPORTEDANCH{\char"0062\char"0302} \SUPPORTEDANCH{\char"0183\char"0302} \SUPPORTEDPREC{\char"0063\char"0302} \SUPPORTEDANCH{\char"0188\char"0302} \SUPPORTEDANCH{\char"0064\char"0302} \SUPPORTEDANCH{\char"018C\char"0302} \SUPPORTEDANCH{\char"0221\char"0302} \SUPPORTEDPREC{\char"0068\char"0302} \SUPPORTEDANCH{\char"2C68\char"0302} \SUPPORTEDANCH{\char"006B\char"0302} \SUPPORTEDANCH{\char"0199\char"0302} \SUPPORTEDANCH{\char"2C6A\char"0302} \SUPPORTEDANCH{\char"006C\char"0302} \SUPPORTEDANCH{\char"0234\char"0302} \SUPPORTEDANCH{\char"006E\char"0302} \SUPPORTEDANCH{\char"0235\char"0302} \SUPPORTEDANCH{\char"0070\char"0302} \SUPPORTEDANCH{\char"01A5\char"0302} \SUPPORTEDANCH{\char"0071\char"0302} \SUPPORTEDANCH{\char"02A0\char"0302} \SUPPORTEDANCH{\char"0074\char"0302} \SUPPORTEDANCH{\char"01AB\char"0302} \SUPPORTEDANCH{\char"01AD\char"0302} \SUPPORTEDANCH{\char"0236\char"0302} \SUPPORTEDPREC{\char"0075\char"0302} \MISSINGGLYPH{\char"AB52\char"0302} \SUPPORTEDPREC{\char"0079\char"0302} \SUPPORTEDANCH{\char"01B4\char"0302} \MISSINGGLYPH{\char"1EFF\char"0302}

\SUPPORTEDANCH{\char"0048\char"0303} \SUPPORTEDANCH{\char"2C67\char"0303} \SUPPORTEDANCH{\char"004B\char"0303} \SUPPORTEDANCH{\char"2C69\char"0303} \SUPPORTEDANCH{\char"004D\char"0303} \MISSINGGLYPH{\char"2C6E\char"0303} \SUPPORTEDANCH{\char"0050\char"0303} \SUPPORTEDANCH{\char"01A4\char"0303} \SUPPORTEDANCH{\char"0052\char"0303} \SUPPORTEDANCH{\char"2C64\char"0303} \SUPPORTEDANCH{\char"0054\char"0303} \SUPPORTEDANCH{\char"01AC\char"0303} \SUPPORTEDANCH{\char"01AE\char"0303} \SUPPORTEDPREC{\char"0059\char"0303} \SUPPORTEDANCH{\char"01B3\char"0303} \SUPPORTEDPREC{\char"0061\char"0303} \SUPPORTEDANCH{\char"1E9A\char"0303} \SUPPORTEDANCH{\char"0062\char"0303} \SUPPORTEDANCH{\char"0183\char"0303} \SUPPORTEDANCH{\char"0063\char"0303} \SUPPORTEDANCH{\char"0188\char"0303} \SUPPORTEDANCH{\char"0064\char"0303} \SUPPORTEDANCH{\char"018C\char"0303} \SUPPORTEDANCH{\char"0221\char"0303} \SUPPORTEDANCH{\char"0068\char"0303} \SUPPORTEDANCH{\char"2C68\char"0303} \SUPPORTEDANCH{\char"006B\char"0303} \SUPPORTEDANCH{\char"0199\char"0303} \SUPPORTEDANCH{\char"2C6A\char"0303} \SUPPORTEDANCH{\char"006C\char"0303} \SUPPORTEDANCH{\char"0234\char"0303} \SUPPORTEDPREC{\char"006E\char"0303} \SUPPORTEDANCH{\char"0235\char"0303} \SUPPORTEDANCH{\char"0070\char"0303} \SUPPORTEDANCH{\char"01A5\char"0303} \SUPPORTEDANCH{\char"0071\char"0303} \SUPPORTEDANCH{\char"02A0\char"0303} \SUPPORTEDANCH{\char"0074\char"0303} \SUPPORTEDANCH{\char"01AB\char"0303} \SUPPORTEDANCH{\char"01AD\char"0303} \SUPPORTEDANCH{\char"0236\char"0303} \SUPPORTEDPREC{\char"0075\char"0303} \MISSINGGLYPH{\char"AB52\char"0303} \SUPPORTEDPREC{\char"0079\char"0303} \SUPPORTEDANCH{\char"01B4\char"0303} \MISSINGGLYPH{\char"1EFF\char"0303}

\SUPPORTEDANCH{\char"0048\char"0304} \SUPPORTEDANCH{\char"2C67\char"0304} \SUPPORTEDANCH{\char"004B\char"0304} \SUPPORTEDANCH{\char"2C69\char"0304} \SUPPORTEDANCH{\char"004D\char"0304} \MISSINGGLYPH{\char"2C6E\char"0304} \SUPPORTEDANCH{\char"0050\char"0304} \SUPPORTEDANCH{\char"01A4\char"0304} \SUPPORTEDANCH{\char"0052\char"0304} \SUPPORTEDANCH{\char"2C64\char"0304} \SUPPORTEDANCH{\char"0054\char"0304} \SUPPORTEDANCH{\char"01AC\char"0304} \SUPPORTEDANCH{\char"01AE\char"0304} \SUPPORTEDPREC{\char"0059\char"0304} \SUPPORTEDANCH{\char"01B3\char"0304} \SUPPORTEDPREC{\char"0061\char"0304} \SUPPORTEDANCH{\char"1E9A\char"0304} \SUPPORTEDANCH{\char"0062\char"0304} \SUPPORTEDANCH{\char"0183\char"0304} \SUPPORTEDANCH{\char"0063\char"0304} \SUPPORTEDANCH{\char"0188\char"0304} \SUPPORTEDANCH{\char"0064\char"0304} \SUPPORTEDANCH{\char"018C\char"0304} \SUPPORTEDANCH{\char"0221\char"0304} \SUPPORTEDANCH{\char"0068\char"0304} \SUPPORTEDANCH{\char"2C68\char"0304} \SUPPORTEDANCH{\char"006B\char"0304} \SUPPORTEDANCH{\char"0199\char"0304} \SUPPORTEDANCH{\char"2C6A\char"0304} \SUPPORTEDANCH{\char"006C\char"0304} \SUPPORTEDANCH{\char"0234\char"0304} \SUPPORTEDANCH{\char"006E\char"0304} \SUPPORTEDANCH{\char"0235\char"0304} \SUPPORTEDANCH{\char"0070\char"0304} \SUPPORTEDANCH{\char"01A5\char"0304} \SUPPORTEDANCH{\char"0071\char"0304} \SUPPORTEDANCH{\char"02A0\char"0304} \SUPPORTEDANCH{\char"0074\char"0304} \SUPPORTEDANCH{\char"01AB\char"0304} \SUPPORTEDANCH{\char"01AD\char"0304} \SUPPORTEDANCH{\char"0236\char"0304} \SUPPORTEDPREC{\char"0075\char"0304} \MISSINGGLYPH{\char"AB52\char"0304} \SUPPORTEDPREC{\char"0079\char"0304} \SUPPORTEDANCH{\char"01B4\char"0304} \MISSINGGLYPH{\char"1EFF\char"0304}

\SUPPORTEDANCH{\char"0048\char"0306} \SUPPORTEDANCH{\char"2C67\char"0306} \SUPPORTEDANCH{\char"004B\char"0306} \SUPPORTEDANCH{\char"2C69\char"0306} \SUPPORTEDANCH{\char"004D\char"0306} \MISSINGGLYPH{\char"2C6E\char"0306} \SUPPORTEDANCH{\char"0050\char"0306} \SUPPORTEDANCH{\char"01A4\char"0306} \SUPPORTEDANCH{\char"0052\char"0306} \SUPPORTEDANCH{\char"2C64\char"0306} \SUPPORTEDANCH{\char"0054\char"0306} \SUPPORTEDANCH{\char"01AC\char"0306} \SUPPORTEDANCH{\char"01AE\char"0306} \SUPPORTEDANCH{\char"0059\char"0306} \SUPPORTEDANCH{\char"01B3\char"0306} \SUPPORTEDPREC{\char"0061\char"0306} \SUPPORTEDANCH{\char"1E9A\char"0306} \SUPPORTEDANCH{\char"0062\char"0306} \SUPPORTEDANCH{\char"0183\char"0306} \SUPPORTEDANCH{\char"0063\char"0306} \SUPPORTEDANCH{\char"0188\char"0306} \SUPPORTEDANCH{\char"0064\char"0306} \SUPPORTEDANCH{\char"018C\char"0306} \SUPPORTEDANCH{\char"0221\char"0306} \SUPPORTEDANCH{\char"0068\char"0306} \SUPPORTEDANCH{\char"2C68\char"0306} \SUPPORTEDANCH{\char"006B\char"0306} \SUPPORTEDANCH{\char"0199\char"0306} \SUPPORTEDANCH{\char"2C6A\char"0306} \SUPPORTEDANCH{\char"006C\char"0306} \SUPPORTEDANCH{\char"0234\char"0306} \SUPPORTEDANCH{\char"006E\char"0306} \SUPPORTEDANCH{\char"0235\char"0306} \SUPPORTEDANCH{\char"0070\char"0306} \SUPPORTEDANCH{\char"01A5\char"0306} \SUPPORTEDANCH{\char"0071\char"0306} \SUPPORTEDANCH{\char"02A0\char"0306} \SUPPORTEDANCH{\char"0074\char"0306} \SUPPORTEDANCH{\char"01AB\char"0306} \SUPPORTEDANCH{\char"01AD\char"0306} \SUPPORTEDANCH{\char"0236\char"0306} \SUPPORTEDPREC{\char"0075\char"0306} \MISSINGGLYPH{\char"AB52\char"0306} \SUPPORTEDANCH{\char"0079\char"0306} \SUPPORTEDANCH{\char"01B4\char"0306} \MISSINGGLYPH{\char"1EFF\char"0306}

\SUPPORTEDPREC{\char"0048\char"0307} \SUPPORTEDANCH{\char"2C67\char"0307} \SUPPORTEDANCH{\char"004B\char"0307} \SUPPORTEDANCH{\char"2C69\char"0307} \SUPPORTEDPREC{\char"004D\char"0307} \MISSINGGLYPH{\char"2C6E\char"0307} \SUPPORTEDPREC{\char"0050\char"0307} \SUPPORTEDANCH{\char"01A4\char"0307} \SUPPORTEDPREC{\char"0052\char"0307} \SUPPORTEDANCH{\char"2C64\char"0307} \SUPPORTEDPREC{\char"0054\char"0307} \SUPPORTEDANCH{\char"01AC\char"0307} \SUPPORTEDANCH{\char"01AE\char"0307} \SUPPORTEDPREC{\char"0059\char"0307} \SUPPORTEDANCH{\char"01B3\char"0307} \SUPPORTEDPREC{\char"0061\char"0307} \SUPPORTEDANCH{\char"1E9A\char"0307} \SUPPORTEDPREC{\char"0062\char"0307} \SUPPORTEDANCH{\char"0183\char"0307} \SUPPORTEDPREC{\char"0063\char"0307} \SUPPORTEDANCH{\char"0188\char"0307} \SUPPORTEDPREC{\char"0064\char"0307} \SUPPORTEDANCH{\char"018C\char"0307} \SUPPORTEDANCH{\char"0221\char"0307} \SUPPORTEDPREC{\char"0068\char"0307} \SUPPORTEDANCH{\char"2C68\char"0307} \SUPPORTEDANCH{\char"006B\char"0307} \SUPPORTEDANCH{\char"0199\char"0307} \SUPPORTEDANCH{\char"2C6A\char"0307} \SUPPORTEDANCH{\char"006C\char"0307} \SUPPORTEDANCH{\char"0234\char"0307} \SUPPORTEDPREC{\char"006E\char"0307} \SUPPORTEDANCH{\char"0235\char"0307} \SUPPORTEDPREC{\char"0070\char"0307} \SUPPORTEDANCH{\char"01A5\char"0307} \SUPPORTEDANCH{\char"0071\char"0307} \SUPPORTEDANCH{\char"02A0\char"0307} \SUPPORTEDPREC{\char"0074\char"0307} \SUPPORTEDANCH{\char"01AB\char"0307} \SUPPORTEDANCH{\char"01AD\char"0307} \SUPPORTEDANCH{\char"0236\char"0307} \SUPPORTEDANCH{\char"0075\char"0307} \MISSINGGLYPH{\char"AB52\char"0307} \SUPPORTEDPREC{\char"0079\char"0307} \SUPPORTEDANCH{\char"01B4\char"0307} \MISSINGGLYPH{\char"1EFF\char"0307}

\SUPPORTEDPREC{\char"0048\char"0308} \SUPPORTEDANCH{\char"2C67\char"0308} \SUPPORTEDANCH{\char"004B\char"0308} \SUPPORTEDANCH{\char"2C69\char"0308} \SUPPORTEDANCH{\char"004D\char"0308} \MISSINGGLYPH{\char"2C6E\char"0308} \SUPPORTEDANCH{\char"0050\char"0308} \SUPPORTEDANCH{\char"01A4\char"0308} \SUPPORTEDANCH{\char"0052\char"0308} \SUPPORTEDANCH{\char"2C64\char"0308} \SUPPORTEDANCH{\char"0054\char"0308} \SUPPORTEDANCH{\char"01AC\char"0308} \SUPPORTEDANCH{\char"01AE\char"0308} \SUPPORTEDPREC{\char"0059\char"0308} \SUPPORTEDANCH{\char"01B3\char"0308} \SUPPORTEDPREC{\char"0061\char"0308} \SUPPORTEDANCH{\char"1E9A\char"0308} \SUPPORTEDANCH{\char"0062\char"0308} \SUPPORTEDANCH{\char"0183\char"0308} \SUPPORTEDANCH{\char"0063\char"0308} \SUPPORTEDANCH{\char"0188\char"0308} \SUPPORTEDANCH{\char"0064\char"0308} \SUPPORTEDANCH{\char"018C\char"0308} \SUPPORTEDANCH{\char"0221\char"0308} \SUPPORTEDPREC{\char"0068\char"0308} \SUPPORTEDANCH{\char"2C68\char"0308} \SUPPORTEDANCH{\char"006B\char"0308} \SUPPORTEDANCH{\char"0199\char"0308} \SUPPORTEDANCH{\char"2C6A\char"0308} \SUPPORTEDANCH{\char"006C\char"0308} \SUPPORTEDANCH{\char"0234\char"0308} \SUPPORTEDANCH{\char"006E\char"0308} \SUPPORTEDANCH{\char"0235\char"0308} \SUPPORTEDANCH{\char"0070\char"0308} \SUPPORTEDANCH{\char"01A5\char"0308} \SUPPORTEDANCH{\char"0071\char"0308} \SUPPORTEDANCH{\char"02A0\char"0308} \SUPPORTEDPREC{\char"0074\char"0308} \SUPPORTEDANCH{\char"01AB\char"0308} \SUPPORTEDANCH{\char"01AD\char"0308} \SUPPORTEDANCH{\char"0236\char"0308} \SUPPORTEDPREC{\char"0075\char"0308} \MISSINGGLYPH{\char"AB52\char"0308} \SUPPORTEDPREC{\char"0079\char"0308} \SUPPORTEDANCH{\char"01B4\char"0308} \MISSINGGLYPH{\char"1EFF\char"0308}

\SUPPORTEDANCH{\char"0048\char"0309} \SUPPORTEDANCH{\char"2C67\char"0309} \SUPPORTEDANCH{\char"004B\char"0309} \SUPPORTEDANCH{\char"2C69\char"0309} \SUPPORTEDANCH{\char"004D\char"0309} \MISSINGGLYPH{\char"2C6E\char"0309} \SUPPORTEDANCH{\char"0050\char"0309} \SUPPORTEDANCH{\char"01A4\char"0309} \SUPPORTEDANCH{\char"0052\char"0309} \SUPPORTEDANCH{\char"2C64\char"0309} \SUPPORTEDANCH{\char"0054\char"0309} \SUPPORTEDANCH{\char"01AC\char"0309} \SUPPORTEDANCH{\char"01AE\char"0309} \SUPPORTEDPREC{\char"0059\char"0309} \SUPPORTEDANCH{\char"01B3\char"0309} \SUPPORTEDPREC{\char"0061\char"0309} \SUPPORTEDANCH{\char"1E9A\char"0309} \SUPPORTEDANCH{\char"0062\char"0309} \SUPPORTEDANCH{\char"0183\char"0309} \SUPPORTEDANCH{\char"0063\char"0309} \SUPPORTEDANCH{\char"0188\char"0309} \SUPPORTEDANCH{\char"0064\char"0309} \SUPPORTEDANCH{\char"018C\char"0309} \SUPPORTEDANCH{\char"0221\char"0309} \SUPPORTEDANCH{\char"0068\char"0309} \SUPPORTEDANCH{\char"2C68\char"0309} \SUPPORTEDANCH{\char"006B\char"0309} \SUPPORTEDANCH{\char"0199\char"0309} \SUPPORTEDANCH{\char"2C6A\char"0309} \SUPPORTEDANCH{\char"006C\char"0309} \SUPPORTEDANCH{\char"0234\char"0309} \SUPPORTEDANCH{\char"006E\char"0309} \SUPPORTEDANCH{\char"0235\char"0309} \SUPPORTEDANCH{\char"0070\char"0309} \SUPPORTEDANCH{\char"01A5\char"0309} \SUPPORTEDANCH{\char"0071\char"0309} \SUPPORTEDANCH{\char"02A0\char"0309} \SUPPORTEDANCH{\char"0074\char"0309} \SUPPORTEDANCH{\char"01AB\char"0309} \SUPPORTEDANCH{\char"01AD\char"0309} \SUPPORTEDANCH{\char"0236\char"0309} \SUPPORTEDPREC{\char"0075\char"0309} \MISSINGGLYPH{\char"AB52\char"0309} \SUPPORTEDPREC{\char"0079\char"0309} \SUPPORTEDANCH{\char"01B4\char"0309} \MISSINGGLYPH{\char"1EFF\char"0309}

\SUPPORTEDANCH{\char"0048\char"030A} \SUPPORTEDANCH{\char"2C67\char"030A} \SUPPORTEDANCH{\char"004B\char"030A} \SUPPORTEDANCH{\char"2C69\char"030A} \SUPPORTEDANCH{\char"004D\char"030A} \MISSINGGLYPH{\char"2C6E\char"030A} \SUPPORTEDANCH{\char"0050\char"030A} \SUPPORTEDANCH{\char"01A4\char"030A} \SUPPORTEDANCH{\char"0052\char"030A} \SUPPORTEDANCH{\char"2C64\char"030A} \SUPPORTEDANCH{\char"0054\char"030A} \SUPPORTEDANCH{\char"01AC\char"030A} \SUPPORTEDANCH{\char"01AE\char"030A} \SUPPORTEDANCH{\char"0059\char"030A} \SUPPORTEDANCH{\char"01B3\char"030A} \SUPPORTEDPREC{\char"0061\char"030A} \SUPPORTEDANCH{\char"1E9A\char"030A} \SUPPORTEDANCH{\char"0062\char"030A} \SUPPORTEDANCH{\char"0183\char"030A} \SUPPORTEDANCH{\char"0063\char"030A} \SUPPORTEDANCH{\char"0188\char"030A} \SUPPORTEDANCH{\char"0064\char"030A} \SUPPORTEDANCH{\char"018C\char"030A} \SUPPORTEDANCH{\char"0221\char"030A} \SUPPORTEDANCH{\char"0068\char"030A} \SUPPORTEDANCH{\char"2C68\char"030A} \SUPPORTEDANCH{\char"006B\char"030A} \SUPPORTEDANCH{\char"0199\char"030A} \SUPPORTEDANCH{\char"2C6A\char"030A} \SUPPORTEDANCH{\char"006C\char"030A} \SUPPORTEDANCH{\char"0234\char"030A} \SUPPORTEDANCH{\char"006E\char"030A} \SUPPORTEDANCH{\char"0235\char"030A} \SUPPORTEDANCH{\char"0070\char"030A} \SUPPORTEDANCH{\char"01A5\char"030A} \SUPPORTEDANCH{\char"0071\char"030A} \SUPPORTEDANCH{\char"02A0\char"030A} \SUPPORTEDANCH{\char"0074\char"030A} \SUPPORTEDANCH{\char"01AB\char"030A} \SUPPORTEDANCH{\char"01AD\char"030A} \SUPPORTEDANCH{\char"0236\char"030A} \SUPPORTEDPREC{\char"0075\char"030A} \MISSINGGLYPH{\char"AB52\char"030A} \SUPPORTEDPREC{\char"0079\char"030A} \SUPPORTEDANCH{\char"01B4\char"030A} \MISSINGGLYPH{\char"1EFF\char"030A}

\SUPPORTEDANCH{\char"0048\char"030B} \SUPPORTEDANCH{\char"2C67\char"030B} \SUPPORTEDANCH{\char"004B\char"030B} \SUPPORTEDANCH{\char"2C69\char"030B} \SUPPORTEDANCH{\char"004D\char"030B} \MISSINGGLYPH{\char"2C6E\char"030B} \SUPPORTEDANCH{\char"0050\char"030B} \SUPPORTEDANCH{\char"01A4\char"030B} \SUPPORTEDANCH{\char"0052\char"030B} \SUPPORTEDANCH{\char"2C64\char"030B} \SUPPORTEDANCH{\char"0054\char"030B} \SUPPORTEDANCH{\char"01AC\char"030B} \SUPPORTEDANCH{\char"01AE\char"030B} \SUPPORTEDANCH{\char"0059\char"030B} \SUPPORTEDANCH{\char"01B3\char"030B} \SUPPORTEDANCH{\char"0061\char"030B} \SUPPORTEDANCH{\char"1E9A\char"030B} \SUPPORTEDANCH{\char"0062\char"030B} \SUPPORTEDANCH{\char"0183\char"030B} \SUPPORTEDANCH{\char"0063\char"030B} \SUPPORTEDANCH{\char"0188\char"030B} \SUPPORTEDANCH{\char"0064\char"030B} \SUPPORTEDANCH{\char"018C\char"030B} \SUPPORTEDANCH{\char"0221\char"030B} \SUPPORTEDANCH{\char"0068\char"030B} \SUPPORTEDANCH{\char"2C68\char"030B} \SUPPORTEDANCH{\char"006B\char"030B} \SUPPORTEDANCH{\char"0199\char"030B} \SUPPORTEDANCH{\char"2C6A\char"030B} \SUPPORTEDANCH{\char"006C\char"030B} \SUPPORTEDANCH{\char"0234\char"030B} \SUPPORTEDANCH{\char"006E\char"030B} \SUPPORTEDANCH{\char"0235\char"030B} \SUPPORTEDANCH{\char"0070\char"030B} \SUPPORTEDANCH{\char"01A5\char"030B} \SUPPORTEDANCH{\char"0071\char"030B} \SUPPORTEDANCH{\char"02A0\char"030B} \SUPPORTEDANCH{\char"0074\char"030B} \SUPPORTEDANCH{\char"01AB\char"030B} \SUPPORTEDANCH{\char"01AD\char"030B} \SUPPORTEDANCH{\char"0236\char"030B} \SUPPORTEDPREC{\char"0075\char"030B} \MISSINGGLYPH{\char"AB52\char"030B} \SUPPORTEDANCH{\char"0079\char"030B} \SUPPORTEDANCH{\char"01B4\char"030B} \MISSINGGLYPH{\char"1EFF\char"030B}

\SUPPORTEDPREC{\char"0048\char"030C} \SUPPORTEDANCH{\char"2C67\char"030C} \SUPPORTEDPREC{\char"004B\char"030C} \SUPPORTEDANCH{\char"2C69\char"030C} \SUPPORTEDANCH{\char"004D\char"030C} \MISSINGGLYPH{\char"2C6E\char"030C} \SUPPORTEDANCH{\char"0050\char"030C} \SUPPORTEDANCH{\char"01A4\char"030C} \SUPPORTEDPREC{\char"0052\char"030C} \SUPPORTEDANCH{\char"2C64\char"030C} \SUPPORTEDPREC{\char"0054\char"030C} \SUPPORTEDANCH{\char"01AC\char"030C} \SUPPORTEDANCH{\char"01AE\char"030C} \SUPPORTEDANCH{\char"0059\char"030C} \SUPPORTEDANCH{\char"01B3\char"030C} \SUPPORTEDPREC{\char"0061\char"030C} \SUPPORTEDANCH{\char"1E9A\char"030C} \SUPPORTEDANCH{\char"0062\char"030C} \SUPPORTEDANCH{\char"0183\char"030C} \SUPPORTEDPREC{\char"0063\char"030C} \SUPPORTEDANCH{\char"0188\char"030C} \SUPPORTEDPREC{\char"0064\char"030C} \SUPPORTEDANCH{\char"018C\char"030C} \SUPPORTEDANCH{\char"0221\char"030C} \SUPPORTEDPREC{\char"0068\char"030C} \SUPPORTEDANCH{\char"2C68\char"030C} \SUPPORTEDPREC{\char"006B\char"030C} \SUPPORTEDANCH{\char"0199\char"030C} \SUPPORTEDANCH{\char"2C6A\char"030C} \SUPPORTEDPREC{\char"006C\char"030C} \SUPPORTEDANCH{\char"0234\char"030C} \SUPPORTEDPREC{\char"006E\char"030C} \SUPPORTEDANCH{\char"0235\char"030C} \SUPPORTEDANCH{\char"0070\char"030C} \SUPPORTEDANCH{\char"01A5\char"030C} \SUPPORTEDANCH{\char"0071\char"030C} \SUPPORTEDANCH{\char"02A0\char"030C} \SUPPORTEDPREC{\char"0074\char"030C} \SUPPORTEDANCH{\char"01AB\char"030C} \SUPPORTEDANCH{\char"01AD\char"030C} \SUPPORTEDANCH{\char"0236\char"030C} \SUPPORTEDPREC{\char"0075\char"030C} \MISSINGGLYPH{\char"AB52\char"030C} \SUPPORTEDANCH{\char"0079\char"030C} \SUPPORTEDANCH{\char"01B4\char"030C} \MISSINGGLYPH{\char"1EFF\char"030C}

\SUPPORTEDANCH{\char"0048\char"030D} \SUPPORTEDANCH{\char"2C67\char"030D} \SUPPORTEDANCH{\char"004B\char"030D} \SUPPORTEDANCH{\char"2C69\char"030D} \SUPPORTEDANCH{\char"004D\char"030D} \MISSINGGLYPH{\char"2C6E\char"030D} \SUPPORTEDANCH{\char"0050\char"030D} \SUPPORTEDANCH{\char"01A4\char"030D} \SUPPORTEDANCH{\char"0052\char"030D} \SUPPORTEDANCH{\char"2C64\char"030D} \SUPPORTEDANCH{\char"0054\char"030D} \SUPPORTEDANCH{\char"01AC\char"030D} \SUPPORTEDANCH{\char"01AE\char"030D} \SUPPORTEDANCH{\char"0059\char"030D} \SUPPORTEDANCH{\char"01B3\char"030D} \SUPPORTEDANCH{\char"0061\char"030D} \SUPPORTEDANCH{\char"1E9A\char"030D} \SUPPORTEDANCH{\char"0062\char"030D} \SUPPORTEDANCH{\char"0183\char"030D} \SUPPORTEDANCH{\char"0063\char"030D} \SUPPORTEDANCH{\char"0188\char"030D} \SUPPORTEDANCH{\char"0064\char"030D} \SUPPORTEDANCH{\char"018C\char"030D} \SUPPORTEDANCH{\char"0221\char"030D} \SUPPORTEDANCH{\char"0068\char"030D} \SUPPORTEDANCH{\char"2C68\char"030D} \SUPPORTEDANCH{\char"006B\char"030D} \SUPPORTEDANCH{\char"0199\char"030D} \SUPPORTEDANCH{\char"2C6A\char"030D} \SUPPORTEDANCH{\char"006C\char"030D} \SUPPORTEDANCH{\char"0234\char"030D} \SUPPORTEDANCH{\char"006E\char"030D} \SUPPORTEDANCH{\char"0235\char"030D} \SUPPORTEDANCH{\char"0070\char"030D} \SUPPORTEDANCH{\char"01A5\char"030D} \SUPPORTEDANCH{\char"0071\char"030D} \SUPPORTEDANCH{\char"02A0\char"030D} \SUPPORTEDANCH{\char"0074\char"030D} \SUPPORTEDANCH{\char"01AB\char"030D} \SUPPORTEDANCH{\char"01AD\char"030D} \SUPPORTEDANCH{\char"0236\char"030D} \SUPPORTEDANCH{\char"0075\char"030D} \MISSINGGLYPH{\char"AB52\char"030D} \SUPPORTEDANCH{\char"0079\char"030D} \SUPPORTEDANCH{\char"01B4\char"030D} \MISSINGGLYPH{\char"1EFF\char"030D}

\SUPPORTEDANCH{\char"0048\char"030F} \SUPPORTEDANCH{\char"2C67\char"030F} \SUPPORTEDANCH{\char"004B\char"030F} \SUPPORTEDANCH{\char"2C69\char"030F} \SUPPORTEDANCH{\char"004D\char"030F} \MISSINGGLYPH{\char"2C6E\char"030F} \SUPPORTEDANCH{\char"0050\char"030F} \SUPPORTEDANCH{\char"01A4\char"030F} \SUPPORTEDPREC{\char"0052\char"030F} \SUPPORTEDANCH{\char"2C64\char"030F} \SUPPORTEDANCH{\char"0054\char"030F} \SUPPORTEDANCH{\char"01AC\char"030F} \SUPPORTEDANCH{\char"01AE\char"030F} \SUPPORTEDANCH{\char"0059\char"030F} \SUPPORTEDANCH{\char"01B3\char"030F} \SUPPORTEDPREC{\char"0061\char"030F} \SUPPORTEDANCH{\char"1E9A\char"030F} \SUPPORTEDANCH{\char"0062\char"030F} \SUPPORTEDANCH{\char"0183\char"030F} \SUPPORTEDANCH{\char"0063\char"030F} \SUPPORTEDANCH{\char"0188\char"030F} \SUPPORTEDANCH{\char"0064\char"030F} \SUPPORTEDANCH{\char"018C\char"030F} \SUPPORTEDANCH{\char"0221\char"030F} \SUPPORTEDANCH{\char"0068\char"030F} \SUPPORTEDANCH{\char"2C68\char"030F} \SUPPORTEDANCH{\char"006B\char"030F} \SUPPORTEDANCH{\char"0199\char"030F} \SUPPORTEDANCH{\char"2C6A\char"030F} \SUPPORTEDANCH{\char"006C\char"030F} \SUPPORTEDANCH{\char"0234\char"030F} \SUPPORTEDANCH{\char"006E\char"030F} \SUPPORTEDANCH{\char"0235\char"030F} \SUPPORTEDANCH{\char"0070\char"030F} \SUPPORTEDANCH{\char"01A5\char"030F} \SUPPORTEDANCH{\char"0071\char"030F} \SUPPORTEDANCH{\char"02A0\char"030F} \SUPPORTEDANCH{\char"0074\char"030F} \SUPPORTEDANCH{\char"01AB\char"030F} \SUPPORTEDANCH{\char"01AD\char"030F} \SUPPORTEDANCH{\char"0236\char"030F} \SUPPORTEDPREC{\char"0075\char"030F} \MISSINGGLYPH{\char"AB52\char"030F} \SUPPORTEDANCH{\char"0079\char"030F} \SUPPORTEDANCH{\char"01B4\char"030F} \MISSINGGLYPH{\char"1EFF\char"030F}

\SUPPORTEDANCH{\char"0048\char"0310} \SUPPORTEDANCH{\char"2C67\char"0310} \SUPPORTEDANCH{\char"004B\char"0310} \SUPPORTEDANCH{\char"2C69\char"0310} \SUPPORTEDANCH{\char"004D\char"0310} \MISSINGGLYPH{\char"2C6E\char"0310} \SUPPORTEDANCH{\char"0050\char"0310} \SUPPORTEDANCH{\char"01A4\char"0310} \SUPPORTEDANCH{\char"0052\char"0310} \SUPPORTEDANCH{\char"2C64\char"0310} \SUPPORTEDANCH{\char"0054\char"0310} \SUPPORTEDANCH{\char"01AC\char"0310} \SUPPORTEDANCH{\char"01AE\char"0310} \SUPPORTEDANCH{\char"0059\char"0310} \SUPPORTEDANCH{\char"01B3\char"0310} \SUPPORTEDANCH{\char"0061\char"0310} \SUPPORTEDANCH{\char"1E9A\char"0310} \SUPPORTEDANCH{\char"0062\char"0310} \SUPPORTEDANCH{\char"0183\char"0310} \SUPPORTEDANCH{\char"0063\char"0310} \SUPPORTEDANCH{\char"0188\char"0310} \SUPPORTEDANCH{\char"0064\char"0310} \SUPPORTEDANCH{\char"018C\char"0310} \SUPPORTEDANCH{\char"0221\char"0310} \SUPPORTEDANCH{\char"0068\char"0310} \SUPPORTEDANCH{\char"2C68\char"0310} \SUPPORTEDANCH{\char"006B\char"0310} \SUPPORTEDANCH{\char"0199\char"0310} \SUPPORTEDANCH{\char"2C6A\char"0310} \SUPPORTEDANCH{\char"006C\char"0310} \SUPPORTEDANCH{\char"0234\char"0310} \SUPPORTEDANCH{\char"006E\char"0310} \SUPPORTEDANCH{\char"0235\char"0310} \SUPPORTEDANCH{\char"0070\char"0310} \SUPPORTEDANCH{\char"01A5\char"0310} \SUPPORTEDANCH{\char"0071\char"0310} \SUPPORTEDANCH{\char"02A0\char"0310} \SUPPORTEDANCH{\char"0074\char"0310} \SUPPORTEDANCH{\char"01AB\char"0310} \SUPPORTEDANCH{\char"01AD\char"0310} \SUPPORTEDANCH{\char"0236\char"0310} \SUPPORTEDANCH{\char"0075\char"0310} \MISSINGGLYPH{\char"AB52\char"0310} \SUPPORTEDANCH{\char"0079\char"0310} \SUPPORTEDANCH{\char"01B4\char"0310} \MISSINGGLYPH{\char"1EFF\char"0310}

\SUPPORTEDANCH{\char"0048\char"0311} \SUPPORTEDANCH{\char"2C67\char"0311} \SUPPORTEDANCH{\char"004B\char"0311} \SUPPORTEDANCH{\char"2C69\char"0311} \SUPPORTEDANCH{\char"004D\char"0311} \MISSINGGLYPH{\char"2C6E\char"0311} \SUPPORTEDANCH{\char"0050\char"0311} \SUPPORTEDANCH{\char"01A4\char"0311} \SUPPORTEDPREC{\char"0052\char"0311} \SUPPORTEDANCH{\char"2C64\char"0311} \SUPPORTEDANCH{\char"0054\char"0311} \SUPPORTEDANCH{\char"01AC\char"0311} \SUPPORTEDANCH{\char"01AE\char"0311} \SUPPORTEDANCH{\char"0059\char"0311} \SUPPORTEDANCH{\char"01B3\char"0311} \SUPPORTEDPREC{\char"0061\char"0311} \SUPPORTEDANCH{\char"1E9A\char"0311} \SUPPORTEDANCH{\char"0062\char"0311} \SUPPORTEDANCH{\char"0183\char"0311} \SUPPORTEDANCH{\char"0063\char"0311} \SUPPORTEDANCH{\char"0188\char"0311} \SUPPORTEDANCH{\char"0064\char"0311} \SUPPORTEDANCH{\char"018C\char"0311} \SUPPORTEDANCH{\char"0221\char"0311} \SUPPORTEDANCH{\char"0068\char"0311} \SUPPORTEDANCH{\char"2C68\char"0311} \SUPPORTEDANCH{\char"006B\char"0311} \SUPPORTEDANCH{\char"0199\char"0311} \SUPPORTEDANCH{\char"2C6A\char"0311} \SUPPORTEDANCH{\char"006C\char"0311} \SUPPORTEDANCH{\char"0234\char"0311} \SUPPORTEDANCH{\char"006E\char"0311} \SUPPORTEDANCH{\char"0235\char"0311} \SUPPORTEDANCH{\char"0070\char"0311} \SUPPORTEDANCH{\char"01A5\char"0311} \SUPPORTEDANCH{\char"0071\char"0311} \SUPPORTEDANCH{\char"02A0\char"0311} \SUPPORTEDANCH{\char"0074\char"0311} \SUPPORTEDANCH{\char"01AB\char"0311} \SUPPORTEDANCH{\char"01AD\char"0311} \SUPPORTEDANCH{\char"0236\char"0311} \SUPPORTEDPREC{\char"0075\char"0311} \MISSINGGLYPH{\char"AB52\char"0311} \SUPPORTEDANCH{\char"0079\char"0311} \SUPPORTEDANCH{\char"01B4\char"0311} \MISSINGGLYPH{\char"1EFF\char"0311}

\SUPPORTEDANCH{\char"0048\char"0312} \SUPPORTEDANCH{\char"2C67\char"0312} \SUPPORTEDANCH{\char"004B\char"0312} \SUPPORTEDANCH{\char"2C69\char"0312} \SUPPORTEDANCH{\char"004D\char"0312} \MISSINGGLYPH{\char"2C6E\char"0312} \SUPPORTEDANCH{\char"0050\char"0312} \SUPPORTEDANCH{\char"01A4\char"0312} \SUPPORTEDANCH{\char"0052\char"0312} \SUPPORTEDANCH{\char"2C64\char"0312} \SUPPORTEDANCH{\char"0054\char"0312} \SUPPORTEDANCH{\char"01AC\char"0312} \SUPPORTEDANCH{\char"01AE\char"0312} \SUPPORTEDANCH{\char"0059\char"0312} \SUPPORTEDANCH{\char"01B3\char"0312} \SUPPORTEDANCH{\char"0061\char"0312} \SUPPORTEDANCH{\char"1E9A\char"0312} \SUPPORTEDANCH{\char"0062\char"0312} \SUPPORTEDANCH{\char"0183\char"0312} \SUPPORTEDANCH{\char"0063\char"0312} \SUPPORTEDANCH{\char"0188\char"0312} \SUPPORTEDANCH{\char"0064\char"0312} \SUPPORTEDANCH{\char"018C\char"0312} \SUPPORTEDANCH{\char"0221\char"0312} \SUPPORTEDANCH{\char"0068\char"0312} \SUPPORTEDANCH{\char"2C68\char"0312} \SUPPORTEDANCH{\char"006B\char"0312} \SUPPORTEDANCH{\char"0199\char"0312} \SUPPORTEDANCH{\char"2C6A\char"0312} \SUPPORTEDANCH{\char"006C\char"0312} \SUPPORTEDANCH{\char"0234\char"0312} \SUPPORTEDANCH{\char"006E\char"0312} \SUPPORTEDANCH{\char"0235\char"0312} \SUPPORTEDANCH{\char"0070\char"0312} \SUPPORTEDANCH{\char"01A5\char"0312} \SUPPORTEDANCH{\char"0071\char"0312} \SUPPORTEDANCH{\char"02A0\char"0312} \SUPPORTEDANCH{\char"0074\char"0312} \SUPPORTEDANCH{\char"01AB\char"0312} \SUPPORTEDANCH{\char"01AD\char"0312} \SUPPORTEDANCH{\char"0236\char"0312} \SUPPORTEDANCH{\char"0075\char"0312} \MISSINGGLYPH{\char"AB52\char"0312} \SUPPORTEDANCH{\char"0079\char"0312} \SUPPORTEDANCH{\char"01B4\char"0312} \MISSINGGLYPH{\char"1EFF\char"0312}

\SUPPORTEDANCH{\char"0048\char"033D} \SUPPORTEDANCH{\char"2C67\char"033D} \SUPPORTEDANCH{\char"004B\char"033D} \SUPPORTEDANCH{\char"2C69\char"033D} \SUPPORTEDANCH{\char"004D\char"033D} \MISSINGGLYPH{\char"2C6E\char"033D} \SUPPORTEDANCH{\char"0050\char"033D} \SUPPORTEDANCH{\char"01A4\char"033D} \SUPPORTEDANCH{\char"0052\char"033D} \SUPPORTEDANCH{\char"2C64\char"033D} \SUPPORTEDANCH{\char"0054\char"033D} \SUPPORTEDANCH{\char"01AC\char"033D} \SUPPORTEDANCH{\char"01AE\char"033D} \SUPPORTEDANCH{\char"0059\char"033D} \SUPPORTEDANCH{\char"01B3\char"033D} \SUPPORTEDANCH{\char"0061\char"033D} \SUPPORTEDANCH{\char"1E9A\char"033D} \SUPPORTEDANCH{\char"0062\char"033D} \SUPPORTEDANCH{\char"0183\char"033D} \SUPPORTEDANCH{\char"0063\char"033D} \SUPPORTEDANCH{\char"0188\char"033D} \SUPPORTEDANCH{\char"0064\char"033D} \SUPPORTEDANCH{\char"018C\char"033D} \SUPPORTEDANCH{\char"0221\char"033D} \SUPPORTEDANCH{\char"0068\char"033D} \SUPPORTEDANCH{\char"2C68\char"033D} \SUPPORTEDANCH{\char"006B\char"033D} \SUPPORTEDANCH{\char"0199\char"033D} \SUPPORTEDANCH{\char"2C6A\char"033D} \SUPPORTEDANCH{\char"006C\char"033D} \SUPPORTEDANCH{\char"0234\char"033D} \SUPPORTEDANCH{\char"006E\char"033D} \SUPPORTEDANCH{\char"0235\char"033D} \SUPPORTEDANCH{\char"0070\char"033D} \SUPPORTEDANCH{\char"01A5\char"033D} \SUPPORTEDANCH{\char"0071\char"033D} \SUPPORTEDANCH{\char"02A0\char"033D} \SUPPORTEDANCH{\char"0074\char"033D} \SUPPORTEDANCH{\char"01AB\char"033D} \SUPPORTEDANCH{\char"01AD\char"033D} \SUPPORTEDANCH{\char"0236\char"033D} \SUPPORTEDANCH{\char"0075\char"033D} \MISSINGGLYPH{\char"AB52\char"033D} \SUPPORTEDANCH{\char"0079\char"033D} \SUPPORTEDANCH{\char"01B4\char"033D} \MISSINGGLYPH{\char"1EFF\char"033D}

\SUPPORTEDANCH{\char"0048\char"033E} \SUPPORTEDANCH{\char"2C67\char"033E} \SUPPORTEDANCH{\char"004B\char"033E} \SUPPORTEDANCH{\char"2C69\char"033E} \SUPPORTEDANCH{\char"004D\char"033E} \MISSINGGLYPH{\char"2C6E\char"033E} \SUPPORTEDANCH{\char"0050\char"033E} \SUPPORTEDANCH{\char"01A4\char"033E} \SUPPORTEDANCH{\char"0052\char"033E} \SUPPORTEDANCH{\char"2C64\char"033E} \SUPPORTEDANCH{\char"0054\char"033E} \SUPPORTEDANCH{\char"01AC\char"033E} \SUPPORTEDANCH{\char"01AE\char"033E} \SUPPORTEDANCH{\char"0059\char"033E} \SUPPORTEDANCH{\char"01B3\char"033E} \SUPPORTEDANCH{\char"0061\char"033E} \SUPPORTEDANCH{\char"1E9A\char"033E} \SUPPORTEDANCH{\char"0062\char"033E} \SUPPORTEDANCH{\char"0183\char"033E} \SUPPORTEDANCH{\char"0063\char"033E} \SUPPORTEDANCH{\char"0188\char"033E} \SUPPORTEDANCH{\char"0064\char"033E} \SUPPORTEDANCH{\char"018C\char"033E} \SUPPORTEDANCH{\char"0221\char"033E} \SUPPORTEDANCH{\char"0068\char"033E} \SUPPORTEDANCH{\char"2C68\char"033E} \SUPPORTEDANCH{\char"006B\char"033E} \SUPPORTEDANCH{\char"0199\char"033E} \SUPPORTEDANCH{\char"2C6A\char"033E} \SUPPORTEDANCH{\char"006C\char"033E} \SUPPORTEDANCH{\char"0234\char"033E} \SUPPORTEDANCH{\char"006E\char"033E} \SUPPORTEDANCH{\char"0235\char"033E} \SUPPORTEDANCH{\char"0070\char"033E} \SUPPORTEDANCH{\char"01A5\char"033E} \SUPPORTEDANCH{\char"0071\char"033E} \SUPPORTEDANCH{\char"02A0\char"033E} \SUPPORTEDANCH{\char"0074\char"033E} \SUPPORTEDANCH{\char"01AB\char"033E} \SUPPORTEDANCH{\char"01AD\char"033E} \SUPPORTEDANCH{\char"0236\char"033E} \SUPPORTEDANCH{\char"0075\char"033E} \MISSINGGLYPH{\char"AB52\char"033E} \SUPPORTEDANCH{\char"0079\char"033E} \SUPPORTEDANCH{\char"01B4\char"033E} \MISSINGGLYPH{\char"1EFF\char"033E}

\SUPPORTEDANCH{\char"0048\char"033F} \SUPPORTEDANCH{\char"2C67\char"033F} \SUPPORTEDANCH{\char"004B\char"033F} \SUPPORTEDANCH{\char"2C69\char"033F} \SUPPORTEDANCH{\char"004D\char"033F} \MISSINGGLYPH{\char"2C6E\char"033F} \SUPPORTEDANCH{\char"0050\char"033F} \SUPPORTEDANCH{\char"01A4\char"033F} \SUPPORTEDANCH{\char"0052\char"033F} \SUPPORTEDANCH{\char"2C64\char"033F} \SUPPORTEDANCH{\char"0054\char"033F} \SUPPORTEDANCH{\char"01AC\char"033F} \SUPPORTEDANCH{\char"01AE\char"033F} \SUPPORTEDANCH{\char"0059\char"033F} \SUPPORTEDANCH{\char"01B3\char"033F} \SUPPORTEDANCH{\char"0061\char"033F} \SUPPORTEDANCH{\char"1E9A\char"033F} \SUPPORTEDANCH{\char"0062\char"033F} \SUPPORTEDANCH{\char"0183\char"033F} \SUPPORTEDANCH{\char"0063\char"033F} \SUPPORTEDANCH{\char"0188\char"033F} \SUPPORTEDANCH{\char"0064\char"033F} \SUPPORTEDANCH{\char"018C\char"033F} \SUPPORTEDANCH{\char"0221\char"033F} \SUPPORTEDANCH{\char"0068\char"033F} \SUPPORTEDANCH{\char"2C68\char"033F} \SUPPORTEDANCH{\char"006B\char"033F} \SUPPORTEDANCH{\char"0199\char"033F} \SUPPORTEDANCH{\char"2C6A\char"033F} \SUPPORTEDANCH{\char"006C\char"033F} \SUPPORTEDANCH{\char"0234\char"033F} \SUPPORTEDANCH{\char"006E\char"033F} \SUPPORTEDANCH{\char"0235\char"033F} \SUPPORTEDANCH{\char"0070\char"033F} \SUPPORTEDANCH{\char"01A5\char"033F} \SUPPORTEDANCH{\char"0071\char"033F} \SUPPORTEDANCH{\char"02A0\char"033F} \SUPPORTEDANCH{\char"0074\char"033F} \SUPPORTEDANCH{\char"01AB\char"033F} \SUPPORTEDANCH{\char"01AD\char"033F} \SUPPORTEDANCH{\char"0236\char"033F} \SUPPORTEDANCH{\char"0075\char"033F} \MISSINGGLYPH{\char"AB52\char"033F} \SUPPORTEDANCH{\char"0079\char"033F} \SUPPORTEDANCH{\char"01B4\char"033F} \MISSINGGLYPH{\char"1EFF\char"033F}

\SUPPORTEDANCH{\char"0048\char"030E} \SUPPORTEDANCH{\char"2C67\char"030E} \SUPPORTEDANCH{\char"004B\char"030E} \SUPPORTEDANCH{\char"2C69\char"030E} \SUPPORTEDANCH{\char"004D\char"030E} \MISSINGGLYPH{\char"2C6E\char"030E} \SUPPORTEDANCH{\char"0050\char"030E} \SUPPORTEDANCH{\char"01A4\char"030E} \SUPPORTEDANCH{\char"0052\char"030E} \SUPPORTEDANCH{\char"2C64\char"030E} \SUPPORTEDANCH{\char"0054\char"030E} \SUPPORTEDANCH{\char"01AC\char"030E} \SUPPORTEDANCH{\char"01AE\char"030E} \SUPPORTEDANCH{\char"0059\char"030E} \SUPPORTEDANCH{\char"01B3\char"030E} \SUPPORTEDANCH{\char"0061\char"030E} \SUPPORTEDANCH{\char"1E9A\char"030E} \SUPPORTEDANCH{\char"0062\char"030E} \SUPPORTEDANCH{\char"0183\char"030E} \SUPPORTEDANCH{\char"0063\char"030E} \SUPPORTEDANCH{\char"0188\char"030E} \SUPPORTEDANCH{\char"0064\char"030E} \SUPPORTEDANCH{\char"018C\char"030E} \SUPPORTEDANCH{\char"0221\char"030E} \SUPPORTEDANCH{\char"0068\char"030E} \SUPPORTEDANCH{\char"2C68\char"030E} \SUPPORTEDANCH{\char"006B\char"030E} \SUPPORTEDANCH{\char"0199\char"030E} \SUPPORTEDANCH{\char"2C6A\char"030E} \SUPPORTEDANCH{\char"006C\char"030E} \SUPPORTEDANCH{\char"0234\char"030E} \SUPPORTEDANCH{\char"006E\char"030E} \SUPPORTEDANCH{\char"0235\char"030E} \SUPPORTEDANCH{\char"0070\char"030E} \SUPPORTEDANCH{\char"01A5\char"030E} \SUPPORTEDANCH{\char"0071\char"030E} \SUPPORTEDANCH{\char"02A0\char"030E} \SUPPORTEDANCH{\char"0074\char"030E} \SUPPORTEDANCH{\char"01AB\char"030E} \SUPPORTEDANCH{\char"01AD\char"030E} \SUPPORTEDANCH{\char"0236\char"030E} \SUPPORTEDANCH{\char"0075\char"030E} \MISSINGGLYPH{\char"AB52\char"030E} \SUPPORTEDANCH{\char"0079\char"030E} \SUPPORTEDANCH{\char"01B4\char"030E} \MISSINGGLYPH{\char"1EFF\char"030E}

\SUPPORTEDANCH{\char"0048\char"0346} \SUPPORTEDANCH{\char"2C67\char"0346} \SUPPORTEDANCH{\char"004B\char"0346} \SUPPORTEDANCH{\char"2C69\char"0346} \SUPPORTEDANCH{\char"004D\char"0346} \MISSINGGLYPH{\char"2C6E\char"0346} \SUPPORTEDANCH{\char"0050\char"0346} \SUPPORTEDANCH{\char"01A4\char"0346} \SUPPORTEDANCH{\char"0052\char"0346} \SUPPORTEDANCH{\char"2C64\char"0346} \SUPPORTEDANCH{\char"0054\char"0346} \SUPPORTEDANCH{\char"01AC\char"0346} \SUPPORTEDANCH{\char"01AE\char"0346} \SUPPORTEDANCH{\char"0059\char"0346} \SUPPORTEDANCH{\char"01B3\char"0346} \SUPPORTEDANCH{\char"0061\char"0346} \SUPPORTEDANCH{\char"1E9A\char"0346} \SUPPORTEDANCH{\char"0062\char"0346} \SUPPORTEDANCH{\char"0183\char"0346} \SUPPORTEDANCH{\char"0063\char"0346} \SUPPORTEDANCH{\char"0188\char"0346} \SUPPORTEDANCH{\char"0064\char"0346} \SUPPORTEDANCH{\char"018C\char"0346} \SUPPORTEDANCH{\char"0221\char"0346} \SUPPORTEDANCH{\char"0068\char"0346} \SUPPORTEDANCH{\char"2C68\char"0346} \SUPPORTEDANCH{\char"006B\char"0346} \SUPPORTEDANCH{\char"0199\char"0346} \SUPPORTEDANCH{\char"2C6A\char"0346} \SUPPORTEDANCH{\char"006C\char"0346} \SUPPORTEDANCH{\char"0234\char"0346} \SUPPORTEDANCH{\char"006E\char"0346} \SUPPORTEDANCH{\char"0235\char"0346} \SUPPORTEDANCH{\char"0070\char"0346} \SUPPORTEDANCH{\char"01A5\char"0346} \SUPPORTEDANCH{\char"0071\char"0346} \SUPPORTEDANCH{\char"02A0\char"0346} \SUPPORTEDANCH{\char"0074\char"0346} \SUPPORTEDANCH{\char"01AB\char"0346} \SUPPORTEDANCH{\char"01AD\char"0346} \SUPPORTEDANCH{\char"0236\char"0346} \SUPPORTEDANCH{\char"0075\char"0346} \MISSINGGLYPH{\char"AB52\char"0346} \SUPPORTEDANCH{\char"0079\char"0346} \SUPPORTEDANCH{\char"01B4\char"0346} \MISSINGGLYPH{\char"1EFF\char"0346}

\SUPPORTEDANCH{\char"0048\char"034A} \SUPPORTEDANCH{\char"2C67\char"034A} \SUPPORTEDANCH{\char"004B\char"034A} \SUPPORTEDANCH{\char"2C69\char"034A} \SUPPORTEDANCH{\char"004D\char"034A} \MISSINGGLYPH{\char"2C6E\char"034A} \SUPPORTEDANCH{\char"0050\char"034A} \SUPPORTEDANCH{\char"01A4\char"034A} \SUPPORTEDANCH{\char"0052\char"034A} \SUPPORTEDANCH{\char"2C64\char"034A} \SUPPORTEDANCH{\char"0054\char"034A} \SUPPORTEDANCH{\char"01AC\char"034A} \SUPPORTEDANCH{\char"01AE\char"034A} \SUPPORTEDANCH{\char"0059\char"034A} \SUPPORTEDANCH{\char"01B3\char"034A} \SUPPORTEDANCH{\char"0061\char"034A} \SUPPORTEDANCH{\char"1E9A\char"034A} \SUPPORTEDANCH{\char"0062\char"034A} \SUPPORTEDANCH{\char"0183\char"034A} \SUPPORTEDANCH{\char"0063\char"034A} \SUPPORTEDANCH{\char"0188\char"034A} \SUPPORTEDANCH{\char"0064\char"034A} \SUPPORTEDANCH{\char"018C\char"034A} \SUPPORTEDANCH{\char"0221\char"034A} \SUPPORTEDANCH{\char"0068\char"034A} \SUPPORTEDANCH{\char"2C68\char"034A} \SUPPORTEDANCH{\char"006B\char"034A} \SUPPORTEDANCH{\char"0199\char"034A} \SUPPORTEDANCH{\char"2C6A\char"034A} \SUPPORTEDANCH{\char"006C\char"034A} \SUPPORTEDANCH{\char"0234\char"034A} \SUPPORTEDANCH{\char"006E\char"034A} \SUPPORTEDANCH{\char"0235\char"034A} \SUPPORTEDANCH{\char"0070\char"034A} \SUPPORTEDANCH{\char"01A5\char"034A} \SUPPORTEDANCH{\char"0071\char"034A} \SUPPORTEDANCH{\char"02A0\char"034A} \SUPPORTEDANCH{\char"0074\char"034A} \SUPPORTEDANCH{\char"01AB\char"034A} \SUPPORTEDANCH{\char"01AD\char"034A} \SUPPORTEDANCH{\char"0236\char"034A} \SUPPORTEDANCH{\char"0075\char"034A} \MISSINGGLYPH{\char"AB52\char"034A} \SUPPORTEDANCH{\char"0079\char"034A} \SUPPORTEDANCH{\char"01B4\char"034A} \MISSINGGLYPH{\char"1EFF\char"034A}

\SUPPORTEDANCH{\char"0048\char"034B} \SUPPORTEDANCH{\char"2C67\char"034B} \SUPPORTEDANCH{\char"004B\char"034B} \SUPPORTEDANCH{\char"2C69\char"034B} \SUPPORTEDANCH{\char"004D\char"034B} \MISSINGGLYPH{\char"2C6E\char"034B} \SUPPORTEDANCH{\char"0050\char"034B} \SUPPORTEDANCH{\char"01A4\char"034B} \SUPPORTEDANCH{\char"0052\char"034B} \SUPPORTEDANCH{\char"2C64\char"034B} \SUPPORTEDANCH{\char"0054\char"034B} \SUPPORTEDANCH{\char"01AC\char"034B} \SUPPORTEDANCH{\char"01AE\char"034B} \SUPPORTEDANCH{\char"0059\char"034B} \SUPPORTEDANCH{\char"01B3\char"034B} \SUPPORTEDANCH{\char"0061\char"034B} \SUPPORTEDANCH{\char"1E9A\char"034B} \SUPPORTEDANCH{\char"0062\char"034B} \SUPPORTEDANCH{\char"0183\char"034B} \SUPPORTEDANCH{\char"0063\char"034B} \SUPPORTEDANCH{\char"0188\char"034B} \SUPPORTEDANCH{\char"0064\char"034B} \SUPPORTEDANCH{\char"018C\char"034B} \SUPPORTEDANCH{\char"0221\char"034B} \SUPPORTEDANCH{\char"0068\char"034B} \SUPPORTEDANCH{\char"2C68\char"034B} \SUPPORTEDANCH{\char"006B\char"034B} \SUPPORTEDANCH{\char"0199\char"034B} \SUPPORTEDANCH{\char"2C6A\char"034B} \SUPPORTEDANCH{\char"006C\char"034B} \SUPPORTEDANCH{\char"0234\char"034B} \SUPPORTEDANCH{\char"006E\char"034B} \SUPPORTEDANCH{\char"0235\char"034B} \SUPPORTEDANCH{\char"0070\char"034B} \SUPPORTEDANCH{\char"01A5\char"034B} \SUPPORTEDANCH{\char"0071\char"034B} \SUPPORTEDANCH{\char"02A0\char"034B} \SUPPORTEDANCH{\char"0074\char"034B} \SUPPORTEDANCH{\char"01AB\char"034B} \SUPPORTEDANCH{\char"01AD\char"034B} \SUPPORTEDANCH{\char"0236\char"034B} \SUPPORTEDANCH{\char"0075\char"034B} \MISSINGGLYPH{\char"AB52\char"034B} \SUPPORTEDANCH{\char"0079\char"034B} \SUPPORTEDANCH{\char"01B4\char"034B} \MISSINGGLYPH{\char"1EFF\char"034B}

\SUPPORTEDANCH{\char"0048\char"034C} \SUPPORTEDANCH{\char"2C67\char"034C} \SUPPORTEDANCH{\char"004B\char"034C} \SUPPORTEDANCH{\char"2C69\char"034C} \SUPPORTEDANCH{\char"004D\char"034C} \MISSINGGLYPH{\char"2C6E\char"034C} \SUPPORTEDANCH{\char"0050\char"034C} \SUPPORTEDANCH{\char"01A4\char"034C} \SUPPORTEDANCH{\char"0052\char"034C} \SUPPORTEDANCH{\char"2C64\char"034C} \SUPPORTEDANCH{\char"0054\char"034C} \SUPPORTEDANCH{\char"01AC\char"034C} \SUPPORTEDANCH{\char"01AE\char"034C} \SUPPORTEDANCH{\char"0059\char"034C} \SUPPORTEDANCH{\char"01B3\char"034C} \SUPPORTEDANCH{\char"0061\char"034C} \SUPPORTEDANCH{\char"1E9A\char"034C} \SUPPORTEDANCH{\char"0062\char"034C} \SUPPORTEDANCH{\char"0183\char"034C} \SUPPORTEDANCH{\char"0063\char"034C} \SUPPORTEDANCH{\char"0188\char"034C} \SUPPORTEDANCH{\char"0064\char"034C} \SUPPORTEDANCH{\char"018C\char"034C} \SUPPORTEDANCH{\char"0221\char"034C} \SUPPORTEDANCH{\char"0068\char"034C} \SUPPORTEDANCH{\char"2C68\char"034C} \SUPPORTEDANCH{\char"006B\char"034C} \SUPPORTEDANCH{\char"0199\char"034C} \SUPPORTEDANCH{\char"2C6A\char"034C} \SUPPORTEDANCH{\char"006C\char"034C} \SUPPORTEDANCH{\char"0234\char"034C} \SUPPORTEDANCH{\char"006E\char"034C} \SUPPORTEDANCH{\char"0235\char"034C} \SUPPORTEDANCH{\char"0070\char"034C} \SUPPORTEDANCH{\char"01A5\char"034C} \SUPPORTEDANCH{\char"0071\char"034C} \SUPPORTEDANCH{\char"02A0\char"034C} \SUPPORTEDANCH{\char"0074\char"034C} \SUPPORTEDANCH{\char"01AB\char"034C} \SUPPORTEDANCH{\char"01AD\char"034C} \SUPPORTEDANCH{\char"0236\char"034C} \SUPPORTEDANCH{\char"0075\char"034C} \MISSINGGLYPH{\char"AB52\char"034C} \SUPPORTEDANCH{\char"0079\char"034C} \SUPPORTEDANCH{\char"01B4\char"034C} \MISSINGGLYPH{\char"1EFF\char"034C}

\SUPPORTEDANCH{\char"0048\char"0350} \SUPPORTEDANCH{\char"2C67\char"0350} \SUPPORTEDANCH{\char"004B\char"0350} \SUPPORTEDANCH{\char"2C69\char"0350} \SUPPORTEDANCH{\char"004D\char"0350} \MISSINGGLYPH{\char"2C6E\char"0350} \SUPPORTEDANCH{\char"0050\char"0350} \SUPPORTEDANCH{\char"01A4\char"0350} \SUPPORTEDANCH{\char"0052\char"0350} \SUPPORTEDANCH{\char"2C64\char"0350} \SUPPORTEDANCH{\char"0054\char"0350} \SUPPORTEDANCH{\char"01AC\char"0350} \SUPPORTEDANCH{\char"01AE\char"0350} \SUPPORTEDANCH{\char"0059\char"0350} \SUPPORTEDANCH{\char"01B3\char"0350} \SUPPORTEDANCH{\char"0061\char"0350} \SUPPORTEDANCH{\char"1E9A\char"0350} \SUPPORTEDANCH{\char"0062\char"0350} \SUPPORTEDANCH{\char"0183\char"0350} \SUPPORTEDANCH{\char"0063\char"0350} \SUPPORTEDANCH{\char"0188\char"0350} \SUPPORTEDANCH{\char"0064\char"0350} \SUPPORTEDANCH{\char"018C\char"0350} \SUPPORTEDANCH{\char"0221\char"0350} \SUPPORTEDANCH{\char"0068\char"0350} \SUPPORTEDANCH{\char"2C68\char"0350} \SUPPORTEDANCH{\char"006B\char"0350} \SUPPORTEDANCH{\char"0199\char"0350} \SUPPORTEDANCH{\char"2C6A\char"0350} \SUPPORTEDANCH{\char"006C\char"0350} \SUPPORTEDANCH{\char"0234\char"0350} \SUPPORTEDANCH{\char"006E\char"0350} \SUPPORTEDANCH{\char"0235\char"0350} \SUPPORTEDANCH{\char"0070\char"0350} \SUPPORTEDANCH{\char"01A5\char"0350} \SUPPORTEDANCH{\char"0071\char"0350} \SUPPORTEDANCH{\char"02A0\char"0350} \SUPPORTEDANCH{\char"0074\char"0350} \SUPPORTEDANCH{\char"01AB\char"0350} \SUPPORTEDANCH{\char"01AD\char"0350} \SUPPORTEDANCH{\char"0236\char"0350} \SUPPORTEDANCH{\char"0075\char"0350} \MISSINGGLYPH{\char"AB52\char"0350} \SUPPORTEDANCH{\char"0079\char"0350} \SUPPORTEDANCH{\char"01B4\char"0350} \MISSINGGLYPH{\char"1EFF\char"0350}

\SUPPORTEDANCH{\char"0048\char"0351} \SUPPORTEDANCH{\char"2C67\char"0351} \SUPPORTEDANCH{\char"004B\char"0351} \SUPPORTEDANCH{\char"2C69\char"0351} \SUPPORTEDANCH{\char"004D\char"0351} \MISSINGGLYPH{\char"2C6E\char"0351} \SUPPORTEDANCH{\char"0050\char"0351} \SUPPORTEDANCH{\char"01A4\char"0351} \SUPPORTEDANCH{\char"0052\char"0351} \SUPPORTEDANCH{\char"2C64\char"0351} \SUPPORTEDANCH{\char"0054\char"0351} \SUPPORTEDANCH{\char"01AC\char"0351} \SUPPORTEDANCH{\char"01AE\char"0351} \SUPPORTEDANCH{\char"0059\char"0351} \SUPPORTEDANCH{\char"01B3\char"0351} \SUPPORTEDANCH{\char"0061\char"0351} \SUPPORTEDANCH{\char"1E9A\char"0351} \SUPPORTEDANCH{\char"0062\char"0351} \SUPPORTEDANCH{\char"0183\char"0351} \SUPPORTEDANCH{\char"0063\char"0351} \SUPPORTEDANCH{\char"0188\char"0351} \SUPPORTEDANCH{\char"0064\char"0351} \SUPPORTEDANCH{\char"018C\char"0351} \SUPPORTEDANCH{\char"0221\char"0351} \SUPPORTEDANCH{\char"0068\char"0351} \SUPPORTEDANCH{\char"2C68\char"0351} \SUPPORTEDANCH{\char"006B\char"0351} \SUPPORTEDANCH{\char"0199\char"0351} \SUPPORTEDANCH{\char"2C6A\char"0351} \SUPPORTEDANCH{\char"006C\char"0351} \SUPPORTEDANCH{\char"0234\char"0351} \SUPPORTEDANCH{\char"006E\char"0351} \SUPPORTEDANCH{\char"0235\char"0351} \SUPPORTEDANCH{\char"0070\char"0351} \SUPPORTEDANCH{\char"01A5\char"0351} \SUPPORTEDANCH{\char"0071\char"0351} \SUPPORTEDANCH{\char"02A0\char"0351} \SUPPORTEDANCH{\char"0074\char"0351} \SUPPORTEDANCH{\char"01AB\char"0351} \SUPPORTEDANCH{\char"01AD\char"0351} \SUPPORTEDANCH{\char"0236\char"0351} \SUPPORTEDANCH{\char"0075\char"0351} \MISSINGGLYPH{\char"AB52\char"0351} \SUPPORTEDANCH{\char"0079\char"0351} \SUPPORTEDANCH{\char"01B4\char"0351} \MISSINGGLYPH{\char"1EFF\char"0351}

\SUPPORTEDANCH{\char"0048\char"0352} \SUPPORTEDANCH{\char"2C67\char"0352} \SUPPORTEDANCH{\char"004B\char"0352} \SUPPORTEDANCH{\char"2C69\char"0352} \SUPPORTEDANCH{\char"004D\char"0352} \MISSINGGLYPH{\char"2C6E\char"0352} \SUPPORTEDANCH{\char"0050\char"0352} \SUPPORTEDANCH{\char"01A4\char"0352} \SUPPORTEDANCH{\char"0052\char"0352} \SUPPORTEDANCH{\char"2C64\char"0352} \SUPPORTEDANCH{\char"0054\char"0352} \SUPPORTEDANCH{\char"01AC\char"0352} \SUPPORTEDANCH{\char"01AE\char"0352} \SUPPORTEDANCH{\char"0059\char"0352} \SUPPORTEDANCH{\char"01B3\char"0352} \SUPPORTEDANCH{\char"0061\char"0352} \SUPPORTEDANCH{\char"1E9A\char"0352} \SUPPORTEDANCH{\char"0062\char"0352} \SUPPORTEDANCH{\char"0183\char"0352} \SUPPORTEDANCH{\char"0063\char"0352} \SUPPORTEDANCH{\char"0188\char"0352} \SUPPORTEDANCH{\char"0064\char"0352} \SUPPORTEDANCH{\char"018C\char"0352} \SUPPORTEDANCH{\char"0221\char"0352} \SUPPORTEDANCH{\char"0068\char"0352} \SUPPORTEDANCH{\char"2C68\char"0352} \SUPPORTEDANCH{\char"006B\char"0352} \SUPPORTEDANCH{\char"0199\char"0352} \SUPPORTEDANCH{\char"2C6A\char"0352} \SUPPORTEDANCH{\char"006C\char"0352} \SUPPORTEDANCH{\char"0234\char"0352} \SUPPORTEDANCH{\char"006E\char"0352} \SUPPORTEDANCH{\char"0235\char"0352} \SUPPORTEDANCH{\char"0070\char"0352} \SUPPORTEDANCH{\char"01A5\char"0352} \SUPPORTEDANCH{\char"0071\char"0352} \SUPPORTEDANCH{\char"02A0\char"0352} \SUPPORTEDANCH{\char"0074\char"0352} \SUPPORTEDANCH{\char"01AB\char"0352} \SUPPORTEDANCH{\char"01AD\char"0352} \SUPPORTEDANCH{\char"0236\char"0352} \SUPPORTEDANCH{\char"0075\char"0352} \MISSINGGLYPH{\char"AB52\char"0352} \SUPPORTEDANCH{\char"0079\char"0352} \SUPPORTEDANCH{\char"01B4\char"0352} \MISSINGGLYPH{\char"1EFF\char"0352}

\SUPPORTEDANCH{\char"0048\char"0357} \SUPPORTEDANCH{\char"2C67\char"0357} \SUPPORTEDANCH{\char"004B\char"0357} \SUPPORTEDANCH{\char"2C69\char"0357} \SUPPORTEDANCH{\char"004D\char"0357} \MISSINGGLYPH{\char"2C6E\char"0357} \SUPPORTEDANCH{\char"0050\char"0357} \SUPPORTEDANCH{\char"01A4\char"0357} \SUPPORTEDANCH{\char"0052\char"0357} \SUPPORTEDANCH{\char"2C64\char"0357} \SUPPORTEDANCH{\char"0054\char"0357} \SUPPORTEDANCH{\char"01AC\char"0357} \SUPPORTEDANCH{\char"01AE\char"0357} \SUPPORTEDANCH{\char"0059\char"0357} \SUPPORTEDANCH{\char"01B3\char"0357} \SUPPORTEDANCH{\char"0061\char"0357} \SUPPORTEDANCH{\char"1E9A\char"0357} \SUPPORTEDANCH{\char"0062\char"0357} \SUPPORTEDANCH{\char"0183\char"0357} \SUPPORTEDANCH{\char"0063\char"0357} \SUPPORTEDANCH{\char"0188\char"0357} \SUPPORTEDANCH{\char"0064\char"0357} \SUPPORTEDANCH{\char"018C\char"0357} \SUPPORTEDANCH{\char"0221\char"0357} \SUPPORTEDANCH{\char"0068\char"0357} \SUPPORTEDANCH{\char"2C68\char"0357} \SUPPORTEDANCH{\char"006B\char"0357} \SUPPORTEDANCH{\char"0199\char"0357} \SUPPORTEDANCH{\char"2C6A\char"0357} \SUPPORTEDANCH{\char"006C\char"0357} \SUPPORTEDANCH{\char"0234\char"0357} \SUPPORTEDANCH{\char"006E\char"0357} \SUPPORTEDANCH{\char"0235\char"0357} \SUPPORTEDANCH{\char"0070\char"0357} \SUPPORTEDANCH{\char"01A5\char"0357} \SUPPORTEDANCH{\char"0071\char"0357} \SUPPORTEDANCH{\char"02A0\char"0357} \SUPPORTEDANCH{\char"0074\char"0357} \SUPPORTEDANCH{\char"01AB\char"0357} \SUPPORTEDANCH{\char"01AD\char"0357} \SUPPORTEDANCH{\char"0236\char"0357} \SUPPORTEDANCH{\char"0075\char"0357} \MISSINGGLYPH{\char"AB52\char"0357} \SUPPORTEDANCH{\char"0079\char"0357} \SUPPORTEDANCH{\char"01B4\char"0357} \MISSINGGLYPH{\char"1EFF\char"0357}

\SUPPORTEDANCH{\char"0048\char"035B} \SUPPORTEDANCH{\char"2C67\char"035B} \SUPPORTEDANCH{\char"004B\char"035B} \SUPPORTEDANCH{\char"2C69\char"035B} \SUPPORTEDANCH{\char"004D\char"035B} \MISSINGGLYPH{\char"2C6E\char"035B} \SUPPORTEDANCH{\char"0050\char"035B} \SUPPORTEDANCH{\char"01A4\char"035B} \SUPPORTEDANCH{\char"0052\char"035B} \SUPPORTEDANCH{\char"2C64\char"035B} \SUPPORTEDANCH{\char"0054\char"035B} \SUPPORTEDANCH{\char"01AC\char"035B} \SUPPORTEDANCH{\char"01AE\char"035B} \SUPPORTEDANCH{\char"0059\char"035B} \SUPPORTEDANCH{\char"01B3\char"035B} \SUPPORTEDANCH{\char"0061\char"035B} \SUPPORTEDANCH{\char"1E9A\char"035B} \SUPPORTEDANCH{\char"0062\char"035B} \SUPPORTEDANCH{\char"0183\char"035B} \SUPPORTEDANCH{\char"0063\char"035B} \SUPPORTEDANCH{\char"0188\char"035B} \SUPPORTEDANCH{\char"0064\char"035B} \SUPPORTEDANCH{\char"018C\char"035B} \SUPPORTEDANCH{\char"0221\char"035B} \SUPPORTEDANCH{\char"0068\char"035B} \SUPPORTEDANCH{\char"2C68\char"035B} \SUPPORTEDANCH{\char"006B\char"035B} \SUPPORTEDANCH{\char"0199\char"035B} \SUPPORTEDANCH{\char"2C6A\char"035B} \SUPPORTEDANCH{\char"006C\char"035B} \SUPPORTEDANCH{\char"0234\char"035B} \SUPPORTEDANCH{\char"006E\char"035B} \SUPPORTEDANCH{\char"0235\char"035B} \SUPPORTEDANCH{\char"0070\char"035B} \SUPPORTEDANCH{\char"01A5\char"035B} \SUPPORTEDANCH{\char"0071\char"035B} \SUPPORTEDANCH{\char"02A0\char"035B} \SUPPORTEDANCH{\char"0074\char"035B} \SUPPORTEDANCH{\char"01AB\char"035B} \SUPPORTEDANCH{\char"01AD\char"035B} \SUPPORTEDANCH{\char"0236\char"035B} \SUPPORTEDANCH{\char"0075\char"035B} \MISSINGGLYPH{\char"AB52\char"035B} \SUPPORTEDANCH{\char"0079\char"035B} \SUPPORTEDANCH{\char"01B4\char"035B} \MISSINGGLYPH{\char"1EFF\char"035B}

}

\newpage
\subsection*{Semibold, Group 4, above, Above}

{\bfseries
% grid. columns are bases, rows are marks.
\SUPPORTEDFALL{\char"0048\char"0300} \SUPPORTEDFALL{\char"2C67\char"0300} \SUPPORTEDFALL{\char"004B\char"0300} \SUPPORTEDFALL{\char"2C69\char"0300} \SUPPORTEDFALL{\char"004D\char"0300} \MISSINGGLYPH{\char"2C6E\char"0300} \SUPPORTEDFALL{\char"0050\char"0300} \SUPPORTEDFALL{\char"01A4\char"0300} \SUPPORTEDANCH{\char"0052\char"0300} \SUPPORTEDFALL{\char"2C64\char"0300} \SUPPORTEDANCH{\char"0054\char"0300} \SUPPORTEDFALL{\char"01AC\char"0300} \SUPPORTEDFALL{\char"01AE\char"0300} \SUPPORTEDPREC{\char"0059\char"0300} \SUPPORTEDFALL{\char"01B3\char"0300} \SUPPORTEDPREC{\char"0061\char"0300} \SUPPORTEDFALL{\char"1E9A\char"0300} \SUPPORTEDFALL{\char"0062\char"0300} \SUPPORTEDFALL{\char"0183\char"0300} \SUPPORTEDFALL{\char"0063\char"0300} \SUPPORTEDFALL{\char"0188\char"0300} \SUPPORTEDFALL{\char"0064\char"0300} \SUPPORTEDFALL{\char"018C\char"0300} \SUPPORTEDFALL{\char"0221\char"0300} \SUPPORTEDFALL{\char"0068\char"0300} \SUPPORTEDFALL{\char"2C68\char"0300} \SUPPORTEDFALL{\char"006B\char"0300} \SUPPORTEDFALL{\char"0199\char"0300} \SUPPORTEDFALL{\char"2C6A\char"0300} \SUPPORTEDANCH{\char"006C\char"0300} \SUPPORTEDFALL{\char"0234\char"0300} \SUPPORTEDPREC{\char"006E\char"0300} \SUPPORTEDFALL{\char"0235\char"0300} \SUPPORTEDFALL{\char"0070\char"0300} \SUPPORTEDFALL{\char"01A5\char"0300} \SUPPORTEDFALL{\char"0071\char"0300} \SUPPORTEDFALL{\char"02A0\char"0300} \SUPPORTEDFALL{\char"0074\char"0300} \SUPPORTEDFALL{\char"01AB\char"0300} \SUPPORTEDFALL{\char"01AD\char"0300} \SUPPORTEDFALL{\char"0236\char"0300} \SUPPORTEDPREC{\char"0075\char"0300} \MISSINGGLYPH{\char"AB52\char"0300} \SUPPORTEDPREC{\char"0079\char"0300} \SUPPORTEDFALL{\char"01B4\char"0300} \MISSINGGLYPH{\char"1EFF\char"0300}

\SUPPORTEDFALL{\char"0048\char"0301} \SUPPORTEDFALL{\char"2C67\char"0301} \SUPPORTEDPREC{\char"004B\char"0301} \SUPPORTEDFALL{\char"2C69\char"0301} \SUPPORTEDPREC{\char"004D\char"0301} \MISSINGGLYPH{\char"2C6E\char"0301} \SUPPORTEDPREC{\char"0050\char"0301} \SUPPORTEDFALL{\char"01A4\char"0301} \SUPPORTEDPREC{\char"0052\char"0301} \SUPPORTEDFALL{\char"2C64\char"0301} \SUPPORTEDANCH{\char"0054\char"0301} \SUPPORTEDFALL{\char"01AC\char"0301} \SUPPORTEDFALL{\char"01AE\char"0301} \SUPPORTEDPREC{\char"0059\char"0301} \SUPPORTEDFALL{\char"01B3\char"0301} \SUPPORTEDPREC{\char"0061\char"0301} \SUPPORTEDFALL{\char"1E9A\char"0301} \SUPPORTEDFALL{\char"0062\char"0301} \SUPPORTEDFALL{\char"0183\char"0301} \SUPPORTEDPREC{\char"0063\char"0301} \SUPPORTEDFALL{\char"0188\char"0301} \SUPPORTEDFALL{\char"0064\char"0301} \SUPPORTEDFALL{\char"018C\char"0301} \SUPPORTEDFALL{\char"0221\char"0301} \SUPPORTEDFALL{\char"0068\char"0301} \SUPPORTEDFALL{\char"2C68\char"0301} \SUPPORTEDPREC{\char"006B\char"0301} \SUPPORTEDFALL{\char"0199\char"0301} \SUPPORTEDFALL{\char"2C6A\char"0301} \SUPPORTEDPREC{\char"006C\char"0301} \SUPPORTEDFALL{\char"0234\char"0301} \SUPPORTEDPREC{\char"006E\char"0301} \SUPPORTEDFALL{\char"0235\char"0301} \SUPPORTEDPREC{\char"0070\char"0301} \SUPPORTEDFALL{\char"01A5\char"0301} \SUPPORTEDFALL{\char"0071\char"0301} \SUPPORTEDFALL{\char"02A0\char"0301} \SUPPORTEDFALL{\char"0074\char"0301} \SUPPORTEDFALL{\char"01AB\char"0301} \SUPPORTEDFALL{\char"01AD\char"0301} \SUPPORTEDFALL{\char"0236\char"0301} \SUPPORTEDPREC{\char"0075\char"0301} \MISSINGGLYPH{\char"AB52\char"0301} \SUPPORTEDPREC{\char"0079\char"0301} \SUPPORTEDFALL{\char"01B4\char"0301} \MISSINGGLYPH{\char"1EFF\char"0301}

\SUPPORTEDPREC{\char"0048\char"0302} \SUPPORTEDFALL{\char"2C67\char"0302} \SUPPORTEDFALL{\char"004B\char"0302} \SUPPORTEDFALL{\char"2C69\char"0302} \SUPPORTEDFALL{\char"004D\char"0302} \MISSINGGLYPH{\char"2C6E\char"0302} \SUPPORTEDFALL{\char"0050\char"0302} \SUPPORTEDFALL{\char"01A4\char"0302} \SUPPORTEDANCH{\char"0052\char"0302} \SUPPORTEDFALL{\char"2C64\char"0302} \SUPPORTEDANCH{\char"0054\char"0302} \SUPPORTEDFALL{\char"01AC\char"0302} \SUPPORTEDFALL{\char"01AE\char"0302} \SUPPORTEDPREC{\char"0059\char"0302} \SUPPORTEDFALL{\char"01B3\char"0302} \SUPPORTEDPREC{\char"0061\char"0302} \SUPPORTEDFALL{\char"1E9A\char"0302} \SUPPORTEDFALL{\char"0062\char"0302} \SUPPORTEDFALL{\char"0183\char"0302} \SUPPORTEDPREC{\char"0063\char"0302} \SUPPORTEDFALL{\char"0188\char"0302} \SUPPORTEDFALL{\char"0064\char"0302} \SUPPORTEDFALL{\char"018C\char"0302} \SUPPORTEDFALL{\char"0221\char"0302} \SUPPORTEDPREC{\char"0068\char"0302} \SUPPORTEDFALL{\char"2C68\char"0302} \SUPPORTEDFALL{\char"006B\char"0302} \SUPPORTEDFALL{\char"0199\char"0302} \SUPPORTEDFALL{\char"2C6A\char"0302} \SUPPORTEDANCH{\char"006C\char"0302} \SUPPORTEDFALL{\char"0234\char"0302} \SUPPORTEDANCH{\char"006E\char"0302} \SUPPORTEDFALL{\char"0235\char"0302} \SUPPORTEDFALL{\char"0070\char"0302} \SUPPORTEDFALL{\char"01A5\char"0302} \SUPPORTEDFALL{\char"0071\char"0302} \SUPPORTEDFALL{\char"02A0\char"0302} \SUPPORTEDFALL{\char"0074\char"0302} \SUPPORTEDFALL{\char"01AB\char"0302} \SUPPORTEDFALL{\char"01AD\char"0302} \SUPPORTEDFALL{\char"0236\char"0302} \SUPPORTEDPREC{\char"0075\char"0302} \MISSINGGLYPH{\char"AB52\char"0302} \SUPPORTEDPREC{\char"0079\char"0302} \SUPPORTEDFALL{\char"01B4\char"0302} \MISSINGGLYPH{\char"1EFF\char"0302}

\SUPPORTEDFALL{\char"0048\char"0303} \SUPPORTEDFALL{\char"2C67\char"0303} \SUPPORTEDFALL{\char"004B\char"0303} \SUPPORTEDFALL{\char"2C69\char"0303} \SUPPORTEDFALL{\char"004D\char"0303} \MISSINGGLYPH{\char"2C6E\char"0303} \SUPPORTEDFALL{\char"0050\char"0303} \SUPPORTEDFALL{\char"01A4\char"0303} \SUPPORTEDANCH{\char"0052\char"0303} \SUPPORTEDFALL{\char"2C64\char"0303} \SUPPORTEDANCH{\char"0054\char"0303} \SUPPORTEDFALL{\char"01AC\char"0303} \SUPPORTEDFALL{\char"01AE\char"0303} \SUPPORTEDPREC{\char"0059\char"0303} \SUPPORTEDFALL{\char"01B3\char"0303} \SUPPORTEDPREC{\char"0061\char"0303} \SUPPORTEDFALL{\char"1E9A\char"0303} \SUPPORTEDFALL{\char"0062\char"0303} \SUPPORTEDFALL{\char"0183\char"0303} \SUPPORTEDFALL{\char"0063\char"0303} \SUPPORTEDFALL{\char"0188\char"0303} \SUPPORTEDFALL{\char"0064\char"0303} \SUPPORTEDFALL{\char"018C\char"0303} \SUPPORTEDFALL{\char"0221\char"0303} \SUPPORTEDFALL{\char"0068\char"0303} \SUPPORTEDFALL{\char"2C68\char"0303} \SUPPORTEDFALL{\char"006B\char"0303} \SUPPORTEDFALL{\char"0199\char"0303} \SUPPORTEDFALL{\char"2C6A\char"0303} \SUPPORTEDANCH{\char"006C\char"0303} \SUPPORTEDFALL{\char"0234\char"0303} \SUPPORTEDPREC{\char"006E\char"0303} \SUPPORTEDFALL{\char"0235\char"0303} \SUPPORTEDFALL{\char"0070\char"0303} \SUPPORTEDFALL{\char"01A5\char"0303} \SUPPORTEDFALL{\char"0071\char"0303} \SUPPORTEDFALL{\char"02A0\char"0303} \SUPPORTEDFALL{\char"0074\char"0303} \SUPPORTEDFALL{\char"01AB\char"0303} \SUPPORTEDFALL{\char"01AD\char"0303} \SUPPORTEDFALL{\char"0236\char"0303} \SUPPORTEDPREC{\char"0075\char"0303} \MISSINGGLYPH{\char"AB52\char"0303} \SUPPORTEDPREC{\char"0079\char"0303} \SUPPORTEDFALL{\char"01B4\char"0303} \MISSINGGLYPH{\char"1EFF\char"0303}

\SUPPORTEDFALL{\char"0048\char"0304} \SUPPORTEDFALL{\char"2C67\char"0304} \SUPPORTEDFALL{\char"004B\char"0304} \SUPPORTEDFALL{\char"2C69\char"0304} \SUPPORTEDFALL{\char"004D\char"0304} \MISSINGGLYPH{\char"2C6E\char"0304} \SUPPORTEDFALL{\char"0050\char"0304} \SUPPORTEDFALL{\char"01A4\char"0304} \SUPPORTEDANCH{\char"0052\char"0304} \SUPPORTEDFALL{\char"2C64\char"0304} \SUPPORTEDANCH{\char"0054\char"0304} \SUPPORTEDFALL{\char"01AC\char"0304} \SUPPORTEDFALL{\char"01AE\char"0304} \SUPPORTEDPREC{\char"0059\char"0304} \SUPPORTEDFALL{\char"01B3\char"0304} \SUPPORTEDPREC{\char"0061\char"0304} \SUPPORTEDFALL{\char"1E9A\char"0304} \SUPPORTEDFALL{\char"0062\char"0304} \SUPPORTEDFALL{\char"0183\char"0304} \SUPPORTEDFALL{\char"0063\char"0304} \SUPPORTEDFALL{\char"0188\char"0304} \SUPPORTEDFALL{\char"0064\char"0304} \SUPPORTEDFALL{\char"018C\char"0304} \SUPPORTEDFALL{\char"0221\char"0304} \SUPPORTEDFALL{\char"0068\char"0304} \SUPPORTEDFALL{\char"2C68\char"0304} \SUPPORTEDFALL{\char"006B\char"0304} \SUPPORTEDFALL{\char"0199\char"0304} \SUPPORTEDFALL{\char"2C6A\char"0304} \SUPPORTEDANCH{\char"006C\char"0304} \SUPPORTEDFALL{\char"0234\char"0304} \SUPPORTEDANCH{\char"006E\char"0304} \SUPPORTEDFALL{\char"0235\char"0304} \SUPPORTEDFALL{\char"0070\char"0304} \SUPPORTEDFALL{\char"01A5\char"0304} \SUPPORTEDFALL{\char"0071\char"0304} \SUPPORTEDFALL{\char"02A0\char"0304} \SUPPORTEDFALL{\char"0074\char"0304} \SUPPORTEDFALL{\char"01AB\char"0304} \SUPPORTEDFALL{\char"01AD\char"0304} \SUPPORTEDFALL{\char"0236\char"0304} \SUPPORTEDPREC{\char"0075\char"0304} \MISSINGGLYPH{\char"AB52\char"0304} \SUPPORTEDPREC{\char"0079\char"0304} \SUPPORTEDFALL{\char"01B4\char"0304} \MISSINGGLYPH{\char"1EFF\char"0304}

\SUPPORTEDFALL{\char"0048\char"0306} \SUPPORTEDFALL{\char"2C67\char"0306} \SUPPORTEDFALL{\char"004B\char"0306} \SUPPORTEDFALL{\char"2C69\char"0306} \SUPPORTEDFALL{\char"004D\char"0306} \MISSINGGLYPH{\char"2C6E\char"0306} \SUPPORTEDFALL{\char"0050\char"0306} \SUPPORTEDFALL{\char"01A4\char"0306} \SUPPORTEDANCH{\char"0052\char"0306} \SUPPORTEDFALL{\char"2C64\char"0306} \SUPPORTEDANCH{\char"0054\char"0306} \SUPPORTEDFALL{\char"01AC\char"0306} \SUPPORTEDFALL{\char"01AE\char"0306} \SUPPORTEDANCH{\char"0059\char"0306} \SUPPORTEDFALL{\char"01B3\char"0306} \SUPPORTEDPREC{\char"0061\char"0306} \SUPPORTEDFALL{\char"1E9A\char"0306} \SUPPORTEDFALL{\char"0062\char"0306} \SUPPORTEDFALL{\char"0183\char"0306} \SUPPORTEDFALL{\char"0063\char"0306} \SUPPORTEDFALL{\char"0188\char"0306} \SUPPORTEDFALL{\char"0064\char"0306} \SUPPORTEDFALL{\char"018C\char"0306} \SUPPORTEDFALL{\char"0221\char"0306} \SUPPORTEDFALL{\char"0068\char"0306} \SUPPORTEDFALL{\char"2C68\char"0306} \SUPPORTEDFALL{\char"006B\char"0306} \SUPPORTEDFALL{\char"0199\char"0306} \SUPPORTEDFALL{\char"2C6A\char"0306} \SUPPORTEDANCH{\char"006C\char"0306} \SUPPORTEDFALL{\char"0234\char"0306} \SUPPORTEDANCH{\char"006E\char"0306} \SUPPORTEDFALL{\char"0235\char"0306} \SUPPORTEDFALL{\char"0070\char"0306} \SUPPORTEDFALL{\char"01A5\char"0306} \SUPPORTEDFALL{\char"0071\char"0306} \SUPPORTEDFALL{\char"02A0\char"0306} \SUPPORTEDFALL{\char"0074\char"0306} \SUPPORTEDFALL{\char"01AB\char"0306} \SUPPORTEDFALL{\char"01AD\char"0306} \SUPPORTEDFALL{\char"0236\char"0306} \SUPPORTEDPREC{\char"0075\char"0306} \MISSINGGLYPH{\char"AB52\char"0306} \SUPPORTEDFALL{\char"0079\char"0306} \SUPPORTEDFALL{\char"01B4\char"0306} \MISSINGGLYPH{\char"1EFF\char"0306}

\SUPPORTEDPREC{\char"0048\char"0307} \SUPPORTEDFALL{\char"2C67\char"0307} \SUPPORTEDFALL{\char"004B\char"0307} \SUPPORTEDFALL{\char"2C69\char"0307} \SUPPORTEDPREC{\char"004D\char"0307} \MISSINGGLYPH{\char"2C6E\char"0307} \SUPPORTEDPREC{\char"0050\char"0307} \SUPPORTEDFALL{\char"01A4\char"0307} \SUPPORTEDPREC{\char"0052\char"0307} \SUPPORTEDFALL{\char"2C64\char"0307} \SUPPORTEDPREC{\char"0054\char"0307} \SUPPORTEDFALL{\char"01AC\char"0307} \SUPPORTEDFALL{\char"01AE\char"0307} \SUPPORTEDPREC{\char"0059\char"0307} \SUPPORTEDFALL{\char"01B3\char"0307} \SUPPORTEDPREC{\char"0061\char"0307} \SUPPORTEDFALL{\char"1E9A\char"0307} \SUPPORTEDPREC{\char"0062\char"0307} \SUPPORTEDFALL{\char"0183\char"0307} \SUPPORTEDPREC{\char"0063\char"0307} \SUPPORTEDFALL{\char"0188\char"0307} \SUPPORTEDPREC{\char"0064\char"0307} \SUPPORTEDFALL{\char"018C\char"0307} \SUPPORTEDFALL{\char"0221\char"0307} \SUPPORTEDPREC{\char"0068\char"0307} \SUPPORTEDFALL{\char"2C68\char"0307} \SUPPORTEDFALL{\char"006B\char"0307} \SUPPORTEDFALL{\char"0199\char"0307} \SUPPORTEDFALL{\char"2C6A\char"0307} \SUPPORTEDANCH{\char"006C\char"0307} \SUPPORTEDFALL{\char"0234\char"0307} \SUPPORTEDPREC{\char"006E\char"0307} \SUPPORTEDFALL{\char"0235\char"0307} \SUPPORTEDPREC{\char"0070\char"0307} \SUPPORTEDFALL{\char"01A5\char"0307} \SUPPORTEDFALL{\char"0071\char"0307} \SUPPORTEDFALL{\char"02A0\char"0307} \SUPPORTEDPREC{\char"0074\char"0307} \SUPPORTEDFALL{\char"01AB\char"0307} \SUPPORTEDFALL{\char"01AD\char"0307} \SUPPORTEDFALL{\char"0236\char"0307} \SUPPORTEDANCH{\char"0075\char"0307} \MISSINGGLYPH{\char"AB52\char"0307} \SUPPORTEDPREC{\char"0079\char"0307} \SUPPORTEDFALL{\char"01B4\char"0307} \MISSINGGLYPH{\char"1EFF\char"0307}

\SUPPORTEDPREC{\char"0048\char"0308} \SUPPORTEDFALL{\char"2C67\char"0308} \SUPPORTEDFALL{\char"004B\char"0308} \SUPPORTEDFALL{\char"2C69\char"0308} \SUPPORTEDFALL{\char"004D\char"0308} \MISSINGGLYPH{\char"2C6E\char"0308} \SUPPORTEDFALL{\char"0050\char"0308} \SUPPORTEDFALL{\char"01A4\char"0308} \SUPPORTEDANCH{\char"0052\char"0308} \SUPPORTEDFALL{\char"2C64\char"0308} \SUPPORTEDANCH{\char"0054\char"0308} \SUPPORTEDFALL{\char"01AC\char"0308} \SUPPORTEDFALL{\char"01AE\char"0308} \SUPPORTEDPREC{\char"0059\char"0308} \SUPPORTEDFALL{\char"01B3\char"0308} \SUPPORTEDPREC{\char"0061\char"0308} \SUPPORTEDFALL{\char"1E9A\char"0308} \SUPPORTEDFALL{\char"0062\char"0308} \SUPPORTEDFALL{\char"0183\char"0308} \SUPPORTEDFALL{\char"0063\char"0308} \SUPPORTEDFALL{\char"0188\char"0308} \SUPPORTEDFALL{\char"0064\char"0308} \SUPPORTEDFALL{\char"018C\char"0308} \SUPPORTEDFALL{\char"0221\char"0308} \SUPPORTEDPREC{\char"0068\char"0308} \SUPPORTEDFALL{\char"2C68\char"0308} \SUPPORTEDFALL{\char"006B\char"0308} \SUPPORTEDFALL{\char"0199\char"0308} \SUPPORTEDFALL{\char"2C6A\char"0308} \SUPPORTEDANCH{\char"006C\char"0308} \SUPPORTEDFALL{\char"0234\char"0308} \SUPPORTEDANCH{\char"006E\char"0308} \SUPPORTEDFALL{\char"0235\char"0308} \SUPPORTEDFALL{\char"0070\char"0308} \SUPPORTEDFALL{\char"01A5\char"0308} \SUPPORTEDFALL{\char"0071\char"0308} \SUPPORTEDFALL{\char"02A0\char"0308} \SUPPORTEDPREC{\char"0074\char"0308} \SUPPORTEDFALL{\char"01AB\char"0308} \SUPPORTEDFALL{\char"01AD\char"0308} \SUPPORTEDFALL{\char"0236\char"0308} \SUPPORTEDPREC{\char"0075\char"0308} \MISSINGGLYPH{\char"AB52\char"0308} \SUPPORTEDPREC{\char"0079\char"0308} \SUPPORTEDFALL{\char"01B4\char"0308} \MISSINGGLYPH{\char"1EFF\char"0308}

\SUPPORTEDFALL{\char"0048\char"0309} \SUPPORTEDFALL{\char"2C67\char"0309} \SUPPORTEDFALL{\char"004B\char"0309} \SUPPORTEDFALL{\char"2C69\char"0309} \SUPPORTEDFALL{\char"004D\char"0309} \MISSINGGLYPH{\char"2C6E\char"0309} \SUPPORTEDFALL{\char"0050\char"0309} \SUPPORTEDFALL{\char"01A4\char"0309} \SUPPORTEDANCH{\char"0052\char"0309} \SUPPORTEDFALL{\char"2C64\char"0309} \SUPPORTEDANCH{\char"0054\char"0309} \SUPPORTEDFALL{\char"01AC\char"0309} \SUPPORTEDFALL{\char"01AE\char"0309} \SUPPORTEDPREC{\char"0059\char"0309} \SUPPORTEDFALL{\char"01B3\char"0309} \SUPPORTEDPREC{\char"0061\char"0309} \SUPPORTEDFALL{\char"1E9A\char"0309} \SUPPORTEDFALL{\char"0062\char"0309} \SUPPORTEDFALL{\char"0183\char"0309} \SUPPORTEDFALL{\char"0063\char"0309} \SUPPORTEDFALL{\char"0188\char"0309} \SUPPORTEDFALL{\char"0064\char"0309} \SUPPORTEDFALL{\char"018C\char"0309} \SUPPORTEDFALL{\char"0221\char"0309} \SUPPORTEDFALL{\char"0068\char"0309} \SUPPORTEDFALL{\char"2C68\char"0309} \SUPPORTEDFALL{\char"006B\char"0309} \SUPPORTEDFALL{\char"0199\char"0309} \SUPPORTEDFALL{\char"2C6A\char"0309} \SUPPORTEDANCH{\char"006C\char"0309} \SUPPORTEDFALL{\char"0234\char"0309} \SUPPORTEDANCH{\char"006E\char"0309} \SUPPORTEDFALL{\char"0235\char"0309} \SUPPORTEDFALL{\char"0070\char"0309} \SUPPORTEDFALL{\char"01A5\char"0309} \SUPPORTEDFALL{\char"0071\char"0309} \SUPPORTEDFALL{\char"02A0\char"0309} \SUPPORTEDFALL{\char"0074\char"0309} \SUPPORTEDFALL{\char"01AB\char"0309} \SUPPORTEDFALL{\char"01AD\char"0309} \SUPPORTEDFALL{\char"0236\char"0309} \SUPPORTEDPREC{\char"0075\char"0309} \MISSINGGLYPH{\char"AB52\char"0309} \SUPPORTEDPREC{\char"0079\char"0309} \SUPPORTEDFALL{\char"01B4\char"0309} \MISSINGGLYPH{\char"1EFF\char"0309}

\SUPPORTEDFALL{\char"0048\char"030A} \SUPPORTEDFALL{\char"2C67\char"030A} \SUPPORTEDFALL{\char"004B\char"030A} \SUPPORTEDFALL{\char"2C69\char"030A} \SUPPORTEDFALL{\char"004D\char"030A} \MISSINGGLYPH{\char"2C6E\char"030A} \SUPPORTEDFALL{\char"0050\char"030A} \SUPPORTEDFALL{\char"01A4\char"030A} \SUPPORTEDANCH{\char"0052\char"030A} \SUPPORTEDFALL{\char"2C64\char"030A} \SUPPORTEDANCH{\char"0054\char"030A} \SUPPORTEDFALL{\char"01AC\char"030A} \SUPPORTEDFALL{\char"01AE\char"030A} \SUPPORTEDANCH{\char"0059\char"030A} \SUPPORTEDFALL{\char"01B3\char"030A} \SUPPORTEDPREC{\char"0061\char"030A} \SUPPORTEDFALL{\char"1E9A\char"030A} \SUPPORTEDFALL{\char"0062\char"030A} \SUPPORTEDFALL{\char"0183\char"030A} \SUPPORTEDFALL{\char"0063\char"030A} \SUPPORTEDFALL{\char"0188\char"030A} \SUPPORTEDFALL{\char"0064\char"030A} \SUPPORTEDFALL{\char"018C\char"030A} \SUPPORTEDFALL{\char"0221\char"030A} \SUPPORTEDFALL{\char"0068\char"030A} \SUPPORTEDFALL{\char"2C68\char"030A} \SUPPORTEDFALL{\char"006B\char"030A} \SUPPORTEDFALL{\char"0199\char"030A} \SUPPORTEDFALL{\char"2C6A\char"030A} \SUPPORTEDANCH{\char"006C\char"030A} \SUPPORTEDFALL{\char"0234\char"030A} \SUPPORTEDANCH{\char"006E\char"030A} \SUPPORTEDFALL{\char"0235\char"030A} \SUPPORTEDFALL{\char"0070\char"030A} \SUPPORTEDFALL{\char"01A5\char"030A} \SUPPORTEDFALL{\char"0071\char"030A} \SUPPORTEDFALL{\char"02A0\char"030A} \SUPPORTEDFALL{\char"0074\char"030A} \SUPPORTEDFALL{\char"01AB\char"030A} \SUPPORTEDFALL{\char"01AD\char"030A} \SUPPORTEDFALL{\char"0236\char"030A} \SUPPORTEDPREC{\char"0075\char"030A} \MISSINGGLYPH{\char"AB52\char"030A} \SUPPORTEDPREC{\char"0079\char"030A} \SUPPORTEDFALL{\char"01B4\char"030A} \MISSINGGLYPH{\char"1EFF\char"030A}

\SUPPORTEDFALL{\char"0048\char"030B} \SUPPORTEDFALL{\char"2C67\char"030B} \SUPPORTEDFALL{\char"004B\char"030B} \SUPPORTEDFALL{\char"2C69\char"030B} \SUPPORTEDFALL{\char"004D\char"030B} \MISSINGGLYPH{\char"2C6E\char"030B} \SUPPORTEDFALL{\char"0050\char"030B} \SUPPORTEDFALL{\char"01A4\char"030B} \SUPPORTEDANCH{\char"0052\char"030B} \SUPPORTEDFALL{\char"2C64\char"030B} \SUPPORTEDANCH{\char"0054\char"030B} \SUPPORTEDFALL{\char"01AC\char"030B} \SUPPORTEDFALL{\char"01AE\char"030B} \SUPPORTEDANCH{\char"0059\char"030B} \SUPPORTEDFALL{\char"01B3\char"030B} \SUPPORTEDFALL{\char"0061\char"030B} \SUPPORTEDFALL{\char"1E9A\char"030B} \SUPPORTEDFALL{\char"0062\char"030B} \SUPPORTEDFALL{\char"0183\char"030B} \SUPPORTEDFALL{\char"0063\char"030B} \SUPPORTEDFALL{\char"0188\char"030B} \SUPPORTEDFALL{\char"0064\char"030B} \SUPPORTEDFALL{\char"018C\char"030B} \SUPPORTEDFALL{\char"0221\char"030B} \SUPPORTEDFALL{\char"0068\char"030B} \SUPPORTEDFALL{\char"2C68\char"030B} \SUPPORTEDFALL{\char"006B\char"030B} \SUPPORTEDFALL{\char"0199\char"030B} \SUPPORTEDFALL{\char"2C6A\char"030B} \SUPPORTEDANCH{\char"006C\char"030B} \SUPPORTEDFALL{\char"0234\char"030B} \SUPPORTEDANCH{\char"006E\char"030B} \SUPPORTEDFALL{\char"0235\char"030B} \SUPPORTEDFALL{\char"0070\char"030B} \SUPPORTEDFALL{\char"01A5\char"030B} \SUPPORTEDFALL{\char"0071\char"030B} \SUPPORTEDFALL{\char"02A0\char"030B} \SUPPORTEDFALL{\char"0074\char"030B} \SUPPORTEDFALL{\char"01AB\char"030B} \SUPPORTEDFALL{\char"01AD\char"030B} \SUPPORTEDFALL{\char"0236\char"030B} \SUPPORTEDPREC{\char"0075\char"030B} \MISSINGGLYPH{\char"AB52\char"030B} \SUPPORTEDFALL{\char"0079\char"030B} \SUPPORTEDFALL{\char"01B4\char"030B} \MISSINGGLYPH{\char"1EFF\char"030B}

\SUPPORTEDPREC{\char"0048\char"030C} \SUPPORTEDFALL{\char"2C67\char"030C} \SUPPORTEDPREC{\char"004B\char"030C} \SUPPORTEDFALL{\char"2C69\char"030C} \SUPPORTEDFALL{\char"004D\char"030C} \MISSINGGLYPH{\char"2C6E\char"030C} \SUPPORTEDFALL{\char"0050\char"030C} \SUPPORTEDFALL{\char"01A4\char"030C} \SUPPORTEDPREC{\char"0052\char"030C} \SUPPORTEDFALL{\char"2C64\char"030C} \SUPPORTEDPREC{\char"0054\char"030C} \SUPPORTEDFALL{\char"01AC\char"030C} \SUPPORTEDFALL{\char"01AE\char"030C} \SUPPORTEDANCH{\char"0059\char"030C} \SUPPORTEDFALL{\char"01B3\char"030C} \SUPPORTEDPREC{\char"0061\char"030C} \SUPPORTEDFALL{\char"1E9A\char"030C} \SUPPORTEDFALL{\char"0062\char"030C} \SUPPORTEDFALL{\char"0183\char"030C} \SUPPORTEDPREC{\char"0063\char"030C} \SUPPORTEDFALL{\char"0188\char"030C} \SUPPORTEDPREC{\char"0064\char"030C} \SUPPORTEDFALL{\char"018C\char"030C} \SUPPORTEDFALL{\char"0221\char"030C} \SUPPORTEDPREC{\char"0068\char"030C} \SUPPORTEDFALL{\char"2C68\char"030C} \SUPPORTEDPREC{\char"006B\char"030C} \SUPPORTEDFALL{\char"0199\char"030C} \SUPPORTEDFALL{\char"2C6A\char"030C} \SUPPORTEDPREC{\char"006C\char"030C} \SUPPORTEDFALL{\char"0234\char"030C} \SUPPORTEDPREC{\char"006E\char"030C} \SUPPORTEDFALL{\char"0235\char"030C} \SUPPORTEDFALL{\char"0070\char"030C} \SUPPORTEDFALL{\char"01A5\char"030C} \SUPPORTEDFALL{\char"0071\char"030C} \SUPPORTEDFALL{\char"02A0\char"030C} \SUPPORTEDPREC{\char"0074\char"030C} \SUPPORTEDFALL{\char"01AB\char"030C} \SUPPORTEDFALL{\char"01AD\char"030C} \SUPPORTEDFALL{\char"0236\char"030C} \SUPPORTEDPREC{\char"0075\char"030C} \MISSINGGLYPH{\char"AB52\char"030C} \SUPPORTEDFALL{\char"0079\char"030C} \SUPPORTEDFALL{\char"01B4\char"030C} \MISSINGGLYPH{\char"1EFF\char"030C}

\SUPPORTEDFALL{\char"0048\char"030D} \SUPPORTEDFALL{\char"2C67\char"030D} \SUPPORTEDFALL{\char"004B\char"030D} \SUPPORTEDFALL{\char"2C69\char"030D} \SUPPORTEDFALL{\char"004D\char"030D} \MISSINGGLYPH{\char"2C6E\char"030D} \SUPPORTEDFALL{\char"0050\char"030D} \SUPPORTEDFALL{\char"01A4\char"030D} \SUPPORTEDANCH{\char"0052\char"030D} \SUPPORTEDFALL{\char"2C64\char"030D} \SUPPORTEDANCH{\char"0054\char"030D} \SUPPORTEDFALL{\char"01AC\char"030D} \SUPPORTEDFALL{\char"01AE\char"030D} \SUPPORTEDANCH{\char"0059\char"030D} \SUPPORTEDFALL{\char"01B3\char"030D} \SUPPORTEDFALL{\char"0061\char"030D} \SUPPORTEDFALL{\char"1E9A\char"030D} \SUPPORTEDFALL{\char"0062\char"030D} \SUPPORTEDFALL{\char"0183\char"030D} \SUPPORTEDFALL{\char"0063\char"030D} \SUPPORTEDFALL{\char"0188\char"030D} \SUPPORTEDFALL{\char"0064\char"030D} \SUPPORTEDFALL{\char"018C\char"030D} \SUPPORTEDFALL{\char"0221\char"030D} \SUPPORTEDFALL{\char"0068\char"030D} \SUPPORTEDFALL{\char"2C68\char"030D} \SUPPORTEDFALL{\char"006B\char"030D} \SUPPORTEDFALL{\char"0199\char"030D} \SUPPORTEDFALL{\char"2C6A\char"030D} \SUPPORTEDANCH{\char"006C\char"030D} \SUPPORTEDFALL{\char"0234\char"030D} \SUPPORTEDANCH{\char"006E\char"030D} \SUPPORTEDFALL{\char"0235\char"030D} \SUPPORTEDFALL{\char"0070\char"030D} \SUPPORTEDFALL{\char"01A5\char"030D} \SUPPORTEDFALL{\char"0071\char"030D} \SUPPORTEDFALL{\char"02A0\char"030D} \SUPPORTEDFALL{\char"0074\char"030D} \SUPPORTEDFALL{\char"01AB\char"030D} \SUPPORTEDFALL{\char"01AD\char"030D} \SUPPORTEDFALL{\char"0236\char"030D} \SUPPORTEDANCH{\char"0075\char"030D} \MISSINGGLYPH{\char"AB52\char"030D} \SUPPORTEDFALL{\char"0079\char"030D} \SUPPORTEDFALL{\char"01B4\char"030D} \MISSINGGLYPH{\char"1EFF\char"030D}

\SUPPORTEDFALL{\char"0048\char"030F} \SUPPORTEDFALL{\char"2C67\char"030F} \SUPPORTEDFALL{\char"004B\char"030F} \SUPPORTEDFALL{\char"2C69\char"030F} \SUPPORTEDFALL{\char"004D\char"030F} \MISSINGGLYPH{\char"2C6E\char"030F} \SUPPORTEDFALL{\char"0050\char"030F} \SUPPORTEDFALL{\char"01A4\char"030F} \SUPPORTEDPREC{\char"0052\char"030F} \SUPPORTEDFALL{\char"2C64\char"030F} \SUPPORTEDANCH{\char"0054\char"030F} \SUPPORTEDFALL{\char"01AC\char"030F} \SUPPORTEDFALL{\char"01AE\char"030F} \SUPPORTEDANCH{\char"0059\char"030F} \SUPPORTEDFALL{\char"01B3\char"030F} \SUPPORTEDPREC{\char"0061\char"030F} \SUPPORTEDFALL{\char"1E9A\char"030F} \SUPPORTEDFALL{\char"0062\char"030F} \SUPPORTEDFALL{\char"0183\char"030F} \SUPPORTEDFALL{\char"0063\char"030F} \SUPPORTEDFALL{\char"0188\char"030F} \SUPPORTEDFALL{\char"0064\char"030F} \SUPPORTEDFALL{\char"018C\char"030F} \SUPPORTEDFALL{\char"0221\char"030F} \SUPPORTEDFALL{\char"0068\char"030F} \SUPPORTEDFALL{\char"2C68\char"030F} \SUPPORTEDFALL{\char"006B\char"030F} \SUPPORTEDFALL{\char"0199\char"030F} \SUPPORTEDFALL{\char"2C6A\char"030F} \SUPPORTEDANCH{\char"006C\char"030F} \SUPPORTEDFALL{\char"0234\char"030F} \SUPPORTEDANCH{\char"006E\char"030F} \SUPPORTEDFALL{\char"0235\char"030F} \SUPPORTEDFALL{\char"0070\char"030F} \SUPPORTEDFALL{\char"01A5\char"030F} \SUPPORTEDFALL{\char"0071\char"030F} \SUPPORTEDFALL{\char"02A0\char"030F} \SUPPORTEDFALL{\char"0074\char"030F} \SUPPORTEDFALL{\char"01AB\char"030F} \SUPPORTEDFALL{\char"01AD\char"030F} \SUPPORTEDFALL{\char"0236\char"030F} \SUPPORTEDPREC{\char"0075\char"030F} \MISSINGGLYPH{\char"AB52\char"030F} \SUPPORTEDFALL{\char"0079\char"030F} \SUPPORTEDFALL{\char"01B4\char"030F} \MISSINGGLYPH{\char"1EFF\char"030F}

\SUPPORTEDFALL{\char"0048\char"0310} \SUPPORTEDFALL{\char"2C67\char"0310} \SUPPORTEDFALL{\char"004B\char"0310} \SUPPORTEDFALL{\char"2C69\char"0310} \SUPPORTEDFALL{\char"004D\char"0310} \MISSINGGLYPH{\char"2C6E\char"0310} \SUPPORTEDFALL{\char"0050\char"0310} \SUPPORTEDFALL{\char"01A4\char"0310} \SUPPORTEDANCH{\char"0052\char"0310} \SUPPORTEDFALL{\char"2C64\char"0310} \SUPPORTEDANCH{\char"0054\char"0310} \SUPPORTEDFALL{\char"01AC\char"0310} \SUPPORTEDFALL{\char"01AE\char"0310} \SUPPORTEDANCH{\char"0059\char"0310} \SUPPORTEDFALL{\char"01B3\char"0310} \SUPPORTEDFALL{\char"0061\char"0310} \SUPPORTEDFALL{\char"1E9A\char"0310} \SUPPORTEDFALL{\char"0062\char"0310} \SUPPORTEDFALL{\char"0183\char"0310} \SUPPORTEDFALL{\char"0063\char"0310} \SUPPORTEDFALL{\char"0188\char"0310} \SUPPORTEDFALL{\char"0064\char"0310} \SUPPORTEDFALL{\char"018C\char"0310} \SUPPORTEDFALL{\char"0221\char"0310} \SUPPORTEDFALL{\char"0068\char"0310} \SUPPORTEDFALL{\char"2C68\char"0310} \SUPPORTEDFALL{\char"006B\char"0310} \SUPPORTEDFALL{\char"0199\char"0310} \SUPPORTEDFALL{\char"2C6A\char"0310} \SUPPORTEDANCH{\char"006C\char"0310} \SUPPORTEDFALL{\char"0234\char"0310} \SUPPORTEDANCH{\char"006E\char"0310} \SUPPORTEDFALL{\char"0235\char"0310} \SUPPORTEDFALL{\char"0070\char"0310} \SUPPORTEDFALL{\char"01A5\char"0310} \SUPPORTEDFALL{\char"0071\char"0310} \SUPPORTEDFALL{\char"02A0\char"0310} \SUPPORTEDFALL{\char"0074\char"0310} \SUPPORTEDFALL{\char"01AB\char"0310} \SUPPORTEDFALL{\char"01AD\char"0310} \SUPPORTEDFALL{\char"0236\char"0310} \SUPPORTEDANCH{\char"0075\char"0310} \MISSINGGLYPH{\char"AB52\char"0310} \SUPPORTEDFALL{\char"0079\char"0310} \SUPPORTEDFALL{\char"01B4\char"0310} \MISSINGGLYPH{\char"1EFF\char"0310}

\SUPPORTEDFALL{\char"0048\char"0311} \SUPPORTEDFALL{\char"2C67\char"0311} \SUPPORTEDFALL{\char"004B\char"0311} \SUPPORTEDFALL{\char"2C69\char"0311} \SUPPORTEDFALL{\char"004D\char"0311} \MISSINGGLYPH{\char"2C6E\char"0311} \SUPPORTEDFALL{\char"0050\char"0311} \SUPPORTEDFALL{\char"01A4\char"0311} \SUPPORTEDPREC{\char"0052\char"0311} \SUPPORTEDFALL{\char"2C64\char"0311} \SUPPORTEDANCH{\char"0054\char"0311} \SUPPORTEDFALL{\char"01AC\char"0311} \SUPPORTEDFALL{\char"01AE\char"0311} \SUPPORTEDANCH{\char"0059\char"0311} \SUPPORTEDFALL{\char"01B3\char"0311} \SUPPORTEDPREC{\char"0061\char"0311} \SUPPORTEDFALL{\char"1E9A\char"0311} \SUPPORTEDFALL{\char"0062\char"0311} \SUPPORTEDFALL{\char"0183\char"0311} \SUPPORTEDFALL{\char"0063\char"0311} \SUPPORTEDFALL{\char"0188\char"0311} \SUPPORTEDFALL{\char"0064\char"0311} \SUPPORTEDFALL{\char"018C\char"0311} \SUPPORTEDFALL{\char"0221\char"0311} \SUPPORTEDFALL{\char"0068\char"0311} \SUPPORTEDFALL{\char"2C68\char"0311} \SUPPORTEDFALL{\char"006B\char"0311} \SUPPORTEDFALL{\char"0199\char"0311} \SUPPORTEDFALL{\char"2C6A\char"0311} \SUPPORTEDANCH{\char"006C\char"0311} \SUPPORTEDFALL{\char"0234\char"0311} \SUPPORTEDANCH{\char"006E\char"0311} \SUPPORTEDFALL{\char"0235\char"0311} \SUPPORTEDFALL{\char"0070\char"0311} \SUPPORTEDFALL{\char"01A5\char"0311} \SUPPORTEDFALL{\char"0071\char"0311} \SUPPORTEDFALL{\char"02A0\char"0311} \SUPPORTEDFALL{\char"0074\char"0311} \SUPPORTEDFALL{\char"01AB\char"0311} \SUPPORTEDFALL{\char"01AD\char"0311} \SUPPORTEDFALL{\char"0236\char"0311} \SUPPORTEDPREC{\char"0075\char"0311} \MISSINGGLYPH{\char"AB52\char"0311} \SUPPORTEDFALL{\char"0079\char"0311} \SUPPORTEDFALL{\char"01B4\char"0311} \MISSINGGLYPH{\char"1EFF\char"0311}

\SUPPORTEDFALL{\char"0048\char"0312} \SUPPORTEDFALL{\char"2C67\char"0312} \SUPPORTEDFALL{\char"004B\char"0312} \SUPPORTEDFALL{\char"2C69\char"0312} \SUPPORTEDFALL{\char"004D\char"0312} \MISSINGGLYPH{\char"2C6E\char"0312} \SUPPORTEDFALL{\char"0050\char"0312} \SUPPORTEDFALL{\char"01A4\char"0312} \SUPPORTEDANCH{\char"0052\char"0312} \SUPPORTEDFALL{\char"2C64\char"0312} \SUPPORTEDANCH{\char"0054\char"0312} \SUPPORTEDFALL{\char"01AC\char"0312} \SUPPORTEDFALL{\char"01AE\char"0312} \SUPPORTEDANCH{\char"0059\char"0312} \SUPPORTEDFALL{\char"01B3\char"0312} \SUPPORTEDFALL{\char"0061\char"0312} \SUPPORTEDFALL{\char"1E9A\char"0312} \SUPPORTEDFALL{\char"0062\char"0312} \SUPPORTEDFALL{\char"0183\char"0312} \SUPPORTEDFALL{\char"0063\char"0312} \SUPPORTEDFALL{\char"0188\char"0312} \SUPPORTEDFALL{\char"0064\char"0312} \SUPPORTEDFALL{\char"018C\char"0312} \SUPPORTEDFALL{\char"0221\char"0312} \SUPPORTEDFALL{\char"0068\char"0312} \SUPPORTEDFALL{\char"2C68\char"0312} \SUPPORTEDFALL{\char"006B\char"0312} \SUPPORTEDFALL{\char"0199\char"0312} \SUPPORTEDFALL{\char"2C6A\char"0312} \SUPPORTEDANCH{\char"006C\char"0312} \SUPPORTEDFALL{\char"0234\char"0312} \SUPPORTEDANCH{\char"006E\char"0312} \SUPPORTEDFALL{\char"0235\char"0312} \SUPPORTEDFALL{\char"0070\char"0312} \SUPPORTEDFALL{\char"01A5\char"0312} \SUPPORTEDFALL{\char"0071\char"0312} \SUPPORTEDFALL{\char"02A0\char"0312} \SUPPORTEDFALL{\char"0074\char"0312} \SUPPORTEDFALL{\char"01AB\char"0312} \SUPPORTEDFALL{\char"01AD\char"0312} \SUPPORTEDFALL{\char"0236\char"0312} \SUPPORTEDANCH{\char"0075\char"0312} \MISSINGGLYPH{\char"AB52\char"0312} \SUPPORTEDFALL{\char"0079\char"0312} \SUPPORTEDFALL{\char"01B4\char"0312} \MISSINGGLYPH{\char"1EFF\char"0312}

\MISSINGGLYPH{\char"0048} \MISSINGGLYPH{\char"2C67} \MISSINGGLYPH{\char"004B} \MISSINGGLYPH{\char"2C69} \MISSINGGLYPH{\char"004D} \MISSINGGLYPH{\char"2C6E} \MISSINGGLYPH{\char"0050} \MISSINGGLYPH{\char"01A4} \MISSINGGLYPH{\char"0052} \MISSINGGLYPH{\char"2C64} \MISSINGGLYPH{\char"0054} \MISSINGGLYPH{\char"01AC} \MISSINGGLYPH{\char"01AE} \MISSINGGLYPH{\char"0059} \MISSINGGLYPH{\char"01B3} \MISSINGGLYPH{\char"0061} \MISSINGGLYPH{\char"1E9A} \MISSINGGLYPH{\char"0062} \MISSINGGLYPH{\char"0183} \MISSINGGLYPH{\char"0063} \MISSINGGLYPH{\char"0188} \MISSINGGLYPH{\char"0064} \MISSINGGLYPH{\char"018C} \MISSINGGLYPH{\char"0221} \MISSINGGLYPH{\char"0068} \MISSINGGLYPH{\char"2C68} \MISSINGGLYPH{\char"006B} \MISSINGGLYPH{\char"0199} \MISSINGGLYPH{\char"2C6A} \MISSINGGLYPH{\char"006C} \MISSINGGLYPH{\char"0234} \MISSINGGLYPH{\char"006E} \MISSINGGLYPH{\char"0235} \MISSINGGLYPH{\char"0070} \MISSINGGLYPH{\char"01A5} \MISSINGGLYPH{\char"0071} \MISSINGGLYPH{\char"02A0} \MISSINGGLYPH{\char"0074} \MISSINGGLYPH{\char"01AB} \MISSINGGLYPH{\char"01AD} \MISSINGGLYPH{\char"0236} \MISSINGGLYPH{\char"0075} \MISSINGGLYPH{\char"AB52} \MISSINGGLYPH{\char"0079} \MISSINGGLYPH{\char"01B4} \MISSINGGLYPH{\char"1EFF}

\MISSINGGLYPH{\char"0048} \MISSINGGLYPH{\char"2C67} \MISSINGGLYPH{\char"004B} \MISSINGGLYPH{\char"2C69} \MISSINGGLYPH{\char"004D} \MISSINGGLYPH{\char"2C6E} \MISSINGGLYPH{\char"0050} \MISSINGGLYPH{\char"01A4} \MISSINGGLYPH{\char"0052} \MISSINGGLYPH{\char"2C64} \MISSINGGLYPH{\char"0054} \MISSINGGLYPH{\char"01AC} \MISSINGGLYPH{\char"01AE} \MISSINGGLYPH{\char"0059} \MISSINGGLYPH{\char"01B3} \MISSINGGLYPH{\char"0061} \MISSINGGLYPH{\char"1E9A} \MISSINGGLYPH{\char"0062} \MISSINGGLYPH{\char"0183} \MISSINGGLYPH{\char"0063} \MISSINGGLYPH{\char"0188} \MISSINGGLYPH{\char"0064} \MISSINGGLYPH{\char"018C} \MISSINGGLYPH{\char"0221} \MISSINGGLYPH{\char"0068} \MISSINGGLYPH{\char"2C68} \MISSINGGLYPH{\char"006B} \MISSINGGLYPH{\char"0199} \MISSINGGLYPH{\char"2C6A} \MISSINGGLYPH{\char"006C} \MISSINGGLYPH{\char"0234} \MISSINGGLYPH{\char"006E} \MISSINGGLYPH{\char"0235} \MISSINGGLYPH{\char"0070} \MISSINGGLYPH{\char"01A5} \MISSINGGLYPH{\char"0071} \MISSINGGLYPH{\char"02A0} \MISSINGGLYPH{\char"0074} \MISSINGGLYPH{\char"01AB} \MISSINGGLYPH{\char"01AD} \MISSINGGLYPH{\char"0236} \MISSINGGLYPH{\char"0075} \MISSINGGLYPH{\char"AB52} \MISSINGGLYPH{\char"0079} \MISSINGGLYPH{\char"01B4} \MISSINGGLYPH{\char"1EFF}

\MISSINGGLYPH{\char"0048} \MISSINGGLYPH{\char"2C67} \MISSINGGLYPH{\char"004B} \MISSINGGLYPH{\char"2C69} \MISSINGGLYPH{\char"004D} \MISSINGGLYPH{\char"2C6E} \MISSINGGLYPH{\char"0050} \MISSINGGLYPH{\char"01A4} \MISSINGGLYPH{\char"0052} \MISSINGGLYPH{\char"2C64} \MISSINGGLYPH{\char"0054} \MISSINGGLYPH{\char"01AC} \MISSINGGLYPH{\char"01AE} \MISSINGGLYPH{\char"0059} \MISSINGGLYPH{\char"01B3} \MISSINGGLYPH{\char"0061} \MISSINGGLYPH{\char"1E9A} \MISSINGGLYPH{\char"0062} \MISSINGGLYPH{\char"0183} \MISSINGGLYPH{\char"0063} \MISSINGGLYPH{\char"0188} \MISSINGGLYPH{\char"0064} \MISSINGGLYPH{\char"018C} \MISSINGGLYPH{\char"0221} \MISSINGGLYPH{\char"0068} \MISSINGGLYPH{\char"2C68} \MISSINGGLYPH{\char"006B} \MISSINGGLYPH{\char"0199} \MISSINGGLYPH{\char"2C6A} \MISSINGGLYPH{\char"006C} \MISSINGGLYPH{\char"0234} \MISSINGGLYPH{\char"006E} \MISSINGGLYPH{\char"0235} \MISSINGGLYPH{\char"0070} \MISSINGGLYPH{\char"01A5} \MISSINGGLYPH{\char"0071} \MISSINGGLYPH{\char"02A0} \MISSINGGLYPH{\char"0074} \MISSINGGLYPH{\char"01AB} \MISSINGGLYPH{\char"01AD} \MISSINGGLYPH{\char"0236} \MISSINGGLYPH{\char"0075} \MISSINGGLYPH{\char"AB52} \MISSINGGLYPH{\char"0079} \MISSINGGLYPH{\char"01B4} \MISSINGGLYPH{\char"1EFF}

\SUPPORTEDFALL{\char"0048\char"030E} \SUPPORTEDFALL{\char"2C67\char"030E} \SUPPORTEDFALL{\char"004B\char"030E} \SUPPORTEDFALL{\char"2C69\char"030E} \SUPPORTEDFALL{\char"004D\char"030E} \MISSINGGLYPH{\char"2C6E\char"030E} \SUPPORTEDFALL{\char"0050\char"030E} \SUPPORTEDFALL{\char"01A4\char"030E} \SUPPORTEDANCH{\char"0052\char"030E} \SUPPORTEDFALL{\char"2C64\char"030E} \SUPPORTEDANCH{\char"0054\char"030E} \SUPPORTEDFALL{\char"01AC\char"030E} \SUPPORTEDFALL{\char"01AE\char"030E} \SUPPORTEDANCH{\char"0059\char"030E} \SUPPORTEDFALL{\char"01B3\char"030E} \SUPPORTEDFALL{\char"0061\char"030E} \SUPPORTEDFALL{\char"1E9A\char"030E} \SUPPORTEDFALL{\char"0062\char"030E} \SUPPORTEDFALL{\char"0183\char"030E} \SUPPORTEDFALL{\char"0063\char"030E} \SUPPORTEDFALL{\char"0188\char"030E} \SUPPORTEDFALL{\char"0064\char"030E} \SUPPORTEDFALL{\char"018C\char"030E} \SUPPORTEDFALL{\char"0221\char"030E} \SUPPORTEDFALL{\char"0068\char"030E} \SUPPORTEDFALL{\char"2C68\char"030E} \SUPPORTEDFALL{\char"006B\char"030E} \SUPPORTEDFALL{\char"0199\char"030E} \SUPPORTEDFALL{\char"2C6A\char"030E} \SUPPORTEDANCH{\char"006C\char"030E} \SUPPORTEDFALL{\char"0234\char"030E} \SUPPORTEDANCH{\char"006E\char"030E} \SUPPORTEDFALL{\char"0235\char"030E} \SUPPORTEDFALL{\char"0070\char"030E} \SUPPORTEDFALL{\char"01A5\char"030E} \SUPPORTEDFALL{\char"0071\char"030E} \SUPPORTEDFALL{\char"02A0\char"030E} \SUPPORTEDFALL{\char"0074\char"030E} \SUPPORTEDFALL{\char"01AB\char"030E} \SUPPORTEDFALL{\char"01AD\char"030E} \SUPPORTEDFALL{\char"0236\char"030E} \SUPPORTEDANCH{\char"0075\char"030E} \MISSINGGLYPH{\char"AB52\char"030E} \SUPPORTEDFALL{\char"0079\char"030E} \SUPPORTEDFALL{\char"01B4\char"030E} \MISSINGGLYPH{\char"1EFF\char"030E}

\MISSINGGLYPH{\char"0048} \MISSINGGLYPH{\char"2C67} \MISSINGGLYPH{\char"004B} \MISSINGGLYPH{\char"2C69} \MISSINGGLYPH{\char"004D} \MISSINGGLYPH{\char"2C6E} \MISSINGGLYPH{\char"0050} \MISSINGGLYPH{\char"01A4} \MISSINGGLYPH{\char"0052} \MISSINGGLYPH{\char"2C64} \MISSINGGLYPH{\char"0054} \MISSINGGLYPH{\char"01AC} \MISSINGGLYPH{\char"01AE} \MISSINGGLYPH{\char"0059} \MISSINGGLYPH{\char"01B3} \MISSINGGLYPH{\char"0061} \MISSINGGLYPH{\char"1E9A} \MISSINGGLYPH{\char"0062} \MISSINGGLYPH{\char"0183} \MISSINGGLYPH{\char"0063} \MISSINGGLYPH{\char"0188} \MISSINGGLYPH{\char"0064} \MISSINGGLYPH{\char"018C} \MISSINGGLYPH{\char"0221} \MISSINGGLYPH{\char"0068} \MISSINGGLYPH{\char"2C68} \MISSINGGLYPH{\char"006B} \MISSINGGLYPH{\char"0199} \MISSINGGLYPH{\char"2C6A} \MISSINGGLYPH{\char"006C} \MISSINGGLYPH{\char"0234} \MISSINGGLYPH{\char"006E} \MISSINGGLYPH{\char"0235} \MISSINGGLYPH{\char"0070} \MISSINGGLYPH{\char"01A5} \MISSINGGLYPH{\char"0071} \MISSINGGLYPH{\char"02A0} \MISSINGGLYPH{\char"0074} \MISSINGGLYPH{\char"01AB} \MISSINGGLYPH{\char"01AD} \MISSINGGLYPH{\char"0236} \MISSINGGLYPH{\char"0075} \MISSINGGLYPH{\char"AB52} \MISSINGGLYPH{\char"0079} \MISSINGGLYPH{\char"01B4} \MISSINGGLYPH{\char"1EFF}

\MISSINGGLYPH{\char"0048} \MISSINGGLYPH{\char"2C67} \MISSINGGLYPH{\char"004B} \MISSINGGLYPH{\char"2C69} \MISSINGGLYPH{\char"004D} \MISSINGGLYPH{\char"2C6E} \MISSINGGLYPH{\char"0050} \MISSINGGLYPH{\char"01A4} \MISSINGGLYPH{\char"0052} \MISSINGGLYPH{\char"2C64} \MISSINGGLYPH{\char"0054} \MISSINGGLYPH{\char"01AC} \MISSINGGLYPH{\char"01AE} \MISSINGGLYPH{\char"0059} \MISSINGGLYPH{\char"01B3} \MISSINGGLYPH{\char"0061} \MISSINGGLYPH{\char"1E9A} \MISSINGGLYPH{\char"0062} \MISSINGGLYPH{\char"0183} \MISSINGGLYPH{\char"0063} \MISSINGGLYPH{\char"0188} \MISSINGGLYPH{\char"0064} \MISSINGGLYPH{\char"018C} \MISSINGGLYPH{\char"0221} \MISSINGGLYPH{\char"0068} \MISSINGGLYPH{\char"2C68} \MISSINGGLYPH{\char"006B} \MISSINGGLYPH{\char"0199} \MISSINGGLYPH{\char"2C6A} \MISSINGGLYPH{\char"006C} \MISSINGGLYPH{\char"0234} \MISSINGGLYPH{\char"006E} \MISSINGGLYPH{\char"0235} \MISSINGGLYPH{\char"0070} \MISSINGGLYPH{\char"01A5} \MISSINGGLYPH{\char"0071} \MISSINGGLYPH{\char"02A0} \MISSINGGLYPH{\char"0074} \MISSINGGLYPH{\char"01AB} \MISSINGGLYPH{\char"01AD} \MISSINGGLYPH{\char"0236} \MISSINGGLYPH{\char"0075} \MISSINGGLYPH{\char"AB52} \MISSINGGLYPH{\char"0079} \MISSINGGLYPH{\char"01B4} \MISSINGGLYPH{\char"1EFF}

\MISSINGGLYPH{\char"0048} \MISSINGGLYPH{\char"2C67} \MISSINGGLYPH{\char"004B} \MISSINGGLYPH{\char"2C69} \MISSINGGLYPH{\char"004D} \MISSINGGLYPH{\char"2C6E} \MISSINGGLYPH{\char"0050} \MISSINGGLYPH{\char"01A4} \MISSINGGLYPH{\char"0052} \MISSINGGLYPH{\char"2C64} \MISSINGGLYPH{\char"0054} \MISSINGGLYPH{\char"01AC} \MISSINGGLYPH{\char"01AE} \MISSINGGLYPH{\char"0059} \MISSINGGLYPH{\char"01B3} \MISSINGGLYPH{\char"0061} \MISSINGGLYPH{\char"1E9A} \MISSINGGLYPH{\char"0062} \MISSINGGLYPH{\char"0183} \MISSINGGLYPH{\char"0063} \MISSINGGLYPH{\char"0188} \MISSINGGLYPH{\char"0064} \MISSINGGLYPH{\char"018C} \MISSINGGLYPH{\char"0221} \MISSINGGLYPH{\char"0068} \MISSINGGLYPH{\char"2C68} \MISSINGGLYPH{\char"006B} \MISSINGGLYPH{\char"0199} \MISSINGGLYPH{\char"2C6A} \MISSINGGLYPH{\char"006C} \MISSINGGLYPH{\char"0234} \MISSINGGLYPH{\char"006E} \MISSINGGLYPH{\char"0235} \MISSINGGLYPH{\char"0070} \MISSINGGLYPH{\char"01A5} \MISSINGGLYPH{\char"0071} \MISSINGGLYPH{\char"02A0} \MISSINGGLYPH{\char"0074} \MISSINGGLYPH{\char"01AB} \MISSINGGLYPH{\char"01AD} \MISSINGGLYPH{\char"0236} \MISSINGGLYPH{\char"0075} \MISSINGGLYPH{\char"AB52} \MISSINGGLYPH{\char"0079} \MISSINGGLYPH{\char"01B4} \MISSINGGLYPH{\char"1EFF}

\MISSINGGLYPH{\char"0048} \MISSINGGLYPH{\char"2C67} \MISSINGGLYPH{\char"004B} \MISSINGGLYPH{\char"2C69} \MISSINGGLYPH{\char"004D} \MISSINGGLYPH{\char"2C6E} \MISSINGGLYPH{\char"0050} \MISSINGGLYPH{\char"01A4} \MISSINGGLYPH{\char"0052} \MISSINGGLYPH{\char"2C64} \MISSINGGLYPH{\char"0054} \MISSINGGLYPH{\char"01AC} \MISSINGGLYPH{\char"01AE} \MISSINGGLYPH{\char"0059} \MISSINGGLYPH{\char"01B3} \MISSINGGLYPH{\char"0061} \MISSINGGLYPH{\char"1E9A} \MISSINGGLYPH{\char"0062} \MISSINGGLYPH{\char"0183} \MISSINGGLYPH{\char"0063} \MISSINGGLYPH{\char"0188} \MISSINGGLYPH{\char"0064} \MISSINGGLYPH{\char"018C} \MISSINGGLYPH{\char"0221} \MISSINGGLYPH{\char"0068} \MISSINGGLYPH{\char"2C68} \MISSINGGLYPH{\char"006B} \MISSINGGLYPH{\char"0199} \MISSINGGLYPH{\char"2C6A} \MISSINGGLYPH{\char"006C} \MISSINGGLYPH{\char"0234} \MISSINGGLYPH{\char"006E} \MISSINGGLYPH{\char"0235} \MISSINGGLYPH{\char"0070} \MISSINGGLYPH{\char"01A5} \MISSINGGLYPH{\char"0071} \MISSINGGLYPH{\char"02A0} \MISSINGGLYPH{\char"0074} \MISSINGGLYPH{\char"01AB} \MISSINGGLYPH{\char"01AD} \MISSINGGLYPH{\char"0236} \MISSINGGLYPH{\char"0075} \MISSINGGLYPH{\char"AB52} \MISSINGGLYPH{\char"0079} \MISSINGGLYPH{\char"01B4} \MISSINGGLYPH{\char"1EFF}

\MISSINGGLYPH{\char"0048} \MISSINGGLYPH{\char"2C67} \MISSINGGLYPH{\char"004B} \MISSINGGLYPH{\char"2C69} \MISSINGGLYPH{\char"004D} \MISSINGGLYPH{\char"2C6E} \MISSINGGLYPH{\char"0050} \MISSINGGLYPH{\char"01A4} \MISSINGGLYPH{\char"0052} \MISSINGGLYPH{\char"2C64} \MISSINGGLYPH{\char"0054} \MISSINGGLYPH{\char"01AC} \MISSINGGLYPH{\char"01AE} \MISSINGGLYPH{\char"0059} \MISSINGGLYPH{\char"01B3} \MISSINGGLYPH{\char"0061} \MISSINGGLYPH{\char"1E9A} \MISSINGGLYPH{\char"0062} \MISSINGGLYPH{\char"0183} \MISSINGGLYPH{\char"0063} \MISSINGGLYPH{\char"0188} \MISSINGGLYPH{\char"0064} \MISSINGGLYPH{\char"018C} \MISSINGGLYPH{\char"0221} \MISSINGGLYPH{\char"0068} \MISSINGGLYPH{\char"2C68} \MISSINGGLYPH{\char"006B} \MISSINGGLYPH{\char"0199} \MISSINGGLYPH{\char"2C6A} \MISSINGGLYPH{\char"006C} \MISSINGGLYPH{\char"0234} \MISSINGGLYPH{\char"006E} \MISSINGGLYPH{\char"0235} \MISSINGGLYPH{\char"0070} \MISSINGGLYPH{\char"01A5} \MISSINGGLYPH{\char"0071} \MISSINGGLYPH{\char"02A0} \MISSINGGLYPH{\char"0074} \MISSINGGLYPH{\char"01AB} \MISSINGGLYPH{\char"01AD} \MISSINGGLYPH{\char"0236} \MISSINGGLYPH{\char"0075} \MISSINGGLYPH{\char"AB52} \MISSINGGLYPH{\char"0079} \MISSINGGLYPH{\char"01B4} \MISSINGGLYPH{\char"1EFF}

\SUPPORTEDFALL{\char"0048\char"0351} \SUPPORTEDFALL{\char"2C67\char"0351} \SUPPORTEDFALL{\char"004B\char"0351} \SUPPORTEDFALL{\char"2C69\char"0351} \SUPPORTEDFALL{\char"004D\char"0351} \MISSINGGLYPH{\char"2C6E\char"0351} \SUPPORTEDFALL{\char"0050\char"0351} \SUPPORTEDFALL{\char"01A4\char"0351} \SUPPORTEDANCH{\char"0052\char"0351} \SUPPORTEDFALL{\char"2C64\char"0351} \SUPPORTEDANCH{\char"0054\char"0351} \SUPPORTEDFALL{\char"01AC\char"0351} \SUPPORTEDFALL{\char"01AE\char"0351} \SUPPORTEDANCH{\char"0059\char"0351} \SUPPORTEDFALL{\char"01B3\char"0351} \SUPPORTEDFALL{\char"0061\char"0351} \SUPPORTEDFALL{\char"1E9A\char"0351} \SUPPORTEDFALL{\char"0062\char"0351} \SUPPORTEDFALL{\char"0183\char"0351} \SUPPORTEDFALL{\char"0063\char"0351} \SUPPORTEDFALL{\char"0188\char"0351} \SUPPORTEDFALL{\char"0064\char"0351} \SUPPORTEDFALL{\char"018C\char"0351} \SUPPORTEDFALL{\char"0221\char"0351} \SUPPORTEDFALL{\char"0068\char"0351} \SUPPORTEDFALL{\char"2C68\char"0351} \SUPPORTEDFALL{\char"006B\char"0351} \SUPPORTEDFALL{\char"0199\char"0351} \SUPPORTEDFALL{\char"2C6A\char"0351} \SUPPORTEDANCH{\char"006C\char"0351} \SUPPORTEDFALL{\char"0234\char"0351} \SUPPORTEDANCH{\char"006E\char"0351} \SUPPORTEDFALL{\char"0235\char"0351} \SUPPORTEDFALL{\char"0070\char"0351} \SUPPORTEDFALL{\char"01A5\char"0351} \SUPPORTEDFALL{\char"0071\char"0351} \SUPPORTEDFALL{\char"02A0\char"0351} \SUPPORTEDFALL{\char"0074\char"0351} \SUPPORTEDFALL{\char"01AB\char"0351} \SUPPORTEDFALL{\char"01AD\char"0351} \SUPPORTEDFALL{\char"0236\char"0351} \SUPPORTEDANCH{\char"0075\char"0351} \MISSINGGLYPH{\char"AB52\char"0351} \SUPPORTEDFALL{\char"0079\char"0351} \SUPPORTEDFALL{\char"01B4\char"0351} \MISSINGGLYPH{\char"1EFF\char"0351}

\MISSINGGLYPH{\char"0048} \MISSINGGLYPH{\char"2C67} \MISSINGGLYPH{\char"004B} \MISSINGGLYPH{\char"2C69} \MISSINGGLYPH{\char"004D} \MISSINGGLYPH{\char"2C6E} \MISSINGGLYPH{\char"0050} \MISSINGGLYPH{\char"01A4} \MISSINGGLYPH{\char"0052} \MISSINGGLYPH{\char"2C64} \MISSINGGLYPH{\char"0054} \MISSINGGLYPH{\char"01AC} \MISSINGGLYPH{\char"01AE} \MISSINGGLYPH{\char"0059} \MISSINGGLYPH{\char"01B3} \MISSINGGLYPH{\char"0061} \MISSINGGLYPH{\char"1E9A} \MISSINGGLYPH{\char"0062} \MISSINGGLYPH{\char"0183} \MISSINGGLYPH{\char"0063} \MISSINGGLYPH{\char"0188} \MISSINGGLYPH{\char"0064} \MISSINGGLYPH{\char"018C} \MISSINGGLYPH{\char"0221} \MISSINGGLYPH{\char"0068} \MISSINGGLYPH{\char"2C68} \MISSINGGLYPH{\char"006B} \MISSINGGLYPH{\char"0199} \MISSINGGLYPH{\char"2C6A} \MISSINGGLYPH{\char"006C} \MISSINGGLYPH{\char"0234} \MISSINGGLYPH{\char"006E} \MISSINGGLYPH{\char"0235} \MISSINGGLYPH{\char"0070} \MISSINGGLYPH{\char"01A5} \MISSINGGLYPH{\char"0071} \MISSINGGLYPH{\char"02A0} \MISSINGGLYPH{\char"0074} \MISSINGGLYPH{\char"01AB} \MISSINGGLYPH{\char"01AD} \MISSINGGLYPH{\char"0236} \MISSINGGLYPH{\char"0075} \MISSINGGLYPH{\char"AB52} \MISSINGGLYPH{\char"0079} \MISSINGGLYPH{\char"01B4} \MISSINGGLYPH{\char"1EFF}

\SUPPORTEDFALL{\char"0048\char"0357} \SUPPORTEDFALL{\char"2C67\char"0357} \SUPPORTEDFALL{\char"004B\char"0357} \SUPPORTEDFALL{\char"2C69\char"0357} \SUPPORTEDFALL{\char"004D\char"0357} \MISSINGGLYPH{\char"2C6E\char"0357} \SUPPORTEDFALL{\char"0050\char"0357} \SUPPORTEDFALL{\char"01A4\char"0357} \SUPPORTEDANCH{\char"0052\char"0357} \SUPPORTEDFALL{\char"2C64\char"0357} \SUPPORTEDANCH{\char"0054\char"0357} \SUPPORTEDFALL{\char"01AC\char"0357} \SUPPORTEDFALL{\char"01AE\char"0357} \SUPPORTEDANCH{\char"0059\char"0357} \SUPPORTEDFALL{\char"01B3\char"0357} \SUPPORTEDFALL{\char"0061\char"0357} \SUPPORTEDFALL{\char"1E9A\char"0357} \SUPPORTEDFALL{\char"0062\char"0357} \SUPPORTEDFALL{\char"0183\char"0357} \SUPPORTEDFALL{\char"0063\char"0357} \SUPPORTEDFALL{\char"0188\char"0357} \SUPPORTEDFALL{\char"0064\char"0357} \SUPPORTEDFALL{\char"018C\char"0357} \SUPPORTEDFALL{\char"0221\char"0357} \SUPPORTEDFALL{\char"0068\char"0357} \SUPPORTEDFALL{\char"2C68\char"0357} \SUPPORTEDFALL{\char"006B\char"0357} \SUPPORTEDFALL{\char"0199\char"0357} \SUPPORTEDFALL{\char"2C6A\char"0357} \SUPPORTEDANCH{\char"006C\char"0357} \SUPPORTEDFALL{\char"0234\char"0357} \SUPPORTEDANCH{\char"006E\char"0357} \SUPPORTEDFALL{\char"0235\char"0357} \SUPPORTEDFALL{\char"0070\char"0357} \SUPPORTEDFALL{\char"01A5\char"0357} \SUPPORTEDFALL{\char"0071\char"0357} \SUPPORTEDFALL{\char"02A0\char"0357} \SUPPORTEDFALL{\char"0074\char"0357} \SUPPORTEDFALL{\char"01AB\char"0357} \SUPPORTEDFALL{\char"01AD\char"0357} \SUPPORTEDFALL{\char"0236\char"0357} \SUPPORTEDANCH{\char"0075\char"0357} \MISSINGGLYPH{\char"AB52\char"0357} \SUPPORTEDFALL{\char"0079\char"0357} \SUPPORTEDFALL{\char"01B4\char"0357} \MISSINGGLYPH{\char"1EFF\char"0357}

\SUPPORTEDFALL{\char"0048\char"035B} \SUPPORTEDFALL{\char"2C67\char"035B} \SUPPORTEDFALL{\char"004B\char"035B} \SUPPORTEDFALL{\char"2C69\char"035B} \SUPPORTEDFALL{\char"004D\char"035B} \MISSINGGLYPH{\char"2C6E\char"035B} \SUPPORTEDFALL{\char"0050\char"035B} \SUPPORTEDFALL{\char"01A4\char"035B} \SUPPORTEDANCH{\char"0052\char"035B} \SUPPORTEDFALL{\char"2C64\char"035B} \SUPPORTEDANCH{\char"0054\char"035B} \SUPPORTEDFALL{\char"01AC\char"035B} \SUPPORTEDFALL{\char"01AE\char"035B} \SUPPORTEDANCH{\char"0059\char"035B} \SUPPORTEDFALL{\char"01B3\char"035B} \SUPPORTEDFALL{\char"0061\char"035B} \SUPPORTEDFALL{\char"1E9A\char"035B} \SUPPORTEDFALL{\char"0062\char"035B} \SUPPORTEDFALL{\char"0183\char"035B} \SUPPORTEDFALL{\char"0063\char"035B} \SUPPORTEDFALL{\char"0188\char"035B} \SUPPORTEDFALL{\char"0064\char"035B} \SUPPORTEDFALL{\char"018C\char"035B} \SUPPORTEDFALL{\char"0221\char"035B} \SUPPORTEDFALL{\char"0068\char"035B} \SUPPORTEDFALL{\char"2C68\char"035B} \SUPPORTEDFALL{\char"006B\char"035B} \SUPPORTEDFALL{\char"0199\char"035B} \SUPPORTEDFALL{\char"2C6A\char"035B} \SUPPORTEDANCH{\char"006C\char"035B} \SUPPORTEDFALL{\char"0234\char"035B} \SUPPORTEDANCH{\char"006E\char"035B} \SUPPORTEDFALL{\char"0235\char"035B} \SUPPORTEDFALL{\char"0070\char"035B} \SUPPORTEDFALL{\char"01A5\char"035B} \SUPPORTEDFALL{\char"0071\char"035B} \SUPPORTEDFALL{\char"02A0\char"035B} \SUPPORTEDFALL{\char"0074\char"035B} \SUPPORTEDFALL{\char"01AB\char"035B} \SUPPORTEDFALL{\char"01AD\char"035B} \SUPPORTEDFALL{\char"0236\char"035B} \SUPPORTEDANCH{\char"0075\char"035B} \MISSINGGLYPH{\char"AB52\char"035B} \SUPPORTEDFALL{\char"0079\char"035B} \SUPPORTEDFALL{\char"01B4\char"035B} \MISSINGGLYPH{\char"1EFF\char"035B}

}

\newpage
\subsection*{Semibold patched, Group 4, above, Above}

{\LibertinusSerifPatch\bfseries
% grid. columns are bases, rows are marks.
\SUPPORTEDANCH{\char"0048\char"0300} \SUPPORTEDANCH{\char"2C67\char"0300} \SUPPORTEDANCH{\char"004B\char"0300} \SUPPORTEDANCH{\char"2C69\char"0300} \SUPPORTEDANCH{\char"004D\char"0300} \MISSINGGLYPH{\char"2C6E\char"0300} \SUPPORTEDANCH{\char"0050\char"0300} \SUPPORTEDANCH{\char"01A4\char"0300} \SUPPORTEDANCH{\char"0052\char"0300} \SUPPORTEDANCH{\char"2C64\char"0300} \SUPPORTEDANCH{\char"0054\char"0300} \SUPPORTEDANCH{\char"01AC\char"0300} \SUPPORTEDANCH{\char"01AE\char"0300} \SUPPORTEDPREC{\char"0059\char"0300} \SUPPORTEDANCH{\char"01B3\char"0300} \SUPPORTEDPREC{\char"0061\char"0300} \SUPPORTEDANCH{\char"1E9A\char"0300} \SUPPORTEDANCH{\char"0062\char"0300} \SUPPORTEDANCH{\char"0183\char"0300} \SUPPORTEDANCH{\char"0063\char"0300} \SUPPORTEDANCH{\char"0188\char"0300} \SUPPORTEDANCH{\char"0064\char"0300} \SUPPORTEDANCH{\char"018C\char"0300} \SUPPORTEDANCH{\char"0221\char"0300} \SUPPORTEDANCH{\char"0068\char"0300} \SUPPORTEDANCH{\char"2C68\char"0300} \SUPPORTEDANCH{\char"006B\char"0300} \SUPPORTEDANCH{\char"0199\char"0300} \SUPPORTEDANCH{\char"2C6A\char"0300} \SUPPORTEDANCH{\char"006C\char"0300} \SUPPORTEDANCH{\char"0234\char"0300} \SUPPORTEDPREC{\char"006E\char"0300} \SUPPORTEDANCH{\char"0235\char"0300} \SUPPORTEDANCH{\char"0070\char"0300} \SUPPORTEDANCH{\char"01A5\char"0300} \SUPPORTEDANCH{\char"0071\char"0300} \SUPPORTEDANCH{\char"02A0\char"0300} \SUPPORTEDANCH{\char"0074\char"0300} \SUPPORTEDANCH{\char"01AB\char"0300} \SUPPORTEDANCH{\char"01AD\char"0300} \SUPPORTEDANCH{\char"0236\char"0300} \SUPPORTEDPREC{\char"0075\char"0300} \MISSINGGLYPH{\char"AB52\char"0300} \SUPPORTEDPREC{\char"0079\char"0300} \SUPPORTEDANCH{\char"01B4\char"0300} \MISSINGGLYPH{\char"1EFF\char"0300}

\SUPPORTEDANCH{\char"0048\char"0301} \SUPPORTEDANCH{\char"2C67\char"0301} \SUPPORTEDPREC{\char"004B\char"0301} \SUPPORTEDANCH{\char"2C69\char"0301} \SUPPORTEDPREC{\char"004D\char"0301} \MISSINGGLYPH{\char"2C6E\char"0301} \SUPPORTEDPREC{\char"0050\char"0301} \SUPPORTEDANCH{\char"01A4\char"0301} \SUPPORTEDPREC{\char"0052\char"0301} \SUPPORTEDANCH{\char"2C64\char"0301} \SUPPORTEDANCH{\char"0054\char"0301} \SUPPORTEDANCH{\char"01AC\char"0301} \SUPPORTEDANCH{\char"01AE\char"0301} \SUPPORTEDPREC{\char"0059\char"0301} \SUPPORTEDANCH{\char"01B3\char"0301} \SUPPORTEDPREC{\char"0061\char"0301} \SUPPORTEDANCH{\char"1E9A\char"0301} \SUPPORTEDANCH{\char"0062\char"0301} \SUPPORTEDANCH{\char"0183\char"0301} \SUPPORTEDPREC{\char"0063\char"0301} \SUPPORTEDANCH{\char"0188\char"0301} \SUPPORTEDANCH{\char"0064\char"0301} \SUPPORTEDANCH{\char"018C\char"0301} \SUPPORTEDANCH{\char"0221\char"0301} \SUPPORTEDANCH{\char"0068\char"0301} \SUPPORTEDANCH{\char"2C68\char"0301} \SUPPORTEDPREC{\char"006B\char"0301} \SUPPORTEDANCH{\char"0199\char"0301} \SUPPORTEDANCH{\char"2C6A\char"0301} \SUPPORTEDPREC{\char"006C\char"0301} \SUPPORTEDANCH{\char"0234\char"0301} \SUPPORTEDPREC{\char"006E\char"0301} \SUPPORTEDANCH{\char"0235\char"0301} \SUPPORTEDPREC{\char"0070\char"0301} \SUPPORTEDANCH{\char"01A5\char"0301} \SUPPORTEDANCH{\char"0071\char"0301} \SUPPORTEDANCH{\char"02A0\char"0301} \SUPPORTEDANCH{\char"0074\char"0301} \SUPPORTEDANCH{\char"01AB\char"0301} \SUPPORTEDANCH{\char"01AD\char"0301} \SUPPORTEDANCH{\char"0236\char"0301} \SUPPORTEDPREC{\char"0075\char"0301} \MISSINGGLYPH{\char"AB52\char"0301} \SUPPORTEDPREC{\char"0079\char"0301} \SUPPORTEDANCH{\char"01B4\char"0301} \MISSINGGLYPH{\char"1EFF\char"0301}

\SUPPORTEDPREC{\char"0048\char"0302} \SUPPORTEDANCH{\char"2C67\char"0302} \SUPPORTEDANCH{\char"004B\char"0302} \SUPPORTEDANCH{\char"2C69\char"0302} \SUPPORTEDANCH{\char"004D\char"0302} \MISSINGGLYPH{\char"2C6E\char"0302} \SUPPORTEDANCH{\char"0050\char"0302} \SUPPORTEDANCH{\char"01A4\char"0302} \SUPPORTEDANCH{\char"0052\char"0302} \SUPPORTEDANCH{\char"2C64\char"0302} \SUPPORTEDANCH{\char"0054\char"0302} \SUPPORTEDANCH{\char"01AC\char"0302} \SUPPORTEDANCH{\char"01AE\char"0302} \SUPPORTEDPREC{\char"0059\char"0302} \SUPPORTEDANCH{\char"01B3\char"0302} \SUPPORTEDPREC{\char"0061\char"0302} \SUPPORTEDANCH{\char"1E9A\char"0302} \SUPPORTEDANCH{\char"0062\char"0302} \SUPPORTEDANCH{\char"0183\char"0302} \SUPPORTEDPREC{\char"0063\char"0302} \SUPPORTEDANCH{\char"0188\char"0302} \SUPPORTEDANCH{\char"0064\char"0302} \SUPPORTEDANCH{\char"018C\char"0302} \SUPPORTEDANCH{\char"0221\char"0302} \SUPPORTEDPREC{\char"0068\char"0302} \SUPPORTEDANCH{\char"2C68\char"0302} \SUPPORTEDANCH{\char"006B\char"0302} \SUPPORTEDANCH{\char"0199\char"0302} \SUPPORTEDANCH{\char"2C6A\char"0302} \SUPPORTEDANCH{\char"006C\char"0302} \SUPPORTEDANCH{\char"0234\char"0302} \SUPPORTEDANCH{\char"006E\char"0302} \SUPPORTEDANCH{\char"0235\char"0302} \SUPPORTEDANCH{\char"0070\char"0302} \SUPPORTEDANCH{\char"01A5\char"0302} \SUPPORTEDANCH{\char"0071\char"0302} \SUPPORTEDANCH{\char"02A0\char"0302} \SUPPORTEDANCH{\char"0074\char"0302} \SUPPORTEDANCH{\char"01AB\char"0302} \SUPPORTEDANCH{\char"01AD\char"0302} \SUPPORTEDANCH{\char"0236\char"0302} \SUPPORTEDPREC{\char"0075\char"0302} \MISSINGGLYPH{\char"AB52\char"0302} \SUPPORTEDPREC{\char"0079\char"0302} \SUPPORTEDANCH{\char"01B4\char"0302} \MISSINGGLYPH{\char"1EFF\char"0302}

\SUPPORTEDANCH{\char"0048\char"0303} \SUPPORTEDANCH{\char"2C67\char"0303} \SUPPORTEDANCH{\char"004B\char"0303} \SUPPORTEDANCH{\char"2C69\char"0303} \SUPPORTEDANCH{\char"004D\char"0303} \MISSINGGLYPH{\char"2C6E\char"0303} \SUPPORTEDANCH{\char"0050\char"0303} \SUPPORTEDANCH{\char"01A4\char"0303} \SUPPORTEDANCH{\char"0052\char"0303} \SUPPORTEDANCH{\char"2C64\char"0303} \SUPPORTEDANCH{\char"0054\char"0303} \SUPPORTEDANCH{\char"01AC\char"0303} \SUPPORTEDANCH{\char"01AE\char"0303} \SUPPORTEDPREC{\char"0059\char"0303} \SUPPORTEDANCH{\char"01B3\char"0303} \SUPPORTEDPREC{\char"0061\char"0303} \SUPPORTEDANCH{\char"1E9A\char"0303} \SUPPORTEDANCH{\char"0062\char"0303} \SUPPORTEDANCH{\char"0183\char"0303} \SUPPORTEDANCH{\char"0063\char"0303} \SUPPORTEDANCH{\char"0188\char"0303} \SUPPORTEDANCH{\char"0064\char"0303} \SUPPORTEDANCH{\char"018C\char"0303} \SUPPORTEDANCH{\char"0221\char"0303} \SUPPORTEDANCH{\char"0068\char"0303} \SUPPORTEDANCH{\char"2C68\char"0303} \SUPPORTEDANCH{\char"006B\char"0303} \SUPPORTEDANCH{\char"0199\char"0303} \SUPPORTEDANCH{\char"2C6A\char"0303} \SUPPORTEDANCH{\char"006C\char"0303} \SUPPORTEDANCH{\char"0234\char"0303} \SUPPORTEDPREC{\char"006E\char"0303} \SUPPORTEDANCH{\char"0235\char"0303} \SUPPORTEDANCH{\char"0070\char"0303} \SUPPORTEDANCH{\char"01A5\char"0303} \SUPPORTEDANCH{\char"0071\char"0303} \SUPPORTEDANCH{\char"02A0\char"0303} \SUPPORTEDANCH{\char"0074\char"0303} \SUPPORTEDANCH{\char"01AB\char"0303} \SUPPORTEDANCH{\char"01AD\char"0303} \SUPPORTEDANCH{\char"0236\char"0303} \SUPPORTEDPREC{\char"0075\char"0303} \MISSINGGLYPH{\char"AB52\char"0303} \SUPPORTEDPREC{\char"0079\char"0303} \SUPPORTEDANCH{\char"01B4\char"0303} \MISSINGGLYPH{\char"1EFF\char"0303}

\SUPPORTEDANCH{\char"0048\char"0304} \SUPPORTEDANCH{\char"2C67\char"0304} \SUPPORTEDANCH{\char"004B\char"0304} \SUPPORTEDANCH{\char"2C69\char"0304} \SUPPORTEDANCH{\char"004D\char"0304} \MISSINGGLYPH{\char"2C6E\char"0304} \SUPPORTEDANCH{\char"0050\char"0304} \SUPPORTEDANCH{\char"01A4\char"0304} \SUPPORTEDANCH{\char"0052\char"0304} \SUPPORTEDANCH{\char"2C64\char"0304} \SUPPORTEDANCH{\char"0054\char"0304} \SUPPORTEDANCH{\char"01AC\char"0304} \SUPPORTEDANCH{\char"01AE\char"0304} \SUPPORTEDPREC{\char"0059\char"0304} \SUPPORTEDANCH{\char"01B3\char"0304} \SUPPORTEDPREC{\char"0061\char"0304} \SUPPORTEDANCH{\char"1E9A\char"0304} \SUPPORTEDANCH{\char"0062\char"0304} \SUPPORTEDANCH{\char"0183\char"0304} \SUPPORTEDANCH{\char"0063\char"0304} \SUPPORTEDANCH{\char"0188\char"0304} \SUPPORTEDANCH{\char"0064\char"0304} \SUPPORTEDANCH{\char"018C\char"0304} \SUPPORTEDANCH{\char"0221\char"0304} \SUPPORTEDANCH{\char"0068\char"0304} \SUPPORTEDANCH{\char"2C68\char"0304} \SUPPORTEDANCH{\char"006B\char"0304} \SUPPORTEDANCH{\char"0199\char"0304} \SUPPORTEDANCH{\char"2C6A\char"0304} \SUPPORTEDANCH{\char"006C\char"0304} \SUPPORTEDANCH{\char"0234\char"0304} \SUPPORTEDANCH{\char"006E\char"0304} \SUPPORTEDANCH{\char"0235\char"0304} \SUPPORTEDANCH{\char"0070\char"0304} \SUPPORTEDANCH{\char"01A5\char"0304} \SUPPORTEDANCH{\char"0071\char"0304} \SUPPORTEDANCH{\char"02A0\char"0304} \SUPPORTEDANCH{\char"0074\char"0304} \SUPPORTEDANCH{\char"01AB\char"0304} \SUPPORTEDANCH{\char"01AD\char"0304} \SUPPORTEDANCH{\char"0236\char"0304} \SUPPORTEDPREC{\char"0075\char"0304} \MISSINGGLYPH{\char"AB52\char"0304} \SUPPORTEDPREC{\char"0079\char"0304} \SUPPORTEDANCH{\char"01B4\char"0304} \MISSINGGLYPH{\char"1EFF\char"0304}

\SUPPORTEDANCH{\char"0048\char"0306} \SUPPORTEDANCH{\char"2C67\char"0306} \SUPPORTEDANCH{\char"004B\char"0306} \SUPPORTEDANCH{\char"2C69\char"0306} \SUPPORTEDANCH{\char"004D\char"0306} \MISSINGGLYPH{\char"2C6E\char"0306} \SUPPORTEDANCH{\char"0050\char"0306} \SUPPORTEDANCH{\char"01A4\char"0306} \SUPPORTEDANCH{\char"0052\char"0306} \SUPPORTEDANCH{\char"2C64\char"0306} \SUPPORTEDANCH{\char"0054\char"0306} \SUPPORTEDANCH{\char"01AC\char"0306} \SUPPORTEDANCH{\char"01AE\char"0306} \SUPPORTEDANCH{\char"0059\char"0306} \SUPPORTEDANCH{\char"01B3\char"0306} \SUPPORTEDPREC{\char"0061\char"0306} \SUPPORTEDANCH{\char"1E9A\char"0306} \SUPPORTEDANCH{\char"0062\char"0306} \SUPPORTEDANCH{\char"0183\char"0306} \SUPPORTEDANCH{\char"0063\char"0306} \SUPPORTEDANCH{\char"0188\char"0306} \SUPPORTEDANCH{\char"0064\char"0306} \SUPPORTEDANCH{\char"018C\char"0306} \SUPPORTEDANCH{\char"0221\char"0306} \SUPPORTEDANCH{\char"0068\char"0306} \SUPPORTEDANCH{\char"2C68\char"0306} \SUPPORTEDANCH{\char"006B\char"0306} \SUPPORTEDANCH{\char"0199\char"0306} \SUPPORTEDANCH{\char"2C6A\char"0306} \SUPPORTEDANCH{\char"006C\char"0306} \SUPPORTEDANCH{\char"0234\char"0306} \SUPPORTEDANCH{\char"006E\char"0306} \SUPPORTEDANCH{\char"0235\char"0306} \SUPPORTEDANCH{\char"0070\char"0306} \SUPPORTEDANCH{\char"01A5\char"0306} \SUPPORTEDANCH{\char"0071\char"0306} \SUPPORTEDANCH{\char"02A0\char"0306} \SUPPORTEDANCH{\char"0074\char"0306} \SUPPORTEDANCH{\char"01AB\char"0306} \SUPPORTEDANCH{\char"01AD\char"0306} \SUPPORTEDANCH{\char"0236\char"0306} \SUPPORTEDPREC{\char"0075\char"0306} \MISSINGGLYPH{\char"AB52\char"0306} \SUPPORTEDANCH{\char"0079\char"0306} \SUPPORTEDANCH{\char"01B4\char"0306} \MISSINGGLYPH{\char"1EFF\char"0306}

\SUPPORTEDPREC{\char"0048\char"0307} \SUPPORTEDANCH{\char"2C67\char"0307} \SUPPORTEDANCH{\char"004B\char"0307} \SUPPORTEDANCH{\char"2C69\char"0307} \SUPPORTEDPREC{\char"004D\char"0307} \MISSINGGLYPH{\char"2C6E\char"0307} \SUPPORTEDPREC{\char"0050\char"0307} \SUPPORTEDANCH{\char"01A4\char"0307} \SUPPORTEDPREC{\char"0052\char"0307} \SUPPORTEDANCH{\char"2C64\char"0307} \SUPPORTEDPREC{\char"0054\char"0307} \SUPPORTEDANCH{\char"01AC\char"0307} \SUPPORTEDANCH{\char"01AE\char"0307} \SUPPORTEDPREC{\char"0059\char"0307} \SUPPORTEDANCH{\char"01B3\char"0307} \SUPPORTEDPREC{\char"0061\char"0307} \SUPPORTEDANCH{\char"1E9A\char"0307} \SUPPORTEDPREC{\char"0062\char"0307} \SUPPORTEDANCH{\char"0183\char"0307} \SUPPORTEDPREC{\char"0063\char"0307} \SUPPORTEDANCH{\char"0188\char"0307} \SUPPORTEDPREC{\char"0064\char"0307} \SUPPORTEDANCH{\char"018C\char"0307} \SUPPORTEDANCH{\char"0221\char"0307} \SUPPORTEDPREC{\char"0068\char"0307} \SUPPORTEDANCH{\char"2C68\char"0307} \SUPPORTEDANCH{\char"006B\char"0307} \SUPPORTEDANCH{\char"0199\char"0307} \SUPPORTEDANCH{\char"2C6A\char"0307} \SUPPORTEDANCH{\char"006C\char"0307} \SUPPORTEDANCH{\char"0234\char"0307} \SUPPORTEDPREC{\char"006E\char"0307} \SUPPORTEDANCH{\char"0235\char"0307} \SUPPORTEDPREC{\char"0070\char"0307} \SUPPORTEDANCH{\char"01A5\char"0307} \SUPPORTEDANCH{\char"0071\char"0307} \SUPPORTEDANCH{\char"02A0\char"0307} \SUPPORTEDPREC{\char"0074\char"0307} \SUPPORTEDANCH{\char"01AB\char"0307} \SUPPORTEDANCH{\char"01AD\char"0307} \SUPPORTEDANCH{\char"0236\char"0307} \SUPPORTEDANCH{\char"0075\char"0307} \MISSINGGLYPH{\char"AB52\char"0307} \SUPPORTEDPREC{\char"0079\char"0307} \SUPPORTEDANCH{\char"01B4\char"0307} \MISSINGGLYPH{\char"1EFF\char"0307}

\SUPPORTEDPREC{\char"0048\char"0308} \SUPPORTEDANCH{\char"2C67\char"0308} \SUPPORTEDANCH{\char"004B\char"0308} \SUPPORTEDANCH{\char"2C69\char"0308} \SUPPORTEDANCH{\char"004D\char"0308} \MISSINGGLYPH{\char"2C6E\char"0308} \SUPPORTEDANCH{\char"0050\char"0308} \SUPPORTEDANCH{\char"01A4\char"0308} \SUPPORTEDANCH{\char"0052\char"0308} \SUPPORTEDANCH{\char"2C64\char"0308} \SUPPORTEDANCH{\char"0054\char"0308} \SUPPORTEDANCH{\char"01AC\char"0308} \SUPPORTEDANCH{\char"01AE\char"0308} \SUPPORTEDPREC{\char"0059\char"0308} \SUPPORTEDANCH{\char"01B3\char"0308} \SUPPORTEDPREC{\char"0061\char"0308} \SUPPORTEDANCH{\char"1E9A\char"0308} \SUPPORTEDANCH{\char"0062\char"0308} \SUPPORTEDANCH{\char"0183\char"0308} \SUPPORTEDANCH{\char"0063\char"0308} \SUPPORTEDANCH{\char"0188\char"0308} \SUPPORTEDANCH{\char"0064\char"0308} \SUPPORTEDANCH{\char"018C\char"0308} \SUPPORTEDANCH{\char"0221\char"0308} \SUPPORTEDPREC{\char"0068\char"0308} \SUPPORTEDANCH{\char"2C68\char"0308} \SUPPORTEDANCH{\char"006B\char"0308} \SUPPORTEDANCH{\char"0199\char"0308} \SUPPORTEDANCH{\char"2C6A\char"0308} \SUPPORTEDANCH{\char"006C\char"0308} \SUPPORTEDANCH{\char"0234\char"0308} \SUPPORTEDANCH{\char"006E\char"0308} \SUPPORTEDANCH{\char"0235\char"0308} \SUPPORTEDANCH{\char"0070\char"0308} \SUPPORTEDANCH{\char"01A5\char"0308} \SUPPORTEDANCH{\char"0071\char"0308} \SUPPORTEDANCH{\char"02A0\char"0308} \SUPPORTEDPREC{\char"0074\char"0308} \SUPPORTEDANCH{\char"01AB\char"0308} \SUPPORTEDANCH{\char"01AD\char"0308} \SUPPORTEDANCH{\char"0236\char"0308} \SUPPORTEDPREC{\char"0075\char"0308} \MISSINGGLYPH{\char"AB52\char"0308} \SUPPORTEDPREC{\char"0079\char"0308} \SUPPORTEDANCH{\char"01B4\char"0308} \MISSINGGLYPH{\char"1EFF\char"0308}

\SUPPORTEDANCH{\char"0048\char"0309} \SUPPORTEDANCH{\char"2C67\char"0309} \SUPPORTEDANCH{\char"004B\char"0309} \SUPPORTEDANCH{\char"2C69\char"0309} \SUPPORTEDANCH{\char"004D\char"0309} \MISSINGGLYPH{\char"2C6E\char"0309} \SUPPORTEDANCH{\char"0050\char"0309} \SUPPORTEDANCH{\char"01A4\char"0309} \SUPPORTEDANCH{\char"0052\char"0309} \SUPPORTEDANCH{\char"2C64\char"0309} \SUPPORTEDANCH{\char"0054\char"0309} \SUPPORTEDANCH{\char"01AC\char"0309} \SUPPORTEDANCH{\char"01AE\char"0309} \SUPPORTEDPREC{\char"0059\char"0309} \SUPPORTEDANCH{\char"01B3\char"0309} \SUPPORTEDPREC{\char"0061\char"0309} \SUPPORTEDANCH{\char"1E9A\char"0309} \SUPPORTEDANCH{\char"0062\char"0309} \SUPPORTEDANCH{\char"0183\char"0309} \SUPPORTEDANCH{\char"0063\char"0309} \SUPPORTEDANCH{\char"0188\char"0309} \SUPPORTEDANCH{\char"0064\char"0309} \SUPPORTEDANCH{\char"018C\char"0309} \SUPPORTEDANCH{\char"0221\char"0309} \SUPPORTEDANCH{\char"0068\char"0309} \SUPPORTEDANCH{\char"2C68\char"0309} \SUPPORTEDANCH{\char"006B\char"0309} \SUPPORTEDANCH{\char"0199\char"0309} \SUPPORTEDANCH{\char"2C6A\char"0309} \SUPPORTEDANCH{\char"006C\char"0309} \SUPPORTEDANCH{\char"0234\char"0309} \SUPPORTEDANCH{\char"006E\char"0309} \SUPPORTEDANCH{\char"0235\char"0309} \SUPPORTEDANCH{\char"0070\char"0309} \SUPPORTEDANCH{\char"01A5\char"0309} \SUPPORTEDANCH{\char"0071\char"0309} \SUPPORTEDANCH{\char"02A0\char"0309} \SUPPORTEDANCH{\char"0074\char"0309} \SUPPORTEDANCH{\char"01AB\char"0309} \SUPPORTEDANCH{\char"01AD\char"0309} \SUPPORTEDANCH{\char"0236\char"0309} \SUPPORTEDPREC{\char"0075\char"0309} \MISSINGGLYPH{\char"AB52\char"0309} \SUPPORTEDPREC{\char"0079\char"0309} \SUPPORTEDANCH{\char"01B4\char"0309} \MISSINGGLYPH{\char"1EFF\char"0309}

\SUPPORTEDANCH{\char"0048\char"030A} \SUPPORTEDANCH{\char"2C67\char"030A} \SUPPORTEDANCH{\char"004B\char"030A} \SUPPORTEDANCH{\char"2C69\char"030A} \SUPPORTEDANCH{\char"004D\char"030A} \MISSINGGLYPH{\char"2C6E\char"030A} \SUPPORTEDANCH{\char"0050\char"030A} \SUPPORTEDANCH{\char"01A4\char"030A} \SUPPORTEDANCH{\char"0052\char"030A} \SUPPORTEDANCH{\char"2C64\char"030A} \SUPPORTEDANCH{\char"0054\char"030A} \SUPPORTEDANCH{\char"01AC\char"030A} \SUPPORTEDANCH{\char"01AE\char"030A} \SUPPORTEDANCH{\char"0059\char"030A} \SUPPORTEDANCH{\char"01B3\char"030A} \SUPPORTEDPREC{\char"0061\char"030A} \SUPPORTEDANCH{\char"1E9A\char"030A} \SUPPORTEDANCH{\char"0062\char"030A} \SUPPORTEDANCH{\char"0183\char"030A} \SUPPORTEDANCH{\char"0063\char"030A} \SUPPORTEDANCH{\char"0188\char"030A} \SUPPORTEDANCH{\char"0064\char"030A} \SUPPORTEDANCH{\char"018C\char"030A} \SUPPORTEDANCH{\char"0221\char"030A} \SUPPORTEDANCH{\char"0068\char"030A} \SUPPORTEDANCH{\char"2C68\char"030A} \SUPPORTEDANCH{\char"006B\char"030A} \SUPPORTEDANCH{\char"0199\char"030A} \SUPPORTEDANCH{\char"2C6A\char"030A} \SUPPORTEDANCH{\char"006C\char"030A} \SUPPORTEDANCH{\char"0234\char"030A} \SUPPORTEDANCH{\char"006E\char"030A} \SUPPORTEDANCH{\char"0235\char"030A} \SUPPORTEDANCH{\char"0070\char"030A} \SUPPORTEDANCH{\char"01A5\char"030A} \SUPPORTEDANCH{\char"0071\char"030A} \SUPPORTEDANCH{\char"02A0\char"030A} \SUPPORTEDANCH{\char"0074\char"030A} \SUPPORTEDANCH{\char"01AB\char"030A} \SUPPORTEDANCH{\char"01AD\char"030A} \SUPPORTEDANCH{\char"0236\char"030A} \SUPPORTEDPREC{\char"0075\char"030A} \MISSINGGLYPH{\char"AB52\char"030A} \SUPPORTEDPREC{\char"0079\char"030A} \SUPPORTEDANCH{\char"01B4\char"030A} \MISSINGGLYPH{\char"1EFF\char"030A}

\SUPPORTEDANCH{\char"0048\char"030B} \SUPPORTEDANCH{\char"2C67\char"030B} \SUPPORTEDANCH{\char"004B\char"030B} \SUPPORTEDANCH{\char"2C69\char"030B} \SUPPORTEDANCH{\char"004D\char"030B} \MISSINGGLYPH{\char"2C6E\char"030B} \SUPPORTEDANCH{\char"0050\char"030B} \SUPPORTEDANCH{\char"01A4\char"030B} \SUPPORTEDANCH{\char"0052\char"030B} \SUPPORTEDANCH{\char"2C64\char"030B} \SUPPORTEDANCH{\char"0054\char"030B} \SUPPORTEDANCH{\char"01AC\char"030B} \SUPPORTEDANCH{\char"01AE\char"030B} \SUPPORTEDANCH{\char"0059\char"030B} \SUPPORTEDANCH{\char"01B3\char"030B} \SUPPORTEDANCH{\char"0061\char"030B} \SUPPORTEDANCH{\char"1E9A\char"030B} \SUPPORTEDANCH{\char"0062\char"030B} \SUPPORTEDANCH{\char"0183\char"030B} \SUPPORTEDANCH{\char"0063\char"030B} \SUPPORTEDANCH{\char"0188\char"030B} \SUPPORTEDANCH{\char"0064\char"030B} \SUPPORTEDANCH{\char"018C\char"030B} \SUPPORTEDANCH{\char"0221\char"030B} \SUPPORTEDANCH{\char"0068\char"030B} \SUPPORTEDANCH{\char"2C68\char"030B} \SUPPORTEDANCH{\char"006B\char"030B} \SUPPORTEDANCH{\char"0199\char"030B} \SUPPORTEDANCH{\char"2C6A\char"030B} \SUPPORTEDANCH{\char"006C\char"030B} \SUPPORTEDANCH{\char"0234\char"030B} \SUPPORTEDANCH{\char"006E\char"030B} \SUPPORTEDANCH{\char"0235\char"030B} \SUPPORTEDANCH{\char"0070\char"030B} \SUPPORTEDANCH{\char"01A5\char"030B} \SUPPORTEDANCH{\char"0071\char"030B} \SUPPORTEDANCH{\char"02A0\char"030B} \SUPPORTEDANCH{\char"0074\char"030B} \SUPPORTEDANCH{\char"01AB\char"030B} \SUPPORTEDANCH{\char"01AD\char"030B} \SUPPORTEDANCH{\char"0236\char"030B} \SUPPORTEDPREC{\char"0075\char"030B} \MISSINGGLYPH{\char"AB52\char"030B} \SUPPORTEDANCH{\char"0079\char"030B} \SUPPORTEDANCH{\char"01B4\char"030B} \MISSINGGLYPH{\char"1EFF\char"030B}

\SUPPORTEDPREC{\char"0048\char"030C} \SUPPORTEDANCH{\char"2C67\char"030C} \SUPPORTEDPREC{\char"004B\char"030C} \SUPPORTEDANCH{\char"2C69\char"030C} \SUPPORTEDANCH{\char"004D\char"030C} \MISSINGGLYPH{\char"2C6E\char"030C} \SUPPORTEDANCH{\char"0050\char"030C} \SUPPORTEDANCH{\char"01A4\char"030C} \SUPPORTEDPREC{\char"0052\char"030C} \SUPPORTEDANCH{\char"2C64\char"030C} \SUPPORTEDPREC{\char"0054\char"030C} \SUPPORTEDANCH{\char"01AC\char"030C} \SUPPORTEDANCH{\char"01AE\char"030C} \SUPPORTEDANCH{\char"0059\char"030C} \SUPPORTEDANCH{\char"01B3\char"030C} \SUPPORTEDPREC{\char"0061\char"030C} \SUPPORTEDANCH{\char"1E9A\char"030C} \SUPPORTEDANCH{\char"0062\char"030C} \SUPPORTEDANCH{\char"0183\char"030C} \SUPPORTEDPREC{\char"0063\char"030C} \SUPPORTEDANCH{\char"0188\char"030C} \SUPPORTEDPREC{\char"0064\char"030C} \SUPPORTEDANCH{\char"018C\char"030C} \SUPPORTEDANCH{\char"0221\char"030C} \SUPPORTEDPREC{\char"0068\char"030C} \SUPPORTEDANCH{\char"2C68\char"030C} \SUPPORTEDPREC{\char"006B\char"030C} \SUPPORTEDANCH{\char"0199\char"030C} \SUPPORTEDANCH{\char"2C6A\char"030C} \SUPPORTEDPREC{\char"006C\char"030C} \SUPPORTEDANCH{\char"0234\char"030C} \SUPPORTEDPREC{\char"006E\char"030C} \SUPPORTEDANCH{\char"0235\char"030C} \SUPPORTEDANCH{\char"0070\char"030C} \SUPPORTEDANCH{\char"01A5\char"030C} \SUPPORTEDANCH{\char"0071\char"030C} \SUPPORTEDANCH{\char"02A0\char"030C} \SUPPORTEDPREC{\char"0074\char"030C} \SUPPORTEDANCH{\char"01AB\char"030C} \SUPPORTEDANCH{\char"01AD\char"030C} \SUPPORTEDANCH{\char"0236\char"030C} \SUPPORTEDPREC{\char"0075\char"030C} \MISSINGGLYPH{\char"AB52\char"030C} \SUPPORTEDANCH{\char"0079\char"030C} \SUPPORTEDANCH{\char"01B4\char"030C} \MISSINGGLYPH{\char"1EFF\char"030C}

\SUPPORTEDANCH{\char"0048\char"030D} \SUPPORTEDANCH{\char"2C67\char"030D} \SUPPORTEDANCH{\char"004B\char"030D} \SUPPORTEDANCH{\char"2C69\char"030D} \SUPPORTEDANCH{\char"004D\char"030D} \MISSINGGLYPH{\char"2C6E\char"030D} \SUPPORTEDANCH{\char"0050\char"030D} \SUPPORTEDANCH{\char"01A4\char"030D} \SUPPORTEDANCH{\char"0052\char"030D} \SUPPORTEDANCH{\char"2C64\char"030D} \SUPPORTEDANCH{\char"0054\char"030D} \SUPPORTEDANCH{\char"01AC\char"030D} \SUPPORTEDANCH{\char"01AE\char"030D} \SUPPORTEDANCH{\char"0059\char"030D} \SUPPORTEDANCH{\char"01B3\char"030D} \SUPPORTEDANCH{\char"0061\char"030D} \SUPPORTEDANCH{\char"1E9A\char"030D} \SUPPORTEDANCH{\char"0062\char"030D} \SUPPORTEDANCH{\char"0183\char"030D} \SUPPORTEDANCH{\char"0063\char"030D} \SUPPORTEDANCH{\char"0188\char"030D} \SUPPORTEDANCH{\char"0064\char"030D} \SUPPORTEDANCH{\char"018C\char"030D} \SUPPORTEDANCH{\char"0221\char"030D} \SUPPORTEDANCH{\char"0068\char"030D} \SUPPORTEDANCH{\char"2C68\char"030D} \SUPPORTEDANCH{\char"006B\char"030D} \SUPPORTEDANCH{\char"0199\char"030D} \SUPPORTEDANCH{\char"2C6A\char"030D} \SUPPORTEDANCH{\char"006C\char"030D} \SUPPORTEDANCH{\char"0234\char"030D} \SUPPORTEDANCH{\char"006E\char"030D} \SUPPORTEDANCH{\char"0235\char"030D} \SUPPORTEDANCH{\char"0070\char"030D} \SUPPORTEDANCH{\char"01A5\char"030D} \SUPPORTEDANCH{\char"0071\char"030D} \SUPPORTEDANCH{\char"02A0\char"030D} \SUPPORTEDANCH{\char"0074\char"030D} \SUPPORTEDANCH{\char"01AB\char"030D} \SUPPORTEDANCH{\char"01AD\char"030D} \SUPPORTEDANCH{\char"0236\char"030D} \SUPPORTEDANCH{\char"0075\char"030D} \MISSINGGLYPH{\char"AB52\char"030D} \SUPPORTEDANCH{\char"0079\char"030D} \SUPPORTEDANCH{\char"01B4\char"030D} \MISSINGGLYPH{\char"1EFF\char"030D}

\SUPPORTEDANCH{\char"0048\char"030F} \SUPPORTEDANCH{\char"2C67\char"030F} \SUPPORTEDANCH{\char"004B\char"030F} \SUPPORTEDANCH{\char"2C69\char"030F} \SUPPORTEDANCH{\char"004D\char"030F} \MISSINGGLYPH{\char"2C6E\char"030F} \SUPPORTEDANCH{\char"0050\char"030F} \SUPPORTEDANCH{\char"01A4\char"030F} \SUPPORTEDPREC{\char"0052\char"030F} \SUPPORTEDANCH{\char"2C64\char"030F} \SUPPORTEDANCH{\char"0054\char"030F} \SUPPORTEDANCH{\char"01AC\char"030F} \SUPPORTEDANCH{\char"01AE\char"030F} \SUPPORTEDANCH{\char"0059\char"030F} \SUPPORTEDANCH{\char"01B3\char"030F} \SUPPORTEDPREC{\char"0061\char"030F} \SUPPORTEDANCH{\char"1E9A\char"030F} \SUPPORTEDANCH{\char"0062\char"030F} \SUPPORTEDANCH{\char"0183\char"030F} \SUPPORTEDANCH{\char"0063\char"030F} \SUPPORTEDANCH{\char"0188\char"030F} \SUPPORTEDANCH{\char"0064\char"030F} \SUPPORTEDANCH{\char"018C\char"030F} \SUPPORTEDANCH{\char"0221\char"030F} \SUPPORTEDANCH{\char"0068\char"030F} \SUPPORTEDANCH{\char"2C68\char"030F} \SUPPORTEDANCH{\char"006B\char"030F} \SUPPORTEDANCH{\char"0199\char"030F} \SUPPORTEDANCH{\char"2C6A\char"030F} \SUPPORTEDANCH{\char"006C\char"030F} \SUPPORTEDANCH{\char"0234\char"030F} \SUPPORTEDANCH{\char"006E\char"030F} \SUPPORTEDANCH{\char"0235\char"030F} \SUPPORTEDANCH{\char"0070\char"030F} \SUPPORTEDANCH{\char"01A5\char"030F} \SUPPORTEDANCH{\char"0071\char"030F} \SUPPORTEDANCH{\char"02A0\char"030F} \SUPPORTEDANCH{\char"0074\char"030F} \SUPPORTEDANCH{\char"01AB\char"030F} \SUPPORTEDANCH{\char"01AD\char"030F} \SUPPORTEDANCH{\char"0236\char"030F} \SUPPORTEDPREC{\char"0075\char"030F} \MISSINGGLYPH{\char"AB52\char"030F} \SUPPORTEDANCH{\char"0079\char"030F} \SUPPORTEDANCH{\char"01B4\char"030F} \MISSINGGLYPH{\char"1EFF\char"030F}

\SUPPORTEDANCH{\char"0048\char"0310} \SUPPORTEDANCH{\char"2C67\char"0310} \SUPPORTEDANCH{\char"004B\char"0310} \SUPPORTEDANCH{\char"2C69\char"0310} \SUPPORTEDANCH{\char"004D\char"0310} \MISSINGGLYPH{\char"2C6E\char"0310} \SUPPORTEDANCH{\char"0050\char"0310} \SUPPORTEDANCH{\char"01A4\char"0310} \SUPPORTEDANCH{\char"0052\char"0310} \SUPPORTEDANCH{\char"2C64\char"0310} \SUPPORTEDANCH{\char"0054\char"0310} \SUPPORTEDANCH{\char"01AC\char"0310} \SUPPORTEDANCH{\char"01AE\char"0310} \SUPPORTEDANCH{\char"0059\char"0310} \SUPPORTEDANCH{\char"01B3\char"0310} \SUPPORTEDANCH{\char"0061\char"0310} \SUPPORTEDANCH{\char"1E9A\char"0310} \SUPPORTEDANCH{\char"0062\char"0310} \SUPPORTEDANCH{\char"0183\char"0310} \SUPPORTEDANCH{\char"0063\char"0310} \SUPPORTEDANCH{\char"0188\char"0310} \SUPPORTEDANCH{\char"0064\char"0310} \SUPPORTEDANCH{\char"018C\char"0310} \SUPPORTEDANCH{\char"0221\char"0310} \SUPPORTEDANCH{\char"0068\char"0310} \SUPPORTEDANCH{\char"2C68\char"0310} \SUPPORTEDANCH{\char"006B\char"0310} \SUPPORTEDANCH{\char"0199\char"0310} \SUPPORTEDANCH{\char"2C6A\char"0310} \SUPPORTEDANCH{\char"006C\char"0310} \SUPPORTEDANCH{\char"0234\char"0310} \SUPPORTEDANCH{\char"006E\char"0310} \SUPPORTEDANCH{\char"0235\char"0310} \SUPPORTEDANCH{\char"0070\char"0310} \SUPPORTEDANCH{\char"01A5\char"0310} \SUPPORTEDANCH{\char"0071\char"0310} \SUPPORTEDANCH{\char"02A0\char"0310} \SUPPORTEDANCH{\char"0074\char"0310} \SUPPORTEDANCH{\char"01AB\char"0310} \SUPPORTEDANCH{\char"01AD\char"0310} \SUPPORTEDANCH{\char"0236\char"0310} \SUPPORTEDANCH{\char"0075\char"0310} \MISSINGGLYPH{\char"AB52\char"0310} \SUPPORTEDANCH{\char"0079\char"0310} \SUPPORTEDANCH{\char"01B4\char"0310} \MISSINGGLYPH{\char"1EFF\char"0310}

\SUPPORTEDANCH{\char"0048\char"0311} \SUPPORTEDANCH{\char"2C67\char"0311} \SUPPORTEDANCH{\char"004B\char"0311} \SUPPORTEDANCH{\char"2C69\char"0311} \SUPPORTEDANCH{\char"004D\char"0311} \MISSINGGLYPH{\char"2C6E\char"0311} \SUPPORTEDANCH{\char"0050\char"0311} \SUPPORTEDANCH{\char"01A4\char"0311} \SUPPORTEDPREC{\char"0052\char"0311} \SUPPORTEDANCH{\char"2C64\char"0311} \SUPPORTEDANCH{\char"0054\char"0311} \SUPPORTEDANCH{\char"01AC\char"0311} \SUPPORTEDANCH{\char"01AE\char"0311} \SUPPORTEDANCH{\char"0059\char"0311} \SUPPORTEDANCH{\char"01B3\char"0311} \SUPPORTEDPREC{\char"0061\char"0311} \SUPPORTEDANCH{\char"1E9A\char"0311} \SUPPORTEDANCH{\char"0062\char"0311} \SUPPORTEDANCH{\char"0183\char"0311} \SUPPORTEDANCH{\char"0063\char"0311} \SUPPORTEDANCH{\char"0188\char"0311} \SUPPORTEDANCH{\char"0064\char"0311} \SUPPORTEDANCH{\char"018C\char"0311} \SUPPORTEDANCH{\char"0221\char"0311} \SUPPORTEDANCH{\char"0068\char"0311} \SUPPORTEDANCH{\char"2C68\char"0311} \SUPPORTEDANCH{\char"006B\char"0311} \SUPPORTEDANCH{\char"0199\char"0311} \SUPPORTEDANCH{\char"2C6A\char"0311} \SUPPORTEDANCH{\char"006C\char"0311} \SUPPORTEDANCH{\char"0234\char"0311} \SUPPORTEDANCH{\char"006E\char"0311} \SUPPORTEDANCH{\char"0235\char"0311} \SUPPORTEDANCH{\char"0070\char"0311} \SUPPORTEDANCH{\char"01A5\char"0311} \SUPPORTEDANCH{\char"0071\char"0311} \SUPPORTEDANCH{\char"02A0\char"0311} \SUPPORTEDANCH{\char"0074\char"0311} \SUPPORTEDANCH{\char"01AB\char"0311} \SUPPORTEDANCH{\char"01AD\char"0311} \SUPPORTEDANCH{\char"0236\char"0311} \SUPPORTEDPREC{\char"0075\char"0311} \MISSINGGLYPH{\char"AB52\char"0311} \SUPPORTEDANCH{\char"0079\char"0311} \SUPPORTEDANCH{\char"01B4\char"0311} \MISSINGGLYPH{\char"1EFF\char"0311}

\SUPPORTEDANCH{\char"0048\char"0312} \SUPPORTEDANCH{\char"2C67\char"0312} \SUPPORTEDANCH{\char"004B\char"0312} \SUPPORTEDANCH{\char"2C69\char"0312} \SUPPORTEDANCH{\char"004D\char"0312} \MISSINGGLYPH{\char"2C6E\char"0312} \SUPPORTEDANCH{\char"0050\char"0312} \SUPPORTEDANCH{\char"01A4\char"0312} \SUPPORTEDANCH{\char"0052\char"0312} \SUPPORTEDANCH{\char"2C64\char"0312} \SUPPORTEDANCH{\char"0054\char"0312} \SUPPORTEDANCH{\char"01AC\char"0312} \SUPPORTEDANCH{\char"01AE\char"0312} \SUPPORTEDANCH{\char"0059\char"0312} \SUPPORTEDANCH{\char"01B3\char"0312} \SUPPORTEDANCH{\char"0061\char"0312} \SUPPORTEDANCH{\char"1E9A\char"0312} \SUPPORTEDANCH{\char"0062\char"0312} \SUPPORTEDANCH{\char"0183\char"0312} \SUPPORTEDANCH{\char"0063\char"0312} \SUPPORTEDANCH{\char"0188\char"0312} \SUPPORTEDANCH{\char"0064\char"0312} \SUPPORTEDANCH{\char"018C\char"0312} \SUPPORTEDANCH{\char"0221\char"0312} \SUPPORTEDANCH{\char"0068\char"0312} \SUPPORTEDANCH{\char"2C68\char"0312} \SUPPORTEDANCH{\char"006B\char"0312} \SUPPORTEDANCH{\char"0199\char"0312} \SUPPORTEDANCH{\char"2C6A\char"0312} \SUPPORTEDANCH{\char"006C\char"0312} \SUPPORTEDANCH{\char"0234\char"0312} \SUPPORTEDANCH{\char"006E\char"0312} \SUPPORTEDANCH{\char"0235\char"0312} \SUPPORTEDANCH{\char"0070\char"0312} \SUPPORTEDANCH{\char"01A5\char"0312} \SUPPORTEDANCH{\char"0071\char"0312} \SUPPORTEDANCH{\char"02A0\char"0312} \SUPPORTEDANCH{\char"0074\char"0312} \SUPPORTEDANCH{\char"01AB\char"0312} \SUPPORTEDANCH{\char"01AD\char"0312} \SUPPORTEDANCH{\char"0236\char"0312} \SUPPORTEDANCH{\char"0075\char"0312} \MISSINGGLYPH{\char"AB52\char"0312} \SUPPORTEDANCH{\char"0079\char"0312} \SUPPORTEDANCH{\char"01B4\char"0312} \MISSINGGLYPH{\char"1EFF\char"0312}

\SUPPORTEDANCH{\char"0048\char"033D} \SUPPORTEDANCH{\char"2C67\char"033D} \SUPPORTEDANCH{\char"004B\char"033D} \SUPPORTEDANCH{\char"2C69\char"033D} \SUPPORTEDANCH{\char"004D\char"033D} \MISSINGGLYPH{\char"2C6E\char"033D} \SUPPORTEDANCH{\char"0050\char"033D} \SUPPORTEDANCH{\char"01A4\char"033D} \SUPPORTEDANCH{\char"0052\char"033D} \SUPPORTEDANCH{\char"2C64\char"033D} \SUPPORTEDANCH{\char"0054\char"033D} \SUPPORTEDANCH{\char"01AC\char"033D} \SUPPORTEDANCH{\char"01AE\char"033D} \SUPPORTEDANCH{\char"0059\char"033D} \SUPPORTEDANCH{\char"01B3\char"033D} \SUPPORTEDANCH{\char"0061\char"033D} \SUPPORTEDANCH{\char"1E9A\char"033D} \SUPPORTEDANCH{\char"0062\char"033D} \SUPPORTEDANCH{\char"0183\char"033D} \SUPPORTEDANCH{\char"0063\char"033D} \SUPPORTEDANCH{\char"0188\char"033D} \SUPPORTEDANCH{\char"0064\char"033D} \SUPPORTEDANCH{\char"018C\char"033D} \SUPPORTEDANCH{\char"0221\char"033D} \SUPPORTEDANCH{\char"0068\char"033D} \SUPPORTEDANCH{\char"2C68\char"033D} \SUPPORTEDANCH{\char"006B\char"033D} \SUPPORTEDANCH{\char"0199\char"033D} \SUPPORTEDANCH{\char"2C6A\char"033D} \SUPPORTEDANCH{\char"006C\char"033D} \SUPPORTEDANCH{\char"0234\char"033D} \SUPPORTEDANCH{\char"006E\char"033D} \SUPPORTEDANCH{\char"0235\char"033D} \SUPPORTEDANCH{\char"0070\char"033D} \SUPPORTEDANCH{\char"01A5\char"033D} \SUPPORTEDANCH{\char"0071\char"033D} \SUPPORTEDANCH{\char"02A0\char"033D} \SUPPORTEDANCH{\char"0074\char"033D} \SUPPORTEDANCH{\char"01AB\char"033D} \SUPPORTEDANCH{\char"01AD\char"033D} \SUPPORTEDANCH{\char"0236\char"033D} \SUPPORTEDANCH{\char"0075\char"033D} \MISSINGGLYPH{\char"AB52\char"033D} \SUPPORTEDANCH{\char"0079\char"033D} \SUPPORTEDANCH{\char"01B4\char"033D} \MISSINGGLYPH{\char"1EFF\char"033D}

\SUPPORTEDANCH{\char"0048\char"033E} \SUPPORTEDANCH{\char"2C67\char"033E} \SUPPORTEDANCH{\char"004B\char"033E} \SUPPORTEDANCH{\char"2C69\char"033E} \SUPPORTEDANCH{\char"004D\char"033E} \MISSINGGLYPH{\char"2C6E\char"033E} \SUPPORTEDANCH{\char"0050\char"033E} \SUPPORTEDANCH{\char"01A4\char"033E} \SUPPORTEDANCH{\char"0052\char"033E} \SUPPORTEDANCH{\char"2C64\char"033E} \SUPPORTEDANCH{\char"0054\char"033E} \SUPPORTEDANCH{\char"01AC\char"033E} \SUPPORTEDANCH{\char"01AE\char"033E} \SUPPORTEDANCH{\char"0059\char"033E} \SUPPORTEDANCH{\char"01B3\char"033E} \SUPPORTEDANCH{\char"0061\char"033E} \SUPPORTEDANCH{\char"1E9A\char"033E} \SUPPORTEDANCH{\char"0062\char"033E} \SUPPORTEDANCH{\char"0183\char"033E} \SUPPORTEDANCH{\char"0063\char"033E} \SUPPORTEDANCH{\char"0188\char"033E} \SUPPORTEDANCH{\char"0064\char"033E} \SUPPORTEDANCH{\char"018C\char"033E} \SUPPORTEDANCH{\char"0221\char"033E} \SUPPORTEDANCH{\char"0068\char"033E} \SUPPORTEDANCH{\char"2C68\char"033E} \SUPPORTEDANCH{\char"006B\char"033E} \SUPPORTEDANCH{\char"0199\char"033E} \SUPPORTEDANCH{\char"2C6A\char"033E} \SUPPORTEDANCH{\char"006C\char"033E} \SUPPORTEDANCH{\char"0234\char"033E} \SUPPORTEDANCH{\char"006E\char"033E} \SUPPORTEDANCH{\char"0235\char"033E} \SUPPORTEDANCH{\char"0070\char"033E} \SUPPORTEDANCH{\char"01A5\char"033E} \SUPPORTEDANCH{\char"0071\char"033E} \SUPPORTEDANCH{\char"02A0\char"033E} \SUPPORTEDANCH{\char"0074\char"033E} \SUPPORTEDANCH{\char"01AB\char"033E} \SUPPORTEDANCH{\char"01AD\char"033E} \SUPPORTEDANCH{\char"0236\char"033E} \SUPPORTEDANCH{\char"0075\char"033E} \MISSINGGLYPH{\char"AB52\char"033E} \SUPPORTEDANCH{\char"0079\char"033E} \SUPPORTEDANCH{\char"01B4\char"033E} \MISSINGGLYPH{\char"1EFF\char"033E}

\SUPPORTEDANCH{\char"0048\char"033F} \SUPPORTEDANCH{\char"2C67\char"033F} \SUPPORTEDANCH{\char"004B\char"033F} \SUPPORTEDANCH{\char"2C69\char"033F} \SUPPORTEDANCH{\char"004D\char"033F} \MISSINGGLYPH{\char"2C6E\char"033F} \SUPPORTEDANCH{\char"0050\char"033F} \SUPPORTEDANCH{\char"01A4\char"033F} \SUPPORTEDANCH{\char"0052\char"033F} \SUPPORTEDANCH{\char"2C64\char"033F} \SUPPORTEDANCH{\char"0054\char"033F} \SUPPORTEDANCH{\char"01AC\char"033F} \SUPPORTEDANCH{\char"01AE\char"033F} \SUPPORTEDANCH{\char"0059\char"033F} \SUPPORTEDANCH{\char"01B3\char"033F} \SUPPORTEDANCH{\char"0061\char"033F} \SUPPORTEDANCH{\char"1E9A\char"033F} \SUPPORTEDANCH{\char"0062\char"033F} \SUPPORTEDANCH{\char"0183\char"033F} \SUPPORTEDANCH{\char"0063\char"033F} \SUPPORTEDANCH{\char"0188\char"033F} \SUPPORTEDANCH{\char"0064\char"033F} \SUPPORTEDANCH{\char"018C\char"033F} \SUPPORTEDANCH{\char"0221\char"033F} \SUPPORTEDANCH{\char"0068\char"033F} \SUPPORTEDANCH{\char"2C68\char"033F} \SUPPORTEDANCH{\char"006B\char"033F} \SUPPORTEDANCH{\char"0199\char"033F} \SUPPORTEDANCH{\char"2C6A\char"033F} \SUPPORTEDANCH{\char"006C\char"033F} \SUPPORTEDANCH{\char"0234\char"033F} \SUPPORTEDANCH{\char"006E\char"033F} \SUPPORTEDANCH{\char"0235\char"033F} \SUPPORTEDANCH{\char"0070\char"033F} \SUPPORTEDANCH{\char"01A5\char"033F} \SUPPORTEDANCH{\char"0071\char"033F} \SUPPORTEDANCH{\char"02A0\char"033F} \SUPPORTEDANCH{\char"0074\char"033F} \SUPPORTEDANCH{\char"01AB\char"033F} \SUPPORTEDANCH{\char"01AD\char"033F} \SUPPORTEDANCH{\char"0236\char"033F} \SUPPORTEDANCH{\char"0075\char"033F} \MISSINGGLYPH{\char"AB52\char"033F} \SUPPORTEDANCH{\char"0079\char"033F} \SUPPORTEDANCH{\char"01B4\char"033F} \MISSINGGLYPH{\char"1EFF\char"033F}

\SUPPORTEDANCH{\char"0048\char"030E} \SUPPORTEDANCH{\char"2C67\char"030E} \SUPPORTEDANCH{\char"004B\char"030E} \SUPPORTEDANCH{\char"2C69\char"030E} \SUPPORTEDANCH{\char"004D\char"030E} \MISSINGGLYPH{\char"2C6E\char"030E} \SUPPORTEDANCH{\char"0050\char"030E} \SUPPORTEDANCH{\char"01A4\char"030E} \SUPPORTEDANCH{\char"0052\char"030E} \SUPPORTEDANCH{\char"2C64\char"030E} \SUPPORTEDANCH{\char"0054\char"030E} \SUPPORTEDANCH{\char"01AC\char"030E} \SUPPORTEDANCH{\char"01AE\char"030E} \SUPPORTEDANCH{\char"0059\char"030E} \SUPPORTEDANCH{\char"01B3\char"030E} \SUPPORTEDANCH{\char"0061\char"030E} \SUPPORTEDANCH{\char"1E9A\char"030E} \SUPPORTEDANCH{\char"0062\char"030E} \SUPPORTEDANCH{\char"0183\char"030E} \SUPPORTEDANCH{\char"0063\char"030E} \SUPPORTEDANCH{\char"0188\char"030E} \SUPPORTEDANCH{\char"0064\char"030E} \SUPPORTEDANCH{\char"018C\char"030E} \SUPPORTEDANCH{\char"0221\char"030E} \SUPPORTEDANCH{\char"0068\char"030E} \SUPPORTEDANCH{\char"2C68\char"030E} \SUPPORTEDANCH{\char"006B\char"030E} \SUPPORTEDANCH{\char"0199\char"030E} \SUPPORTEDANCH{\char"2C6A\char"030E} \SUPPORTEDANCH{\char"006C\char"030E} \SUPPORTEDANCH{\char"0234\char"030E} \SUPPORTEDANCH{\char"006E\char"030E} \SUPPORTEDANCH{\char"0235\char"030E} \SUPPORTEDANCH{\char"0070\char"030E} \SUPPORTEDANCH{\char"01A5\char"030E} \SUPPORTEDANCH{\char"0071\char"030E} \SUPPORTEDANCH{\char"02A0\char"030E} \SUPPORTEDANCH{\char"0074\char"030E} \SUPPORTEDANCH{\char"01AB\char"030E} \SUPPORTEDANCH{\char"01AD\char"030E} \SUPPORTEDANCH{\char"0236\char"030E} \SUPPORTEDANCH{\char"0075\char"030E} \MISSINGGLYPH{\char"AB52\char"030E} \SUPPORTEDANCH{\char"0079\char"030E} \SUPPORTEDANCH{\char"01B4\char"030E} \MISSINGGLYPH{\char"1EFF\char"030E}

\SUPPORTEDANCH{\char"0048\char"0346} \SUPPORTEDANCH{\char"2C67\char"0346} \SUPPORTEDANCH{\char"004B\char"0346} \SUPPORTEDANCH{\char"2C69\char"0346} \SUPPORTEDANCH{\char"004D\char"0346} \MISSINGGLYPH{\char"2C6E\char"0346} \SUPPORTEDANCH{\char"0050\char"0346} \SUPPORTEDANCH{\char"01A4\char"0346} \SUPPORTEDANCH{\char"0052\char"0346} \SUPPORTEDANCH{\char"2C64\char"0346} \SUPPORTEDANCH{\char"0054\char"0346} \SUPPORTEDANCH{\char"01AC\char"0346} \SUPPORTEDANCH{\char"01AE\char"0346} \SUPPORTEDANCH{\char"0059\char"0346} \SUPPORTEDANCH{\char"01B3\char"0346} \SUPPORTEDANCH{\char"0061\char"0346} \SUPPORTEDANCH{\char"1E9A\char"0346} \SUPPORTEDANCH{\char"0062\char"0346} \SUPPORTEDANCH{\char"0183\char"0346} \SUPPORTEDANCH{\char"0063\char"0346} \SUPPORTEDANCH{\char"0188\char"0346} \SUPPORTEDANCH{\char"0064\char"0346} \SUPPORTEDANCH{\char"018C\char"0346} \SUPPORTEDANCH{\char"0221\char"0346} \SUPPORTEDANCH{\char"0068\char"0346} \SUPPORTEDANCH{\char"2C68\char"0346} \SUPPORTEDANCH{\char"006B\char"0346} \SUPPORTEDANCH{\char"0199\char"0346} \SUPPORTEDANCH{\char"2C6A\char"0346} \SUPPORTEDANCH{\char"006C\char"0346} \SUPPORTEDANCH{\char"0234\char"0346} \SUPPORTEDANCH{\char"006E\char"0346} \SUPPORTEDANCH{\char"0235\char"0346} \SUPPORTEDANCH{\char"0070\char"0346} \SUPPORTEDANCH{\char"01A5\char"0346} \SUPPORTEDANCH{\char"0071\char"0346} \SUPPORTEDANCH{\char"02A0\char"0346} \SUPPORTEDANCH{\char"0074\char"0346} \SUPPORTEDANCH{\char"01AB\char"0346} \SUPPORTEDANCH{\char"01AD\char"0346} \SUPPORTEDANCH{\char"0236\char"0346} \SUPPORTEDANCH{\char"0075\char"0346} \MISSINGGLYPH{\char"AB52\char"0346} \SUPPORTEDANCH{\char"0079\char"0346} \SUPPORTEDANCH{\char"01B4\char"0346} \MISSINGGLYPH{\char"1EFF\char"0346}

\SUPPORTEDANCH{\char"0048\char"034A} \SUPPORTEDANCH{\char"2C67\char"034A} \SUPPORTEDANCH{\char"004B\char"034A} \SUPPORTEDANCH{\char"2C69\char"034A} \SUPPORTEDANCH{\char"004D\char"034A} \MISSINGGLYPH{\char"2C6E\char"034A} \SUPPORTEDANCH{\char"0050\char"034A} \SUPPORTEDANCH{\char"01A4\char"034A} \SUPPORTEDANCH{\char"0052\char"034A} \SUPPORTEDANCH{\char"2C64\char"034A} \SUPPORTEDANCH{\char"0054\char"034A} \SUPPORTEDANCH{\char"01AC\char"034A} \SUPPORTEDANCH{\char"01AE\char"034A} \SUPPORTEDANCH{\char"0059\char"034A} \SUPPORTEDANCH{\char"01B3\char"034A} \SUPPORTEDANCH{\char"0061\char"034A} \SUPPORTEDANCH{\char"1E9A\char"034A} \SUPPORTEDANCH{\char"0062\char"034A} \SUPPORTEDANCH{\char"0183\char"034A} \SUPPORTEDANCH{\char"0063\char"034A} \SUPPORTEDANCH{\char"0188\char"034A} \SUPPORTEDANCH{\char"0064\char"034A} \SUPPORTEDANCH{\char"018C\char"034A} \SUPPORTEDANCH{\char"0221\char"034A} \SUPPORTEDANCH{\char"0068\char"034A} \SUPPORTEDANCH{\char"2C68\char"034A} \SUPPORTEDANCH{\char"006B\char"034A} \SUPPORTEDANCH{\char"0199\char"034A} \SUPPORTEDANCH{\char"2C6A\char"034A} \SUPPORTEDANCH{\char"006C\char"034A} \SUPPORTEDANCH{\char"0234\char"034A} \SUPPORTEDANCH{\char"006E\char"034A} \SUPPORTEDANCH{\char"0235\char"034A} \SUPPORTEDANCH{\char"0070\char"034A} \SUPPORTEDANCH{\char"01A5\char"034A} \SUPPORTEDANCH{\char"0071\char"034A} \SUPPORTEDANCH{\char"02A0\char"034A} \SUPPORTEDANCH{\char"0074\char"034A} \SUPPORTEDANCH{\char"01AB\char"034A} \SUPPORTEDANCH{\char"01AD\char"034A} \SUPPORTEDANCH{\char"0236\char"034A} \SUPPORTEDANCH{\char"0075\char"034A} \MISSINGGLYPH{\char"AB52\char"034A} \SUPPORTEDANCH{\char"0079\char"034A} \SUPPORTEDANCH{\char"01B4\char"034A} \MISSINGGLYPH{\char"1EFF\char"034A}

\SUPPORTEDANCH{\char"0048\char"034B} \SUPPORTEDANCH{\char"2C67\char"034B} \SUPPORTEDANCH{\char"004B\char"034B} \SUPPORTEDANCH{\char"2C69\char"034B} \SUPPORTEDANCH{\char"004D\char"034B} \MISSINGGLYPH{\char"2C6E\char"034B} \SUPPORTEDANCH{\char"0050\char"034B} \SUPPORTEDANCH{\char"01A4\char"034B} \SUPPORTEDANCH{\char"0052\char"034B} \SUPPORTEDANCH{\char"2C64\char"034B} \SUPPORTEDANCH{\char"0054\char"034B} \SUPPORTEDANCH{\char"01AC\char"034B} \SUPPORTEDANCH{\char"01AE\char"034B} \SUPPORTEDANCH{\char"0059\char"034B} \SUPPORTEDANCH{\char"01B3\char"034B} \SUPPORTEDANCH{\char"0061\char"034B} \SUPPORTEDANCH{\char"1E9A\char"034B} \SUPPORTEDANCH{\char"0062\char"034B} \SUPPORTEDANCH{\char"0183\char"034B} \SUPPORTEDANCH{\char"0063\char"034B} \SUPPORTEDANCH{\char"0188\char"034B} \SUPPORTEDANCH{\char"0064\char"034B} \SUPPORTEDANCH{\char"018C\char"034B} \SUPPORTEDANCH{\char"0221\char"034B} \SUPPORTEDANCH{\char"0068\char"034B} \SUPPORTEDANCH{\char"2C68\char"034B} \SUPPORTEDANCH{\char"006B\char"034B} \SUPPORTEDANCH{\char"0199\char"034B} \SUPPORTEDANCH{\char"2C6A\char"034B} \SUPPORTEDANCH{\char"006C\char"034B} \SUPPORTEDANCH{\char"0234\char"034B} \SUPPORTEDANCH{\char"006E\char"034B} \SUPPORTEDANCH{\char"0235\char"034B} \SUPPORTEDANCH{\char"0070\char"034B} \SUPPORTEDANCH{\char"01A5\char"034B} \SUPPORTEDANCH{\char"0071\char"034B} \SUPPORTEDANCH{\char"02A0\char"034B} \SUPPORTEDANCH{\char"0074\char"034B} \SUPPORTEDANCH{\char"01AB\char"034B} \SUPPORTEDANCH{\char"01AD\char"034B} \SUPPORTEDANCH{\char"0236\char"034B} \SUPPORTEDANCH{\char"0075\char"034B} \MISSINGGLYPH{\char"AB52\char"034B} \SUPPORTEDANCH{\char"0079\char"034B} \SUPPORTEDANCH{\char"01B4\char"034B} \MISSINGGLYPH{\char"1EFF\char"034B}

\SUPPORTEDANCH{\char"0048\char"034C} \SUPPORTEDANCH{\char"2C67\char"034C} \SUPPORTEDANCH{\char"004B\char"034C} \SUPPORTEDANCH{\char"2C69\char"034C} \SUPPORTEDANCH{\char"004D\char"034C} \MISSINGGLYPH{\char"2C6E\char"034C} \SUPPORTEDANCH{\char"0050\char"034C} \SUPPORTEDANCH{\char"01A4\char"034C} \SUPPORTEDANCH{\char"0052\char"034C} \SUPPORTEDANCH{\char"2C64\char"034C} \SUPPORTEDANCH{\char"0054\char"034C} \SUPPORTEDANCH{\char"01AC\char"034C} \SUPPORTEDANCH{\char"01AE\char"034C} \SUPPORTEDANCH{\char"0059\char"034C} \SUPPORTEDANCH{\char"01B3\char"034C} \SUPPORTEDANCH{\char"0061\char"034C} \SUPPORTEDANCH{\char"1E9A\char"034C} \SUPPORTEDANCH{\char"0062\char"034C} \SUPPORTEDANCH{\char"0183\char"034C} \SUPPORTEDANCH{\char"0063\char"034C} \SUPPORTEDANCH{\char"0188\char"034C} \SUPPORTEDANCH{\char"0064\char"034C} \SUPPORTEDANCH{\char"018C\char"034C} \SUPPORTEDANCH{\char"0221\char"034C} \SUPPORTEDANCH{\char"0068\char"034C} \SUPPORTEDANCH{\char"2C68\char"034C} \SUPPORTEDANCH{\char"006B\char"034C} \SUPPORTEDANCH{\char"0199\char"034C} \SUPPORTEDANCH{\char"2C6A\char"034C} \SUPPORTEDANCH{\char"006C\char"034C} \SUPPORTEDANCH{\char"0234\char"034C} \SUPPORTEDANCH{\char"006E\char"034C} \SUPPORTEDANCH{\char"0235\char"034C} \SUPPORTEDANCH{\char"0070\char"034C} \SUPPORTEDANCH{\char"01A5\char"034C} \SUPPORTEDANCH{\char"0071\char"034C} \SUPPORTEDANCH{\char"02A0\char"034C} \SUPPORTEDANCH{\char"0074\char"034C} \SUPPORTEDANCH{\char"01AB\char"034C} \SUPPORTEDANCH{\char"01AD\char"034C} \SUPPORTEDANCH{\char"0236\char"034C} \SUPPORTEDANCH{\char"0075\char"034C} \MISSINGGLYPH{\char"AB52\char"034C} \SUPPORTEDANCH{\char"0079\char"034C} \SUPPORTEDANCH{\char"01B4\char"034C} \MISSINGGLYPH{\char"1EFF\char"034C}

\SUPPORTEDANCH{\char"0048\char"0350} \SUPPORTEDANCH{\char"2C67\char"0350} \SUPPORTEDANCH{\char"004B\char"0350} \SUPPORTEDANCH{\char"2C69\char"0350} \SUPPORTEDANCH{\char"004D\char"0350} \MISSINGGLYPH{\char"2C6E\char"0350} \SUPPORTEDANCH{\char"0050\char"0350} \SUPPORTEDANCH{\char"01A4\char"0350} \SUPPORTEDANCH{\char"0052\char"0350} \SUPPORTEDANCH{\char"2C64\char"0350} \SUPPORTEDANCH{\char"0054\char"0350} \SUPPORTEDANCH{\char"01AC\char"0350} \SUPPORTEDANCH{\char"01AE\char"0350} \SUPPORTEDANCH{\char"0059\char"0350} \SUPPORTEDANCH{\char"01B3\char"0350} \SUPPORTEDANCH{\char"0061\char"0350} \SUPPORTEDANCH{\char"1E9A\char"0350} \SUPPORTEDANCH{\char"0062\char"0350} \SUPPORTEDANCH{\char"0183\char"0350} \SUPPORTEDANCH{\char"0063\char"0350} \SUPPORTEDANCH{\char"0188\char"0350} \SUPPORTEDANCH{\char"0064\char"0350} \SUPPORTEDANCH{\char"018C\char"0350} \SUPPORTEDANCH{\char"0221\char"0350} \SUPPORTEDANCH{\char"0068\char"0350} \SUPPORTEDANCH{\char"2C68\char"0350} \SUPPORTEDANCH{\char"006B\char"0350} \SUPPORTEDANCH{\char"0199\char"0350} \SUPPORTEDANCH{\char"2C6A\char"0350} \SUPPORTEDANCH{\char"006C\char"0350} \SUPPORTEDANCH{\char"0234\char"0350} \SUPPORTEDANCH{\char"006E\char"0350} \SUPPORTEDANCH{\char"0235\char"0350} \SUPPORTEDANCH{\char"0070\char"0350} \SUPPORTEDANCH{\char"01A5\char"0350} \SUPPORTEDANCH{\char"0071\char"0350} \SUPPORTEDANCH{\char"02A0\char"0350} \SUPPORTEDANCH{\char"0074\char"0350} \SUPPORTEDANCH{\char"01AB\char"0350} \SUPPORTEDANCH{\char"01AD\char"0350} \SUPPORTEDANCH{\char"0236\char"0350} \SUPPORTEDANCH{\char"0075\char"0350} \MISSINGGLYPH{\char"AB52\char"0350} \SUPPORTEDANCH{\char"0079\char"0350} \SUPPORTEDANCH{\char"01B4\char"0350} \MISSINGGLYPH{\char"1EFF\char"0350}

\SUPPORTEDANCH{\char"0048\char"0351} \SUPPORTEDANCH{\char"2C67\char"0351} \SUPPORTEDANCH{\char"004B\char"0351} \SUPPORTEDANCH{\char"2C69\char"0351} \SUPPORTEDANCH{\char"004D\char"0351} \MISSINGGLYPH{\char"2C6E\char"0351} \SUPPORTEDANCH{\char"0050\char"0351} \SUPPORTEDANCH{\char"01A4\char"0351} \SUPPORTEDANCH{\char"0052\char"0351} \SUPPORTEDANCH{\char"2C64\char"0351} \SUPPORTEDANCH{\char"0054\char"0351} \SUPPORTEDANCH{\char"01AC\char"0351} \SUPPORTEDANCH{\char"01AE\char"0351} \SUPPORTEDANCH{\char"0059\char"0351} \SUPPORTEDANCH{\char"01B3\char"0351} \SUPPORTEDANCH{\char"0061\char"0351} \SUPPORTEDANCH{\char"1E9A\char"0351} \SUPPORTEDANCH{\char"0062\char"0351} \SUPPORTEDANCH{\char"0183\char"0351} \SUPPORTEDANCH{\char"0063\char"0351} \SUPPORTEDANCH{\char"0188\char"0351} \SUPPORTEDANCH{\char"0064\char"0351} \SUPPORTEDANCH{\char"018C\char"0351} \SUPPORTEDANCH{\char"0221\char"0351} \SUPPORTEDANCH{\char"0068\char"0351} \SUPPORTEDANCH{\char"2C68\char"0351} \SUPPORTEDANCH{\char"006B\char"0351} \SUPPORTEDANCH{\char"0199\char"0351} \SUPPORTEDANCH{\char"2C6A\char"0351} \SUPPORTEDANCH{\char"006C\char"0351} \SUPPORTEDANCH{\char"0234\char"0351} \SUPPORTEDANCH{\char"006E\char"0351} \SUPPORTEDANCH{\char"0235\char"0351} \SUPPORTEDANCH{\char"0070\char"0351} \SUPPORTEDANCH{\char"01A5\char"0351} \SUPPORTEDANCH{\char"0071\char"0351} \SUPPORTEDANCH{\char"02A0\char"0351} \SUPPORTEDANCH{\char"0074\char"0351} \SUPPORTEDANCH{\char"01AB\char"0351} \SUPPORTEDANCH{\char"01AD\char"0351} \SUPPORTEDANCH{\char"0236\char"0351} \SUPPORTEDANCH{\char"0075\char"0351} \MISSINGGLYPH{\char"AB52\char"0351} \SUPPORTEDANCH{\char"0079\char"0351} \SUPPORTEDANCH{\char"01B4\char"0351} \MISSINGGLYPH{\char"1EFF\char"0351}

\SUPPORTEDANCH{\char"0048\char"0352} \SUPPORTEDANCH{\char"2C67\char"0352} \SUPPORTEDANCH{\char"004B\char"0352} \SUPPORTEDANCH{\char"2C69\char"0352} \SUPPORTEDANCH{\char"004D\char"0352} \MISSINGGLYPH{\char"2C6E\char"0352} \SUPPORTEDANCH{\char"0050\char"0352} \SUPPORTEDANCH{\char"01A4\char"0352} \SUPPORTEDANCH{\char"0052\char"0352} \SUPPORTEDANCH{\char"2C64\char"0352} \SUPPORTEDANCH{\char"0054\char"0352} \SUPPORTEDANCH{\char"01AC\char"0352} \SUPPORTEDANCH{\char"01AE\char"0352} \SUPPORTEDANCH{\char"0059\char"0352} \SUPPORTEDANCH{\char"01B3\char"0352} \SUPPORTEDANCH{\char"0061\char"0352} \SUPPORTEDANCH{\char"1E9A\char"0352} \SUPPORTEDANCH{\char"0062\char"0352} \SUPPORTEDANCH{\char"0183\char"0352} \SUPPORTEDANCH{\char"0063\char"0352} \SUPPORTEDANCH{\char"0188\char"0352} \SUPPORTEDANCH{\char"0064\char"0352} \SUPPORTEDANCH{\char"018C\char"0352} \SUPPORTEDANCH{\char"0221\char"0352} \SUPPORTEDANCH{\char"0068\char"0352} \SUPPORTEDANCH{\char"2C68\char"0352} \SUPPORTEDANCH{\char"006B\char"0352} \SUPPORTEDANCH{\char"0199\char"0352} \SUPPORTEDANCH{\char"2C6A\char"0352} \SUPPORTEDANCH{\char"006C\char"0352} \SUPPORTEDANCH{\char"0234\char"0352} \SUPPORTEDANCH{\char"006E\char"0352} \SUPPORTEDANCH{\char"0235\char"0352} \SUPPORTEDANCH{\char"0070\char"0352} \SUPPORTEDANCH{\char"01A5\char"0352} \SUPPORTEDANCH{\char"0071\char"0352} \SUPPORTEDANCH{\char"02A0\char"0352} \SUPPORTEDANCH{\char"0074\char"0352} \SUPPORTEDANCH{\char"01AB\char"0352} \SUPPORTEDANCH{\char"01AD\char"0352} \SUPPORTEDANCH{\char"0236\char"0352} \SUPPORTEDANCH{\char"0075\char"0352} \MISSINGGLYPH{\char"AB52\char"0352} \SUPPORTEDANCH{\char"0079\char"0352} \SUPPORTEDANCH{\char"01B4\char"0352} \MISSINGGLYPH{\char"1EFF\char"0352}

\SUPPORTEDANCH{\char"0048\char"0357} \SUPPORTEDANCH{\char"2C67\char"0357} \SUPPORTEDANCH{\char"004B\char"0357} \SUPPORTEDANCH{\char"2C69\char"0357} \SUPPORTEDANCH{\char"004D\char"0357} \MISSINGGLYPH{\char"2C6E\char"0357} \SUPPORTEDANCH{\char"0050\char"0357} \SUPPORTEDANCH{\char"01A4\char"0357} \SUPPORTEDANCH{\char"0052\char"0357} \SUPPORTEDANCH{\char"2C64\char"0357} \SUPPORTEDANCH{\char"0054\char"0357} \SUPPORTEDANCH{\char"01AC\char"0357} \SUPPORTEDANCH{\char"01AE\char"0357} \SUPPORTEDANCH{\char"0059\char"0357} \SUPPORTEDANCH{\char"01B3\char"0357} \SUPPORTEDANCH{\char"0061\char"0357} \SUPPORTEDANCH{\char"1E9A\char"0357} \SUPPORTEDANCH{\char"0062\char"0357} \SUPPORTEDANCH{\char"0183\char"0357} \SUPPORTEDANCH{\char"0063\char"0357} \SUPPORTEDANCH{\char"0188\char"0357} \SUPPORTEDANCH{\char"0064\char"0357} \SUPPORTEDANCH{\char"018C\char"0357} \SUPPORTEDANCH{\char"0221\char"0357} \SUPPORTEDANCH{\char"0068\char"0357} \SUPPORTEDANCH{\char"2C68\char"0357} \SUPPORTEDANCH{\char"006B\char"0357} \SUPPORTEDANCH{\char"0199\char"0357} \SUPPORTEDANCH{\char"2C6A\char"0357} \SUPPORTEDANCH{\char"006C\char"0357} \SUPPORTEDANCH{\char"0234\char"0357} \SUPPORTEDANCH{\char"006E\char"0357} \SUPPORTEDANCH{\char"0235\char"0357} \SUPPORTEDANCH{\char"0070\char"0357} \SUPPORTEDANCH{\char"01A5\char"0357} \SUPPORTEDANCH{\char"0071\char"0357} \SUPPORTEDANCH{\char"02A0\char"0357} \SUPPORTEDANCH{\char"0074\char"0357} \SUPPORTEDANCH{\char"01AB\char"0357} \SUPPORTEDANCH{\char"01AD\char"0357} \SUPPORTEDANCH{\char"0236\char"0357} \SUPPORTEDANCH{\char"0075\char"0357} \MISSINGGLYPH{\char"AB52\char"0357} \SUPPORTEDANCH{\char"0079\char"0357} \SUPPORTEDANCH{\char"01B4\char"0357} \MISSINGGLYPH{\char"1EFF\char"0357}

\SUPPORTEDANCH{\char"0048\char"035B} \SUPPORTEDANCH{\char"2C67\char"035B} \SUPPORTEDANCH{\char"004B\char"035B} \SUPPORTEDANCH{\char"2C69\char"035B} \SUPPORTEDANCH{\char"004D\char"035B} \MISSINGGLYPH{\char"2C6E\char"035B} \SUPPORTEDANCH{\char"0050\char"035B} \SUPPORTEDANCH{\char"01A4\char"035B} \SUPPORTEDANCH{\char"0052\char"035B} \SUPPORTEDANCH{\char"2C64\char"035B} \SUPPORTEDANCH{\char"0054\char"035B} \SUPPORTEDANCH{\char"01AC\char"035B} \SUPPORTEDANCH{\char"01AE\char"035B} \SUPPORTEDANCH{\char"0059\char"035B} \SUPPORTEDANCH{\char"01B3\char"035B} \SUPPORTEDANCH{\char"0061\char"035B} \SUPPORTEDANCH{\char"1E9A\char"035B} \SUPPORTEDANCH{\char"0062\char"035B} \SUPPORTEDANCH{\char"0183\char"035B} \SUPPORTEDANCH{\char"0063\char"035B} \SUPPORTEDANCH{\char"0188\char"035B} \SUPPORTEDANCH{\char"0064\char"035B} \SUPPORTEDANCH{\char"018C\char"035B} \SUPPORTEDANCH{\char"0221\char"035B} \SUPPORTEDANCH{\char"0068\char"035B} \SUPPORTEDANCH{\char"2C68\char"035B} \SUPPORTEDANCH{\char"006B\char"035B} \SUPPORTEDANCH{\char"0199\char"035B} \SUPPORTEDANCH{\char"2C6A\char"035B} \SUPPORTEDANCH{\char"006C\char"035B} \SUPPORTEDANCH{\char"0234\char"035B} \SUPPORTEDANCH{\char"006E\char"035B} \SUPPORTEDANCH{\char"0235\char"035B} \SUPPORTEDANCH{\char"0070\char"035B} \SUPPORTEDANCH{\char"01A5\char"035B} \SUPPORTEDANCH{\char"0071\char"035B} \SUPPORTEDANCH{\char"02A0\char"035B} \SUPPORTEDANCH{\char"0074\char"035B} \SUPPORTEDANCH{\char"01AB\char"035B} \SUPPORTEDANCH{\char"01AD\char"035B} \SUPPORTEDANCH{\char"0236\char"035B} \SUPPORTEDANCH{\char"0075\char"035B} \MISSINGGLYPH{\char"AB52\char"035B} \SUPPORTEDANCH{\char"0079\char"035B} \SUPPORTEDANCH{\char"01B4\char"035B} \MISSINGGLYPH{\char"1EFF\char"035B}

}

\newpage
\subsection*{Semibold italic, Group 4, above, Above}

{\bfseries\itshape
% grid. columns are bases, rows are marks.
\SUPPORTEDANCH{\char"0048\char"0300} \SUPPORTEDFALL{\char"2C67\char"0300} \SUPPORTEDANCH{\char"004B\char"0300} \SUPPORTEDANCH{\char"2C69\char"0300} \SUPPORTEDANCH{\char"004D\char"0300} \MISSINGGLYPH{\char"2C6E\char"0300} \SUPPORTEDANCH{\char"0050\char"0300} \SUPPORTEDANCH{\char"01A4\char"0300} \SUPPORTEDANCH{\char"0052\char"0300} \SUPPORTEDFALL{\char"2C64\char"0300} \SUPPORTEDANCH{\char"0054\char"0300} \SUPPORTEDFALL{\char"01AC\char"0300} \SUPPORTEDFALL{\char"01AE\char"0300} \SUPPORTEDPREC{\char"0059\char"0300} \SUPPORTEDFALL{\char"01B3\char"0300} \SUPPORTEDPREC{\char"0061\char"0300} \SUPPORTEDFALL{\char"1E9A\char"0300} \SUPPORTEDFALL{\char"0062\char"0300} \SUPPORTEDFALL{\char"0183\char"0300} \SUPPORTEDANCH{\char"0063\char"0300} \SUPPORTEDFALL{\char"0188\char"0300} \SUPPORTEDANCH{\char"0064\char"0300} \SUPPORTEDANCH{\char"018C\char"0300} \SUPPORTEDFALL{\char"0221\char"0300} \SUPPORTEDANCH{\char"0068\char"0300} \SUPPORTEDANCH{\char"2C68\char"0300} \SUPPORTEDANCH{\char"006B\char"0300} \SUPPORTEDFALL{\char"0199\char"0300} \SUPPORTEDANCH{\char"2C6A\char"0300} \SUPPORTEDANCH{\char"006C\char"0300} \SUPPORTEDFALL{\char"0234\char"0300} \SUPPORTEDPREC{\char"006E\char"0300} \SUPPORTEDANCH{\char"0235\char"0300} \SUPPORTEDANCH{\char"0070\char"0300} \SUPPORTEDFALL{\char"01A5\char"0300} \SUPPORTEDANCH{\char"0071\char"0300} \SUPPORTEDANCH{\char"02A0\char"0300} \SUPPORTEDFALL{\char"0074\char"0300} \SUPPORTEDFALL{\char"01AB\char"0300} \SUPPORTEDFALL{\char"01AD\char"0300} \SUPPORTEDFALL{\char"0236\char"0300} \SUPPORTEDPREC{\char"0075\char"0300} \MISSINGGLYPH{\char"AB52\char"0300} \SUPPORTEDPREC{\char"0079\char"0300} \SUPPORTEDANCH{\char"01B4\char"0300} \MISSINGGLYPH{\char"1EFF\char"0300}

\SUPPORTEDANCH{\char"0048\char"0301} \SUPPORTEDFALL{\char"2C67\char"0301} \SUPPORTEDPREC{\char"004B\char"0301} \SUPPORTEDANCH{\char"2C69\char"0301} \SUPPORTEDPREC{\char"004D\char"0301} \MISSINGGLYPH{\char"2C6E\char"0301} \SUPPORTEDPREC{\char"0050\char"0301} \SUPPORTEDANCH{\char"01A4\char"0301} \SUPPORTEDPREC{\char"0052\char"0301} \SUPPORTEDFALL{\char"2C64\char"0301} \SUPPORTEDANCH{\char"0054\char"0301} \SUPPORTEDFALL{\char"01AC\char"0301} \SUPPORTEDFALL{\char"01AE\char"0301} \SUPPORTEDPREC{\char"0059\char"0301} \SUPPORTEDFALL{\char"01B3\char"0301} \SUPPORTEDPREC{\char"0061\char"0301} \SUPPORTEDFALL{\char"1E9A\char"0301} \SUPPORTEDFALL{\char"0062\char"0301} \SUPPORTEDFALL{\char"0183\char"0301} \SUPPORTEDPREC{\char"0063\char"0301} \SUPPORTEDFALL{\char"0188\char"0301} \SUPPORTEDANCH{\char"0064\char"0301} \SUPPORTEDANCH{\char"018C\char"0301} \SUPPORTEDFALL{\char"0221\char"0301} \SUPPORTEDANCH{\char"0068\char"0301} \SUPPORTEDANCH{\char"2C68\char"0301} \SUPPORTEDPREC{\char"006B\char"0301} \SUPPORTEDFALL{\char"0199\char"0301} \SUPPORTEDANCH{\char"2C6A\char"0301} \SUPPORTEDPREC{\char"006C\char"0301} \SUPPORTEDFALL{\char"0234\char"0301} \SUPPORTEDPREC{\char"006E\char"0301} \SUPPORTEDANCH{\char"0235\char"0301} \SUPPORTEDPREC{\char"0070\char"0301} \SUPPORTEDFALL{\char"01A5\char"0301} \SUPPORTEDANCH{\char"0071\char"0301} \SUPPORTEDANCH{\char"02A0\char"0301} \SUPPORTEDFALL{\char"0074\char"0301} \SUPPORTEDFALL{\char"01AB\char"0301} \SUPPORTEDFALL{\char"01AD\char"0301} \SUPPORTEDFALL{\char"0236\char"0301} \SUPPORTEDPREC{\char"0075\char"0301} \MISSINGGLYPH{\char"AB52\char"0301} \SUPPORTEDPREC{\char"0079\char"0301} \SUPPORTEDANCH{\char"01B4\char"0301} \MISSINGGLYPH{\char"1EFF\char"0301}

\SUPPORTEDPREC{\char"0048\char"0302} \SUPPORTEDFALL{\char"2C67\char"0302} \SUPPORTEDANCH{\char"004B\char"0302} \SUPPORTEDANCH{\char"2C69\char"0302} \SUPPORTEDANCH{\char"004D\char"0302} \MISSINGGLYPH{\char"2C6E\char"0302} \SUPPORTEDANCH{\char"0050\char"0302} \SUPPORTEDANCH{\char"01A4\char"0302} \SUPPORTEDANCH{\char"0052\char"0302} \SUPPORTEDFALL{\char"2C64\char"0302} \SUPPORTEDANCH{\char"0054\char"0302} \SUPPORTEDFALL{\char"01AC\char"0302} \SUPPORTEDFALL{\char"01AE\char"0302} \SUPPORTEDPREC{\char"0059\char"0302} \SUPPORTEDFALL{\char"01B3\char"0302} \SUPPORTEDPREC{\char"0061\char"0302} \SUPPORTEDFALL{\char"1E9A\char"0302} \SUPPORTEDFALL{\char"0062\char"0302} \SUPPORTEDFALL{\char"0183\char"0302} \SUPPORTEDPREC{\char"0063\char"0302} \SUPPORTEDFALL{\char"0188\char"0302} \SUPPORTEDANCH{\char"0064\char"0302} \SUPPORTEDANCH{\char"018C\char"0302} \SUPPORTEDFALL{\char"0221\char"0302} \SUPPORTEDPREC{\char"0068\char"0302} \SUPPORTEDANCH{\char"2C68\char"0302} \SUPPORTEDANCH{\char"006B\char"0302} \SUPPORTEDFALL{\char"0199\char"0302} \SUPPORTEDANCH{\char"2C6A\char"0302} \SUPPORTEDANCH{\char"006C\char"0302} \SUPPORTEDFALL{\char"0234\char"0302} \SUPPORTEDANCH{\char"006E\char"0302} \SUPPORTEDANCH{\char"0235\char"0302} \SUPPORTEDANCH{\char"0070\char"0302} \SUPPORTEDFALL{\char"01A5\char"0302} \SUPPORTEDANCH{\char"0071\char"0302} \SUPPORTEDANCH{\char"02A0\char"0302} \SUPPORTEDFALL{\char"0074\char"0302} \SUPPORTEDFALL{\char"01AB\char"0302} \SUPPORTEDFALL{\char"01AD\char"0302} \SUPPORTEDFALL{\char"0236\char"0302} \SUPPORTEDPREC{\char"0075\char"0302} \MISSINGGLYPH{\char"AB52\char"0302} \SUPPORTEDPREC{\char"0079\char"0302} \SUPPORTEDANCH{\char"01B4\char"0302} \MISSINGGLYPH{\char"1EFF\char"0302}

\SUPPORTEDANCH{\char"0048\char"0303} \SUPPORTEDFALL{\char"2C67\char"0303} \SUPPORTEDANCH{\char"004B\char"0303} \SUPPORTEDANCH{\char"2C69\char"0303} \SUPPORTEDANCH{\char"004D\char"0303} \MISSINGGLYPH{\char"2C6E\char"0303} \SUPPORTEDANCH{\char"0050\char"0303} \SUPPORTEDANCH{\char"01A4\char"0303} \SUPPORTEDANCH{\char"0052\char"0303} \SUPPORTEDFALL{\char"2C64\char"0303} \SUPPORTEDANCH{\char"0054\char"0303} \SUPPORTEDFALL{\char"01AC\char"0303} \SUPPORTEDFALL{\char"01AE\char"0303} \SUPPORTEDPREC{\char"0059\char"0303} \SUPPORTEDFALL{\char"01B3\char"0303} \SUPPORTEDPREC{\char"0061\char"0303} \SUPPORTEDFALL{\char"1E9A\char"0303} \SUPPORTEDFALL{\char"0062\char"0303} \SUPPORTEDFALL{\char"0183\char"0303} \SUPPORTEDANCH{\char"0063\char"0303} \SUPPORTEDFALL{\char"0188\char"0303} \SUPPORTEDANCH{\char"0064\char"0303} \SUPPORTEDANCH{\char"018C\char"0303} \SUPPORTEDFALL{\char"0221\char"0303} \SUPPORTEDANCH{\char"0068\char"0303} \SUPPORTEDANCH{\char"2C68\char"0303} \SUPPORTEDANCH{\char"006B\char"0303} \SUPPORTEDFALL{\char"0199\char"0303} \SUPPORTEDANCH{\char"2C6A\char"0303} \SUPPORTEDANCH{\char"006C\char"0303} \SUPPORTEDFALL{\char"0234\char"0303} \SUPPORTEDPREC{\char"006E\char"0303} \SUPPORTEDANCH{\char"0235\char"0303} \SUPPORTEDANCH{\char"0070\char"0303} \SUPPORTEDFALL{\char"01A5\char"0303} \SUPPORTEDANCH{\char"0071\char"0303} \SUPPORTEDANCH{\char"02A0\char"0303} \SUPPORTEDFALL{\char"0074\char"0303} \SUPPORTEDFALL{\char"01AB\char"0303} \SUPPORTEDFALL{\char"01AD\char"0303} \SUPPORTEDFALL{\char"0236\char"0303} \SUPPORTEDPREC{\char"0075\char"0303} \MISSINGGLYPH{\char"AB52\char"0303} \SUPPORTEDPREC{\char"0079\char"0303} \SUPPORTEDANCH{\char"01B4\char"0303} \MISSINGGLYPH{\char"1EFF\char"0303}

\SUPPORTEDANCH{\char"0048\char"0304} \SUPPORTEDFALL{\char"2C67\char"0304} \SUPPORTEDANCH{\char"004B\char"0304} \SUPPORTEDANCH{\char"2C69\char"0304} \SUPPORTEDANCH{\char"004D\char"0304} \MISSINGGLYPH{\char"2C6E\char"0304} \SUPPORTEDANCH{\char"0050\char"0304} \SUPPORTEDANCH{\char"01A4\char"0304} \SUPPORTEDANCH{\char"0052\char"0304} \SUPPORTEDFALL{\char"2C64\char"0304} \SUPPORTEDANCH{\char"0054\char"0304} \SUPPORTEDFALL{\char"01AC\char"0304} \SUPPORTEDFALL{\char"01AE\char"0304} \SUPPORTEDPREC{\char"0059\char"0304} \SUPPORTEDFALL{\char"01B3\char"0304} \SUPPORTEDPREC{\char"0061\char"0304} \SUPPORTEDFALL{\char"1E9A\char"0304} \SUPPORTEDFALL{\char"0062\char"0304} \SUPPORTEDFALL{\char"0183\char"0304} \SUPPORTEDANCH{\char"0063\char"0304} \SUPPORTEDFALL{\char"0188\char"0304} \SUPPORTEDANCH{\char"0064\char"0304} \SUPPORTEDANCH{\char"018C\char"0304} \SUPPORTEDFALL{\char"0221\char"0304} \SUPPORTEDANCH{\char"0068\char"0304} \SUPPORTEDANCH{\char"2C68\char"0304} \SUPPORTEDANCH{\char"006B\char"0304} \SUPPORTEDFALL{\char"0199\char"0304} \SUPPORTEDANCH{\char"2C6A\char"0304} \SUPPORTEDANCH{\char"006C\char"0304} \SUPPORTEDFALL{\char"0234\char"0304} \SUPPORTEDANCH{\char"006E\char"0304} \SUPPORTEDANCH{\char"0235\char"0304} \SUPPORTEDANCH{\char"0070\char"0304} \SUPPORTEDFALL{\char"01A5\char"0304} \SUPPORTEDANCH{\char"0071\char"0304} \SUPPORTEDANCH{\char"02A0\char"0304} \SUPPORTEDFALL{\char"0074\char"0304} \SUPPORTEDFALL{\char"01AB\char"0304} \SUPPORTEDFALL{\char"01AD\char"0304} \SUPPORTEDFALL{\char"0236\char"0304} \SUPPORTEDPREC{\char"0075\char"0304} \MISSINGGLYPH{\char"AB52\char"0304} \SUPPORTEDPREC{\char"0079\char"0304} \SUPPORTEDANCH{\char"01B4\char"0304} \MISSINGGLYPH{\char"1EFF\char"0304}

\SUPPORTEDANCH{\char"0048\char"0306} \SUPPORTEDFALL{\char"2C67\char"0306} \SUPPORTEDANCH{\char"004B\char"0306} \SUPPORTEDANCH{\char"2C69\char"0306} \SUPPORTEDANCH{\char"004D\char"0306} \MISSINGGLYPH{\char"2C6E\char"0306} \SUPPORTEDANCH{\char"0050\char"0306} \SUPPORTEDANCH{\char"01A4\char"0306} \SUPPORTEDANCH{\char"0052\char"0306} \SUPPORTEDFALL{\char"2C64\char"0306} \SUPPORTEDANCH{\char"0054\char"0306} \SUPPORTEDFALL{\char"01AC\char"0306} \SUPPORTEDFALL{\char"01AE\char"0306} \SUPPORTEDANCH{\char"0059\char"0306} \SUPPORTEDFALL{\char"01B3\char"0306} \SUPPORTEDPREC{\char"0061\char"0306} \SUPPORTEDFALL{\char"1E9A\char"0306} \SUPPORTEDFALL{\char"0062\char"0306} \SUPPORTEDFALL{\char"0183\char"0306} \SUPPORTEDANCH{\char"0063\char"0306} \SUPPORTEDFALL{\char"0188\char"0306} \SUPPORTEDANCH{\char"0064\char"0306} \SUPPORTEDANCH{\char"018C\char"0306} \SUPPORTEDFALL{\char"0221\char"0306} \SUPPORTEDANCH{\char"0068\char"0306} \SUPPORTEDANCH{\char"2C68\char"0306} \SUPPORTEDANCH{\char"006B\char"0306} \SUPPORTEDFALL{\char"0199\char"0306} \SUPPORTEDANCH{\char"2C6A\char"0306} \SUPPORTEDANCH{\char"006C\char"0306} \SUPPORTEDFALL{\char"0234\char"0306} \SUPPORTEDANCH{\char"006E\char"0306} \SUPPORTEDANCH{\char"0235\char"0306} \SUPPORTEDANCH{\char"0070\char"0306} \SUPPORTEDFALL{\char"01A5\char"0306} \SUPPORTEDANCH{\char"0071\char"0306} \SUPPORTEDANCH{\char"02A0\char"0306} \SUPPORTEDFALL{\char"0074\char"0306} \SUPPORTEDFALL{\char"01AB\char"0306} \SUPPORTEDFALL{\char"01AD\char"0306} \SUPPORTEDFALL{\char"0236\char"0306} \SUPPORTEDPREC{\char"0075\char"0306} \MISSINGGLYPH{\char"AB52\char"0306} \SUPPORTEDANCH{\char"0079\char"0306} \SUPPORTEDANCH{\char"01B4\char"0306} \MISSINGGLYPH{\char"1EFF\char"0306}

\SUPPORTEDPREC{\char"0048\char"0307} \SUPPORTEDFALL{\char"2C67\char"0307} \SUPPORTEDANCH{\char"004B\char"0307} \SUPPORTEDANCH{\char"2C69\char"0307} \SUPPORTEDPREC{\char"004D\char"0307} \MISSINGGLYPH{\char"2C6E\char"0307} \SUPPORTEDPREC{\char"0050\char"0307} \SUPPORTEDANCH{\char"01A4\char"0307} \SUPPORTEDPREC{\char"0052\char"0307} \SUPPORTEDFALL{\char"2C64\char"0307} \SUPPORTEDPREC{\char"0054\char"0307} \SUPPORTEDFALL{\char"01AC\char"0307} \SUPPORTEDFALL{\char"01AE\char"0307} \SUPPORTEDPREC{\char"0059\char"0307} \SUPPORTEDFALL{\char"01B3\char"0307} \SUPPORTEDPREC{\char"0061\char"0307} \SUPPORTEDFALL{\char"1E9A\char"0307} \SUPPORTEDPREC{\char"0062\char"0307} \SUPPORTEDFALL{\char"0183\char"0307} \SUPPORTEDPREC{\char"0063\char"0307} \SUPPORTEDFALL{\char"0188\char"0307} \SUPPORTEDPREC{\char"0064\char"0307} \SUPPORTEDANCH{\char"018C\char"0307} \SUPPORTEDFALL{\char"0221\char"0307} \SUPPORTEDPREC{\char"0068\char"0307} \SUPPORTEDANCH{\char"2C68\char"0307} \SUPPORTEDANCH{\char"006B\char"0307} \SUPPORTEDFALL{\char"0199\char"0307} \SUPPORTEDANCH{\char"2C6A\char"0307} \SUPPORTEDANCH{\char"006C\char"0307} \SUPPORTEDFALL{\char"0234\char"0307} \SUPPORTEDPREC{\char"006E\char"0307} \SUPPORTEDANCH{\char"0235\char"0307} \SUPPORTEDPREC{\char"0070\char"0307} \SUPPORTEDFALL{\char"01A5\char"0307} \SUPPORTEDANCH{\char"0071\char"0307} \SUPPORTEDANCH{\char"02A0\char"0307} \SUPPORTEDPREC{\char"0074\char"0307} \SUPPORTEDFALL{\char"01AB\char"0307} \SUPPORTEDFALL{\char"01AD\char"0307} \SUPPORTEDFALL{\char"0236\char"0307} \SUPPORTEDANCH{\char"0075\char"0307} \MISSINGGLYPH{\char"AB52\char"0307} \SUPPORTEDPREC{\char"0079\char"0307} \SUPPORTEDANCH{\char"01B4\char"0307} \MISSINGGLYPH{\char"1EFF\char"0307}

\SUPPORTEDPREC{\char"0048\char"0308} \SUPPORTEDFALL{\char"2C67\char"0308} \SUPPORTEDANCH{\char"004B\char"0308} \SUPPORTEDANCH{\char"2C69\char"0308} \SUPPORTEDANCH{\char"004D\char"0308} \MISSINGGLYPH{\char"2C6E\char"0308} \SUPPORTEDANCH{\char"0050\char"0308} \SUPPORTEDANCH{\char"01A4\char"0308} \SUPPORTEDANCH{\char"0052\char"0308} \SUPPORTEDFALL{\char"2C64\char"0308} \SUPPORTEDANCH{\char"0054\char"0308} \SUPPORTEDFALL{\char"01AC\char"0308} \SUPPORTEDFALL{\char"01AE\char"0308} \SUPPORTEDPREC{\char"0059\char"0308} \SUPPORTEDFALL{\char"01B3\char"0308} \SUPPORTEDPREC{\char"0061\char"0308} \SUPPORTEDFALL{\char"1E9A\char"0308} \SUPPORTEDFALL{\char"0062\char"0308} \SUPPORTEDFALL{\char"0183\char"0308} \SUPPORTEDANCH{\char"0063\char"0308} \SUPPORTEDFALL{\char"0188\char"0308} \SUPPORTEDANCH{\char"0064\char"0308} \SUPPORTEDANCH{\char"018C\char"0308} \SUPPORTEDFALL{\char"0221\char"0308} \SUPPORTEDPREC{\char"0068\char"0308} \SUPPORTEDANCH{\char"2C68\char"0308} \SUPPORTEDANCH{\char"006B\char"0308} \SUPPORTEDFALL{\char"0199\char"0308} \SUPPORTEDANCH{\char"2C6A\char"0308} \SUPPORTEDANCH{\char"006C\char"0308} \SUPPORTEDFALL{\char"0234\char"0308} \SUPPORTEDANCH{\char"006E\char"0308} \SUPPORTEDANCH{\char"0235\char"0308} \SUPPORTEDANCH{\char"0070\char"0308} \SUPPORTEDFALL{\char"01A5\char"0308} \SUPPORTEDANCH{\char"0071\char"0308} \SUPPORTEDANCH{\char"02A0\char"0308} \SUPPORTEDPREC{\char"0074\char"0308} \SUPPORTEDFALL{\char"01AB\char"0308} \SUPPORTEDFALL{\char"01AD\char"0308} \SUPPORTEDFALL{\char"0236\char"0308} \SUPPORTEDPREC{\char"0075\char"0308} \MISSINGGLYPH{\char"AB52\char"0308} \SUPPORTEDPREC{\char"0079\char"0308} \SUPPORTEDANCH{\char"01B4\char"0308} \MISSINGGLYPH{\char"1EFF\char"0308}

\SUPPORTEDANCH{\char"0048\char"0309} \SUPPORTEDFALL{\char"2C67\char"0309} \SUPPORTEDANCH{\char"004B\char"0309} \SUPPORTEDANCH{\char"2C69\char"0309} \SUPPORTEDANCH{\char"004D\char"0309} \MISSINGGLYPH{\char"2C6E\char"0309} \SUPPORTEDANCH{\char"0050\char"0309} \SUPPORTEDANCH{\char"01A4\char"0309} \SUPPORTEDANCH{\char"0052\char"0309} \SUPPORTEDFALL{\char"2C64\char"0309} \SUPPORTEDANCH{\char"0054\char"0309} \SUPPORTEDFALL{\char"01AC\char"0309} \SUPPORTEDFALL{\char"01AE\char"0309} \SUPPORTEDPREC{\char"0059\char"0309} \SUPPORTEDFALL{\char"01B3\char"0309} \SUPPORTEDPREC{\char"0061\char"0309} \SUPPORTEDFALL{\char"1E9A\char"0309} \SUPPORTEDFALL{\char"0062\char"0309} \SUPPORTEDFALL{\char"0183\char"0309} \SUPPORTEDANCH{\char"0063\char"0309} \SUPPORTEDFALL{\char"0188\char"0309} \SUPPORTEDANCH{\char"0064\char"0309} \SUPPORTEDANCH{\char"018C\char"0309} \SUPPORTEDFALL{\char"0221\char"0309} \SUPPORTEDANCH{\char"0068\char"0309} \SUPPORTEDANCH{\char"2C68\char"0309} \SUPPORTEDANCH{\char"006B\char"0309} \SUPPORTEDFALL{\char"0199\char"0309} \SUPPORTEDANCH{\char"2C6A\char"0309} \SUPPORTEDANCH{\char"006C\char"0309} \SUPPORTEDFALL{\char"0234\char"0309} \SUPPORTEDANCH{\char"006E\char"0309} \SUPPORTEDANCH{\char"0235\char"0309} \SUPPORTEDANCH{\char"0070\char"0309} \SUPPORTEDFALL{\char"01A5\char"0309} \SUPPORTEDANCH{\char"0071\char"0309} \SUPPORTEDANCH{\char"02A0\char"0309} \SUPPORTEDFALL{\char"0074\char"0309} \SUPPORTEDFALL{\char"01AB\char"0309} \SUPPORTEDFALL{\char"01AD\char"0309} \SUPPORTEDFALL{\char"0236\char"0309} \SUPPORTEDPREC{\char"0075\char"0309} \MISSINGGLYPH{\char"AB52\char"0309} \SUPPORTEDPREC{\char"0079\char"0309} \SUPPORTEDANCH{\char"01B4\char"0309} \MISSINGGLYPH{\char"1EFF\char"0309}

\SUPPORTEDANCH{\char"0048\char"030A} \SUPPORTEDFALL{\char"2C67\char"030A} \SUPPORTEDANCH{\char"004B\char"030A} \SUPPORTEDANCH{\char"2C69\char"030A} \SUPPORTEDANCH{\char"004D\char"030A} \MISSINGGLYPH{\char"2C6E\char"030A} \SUPPORTEDANCH{\char"0050\char"030A} \SUPPORTEDANCH{\char"01A4\char"030A} \SUPPORTEDANCH{\char"0052\char"030A} \SUPPORTEDFALL{\char"2C64\char"030A} \SUPPORTEDANCH{\char"0054\char"030A} \SUPPORTEDFALL{\char"01AC\char"030A} \SUPPORTEDFALL{\char"01AE\char"030A} \SUPPORTEDANCH{\char"0059\char"030A} \SUPPORTEDFALL{\char"01B3\char"030A} \SUPPORTEDPREC{\char"0061\char"030A} \SUPPORTEDFALL{\char"1E9A\char"030A} \SUPPORTEDFALL{\char"0062\char"030A} \SUPPORTEDFALL{\char"0183\char"030A} \SUPPORTEDANCH{\char"0063\char"030A} \SUPPORTEDFALL{\char"0188\char"030A} \SUPPORTEDANCH{\char"0064\char"030A} \SUPPORTEDANCH{\char"018C\char"030A} \SUPPORTEDFALL{\char"0221\char"030A} \SUPPORTEDANCH{\char"0068\char"030A} \SUPPORTEDANCH{\char"2C68\char"030A} \SUPPORTEDANCH{\char"006B\char"030A} \SUPPORTEDFALL{\char"0199\char"030A} \SUPPORTEDANCH{\char"2C6A\char"030A} \SUPPORTEDANCH{\char"006C\char"030A} \SUPPORTEDFALL{\char"0234\char"030A} \SUPPORTEDANCH{\char"006E\char"030A} \SUPPORTEDANCH{\char"0235\char"030A} \SUPPORTEDANCH{\char"0070\char"030A} \SUPPORTEDFALL{\char"01A5\char"030A} \SUPPORTEDANCH{\char"0071\char"030A} \SUPPORTEDANCH{\char"02A0\char"030A} \SUPPORTEDFALL{\char"0074\char"030A} \SUPPORTEDFALL{\char"01AB\char"030A} \SUPPORTEDFALL{\char"01AD\char"030A} \SUPPORTEDFALL{\char"0236\char"030A} \SUPPORTEDPREC{\char"0075\char"030A} \MISSINGGLYPH{\char"AB52\char"030A} \SUPPORTEDPREC{\char"0079\char"030A} \SUPPORTEDANCH{\char"01B4\char"030A} \MISSINGGLYPH{\char"1EFF\char"030A}

\SUPPORTEDANCH{\char"0048\char"030B} \SUPPORTEDFALL{\char"2C67\char"030B} \SUPPORTEDANCH{\char"004B\char"030B} \SUPPORTEDANCH{\char"2C69\char"030B} \SUPPORTEDANCH{\char"004D\char"030B} \MISSINGGLYPH{\char"2C6E\char"030B} \SUPPORTEDANCH{\char"0050\char"030B} \SUPPORTEDANCH{\char"01A4\char"030B} \SUPPORTEDANCH{\char"0052\char"030B} \SUPPORTEDFALL{\char"2C64\char"030B} \SUPPORTEDANCH{\char"0054\char"030B} \SUPPORTEDFALL{\char"01AC\char"030B} \SUPPORTEDFALL{\char"01AE\char"030B} \SUPPORTEDANCH{\char"0059\char"030B} \SUPPORTEDFALL{\char"01B3\char"030B} \SUPPORTEDANCH{\char"0061\char"030B} \SUPPORTEDFALL{\char"1E9A\char"030B} \SUPPORTEDFALL{\char"0062\char"030B} \SUPPORTEDFALL{\char"0183\char"030B} \SUPPORTEDANCH{\char"0063\char"030B} \SUPPORTEDFALL{\char"0188\char"030B} \SUPPORTEDANCH{\char"0064\char"030B} \SUPPORTEDANCH{\char"018C\char"030B} \SUPPORTEDFALL{\char"0221\char"030B} \SUPPORTEDANCH{\char"0068\char"030B} \SUPPORTEDANCH{\char"2C68\char"030B} \SUPPORTEDANCH{\char"006B\char"030B} \SUPPORTEDFALL{\char"0199\char"030B} \SUPPORTEDANCH{\char"2C6A\char"030B} \SUPPORTEDANCH{\char"006C\char"030B} \SUPPORTEDFALL{\char"0234\char"030B} \SUPPORTEDANCH{\char"006E\char"030B} \SUPPORTEDANCH{\char"0235\char"030B} \SUPPORTEDANCH{\char"0070\char"030B} \SUPPORTEDFALL{\char"01A5\char"030B} \SUPPORTEDANCH{\char"0071\char"030B} \SUPPORTEDANCH{\char"02A0\char"030B} \SUPPORTEDFALL{\char"0074\char"030B} \SUPPORTEDFALL{\char"01AB\char"030B} \SUPPORTEDFALL{\char"01AD\char"030B} \SUPPORTEDFALL{\char"0236\char"030B} \SUPPORTEDPREC{\char"0075\char"030B} \MISSINGGLYPH{\char"AB52\char"030B} \SUPPORTEDANCH{\char"0079\char"030B} \SUPPORTEDANCH{\char"01B4\char"030B} \MISSINGGLYPH{\char"1EFF\char"030B}

\SUPPORTEDPREC{\char"0048\char"030C} \SUPPORTEDFALL{\char"2C67\char"030C} \SUPPORTEDPREC{\char"004B\char"030C} \SUPPORTEDANCH{\char"2C69\char"030C} \SUPPORTEDANCH{\char"004D\char"030C} \MISSINGGLYPH{\char"2C6E\char"030C} \SUPPORTEDANCH{\char"0050\char"030C} \SUPPORTEDANCH{\char"01A4\char"030C} \SUPPORTEDPREC{\char"0052\char"030C} \SUPPORTEDFALL{\char"2C64\char"030C} \SUPPORTEDPREC{\char"0054\char"030C} \SUPPORTEDFALL{\char"01AC\char"030C} \SUPPORTEDFALL{\char"01AE\char"030C} \SUPPORTEDANCH{\char"0059\char"030C} \SUPPORTEDFALL{\char"01B3\char"030C} \SUPPORTEDPREC{\char"0061\char"030C} \SUPPORTEDFALL{\char"1E9A\char"030C} \SUPPORTEDFALL{\char"0062\char"030C} \SUPPORTEDFALL{\char"0183\char"030C} \SUPPORTEDPREC{\char"0063\char"030C} \SUPPORTEDFALL{\char"0188\char"030C} \SUPPORTEDPREC{\char"0064\char"030C} \SUPPORTEDANCH{\char"018C\char"030C} \SUPPORTEDFALL{\char"0221\char"030C} \SUPPORTEDPREC{\char"0068\char"030C} \SUPPORTEDANCH{\char"2C68\char"030C} \SUPPORTEDPREC{\char"006B\char"030C} \SUPPORTEDFALL{\char"0199\char"030C} \SUPPORTEDANCH{\char"2C6A\char"030C} \SUPPORTEDPREC{\char"006C\char"030C} \SUPPORTEDFALL{\char"0234\char"030C} \SUPPORTEDPREC{\char"006E\char"030C} \SUPPORTEDANCH{\char"0235\char"030C} \SUPPORTEDANCH{\char"0070\char"030C} \SUPPORTEDFALL{\char"01A5\char"030C} \SUPPORTEDANCH{\char"0071\char"030C} \SUPPORTEDANCH{\char"02A0\char"030C} \SUPPORTEDPREC{\char"0074\char"030C} \SUPPORTEDFALL{\char"01AB\char"030C} \SUPPORTEDFALL{\char"01AD\char"030C} \SUPPORTEDFALL{\char"0236\char"030C} \SUPPORTEDPREC{\char"0075\char"030C} \MISSINGGLYPH{\char"AB52\char"030C} \SUPPORTEDANCH{\char"0079\char"030C} \SUPPORTEDANCH{\char"01B4\char"030C} \MISSINGGLYPH{\char"1EFF\char"030C}

\SUPPORTEDANCH{\char"0048\char"030D} \SUPPORTEDFALL{\char"2C67\char"030D} \SUPPORTEDANCH{\char"004B\char"030D} \SUPPORTEDANCH{\char"2C69\char"030D} \SUPPORTEDANCH{\char"004D\char"030D} \MISSINGGLYPH{\char"2C6E\char"030D} \SUPPORTEDANCH{\char"0050\char"030D} \SUPPORTEDANCH{\char"01A4\char"030D} \SUPPORTEDANCH{\char"0052\char"030D} \SUPPORTEDFALL{\char"2C64\char"030D} \SUPPORTEDANCH{\char"0054\char"030D} \SUPPORTEDFALL{\char"01AC\char"030D} \SUPPORTEDFALL{\char"01AE\char"030D} \SUPPORTEDANCH{\char"0059\char"030D} \SUPPORTEDFALL{\char"01B3\char"030D} \SUPPORTEDANCH{\char"0061\char"030D} \SUPPORTEDFALL{\char"1E9A\char"030D} \SUPPORTEDFALL{\char"0062\char"030D} \SUPPORTEDFALL{\char"0183\char"030D} \SUPPORTEDANCH{\char"0063\char"030D} \SUPPORTEDFALL{\char"0188\char"030D} \SUPPORTEDANCH{\char"0064\char"030D} \SUPPORTEDANCH{\char"018C\char"030D} \SUPPORTEDFALL{\char"0221\char"030D} \SUPPORTEDANCH{\char"0068\char"030D} \SUPPORTEDANCH{\char"2C68\char"030D} \SUPPORTEDANCH{\char"006B\char"030D} \SUPPORTEDFALL{\char"0199\char"030D} \SUPPORTEDANCH{\char"2C6A\char"030D} \SUPPORTEDANCH{\char"006C\char"030D} \SUPPORTEDFALL{\char"0234\char"030D} \SUPPORTEDANCH{\char"006E\char"030D} \SUPPORTEDANCH{\char"0235\char"030D} \SUPPORTEDANCH{\char"0070\char"030D} \SUPPORTEDFALL{\char"01A5\char"030D} \SUPPORTEDANCH{\char"0071\char"030D} \SUPPORTEDANCH{\char"02A0\char"030D} \SUPPORTEDFALL{\char"0074\char"030D} \SUPPORTEDFALL{\char"01AB\char"030D} \SUPPORTEDFALL{\char"01AD\char"030D} \SUPPORTEDFALL{\char"0236\char"030D} \SUPPORTEDANCH{\char"0075\char"030D} \MISSINGGLYPH{\char"AB52\char"030D} \SUPPORTEDANCH{\char"0079\char"030D} \SUPPORTEDANCH{\char"01B4\char"030D} \MISSINGGLYPH{\char"1EFF\char"030D}

\SUPPORTEDANCH{\char"0048\char"030F} \SUPPORTEDFALL{\char"2C67\char"030F} \SUPPORTEDANCH{\char"004B\char"030F} \SUPPORTEDANCH{\char"2C69\char"030F} \SUPPORTEDANCH{\char"004D\char"030F} \MISSINGGLYPH{\char"2C6E\char"030F} \SUPPORTEDANCH{\char"0050\char"030F} \SUPPORTEDANCH{\char"01A4\char"030F} \SUPPORTEDPREC{\char"0052\char"030F} \SUPPORTEDFALL{\char"2C64\char"030F} \SUPPORTEDANCH{\char"0054\char"030F} \SUPPORTEDFALL{\char"01AC\char"030F} \SUPPORTEDFALL{\char"01AE\char"030F} \SUPPORTEDANCH{\char"0059\char"030F} \SUPPORTEDFALL{\char"01B3\char"030F} \SUPPORTEDPREC{\char"0061\char"030F} \SUPPORTEDFALL{\char"1E9A\char"030F} \SUPPORTEDFALL{\char"0062\char"030F} \SUPPORTEDFALL{\char"0183\char"030F} \SUPPORTEDANCH{\char"0063\char"030F} \SUPPORTEDFALL{\char"0188\char"030F} \SUPPORTEDANCH{\char"0064\char"030F} \SUPPORTEDANCH{\char"018C\char"030F} \SUPPORTEDFALL{\char"0221\char"030F} \SUPPORTEDANCH{\char"0068\char"030F} \SUPPORTEDANCH{\char"2C68\char"030F} \SUPPORTEDANCH{\char"006B\char"030F} \SUPPORTEDFALL{\char"0199\char"030F} \SUPPORTEDANCH{\char"2C6A\char"030F} \SUPPORTEDANCH{\char"006C\char"030F} \SUPPORTEDFALL{\char"0234\char"030F} \SUPPORTEDANCH{\char"006E\char"030F} \SUPPORTEDANCH{\char"0235\char"030F} \SUPPORTEDANCH{\char"0070\char"030F} \SUPPORTEDFALL{\char"01A5\char"030F} \SUPPORTEDANCH{\char"0071\char"030F} \SUPPORTEDANCH{\char"02A0\char"030F} \SUPPORTEDFALL{\char"0074\char"030F} \SUPPORTEDFALL{\char"01AB\char"030F} \SUPPORTEDFALL{\char"01AD\char"030F} \SUPPORTEDFALL{\char"0236\char"030F} \SUPPORTEDPREC{\char"0075\char"030F} \MISSINGGLYPH{\char"AB52\char"030F} \SUPPORTEDANCH{\char"0079\char"030F} \SUPPORTEDANCH{\char"01B4\char"030F} \MISSINGGLYPH{\char"1EFF\char"030F}

\SUPPORTEDANCH{\char"0048\char"0310} \SUPPORTEDFALL{\char"2C67\char"0310} \SUPPORTEDANCH{\char"004B\char"0310} \SUPPORTEDANCH{\char"2C69\char"0310} \SUPPORTEDANCH{\char"004D\char"0310} \MISSINGGLYPH{\char"2C6E\char"0310} \SUPPORTEDANCH{\char"0050\char"0310} \SUPPORTEDANCH{\char"01A4\char"0310} \SUPPORTEDANCH{\char"0052\char"0310} \SUPPORTEDFALL{\char"2C64\char"0310} \SUPPORTEDANCH{\char"0054\char"0310} \SUPPORTEDFALL{\char"01AC\char"0310} \SUPPORTEDFALL{\char"01AE\char"0310} \SUPPORTEDANCH{\char"0059\char"0310} \SUPPORTEDFALL{\char"01B3\char"0310} \SUPPORTEDANCH{\char"0061\char"0310} \SUPPORTEDFALL{\char"1E9A\char"0310} \SUPPORTEDFALL{\char"0062\char"0310} \SUPPORTEDFALL{\char"0183\char"0310} \SUPPORTEDANCH{\char"0063\char"0310} \SUPPORTEDFALL{\char"0188\char"0310} \SUPPORTEDANCH{\char"0064\char"0310} \SUPPORTEDANCH{\char"018C\char"0310} \SUPPORTEDFALL{\char"0221\char"0310} \SUPPORTEDANCH{\char"0068\char"0310} \SUPPORTEDANCH{\char"2C68\char"0310} \SUPPORTEDANCH{\char"006B\char"0310} \SUPPORTEDFALL{\char"0199\char"0310} \SUPPORTEDANCH{\char"2C6A\char"0310} \SUPPORTEDANCH{\char"006C\char"0310} \SUPPORTEDFALL{\char"0234\char"0310} \SUPPORTEDANCH{\char"006E\char"0310} \SUPPORTEDANCH{\char"0235\char"0310} \SUPPORTEDANCH{\char"0070\char"0310} \SUPPORTEDFALL{\char"01A5\char"0310} \SUPPORTEDANCH{\char"0071\char"0310} \SUPPORTEDANCH{\char"02A0\char"0310} \SUPPORTEDFALL{\char"0074\char"0310} \SUPPORTEDFALL{\char"01AB\char"0310} \SUPPORTEDFALL{\char"01AD\char"0310} \SUPPORTEDFALL{\char"0236\char"0310} \SUPPORTEDANCH{\char"0075\char"0310} \MISSINGGLYPH{\char"AB52\char"0310} \SUPPORTEDANCH{\char"0079\char"0310} \SUPPORTEDANCH{\char"01B4\char"0310} \MISSINGGLYPH{\char"1EFF\char"0310}

\SUPPORTEDANCH{\char"0048\char"0311} \SUPPORTEDFALL{\char"2C67\char"0311} \SUPPORTEDANCH{\char"004B\char"0311} \SUPPORTEDANCH{\char"2C69\char"0311} \SUPPORTEDANCH{\char"004D\char"0311} \MISSINGGLYPH{\char"2C6E\char"0311} \SUPPORTEDANCH{\char"0050\char"0311} \SUPPORTEDANCH{\char"01A4\char"0311} \SUPPORTEDPREC{\char"0052\char"0311} \SUPPORTEDFALL{\char"2C64\char"0311} \SUPPORTEDANCH{\char"0054\char"0311} \SUPPORTEDFALL{\char"01AC\char"0311} \SUPPORTEDFALL{\char"01AE\char"0311} \SUPPORTEDANCH{\char"0059\char"0311} \SUPPORTEDFALL{\char"01B3\char"0311} \SUPPORTEDPREC{\char"0061\char"0311} \SUPPORTEDFALL{\char"1E9A\char"0311} \SUPPORTEDFALL{\char"0062\char"0311} \SUPPORTEDFALL{\char"0183\char"0311} \SUPPORTEDANCH{\char"0063\char"0311} \SUPPORTEDFALL{\char"0188\char"0311} \SUPPORTEDANCH{\char"0064\char"0311} \SUPPORTEDANCH{\char"018C\char"0311} \SUPPORTEDFALL{\char"0221\char"0311} \SUPPORTEDANCH{\char"0068\char"0311} \SUPPORTEDANCH{\char"2C68\char"0311} \SUPPORTEDANCH{\char"006B\char"0311} \SUPPORTEDFALL{\char"0199\char"0311} \SUPPORTEDANCH{\char"2C6A\char"0311} \SUPPORTEDANCH{\char"006C\char"0311} \SUPPORTEDFALL{\char"0234\char"0311} \SUPPORTEDANCH{\char"006E\char"0311} \SUPPORTEDANCH{\char"0235\char"0311} \SUPPORTEDANCH{\char"0070\char"0311} \SUPPORTEDFALL{\char"01A5\char"0311} \SUPPORTEDANCH{\char"0071\char"0311} \SUPPORTEDANCH{\char"02A0\char"0311} \SUPPORTEDFALL{\char"0074\char"0311} \SUPPORTEDFALL{\char"01AB\char"0311} \SUPPORTEDFALL{\char"01AD\char"0311} \SUPPORTEDFALL{\char"0236\char"0311} \SUPPORTEDPREC{\char"0075\char"0311} \MISSINGGLYPH{\char"AB52\char"0311} \SUPPORTEDANCH{\char"0079\char"0311} \SUPPORTEDANCH{\char"01B4\char"0311} \MISSINGGLYPH{\char"1EFF\char"0311}

\SUPPORTEDANCH{\char"0048\char"0312} \SUPPORTEDFALL{\char"2C67\char"0312} \SUPPORTEDANCH{\char"004B\char"0312} \SUPPORTEDANCH{\char"2C69\char"0312} \SUPPORTEDANCH{\char"004D\char"0312} \MISSINGGLYPH{\char"2C6E\char"0312} \SUPPORTEDANCH{\char"0050\char"0312} \SUPPORTEDANCH{\char"01A4\char"0312} \SUPPORTEDANCH{\char"0052\char"0312} \SUPPORTEDFALL{\char"2C64\char"0312} \SUPPORTEDANCH{\char"0054\char"0312} \SUPPORTEDFALL{\char"01AC\char"0312} \SUPPORTEDFALL{\char"01AE\char"0312} \SUPPORTEDANCH{\char"0059\char"0312} \SUPPORTEDFALL{\char"01B3\char"0312} \SUPPORTEDANCH{\char"0061\char"0312} \SUPPORTEDFALL{\char"1E9A\char"0312} \SUPPORTEDFALL{\char"0062\char"0312} \SUPPORTEDFALL{\char"0183\char"0312} \SUPPORTEDANCH{\char"0063\char"0312} \SUPPORTEDFALL{\char"0188\char"0312} \SUPPORTEDANCH{\char"0064\char"0312} \SUPPORTEDANCH{\char"018C\char"0312} \SUPPORTEDFALL{\char"0221\char"0312} \SUPPORTEDANCH{\char"0068\char"0312} \SUPPORTEDANCH{\char"2C68\char"0312} \SUPPORTEDANCH{\char"006B\char"0312} \SUPPORTEDFALL{\char"0199\char"0312} \SUPPORTEDANCH{\char"2C6A\char"0312} \SUPPORTEDANCH{\char"006C\char"0312} \SUPPORTEDFALL{\char"0234\char"0312} \SUPPORTEDANCH{\char"006E\char"0312} \SUPPORTEDANCH{\char"0235\char"0312} \SUPPORTEDANCH{\char"0070\char"0312} \SUPPORTEDFALL{\char"01A5\char"0312} \SUPPORTEDANCH{\char"0071\char"0312} \SUPPORTEDANCH{\char"02A0\char"0312} \SUPPORTEDFALL{\char"0074\char"0312} \SUPPORTEDFALL{\char"01AB\char"0312} \SUPPORTEDFALL{\char"01AD\char"0312} \SUPPORTEDFALL{\char"0236\char"0312} \SUPPORTEDANCH{\char"0075\char"0312} \MISSINGGLYPH{\char"AB52\char"0312} \SUPPORTEDANCH{\char"0079\char"0312} \SUPPORTEDANCH{\char"01B4\char"0312} \MISSINGGLYPH{\char"1EFF\char"0312}

\MISSINGGLYPH{\char"0048} \MISSINGGLYPH{\char"2C67} \MISSINGGLYPH{\char"004B} \MISSINGGLYPH{\char"2C69} \MISSINGGLYPH{\char"004D} \MISSINGGLYPH{\char"2C6E} \MISSINGGLYPH{\char"0050} \MISSINGGLYPH{\char"01A4} \MISSINGGLYPH{\char"0052} \MISSINGGLYPH{\char"2C64} \MISSINGGLYPH{\char"0054} \MISSINGGLYPH{\char"01AC} \MISSINGGLYPH{\char"01AE} \MISSINGGLYPH{\char"0059} \MISSINGGLYPH{\char"01B3} \MISSINGGLYPH{\char"0061} \MISSINGGLYPH{\char"1E9A} \MISSINGGLYPH{\char"0062} \MISSINGGLYPH{\char"0183} \MISSINGGLYPH{\char"0063} \MISSINGGLYPH{\char"0188} \MISSINGGLYPH{\char"0064} \MISSINGGLYPH{\char"018C} \MISSINGGLYPH{\char"0221} \MISSINGGLYPH{\char"0068} \MISSINGGLYPH{\char"2C68} \MISSINGGLYPH{\char"006B} \MISSINGGLYPH{\char"0199} \MISSINGGLYPH{\char"2C6A} \MISSINGGLYPH{\char"006C} \MISSINGGLYPH{\char"0234} \MISSINGGLYPH{\char"006E} \MISSINGGLYPH{\char"0235} \MISSINGGLYPH{\char"0070} \MISSINGGLYPH{\char"01A5} \MISSINGGLYPH{\char"0071} \MISSINGGLYPH{\char"02A0} \MISSINGGLYPH{\char"0074} \MISSINGGLYPH{\char"01AB} \MISSINGGLYPH{\char"01AD} \MISSINGGLYPH{\char"0236} \MISSINGGLYPH{\char"0075} \MISSINGGLYPH{\char"AB52} \MISSINGGLYPH{\char"0079} \MISSINGGLYPH{\char"01B4} \MISSINGGLYPH{\char"1EFF}

\MISSINGGLYPH{\char"0048} \MISSINGGLYPH{\char"2C67} \MISSINGGLYPH{\char"004B} \MISSINGGLYPH{\char"2C69} \MISSINGGLYPH{\char"004D} \MISSINGGLYPH{\char"2C6E} \MISSINGGLYPH{\char"0050} \MISSINGGLYPH{\char"01A4} \MISSINGGLYPH{\char"0052} \MISSINGGLYPH{\char"2C64} \MISSINGGLYPH{\char"0054} \MISSINGGLYPH{\char"01AC} \MISSINGGLYPH{\char"01AE} \MISSINGGLYPH{\char"0059} \MISSINGGLYPH{\char"01B3} \MISSINGGLYPH{\char"0061} \MISSINGGLYPH{\char"1E9A} \MISSINGGLYPH{\char"0062} \MISSINGGLYPH{\char"0183} \MISSINGGLYPH{\char"0063} \MISSINGGLYPH{\char"0188} \MISSINGGLYPH{\char"0064} \MISSINGGLYPH{\char"018C} \MISSINGGLYPH{\char"0221} \MISSINGGLYPH{\char"0068} \MISSINGGLYPH{\char"2C68} \MISSINGGLYPH{\char"006B} \MISSINGGLYPH{\char"0199} \MISSINGGLYPH{\char"2C6A} \MISSINGGLYPH{\char"006C} \MISSINGGLYPH{\char"0234} \MISSINGGLYPH{\char"006E} \MISSINGGLYPH{\char"0235} \MISSINGGLYPH{\char"0070} \MISSINGGLYPH{\char"01A5} \MISSINGGLYPH{\char"0071} \MISSINGGLYPH{\char"02A0} \MISSINGGLYPH{\char"0074} \MISSINGGLYPH{\char"01AB} \MISSINGGLYPH{\char"01AD} \MISSINGGLYPH{\char"0236} \MISSINGGLYPH{\char"0075} \MISSINGGLYPH{\char"AB52} \MISSINGGLYPH{\char"0079} \MISSINGGLYPH{\char"01B4} \MISSINGGLYPH{\char"1EFF}

\MISSINGGLYPH{\char"0048} \MISSINGGLYPH{\char"2C67} \MISSINGGLYPH{\char"004B} \MISSINGGLYPH{\char"2C69} \MISSINGGLYPH{\char"004D} \MISSINGGLYPH{\char"2C6E} \MISSINGGLYPH{\char"0050} \MISSINGGLYPH{\char"01A4} \MISSINGGLYPH{\char"0052} \MISSINGGLYPH{\char"2C64} \MISSINGGLYPH{\char"0054} \MISSINGGLYPH{\char"01AC} \MISSINGGLYPH{\char"01AE} \MISSINGGLYPH{\char"0059} \MISSINGGLYPH{\char"01B3} \MISSINGGLYPH{\char"0061} \MISSINGGLYPH{\char"1E9A} \MISSINGGLYPH{\char"0062} \MISSINGGLYPH{\char"0183} \MISSINGGLYPH{\char"0063} \MISSINGGLYPH{\char"0188} \MISSINGGLYPH{\char"0064} \MISSINGGLYPH{\char"018C} \MISSINGGLYPH{\char"0221} \MISSINGGLYPH{\char"0068} \MISSINGGLYPH{\char"2C68} \MISSINGGLYPH{\char"006B} \MISSINGGLYPH{\char"0199} \MISSINGGLYPH{\char"2C6A} \MISSINGGLYPH{\char"006C} \MISSINGGLYPH{\char"0234} \MISSINGGLYPH{\char"006E} \MISSINGGLYPH{\char"0235} \MISSINGGLYPH{\char"0070} \MISSINGGLYPH{\char"01A5} \MISSINGGLYPH{\char"0071} \MISSINGGLYPH{\char"02A0} \MISSINGGLYPH{\char"0074} \MISSINGGLYPH{\char"01AB} \MISSINGGLYPH{\char"01AD} \MISSINGGLYPH{\char"0236} \MISSINGGLYPH{\char"0075} \MISSINGGLYPH{\char"AB52} \MISSINGGLYPH{\char"0079} \MISSINGGLYPH{\char"01B4} \MISSINGGLYPH{\char"1EFF}

\SUPPORTEDANCH{\char"0048\char"030E} \SUPPORTEDFALL{\char"2C67\char"030E} \SUPPORTEDANCH{\char"004B\char"030E} \SUPPORTEDANCH{\char"2C69\char"030E} \SUPPORTEDANCH{\char"004D\char"030E} \MISSINGGLYPH{\char"2C6E\char"030E} \SUPPORTEDANCH{\char"0050\char"030E} \SUPPORTEDANCH{\char"01A4\char"030E} \SUPPORTEDANCH{\char"0052\char"030E} \SUPPORTEDFALL{\char"2C64\char"030E} \SUPPORTEDANCH{\char"0054\char"030E} \SUPPORTEDFALL{\char"01AC\char"030E} \SUPPORTEDFALL{\char"01AE\char"030E} \SUPPORTEDANCH{\char"0059\char"030E} \SUPPORTEDFALL{\char"01B3\char"030E} \SUPPORTEDANCH{\char"0061\char"030E} \SUPPORTEDFALL{\char"1E9A\char"030E} \SUPPORTEDFALL{\char"0062\char"030E} \SUPPORTEDFALL{\char"0183\char"030E} \SUPPORTEDANCH{\char"0063\char"030E} \SUPPORTEDFALL{\char"0188\char"030E} \SUPPORTEDANCH{\char"0064\char"030E} \SUPPORTEDANCH{\char"018C\char"030E} \SUPPORTEDFALL{\char"0221\char"030E} \SUPPORTEDANCH{\char"0068\char"030E} \SUPPORTEDANCH{\char"2C68\char"030E} \SUPPORTEDANCH{\char"006B\char"030E} \SUPPORTEDFALL{\char"0199\char"030E} \SUPPORTEDANCH{\char"2C6A\char"030E} \SUPPORTEDANCH{\char"006C\char"030E} \SUPPORTEDFALL{\char"0234\char"030E} \SUPPORTEDANCH{\char"006E\char"030E} \SUPPORTEDANCH{\char"0235\char"030E} \SUPPORTEDANCH{\char"0070\char"030E} \SUPPORTEDFALL{\char"01A5\char"030E} \SUPPORTEDANCH{\char"0071\char"030E} \SUPPORTEDANCH{\char"02A0\char"030E} \SUPPORTEDFALL{\char"0074\char"030E} \SUPPORTEDFALL{\char"01AB\char"030E} \SUPPORTEDFALL{\char"01AD\char"030E} \SUPPORTEDFALL{\char"0236\char"030E} \SUPPORTEDANCH{\char"0075\char"030E} \MISSINGGLYPH{\char"AB52\char"030E} \SUPPORTEDANCH{\char"0079\char"030E} \SUPPORTEDANCH{\char"01B4\char"030E} \MISSINGGLYPH{\char"1EFF\char"030E}

\MISSINGGLYPH{\char"0048} \MISSINGGLYPH{\char"2C67} \MISSINGGLYPH{\char"004B} \MISSINGGLYPH{\char"2C69} \MISSINGGLYPH{\char"004D} \MISSINGGLYPH{\char"2C6E} \MISSINGGLYPH{\char"0050} \MISSINGGLYPH{\char"01A4} \MISSINGGLYPH{\char"0052} \MISSINGGLYPH{\char"2C64} \MISSINGGLYPH{\char"0054} \MISSINGGLYPH{\char"01AC} \MISSINGGLYPH{\char"01AE} \MISSINGGLYPH{\char"0059} \MISSINGGLYPH{\char"01B3} \MISSINGGLYPH{\char"0061} \MISSINGGLYPH{\char"1E9A} \MISSINGGLYPH{\char"0062} \MISSINGGLYPH{\char"0183} \MISSINGGLYPH{\char"0063} \MISSINGGLYPH{\char"0188} \MISSINGGLYPH{\char"0064} \MISSINGGLYPH{\char"018C} \MISSINGGLYPH{\char"0221} \MISSINGGLYPH{\char"0068} \MISSINGGLYPH{\char"2C68} \MISSINGGLYPH{\char"006B} \MISSINGGLYPH{\char"0199} \MISSINGGLYPH{\char"2C6A} \MISSINGGLYPH{\char"006C} \MISSINGGLYPH{\char"0234} \MISSINGGLYPH{\char"006E} \MISSINGGLYPH{\char"0235} \MISSINGGLYPH{\char"0070} \MISSINGGLYPH{\char"01A5} \MISSINGGLYPH{\char"0071} \MISSINGGLYPH{\char"02A0} \MISSINGGLYPH{\char"0074} \MISSINGGLYPH{\char"01AB} \MISSINGGLYPH{\char"01AD} \MISSINGGLYPH{\char"0236} \MISSINGGLYPH{\char"0075} \MISSINGGLYPH{\char"AB52} \MISSINGGLYPH{\char"0079} \MISSINGGLYPH{\char"01B4} \MISSINGGLYPH{\char"1EFF}

\MISSINGGLYPH{\char"0048} \MISSINGGLYPH{\char"2C67} \MISSINGGLYPH{\char"004B} \MISSINGGLYPH{\char"2C69} \MISSINGGLYPH{\char"004D} \MISSINGGLYPH{\char"2C6E} \MISSINGGLYPH{\char"0050} \MISSINGGLYPH{\char"01A4} \MISSINGGLYPH{\char"0052} \MISSINGGLYPH{\char"2C64} \MISSINGGLYPH{\char"0054} \MISSINGGLYPH{\char"01AC} \MISSINGGLYPH{\char"01AE} \MISSINGGLYPH{\char"0059} \MISSINGGLYPH{\char"01B3} \MISSINGGLYPH{\char"0061} \MISSINGGLYPH{\char"1E9A} \MISSINGGLYPH{\char"0062} \MISSINGGLYPH{\char"0183} \MISSINGGLYPH{\char"0063} \MISSINGGLYPH{\char"0188} \MISSINGGLYPH{\char"0064} \MISSINGGLYPH{\char"018C} \MISSINGGLYPH{\char"0221} \MISSINGGLYPH{\char"0068} \MISSINGGLYPH{\char"2C68} \MISSINGGLYPH{\char"006B} \MISSINGGLYPH{\char"0199} \MISSINGGLYPH{\char"2C6A} \MISSINGGLYPH{\char"006C} \MISSINGGLYPH{\char"0234} \MISSINGGLYPH{\char"006E} \MISSINGGLYPH{\char"0235} \MISSINGGLYPH{\char"0070} \MISSINGGLYPH{\char"01A5} \MISSINGGLYPH{\char"0071} \MISSINGGLYPH{\char"02A0} \MISSINGGLYPH{\char"0074} \MISSINGGLYPH{\char"01AB} \MISSINGGLYPH{\char"01AD} \MISSINGGLYPH{\char"0236} \MISSINGGLYPH{\char"0075} \MISSINGGLYPH{\char"AB52} \MISSINGGLYPH{\char"0079} \MISSINGGLYPH{\char"01B4} \MISSINGGLYPH{\char"1EFF}

\MISSINGGLYPH{\char"0048} \MISSINGGLYPH{\char"2C67} \MISSINGGLYPH{\char"004B} \MISSINGGLYPH{\char"2C69} \MISSINGGLYPH{\char"004D} \MISSINGGLYPH{\char"2C6E} \MISSINGGLYPH{\char"0050} \MISSINGGLYPH{\char"01A4} \MISSINGGLYPH{\char"0052} \MISSINGGLYPH{\char"2C64} \MISSINGGLYPH{\char"0054} \MISSINGGLYPH{\char"01AC} \MISSINGGLYPH{\char"01AE} \MISSINGGLYPH{\char"0059} \MISSINGGLYPH{\char"01B3} \MISSINGGLYPH{\char"0061} \MISSINGGLYPH{\char"1E9A} \MISSINGGLYPH{\char"0062} \MISSINGGLYPH{\char"0183} \MISSINGGLYPH{\char"0063} \MISSINGGLYPH{\char"0188} \MISSINGGLYPH{\char"0064} \MISSINGGLYPH{\char"018C} \MISSINGGLYPH{\char"0221} \MISSINGGLYPH{\char"0068} \MISSINGGLYPH{\char"2C68} \MISSINGGLYPH{\char"006B} \MISSINGGLYPH{\char"0199} \MISSINGGLYPH{\char"2C6A} \MISSINGGLYPH{\char"006C} \MISSINGGLYPH{\char"0234} \MISSINGGLYPH{\char"006E} \MISSINGGLYPH{\char"0235} \MISSINGGLYPH{\char"0070} \MISSINGGLYPH{\char"01A5} \MISSINGGLYPH{\char"0071} \MISSINGGLYPH{\char"02A0} \MISSINGGLYPH{\char"0074} \MISSINGGLYPH{\char"01AB} \MISSINGGLYPH{\char"01AD} \MISSINGGLYPH{\char"0236} \MISSINGGLYPH{\char"0075} \MISSINGGLYPH{\char"AB52} \MISSINGGLYPH{\char"0079} \MISSINGGLYPH{\char"01B4} \MISSINGGLYPH{\char"1EFF}

\MISSINGGLYPH{\char"0048} \MISSINGGLYPH{\char"2C67} \MISSINGGLYPH{\char"004B} \MISSINGGLYPH{\char"2C69} \MISSINGGLYPH{\char"004D} \MISSINGGLYPH{\char"2C6E} \MISSINGGLYPH{\char"0050} \MISSINGGLYPH{\char"01A4} \MISSINGGLYPH{\char"0052} \MISSINGGLYPH{\char"2C64} \MISSINGGLYPH{\char"0054} \MISSINGGLYPH{\char"01AC} \MISSINGGLYPH{\char"01AE} \MISSINGGLYPH{\char"0059} \MISSINGGLYPH{\char"01B3} \MISSINGGLYPH{\char"0061} \MISSINGGLYPH{\char"1E9A} \MISSINGGLYPH{\char"0062} \MISSINGGLYPH{\char"0183} \MISSINGGLYPH{\char"0063} \MISSINGGLYPH{\char"0188} \MISSINGGLYPH{\char"0064} \MISSINGGLYPH{\char"018C} \MISSINGGLYPH{\char"0221} \MISSINGGLYPH{\char"0068} \MISSINGGLYPH{\char"2C68} \MISSINGGLYPH{\char"006B} \MISSINGGLYPH{\char"0199} \MISSINGGLYPH{\char"2C6A} \MISSINGGLYPH{\char"006C} \MISSINGGLYPH{\char"0234} \MISSINGGLYPH{\char"006E} \MISSINGGLYPH{\char"0235} \MISSINGGLYPH{\char"0070} \MISSINGGLYPH{\char"01A5} \MISSINGGLYPH{\char"0071} \MISSINGGLYPH{\char"02A0} \MISSINGGLYPH{\char"0074} \MISSINGGLYPH{\char"01AB} \MISSINGGLYPH{\char"01AD} \MISSINGGLYPH{\char"0236} \MISSINGGLYPH{\char"0075} \MISSINGGLYPH{\char"AB52} \MISSINGGLYPH{\char"0079} \MISSINGGLYPH{\char"01B4} \MISSINGGLYPH{\char"1EFF}

\MISSINGGLYPH{\char"0048} \MISSINGGLYPH{\char"2C67} \MISSINGGLYPH{\char"004B} \MISSINGGLYPH{\char"2C69} \MISSINGGLYPH{\char"004D} \MISSINGGLYPH{\char"2C6E} \MISSINGGLYPH{\char"0050} \MISSINGGLYPH{\char"01A4} \MISSINGGLYPH{\char"0052} \MISSINGGLYPH{\char"2C64} \MISSINGGLYPH{\char"0054} \MISSINGGLYPH{\char"01AC} \MISSINGGLYPH{\char"01AE} \MISSINGGLYPH{\char"0059} \MISSINGGLYPH{\char"01B3} \MISSINGGLYPH{\char"0061} \MISSINGGLYPH{\char"1E9A} \MISSINGGLYPH{\char"0062} \MISSINGGLYPH{\char"0183} \MISSINGGLYPH{\char"0063} \MISSINGGLYPH{\char"0188} \MISSINGGLYPH{\char"0064} \MISSINGGLYPH{\char"018C} \MISSINGGLYPH{\char"0221} \MISSINGGLYPH{\char"0068} \MISSINGGLYPH{\char"2C68} \MISSINGGLYPH{\char"006B} \MISSINGGLYPH{\char"0199} \MISSINGGLYPH{\char"2C6A} \MISSINGGLYPH{\char"006C} \MISSINGGLYPH{\char"0234} \MISSINGGLYPH{\char"006E} \MISSINGGLYPH{\char"0235} \MISSINGGLYPH{\char"0070} \MISSINGGLYPH{\char"01A5} \MISSINGGLYPH{\char"0071} \MISSINGGLYPH{\char"02A0} \MISSINGGLYPH{\char"0074} \MISSINGGLYPH{\char"01AB} \MISSINGGLYPH{\char"01AD} \MISSINGGLYPH{\char"0236} \MISSINGGLYPH{\char"0075} \MISSINGGLYPH{\char"AB52} \MISSINGGLYPH{\char"0079} \MISSINGGLYPH{\char"01B4} \MISSINGGLYPH{\char"1EFF}

\MISSINGGLYPH{\char"0048} \MISSINGGLYPH{\char"2C67} \MISSINGGLYPH{\char"004B} \MISSINGGLYPH{\char"2C69} \MISSINGGLYPH{\char"004D} \MISSINGGLYPH{\char"2C6E} \MISSINGGLYPH{\char"0050} \MISSINGGLYPH{\char"01A4} \MISSINGGLYPH{\char"0052} \MISSINGGLYPH{\char"2C64} \MISSINGGLYPH{\char"0054} \MISSINGGLYPH{\char"01AC} \MISSINGGLYPH{\char"01AE} \MISSINGGLYPH{\char"0059} \MISSINGGLYPH{\char"01B3} \MISSINGGLYPH{\char"0061} \MISSINGGLYPH{\char"1E9A} \MISSINGGLYPH{\char"0062} \MISSINGGLYPH{\char"0183} \MISSINGGLYPH{\char"0063} \MISSINGGLYPH{\char"0188} \MISSINGGLYPH{\char"0064} \MISSINGGLYPH{\char"018C} \MISSINGGLYPH{\char"0221} \MISSINGGLYPH{\char"0068} \MISSINGGLYPH{\char"2C68} \MISSINGGLYPH{\char"006B} \MISSINGGLYPH{\char"0199} \MISSINGGLYPH{\char"2C6A} \MISSINGGLYPH{\char"006C} \MISSINGGLYPH{\char"0234} \MISSINGGLYPH{\char"006E} \MISSINGGLYPH{\char"0235} \MISSINGGLYPH{\char"0070} \MISSINGGLYPH{\char"01A5} \MISSINGGLYPH{\char"0071} \MISSINGGLYPH{\char"02A0} \MISSINGGLYPH{\char"0074} \MISSINGGLYPH{\char"01AB} \MISSINGGLYPH{\char"01AD} \MISSINGGLYPH{\char"0236} \MISSINGGLYPH{\char"0075} \MISSINGGLYPH{\char"AB52} \MISSINGGLYPH{\char"0079} \MISSINGGLYPH{\char"01B4} \MISSINGGLYPH{\char"1EFF}

\MISSINGGLYPH{\char"0048} \MISSINGGLYPH{\char"2C67} \MISSINGGLYPH{\char"004B} \MISSINGGLYPH{\char"2C69} \MISSINGGLYPH{\char"004D} \MISSINGGLYPH{\char"2C6E} \MISSINGGLYPH{\char"0050} \MISSINGGLYPH{\char"01A4} \MISSINGGLYPH{\char"0052} \MISSINGGLYPH{\char"2C64} \MISSINGGLYPH{\char"0054} \MISSINGGLYPH{\char"01AC} \MISSINGGLYPH{\char"01AE} \MISSINGGLYPH{\char"0059} \MISSINGGLYPH{\char"01B3} \MISSINGGLYPH{\char"0061} \MISSINGGLYPH{\char"1E9A} \MISSINGGLYPH{\char"0062} \MISSINGGLYPH{\char"0183} \MISSINGGLYPH{\char"0063} \MISSINGGLYPH{\char"0188} \MISSINGGLYPH{\char"0064} \MISSINGGLYPH{\char"018C} \MISSINGGLYPH{\char"0221} \MISSINGGLYPH{\char"0068} \MISSINGGLYPH{\char"2C68} \MISSINGGLYPH{\char"006B} \MISSINGGLYPH{\char"0199} \MISSINGGLYPH{\char"2C6A} \MISSINGGLYPH{\char"006C} \MISSINGGLYPH{\char"0234} \MISSINGGLYPH{\char"006E} \MISSINGGLYPH{\char"0235} \MISSINGGLYPH{\char"0070} \MISSINGGLYPH{\char"01A5} \MISSINGGLYPH{\char"0071} \MISSINGGLYPH{\char"02A0} \MISSINGGLYPH{\char"0074} \MISSINGGLYPH{\char"01AB} \MISSINGGLYPH{\char"01AD} \MISSINGGLYPH{\char"0236} \MISSINGGLYPH{\char"0075} \MISSINGGLYPH{\char"AB52} \MISSINGGLYPH{\char"0079} \MISSINGGLYPH{\char"01B4} \MISSINGGLYPH{\char"1EFF}

\MISSINGGLYPH{\char"0048} \MISSINGGLYPH{\char"2C67} \MISSINGGLYPH{\char"004B} \MISSINGGLYPH{\char"2C69} \MISSINGGLYPH{\char"004D} \MISSINGGLYPH{\char"2C6E} \MISSINGGLYPH{\char"0050} \MISSINGGLYPH{\char"01A4} \MISSINGGLYPH{\char"0052} \MISSINGGLYPH{\char"2C64} \MISSINGGLYPH{\char"0054} \MISSINGGLYPH{\char"01AC} \MISSINGGLYPH{\char"01AE} \MISSINGGLYPH{\char"0059} \MISSINGGLYPH{\char"01B3} \MISSINGGLYPH{\char"0061} \MISSINGGLYPH{\char"1E9A} \MISSINGGLYPH{\char"0062} \MISSINGGLYPH{\char"0183} \MISSINGGLYPH{\char"0063} \MISSINGGLYPH{\char"0188} \MISSINGGLYPH{\char"0064} \MISSINGGLYPH{\char"018C} \MISSINGGLYPH{\char"0221} \MISSINGGLYPH{\char"0068} \MISSINGGLYPH{\char"2C68} \MISSINGGLYPH{\char"006B} \MISSINGGLYPH{\char"0199} \MISSINGGLYPH{\char"2C6A} \MISSINGGLYPH{\char"006C} \MISSINGGLYPH{\char"0234} \MISSINGGLYPH{\char"006E} \MISSINGGLYPH{\char"0235} \MISSINGGLYPH{\char"0070} \MISSINGGLYPH{\char"01A5} \MISSINGGLYPH{\char"0071} \MISSINGGLYPH{\char"02A0} \MISSINGGLYPH{\char"0074} \MISSINGGLYPH{\char"01AB} \MISSINGGLYPH{\char"01AD} \MISSINGGLYPH{\char"0236} \MISSINGGLYPH{\char"0075} \MISSINGGLYPH{\char"AB52} \MISSINGGLYPH{\char"0079} \MISSINGGLYPH{\char"01B4} \MISSINGGLYPH{\char"1EFF}

\MISSINGGLYPH{\char"0048} \MISSINGGLYPH{\char"2C67} \MISSINGGLYPH{\char"004B} \MISSINGGLYPH{\char"2C69} \MISSINGGLYPH{\char"004D} \MISSINGGLYPH{\char"2C6E} \MISSINGGLYPH{\char"0050} \MISSINGGLYPH{\char"01A4} \MISSINGGLYPH{\char"0052} \MISSINGGLYPH{\char"2C64} \MISSINGGLYPH{\char"0054} \MISSINGGLYPH{\char"01AC} \MISSINGGLYPH{\char"01AE} \MISSINGGLYPH{\char"0059} \MISSINGGLYPH{\char"01B3} \MISSINGGLYPH{\char"0061} \MISSINGGLYPH{\char"1E9A} \MISSINGGLYPH{\char"0062} \MISSINGGLYPH{\char"0183} \MISSINGGLYPH{\char"0063} \MISSINGGLYPH{\char"0188} \MISSINGGLYPH{\char"0064} \MISSINGGLYPH{\char"018C} \MISSINGGLYPH{\char"0221} \MISSINGGLYPH{\char"0068} \MISSINGGLYPH{\char"2C68} \MISSINGGLYPH{\char"006B} \MISSINGGLYPH{\char"0199} \MISSINGGLYPH{\char"2C6A} \MISSINGGLYPH{\char"006C} \MISSINGGLYPH{\char"0234} \MISSINGGLYPH{\char"006E} \MISSINGGLYPH{\char"0235} \MISSINGGLYPH{\char"0070} \MISSINGGLYPH{\char"01A5} \MISSINGGLYPH{\char"0071} \MISSINGGLYPH{\char"02A0} \MISSINGGLYPH{\char"0074} \MISSINGGLYPH{\char"01AB} \MISSINGGLYPH{\char"01AD} \MISSINGGLYPH{\char"0236} \MISSINGGLYPH{\char"0075} \MISSINGGLYPH{\char"AB52} \MISSINGGLYPH{\char"0079} \MISSINGGLYPH{\char"01B4} \MISSINGGLYPH{\char"1EFF}

}

\newpage
\subsection*{Semibold italic patched, Group 4, above, Above}

{\LibertinusSerifPatch\bfseries\itshape
% grid. columns are bases, rows are marks.
\SUPPORTEDANCH{\char"0048\char"0300} \SUPPORTEDANCH{\char"2C67\char"0300} \SUPPORTEDANCH{\char"004B\char"0300} \SUPPORTEDANCH{\char"2C69\char"0300} \SUPPORTEDANCH{\char"004D\char"0300} \MISSINGGLYPH{\char"2C6E\char"0300} \SUPPORTEDANCH{\char"0050\char"0300} \SUPPORTEDANCH{\char"01A4\char"0300} \SUPPORTEDANCH{\char"0052\char"0300} \SUPPORTEDANCH{\char"2C64\char"0300} \SUPPORTEDANCH{\char"0054\char"0300} \SUPPORTEDANCH{\char"01AC\char"0300} \SUPPORTEDANCH{\char"01AE\char"0300} \SUPPORTEDPREC{\char"0059\char"0300} \SUPPORTEDANCH{\char"01B3\char"0300} \SUPPORTEDPREC{\char"0061\char"0300} \SUPPORTEDANCH{\char"1E9A\char"0300} \SUPPORTEDANCH{\char"0062\char"0300} \SUPPORTEDANCH{\char"0183\char"0300} \SUPPORTEDANCH{\char"0063\char"0300} \SUPPORTEDANCH{\char"0188\char"0300} \SUPPORTEDANCH{\char"0064\char"0300} \SUPPORTEDANCH{\char"018C\char"0300} \SUPPORTEDANCH{\char"0221\char"0300} \SUPPORTEDANCH{\char"0068\char"0300} \SUPPORTEDANCH{\char"2C68\char"0300} \SUPPORTEDANCH{\char"006B\char"0300} \SUPPORTEDANCH{\char"0199\char"0300} \SUPPORTEDANCH{\char"2C6A\char"0300} \SUPPORTEDANCH{\char"006C\char"0300} \SUPPORTEDANCH{\char"0234\char"0300} \SUPPORTEDPREC{\char"006E\char"0300} \SUPPORTEDANCH{\char"0235\char"0300} \SUPPORTEDANCH{\char"0070\char"0300} \SUPPORTEDANCH{\char"01A5\char"0300} \SUPPORTEDANCH{\char"0071\char"0300} \SUPPORTEDANCH{\char"02A0\char"0300} \SUPPORTEDANCH{\char"0074\char"0300} \SUPPORTEDANCH{\char"01AB\char"0300} \SUPPORTEDANCH{\char"01AD\char"0300} \SUPPORTEDANCH{\char"0236\char"0300} \SUPPORTEDPREC{\char"0075\char"0300} \MISSINGGLYPH{\char"AB52\char"0300} \SUPPORTEDPREC{\char"0079\char"0300} \SUPPORTEDANCH{\char"01B4\char"0300} \MISSINGGLYPH{\char"1EFF\char"0300}

\SUPPORTEDANCH{\char"0048\char"0301} \SUPPORTEDANCH{\char"2C67\char"0301} \SUPPORTEDPREC{\char"004B\char"0301} \SUPPORTEDANCH{\char"2C69\char"0301} \SUPPORTEDPREC{\char"004D\char"0301} \MISSINGGLYPH{\char"2C6E\char"0301} \SUPPORTEDPREC{\char"0050\char"0301} \SUPPORTEDANCH{\char"01A4\char"0301} \SUPPORTEDPREC{\char"0052\char"0301} \SUPPORTEDANCH{\char"2C64\char"0301} \SUPPORTEDANCH{\char"0054\char"0301} \SUPPORTEDANCH{\char"01AC\char"0301} \SUPPORTEDANCH{\char"01AE\char"0301} \SUPPORTEDPREC{\char"0059\char"0301} \SUPPORTEDANCH{\char"01B3\char"0301} \SUPPORTEDPREC{\char"0061\char"0301} \SUPPORTEDANCH{\char"1E9A\char"0301} \SUPPORTEDANCH{\char"0062\char"0301} \SUPPORTEDANCH{\char"0183\char"0301} \SUPPORTEDPREC{\char"0063\char"0301} \SUPPORTEDANCH{\char"0188\char"0301} \SUPPORTEDANCH{\char"0064\char"0301} \SUPPORTEDANCH{\char"018C\char"0301} \SUPPORTEDANCH{\char"0221\char"0301} \SUPPORTEDANCH{\char"0068\char"0301} \SUPPORTEDANCH{\char"2C68\char"0301} \SUPPORTEDPREC{\char"006B\char"0301} \SUPPORTEDANCH{\char"0199\char"0301} \SUPPORTEDANCH{\char"2C6A\char"0301} \SUPPORTEDPREC{\char"006C\char"0301} \SUPPORTEDANCH{\char"0234\char"0301} \SUPPORTEDPREC{\char"006E\char"0301} \SUPPORTEDANCH{\char"0235\char"0301} \SUPPORTEDPREC{\char"0070\char"0301} \SUPPORTEDANCH{\char"01A5\char"0301} \SUPPORTEDANCH{\char"0071\char"0301} \SUPPORTEDANCH{\char"02A0\char"0301} \SUPPORTEDANCH{\char"0074\char"0301} \SUPPORTEDANCH{\char"01AB\char"0301} \SUPPORTEDANCH{\char"01AD\char"0301} \SUPPORTEDANCH{\char"0236\char"0301} \SUPPORTEDPREC{\char"0075\char"0301} \MISSINGGLYPH{\char"AB52\char"0301} \SUPPORTEDPREC{\char"0079\char"0301} \SUPPORTEDANCH{\char"01B4\char"0301} \MISSINGGLYPH{\char"1EFF\char"0301}

\SUPPORTEDPREC{\char"0048\char"0302} \SUPPORTEDANCH{\char"2C67\char"0302} \SUPPORTEDANCH{\char"004B\char"0302} \SUPPORTEDANCH{\char"2C69\char"0302} \SUPPORTEDANCH{\char"004D\char"0302} \MISSINGGLYPH{\char"2C6E\char"0302} \SUPPORTEDANCH{\char"0050\char"0302} \SUPPORTEDANCH{\char"01A4\char"0302} \SUPPORTEDANCH{\char"0052\char"0302} \SUPPORTEDANCH{\char"2C64\char"0302} \SUPPORTEDANCH{\char"0054\char"0302} \SUPPORTEDANCH{\char"01AC\char"0302} \SUPPORTEDANCH{\char"01AE\char"0302} \SUPPORTEDPREC{\char"0059\char"0302} \SUPPORTEDANCH{\char"01B3\char"0302} \SUPPORTEDPREC{\char"0061\char"0302} \SUPPORTEDANCH{\char"1E9A\char"0302} \SUPPORTEDANCH{\char"0062\char"0302} \SUPPORTEDANCH{\char"0183\char"0302} \SUPPORTEDPREC{\char"0063\char"0302} \SUPPORTEDANCH{\char"0188\char"0302} \SUPPORTEDANCH{\char"0064\char"0302} \SUPPORTEDANCH{\char"018C\char"0302} \SUPPORTEDANCH{\char"0221\char"0302} \SUPPORTEDPREC{\char"0068\char"0302} \SUPPORTEDANCH{\char"2C68\char"0302} \SUPPORTEDANCH{\char"006B\char"0302} \SUPPORTEDANCH{\char"0199\char"0302} \SUPPORTEDANCH{\char"2C6A\char"0302} \SUPPORTEDANCH{\char"006C\char"0302} \SUPPORTEDANCH{\char"0234\char"0302} \SUPPORTEDANCH{\char"006E\char"0302} \SUPPORTEDANCH{\char"0235\char"0302} \SUPPORTEDANCH{\char"0070\char"0302} \SUPPORTEDANCH{\char"01A5\char"0302} \SUPPORTEDANCH{\char"0071\char"0302} \SUPPORTEDANCH{\char"02A0\char"0302} \SUPPORTEDANCH{\char"0074\char"0302} \SUPPORTEDANCH{\char"01AB\char"0302} \SUPPORTEDANCH{\char"01AD\char"0302} \SUPPORTEDANCH{\char"0236\char"0302} \SUPPORTEDPREC{\char"0075\char"0302} \MISSINGGLYPH{\char"AB52\char"0302} \SUPPORTEDPREC{\char"0079\char"0302} \SUPPORTEDANCH{\char"01B4\char"0302} \MISSINGGLYPH{\char"1EFF\char"0302}

\SUPPORTEDANCH{\char"0048\char"0303} \SUPPORTEDANCH{\char"2C67\char"0303} \SUPPORTEDANCH{\char"004B\char"0303} \SUPPORTEDANCH{\char"2C69\char"0303} \SUPPORTEDANCH{\char"004D\char"0303} \MISSINGGLYPH{\char"2C6E\char"0303} \SUPPORTEDANCH{\char"0050\char"0303} \SUPPORTEDANCH{\char"01A4\char"0303} \SUPPORTEDANCH{\char"0052\char"0303} \SUPPORTEDANCH{\char"2C64\char"0303} \SUPPORTEDANCH{\char"0054\char"0303} \SUPPORTEDANCH{\char"01AC\char"0303} \SUPPORTEDANCH{\char"01AE\char"0303} \SUPPORTEDPREC{\char"0059\char"0303} \SUPPORTEDANCH{\char"01B3\char"0303} \SUPPORTEDPREC{\char"0061\char"0303} \SUPPORTEDANCH{\char"1E9A\char"0303} \SUPPORTEDANCH{\char"0062\char"0303} \SUPPORTEDANCH{\char"0183\char"0303} \SUPPORTEDANCH{\char"0063\char"0303} \SUPPORTEDANCH{\char"0188\char"0303} \SUPPORTEDANCH{\char"0064\char"0303} \SUPPORTEDANCH{\char"018C\char"0303} \SUPPORTEDANCH{\char"0221\char"0303} \SUPPORTEDANCH{\char"0068\char"0303} \SUPPORTEDANCH{\char"2C68\char"0303} \SUPPORTEDANCH{\char"006B\char"0303} \SUPPORTEDANCH{\char"0199\char"0303} \SUPPORTEDANCH{\char"2C6A\char"0303} \SUPPORTEDANCH{\char"006C\char"0303} \SUPPORTEDANCH{\char"0234\char"0303} \SUPPORTEDPREC{\char"006E\char"0303} \SUPPORTEDANCH{\char"0235\char"0303} \SUPPORTEDANCH{\char"0070\char"0303} \SUPPORTEDANCH{\char"01A5\char"0303} \SUPPORTEDANCH{\char"0071\char"0303} \SUPPORTEDANCH{\char"02A0\char"0303} \SUPPORTEDANCH{\char"0074\char"0303} \SUPPORTEDANCH{\char"01AB\char"0303} \SUPPORTEDANCH{\char"01AD\char"0303} \SUPPORTEDANCH{\char"0236\char"0303} \SUPPORTEDPREC{\char"0075\char"0303} \MISSINGGLYPH{\char"AB52\char"0303} \SUPPORTEDPREC{\char"0079\char"0303} \SUPPORTEDANCH{\char"01B4\char"0303} \MISSINGGLYPH{\char"1EFF\char"0303}

\SUPPORTEDANCH{\char"0048\char"0304} \SUPPORTEDANCH{\char"2C67\char"0304} \SUPPORTEDANCH{\char"004B\char"0304} \SUPPORTEDANCH{\char"2C69\char"0304} \SUPPORTEDANCH{\char"004D\char"0304} \MISSINGGLYPH{\char"2C6E\char"0304} \SUPPORTEDANCH{\char"0050\char"0304} \SUPPORTEDANCH{\char"01A4\char"0304} \SUPPORTEDANCH{\char"0052\char"0304} \SUPPORTEDANCH{\char"2C64\char"0304} \SUPPORTEDANCH{\char"0054\char"0304} \SUPPORTEDANCH{\char"01AC\char"0304} \SUPPORTEDANCH{\char"01AE\char"0304} \SUPPORTEDPREC{\char"0059\char"0304} \SUPPORTEDANCH{\char"01B3\char"0304} \SUPPORTEDPREC{\char"0061\char"0304} \SUPPORTEDANCH{\char"1E9A\char"0304} \SUPPORTEDANCH{\char"0062\char"0304} \SUPPORTEDANCH{\char"0183\char"0304} \SUPPORTEDANCH{\char"0063\char"0304} \SUPPORTEDANCH{\char"0188\char"0304} \SUPPORTEDANCH{\char"0064\char"0304} \SUPPORTEDANCH{\char"018C\char"0304} \SUPPORTEDANCH{\char"0221\char"0304} \SUPPORTEDANCH{\char"0068\char"0304} \SUPPORTEDANCH{\char"2C68\char"0304} \SUPPORTEDANCH{\char"006B\char"0304} \SUPPORTEDANCH{\char"0199\char"0304} \SUPPORTEDANCH{\char"2C6A\char"0304} \SUPPORTEDANCH{\char"006C\char"0304} \SUPPORTEDANCH{\char"0234\char"0304} \SUPPORTEDANCH{\char"006E\char"0304} \SUPPORTEDANCH{\char"0235\char"0304} \SUPPORTEDANCH{\char"0070\char"0304} \SUPPORTEDANCH{\char"01A5\char"0304} \SUPPORTEDANCH{\char"0071\char"0304} \SUPPORTEDANCH{\char"02A0\char"0304} \SUPPORTEDANCH{\char"0074\char"0304} \SUPPORTEDANCH{\char"01AB\char"0304} \SUPPORTEDANCH{\char"01AD\char"0304} \SUPPORTEDANCH{\char"0236\char"0304} \SUPPORTEDPREC{\char"0075\char"0304} \MISSINGGLYPH{\char"AB52\char"0304} \SUPPORTEDPREC{\char"0079\char"0304} \SUPPORTEDANCH{\char"01B4\char"0304} \MISSINGGLYPH{\char"1EFF\char"0304}

\SUPPORTEDANCH{\char"0048\char"0306} \SUPPORTEDANCH{\char"2C67\char"0306} \SUPPORTEDANCH{\char"004B\char"0306} \SUPPORTEDANCH{\char"2C69\char"0306} \SUPPORTEDANCH{\char"004D\char"0306} \MISSINGGLYPH{\char"2C6E\char"0306} \SUPPORTEDANCH{\char"0050\char"0306} \SUPPORTEDANCH{\char"01A4\char"0306} \SUPPORTEDANCH{\char"0052\char"0306} \SUPPORTEDANCH{\char"2C64\char"0306} \SUPPORTEDANCH{\char"0054\char"0306} \SUPPORTEDANCH{\char"01AC\char"0306} \SUPPORTEDANCH{\char"01AE\char"0306} \SUPPORTEDANCH{\char"0059\char"0306} \SUPPORTEDANCH{\char"01B3\char"0306} \SUPPORTEDPREC{\char"0061\char"0306} \SUPPORTEDANCH{\char"1E9A\char"0306} \SUPPORTEDANCH{\char"0062\char"0306} \SUPPORTEDANCH{\char"0183\char"0306} \SUPPORTEDANCH{\char"0063\char"0306} \SUPPORTEDANCH{\char"0188\char"0306} \SUPPORTEDANCH{\char"0064\char"0306} \SUPPORTEDANCH{\char"018C\char"0306} \SUPPORTEDANCH{\char"0221\char"0306} \SUPPORTEDANCH{\char"0068\char"0306} \SUPPORTEDANCH{\char"2C68\char"0306} \SUPPORTEDANCH{\char"006B\char"0306} \SUPPORTEDANCH{\char"0199\char"0306} \SUPPORTEDANCH{\char"2C6A\char"0306} \SUPPORTEDANCH{\char"006C\char"0306} \SUPPORTEDANCH{\char"0234\char"0306} \SUPPORTEDANCH{\char"006E\char"0306} \SUPPORTEDANCH{\char"0235\char"0306} \SUPPORTEDANCH{\char"0070\char"0306} \SUPPORTEDANCH{\char"01A5\char"0306} \SUPPORTEDANCH{\char"0071\char"0306} \SUPPORTEDANCH{\char"02A0\char"0306} \SUPPORTEDANCH{\char"0074\char"0306} \SUPPORTEDANCH{\char"01AB\char"0306} \SUPPORTEDANCH{\char"01AD\char"0306} \SUPPORTEDANCH{\char"0236\char"0306} \SUPPORTEDPREC{\char"0075\char"0306} \MISSINGGLYPH{\char"AB52\char"0306} \SUPPORTEDANCH{\char"0079\char"0306} \SUPPORTEDANCH{\char"01B4\char"0306} \MISSINGGLYPH{\char"1EFF\char"0306}

\SUPPORTEDPREC{\char"0048\char"0307} \SUPPORTEDANCH{\char"2C67\char"0307} \SUPPORTEDANCH{\char"004B\char"0307} \SUPPORTEDANCH{\char"2C69\char"0307} \SUPPORTEDPREC{\char"004D\char"0307} \MISSINGGLYPH{\char"2C6E\char"0307} \SUPPORTEDPREC{\char"0050\char"0307} \SUPPORTEDANCH{\char"01A4\char"0307} \SUPPORTEDPREC{\char"0052\char"0307} \SUPPORTEDANCH{\char"2C64\char"0307} \SUPPORTEDPREC{\char"0054\char"0307} \SUPPORTEDANCH{\char"01AC\char"0307} \SUPPORTEDANCH{\char"01AE\char"0307} \SUPPORTEDPREC{\char"0059\char"0307} \SUPPORTEDANCH{\char"01B3\char"0307} \SUPPORTEDPREC{\char"0061\char"0307} \SUPPORTEDANCH{\char"1E9A\char"0307} \SUPPORTEDPREC{\char"0062\char"0307} \SUPPORTEDANCH{\char"0183\char"0307} \SUPPORTEDPREC{\char"0063\char"0307} \SUPPORTEDANCH{\char"0188\char"0307} \SUPPORTEDPREC{\char"0064\char"0307} \SUPPORTEDANCH{\char"018C\char"0307} \SUPPORTEDANCH{\char"0221\char"0307} \SUPPORTEDPREC{\char"0068\char"0307} \SUPPORTEDANCH{\char"2C68\char"0307} \SUPPORTEDANCH{\char"006B\char"0307} \SUPPORTEDANCH{\char"0199\char"0307} \SUPPORTEDANCH{\char"2C6A\char"0307} \SUPPORTEDANCH{\char"006C\char"0307} \SUPPORTEDANCH{\char"0234\char"0307} \SUPPORTEDPREC{\char"006E\char"0307} \SUPPORTEDANCH{\char"0235\char"0307} \SUPPORTEDPREC{\char"0070\char"0307} \SUPPORTEDANCH{\char"01A5\char"0307} \SUPPORTEDANCH{\char"0071\char"0307} \SUPPORTEDANCH{\char"02A0\char"0307} \SUPPORTEDPREC{\char"0074\char"0307} \SUPPORTEDANCH{\char"01AB\char"0307} \SUPPORTEDANCH{\char"01AD\char"0307} \SUPPORTEDANCH{\char"0236\char"0307} \SUPPORTEDANCH{\char"0075\char"0307} \MISSINGGLYPH{\char"AB52\char"0307} \SUPPORTEDPREC{\char"0079\char"0307} \SUPPORTEDANCH{\char"01B4\char"0307} \MISSINGGLYPH{\char"1EFF\char"0307}

\SUPPORTEDPREC{\char"0048\char"0308} \SUPPORTEDANCH{\char"2C67\char"0308} \SUPPORTEDANCH{\char"004B\char"0308} \SUPPORTEDANCH{\char"2C69\char"0308} \SUPPORTEDANCH{\char"004D\char"0308} \MISSINGGLYPH{\char"2C6E\char"0308} \SUPPORTEDANCH{\char"0050\char"0308} \SUPPORTEDANCH{\char"01A4\char"0308} \SUPPORTEDANCH{\char"0052\char"0308} \SUPPORTEDANCH{\char"2C64\char"0308} \SUPPORTEDANCH{\char"0054\char"0308} \SUPPORTEDANCH{\char"01AC\char"0308} \SUPPORTEDANCH{\char"01AE\char"0308} \SUPPORTEDPREC{\char"0059\char"0308} \SUPPORTEDANCH{\char"01B3\char"0308} \SUPPORTEDPREC{\char"0061\char"0308} \SUPPORTEDANCH{\char"1E9A\char"0308} \SUPPORTEDANCH{\char"0062\char"0308} \SUPPORTEDANCH{\char"0183\char"0308} \SUPPORTEDANCH{\char"0063\char"0308} \SUPPORTEDANCH{\char"0188\char"0308} \SUPPORTEDANCH{\char"0064\char"0308} \SUPPORTEDANCH{\char"018C\char"0308} \SUPPORTEDANCH{\char"0221\char"0308} \SUPPORTEDPREC{\char"0068\char"0308} \SUPPORTEDANCH{\char"2C68\char"0308} \SUPPORTEDANCH{\char"006B\char"0308} \SUPPORTEDANCH{\char"0199\char"0308} \SUPPORTEDANCH{\char"2C6A\char"0308} \SUPPORTEDANCH{\char"006C\char"0308} \SUPPORTEDANCH{\char"0234\char"0308} \SUPPORTEDANCH{\char"006E\char"0308} \SUPPORTEDANCH{\char"0235\char"0308} \SUPPORTEDANCH{\char"0070\char"0308} \SUPPORTEDANCH{\char"01A5\char"0308} \SUPPORTEDANCH{\char"0071\char"0308} \SUPPORTEDANCH{\char"02A0\char"0308} \SUPPORTEDPREC{\char"0074\char"0308} \SUPPORTEDANCH{\char"01AB\char"0308} \SUPPORTEDANCH{\char"01AD\char"0308} \SUPPORTEDANCH{\char"0236\char"0308} \SUPPORTEDPREC{\char"0075\char"0308} \MISSINGGLYPH{\char"AB52\char"0308} \SUPPORTEDPREC{\char"0079\char"0308} \SUPPORTEDANCH{\char"01B4\char"0308} \MISSINGGLYPH{\char"1EFF\char"0308}

\SUPPORTEDANCH{\char"0048\char"0309} \SUPPORTEDANCH{\char"2C67\char"0309} \SUPPORTEDANCH{\char"004B\char"0309} \SUPPORTEDANCH{\char"2C69\char"0309} \SUPPORTEDANCH{\char"004D\char"0309} \MISSINGGLYPH{\char"2C6E\char"0309} \SUPPORTEDANCH{\char"0050\char"0309} \SUPPORTEDANCH{\char"01A4\char"0309} \SUPPORTEDANCH{\char"0052\char"0309} \SUPPORTEDANCH{\char"2C64\char"0309} \SUPPORTEDANCH{\char"0054\char"0309} \SUPPORTEDANCH{\char"01AC\char"0309} \SUPPORTEDANCH{\char"01AE\char"0309} \SUPPORTEDPREC{\char"0059\char"0309} \SUPPORTEDANCH{\char"01B3\char"0309} \SUPPORTEDPREC{\char"0061\char"0309} \SUPPORTEDANCH{\char"1E9A\char"0309} \SUPPORTEDANCH{\char"0062\char"0309} \SUPPORTEDANCH{\char"0183\char"0309} \SUPPORTEDANCH{\char"0063\char"0309} \SUPPORTEDANCH{\char"0188\char"0309} \SUPPORTEDANCH{\char"0064\char"0309} \SUPPORTEDANCH{\char"018C\char"0309} \SUPPORTEDANCH{\char"0221\char"0309} \SUPPORTEDANCH{\char"0068\char"0309} \SUPPORTEDANCH{\char"2C68\char"0309} \SUPPORTEDANCH{\char"006B\char"0309} \SUPPORTEDANCH{\char"0199\char"0309} \SUPPORTEDANCH{\char"2C6A\char"0309} \SUPPORTEDANCH{\char"006C\char"0309} \SUPPORTEDANCH{\char"0234\char"0309} \SUPPORTEDANCH{\char"006E\char"0309} \SUPPORTEDANCH{\char"0235\char"0309} \SUPPORTEDANCH{\char"0070\char"0309} \SUPPORTEDANCH{\char"01A5\char"0309} \SUPPORTEDANCH{\char"0071\char"0309} \SUPPORTEDANCH{\char"02A0\char"0309} \SUPPORTEDANCH{\char"0074\char"0309} \SUPPORTEDANCH{\char"01AB\char"0309} \SUPPORTEDANCH{\char"01AD\char"0309} \SUPPORTEDANCH{\char"0236\char"0309} \SUPPORTEDPREC{\char"0075\char"0309} \MISSINGGLYPH{\char"AB52\char"0309} \SUPPORTEDPREC{\char"0079\char"0309} \SUPPORTEDANCH{\char"01B4\char"0309} \MISSINGGLYPH{\char"1EFF\char"0309}

\SUPPORTEDANCH{\char"0048\char"030A} \SUPPORTEDANCH{\char"2C67\char"030A} \SUPPORTEDANCH{\char"004B\char"030A} \SUPPORTEDANCH{\char"2C69\char"030A} \SUPPORTEDANCH{\char"004D\char"030A} \MISSINGGLYPH{\char"2C6E\char"030A} \SUPPORTEDANCH{\char"0050\char"030A} \SUPPORTEDANCH{\char"01A4\char"030A} \SUPPORTEDANCH{\char"0052\char"030A} \SUPPORTEDANCH{\char"2C64\char"030A} \SUPPORTEDANCH{\char"0054\char"030A} \SUPPORTEDANCH{\char"01AC\char"030A} \SUPPORTEDANCH{\char"01AE\char"030A} \SUPPORTEDANCH{\char"0059\char"030A} \SUPPORTEDANCH{\char"01B3\char"030A} \SUPPORTEDPREC{\char"0061\char"030A} \SUPPORTEDANCH{\char"1E9A\char"030A} \SUPPORTEDANCH{\char"0062\char"030A} \SUPPORTEDANCH{\char"0183\char"030A} \SUPPORTEDANCH{\char"0063\char"030A} \SUPPORTEDANCH{\char"0188\char"030A} \SUPPORTEDANCH{\char"0064\char"030A} \SUPPORTEDANCH{\char"018C\char"030A} \SUPPORTEDANCH{\char"0221\char"030A} \SUPPORTEDANCH{\char"0068\char"030A} \SUPPORTEDANCH{\char"2C68\char"030A} \SUPPORTEDANCH{\char"006B\char"030A} \SUPPORTEDANCH{\char"0199\char"030A} \SUPPORTEDANCH{\char"2C6A\char"030A} \SUPPORTEDANCH{\char"006C\char"030A} \SUPPORTEDANCH{\char"0234\char"030A} \SUPPORTEDANCH{\char"006E\char"030A} \SUPPORTEDANCH{\char"0235\char"030A} \SUPPORTEDANCH{\char"0070\char"030A} \SUPPORTEDANCH{\char"01A5\char"030A} \SUPPORTEDANCH{\char"0071\char"030A} \SUPPORTEDANCH{\char"02A0\char"030A} \SUPPORTEDANCH{\char"0074\char"030A} \SUPPORTEDANCH{\char"01AB\char"030A} \SUPPORTEDANCH{\char"01AD\char"030A} \SUPPORTEDANCH{\char"0236\char"030A} \SUPPORTEDPREC{\char"0075\char"030A} \MISSINGGLYPH{\char"AB52\char"030A} \SUPPORTEDPREC{\char"0079\char"030A} \SUPPORTEDANCH{\char"01B4\char"030A} \MISSINGGLYPH{\char"1EFF\char"030A}

\SUPPORTEDANCH{\char"0048\char"030B} \SUPPORTEDANCH{\char"2C67\char"030B} \SUPPORTEDANCH{\char"004B\char"030B} \SUPPORTEDANCH{\char"2C69\char"030B} \SUPPORTEDANCH{\char"004D\char"030B} \MISSINGGLYPH{\char"2C6E\char"030B} \SUPPORTEDANCH{\char"0050\char"030B} \SUPPORTEDANCH{\char"01A4\char"030B} \SUPPORTEDANCH{\char"0052\char"030B} \SUPPORTEDANCH{\char"2C64\char"030B} \SUPPORTEDANCH{\char"0054\char"030B} \SUPPORTEDANCH{\char"01AC\char"030B} \SUPPORTEDANCH{\char"01AE\char"030B} \SUPPORTEDANCH{\char"0059\char"030B} \SUPPORTEDANCH{\char"01B3\char"030B} \SUPPORTEDANCH{\char"0061\char"030B} \SUPPORTEDANCH{\char"1E9A\char"030B} \SUPPORTEDANCH{\char"0062\char"030B} \SUPPORTEDANCH{\char"0183\char"030B} \SUPPORTEDANCH{\char"0063\char"030B} \SUPPORTEDANCH{\char"0188\char"030B} \SUPPORTEDANCH{\char"0064\char"030B} \SUPPORTEDANCH{\char"018C\char"030B} \SUPPORTEDANCH{\char"0221\char"030B} \SUPPORTEDANCH{\char"0068\char"030B} \SUPPORTEDANCH{\char"2C68\char"030B} \SUPPORTEDANCH{\char"006B\char"030B} \SUPPORTEDANCH{\char"0199\char"030B} \SUPPORTEDANCH{\char"2C6A\char"030B} \SUPPORTEDANCH{\char"006C\char"030B} \SUPPORTEDANCH{\char"0234\char"030B} \SUPPORTEDANCH{\char"006E\char"030B} \SUPPORTEDANCH{\char"0235\char"030B} \SUPPORTEDANCH{\char"0070\char"030B} \SUPPORTEDANCH{\char"01A5\char"030B} \SUPPORTEDANCH{\char"0071\char"030B} \SUPPORTEDANCH{\char"02A0\char"030B} \SUPPORTEDANCH{\char"0074\char"030B} \SUPPORTEDANCH{\char"01AB\char"030B} \SUPPORTEDANCH{\char"01AD\char"030B} \SUPPORTEDANCH{\char"0236\char"030B} \SUPPORTEDPREC{\char"0075\char"030B} \MISSINGGLYPH{\char"AB52\char"030B} \SUPPORTEDANCH{\char"0079\char"030B} \SUPPORTEDANCH{\char"01B4\char"030B} \MISSINGGLYPH{\char"1EFF\char"030B}

\SUPPORTEDPREC{\char"0048\char"030C} \SUPPORTEDANCH{\char"2C67\char"030C} \SUPPORTEDPREC{\char"004B\char"030C} \SUPPORTEDANCH{\char"2C69\char"030C} \SUPPORTEDANCH{\char"004D\char"030C} \MISSINGGLYPH{\char"2C6E\char"030C} \SUPPORTEDANCH{\char"0050\char"030C} \SUPPORTEDANCH{\char"01A4\char"030C} \SUPPORTEDPREC{\char"0052\char"030C} \SUPPORTEDANCH{\char"2C64\char"030C} \SUPPORTEDPREC{\char"0054\char"030C} \SUPPORTEDANCH{\char"01AC\char"030C} \SUPPORTEDANCH{\char"01AE\char"030C} \SUPPORTEDANCH{\char"0059\char"030C} \SUPPORTEDANCH{\char"01B3\char"030C} \SUPPORTEDPREC{\char"0061\char"030C} \SUPPORTEDANCH{\char"1E9A\char"030C} \SUPPORTEDANCH{\char"0062\char"030C} \SUPPORTEDANCH{\char"0183\char"030C} \SUPPORTEDPREC{\char"0063\char"030C} \SUPPORTEDANCH{\char"0188\char"030C} \SUPPORTEDPREC{\char"0064\char"030C} \SUPPORTEDANCH{\char"018C\char"030C} \SUPPORTEDANCH{\char"0221\char"030C} \SUPPORTEDPREC{\char"0068\char"030C} \SUPPORTEDANCH{\char"2C68\char"030C} \SUPPORTEDPREC{\char"006B\char"030C} \SUPPORTEDANCH{\char"0199\char"030C} \SUPPORTEDANCH{\char"2C6A\char"030C} \SUPPORTEDPREC{\char"006C\char"030C} \SUPPORTEDANCH{\char"0234\char"030C} \SUPPORTEDPREC{\char"006E\char"030C} \SUPPORTEDANCH{\char"0235\char"030C} \SUPPORTEDANCH{\char"0070\char"030C} \SUPPORTEDANCH{\char"01A5\char"030C} \SUPPORTEDANCH{\char"0071\char"030C} \SUPPORTEDANCH{\char"02A0\char"030C} \SUPPORTEDPREC{\char"0074\char"030C} \SUPPORTEDANCH{\char"01AB\char"030C} \SUPPORTEDANCH{\char"01AD\char"030C} \SUPPORTEDANCH{\char"0236\char"030C} \SUPPORTEDPREC{\char"0075\char"030C} \MISSINGGLYPH{\char"AB52\char"030C} \SUPPORTEDANCH{\char"0079\char"030C} \SUPPORTEDANCH{\char"01B4\char"030C} \MISSINGGLYPH{\char"1EFF\char"030C}

\SUPPORTEDANCH{\char"0048\char"030D} \SUPPORTEDANCH{\char"2C67\char"030D} \SUPPORTEDANCH{\char"004B\char"030D} \SUPPORTEDANCH{\char"2C69\char"030D} \SUPPORTEDANCH{\char"004D\char"030D} \MISSINGGLYPH{\char"2C6E\char"030D} \SUPPORTEDANCH{\char"0050\char"030D} \SUPPORTEDANCH{\char"01A4\char"030D} \SUPPORTEDANCH{\char"0052\char"030D} \SUPPORTEDANCH{\char"2C64\char"030D} \SUPPORTEDANCH{\char"0054\char"030D} \SUPPORTEDANCH{\char"01AC\char"030D} \SUPPORTEDANCH{\char"01AE\char"030D} \SUPPORTEDANCH{\char"0059\char"030D} \SUPPORTEDANCH{\char"01B3\char"030D} \SUPPORTEDANCH{\char"0061\char"030D} \SUPPORTEDANCH{\char"1E9A\char"030D} \SUPPORTEDANCH{\char"0062\char"030D} \SUPPORTEDANCH{\char"0183\char"030D} \SUPPORTEDANCH{\char"0063\char"030D} \SUPPORTEDANCH{\char"0188\char"030D} \SUPPORTEDANCH{\char"0064\char"030D} \SUPPORTEDANCH{\char"018C\char"030D} \SUPPORTEDANCH{\char"0221\char"030D} \SUPPORTEDANCH{\char"0068\char"030D} \SUPPORTEDANCH{\char"2C68\char"030D} \SUPPORTEDANCH{\char"006B\char"030D} \SUPPORTEDANCH{\char"0199\char"030D} \SUPPORTEDANCH{\char"2C6A\char"030D} \SUPPORTEDANCH{\char"006C\char"030D} \SUPPORTEDANCH{\char"0234\char"030D} \SUPPORTEDANCH{\char"006E\char"030D} \SUPPORTEDANCH{\char"0235\char"030D} \SUPPORTEDANCH{\char"0070\char"030D} \SUPPORTEDANCH{\char"01A5\char"030D} \SUPPORTEDANCH{\char"0071\char"030D} \SUPPORTEDANCH{\char"02A0\char"030D} \SUPPORTEDANCH{\char"0074\char"030D} \SUPPORTEDANCH{\char"01AB\char"030D} \SUPPORTEDANCH{\char"01AD\char"030D} \SUPPORTEDANCH{\char"0236\char"030D} \SUPPORTEDANCH{\char"0075\char"030D} \MISSINGGLYPH{\char"AB52\char"030D} \SUPPORTEDANCH{\char"0079\char"030D} \SUPPORTEDANCH{\char"01B4\char"030D} \MISSINGGLYPH{\char"1EFF\char"030D}

\SUPPORTEDANCH{\char"0048\char"030F} \SUPPORTEDANCH{\char"2C67\char"030F} \SUPPORTEDANCH{\char"004B\char"030F} \SUPPORTEDANCH{\char"2C69\char"030F} \SUPPORTEDANCH{\char"004D\char"030F} \MISSINGGLYPH{\char"2C6E\char"030F} \SUPPORTEDANCH{\char"0050\char"030F} \SUPPORTEDANCH{\char"01A4\char"030F} \SUPPORTEDPREC{\char"0052\char"030F} \SUPPORTEDANCH{\char"2C64\char"030F} \SUPPORTEDANCH{\char"0054\char"030F} \SUPPORTEDANCH{\char"01AC\char"030F} \SUPPORTEDANCH{\char"01AE\char"030F} \SUPPORTEDANCH{\char"0059\char"030F} \SUPPORTEDANCH{\char"01B3\char"030F} \SUPPORTEDPREC{\char"0061\char"030F} \SUPPORTEDANCH{\char"1E9A\char"030F} \SUPPORTEDANCH{\char"0062\char"030F} \SUPPORTEDANCH{\char"0183\char"030F} \SUPPORTEDANCH{\char"0063\char"030F} \SUPPORTEDANCH{\char"0188\char"030F} \SUPPORTEDANCH{\char"0064\char"030F} \SUPPORTEDANCH{\char"018C\char"030F} \SUPPORTEDANCH{\char"0221\char"030F} \SUPPORTEDANCH{\char"0068\char"030F} \SUPPORTEDANCH{\char"2C68\char"030F} \SUPPORTEDANCH{\char"006B\char"030F} \SUPPORTEDANCH{\char"0199\char"030F} \SUPPORTEDANCH{\char"2C6A\char"030F} \SUPPORTEDANCH{\char"006C\char"030F} \SUPPORTEDANCH{\char"0234\char"030F} \SUPPORTEDANCH{\char"006E\char"030F} \SUPPORTEDANCH{\char"0235\char"030F} \SUPPORTEDANCH{\char"0070\char"030F} \SUPPORTEDANCH{\char"01A5\char"030F} \SUPPORTEDANCH{\char"0071\char"030F} \SUPPORTEDANCH{\char"02A0\char"030F} \SUPPORTEDANCH{\char"0074\char"030F} \SUPPORTEDANCH{\char"01AB\char"030F} \SUPPORTEDANCH{\char"01AD\char"030F} \SUPPORTEDANCH{\char"0236\char"030F} \SUPPORTEDPREC{\char"0075\char"030F} \MISSINGGLYPH{\char"AB52\char"030F} \SUPPORTEDANCH{\char"0079\char"030F} \SUPPORTEDANCH{\char"01B4\char"030F} \MISSINGGLYPH{\char"1EFF\char"030F}

\SUPPORTEDANCH{\char"0048\char"0310} \SUPPORTEDANCH{\char"2C67\char"0310} \SUPPORTEDANCH{\char"004B\char"0310} \SUPPORTEDANCH{\char"2C69\char"0310} \SUPPORTEDANCH{\char"004D\char"0310} \MISSINGGLYPH{\char"2C6E\char"0310} \SUPPORTEDANCH{\char"0050\char"0310} \SUPPORTEDANCH{\char"01A4\char"0310} \SUPPORTEDANCH{\char"0052\char"0310} \SUPPORTEDANCH{\char"2C64\char"0310} \SUPPORTEDANCH{\char"0054\char"0310} \SUPPORTEDANCH{\char"01AC\char"0310} \SUPPORTEDANCH{\char"01AE\char"0310} \SUPPORTEDANCH{\char"0059\char"0310} \SUPPORTEDANCH{\char"01B3\char"0310} \SUPPORTEDANCH{\char"0061\char"0310} \SUPPORTEDANCH{\char"1E9A\char"0310} \SUPPORTEDANCH{\char"0062\char"0310} \SUPPORTEDANCH{\char"0183\char"0310} \SUPPORTEDANCH{\char"0063\char"0310} \SUPPORTEDANCH{\char"0188\char"0310} \SUPPORTEDANCH{\char"0064\char"0310} \SUPPORTEDANCH{\char"018C\char"0310} \SUPPORTEDANCH{\char"0221\char"0310} \SUPPORTEDANCH{\char"0068\char"0310} \SUPPORTEDANCH{\char"2C68\char"0310} \SUPPORTEDANCH{\char"006B\char"0310} \SUPPORTEDANCH{\char"0199\char"0310} \SUPPORTEDANCH{\char"2C6A\char"0310} \SUPPORTEDANCH{\char"006C\char"0310} \SUPPORTEDANCH{\char"0234\char"0310} \SUPPORTEDANCH{\char"006E\char"0310} \SUPPORTEDANCH{\char"0235\char"0310} \SUPPORTEDANCH{\char"0070\char"0310} \SUPPORTEDANCH{\char"01A5\char"0310} \SUPPORTEDANCH{\char"0071\char"0310} \SUPPORTEDANCH{\char"02A0\char"0310} \SUPPORTEDANCH{\char"0074\char"0310} \SUPPORTEDANCH{\char"01AB\char"0310} \SUPPORTEDANCH{\char"01AD\char"0310} \SUPPORTEDANCH{\char"0236\char"0310} \SUPPORTEDANCH{\char"0075\char"0310} \MISSINGGLYPH{\char"AB52\char"0310} \SUPPORTEDANCH{\char"0079\char"0310} \SUPPORTEDANCH{\char"01B4\char"0310} \MISSINGGLYPH{\char"1EFF\char"0310}

\SUPPORTEDANCH{\char"0048\char"0311} \SUPPORTEDANCH{\char"2C67\char"0311} \SUPPORTEDANCH{\char"004B\char"0311} \SUPPORTEDANCH{\char"2C69\char"0311} \SUPPORTEDANCH{\char"004D\char"0311} \MISSINGGLYPH{\char"2C6E\char"0311} \SUPPORTEDANCH{\char"0050\char"0311} \SUPPORTEDANCH{\char"01A4\char"0311} \SUPPORTEDPREC{\char"0052\char"0311} \SUPPORTEDANCH{\char"2C64\char"0311} \SUPPORTEDANCH{\char"0054\char"0311} \SUPPORTEDANCH{\char"01AC\char"0311} \SUPPORTEDANCH{\char"01AE\char"0311} \SUPPORTEDANCH{\char"0059\char"0311} \SUPPORTEDANCH{\char"01B3\char"0311} \SUPPORTEDPREC{\char"0061\char"0311} \SUPPORTEDANCH{\char"1E9A\char"0311} \SUPPORTEDANCH{\char"0062\char"0311} \SUPPORTEDANCH{\char"0183\char"0311} \SUPPORTEDANCH{\char"0063\char"0311} \SUPPORTEDANCH{\char"0188\char"0311} \SUPPORTEDANCH{\char"0064\char"0311} \SUPPORTEDANCH{\char"018C\char"0311} \SUPPORTEDANCH{\char"0221\char"0311} \SUPPORTEDANCH{\char"0068\char"0311} \SUPPORTEDANCH{\char"2C68\char"0311} \SUPPORTEDANCH{\char"006B\char"0311} \SUPPORTEDANCH{\char"0199\char"0311} \SUPPORTEDANCH{\char"2C6A\char"0311} \SUPPORTEDANCH{\char"006C\char"0311} \SUPPORTEDANCH{\char"0234\char"0311} \SUPPORTEDANCH{\char"006E\char"0311} \SUPPORTEDANCH{\char"0235\char"0311} \SUPPORTEDANCH{\char"0070\char"0311} \SUPPORTEDANCH{\char"01A5\char"0311} \SUPPORTEDANCH{\char"0071\char"0311} \SUPPORTEDANCH{\char"02A0\char"0311} \SUPPORTEDANCH{\char"0074\char"0311} \SUPPORTEDANCH{\char"01AB\char"0311} \SUPPORTEDANCH{\char"01AD\char"0311} \SUPPORTEDANCH{\char"0236\char"0311} \SUPPORTEDPREC{\char"0075\char"0311} \MISSINGGLYPH{\char"AB52\char"0311} \SUPPORTEDANCH{\char"0079\char"0311} \SUPPORTEDANCH{\char"01B4\char"0311} \MISSINGGLYPH{\char"1EFF\char"0311}

\SUPPORTEDANCH{\char"0048\char"0312} \SUPPORTEDANCH{\char"2C67\char"0312} \SUPPORTEDANCH{\char"004B\char"0312} \SUPPORTEDANCH{\char"2C69\char"0312} \SUPPORTEDANCH{\char"004D\char"0312} \MISSINGGLYPH{\char"2C6E\char"0312} \SUPPORTEDANCH{\char"0050\char"0312} \SUPPORTEDANCH{\char"01A4\char"0312} \SUPPORTEDANCH{\char"0052\char"0312} \SUPPORTEDANCH{\char"2C64\char"0312} \SUPPORTEDANCH{\char"0054\char"0312} \SUPPORTEDANCH{\char"01AC\char"0312} \SUPPORTEDANCH{\char"01AE\char"0312} \SUPPORTEDANCH{\char"0059\char"0312} \SUPPORTEDANCH{\char"01B3\char"0312} \SUPPORTEDANCH{\char"0061\char"0312} \SUPPORTEDANCH{\char"1E9A\char"0312} \SUPPORTEDANCH{\char"0062\char"0312} \SUPPORTEDANCH{\char"0183\char"0312} \SUPPORTEDANCH{\char"0063\char"0312} \SUPPORTEDANCH{\char"0188\char"0312} \SUPPORTEDANCH{\char"0064\char"0312} \SUPPORTEDANCH{\char"018C\char"0312} \SUPPORTEDANCH{\char"0221\char"0312} \SUPPORTEDANCH{\char"0068\char"0312} \SUPPORTEDANCH{\char"2C68\char"0312} \SUPPORTEDANCH{\char"006B\char"0312} \SUPPORTEDANCH{\char"0199\char"0312} \SUPPORTEDANCH{\char"2C6A\char"0312} \SUPPORTEDANCH{\char"006C\char"0312} \SUPPORTEDANCH{\char"0234\char"0312} \SUPPORTEDANCH{\char"006E\char"0312} \SUPPORTEDANCH{\char"0235\char"0312} \SUPPORTEDANCH{\char"0070\char"0312} \SUPPORTEDANCH{\char"01A5\char"0312} \SUPPORTEDANCH{\char"0071\char"0312} \SUPPORTEDANCH{\char"02A0\char"0312} \SUPPORTEDANCH{\char"0074\char"0312} \SUPPORTEDANCH{\char"01AB\char"0312} \SUPPORTEDANCH{\char"01AD\char"0312} \SUPPORTEDANCH{\char"0236\char"0312} \SUPPORTEDANCH{\char"0075\char"0312} \MISSINGGLYPH{\char"AB52\char"0312} \SUPPORTEDANCH{\char"0079\char"0312} \SUPPORTEDANCH{\char"01B4\char"0312} \MISSINGGLYPH{\char"1EFF\char"0312}

\SUPPORTEDANCH{\char"0048\char"033D} \SUPPORTEDANCH{\char"2C67\char"033D} \SUPPORTEDANCH{\char"004B\char"033D} \SUPPORTEDANCH{\char"2C69\char"033D} \SUPPORTEDANCH{\char"004D\char"033D} \MISSINGGLYPH{\char"2C6E\char"033D} \SUPPORTEDANCH{\char"0050\char"033D} \SUPPORTEDANCH{\char"01A4\char"033D} \SUPPORTEDANCH{\char"0052\char"033D} \SUPPORTEDANCH{\char"2C64\char"033D} \SUPPORTEDANCH{\char"0054\char"033D} \SUPPORTEDANCH{\char"01AC\char"033D} \SUPPORTEDANCH{\char"01AE\char"033D} \SUPPORTEDANCH{\char"0059\char"033D} \SUPPORTEDANCH{\char"01B3\char"033D} \SUPPORTEDANCH{\char"0061\char"033D} \SUPPORTEDANCH{\char"1E9A\char"033D} \SUPPORTEDANCH{\char"0062\char"033D} \SUPPORTEDANCH{\char"0183\char"033D} \SUPPORTEDANCH{\char"0063\char"033D} \SUPPORTEDANCH{\char"0188\char"033D} \SUPPORTEDANCH{\char"0064\char"033D} \SUPPORTEDANCH{\char"018C\char"033D} \SUPPORTEDANCH{\char"0221\char"033D} \SUPPORTEDANCH{\char"0068\char"033D} \SUPPORTEDANCH{\char"2C68\char"033D} \SUPPORTEDANCH{\char"006B\char"033D} \SUPPORTEDANCH{\char"0199\char"033D} \SUPPORTEDANCH{\char"2C6A\char"033D} \SUPPORTEDANCH{\char"006C\char"033D} \SUPPORTEDANCH{\char"0234\char"033D} \SUPPORTEDANCH{\char"006E\char"033D} \SUPPORTEDANCH{\char"0235\char"033D} \SUPPORTEDANCH{\char"0070\char"033D} \SUPPORTEDANCH{\char"01A5\char"033D} \SUPPORTEDANCH{\char"0071\char"033D} \SUPPORTEDANCH{\char"02A0\char"033D} \SUPPORTEDANCH{\char"0074\char"033D} \SUPPORTEDANCH{\char"01AB\char"033D} \SUPPORTEDANCH{\char"01AD\char"033D} \SUPPORTEDANCH{\char"0236\char"033D} \SUPPORTEDANCH{\char"0075\char"033D} \MISSINGGLYPH{\char"AB52\char"033D} \SUPPORTEDANCH{\char"0079\char"033D} \SUPPORTEDANCH{\char"01B4\char"033D} \MISSINGGLYPH{\char"1EFF\char"033D}

\SUPPORTEDANCH{\char"0048\char"033E} \SUPPORTEDANCH{\char"2C67\char"033E} \SUPPORTEDANCH{\char"004B\char"033E} \SUPPORTEDANCH{\char"2C69\char"033E} \SUPPORTEDANCH{\char"004D\char"033E} \MISSINGGLYPH{\char"2C6E\char"033E} \SUPPORTEDANCH{\char"0050\char"033E} \SUPPORTEDANCH{\char"01A4\char"033E} \SUPPORTEDANCH{\char"0052\char"033E} \SUPPORTEDANCH{\char"2C64\char"033E} \SUPPORTEDANCH{\char"0054\char"033E} \SUPPORTEDANCH{\char"01AC\char"033E} \SUPPORTEDANCH{\char"01AE\char"033E} \SUPPORTEDANCH{\char"0059\char"033E} \SUPPORTEDANCH{\char"01B3\char"033E} \SUPPORTEDANCH{\char"0061\char"033E} \SUPPORTEDANCH{\char"1E9A\char"033E} \SUPPORTEDANCH{\char"0062\char"033E} \SUPPORTEDANCH{\char"0183\char"033E} \SUPPORTEDANCH{\char"0063\char"033E} \SUPPORTEDANCH{\char"0188\char"033E} \SUPPORTEDANCH{\char"0064\char"033E} \SUPPORTEDANCH{\char"018C\char"033E} \SUPPORTEDANCH{\char"0221\char"033E} \SUPPORTEDANCH{\char"0068\char"033E} \SUPPORTEDANCH{\char"2C68\char"033E} \SUPPORTEDANCH{\char"006B\char"033E} \SUPPORTEDANCH{\char"0199\char"033E} \SUPPORTEDANCH{\char"2C6A\char"033E} \SUPPORTEDANCH{\char"006C\char"033E} \SUPPORTEDANCH{\char"0234\char"033E} \SUPPORTEDANCH{\char"006E\char"033E} \SUPPORTEDANCH{\char"0235\char"033E} \SUPPORTEDANCH{\char"0070\char"033E} \SUPPORTEDANCH{\char"01A5\char"033E} \SUPPORTEDANCH{\char"0071\char"033E} \SUPPORTEDANCH{\char"02A0\char"033E} \SUPPORTEDANCH{\char"0074\char"033E} \SUPPORTEDANCH{\char"01AB\char"033E} \SUPPORTEDANCH{\char"01AD\char"033E} \SUPPORTEDANCH{\char"0236\char"033E} \SUPPORTEDANCH{\char"0075\char"033E} \MISSINGGLYPH{\char"AB52\char"033E} \SUPPORTEDANCH{\char"0079\char"033E} \SUPPORTEDANCH{\char"01B4\char"033E} \MISSINGGLYPH{\char"1EFF\char"033E}

\SUPPORTEDANCH{\char"0048\char"033F} \SUPPORTEDANCH{\char"2C67\char"033F} \SUPPORTEDANCH{\char"004B\char"033F} \SUPPORTEDANCH{\char"2C69\char"033F} \SUPPORTEDANCH{\char"004D\char"033F} \MISSINGGLYPH{\char"2C6E\char"033F} \SUPPORTEDANCH{\char"0050\char"033F} \SUPPORTEDANCH{\char"01A4\char"033F} \SUPPORTEDANCH{\char"0052\char"033F} \SUPPORTEDANCH{\char"2C64\char"033F} \SUPPORTEDANCH{\char"0054\char"033F} \SUPPORTEDANCH{\char"01AC\char"033F} \SUPPORTEDANCH{\char"01AE\char"033F} \SUPPORTEDANCH{\char"0059\char"033F} \SUPPORTEDANCH{\char"01B3\char"033F} \SUPPORTEDANCH{\char"0061\char"033F} \SUPPORTEDANCH{\char"1E9A\char"033F} \SUPPORTEDANCH{\char"0062\char"033F} \SUPPORTEDANCH{\char"0183\char"033F} \SUPPORTEDANCH{\char"0063\char"033F} \SUPPORTEDANCH{\char"0188\char"033F} \SUPPORTEDANCH{\char"0064\char"033F} \SUPPORTEDANCH{\char"018C\char"033F} \SUPPORTEDANCH{\char"0221\char"033F} \SUPPORTEDANCH{\char"0068\char"033F} \SUPPORTEDANCH{\char"2C68\char"033F} \SUPPORTEDANCH{\char"006B\char"033F} \SUPPORTEDANCH{\char"0199\char"033F} \SUPPORTEDANCH{\char"2C6A\char"033F} \SUPPORTEDANCH{\char"006C\char"033F} \SUPPORTEDANCH{\char"0234\char"033F} \SUPPORTEDANCH{\char"006E\char"033F} \SUPPORTEDANCH{\char"0235\char"033F} \SUPPORTEDANCH{\char"0070\char"033F} \SUPPORTEDANCH{\char"01A5\char"033F} \SUPPORTEDANCH{\char"0071\char"033F} \SUPPORTEDANCH{\char"02A0\char"033F} \SUPPORTEDANCH{\char"0074\char"033F} \SUPPORTEDANCH{\char"01AB\char"033F} \SUPPORTEDANCH{\char"01AD\char"033F} \SUPPORTEDANCH{\char"0236\char"033F} \SUPPORTEDANCH{\char"0075\char"033F} \MISSINGGLYPH{\char"AB52\char"033F} \SUPPORTEDANCH{\char"0079\char"033F} \SUPPORTEDANCH{\char"01B4\char"033F} \MISSINGGLYPH{\char"1EFF\char"033F}

\SUPPORTEDANCH{\char"0048\char"030E} \SUPPORTEDANCH{\char"2C67\char"030E} \SUPPORTEDANCH{\char"004B\char"030E} \SUPPORTEDANCH{\char"2C69\char"030E} \SUPPORTEDANCH{\char"004D\char"030E} \MISSINGGLYPH{\char"2C6E\char"030E} \SUPPORTEDANCH{\char"0050\char"030E} \SUPPORTEDANCH{\char"01A4\char"030E} \SUPPORTEDANCH{\char"0052\char"030E} \SUPPORTEDANCH{\char"2C64\char"030E} \SUPPORTEDANCH{\char"0054\char"030E} \SUPPORTEDANCH{\char"01AC\char"030E} \SUPPORTEDANCH{\char"01AE\char"030E} \SUPPORTEDANCH{\char"0059\char"030E} \SUPPORTEDANCH{\char"01B3\char"030E} \SUPPORTEDANCH{\char"0061\char"030E} \SUPPORTEDANCH{\char"1E9A\char"030E} \SUPPORTEDANCH{\char"0062\char"030E} \SUPPORTEDANCH{\char"0183\char"030E} \SUPPORTEDANCH{\char"0063\char"030E} \SUPPORTEDANCH{\char"0188\char"030E} \SUPPORTEDANCH{\char"0064\char"030E} \SUPPORTEDANCH{\char"018C\char"030E} \SUPPORTEDANCH{\char"0221\char"030E} \SUPPORTEDANCH{\char"0068\char"030E} \SUPPORTEDANCH{\char"2C68\char"030E} \SUPPORTEDANCH{\char"006B\char"030E} \SUPPORTEDANCH{\char"0199\char"030E} \SUPPORTEDANCH{\char"2C6A\char"030E} \SUPPORTEDANCH{\char"006C\char"030E} \SUPPORTEDANCH{\char"0234\char"030E} \SUPPORTEDANCH{\char"006E\char"030E} \SUPPORTEDANCH{\char"0235\char"030E} \SUPPORTEDANCH{\char"0070\char"030E} \SUPPORTEDANCH{\char"01A5\char"030E} \SUPPORTEDANCH{\char"0071\char"030E} \SUPPORTEDANCH{\char"02A0\char"030E} \SUPPORTEDANCH{\char"0074\char"030E} \SUPPORTEDANCH{\char"01AB\char"030E} \SUPPORTEDANCH{\char"01AD\char"030E} \SUPPORTEDANCH{\char"0236\char"030E} \SUPPORTEDANCH{\char"0075\char"030E} \MISSINGGLYPH{\char"AB52\char"030E} \SUPPORTEDANCH{\char"0079\char"030E} \SUPPORTEDANCH{\char"01B4\char"030E} \MISSINGGLYPH{\char"1EFF\char"030E}

\SUPPORTEDANCH{\char"0048\char"0346} \SUPPORTEDANCH{\char"2C67\char"0346} \SUPPORTEDANCH{\char"004B\char"0346} \SUPPORTEDANCH{\char"2C69\char"0346} \SUPPORTEDANCH{\char"004D\char"0346} \MISSINGGLYPH{\char"2C6E\char"0346} \SUPPORTEDANCH{\char"0050\char"0346} \SUPPORTEDANCH{\char"01A4\char"0346} \SUPPORTEDANCH{\char"0052\char"0346} \SUPPORTEDANCH{\char"2C64\char"0346} \SUPPORTEDANCH{\char"0054\char"0346} \SUPPORTEDANCH{\char"01AC\char"0346} \SUPPORTEDANCH{\char"01AE\char"0346} \SUPPORTEDANCH{\char"0059\char"0346} \SUPPORTEDANCH{\char"01B3\char"0346} \SUPPORTEDANCH{\char"0061\char"0346} \SUPPORTEDANCH{\char"1E9A\char"0346} \SUPPORTEDANCH{\char"0062\char"0346} \SUPPORTEDANCH{\char"0183\char"0346} \SUPPORTEDANCH{\char"0063\char"0346} \SUPPORTEDANCH{\char"0188\char"0346} \SUPPORTEDANCH{\char"0064\char"0346} \SUPPORTEDANCH{\char"018C\char"0346} \SUPPORTEDANCH{\char"0221\char"0346} \SUPPORTEDANCH{\char"0068\char"0346} \SUPPORTEDANCH{\char"2C68\char"0346} \SUPPORTEDANCH{\char"006B\char"0346} \SUPPORTEDANCH{\char"0199\char"0346} \SUPPORTEDANCH{\char"2C6A\char"0346} \SUPPORTEDANCH{\char"006C\char"0346} \SUPPORTEDANCH{\char"0234\char"0346} \SUPPORTEDANCH{\char"006E\char"0346} \SUPPORTEDANCH{\char"0235\char"0346} \SUPPORTEDANCH{\char"0070\char"0346} \SUPPORTEDANCH{\char"01A5\char"0346} \SUPPORTEDANCH{\char"0071\char"0346} \SUPPORTEDANCH{\char"02A0\char"0346} \SUPPORTEDANCH{\char"0074\char"0346} \SUPPORTEDANCH{\char"01AB\char"0346} \SUPPORTEDANCH{\char"01AD\char"0346} \SUPPORTEDANCH{\char"0236\char"0346} \SUPPORTEDANCH{\char"0075\char"0346} \MISSINGGLYPH{\char"AB52\char"0346} \SUPPORTEDANCH{\char"0079\char"0346} \SUPPORTEDANCH{\char"01B4\char"0346} \MISSINGGLYPH{\char"1EFF\char"0346}

\SUPPORTEDANCH{\char"0048\char"034A} \SUPPORTEDANCH{\char"2C67\char"034A} \SUPPORTEDANCH{\char"004B\char"034A} \SUPPORTEDANCH{\char"2C69\char"034A} \SUPPORTEDANCH{\char"004D\char"034A} \MISSINGGLYPH{\char"2C6E\char"034A} \SUPPORTEDANCH{\char"0050\char"034A} \SUPPORTEDANCH{\char"01A4\char"034A} \SUPPORTEDANCH{\char"0052\char"034A} \SUPPORTEDANCH{\char"2C64\char"034A} \SUPPORTEDANCH{\char"0054\char"034A} \SUPPORTEDANCH{\char"01AC\char"034A} \SUPPORTEDANCH{\char"01AE\char"034A} \SUPPORTEDANCH{\char"0059\char"034A} \SUPPORTEDANCH{\char"01B3\char"034A} \SUPPORTEDANCH{\char"0061\char"034A} \SUPPORTEDANCH{\char"1E9A\char"034A} \SUPPORTEDANCH{\char"0062\char"034A} \SUPPORTEDANCH{\char"0183\char"034A} \SUPPORTEDANCH{\char"0063\char"034A} \SUPPORTEDANCH{\char"0188\char"034A} \SUPPORTEDANCH{\char"0064\char"034A} \SUPPORTEDANCH{\char"018C\char"034A} \SUPPORTEDANCH{\char"0221\char"034A} \SUPPORTEDANCH{\char"0068\char"034A} \SUPPORTEDANCH{\char"2C68\char"034A} \SUPPORTEDANCH{\char"006B\char"034A} \SUPPORTEDANCH{\char"0199\char"034A} \SUPPORTEDANCH{\char"2C6A\char"034A} \SUPPORTEDANCH{\char"006C\char"034A} \SUPPORTEDANCH{\char"0234\char"034A} \SUPPORTEDANCH{\char"006E\char"034A} \SUPPORTEDANCH{\char"0235\char"034A} \SUPPORTEDANCH{\char"0070\char"034A} \SUPPORTEDANCH{\char"01A5\char"034A} \SUPPORTEDANCH{\char"0071\char"034A} \SUPPORTEDANCH{\char"02A0\char"034A} \SUPPORTEDANCH{\char"0074\char"034A} \SUPPORTEDANCH{\char"01AB\char"034A} \SUPPORTEDANCH{\char"01AD\char"034A} \SUPPORTEDANCH{\char"0236\char"034A} \SUPPORTEDANCH{\char"0075\char"034A} \MISSINGGLYPH{\char"AB52\char"034A} \SUPPORTEDANCH{\char"0079\char"034A} \SUPPORTEDANCH{\char"01B4\char"034A} \MISSINGGLYPH{\char"1EFF\char"034A}

\SUPPORTEDANCH{\char"0048\char"034B} \SUPPORTEDANCH{\char"2C67\char"034B} \SUPPORTEDANCH{\char"004B\char"034B} \SUPPORTEDANCH{\char"2C69\char"034B} \SUPPORTEDANCH{\char"004D\char"034B} \MISSINGGLYPH{\char"2C6E\char"034B} \SUPPORTEDANCH{\char"0050\char"034B} \SUPPORTEDANCH{\char"01A4\char"034B} \SUPPORTEDANCH{\char"0052\char"034B} \SUPPORTEDANCH{\char"2C64\char"034B} \SUPPORTEDANCH{\char"0054\char"034B} \SUPPORTEDANCH{\char"01AC\char"034B} \SUPPORTEDANCH{\char"01AE\char"034B} \SUPPORTEDANCH{\char"0059\char"034B} \SUPPORTEDANCH{\char"01B3\char"034B} \SUPPORTEDANCH{\char"0061\char"034B} \SUPPORTEDANCH{\char"1E9A\char"034B} \SUPPORTEDANCH{\char"0062\char"034B} \SUPPORTEDANCH{\char"0183\char"034B} \SUPPORTEDANCH{\char"0063\char"034B} \SUPPORTEDANCH{\char"0188\char"034B} \SUPPORTEDANCH{\char"0064\char"034B} \SUPPORTEDANCH{\char"018C\char"034B} \SUPPORTEDANCH{\char"0221\char"034B} \SUPPORTEDANCH{\char"0068\char"034B} \SUPPORTEDANCH{\char"2C68\char"034B} \SUPPORTEDANCH{\char"006B\char"034B} \SUPPORTEDANCH{\char"0199\char"034B} \SUPPORTEDANCH{\char"2C6A\char"034B} \SUPPORTEDANCH{\char"006C\char"034B} \SUPPORTEDANCH{\char"0234\char"034B} \SUPPORTEDANCH{\char"006E\char"034B} \SUPPORTEDANCH{\char"0235\char"034B} \SUPPORTEDANCH{\char"0070\char"034B} \SUPPORTEDANCH{\char"01A5\char"034B} \SUPPORTEDANCH{\char"0071\char"034B} \SUPPORTEDANCH{\char"02A0\char"034B} \SUPPORTEDANCH{\char"0074\char"034B} \SUPPORTEDANCH{\char"01AB\char"034B} \SUPPORTEDANCH{\char"01AD\char"034B} \SUPPORTEDANCH{\char"0236\char"034B} \SUPPORTEDANCH{\char"0075\char"034B} \MISSINGGLYPH{\char"AB52\char"034B} \SUPPORTEDANCH{\char"0079\char"034B} \SUPPORTEDANCH{\char"01B4\char"034B} \MISSINGGLYPH{\char"1EFF\char"034B}

\SUPPORTEDANCH{\char"0048\char"034C} \SUPPORTEDANCH{\char"2C67\char"034C} \SUPPORTEDANCH{\char"004B\char"034C} \SUPPORTEDANCH{\char"2C69\char"034C} \SUPPORTEDANCH{\char"004D\char"034C} \MISSINGGLYPH{\char"2C6E\char"034C} \SUPPORTEDANCH{\char"0050\char"034C} \SUPPORTEDANCH{\char"01A4\char"034C} \SUPPORTEDANCH{\char"0052\char"034C} \SUPPORTEDANCH{\char"2C64\char"034C} \SUPPORTEDANCH{\char"0054\char"034C} \SUPPORTEDANCH{\char"01AC\char"034C} \SUPPORTEDANCH{\char"01AE\char"034C} \SUPPORTEDANCH{\char"0059\char"034C} \SUPPORTEDANCH{\char"01B3\char"034C} \SUPPORTEDANCH{\char"0061\char"034C} \SUPPORTEDANCH{\char"1E9A\char"034C} \SUPPORTEDANCH{\char"0062\char"034C} \SUPPORTEDANCH{\char"0183\char"034C} \SUPPORTEDANCH{\char"0063\char"034C} \SUPPORTEDANCH{\char"0188\char"034C} \SUPPORTEDANCH{\char"0064\char"034C} \SUPPORTEDANCH{\char"018C\char"034C} \SUPPORTEDANCH{\char"0221\char"034C} \SUPPORTEDANCH{\char"0068\char"034C} \SUPPORTEDANCH{\char"2C68\char"034C} \SUPPORTEDANCH{\char"006B\char"034C} \SUPPORTEDANCH{\char"0199\char"034C} \SUPPORTEDANCH{\char"2C6A\char"034C} \SUPPORTEDANCH{\char"006C\char"034C} \SUPPORTEDANCH{\char"0234\char"034C} \SUPPORTEDANCH{\char"006E\char"034C} \SUPPORTEDANCH{\char"0235\char"034C} \SUPPORTEDANCH{\char"0070\char"034C} \SUPPORTEDANCH{\char"01A5\char"034C} \SUPPORTEDANCH{\char"0071\char"034C} \SUPPORTEDANCH{\char"02A0\char"034C} \SUPPORTEDANCH{\char"0074\char"034C} \SUPPORTEDANCH{\char"01AB\char"034C} \SUPPORTEDANCH{\char"01AD\char"034C} \SUPPORTEDANCH{\char"0236\char"034C} \SUPPORTEDANCH{\char"0075\char"034C} \MISSINGGLYPH{\char"AB52\char"034C} \SUPPORTEDANCH{\char"0079\char"034C} \SUPPORTEDANCH{\char"01B4\char"034C} \MISSINGGLYPH{\char"1EFF\char"034C}

\SUPPORTEDANCH{\char"0048\char"0350} \SUPPORTEDANCH{\char"2C67\char"0350} \SUPPORTEDANCH{\char"004B\char"0350} \SUPPORTEDANCH{\char"2C69\char"0350} \SUPPORTEDANCH{\char"004D\char"0350} \MISSINGGLYPH{\char"2C6E\char"0350} \SUPPORTEDANCH{\char"0050\char"0350} \SUPPORTEDANCH{\char"01A4\char"0350} \SUPPORTEDANCH{\char"0052\char"0350} \SUPPORTEDANCH{\char"2C64\char"0350} \SUPPORTEDANCH{\char"0054\char"0350} \SUPPORTEDANCH{\char"01AC\char"0350} \SUPPORTEDANCH{\char"01AE\char"0350} \SUPPORTEDANCH{\char"0059\char"0350} \SUPPORTEDANCH{\char"01B3\char"0350} \SUPPORTEDANCH{\char"0061\char"0350} \SUPPORTEDANCH{\char"1E9A\char"0350} \SUPPORTEDANCH{\char"0062\char"0350} \SUPPORTEDANCH{\char"0183\char"0350} \SUPPORTEDANCH{\char"0063\char"0350} \SUPPORTEDANCH{\char"0188\char"0350} \SUPPORTEDANCH{\char"0064\char"0350} \SUPPORTEDANCH{\char"018C\char"0350} \SUPPORTEDANCH{\char"0221\char"0350} \SUPPORTEDANCH{\char"0068\char"0350} \SUPPORTEDANCH{\char"2C68\char"0350} \SUPPORTEDANCH{\char"006B\char"0350} \SUPPORTEDANCH{\char"0199\char"0350} \SUPPORTEDANCH{\char"2C6A\char"0350} \SUPPORTEDANCH{\char"006C\char"0350} \SUPPORTEDANCH{\char"0234\char"0350} \SUPPORTEDANCH{\char"006E\char"0350} \SUPPORTEDANCH{\char"0235\char"0350} \SUPPORTEDANCH{\char"0070\char"0350} \SUPPORTEDANCH{\char"01A5\char"0350} \SUPPORTEDANCH{\char"0071\char"0350} \SUPPORTEDANCH{\char"02A0\char"0350} \SUPPORTEDANCH{\char"0074\char"0350} \SUPPORTEDANCH{\char"01AB\char"0350} \SUPPORTEDANCH{\char"01AD\char"0350} \SUPPORTEDANCH{\char"0236\char"0350} \SUPPORTEDANCH{\char"0075\char"0350} \MISSINGGLYPH{\char"AB52\char"0350} \SUPPORTEDANCH{\char"0079\char"0350} \SUPPORTEDANCH{\char"01B4\char"0350} \MISSINGGLYPH{\char"1EFF\char"0350}

\SUPPORTEDANCH{\char"0048\char"0351} \SUPPORTEDANCH{\char"2C67\char"0351} \SUPPORTEDANCH{\char"004B\char"0351} \SUPPORTEDANCH{\char"2C69\char"0351} \SUPPORTEDANCH{\char"004D\char"0351} \MISSINGGLYPH{\char"2C6E\char"0351} \SUPPORTEDANCH{\char"0050\char"0351} \SUPPORTEDANCH{\char"01A4\char"0351} \SUPPORTEDANCH{\char"0052\char"0351} \SUPPORTEDANCH{\char"2C64\char"0351} \SUPPORTEDANCH{\char"0054\char"0351} \SUPPORTEDANCH{\char"01AC\char"0351} \SUPPORTEDANCH{\char"01AE\char"0351} \SUPPORTEDANCH{\char"0059\char"0351} \SUPPORTEDANCH{\char"01B3\char"0351} \SUPPORTEDANCH{\char"0061\char"0351} \SUPPORTEDANCH{\char"1E9A\char"0351} \SUPPORTEDANCH{\char"0062\char"0351} \SUPPORTEDANCH{\char"0183\char"0351} \SUPPORTEDANCH{\char"0063\char"0351} \SUPPORTEDANCH{\char"0188\char"0351} \SUPPORTEDANCH{\char"0064\char"0351} \SUPPORTEDANCH{\char"018C\char"0351} \SUPPORTEDANCH{\char"0221\char"0351} \SUPPORTEDANCH{\char"0068\char"0351} \SUPPORTEDANCH{\char"2C68\char"0351} \SUPPORTEDANCH{\char"006B\char"0351} \SUPPORTEDANCH{\char"0199\char"0351} \SUPPORTEDANCH{\char"2C6A\char"0351} \SUPPORTEDANCH{\char"006C\char"0351} \SUPPORTEDANCH{\char"0234\char"0351} \SUPPORTEDANCH{\char"006E\char"0351} \SUPPORTEDANCH{\char"0235\char"0351} \SUPPORTEDANCH{\char"0070\char"0351} \SUPPORTEDANCH{\char"01A5\char"0351} \SUPPORTEDANCH{\char"0071\char"0351} \SUPPORTEDANCH{\char"02A0\char"0351} \SUPPORTEDANCH{\char"0074\char"0351} \SUPPORTEDANCH{\char"01AB\char"0351} \SUPPORTEDANCH{\char"01AD\char"0351} \SUPPORTEDANCH{\char"0236\char"0351} \SUPPORTEDANCH{\char"0075\char"0351} \MISSINGGLYPH{\char"AB52\char"0351} \SUPPORTEDANCH{\char"0079\char"0351} \SUPPORTEDANCH{\char"01B4\char"0351} \MISSINGGLYPH{\char"1EFF\char"0351}

\SUPPORTEDANCH{\char"0048\char"0352} \SUPPORTEDANCH{\char"2C67\char"0352} \SUPPORTEDANCH{\char"004B\char"0352} \SUPPORTEDANCH{\char"2C69\char"0352} \SUPPORTEDANCH{\char"004D\char"0352} \MISSINGGLYPH{\char"2C6E\char"0352} \SUPPORTEDANCH{\char"0050\char"0352} \SUPPORTEDANCH{\char"01A4\char"0352} \SUPPORTEDANCH{\char"0052\char"0352} \SUPPORTEDANCH{\char"2C64\char"0352} \SUPPORTEDANCH{\char"0054\char"0352} \SUPPORTEDANCH{\char"01AC\char"0352} \SUPPORTEDANCH{\char"01AE\char"0352} \SUPPORTEDANCH{\char"0059\char"0352} \SUPPORTEDANCH{\char"01B3\char"0352} \SUPPORTEDANCH{\char"0061\char"0352} \SUPPORTEDANCH{\char"1E9A\char"0352} \SUPPORTEDANCH{\char"0062\char"0352} \SUPPORTEDANCH{\char"0183\char"0352} \SUPPORTEDANCH{\char"0063\char"0352} \SUPPORTEDANCH{\char"0188\char"0352} \SUPPORTEDANCH{\char"0064\char"0352} \SUPPORTEDANCH{\char"018C\char"0352} \SUPPORTEDANCH{\char"0221\char"0352} \SUPPORTEDANCH{\char"0068\char"0352} \SUPPORTEDANCH{\char"2C68\char"0352} \SUPPORTEDANCH{\char"006B\char"0352} \SUPPORTEDANCH{\char"0199\char"0352} \SUPPORTEDANCH{\char"2C6A\char"0352} \SUPPORTEDANCH{\char"006C\char"0352} \SUPPORTEDANCH{\char"0234\char"0352} \SUPPORTEDANCH{\char"006E\char"0352} \SUPPORTEDANCH{\char"0235\char"0352} \SUPPORTEDANCH{\char"0070\char"0352} \SUPPORTEDANCH{\char"01A5\char"0352} \SUPPORTEDANCH{\char"0071\char"0352} \SUPPORTEDANCH{\char"02A0\char"0352} \SUPPORTEDANCH{\char"0074\char"0352} \SUPPORTEDANCH{\char"01AB\char"0352} \SUPPORTEDANCH{\char"01AD\char"0352} \SUPPORTEDANCH{\char"0236\char"0352} \SUPPORTEDANCH{\char"0075\char"0352} \MISSINGGLYPH{\char"AB52\char"0352} \SUPPORTEDANCH{\char"0079\char"0352} \SUPPORTEDANCH{\char"01B4\char"0352} \MISSINGGLYPH{\char"1EFF\char"0352}

\SUPPORTEDANCH{\char"0048\char"0357} \SUPPORTEDANCH{\char"2C67\char"0357} \SUPPORTEDANCH{\char"004B\char"0357} \SUPPORTEDANCH{\char"2C69\char"0357} \SUPPORTEDANCH{\char"004D\char"0357} \MISSINGGLYPH{\char"2C6E\char"0357} \SUPPORTEDANCH{\char"0050\char"0357} \SUPPORTEDANCH{\char"01A4\char"0357} \SUPPORTEDANCH{\char"0052\char"0357} \SUPPORTEDANCH{\char"2C64\char"0357} \SUPPORTEDANCH{\char"0054\char"0357} \SUPPORTEDANCH{\char"01AC\char"0357} \SUPPORTEDANCH{\char"01AE\char"0357} \SUPPORTEDANCH{\char"0059\char"0357} \SUPPORTEDANCH{\char"01B3\char"0357} \SUPPORTEDANCH{\char"0061\char"0357} \SUPPORTEDANCH{\char"1E9A\char"0357} \SUPPORTEDANCH{\char"0062\char"0357} \SUPPORTEDANCH{\char"0183\char"0357} \SUPPORTEDANCH{\char"0063\char"0357} \SUPPORTEDANCH{\char"0188\char"0357} \SUPPORTEDANCH{\char"0064\char"0357} \SUPPORTEDANCH{\char"018C\char"0357} \SUPPORTEDANCH{\char"0221\char"0357} \SUPPORTEDANCH{\char"0068\char"0357} \SUPPORTEDANCH{\char"2C68\char"0357} \SUPPORTEDANCH{\char"006B\char"0357} \SUPPORTEDANCH{\char"0199\char"0357} \SUPPORTEDANCH{\char"2C6A\char"0357} \SUPPORTEDANCH{\char"006C\char"0357} \SUPPORTEDANCH{\char"0234\char"0357} \SUPPORTEDANCH{\char"006E\char"0357} \SUPPORTEDANCH{\char"0235\char"0357} \SUPPORTEDANCH{\char"0070\char"0357} \SUPPORTEDANCH{\char"01A5\char"0357} \SUPPORTEDANCH{\char"0071\char"0357} \SUPPORTEDANCH{\char"02A0\char"0357} \SUPPORTEDANCH{\char"0074\char"0357} \SUPPORTEDANCH{\char"01AB\char"0357} \SUPPORTEDANCH{\char"01AD\char"0357} \SUPPORTEDANCH{\char"0236\char"0357} \SUPPORTEDANCH{\char"0075\char"0357} \MISSINGGLYPH{\char"AB52\char"0357} \SUPPORTEDANCH{\char"0079\char"0357} \SUPPORTEDANCH{\char"01B4\char"0357} \MISSINGGLYPH{\char"1EFF\char"0357}

\SUPPORTEDANCH{\char"0048\char"035B} \SUPPORTEDANCH{\char"2C67\char"035B} \SUPPORTEDANCH{\char"004B\char"035B} \SUPPORTEDANCH{\char"2C69\char"035B} \SUPPORTEDANCH{\char"004D\char"035B} \MISSINGGLYPH{\char"2C6E\char"035B} \SUPPORTEDANCH{\char"0050\char"035B} \SUPPORTEDANCH{\char"01A4\char"035B} \SUPPORTEDANCH{\char"0052\char"035B} \SUPPORTEDANCH{\char"2C64\char"035B} \SUPPORTEDANCH{\char"0054\char"035B} \SUPPORTEDANCH{\char"01AC\char"035B} \SUPPORTEDANCH{\char"01AE\char"035B} \SUPPORTEDANCH{\char"0059\char"035B} \SUPPORTEDANCH{\char"01B3\char"035B} \SUPPORTEDANCH{\char"0061\char"035B} \SUPPORTEDANCH{\char"1E9A\char"035B} \SUPPORTEDANCH{\char"0062\char"035B} \SUPPORTEDANCH{\char"0183\char"035B} \SUPPORTEDANCH{\char"0063\char"035B} \SUPPORTEDANCH{\char"0188\char"035B} \SUPPORTEDANCH{\char"0064\char"035B} \SUPPORTEDANCH{\char"018C\char"035B} \SUPPORTEDANCH{\char"0221\char"035B} \SUPPORTEDANCH{\char"0068\char"035B} \SUPPORTEDANCH{\char"2C68\char"035B} \SUPPORTEDANCH{\char"006B\char"035B} \SUPPORTEDANCH{\char"0199\char"035B} \SUPPORTEDANCH{\char"2C6A\char"035B} \SUPPORTEDANCH{\char"006C\char"035B} \SUPPORTEDANCH{\char"0234\char"035B} \SUPPORTEDANCH{\char"006E\char"035B} \SUPPORTEDANCH{\char"0235\char"035B} \SUPPORTEDANCH{\char"0070\char"035B} \SUPPORTEDANCH{\char"01A5\char"035B} \SUPPORTEDANCH{\char"0071\char"035B} \SUPPORTEDANCH{\char"02A0\char"035B} \SUPPORTEDANCH{\char"0074\char"035B} \SUPPORTEDANCH{\char"01AB\char"035B} \SUPPORTEDANCH{\char"01AD\char"035B} \SUPPORTEDANCH{\char"0236\char"035B} \SUPPORTEDANCH{\char"0075\char"035B} \MISSINGGLYPH{\char"AB52\char"035B} \SUPPORTEDANCH{\char"0079\char"035B} \SUPPORTEDANCH{\char"01B4\char"035B} \MISSINGGLYPH{\char"1EFF\char"035B}

}

\newpage
\subsection*{Regular, Group 4, below, Below}

% grid. columns are bases, rows are marks.
\SUPPORTEDANCH{\char"0048\char"0316} \SUPPORTEDFALL{\char"2C67\char"0316} \SUPPORTEDANCH{\char"004B\char"0316} \SUPPORTEDANCH{\char"2C69\char"0316} \SUPPORTEDANCH{\char"004D\char"0316} \SUPPORTEDANCH{\char"2C6E\char"0316} \SUPPORTEDANCH{\char"0050\char"0316} \SUPPORTEDANCH{\char"01A4\char"0316} \SUPPORTEDANCH{\char"0052\char"0316} \SUPPORTEDFALL{\char"2C64\char"0316} \SUPPORTEDANCH{\char"0054\char"0316} \SUPPORTEDFALL{\char"01AC\char"0316} \SUPPORTEDFALL{\char"01AE\char"0316} \SUPPORTEDANCH{\char"0059\char"0316} \SUPPORTEDFALL{\char"01B3\char"0316} \SUPPORTEDANCH{\char"0061\char"0316} \SUPPORTEDFALL{\char"1E9A\char"0316} \SUPPORTEDANCH{\char"0062\char"0316} \SUPPORTEDANCH{\char"0183\char"0316} \SUPPORTEDANCH{\char"0063\char"0316} \SUPPORTEDFALL{\char"0188\char"0316} \SUPPORTEDANCH{\char"0064\char"0316} \SUPPORTEDANCH{\char"018C\char"0316} \SUPPORTEDFALL{\char"0221\char"0316} \SUPPORTEDANCH{\char"0068\char"0316} \SUPPORTEDANCH{\char"2C68\char"0316} \SUPPORTEDANCH{\char"006B\char"0316} \SUPPORTEDANCH{\char"0199\char"0316} \SUPPORTEDANCH{\char"2C6A\char"0316} \SUPPORTEDANCH{\char"006C\char"0316} \SUPPORTEDFALL{\char"0234\char"0316} \SUPPORTEDANCH{\char"006E\char"0316} \SUPPORTEDANCH{\char"0235\char"0316} \SUPPORTEDANCH{\char"0070\char"0316} \SUPPORTEDFALL{\char"01A5\char"0316} \SUPPORTEDANCH{\char"0071\char"0316} \SUPPORTEDANCH{\char"02A0\char"0316} \SUPPORTEDANCH{\char"0074\char"0316} \SUPPORTEDFALL{\char"01AB\char"0316} \SUPPORTEDANCH{\char"01AD\char"0316} \SUPPORTEDANCH{\char"0236\char"0316} \SUPPORTEDANCH{\char"0075\char"0316} \SUPPORTEDANCH{\char"AB52\char"0316} \SUPPORTEDANCH{\char"0079\char"0316} \SUPPORTEDFALL{\char"01B4\char"0316} \SUPPORTEDFALL{\char"1EFF\char"0316}

\SUPPORTEDANCH{\char"0048\char"0317} \SUPPORTEDFALL{\char"2C67\char"0317} \SUPPORTEDANCH{\char"004B\char"0317} \SUPPORTEDANCH{\char"2C69\char"0317} \SUPPORTEDANCH{\char"004D\char"0317} \SUPPORTEDANCH{\char"2C6E\char"0317} \SUPPORTEDANCH{\char"0050\char"0317} \SUPPORTEDANCH{\char"01A4\char"0317} \SUPPORTEDANCH{\char"0052\char"0317} \SUPPORTEDFALL{\char"2C64\char"0317} \SUPPORTEDANCH{\char"0054\char"0317} \SUPPORTEDFALL{\char"01AC\char"0317} \SUPPORTEDFALL{\char"01AE\char"0317} \SUPPORTEDANCH{\char"0059\char"0317} \SUPPORTEDFALL{\char"01B3\char"0317} \SUPPORTEDANCH{\char"0061\char"0317} \SUPPORTEDFALL{\char"1E9A\char"0317} \SUPPORTEDANCH{\char"0062\char"0317} \SUPPORTEDANCH{\char"0183\char"0317} \SUPPORTEDANCH{\char"0063\char"0317} \SUPPORTEDFALL{\char"0188\char"0317} \SUPPORTEDANCH{\char"0064\char"0317} \SUPPORTEDANCH{\char"018C\char"0317} \SUPPORTEDFALL{\char"0221\char"0317} \SUPPORTEDANCH{\char"0068\char"0317} \SUPPORTEDANCH{\char"2C68\char"0317} \SUPPORTEDANCH{\char"006B\char"0317} \SUPPORTEDANCH{\char"0199\char"0317} \SUPPORTEDANCH{\char"2C6A\char"0317} \SUPPORTEDANCH{\char"006C\char"0317} \SUPPORTEDFALL{\char"0234\char"0317} \SUPPORTEDANCH{\char"006E\char"0317} \SUPPORTEDANCH{\char"0235\char"0317} \SUPPORTEDANCH{\char"0070\char"0317} \SUPPORTEDFALL{\char"01A5\char"0317} \SUPPORTEDANCH{\char"0071\char"0317} \SUPPORTEDANCH{\char"02A0\char"0317} \SUPPORTEDANCH{\char"0074\char"0317} \SUPPORTEDFALL{\char"01AB\char"0317} \SUPPORTEDANCH{\char"01AD\char"0317} \SUPPORTEDANCH{\char"0236\char"0317} \SUPPORTEDANCH{\char"0075\char"0317} \SUPPORTEDANCH{\char"AB52\char"0317} \SUPPORTEDANCH{\char"0079\char"0317} \SUPPORTEDFALL{\char"01B4\char"0317} \SUPPORTEDFALL{\char"1EFF\char"0317}

\SUPPORTEDANCH{\char"0048\char"0318} \SUPPORTEDFALL{\char"2C67\char"0318} \SUPPORTEDANCH{\char"004B\char"0318} \SUPPORTEDANCH{\char"2C69\char"0318} \SUPPORTEDANCH{\char"004D\char"0318} \SUPPORTEDANCH{\char"2C6E\char"0318} \SUPPORTEDANCH{\char"0050\char"0318} \SUPPORTEDANCH{\char"01A4\char"0318} \SUPPORTEDANCH{\char"0052\char"0318} \SUPPORTEDFALL{\char"2C64\char"0318} \SUPPORTEDANCH{\char"0054\char"0318} \SUPPORTEDFALL{\char"01AC\char"0318} \SUPPORTEDFALL{\char"01AE\char"0318} \SUPPORTEDANCH{\char"0059\char"0318} \SUPPORTEDFALL{\char"01B3\char"0318} \SUPPORTEDANCH{\char"0061\char"0318} \SUPPORTEDFALL{\char"1E9A\char"0318} \SUPPORTEDANCH{\char"0062\char"0318} \SUPPORTEDANCH{\char"0183\char"0318} \SUPPORTEDANCH{\char"0063\char"0318} \SUPPORTEDFALL{\char"0188\char"0318} \SUPPORTEDANCH{\char"0064\char"0318} \SUPPORTEDANCH{\char"018C\char"0318} \SUPPORTEDFALL{\char"0221\char"0318} \SUPPORTEDANCH{\char"0068\char"0318} \SUPPORTEDANCH{\char"2C68\char"0318} \SUPPORTEDANCH{\char"006B\char"0318} \SUPPORTEDANCH{\char"0199\char"0318} \SUPPORTEDANCH{\char"2C6A\char"0318} \SUPPORTEDANCH{\char"006C\char"0318} \SUPPORTEDFALL{\char"0234\char"0318} \SUPPORTEDANCH{\char"006E\char"0318} \SUPPORTEDANCH{\char"0235\char"0318} \SUPPORTEDANCH{\char"0070\char"0318} \SUPPORTEDFALL{\char"01A5\char"0318} \SUPPORTEDANCH{\char"0071\char"0318} \SUPPORTEDANCH{\char"02A0\char"0318} \SUPPORTEDANCH{\char"0074\char"0318} \SUPPORTEDFALL{\char"01AB\char"0318} \SUPPORTEDANCH{\char"01AD\char"0318} \SUPPORTEDANCH{\char"0236\char"0318} \SUPPORTEDANCH{\char"0075\char"0318} \SUPPORTEDANCH{\char"AB52\char"0318} \SUPPORTEDANCH{\char"0079\char"0318} \SUPPORTEDFALL{\char"01B4\char"0318} \SUPPORTEDFALL{\char"1EFF\char"0318}

\SUPPORTEDANCH{\char"0048\char"0319} \SUPPORTEDFALL{\char"2C67\char"0319} \SUPPORTEDANCH{\char"004B\char"0319} \SUPPORTEDANCH{\char"2C69\char"0319} \SUPPORTEDANCH{\char"004D\char"0319} \SUPPORTEDANCH{\char"2C6E\char"0319} \SUPPORTEDANCH{\char"0050\char"0319} \SUPPORTEDANCH{\char"01A4\char"0319} \SUPPORTEDANCH{\char"0052\char"0319} \SUPPORTEDFALL{\char"2C64\char"0319} \SUPPORTEDANCH{\char"0054\char"0319} \SUPPORTEDFALL{\char"01AC\char"0319} \SUPPORTEDFALL{\char"01AE\char"0319} \SUPPORTEDANCH{\char"0059\char"0319} \SUPPORTEDFALL{\char"01B3\char"0319} \SUPPORTEDANCH{\char"0061\char"0319} \SUPPORTEDFALL{\char"1E9A\char"0319} \SUPPORTEDANCH{\char"0062\char"0319} \SUPPORTEDANCH{\char"0183\char"0319} \SUPPORTEDANCH{\char"0063\char"0319} \SUPPORTEDFALL{\char"0188\char"0319} \SUPPORTEDANCH{\char"0064\char"0319} \SUPPORTEDANCH{\char"018C\char"0319} \SUPPORTEDFALL{\char"0221\char"0319} \SUPPORTEDANCH{\char"0068\char"0319} \SUPPORTEDANCH{\char"2C68\char"0319} \SUPPORTEDANCH{\char"006B\char"0319} \SUPPORTEDANCH{\char"0199\char"0319} \SUPPORTEDANCH{\char"2C6A\char"0319} \SUPPORTEDANCH{\char"006C\char"0319} \SUPPORTEDFALL{\char"0234\char"0319} \SUPPORTEDANCH{\char"006E\char"0319} \SUPPORTEDANCH{\char"0235\char"0319} \SUPPORTEDANCH{\char"0070\char"0319} \SUPPORTEDFALL{\char"01A5\char"0319} \SUPPORTEDANCH{\char"0071\char"0319} \SUPPORTEDANCH{\char"02A0\char"0319} \SUPPORTEDANCH{\char"0074\char"0319} \SUPPORTEDFALL{\char"01AB\char"0319} \SUPPORTEDANCH{\char"01AD\char"0319} \SUPPORTEDANCH{\char"0236\char"0319} \SUPPORTEDANCH{\char"0075\char"0319} \SUPPORTEDANCH{\char"AB52\char"0319} \SUPPORTEDANCH{\char"0079\char"0319} \SUPPORTEDFALL{\char"01B4\char"0319} \SUPPORTEDFALL{\char"1EFF\char"0319}

\SUPPORTEDANCH{\char"0048\char"031C} \SUPPORTEDFALL{\char"2C67\char"031C} \SUPPORTEDANCH{\char"004B\char"031C} \SUPPORTEDANCH{\char"2C69\char"031C} \SUPPORTEDANCH{\char"004D\char"031C} \SUPPORTEDANCH{\char"2C6E\char"031C} \SUPPORTEDANCH{\char"0050\char"031C} \SUPPORTEDANCH{\char"01A4\char"031C} \SUPPORTEDANCH{\char"0052\char"031C} \SUPPORTEDFALL{\char"2C64\char"031C} \SUPPORTEDANCH{\char"0054\char"031C} \SUPPORTEDFALL{\char"01AC\char"031C} \SUPPORTEDFALL{\char"01AE\char"031C} \SUPPORTEDANCH{\char"0059\char"031C} \SUPPORTEDFALL{\char"01B3\char"031C} \SUPPORTEDANCH{\char"0061\char"031C} \SUPPORTEDFALL{\char"1E9A\char"031C} \SUPPORTEDANCH{\char"0062\char"031C} \SUPPORTEDANCH{\char"0183\char"031C} \SUPPORTEDANCH{\char"0063\char"031C} \SUPPORTEDFALL{\char"0188\char"031C} \SUPPORTEDANCH{\char"0064\char"031C} \SUPPORTEDANCH{\char"018C\char"031C} \SUPPORTEDFALL{\char"0221\char"031C} \SUPPORTEDANCH{\char"0068\char"031C} \SUPPORTEDANCH{\char"2C68\char"031C} \SUPPORTEDANCH{\char"006B\char"031C} \SUPPORTEDANCH{\char"0199\char"031C} \SUPPORTEDANCH{\char"2C6A\char"031C} \SUPPORTEDANCH{\char"006C\char"031C} \SUPPORTEDFALL{\char"0234\char"031C} \SUPPORTEDANCH{\char"006E\char"031C} \SUPPORTEDANCH{\char"0235\char"031C} \SUPPORTEDANCH{\char"0070\char"031C} \SUPPORTEDFALL{\char"01A5\char"031C} \SUPPORTEDANCH{\char"0071\char"031C} \SUPPORTEDANCH{\char"02A0\char"031C} \SUPPORTEDANCH{\char"0074\char"031C} \SUPPORTEDFALL{\char"01AB\char"031C} \SUPPORTEDANCH{\char"01AD\char"031C} \SUPPORTEDANCH{\char"0236\char"031C} \SUPPORTEDANCH{\char"0075\char"031C} \SUPPORTEDANCH{\char"AB52\char"031C} \SUPPORTEDANCH{\char"0079\char"031C} \SUPPORTEDFALL{\char"01B4\char"031C} \SUPPORTEDFALL{\char"1EFF\char"031C}

\SUPPORTEDANCH{\char"0048\char"0320} \SUPPORTEDFALL{\char"2C67\char"0320} \SUPPORTEDANCH{\char"004B\char"0320} \SUPPORTEDANCH{\char"2C69\char"0320} \SUPPORTEDANCH{\char"004D\char"0320} \SUPPORTEDANCH{\char"2C6E\char"0320} \SUPPORTEDANCH{\char"0050\char"0320} \SUPPORTEDANCH{\char"01A4\char"0320} \SUPPORTEDANCH{\char"0052\char"0320} \SUPPORTEDFALL{\char"2C64\char"0320} \SUPPORTEDANCH{\char"0054\char"0320} \SUPPORTEDFALL{\char"01AC\char"0320} \SUPPORTEDFALL{\char"01AE\char"0320} \SUPPORTEDANCH{\char"0059\char"0320} \SUPPORTEDFALL{\char"01B3\char"0320} \SUPPORTEDANCH{\char"0061\char"0320} \SUPPORTEDFALL{\char"1E9A\char"0320} \SUPPORTEDANCH{\char"0062\char"0320} \SUPPORTEDANCH{\char"0183\char"0320} \SUPPORTEDANCH{\char"0063\char"0320} \SUPPORTEDFALL{\char"0188\char"0320} \SUPPORTEDANCH{\char"0064\char"0320} \SUPPORTEDANCH{\char"018C\char"0320} \SUPPORTEDFALL{\char"0221\char"0320} \SUPPORTEDANCH{\char"0068\char"0320} \SUPPORTEDANCH{\char"2C68\char"0320} \SUPPORTEDANCH{\char"006B\char"0320} \SUPPORTEDANCH{\char"0199\char"0320} \SUPPORTEDANCH{\char"2C6A\char"0320} \SUPPORTEDANCH{\char"006C\char"0320} \SUPPORTEDFALL{\char"0234\char"0320} \SUPPORTEDANCH{\char"006E\char"0320} \SUPPORTEDANCH{\char"0235\char"0320} \SUPPORTEDANCH{\char"0070\char"0320} \SUPPORTEDFALL{\char"01A5\char"0320} \SUPPORTEDANCH{\char"0071\char"0320} \SUPPORTEDANCH{\char"02A0\char"0320} \SUPPORTEDANCH{\char"0074\char"0320} \SUPPORTEDFALL{\char"01AB\char"0320} \SUPPORTEDANCH{\char"01AD\char"0320} \SUPPORTEDANCH{\char"0236\char"0320} \SUPPORTEDANCH{\char"0075\char"0320} \SUPPORTEDANCH{\char"AB52\char"0320} \SUPPORTEDANCH{\char"0079\char"0320} \SUPPORTEDFALL{\char"01B4\char"0320} \SUPPORTEDFALL{\char"1EFF\char"0320}

\SUPPORTEDPREC{\char"0048\char"0323} \SUPPORTEDFALL{\char"2C67\char"0323} \SUPPORTEDPREC{\char"004B\char"0323} \SUPPORTEDANCH{\char"2C69\char"0323} \SUPPORTEDPREC{\char"004D\char"0323} \SUPPORTEDANCH{\char"2C6E\char"0323} \SUPPORTEDANCH{\char"0050\char"0323} \SUPPORTEDANCH{\char"01A4\char"0323} \SUPPORTEDPREC{\char"0052\char"0323} \SUPPORTEDFALL{\char"2C64\char"0323} \SUPPORTEDPREC{\char"0054\char"0323} \SUPPORTEDFALL{\char"01AC\char"0323} \SUPPORTEDFALL{\char"01AE\char"0323} \SUPPORTEDPREC{\char"0059\char"0323} \SUPPORTEDFALL{\char"01B3\char"0323} \SUPPORTEDPREC{\char"0061\char"0323} \SUPPORTEDFALL{\char"1E9A\char"0323} \SUPPORTEDPREC{\char"0062\char"0323} \SUPPORTEDANCH{\char"0183\char"0323} \SUPPORTEDANCH{\char"0063\char"0323} \SUPPORTEDFALL{\char"0188\char"0323} \SUPPORTEDPREC{\char"0064\char"0323} \SUPPORTEDANCH{\char"018C\char"0323} \SUPPORTEDFALL{\char"0221\char"0323} \SUPPORTEDPREC{\char"0068\char"0323} \SUPPORTEDANCH{\char"2C68\char"0323} \SUPPORTEDPREC{\char"006B\char"0323} \SUPPORTEDANCH{\char"0199\char"0323} \SUPPORTEDANCH{\char"2C6A\char"0323} \SUPPORTEDPREC{\char"006C\char"0323} \SUPPORTEDFALL{\char"0234\char"0323} \SUPPORTEDPREC{\char"006E\char"0323} \SUPPORTEDANCH{\char"0235\char"0323} \SUPPORTEDANCH{\char"0070\char"0323} \SUPPORTEDFALL{\char"01A5\char"0323} \SUPPORTEDANCH{\char"0071\char"0323} \SUPPORTEDANCH{\char"02A0\char"0323} \SUPPORTEDPREC{\char"0074\char"0323} \SUPPORTEDFALL{\char"01AB\char"0323} \SUPPORTEDANCH{\char"01AD\char"0323} \SUPPORTEDANCH{\char"0236\char"0323} \SUPPORTEDPREC{\char"0075\char"0323} \SUPPORTEDANCH{\char"AB52\char"0323} \SUPPORTEDPREC{\char"0079\char"0323} \SUPPORTEDFALL{\char"01B4\char"0323} \SUPPORTEDFALL{\char"1EFF\char"0323}

\SUPPORTEDANCH{\char"0048\char"0324} \SUPPORTEDFALL{\char"2C67\char"0324} \SUPPORTEDANCH{\char"004B\char"0324} \SUPPORTEDANCH{\char"2C69\char"0324} \SUPPORTEDANCH{\char"004D\char"0324} \SUPPORTEDANCH{\char"2C6E\char"0324} \SUPPORTEDANCH{\char"0050\char"0324} \SUPPORTEDANCH{\char"01A4\char"0324} \SUPPORTEDANCH{\char"0052\char"0324} \SUPPORTEDFALL{\char"2C64\char"0324} \SUPPORTEDANCH{\char"0054\char"0324} \SUPPORTEDFALL{\char"01AC\char"0324} \SUPPORTEDFALL{\char"01AE\char"0324} \SUPPORTEDANCH{\char"0059\char"0324} \SUPPORTEDFALL{\char"01B3\char"0324} \SUPPORTEDANCH{\char"0061\char"0324} \SUPPORTEDFALL{\char"1E9A\char"0324} \SUPPORTEDANCH{\char"0062\char"0324} \SUPPORTEDANCH{\char"0183\char"0324} \SUPPORTEDANCH{\char"0063\char"0324} \SUPPORTEDFALL{\char"0188\char"0324} \SUPPORTEDANCH{\char"0064\char"0324} \SUPPORTEDANCH{\char"018C\char"0324} \SUPPORTEDFALL{\char"0221\char"0324} \SUPPORTEDANCH{\char"0068\char"0324} \SUPPORTEDANCH{\char"2C68\char"0324} \SUPPORTEDANCH{\char"006B\char"0324} \SUPPORTEDANCH{\char"0199\char"0324} \SUPPORTEDANCH{\char"2C6A\char"0324} \SUPPORTEDANCH{\char"006C\char"0324} \SUPPORTEDFALL{\char"0234\char"0324} \SUPPORTEDANCH{\char"006E\char"0324} \SUPPORTEDANCH{\char"0235\char"0324} \SUPPORTEDANCH{\char"0070\char"0324} \SUPPORTEDFALL{\char"01A5\char"0324} \SUPPORTEDANCH{\char"0071\char"0324} \SUPPORTEDANCH{\char"02A0\char"0324} \SUPPORTEDANCH{\char"0074\char"0324} \SUPPORTEDFALL{\char"01AB\char"0324} \SUPPORTEDANCH{\char"01AD\char"0324} \SUPPORTEDANCH{\char"0236\char"0324} \SUPPORTEDPREC{\char"0075\char"0324} \SUPPORTEDANCH{\char"AB52\char"0324} \SUPPORTEDANCH{\char"0079\char"0324} \SUPPORTEDFALL{\char"01B4\char"0324} \SUPPORTEDFALL{\char"1EFF\char"0324}

\SUPPORTEDANCH{\char"0048\char"0325} \SUPPORTEDFALL{\char"2C67\char"0325} \SUPPORTEDANCH{\char"004B\char"0325} \SUPPORTEDANCH{\char"2C69\char"0325} \SUPPORTEDANCH{\char"004D\char"0325} \SUPPORTEDANCH{\char"2C6E\char"0325} \SUPPORTEDANCH{\char"0050\char"0325} \SUPPORTEDANCH{\char"01A4\char"0325} \SUPPORTEDANCH{\char"0052\char"0325} \SUPPORTEDFALL{\char"2C64\char"0325} \SUPPORTEDANCH{\char"0054\char"0325} \SUPPORTEDFALL{\char"01AC\char"0325} \SUPPORTEDFALL{\char"01AE\char"0325} \SUPPORTEDANCH{\char"0059\char"0325} \SUPPORTEDFALL{\char"01B3\char"0325} \SUPPORTEDPREC{\char"0061\char"0325} \SUPPORTEDFALL{\char"1E9A\char"0325} \SUPPORTEDANCH{\char"0062\char"0325} \SUPPORTEDANCH{\char"0183\char"0325} \SUPPORTEDANCH{\char"0063\char"0325} \SUPPORTEDFALL{\char"0188\char"0325} \SUPPORTEDANCH{\char"0064\char"0325} \SUPPORTEDANCH{\char"018C\char"0325} \SUPPORTEDFALL{\char"0221\char"0325} \SUPPORTEDANCH{\char"0068\char"0325} \SUPPORTEDANCH{\char"2C68\char"0325} \SUPPORTEDANCH{\char"006B\char"0325} \SUPPORTEDANCH{\char"0199\char"0325} \SUPPORTEDANCH{\char"2C6A\char"0325} \SUPPORTEDANCH{\char"006C\char"0325} \SUPPORTEDFALL{\char"0234\char"0325} \SUPPORTEDANCH{\char"006E\char"0325} \SUPPORTEDANCH{\char"0235\char"0325} \SUPPORTEDANCH{\char"0070\char"0325} \SUPPORTEDFALL{\char"01A5\char"0325} \SUPPORTEDANCH{\char"0071\char"0325} \SUPPORTEDANCH{\char"02A0\char"0325} \SUPPORTEDANCH{\char"0074\char"0325} \SUPPORTEDFALL{\char"01AB\char"0325} \SUPPORTEDANCH{\char"01AD\char"0325} \SUPPORTEDANCH{\char"0236\char"0325} \SUPPORTEDANCH{\char"0075\char"0325} \SUPPORTEDANCH{\char"AB52\char"0325} \SUPPORTEDANCH{\char"0079\char"0325} \SUPPORTEDFALL{\char"01B4\char"0325} \SUPPORTEDFALL{\char"1EFF\char"0325}

\SUPPORTEDANCH{\char"0048\char"0326} \SUPPORTEDFALL{\char"2C67\char"0326} \SUPPORTEDANCH{\char"004B\char"0326} \SUPPORTEDANCH{\char"2C69\char"0326} \SUPPORTEDANCH{\char"004D\char"0326} \SUPPORTEDANCH{\char"2C6E\char"0326} \SUPPORTEDANCH{\char"0050\char"0326} \SUPPORTEDANCH{\char"01A4\char"0326} \SUPPORTEDANCH{\char"0052\char"0326} \SUPPORTEDFALL{\char"2C64\char"0326} \SUPPORTEDPREC{\char"0054\char"0326} \SUPPORTEDFALL{\char"01AC\char"0326} \SUPPORTEDFALL{\char"01AE\char"0326} \SUPPORTEDANCH{\char"0059\char"0326} \SUPPORTEDFALL{\char"01B3\char"0326} \SUPPORTEDANCH{\char"0061\char"0326} \SUPPORTEDFALL{\char"1E9A\char"0326} \SUPPORTEDANCH{\char"0062\char"0326} \SUPPORTEDANCH{\char"0183\char"0326} \SUPPORTEDANCH{\char"0063\char"0326} \SUPPORTEDFALL{\char"0188\char"0326} \SUPPORTEDANCH{\char"0064\char"0326} \SUPPORTEDANCH{\char"018C\char"0326} \SUPPORTEDFALL{\char"0221\char"0326} \SUPPORTEDANCH{\char"0068\char"0326} \SUPPORTEDANCH{\char"2C68\char"0326} \SUPPORTEDANCH{\char"006B\char"0326} \SUPPORTEDANCH{\char"0199\char"0326} \SUPPORTEDANCH{\char"2C6A\char"0326} \SUPPORTEDANCH{\char"006C\char"0326} \SUPPORTEDFALL{\char"0234\char"0326} \SUPPORTEDANCH{\char"006E\char"0326} \SUPPORTEDANCH{\char"0235\char"0326} \SUPPORTEDANCH{\char"0070\char"0326} \SUPPORTEDFALL{\char"01A5\char"0326} \SUPPORTEDANCH{\char"0071\char"0326} \SUPPORTEDANCH{\char"02A0\char"0326} \SUPPORTEDPREC{\char"0074\char"0326} \SUPPORTEDFALL{\char"01AB\char"0326} \SUPPORTEDANCH{\char"01AD\char"0326} \SUPPORTEDANCH{\char"0236\char"0326} \SUPPORTEDANCH{\char"0075\char"0326} \SUPPORTEDANCH{\char"AB52\char"0326} \SUPPORTEDANCH{\char"0079\char"0326} \SUPPORTEDFALL{\char"01B4\char"0326} \SUPPORTEDFALL{\char"1EFF\char"0326}

\SUPPORTEDANCH{\char"0048\char"0329} \SUPPORTEDFALL{\char"2C67\char"0329} \SUPPORTEDANCH{\char"004B\char"0329} \SUPPORTEDANCH{\char"2C69\char"0329} \SUPPORTEDANCH{\char"004D\char"0329} \SUPPORTEDANCH{\char"2C6E\char"0329} \SUPPORTEDANCH{\char"0050\char"0329} \SUPPORTEDANCH{\char"01A4\char"0329} \SUPPORTEDANCH{\char"0052\char"0329} \SUPPORTEDFALL{\char"2C64\char"0329} \SUPPORTEDANCH{\char"0054\char"0329} \SUPPORTEDFALL{\char"01AC\char"0329} \SUPPORTEDFALL{\char"01AE\char"0329} \SUPPORTEDANCH{\char"0059\char"0329} \SUPPORTEDFALL{\char"01B3\char"0329} \SUPPORTEDANCH{\char"0061\char"0329} \SUPPORTEDFALL{\char"1E9A\char"0329} \SUPPORTEDANCH{\char"0062\char"0329} \SUPPORTEDANCH{\char"0183\char"0329} \SUPPORTEDANCH{\char"0063\char"0329} \SUPPORTEDFALL{\char"0188\char"0329} \SUPPORTEDANCH{\char"0064\char"0329} \SUPPORTEDANCH{\char"018C\char"0329} \SUPPORTEDFALL{\char"0221\char"0329} \SUPPORTEDANCH{\char"0068\char"0329} \SUPPORTEDANCH{\char"2C68\char"0329} \SUPPORTEDANCH{\char"006B\char"0329} \SUPPORTEDANCH{\char"0199\char"0329} \SUPPORTEDANCH{\char"2C6A\char"0329} \SUPPORTEDANCH{\char"006C\char"0329} \SUPPORTEDFALL{\char"0234\char"0329} \SUPPORTEDANCH{\char"006E\char"0329} \SUPPORTEDANCH{\char"0235\char"0329} \SUPPORTEDANCH{\char"0070\char"0329} \SUPPORTEDFALL{\char"01A5\char"0329} \SUPPORTEDANCH{\char"0071\char"0329} \SUPPORTEDANCH{\char"02A0\char"0329} \SUPPORTEDANCH{\char"0074\char"0329} \SUPPORTEDFALL{\char"01AB\char"0329} \SUPPORTEDANCH{\char"01AD\char"0329} \SUPPORTEDANCH{\char"0236\char"0329} \SUPPORTEDANCH{\char"0075\char"0329} \SUPPORTEDANCH{\char"AB52\char"0329} \SUPPORTEDANCH{\char"0079\char"0329} \SUPPORTEDFALL{\char"01B4\char"0329} \SUPPORTEDFALL{\char"1EFF\char"0329}

\SUPPORTEDANCH{\char"0048\char"032A} \SUPPORTEDFALL{\char"2C67\char"032A} \SUPPORTEDANCH{\char"004B\char"032A} \SUPPORTEDANCH{\char"2C69\char"032A} \SUPPORTEDANCH{\char"004D\char"032A} \SUPPORTEDANCH{\char"2C6E\char"032A} \SUPPORTEDANCH{\char"0050\char"032A} \SUPPORTEDANCH{\char"01A4\char"032A} \SUPPORTEDANCH{\char"0052\char"032A} \SUPPORTEDFALL{\char"2C64\char"032A} \SUPPORTEDANCH{\char"0054\char"032A} \SUPPORTEDFALL{\char"01AC\char"032A} \SUPPORTEDFALL{\char"01AE\char"032A} \SUPPORTEDANCH{\char"0059\char"032A} \SUPPORTEDFALL{\char"01B3\char"032A} \SUPPORTEDANCH{\char"0061\char"032A} \SUPPORTEDFALL{\char"1E9A\char"032A} \SUPPORTEDANCH{\char"0062\char"032A} \SUPPORTEDANCH{\char"0183\char"032A} \SUPPORTEDANCH{\char"0063\char"032A} \SUPPORTEDFALL{\char"0188\char"032A} \SUPPORTEDANCH{\char"0064\char"032A} \SUPPORTEDANCH{\char"018C\char"032A} \SUPPORTEDFALL{\char"0221\char"032A} \SUPPORTEDANCH{\char"0068\char"032A} \SUPPORTEDANCH{\char"2C68\char"032A} \SUPPORTEDANCH{\char"006B\char"032A} \SUPPORTEDANCH{\char"0199\char"032A} \SUPPORTEDANCH{\char"2C6A\char"032A} \SUPPORTEDANCH{\char"006C\char"032A} \SUPPORTEDFALL{\char"0234\char"032A} \SUPPORTEDANCH{\char"006E\char"032A} \SUPPORTEDANCH{\char"0235\char"032A} \SUPPORTEDANCH{\char"0070\char"032A} \SUPPORTEDFALL{\char"01A5\char"032A} \SUPPORTEDANCH{\char"0071\char"032A} \SUPPORTEDANCH{\char"02A0\char"032A} \SUPPORTEDANCH{\char"0074\char"032A} \SUPPORTEDFALL{\char"01AB\char"032A} \SUPPORTEDANCH{\char"01AD\char"032A} \SUPPORTEDANCH{\char"0236\char"032A} \SUPPORTEDANCH{\char"0075\char"032A} \SUPPORTEDANCH{\char"AB52\char"032A} \SUPPORTEDANCH{\char"0079\char"032A} \SUPPORTEDFALL{\char"01B4\char"032A} \SUPPORTEDFALL{\char"1EFF\char"032A}

\SUPPORTEDANCH{\char"0048\char"032B} \SUPPORTEDFALL{\char"2C67\char"032B} \SUPPORTEDANCH{\char"004B\char"032B} \SUPPORTEDANCH{\char"2C69\char"032B} \SUPPORTEDANCH{\char"004D\char"032B} \SUPPORTEDANCH{\char"2C6E\char"032B} \SUPPORTEDANCH{\char"0050\char"032B} \SUPPORTEDANCH{\char"01A4\char"032B} \SUPPORTEDANCH{\char"0052\char"032B} \SUPPORTEDFALL{\char"2C64\char"032B} \SUPPORTEDANCH{\char"0054\char"032B} \SUPPORTEDFALL{\char"01AC\char"032B} \SUPPORTEDFALL{\char"01AE\char"032B} \SUPPORTEDANCH{\char"0059\char"032B} \SUPPORTEDFALL{\char"01B3\char"032B} \SUPPORTEDANCH{\char"0061\char"032B} \SUPPORTEDFALL{\char"1E9A\char"032B} \SUPPORTEDANCH{\char"0062\char"032B} \SUPPORTEDANCH{\char"0183\char"032B} \SUPPORTEDANCH{\char"0063\char"032B} \SUPPORTEDFALL{\char"0188\char"032B} \SUPPORTEDANCH{\char"0064\char"032B} \SUPPORTEDANCH{\char"018C\char"032B} \SUPPORTEDFALL{\char"0221\char"032B} \SUPPORTEDANCH{\char"0068\char"032B} \SUPPORTEDANCH{\char"2C68\char"032B} \SUPPORTEDANCH{\char"006B\char"032B} \SUPPORTEDANCH{\char"0199\char"032B} \SUPPORTEDANCH{\char"2C6A\char"032B} \SUPPORTEDANCH{\char"006C\char"032B} \SUPPORTEDFALL{\char"0234\char"032B} \SUPPORTEDANCH{\char"006E\char"032B} \SUPPORTEDANCH{\char"0235\char"032B} \SUPPORTEDANCH{\char"0070\char"032B} \SUPPORTEDFALL{\char"01A5\char"032B} \SUPPORTEDANCH{\char"0071\char"032B} \SUPPORTEDANCH{\char"02A0\char"032B} \SUPPORTEDANCH{\char"0074\char"032B} \SUPPORTEDFALL{\char"01AB\char"032B} \SUPPORTEDANCH{\char"01AD\char"032B} \SUPPORTEDANCH{\char"0236\char"032B} \SUPPORTEDANCH{\char"0075\char"032B} \SUPPORTEDANCH{\char"AB52\char"032B} \SUPPORTEDANCH{\char"0079\char"032B} \SUPPORTEDFALL{\char"01B4\char"032B} \SUPPORTEDFALL{\char"1EFF\char"032B}

\SUPPORTEDANCH{\char"0048\char"032C} \SUPPORTEDFALL{\char"2C67\char"032C} \SUPPORTEDANCH{\char"004B\char"032C} \SUPPORTEDANCH{\char"2C69\char"032C} \SUPPORTEDANCH{\char"004D\char"032C} \SUPPORTEDANCH{\char"2C6E\char"032C} \SUPPORTEDANCH{\char"0050\char"032C} \SUPPORTEDANCH{\char"01A4\char"032C} \SUPPORTEDANCH{\char"0052\char"032C} \SUPPORTEDFALL{\char"2C64\char"032C} \SUPPORTEDANCH{\char"0054\char"032C} \SUPPORTEDFALL{\char"01AC\char"032C} \SUPPORTEDFALL{\char"01AE\char"032C} \SUPPORTEDANCH{\char"0059\char"032C} \SUPPORTEDFALL{\char"01B3\char"032C} \SUPPORTEDANCH{\char"0061\char"032C} \SUPPORTEDFALL{\char"1E9A\char"032C} \SUPPORTEDANCH{\char"0062\char"032C} \SUPPORTEDANCH{\char"0183\char"032C} \SUPPORTEDANCH{\char"0063\char"032C} \SUPPORTEDFALL{\char"0188\char"032C} \SUPPORTEDANCH{\char"0064\char"032C} \SUPPORTEDANCH{\char"018C\char"032C} \SUPPORTEDFALL{\char"0221\char"032C} \SUPPORTEDANCH{\char"0068\char"032C} \SUPPORTEDANCH{\char"2C68\char"032C} \SUPPORTEDANCH{\char"006B\char"032C} \SUPPORTEDANCH{\char"0199\char"032C} \SUPPORTEDANCH{\char"2C6A\char"032C} \SUPPORTEDANCH{\char"006C\char"032C} \SUPPORTEDFALL{\char"0234\char"032C} \SUPPORTEDANCH{\char"006E\char"032C} \SUPPORTEDANCH{\char"0235\char"032C} \SUPPORTEDANCH{\char"0070\char"032C} \SUPPORTEDFALL{\char"01A5\char"032C} \SUPPORTEDANCH{\char"0071\char"032C} \SUPPORTEDANCH{\char"02A0\char"032C} \SUPPORTEDANCH{\char"0074\char"032C} \SUPPORTEDFALL{\char"01AB\char"032C} \SUPPORTEDANCH{\char"01AD\char"032C} \SUPPORTEDANCH{\char"0236\char"032C} \SUPPORTEDANCH{\char"0075\char"032C} \SUPPORTEDANCH{\char"AB52\char"032C} \SUPPORTEDANCH{\char"0079\char"032C} \SUPPORTEDFALL{\char"01B4\char"032C} \SUPPORTEDFALL{\char"1EFF\char"032C}

\SUPPORTEDANCH{\char"0048\char"032D} \SUPPORTEDFALL{\char"2C67\char"032D} \SUPPORTEDANCH{\char"004B\char"032D} \SUPPORTEDANCH{\char"2C69\char"032D} \SUPPORTEDANCH{\char"004D\char"032D} \SUPPORTEDANCH{\char"2C6E\char"032D} \SUPPORTEDANCH{\char"0050\char"032D} \SUPPORTEDANCH{\char"01A4\char"032D} \SUPPORTEDANCH{\char"0052\char"032D} \SUPPORTEDFALL{\char"2C64\char"032D} \SUPPORTEDPREC{\char"0054\char"032D} \SUPPORTEDFALL{\char"01AC\char"032D} \SUPPORTEDFALL{\char"01AE\char"032D} \SUPPORTEDANCH{\char"0059\char"032D} \SUPPORTEDFALL{\char"01B3\char"032D} \SUPPORTEDANCH{\char"0061\char"032D} \SUPPORTEDFALL{\char"1E9A\char"032D} \SUPPORTEDANCH{\char"0062\char"032D} \SUPPORTEDANCH{\char"0183\char"032D} \SUPPORTEDANCH{\char"0063\char"032D} \SUPPORTEDFALL{\char"0188\char"032D} \SUPPORTEDPREC{\char"0064\char"032D} \SUPPORTEDANCH{\char"018C\char"032D} \SUPPORTEDFALL{\char"0221\char"032D} \SUPPORTEDANCH{\char"0068\char"032D} \SUPPORTEDANCH{\char"2C68\char"032D} \SUPPORTEDANCH{\char"006B\char"032D} \SUPPORTEDANCH{\char"0199\char"032D} \SUPPORTEDANCH{\char"2C6A\char"032D} \SUPPORTEDPREC{\char"006C\char"032D} \SUPPORTEDFALL{\char"0234\char"032D} \SUPPORTEDPREC{\char"006E\char"032D} \SUPPORTEDANCH{\char"0235\char"032D} \SUPPORTEDANCH{\char"0070\char"032D} \SUPPORTEDFALL{\char"01A5\char"032D} \SUPPORTEDANCH{\char"0071\char"032D} \SUPPORTEDANCH{\char"02A0\char"032D} \SUPPORTEDPREC{\char"0074\char"032D} \SUPPORTEDFALL{\char"01AB\char"032D} \SUPPORTEDANCH{\char"01AD\char"032D} \SUPPORTEDANCH{\char"0236\char"032D} \SUPPORTEDPREC{\char"0075\char"032D} \SUPPORTEDANCH{\char"AB52\char"032D} \SUPPORTEDANCH{\char"0079\char"032D} \SUPPORTEDFALL{\char"01B4\char"032D} \SUPPORTEDFALL{\char"1EFF\char"032D}

\SUPPORTEDPREC{\char"0048\char"032E} \SUPPORTEDFALL{\char"2C67\char"032E} \SUPPORTEDANCH{\char"004B\char"032E} \SUPPORTEDANCH{\char"2C69\char"032E} \SUPPORTEDANCH{\char"004D\char"032E} \SUPPORTEDANCH{\char"2C6E\char"032E} \SUPPORTEDANCH{\char"0050\char"032E} \SUPPORTEDANCH{\char"01A4\char"032E} \SUPPORTEDANCH{\char"0052\char"032E} \SUPPORTEDFALL{\char"2C64\char"032E} \SUPPORTEDANCH{\char"0054\char"032E} \SUPPORTEDFALL{\char"01AC\char"032E} \SUPPORTEDFALL{\char"01AE\char"032E} \SUPPORTEDANCH{\char"0059\char"032E} \SUPPORTEDFALL{\char"01B3\char"032E} \SUPPORTEDANCH{\char"0061\char"032E} \SUPPORTEDFALL{\char"1E9A\char"032E} \SUPPORTEDANCH{\char"0062\char"032E} \SUPPORTEDANCH{\char"0183\char"032E} \SUPPORTEDANCH{\char"0063\char"032E} \SUPPORTEDFALL{\char"0188\char"032E} \SUPPORTEDANCH{\char"0064\char"032E} \SUPPORTEDANCH{\char"018C\char"032E} \SUPPORTEDFALL{\char"0221\char"032E} \SUPPORTEDPREC{\char"0068\char"032E} \SUPPORTEDANCH{\char"2C68\char"032E} \SUPPORTEDANCH{\char"006B\char"032E} \SUPPORTEDANCH{\char"0199\char"032E} \SUPPORTEDANCH{\char"2C6A\char"032E} \SUPPORTEDANCH{\char"006C\char"032E} \SUPPORTEDFALL{\char"0234\char"032E} \SUPPORTEDANCH{\char"006E\char"032E} \SUPPORTEDANCH{\char"0235\char"032E} \SUPPORTEDANCH{\char"0070\char"032E} \SUPPORTEDFALL{\char"01A5\char"032E} \SUPPORTEDANCH{\char"0071\char"032E} \SUPPORTEDANCH{\char"02A0\char"032E} \SUPPORTEDANCH{\char"0074\char"032E} \SUPPORTEDFALL{\char"01AB\char"032E} \SUPPORTEDANCH{\char"01AD\char"032E} \SUPPORTEDANCH{\char"0236\char"032E} \SUPPORTEDANCH{\char"0075\char"032E} \SUPPORTEDANCH{\char"AB52\char"032E} \SUPPORTEDANCH{\char"0079\char"032E} \SUPPORTEDFALL{\char"01B4\char"032E} \SUPPORTEDFALL{\char"1EFF\char"032E}

\SUPPORTEDANCH{\char"0048\char"032F} \SUPPORTEDFALL{\char"2C67\char"032F} \SUPPORTEDANCH{\char"004B\char"032F} \SUPPORTEDANCH{\char"2C69\char"032F} \SUPPORTEDANCH{\char"004D\char"032F} \SUPPORTEDANCH{\char"2C6E\char"032F} \SUPPORTEDANCH{\char"0050\char"032F} \SUPPORTEDANCH{\char"01A4\char"032F} \SUPPORTEDANCH{\char"0052\char"032F} \SUPPORTEDFALL{\char"2C64\char"032F} \SUPPORTEDANCH{\char"0054\char"032F} \SUPPORTEDFALL{\char"01AC\char"032F} \SUPPORTEDFALL{\char"01AE\char"032F} \SUPPORTEDANCH{\char"0059\char"032F} \SUPPORTEDFALL{\char"01B3\char"032F} \SUPPORTEDANCH{\char"0061\char"032F} \SUPPORTEDFALL{\char"1E9A\char"032F} \SUPPORTEDANCH{\char"0062\char"032F} \SUPPORTEDANCH{\char"0183\char"032F} \SUPPORTEDANCH{\char"0063\char"032F} \SUPPORTEDFALL{\char"0188\char"032F} \SUPPORTEDANCH{\char"0064\char"032F} \SUPPORTEDANCH{\char"018C\char"032F} \SUPPORTEDFALL{\char"0221\char"032F} \SUPPORTEDANCH{\char"0068\char"032F} \SUPPORTEDANCH{\char"2C68\char"032F} \SUPPORTEDANCH{\char"006B\char"032F} \SUPPORTEDANCH{\char"0199\char"032F} \SUPPORTEDANCH{\char"2C6A\char"032F} \SUPPORTEDANCH{\char"006C\char"032F} \SUPPORTEDFALL{\char"0234\char"032F} \SUPPORTEDANCH{\char"006E\char"032F} \SUPPORTEDANCH{\char"0235\char"032F} \SUPPORTEDANCH{\char"0070\char"032F} \SUPPORTEDFALL{\char"01A5\char"032F} \SUPPORTEDANCH{\char"0071\char"032F} \SUPPORTEDANCH{\char"02A0\char"032F} \SUPPORTEDANCH{\char"0074\char"032F} \SUPPORTEDFALL{\char"01AB\char"032F} \SUPPORTEDANCH{\char"01AD\char"032F} \SUPPORTEDANCH{\char"0236\char"032F} \SUPPORTEDANCH{\char"0075\char"032F} \SUPPORTEDANCH{\char"AB52\char"032F} \SUPPORTEDANCH{\char"0079\char"032F} \SUPPORTEDFALL{\char"01B4\char"032F} \SUPPORTEDFALL{\char"1EFF\char"032F}

\SUPPORTEDANCH{\char"0048\char"031D} \SUPPORTEDFALL{\char"2C67\char"031D} \SUPPORTEDANCH{\char"004B\char"031D} \SUPPORTEDANCH{\char"2C69\char"031D} \SUPPORTEDANCH{\char"004D\char"031D} \SUPPORTEDANCH{\char"2C6E\char"031D} \SUPPORTEDANCH{\char"0050\char"031D} \SUPPORTEDANCH{\char"01A4\char"031D} \SUPPORTEDANCH{\char"0052\char"031D} \SUPPORTEDFALL{\char"2C64\char"031D} \SUPPORTEDANCH{\char"0054\char"031D} \SUPPORTEDFALL{\char"01AC\char"031D} \SUPPORTEDFALL{\char"01AE\char"031D} \SUPPORTEDANCH{\char"0059\char"031D} \SUPPORTEDFALL{\char"01B3\char"031D} \SUPPORTEDANCH{\char"0061\char"031D} \SUPPORTEDFALL{\char"1E9A\char"031D} \SUPPORTEDANCH{\char"0062\char"031D} \SUPPORTEDANCH{\char"0183\char"031D} \SUPPORTEDANCH{\char"0063\char"031D} \SUPPORTEDFALL{\char"0188\char"031D} \SUPPORTEDANCH{\char"0064\char"031D} \SUPPORTEDANCH{\char"018C\char"031D} \SUPPORTEDFALL{\char"0221\char"031D} \SUPPORTEDANCH{\char"0068\char"031D} \SUPPORTEDANCH{\char"2C68\char"031D} \SUPPORTEDANCH{\char"006B\char"031D} \SUPPORTEDANCH{\char"0199\char"031D} \SUPPORTEDANCH{\char"2C6A\char"031D} \SUPPORTEDANCH{\char"006C\char"031D} \SUPPORTEDFALL{\char"0234\char"031D} \SUPPORTEDANCH{\char"006E\char"031D} \SUPPORTEDANCH{\char"0235\char"031D} \SUPPORTEDANCH{\char"0070\char"031D} \SUPPORTEDFALL{\char"01A5\char"031D} \SUPPORTEDANCH{\char"0071\char"031D} \SUPPORTEDANCH{\char"02A0\char"031D} \SUPPORTEDANCH{\char"0074\char"031D} \SUPPORTEDFALL{\char"01AB\char"031D} \SUPPORTEDANCH{\char"01AD\char"031D} \SUPPORTEDANCH{\char"0236\char"031D} \SUPPORTEDANCH{\char"0075\char"031D} \SUPPORTEDANCH{\char"AB52\char"031D} \SUPPORTEDANCH{\char"0079\char"031D} \SUPPORTEDFALL{\char"01B4\char"031D} \SUPPORTEDFALL{\char"1EFF\char"031D}

\SUPPORTEDANCH{\char"0048\char"031E} \SUPPORTEDFALL{\char"2C67\char"031E} \SUPPORTEDANCH{\char"004B\char"031E} \SUPPORTEDANCH{\char"2C69\char"031E} \SUPPORTEDANCH{\char"004D\char"031E} \SUPPORTEDANCH{\char"2C6E\char"031E} \SUPPORTEDANCH{\char"0050\char"031E} \SUPPORTEDANCH{\char"01A4\char"031E} \SUPPORTEDANCH{\char"0052\char"031E} \SUPPORTEDFALL{\char"2C64\char"031E} \SUPPORTEDANCH{\char"0054\char"031E} \SUPPORTEDFALL{\char"01AC\char"031E} \SUPPORTEDFALL{\char"01AE\char"031E} \SUPPORTEDANCH{\char"0059\char"031E} \SUPPORTEDFALL{\char"01B3\char"031E} \SUPPORTEDANCH{\char"0061\char"031E} \SUPPORTEDFALL{\char"1E9A\char"031E} \SUPPORTEDANCH{\char"0062\char"031E} \SUPPORTEDANCH{\char"0183\char"031E} \SUPPORTEDANCH{\char"0063\char"031E} \SUPPORTEDFALL{\char"0188\char"031E} \SUPPORTEDANCH{\char"0064\char"031E} \SUPPORTEDANCH{\char"018C\char"031E} \SUPPORTEDFALL{\char"0221\char"031E} \SUPPORTEDANCH{\char"0068\char"031E} \SUPPORTEDANCH{\char"2C68\char"031E} \SUPPORTEDANCH{\char"006B\char"031E} \SUPPORTEDANCH{\char"0199\char"031E} \SUPPORTEDANCH{\char"2C6A\char"031E} \SUPPORTEDANCH{\char"006C\char"031E} \SUPPORTEDFALL{\char"0234\char"031E} \SUPPORTEDANCH{\char"006E\char"031E} \SUPPORTEDANCH{\char"0235\char"031E} \SUPPORTEDANCH{\char"0070\char"031E} \SUPPORTEDFALL{\char"01A5\char"031E} \SUPPORTEDANCH{\char"0071\char"031E} \SUPPORTEDANCH{\char"02A0\char"031E} \SUPPORTEDANCH{\char"0074\char"031E} \SUPPORTEDFALL{\char"01AB\char"031E} \SUPPORTEDANCH{\char"01AD\char"031E} \SUPPORTEDANCH{\char"0236\char"031E} \SUPPORTEDANCH{\char"0075\char"031E} \SUPPORTEDANCH{\char"AB52\char"031E} \SUPPORTEDANCH{\char"0079\char"031E} \SUPPORTEDFALL{\char"01B4\char"031E} \SUPPORTEDFALL{\char"1EFF\char"031E}

\SUPPORTEDANCH{\char"0048\char"031F} \SUPPORTEDFALL{\char"2C67\char"031F} \SUPPORTEDANCH{\char"004B\char"031F} \SUPPORTEDANCH{\char"2C69\char"031F} \SUPPORTEDANCH{\char"004D\char"031F} \SUPPORTEDANCH{\char"2C6E\char"031F} \SUPPORTEDANCH{\char"0050\char"031F} \SUPPORTEDANCH{\char"01A4\char"031F} \SUPPORTEDANCH{\char"0052\char"031F} \SUPPORTEDFALL{\char"2C64\char"031F} \SUPPORTEDANCH{\char"0054\char"031F} \SUPPORTEDFALL{\char"01AC\char"031F} \SUPPORTEDFALL{\char"01AE\char"031F} \SUPPORTEDANCH{\char"0059\char"031F} \SUPPORTEDFALL{\char"01B3\char"031F} \SUPPORTEDANCH{\char"0061\char"031F} \SUPPORTEDFALL{\char"1E9A\char"031F} \SUPPORTEDANCH{\char"0062\char"031F} \SUPPORTEDANCH{\char"0183\char"031F} \SUPPORTEDANCH{\char"0063\char"031F} \SUPPORTEDFALL{\char"0188\char"031F} \SUPPORTEDANCH{\char"0064\char"031F} \SUPPORTEDANCH{\char"018C\char"031F} \SUPPORTEDFALL{\char"0221\char"031F} \SUPPORTEDANCH{\char"0068\char"031F} \SUPPORTEDANCH{\char"2C68\char"031F} \SUPPORTEDANCH{\char"006B\char"031F} \SUPPORTEDANCH{\char"0199\char"031F} \SUPPORTEDANCH{\char"2C6A\char"031F} \SUPPORTEDANCH{\char"006C\char"031F} \SUPPORTEDFALL{\char"0234\char"031F} \SUPPORTEDANCH{\char"006E\char"031F} \SUPPORTEDANCH{\char"0235\char"031F} \SUPPORTEDANCH{\char"0070\char"031F} \SUPPORTEDFALL{\char"01A5\char"031F} \SUPPORTEDANCH{\char"0071\char"031F} \SUPPORTEDANCH{\char"02A0\char"031F} \SUPPORTEDANCH{\char"0074\char"031F} \SUPPORTEDFALL{\char"01AB\char"031F} \SUPPORTEDANCH{\char"01AD\char"031F} \SUPPORTEDANCH{\char"0236\char"031F} \SUPPORTEDANCH{\char"0075\char"031F} \SUPPORTEDANCH{\char"AB52\char"031F} \SUPPORTEDANCH{\char"0079\char"031F} \SUPPORTEDFALL{\char"01B4\char"031F} \SUPPORTEDFALL{\char"1EFF\char"031F}

\SUPPORTEDANCH{\char"0048\char"0330} \SUPPORTEDFALL{\char"2C67\char"0330} \SUPPORTEDANCH{\char"004B\char"0330} \SUPPORTEDANCH{\char"2C69\char"0330} \SUPPORTEDANCH{\char"004D\char"0330} \SUPPORTEDANCH{\char"2C6E\char"0330} \SUPPORTEDANCH{\char"0050\char"0330} \SUPPORTEDANCH{\char"01A4\char"0330} \SUPPORTEDANCH{\char"0052\char"0330} \SUPPORTEDFALL{\char"2C64\char"0330} \SUPPORTEDANCH{\char"0054\char"0330} \SUPPORTEDFALL{\char"01AC\char"0330} \SUPPORTEDFALL{\char"01AE\char"0330} \SUPPORTEDANCH{\char"0059\char"0330} \SUPPORTEDFALL{\char"01B3\char"0330} \SUPPORTEDANCH{\char"0061\char"0330} \SUPPORTEDFALL{\char"1E9A\char"0330} \SUPPORTEDANCH{\char"0062\char"0330} \SUPPORTEDANCH{\char"0183\char"0330} \SUPPORTEDANCH{\char"0063\char"0330} \SUPPORTEDFALL{\char"0188\char"0330} \SUPPORTEDANCH{\char"0064\char"0330} \SUPPORTEDANCH{\char"018C\char"0330} \SUPPORTEDFALL{\char"0221\char"0330} \SUPPORTEDANCH{\char"0068\char"0330} \SUPPORTEDANCH{\char"2C68\char"0330} \SUPPORTEDANCH{\char"006B\char"0330} \SUPPORTEDANCH{\char"0199\char"0330} \SUPPORTEDANCH{\char"2C6A\char"0330} \SUPPORTEDANCH{\char"006C\char"0330} \SUPPORTEDFALL{\char"0234\char"0330} \SUPPORTEDANCH{\char"006E\char"0330} \SUPPORTEDANCH{\char"0235\char"0330} \SUPPORTEDANCH{\char"0070\char"0330} \SUPPORTEDFALL{\char"01A5\char"0330} \SUPPORTEDANCH{\char"0071\char"0330} \SUPPORTEDANCH{\char"02A0\char"0330} \SUPPORTEDANCH{\char"0074\char"0330} \SUPPORTEDFALL{\char"01AB\char"0330} \SUPPORTEDANCH{\char"01AD\char"0330} \SUPPORTEDANCH{\char"0236\char"0330} \SUPPORTEDPREC{\char"0075\char"0330} \SUPPORTEDANCH{\char"AB52\char"0330} \SUPPORTEDANCH{\char"0079\char"0330} \SUPPORTEDFALL{\char"01B4\char"0330} \SUPPORTEDFALL{\char"1EFF\char"0330}

\SUPPORTEDANCH{\char"0048\char"0331} \SUPPORTEDFALL{\char"2C67\char"0331} \SUPPORTEDPREC{\char"004B\char"0331} \SUPPORTEDANCH{\char"2C69\char"0331} \SUPPORTEDANCH{\char"004D\char"0331} \SUPPORTEDANCH{\char"2C6E\char"0331} \SUPPORTEDANCH{\char"0050\char"0331} \SUPPORTEDANCH{\char"01A4\char"0331} \SUPPORTEDPREC{\char"0052\char"0331} \SUPPORTEDFALL{\char"2C64\char"0331} \SUPPORTEDPREC{\char"0054\char"0331} \SUPPORTEDFALL{\char"01AC\char"0331} \SUPPORTEDFALL{\char"01AE\char"0331} \SUPPORTEDANCH{\char"0059\char"0331} \SUPPORTEDFALL{\char"01B3\char"0331} \SUPPORTEDANCH{\char"0061\char"0331} \SUPPORTEDFALL{\char"1E9A\char"0331} \SUPPORTEDPREC{\char"0062\char"0331} \SUPPORTEDANCH{\char"0183\char"0331} \SUPPORTEDANCH{\char"0063\char"0331} \SUPPORTEDFALL{\char"0188\char"0331} \SUPPORTEDPREC{\char"0064\char"0331} \SUPPORTEDANCH{\char"018C\char"0331} \SUPPORTEDFALL{\char"0221\char"0331} \SUPPORTEDPREC{\char"0068\char"0331} \SUPPORTEDANCH{\char"2C68\char"0331} \SUPPORTEDPREC{\char"006B\char"0331} \SUPPORTEDANCH{\char"0199\char"0331} \SUPPORTEDANCH{\char"2C6A\char"0331} \SUPPORTEDPREC{\char"006C\char"0331} \SUPPORTEDFALL{\char"0234\char"0331} \SUPPORTEDPREC{\char"006E\char"0331} \SUPPORTEDANCH{\char"0235\char"0331} \SUPPORTEDANCH{\char"0070\char"0331} \SUPPORTEDFALL{\char"01A5\char"0331} \SUPPORTEDANCH{\char"0071\char"0331} \SUPPORTEDANCH{\char"02A0\char"0331} \SUPPORTEDPREC{\char"0074\char"0331} \SUPPORTEDFALL{\char"01AB\char"0331} \SUPPORTEDANCH{\char"01AD\char"0331} \SUPPORTEDANCH{\char"0236\char"0331} \SUPPORTEDANCH{\char"0075\char"0331} \SUPPORTEDANCH{\char"AB52\char"0331} \SUPPORTEDANCH{\char"0079\char"0331} \SUPPORTEDFALL{\char"01B4\char"0331} \SUPPORTEDFALL{\char"1EFF\char"0331}

\SUPPORTEDANCH{\char"0048\char"0333} \SUPPORTEDFALL{\char"2C67\char"0333} \SUPPORTEDANCH{\char"004B\char"0333} \SUPPORTEDANCH{\char"2C69\char"0333} \SUPPORTEDANCH{\char"004D\char"0333} \SUPPORTEDANCH{\char"2C6E\char"0333} \SUPPORTEDANCH{\char"0050\char"0333} \SUPPORTEDANCH{\char"01A4\char"0333} \SUPPORTEDANCH{\char"0052\char"0333} \SUPPORTEDFALL{\char"2C64\char"0333} \SUPPORTEDANCH{\char"0054\char"0333} \SUPPORTEDFALL{\char"01AC\char"0333} \SUPPORTEDFALL{\char"01AE\char"0333} \SUPPORTEDANCH{\char"0059\char"0333} \SUPPORTEDFALL{\char"01B3\char"0333} \SUPPORTEDANCH{\char"0061\char"0333} \SUPPORTEDFALL{\char"1E9A\char"0333} \SUPPORTEDANCH{\char"0062\char"0333} \SUPPORTEDANCH{\char"0183\char"0333} \SUPPORTEDANCH{\char"0063\char"0333} \SUPPORTEDFALL{\char"0188\char"0333} \SUPPORTEDANCH{\char"0064\char"0333} \SUPPORTEDANCH{\char"018C\char"0333} \SUPPORTEDFALL{\char"0221\char"0333} \SUPPORTEDANCH{\char"0068\char"0333} \SUPPORTEDANCH{\char"2C68\char"0333} \SUPPORTEDANCH{\char"006B\char"0333} \SUPPORTEDANCH{\char"0199\char"0333} \SUPPORTEDANCH{\char"2C6A\char"0333} \SUPPORTEDANCH{\char"006C\char"0333} \SUPPORTEDFALL{\char"0234\char"0333} \SUPPORTEDANCH{\char"006E\char"0333} \SUPPORTEDANCH{\char"0235\char"0333} \SUPPORTEDANCH{\char"0070\char"0333} \SUPPORTEDFALL{\char"01A5\char"0333} \SUPPORTEDANCH{\char"0071\char"0333} \SUPPORTEDANCH{\char"02A0\char"0333} \SUPPORTEDANCH{\char"0074\char"0333} \SUPPORTEDFALL{\char"01AB\char"0333} \SUPPORTEDANCH{\char"01AD\char"0333} \SUPPORTEDANCH{\char"0236\char"0333} \SUPPORTEDANCH{\char"0075\char"0333} \SUPPORTEDANCH{\char"AB52\char"0333} \SUPPORTEDANCH{\char"0079\char"0333} \SUPPORTEDFALL{\char"01B4\char"0333} \SUPPORTEDFALL{\char"1EFF\char"0333}

\SUPPORTEDANCH{\char"0048\char"0339} \SUPPORTEDFALL{\char"2C67\char"0339} \SUPPORTEDANCH{\char"004B\char"0339} \SUPPORTEDANCH{\char"2C69\char"0339} \SUPPORTEDANCH{\char"004D\char"0339} \SUPPORTEDANCH{\char"2C6E\char"0339} \SUPPORTEDANCH{\char"0050\char"0339} \SUPPORTEDANCH{\char"01A4\char"0339} \SUPPORTEDANCH{\char"0052\char"0339} \SUPPORTEDFALL{\char"2C64\char"0339} \SUPPORTEDANCH{\char"0054\char"0339} \SUPPORTEDFALL{\char"01AC\char"0339} \SUPPORTEDFALL{\char"01AE\char"0339} \SUPPORTEDANCH{\char"0059\char"0339} \SUPPORTEDFALL{\char"01B3\char"0339} \SUPPORTEDANCH{\char"0061\char"0339} \SUPPORTEDFALL{\char"1E9A\char"0339} \SUPPORTEDANCH{\char"0062\char"0339} \SUPPORTEDANCH{\char"0183\char"0339} \SUPPORTEDANCH{\char"0063\char"0339} \SUPPORTEDFALL{\char"0188\char"0339} \SUPPORTEDANCH{\char"0064\char"0339} \SUPPORTEDANCH{\char"018C\char"0339} \SUPPORTEDFALL{\char"0221\char"0339} \SUPPORTEDANCH{\char"0068\char"0339} \SUPPORTEDANCH{\char"2C68\char"0339} \SUPPORTEDANCH{\char"006B\char"0339} \SUPPORTEDANCH{\char"0199\char"0339} \SUPPORTEDANCH{\char"2C6A\char"0339} \SUPPORTEDANCH{\char"006C\char"0339} \SUPPORTEDFALL{\char"0234\char"0339} \SUPPORTEDANCH{\char"006E\char"0339} \SUPPORTEDANCH{\char"0235\char"0339} \SUPPORTEDANCH{\char"0070\char"0339} \SUPPORTEDFALL{\char"01A5\char"0339} \SUPPORTEDANCH{\char"0071\char"0339} \SUPPORTEDANCH{\char"02A0\char"0339} \SUPPORTEDANCH{\char"0074\char"0339} \SUPPORTEDFALL{\char"01AB\char"0339} \SUPPORTEDANCH{\char"01AD\char"0339} \SUPPORTEDANCH{\char"0236\char"0339} \SUPPORTEDANCH{\char"0075\char"0339} \SUPPORTEDANCH{\char"AB52\char"0339} \SUPPORTEDANCH{\char"0079\char"0339} \SUPPORTEDFALL{\char"01B4\char"0339} \SUPPORTEDFALL{\char"1EFF\char"0339}

\SUPPORTEDANCH{\char"0048\char"033A} \SUPPORTEDFALL{\char"2C67\char"033A} \SUPPORTEDANCH{\char"004B\char"033A} \SUPPORTEDANCH{\char"2C69\char"033A} \SUPPORTEDANCH{\char"004D\char"033A} \SUPPORTEDANCH{\char"2C6E\char"033A} \SUPPORTEDANCH{\char"0050\char"033A} \SUPPORTEDANCH{\char"01A4\char"033A} \SUPPORTEDANCH{\char"0052\char"033A} \SUPPORTEDFALL{\char"2C64\char"033A} \SUPPORTEDANCH{\char"0054\char"033A} \SUPPORTEDFALL{\char"01AC\char"033A} \SUPPORTEDFALL{\char"01AE\char"033A} \SUPPORTEDANCH{\char"0059\char"033A} \SUPPORTEDFALL{\char"01B3\char"033A} \SUPPORTEDANCH{\char"0061\char"033A} \SUPPORTEDFALL{\char"1E9A\char"033A} \SUPPORTEDANCH{\char"0062\char"033A} \SUPPORTEDANCH{\char"0183\char"033A} \SUPPORTEDANCH{\char"0063\char"033A} \SUPPORTEDFALL{\char"0188\char"033A} \SUPPORTEDANCH{\char"0064\char"033A} \SUPPORTEDANCH{\char"018C\char"033A} \SUPPORTEDFALL{\char"0221\char"033A} \SUPPORTEDANCH{\char"0068\char"033A} \SUPPORTEDANCH{\char"2C68\char"033A} \SUPPORTEDANCH{\char"006B\char"033A} \SUPPORTEDANCH{\char"0199\char"033A} \SUPPORTEDANCH{\char"2C6A\char"033A} \SUPPORTEDANCH{\char"006C\char"033A} \SUPPORTEDFALL{\char"0234\char"033A} \SUPPORTEDANCH{\char"006E\char"033A} \SUPPORTEDANCH{\char"0235\char"033A} \SUPPORTEDANCH{\char"0070\char"033A} \SUPPORTEDFALL{\char"01A5\char"033A} \SUPPORTEDANCH{\char"0071\char"033A} \SUPPORTEDANCH{\char"02A0\char"033A} \SUPPORTEDANCH{\char"0074\char"033A} \SUPPORTEDFALL{\char"01AB\char"033A} \SUPPORTEDANCH{\char"01AD\char"033A} \SUPPORTEDANCH{\char"0236\char"033A} \SUPPORTEDANCH{\char"0075\char"033A} \SUPPORTEDANCH{\char"AB52\char"033A} \SUPPORTEDANCH{\char"0079\char"033A} \SUPPORTEDFALL{\char"01B4\char"033A} \SUPPORTEDFALL{\char"1EFF\char"033A}

\SUPPORTEDANCH{\char"0048\char"033B} \SUPPORTEDFALL{\char"2C67\char"033B} \SUPPORTEDANCH{\char"004B\char"033B} \SUPPORTEDANCH{\char"2C69\char"033B} \SUPPORTEDANCH{\char"004D\char"033B} \SUPPORTEDANCH{\char"2C6E\char"033B} \SUPPORTEDANCH{\char"0050\char"033B} \SUPPORTEDANCH{\char"01A4\char"033B} \SUPPORTEDANCH{\char"0052\char"033B} \SUPPORTEDFALL{\char"2C64\char"033B} \SUPPORTEDANCH{\char"0054\char"033B} \SUPPORTEDFALL{\char"01AC\char"033B} \SUPPORTEDFALL{\char"01AE\char"033B} \SUPPORTEDANCH{\char"0059\char"033B} \SUPPORTEDFALL{\char"01B3\char"033B} \SUPPORTEDANCH{\char"0061\char"033B} \SUPPORTEDFALL{\char"1E9A\char"033B} \SUPPORTEDANCH{\char"0062\char"033B} \SUPPORTEDANCH{\char"0183\char"033B} \SUPPORTEDANCH{\char"0063\char"033B} \SUPPORTEDFALL{\char"0188\char"033B} \SUPPORTEDANCH{\char"0064\char"033B} \SUPPORTEDANCH{\char"018C\char"033B} \SUPPORTEDFALL{\char"0221\char"033B} \SUPPORTEDANCH{\char"0068\char"033B} \SUPPORTEDANCH{\char"2C68\char"033B} \SUPPORTEDANCH{\char"006B\char"033B} \SUPPORTEDANCH{\char"0199\char"033B} \SUPPORTEDANCH{\char"2C6A\char"033B} \SUPPORTEDANCH{\char"006C\char"033B} \SUPPORTEDFALL{\char"0234\char"033B} \SUPPORTEDANCH{\char"006E\char"033B} \SUPPORTEDANCH{\char"0235\char"033B} \SUPPORTEDANCH{\char"0070\char"033B} \SUPPORTEDFALL{\char"01A5\char"033B} \SUPPORTEDANCH{\char"0071\char"033B} \SUPPORTEDANCH{\char"02A0\char"033B} \SUPPORTEDANCH{\char"0074\char"033B} \SUPPORTEDFALL{\char"01AB\char"033B} \SUPPORTEDANCH{\char"01AD\char"033B} \SUPPORTEDANCH{\char"0236\char"033B} \SUPPORTEDANCH{\char"0075\char"033B} \SUPPORTEDANCH{\char"AB52\char"033B} \SUPPORTEDANCH{\char"0079\char"033B} \SUPPORTEDFALL{\char"01B4\char"033B} \SUPPORTEDFALL{\char"1EFF\char"033B}

\SUPPORTEDANCH{\char"0048\char"033C} \SUPPORTEDFALL{\char"2C67\char"033C} \SUPPORTEDANCH{\char"004B\char"033C} \SUPPORTEDANCH{\char"2C69\char"033C} \SUPPORTEDANCH{\char"004D\char"033C} \SUPPORTEDANCH{\char"2C6E\char"033C} \SUPPORTEDANCH{\char"0050\char"033C} \SUPPORTEDANCH{\char"01A4\char"033C} \SUPPORTEDANCH{\char"0052\char"033C} \SUPPORTEDFALL{\char"2C64\char"033C} \SUPPORTEDANCH{\char"0054\char"033C} \SUPPORTEDFALL{\char"01AC\char"033C} \SUPPORTEDFALL{\char"01AE\char"033C} \SUPPORTEDANCH{\char"0059\char"033C} \SUPPORTEDFALL{\char"01B3\char"033C} \SUPPORTEDANCH{\char"0061\char"033C} \SUPPORTEDFALL{\char"1E9A\char"033C} \SUPPORTEDANCH{\char"0062\char"033C} \SUPPORTEDANCH{\char"0183\char"033C} \SUPPORTEDANCH{\char"0063\char"033C} \SUPPORTEDFALL{\char"0188\char"033C} \SUPPORTEDANCH{\char"0064\char"033C} \SUPPORTEDANCH{\char"018C\char"033C} \SUPPORTEDFALL{\char"0221\char"033C} \SUPPORTEDANCH{\char"0068\char"033C} \SUPPORTEDANCH{\char"2C68\char"033C} \SUPPORTEDANCH{\char"006B\char"033C} \SUPPORTEDANCH{\char"0199\char"033C} \SUPPORTEDANCH{\char"2C6A\char"033C} \SUPPORTEDANCH{\char"006C\char"033C} \SUPPORTEDFALL{\char"0234\char"033C} \SUPPORTEDANCH{\char"006E\char"033C} \SUPPORTEDANCH{\char"0235\char"033C} \SUPPORTEDANCH{\char"0070\char"033C} \SUPPORTEDFALL{\char"01A5\char"033C} \SUPPORTEDANCH{\char"0071\char"033C} \SUPPORTEDANCH{\char"02A0\char"033C} \SUPPORTEDANCH{\char"0074\char"033C} \SUPPORTEDFALL{\char"01AB\char"033C} \SUPPORTEDANCH{\char"01AD\char"033C} \SUPPORTEDANCH{\char"0236\char"033C} \SUPPORTEDANCH{\char"0075\char"033C} \SUPPORTEDANCH{\char"AB52\char"033C} \SUPPORTEDANCH{\char"0079\char"033C} \SUPPORTEDFALL{\char"01B4\char"033C} \SUPPORTEDFALL{\char"1EFF\char"033C}

\SUPPORTEDANCH{\char"0048\char"0347} \SUPPORTEDFALL{\char"2C67\char"0347} \SUPPORTEDANCH{\char"004B\char"0347} \SUPPORTEDANCH{\char"2C69\char"0347} \SUPPORTEDANCH{\char"004D\char"0347} \SUPPORTEDANCH{\char"2C6E\char"0347} \SUPPORTEDANCH{\char"0050\char"0347} \SUPPORTEDANCH{\char"01A4\char"0347} \SUPPORTEDANCH{\char"0052\char"0347} \SUPPORTEDFALL{\char"2C64\char"0347} \SUPPORTEDANCH{\char"0054\char"0347} \SUPPORTEDFALL{\char"01AC\char"0347} \SUPPORTEDFALL{\char"01AE\char"0347} \SUPPORTEDANCH{\char"0059\char"0347} \SUPPORTEDFALL{\char"01B3\char"0347} \SUPPORTEDANCH{\char"0061\char"0347} \SUPPORTEDFALL{\char"1E9A\char"0347} \SUPPORTEDANCH{\char"0062\char"0347} \SUPPORTEDANCH{\char"0183\char"0347} \SUPPORTEDANCH{\char"0063\char"0347} \SUPPORTEDFALL{\char"0188\char"0347} \SUPPORTEDANCH{\char"0064\char"0347} \SUPPORTEDANCH{\char"018C\char"0347} \SUPPORTEDFALL{\char"0221\char"0347} \SUPPORTEDANCH{\char"0068\char"0347} \SUPPORTEDANCH{\char"2C68\char"0347} \SUPPORTEDANCH{\char"006B\char"0347} \SUPPORTEDANCH{\char"0199\char"0347} \SUPPORTEDANCH{\char"2C6A\char"0347} \SUPPORTEDANCH{\char"006C\char"0347} \SUPPORTEDFALL{\char"0234\char"0347} \SUPPORTEDANCH{\char"006E\char"0347} \SUPPORTEDANCH{\char"0235\char"0347} \SUPPORTEDANCH{\char"0070\char"0347} \SUPPORTEDFALL{\char"01A5\char"0347} \SUPPORTEDANCH{\char"0071\char"0347} \SUPPORTEDANCH{\char"02A0\char"0347} \SUPPORTEDANCH{\char"0074\char"0347} \SUPPORTEDFALL{\char"01AB\char"0347} \SUPPORTEDANCH{\char"01AD\char"0347} \SUPPORTEDANCH{\char"0236\char"0347} \SUPPORTEDANCH{\char"0075\char"0347} \SUPPORTEDANCH{\char"AB52\char"0347} \SUPPORTEDANCH{\char"0079\char"0347} \SUPPORTEDFALL{\char"01B4\char"0347} \SUPPORTEDFALL{\char"1EFF\char"0347}

\SUPPORTEDANCH{\char"0048\char"0348} \SUPPORTEDFALL{\char"2C67\char"0348} \SUPPORTEDANCH{\char"004B\char"0348} \SUPPORTEDANCH{\char"2C69\char"0348} \SUPPORTEDANCH{\char"004D\char"0348} \SUPPORTEDANCH{\char"2C6E\char"0348} \SUPPORTEDANCH{\char"0050\char"0348} \SUPPORTEDANCH{\char"01A4\char"0348} \SUPPORTEDANCH{\char"0052\char"0348} \SUPPORTEDFALL{\char"2C64\char"0348} \SUPPORTEDANCH{\char"0054\char"0348} \SUPPORTEDFALL{\char"01AC\char"0348} \SUPPORTEDFALL{\char"01AE\char"0348} \SUPPORTEDANCH{\char"0059\char"0348} \SUPPORTEDFALL{\char"01B3\char"0348} \SUPPORTEDANCH{\char"0061\char"0348} \SUPPORTEDFALL{\char"1E9A\char"0348} \SUPPORTEDANCH{\char"0062\char"0348} \SUPPORTEDANCH{\char"0183\char"0348} \SUPPORTEDANCH{\char"0063\char"0348} \SUPPORTEDFALL{\char"0188\char"0348} \SUPPORTEDANCH{\char"0064\char"0348} \SUPPORTEDANCH{\char"018C\char"0348} \SUPPORTEDFALL{\char"0221\char"0348} \SUPPORTEDANCH{\char"0068\char"0348} \SUPPORTEDANCH{\char"2C68\char"0348} \SUPPORTEDANCH{\char"006B\char"0348} \SUPPORTEDANCH{\char"0199\char"0348} \SUPPORTEDANCH{\char"2C6A\char"0348} \SUPPORTEDANCH{\char"006C\char"0348} \SUPPORTEDFALL{\char"0234\char"0348} \SUPPORTEDANCH{\char"006E\char"0348} \SUPPORTEDANCH{\char"0235\char"0348} \SUPPORTEDANCH{\char"0070\char"0348} \SUPPORTEDFALL{\char"01A5\char"0348} \SUPPORTEDANCH{\char"0071\char"0348} \SUPPORTEDANCH{\char"02A0\char"0348} \SUPPORTEDANCH{\char"0074\char"0348} \SUPPORTEDFALL{\char"01AB\char"0348} \SUPPORTEDANCH{\char"01AD\char"0348} \SUPPORTEDANCH{\char"0236\char"0348} \SUPPORTEDANCH{\char"0075\char"0348} \SUPPORTEDANCH{\char"AB52\char"0348} \SUPPORTEDANCH{\char"0079\char"0348} \SUPPORTEDFALL{\char"01B4\char"0348} \SUPPORTEDFALL{\char"1EFF\char"0348}

\SUPPORTEDANCH{\char"0048\char"0349} \SUPPORTEDFALL{\char"2C67\char"0349} \SUPPORTEDANCH{\char"004B\char"0349} \SUPPORTEDANCH{\char"2C69\char"0349} \SUPPORTEDANCH{\char"004D\char"0349} \SUPPORTEDANCH{\char"2C6E\char"0349} \SUPPORTEDANCH{\char"0050\char"0349} \SUPPORTEDANCH{\char"01A4\char"0349} \SUPPORTEDANCH{\char"0052\char"0349} \SUPPORTEDFALL{\char"2C64\char"0349} \SUPPORTEDANCH{\char"0054\char"0349} \SUPPORTEDFALL{\char"01AC\char"0349} \SUPPORTEDFALL{\char"01AE\char"0349} \SUPPORTEDANCH{\char"0059\char"0349} \SUPPORTEDFALL{\char"01B3\char"0349} \SUPPORTEDANCH{\char"0061\char"0349} \SUPPORTEDFALL{\char"1E9A\char"0349} \SUPPORTEDANCH{\char"0062\char"0349} \SUPPORTEDANCH{\char"0183\char"0349} \SUPPORTEDANCH{\char"0063\char"0349} \SUPPORTEDFALL{\char"0188\char"0349} \SUPPORTEDANCH{\char"0064\char"0349} \SUPPORTEDANCH{\char"018C\char"0349} \SUPPORTEDFALL{\char"0221\char"0349} \SUPPORTEDANCH{\char"0068\char"0349} \SUPPORTEDANCH{\char"2C68\char"0349} \SUPPORTEDANCH{\char"006B\char"0349} \SUPPORTEDANCH{\char"0199\char"0349} \SUPPORTEDANCH{\char"2C6A\char"0349} \SUPPORTEDANCH{\char"006C\char"0349} \SUPPORTEDFALL{\char"0234\char"0349} \SUPPORTEDANCH{\char"006E\char"0349} \SUPPORTEDANCH{\char"0235\char"0349} \SUPPORTEDANCH{\char"0070\char"0349} \SUPPORTEDFALL{\char"01A5\char"0349} \SUPPORTEDANCH{\char"0071\char"0349} \SUPPORTEDANCH{\char"02A0\char"0349} \SUPPORTEDANCH{\char"0074\char"0349} \SUPPORTEDFALL{\char"01AB\char"0349} \SUPPORTEDANCH{\char"01AD\char"0349} \SUPPORTEDANCH{\char"0236\char"0349} \SUPPORTEDANCH{\char"0075\char"0349} \SUPPORTEDANCH{\char"AB52\char"0349} \SUPPORTEDANCH{\char"0079\char"0349} \SUPPORTEDFALL{\char"01B4\char"0349} \SUPPORTEDFALL{\char"1EFF\char"0349}

\SUPPORTEDANCH{\char"0048\char"034D} \SUPPORTEDFALL{\char"2C67\char"034D} \SUPPORTEDANCH{\char"004B\char"034D} \SUPPORTEDANCH{\char"2C69\char"034D} \SUPPORTEDANCH{\char"004D\char"034D} \SUPPORTEDANCH{\char"2C6E\char"034D} \SUPPORTEDANCH{\char"0050\char"034D} \SUPPORTEDANCH{\char"01A4\char"034D} \SUPPORTEDANCH{\char"0052\char"034D} \SUPPORTEDFALL{\char"2C64\char"034D} \SUPPORTEDANCH{\char"0054\char"034D} \SUPPORTEDFALL{\char"01AC\char"034D} \SUPPORTEDFALL{\char"01AE\char"034D} \SUPPORTEDANCH{\char"0059\char"034D} \SUPPORTEDFALL{\char"01B3\char"034D} \SUPPORTEDANCH{\char"0061\char"034D} \SUPPORTEDFALL{\char"1E9A\char"034D} \SUPPORTEDANCH{\char"0062\char"034D} \SUPPORTEDANCH{\char"0183\char"034D} \SUPPORTEDANCH{\char"0063\char"034D} \SUPPORTEDFALL{\char"0188\char"034D} \SUPPORTEDANCH{\char"0064\char"034D} \SUPPORTEDANCH{\char"018C\char"034D} \SUPPORTEDFALL{\char"0221\char"034D} \SUPPORTEDANCH{\char"0068\char"034D} \SUPPORTEDANCH{\char"2C68\char"034D} \SUPPORTEDANCH{\char"006B\char"034D} \SUPPORTEDANCH{\char"0199\char"034D} \SUPPORTEDANCH{\char"2C6A\char"034D} \SUPPORTEDANCH{\char"006C\char"034D} \SUPPORTEDFALL{\char"0234\char"034D} \SUPPORTEDANCH{\char"006E\char"034D} \SUPPORTEDANCH{\char"0235\char"034D} \SUPPORTEDANCH{\char"0070\char"034D} \SUPPORTEDFALL{\char"01A5\char"034D} \SUPPORTEDANCH{\char"0071\char"034D} \SUPPORTEDANCH{\char"02A0\char"034D} \SUPPORTEDANCH{\char"0074\char"034D} \SUPPORTEDFALL{\char"01AB\char"034D} \SUPPORTEDANCH{\char"01AD\char"034D} \SUPPORTEDANCH{\char"0236\char"034D} \SUPPORTEDANCH{\char"0075\char"034D} \SUPPORTEDANCH{\char"AB52\char"034D} \SUPPORTEDANCH{\char"0079\char"034D} \SUPPORTEDFALL{\char"01B4\char"034D} \SUPPORTEDFALL{\char"1EFF\char"034D}

\SUPPORTEDANCH{\char"0048\char"034E} \SUPPORTEDFALL{\char"2C67\char"034E} \SUPPORTEDANCH{\char"004B\char"034E} \SUPPORTEDANCH{\char"2C69\char"034E} \SUPPORTEDANCH{\char"004D\char"034E} \SUPPORTEDANCH{\char"2C6E\char"034E} \SUPPORTEDANCH{\char"0050\char"034E} \SUPPORTEDANCH{\char"01A4\char"034E} \SUPPORTEDANCH{\char"0052\char"034E} \SUPPORTEDFALL{\char"2C64\char"034E} \SUPPORTEDANCH{\char"0054\char"034E} \SUPPORTEDFALL{\char"01AC\char"034E} \SUPPORTEDFALL{\char"01AE\char"034E} \SUPPORTEDANCH{\char"0059\char"034E} \SUPPORTEDFALL{\char"01B3\char"034E} \SUPPORTEDANCH{\char"0061\char"034E} \SUPPORTEDFALL{\char"1E9A\char"034E} \SUPPORTEDANCH{\char"0062\char"034E} \SUPPORTEDANCH{\char"0183\char"034E} \SUPPORTEDANCH{\char"0063\char"034E} \SUPPORTEDFALL{\char"0188\char"034E} \SUPPORTEDANCH{\char"0064\char"034E} \SUPPORTEDANCH{\char"018C\char"034E} \SUPPORTEDFALL{\char"0221\char"034E} \SUPPORTEDANCH{\char"0068\char"034E} \SUPPORTEDANCH{\char"2C68\char"034E} \SUPPORTEDANCH{\char"006B\char"034E} \SUPPORTEDANCH{\char"0199\char"034E} \SUPPORTEDANCH{\char"2C6A\char"034E} \SUPPORTEDANCH{\char"006C\char"034E} \SUPPORTEDFALL{\char"0234\char"034E} \SUPPORTEDANCH{\char"006E\char"034E} \SUPPORTEDANCH{\char"0235\char"034E} \SUPPORTEDANCH{\char"0070\char"034E} \SUPPORTEDFALL{\char"01A5\char"034E} \SUPPORTEDANCH{\char"0071\char"034E} \SUPPORTEDANCH{\char"02A0\char"034E} \SUPPORTEDANCH{\char"0074\char"034E} \SUPPORTEDFALL{\char"01AB\char"034E} \SUPPORTEDANCH{\char"01AD\char"034E} \SUPPORTEDANCH{\char"0236\char"034E} \SUPPORTEDANCH{\char"0075\char"034E} \SUPPORTEDANCH{\char"AB52\char"034E} \SUPPORTEDANCH{\char"0079\char"034E} \SUPPORTEDFALL{\char"01B4\char"034E} \SUPPORTEDFALL{\char"1EFF\char"034E}


\newpage
\subsection*{Regular patched, Group 4, below, Below}

{\LibertinusSerifPatch
% grid. columns are bases, rows are marks.
\SUPPORTEDANCH{\char"0048\char"0316} \SUPPORTEDANCH{\char"2C67\char"0316} \SUPPORTEDANCH{\char"004B\char"0316} \SUPPORTEDANCH{\char"2C69\char"0316} \SUPPORTEDANCH{\char"004D\char"0316} \SUPPORTEDANCH{\char"2C6E\char"0316} \SUPPORTEDANCH{\char"0050\char"0316} \SUPPORTEDANCH{\char"01A4\char"0316} \SUPPORTEDANCH{\char"0052\char"0316} \SUPPORTEDANCH{\char"2C64\char"0316} \SUPPORTEDANCH{\char"0054\char"0316} \SUPPORTEDANCH{\char"01AC\char"0316} \SUPPORTEDANCH{\char"01AE\char"0316} \SUPPORTEDANCH{\char"0059\char"0316} \SUPPORTEDANCH{\char"01B3\char"0316} \SUPPORTEDANCH{\char"0061\char"0316} \SUPPORTEDANCH{\char"1E9A\char"0316} \SUPPORTEDANCH{\char"0062\char"0316} \SUPPORTEDANCH{\char"0183\char"0316} \SUPPORTEDANCH{\char"0063\char"0316} \SUPPORTEDANCH{\char"0188\char"0316} \SUPPORTEDANCH{\char"0064\char"0316} \SUPPORTEDANCH{\char"018C\char"0316} \SUPPORTEDANCH{\char"0221\char"0316} \SUPPORTEDANCH{\char"0068\char"0316} \SUPPORTEDANCH{\char"2C68\char"0316} \SUPPORTEDANCH{\char"006B\char"0316} \SUPPORTEDANCH{\char"0199\char"0316} \SUPPORTEDANCH{\char"2C6A\char"0316} \SUPPORTEDANCH{\char"006C\char"0316} \SUPPORTEDANCH{\char"0234\char"0316} \SUPPORTEDANCH{\char"006E\char"0316} \SUPPORTEDANCH{\char"0235\char"0316} \SUPPORTEDANCH{\char"0070\char"0316} \SUPPORTEDANCH{\char"01A5\char"0316} \SUPPORTEDANCH{\char"0071\char"0316} \SUPPORTEDANCH{\char"02A0\char"0316} \SUPPORTEDANCH{\char"0074\char"0316} \SUPPORTEDANCH{\char"01AB\char"0316} \SUPPORTEDANCH{\char"01AD\char"0316} \SUPPORTEDANCH{\char"0236\char"0316} \SUPPORTEDANCH{\char"0075\char"0316} \SUPPORTEDANCH{\char"AB52\char"0316} \SUPPORTEDANCH{\char"0079\char"0316} \SUPPORTEDANCH{\char"01B4\char"0316} \SUPPORTEDANCH{\char"1EFF\char"0316}

\SUPPORTEDANCH{\char"0048\char"0317} \SUPPORTEDANCH{\char"2C67\char"0317} \SUPPORTEDANCH{\char"004B\char"0317} \SUPPORTEDANCH{\char"2C69\char"0317} \SUPPORTEDANCH{\char"004D\char"0317} \SUPPORTEDANCH{\char"2C6E\char"0317} \SUPPORTEDANCH{\char"0050\char"0317} \SUPPORTEDANCH{\char"01A4\char"0317} \SUPPORTEDANCH{\char"0052\char"0317} \SUPPORTEDANCH{\char"2C64\char"0317} \SUPPORTEDANCH{\char"0054\char"0317} \SUPPORTEDANCH{\char"01AC\char"0317} \SUPPORTEDANCH{\char"01AE\char"0317} \SUPPORTEDANCH{\char"0059\char"0317} \SUPPORTEDANCH{\char"01B3\char"0317} \SUPPORTEDANCH{\char"0061\char"0317} \SUPPORTEDANCH{\char"1E9A\char"0317} \SUPPORTEDANCH{\char"0062\char"0317} \SUPPORTEDANCH{\char"0183\char"0317} \SUPPORTEDANCH{\char"0063\char"0317} \SUPPORTEDANCH{\char"0188\char"0317} \SUPPORTEDANCH{\char"0064\char"0317} \SUPPORTEDANCH{\char"018C\char"0317} \SUPPORTEDANCH{\char"0221\char"0317} \SUPPORTEDANCH{\char"0068\char"0317} \SUPPORTEDANCH{\char"2C68\char"0317} \SUPPORTEDANCH{\char"006B\char"0317} \SUPPORTEDANCH{\char"0199\char"0317} \SUPPORTEDANCH{\char"2C6A\char"0317} \SUPPORTEDANCH{\char"006C\char"0317} \SUPPORTEDANCH{\char"0234\char"0317} \SUPPORTEDANCH{\char"006E\char"0317} \SUPPORTEDANCH{\char"0235\char"0317} \SUPPORTEDANCH{\char"0070\char"0317} \SUPPORTEDANCH{\char"01A5\char"0317} \SUPPORTEDANCH{\char"0071\char"0317} \SUPPORTEDANCH{\char"02A0\char"0317} \SUPPORTEDANCH{\char"0074\char"0317} \SUPPORTEDANCH{\char"01AB\char"0317} \SUPPORTEDANCH{\char"01AD\char"0317} \SUPPORTEDANCH{\char"0236\char"0317} \SUPPORTEDANCH{\char"0075\char"0317} \SUPPORTEDANCH{\char"AB52\char"0317} \SUPPORTEDANCH{\char"0079\char"0317} \SUPPORTEDANCH{\char"01B4\char"0317} \SUPPORTEDANCH{\char"1EFF\char"0317}

\SUPPORTEDANCH{\char"0048\char"0318} \SUPPORTEDANCH{\char"2C67\char"0318} \SUPPORTEDANCH{\char"004B\char"0318} \SUPPORTEDANCH{\char"2C69\char"0318} \SUPPORTEDANCH{\char"004D\char"0318} \SUPPORTEDANCH{\char"2C6E\char"0318} \SUPPORTEDANCH{\char"0050\char"0318} \SUPPORTEDANCH{\char"01A4\char"0318} \SUPPORTEDANCH{\char"0052\char"0318} \SUPPORTEDANCH{\char"2C64\char"0318} \SUPPORTEDANCH{\char"0054\char"0318} \SUPPORTEDANCH{\char"01AC\char"0318} \SUPPORTEDANCH{\char"01AE\char"0318} \SUPPORTEDANCH{\char"0059\char"0318} \SUPPORTEDANCH{\char"01B3\char"0318} \SUPPORTEDANCH{\char"0061\char"0318} \SUPPORTEDANCH{\char"1E9A\char"0318} \SUPPORTEDANCH{\char"0062\char"0318} \SUPPORTEDANCH{\char"0183\char"0318} \SUPPORTEDANCH{\char"0063\char"0318} \SUPPORTEDANCH{\char"0188\char"0318} \SUPPORTEDANCH{\char"0064\char"0318} \SUPPORTEDANCH{\char"018C\char"0318} \SUPPORTEDANCH{\char"0221\char"0318} \SUPPORTEDANCH{\char"0068\char"0318} \SUPPORTEDANCH{\char"2C68\char"0318} \SUPPORTEDANCH{\char"006B\char"0318} \SUPPORTEDANCH{\char"0199\char"0318} \SUPPORTEDANCH{\char"2C6A\char"0318} \SUPPORTEDANCH{\char"006C\char"0318} \SUPPORTEDANCH{\char"0234\char"0318} \SUPPORTEDANCH{\char"006E\char"0318} \SUPPORTEDANCH{\char"0235\char"0318} \SUPPORTEDANCH{\char"0070\char"0318} \SUPPORTEDANCH{\char"01A5\char"0318} \SUPPORTEDANCH{\char"0071\char"0318} \SUPPORTEDANCH{\char"02A0\char"0318} \SUPPORTEDANCH{\char"0074\char"0318} \SUPPORTEDANCH{\char"01AB\char"0318} \SUPPORTEDANCH{\char"01AD\char"0318} \SUPPORTEDANCH{\char"0236\char"0318} \SUPPORTEDANCH{\char"0075\char"0318} \SUPPORTEDANCH{\char"AB52\char"0318} \SUPPORTEDANCH{\char"0079\char"0318} \SUPPORTEDANCH{\char"01B4\char"0318} \SUPPORTEDANCH{\char"1EFF\char"0318}

\SUPPORTEDANCH{\char"0048\char"0319} \SUPPORTEDANCH{\char"2C67\char"0319} \SUPPORTEDANCH{\char"004B\char"0319} \SUPPORTEDANCH{\char"2C69\char"0319} \SUPPORTEDANCH{\char"004D\char"0319} \SUPPORTEDANCH{\char"2C6E\char"0319} \SUPPORTEDANCH{\char"0050\char"0319} \SUPPORTEDANCH{\char"01A4\char"0319} \SUPPORTEDANCH{\char"0052\char"0319} \SUPPORTEDANCH{\char"2C64\char"0319} \SUPPORTEDANCH{\char"0054\char"0319} \SUPPORTEDANCH{\char"01AC\char"0319} \SUPPORTEDANCH{\char"01AE\char"0319} \SUPPORTEDANCH{\char"0059\char"0319} \SUPPORTEDANCH{\char"01B3\char"0319} \SUPPORTEDANCH{\char"0061\char"0319} \SUPPORTEDANCH{\char"1E9A\char"0319} \SUPPORTEDANCH{\char"0062\char"0319} \SUPPORTEDANCH{\char"0183\char"0319} \SUPPORTEDANCH{\char"0063\char"0319} \SUPPORTEDANCH{\char"0188\char"0319} \SUPPORTEDANCH{\char"0064\char"0319} \SUPPORTEDANCH{\char"018C\char"0319} \SUPPORTEDANCH{\char"0221\char"0319} \SUPPORTEDANCH{\char"0068\char"0319} \SUPPORTEDANCH{\char"2C68\char"0319} \SUPPORTEDANCH{\char"006B\char"0319} \SUPPORTEDANCH{\char"0199\char"0319} \SUPPORTEDANCH{\char"2C6A\char"0319} \SUPPORTEDANCH{\char"006C\char"0319} \SUPPORTEDANCH{\char"0234\char"0319} \SUPPORTEDANCH{\char"006E\char"0319} \SUPPORTEDANCH{\char"0235\char"0319} \SUPPORTEDANCH{\char"0070\char"0319} \SUPPORTEDANCH{\char"01A5\char"0319} \SUPPORTEDANCH{\char"0071\char"0319} \SUPPORTEDANCH{\char"02A0\char"0319} \SUPPORTEDANCH{\char"0074\char"0319} \SUPPORTEDANCH{\char"01AB\char"0319} \SUPPORTEDANCH{\char"01AD\char"0319} \SUPPORTEDANCH{\char"0236\char"0319} \SUPPORTEDANCH{\char"0075\char"0319} \SUPPORTEDANCH{\char"AB52\char"0319} \SUPPORTEDANCH{\char"0079\char"0319} \SUPPORTEDANCH{\char"01B4\char"0319} \SUPPORTEDANCH{\char"1EFF\char"0319}

\SUPPORTEDANCH{\char"0048\char"031C} \SUPPORTEDANCH{\char"2C67\char"031C} \SUPPORTEDANCH{\char"004B\char"031C} \SUPPORTEDANCH{\char"2C69\char"031C} \SUPPORTEDANCH{\char"004D\char"031C} \SUPPORTEDANCH{\char"2C6E\char"031C} \SUPPORTEDANCH{\char"0050\char"031C} \SUPPORTEDANCH{\char"01A4\char"031C} \SUPPORTEDANCH{\char"0052\char"031C} \SUPPORTEDANCH{\char"2C64\char"031C} \SUPPORTEDANCH{\char"0054\char"031C} \SUPPORTEDANCH{\char"01AC\char"031C} \SUPPORTEDANCH{\char"01AE\char"031C} \SUPPORTEDANCH{\char"0059\char"031C} \SUPPORTEDANCH{\char"01B3\char"031C} \SUPPORTEDANCH{\char"0061\char"031C} \SUPPORTEDANCH{\char"1E9A\char"031C} \SUPPORTEDANCH{\char"0062\char"031C} \SUPPORTEDANCH{\char"0183\char"031C} \SUPPORTEDANCH{\char"0063\char"031C} \SUPPORTEDANCH{\char"0188\char"031C} \SUPPORTEDANCH{\char"0064\char"031C} \SUPPORTEDANCH{\char"018C\char"031C} \SUPPORTEDANCH{\char"0221\char"031C} \SUPPORTEDANCH{\char"0068\char"031C} \SUPPORTEDANCH{\char"2C68\char"031C} \SUPPORTEDANCH{\char"006B\char"031C} \SUPPORTEDANCH{\char"0199\char"031C} \SUPPORTEDANCH{\char"2C6A\char"031C} \SUPPORTEDANCH{\char"006C\char"031C} \SUPPORTEDANCH{\char"0234\char"031C} \SUPPORTEDANCH{\char"006E\char"031C} \SUPPORTEDANCH{\char"0235\char"031C} \SUPPORTEDANCH{\char"0070\char"031C} \SUPPORTEDANCH{\char"01A5\char"031C} \SUPPORTEDANCH{\char"0071\char"031C} \SUPPORTEDANCH{\char"02A0\char"031C} \SUPPORTEDANCH{\char"0074\char"031C} \SUPPORTEDANCH{\char"01AB\char"031C} \SUPPORTEDANCH{\char"01AD\char"031C} \SUPPORTEDANCH{\char"0236\char"031C} \SUPPORTEDANCH{\char"0075\char"031C} \SUPPORTEDANCH{\char"AB52\char"031C} \SUPPORTEDANCH{\char"0079\char"031C} \SUPPORTEDANCH{\char"01B4\char"031C} \SUPPORTEDANCH{\char"1EFF\char"031C}

\SUPPORTEDANCH{\char"0048\char"0320} \SUPPORTEDANCH{\char"2C67\char"0320} \SUPPORTEDANCH{\char"004B\char"0320} \SUPPORTEDANCH{\char"2C69\char"0320} \SUPPORTEDANCH{\char"004D\char"0320} \SUPPORTEDANCH{\char"2C6E\char"0320} \SUPPORTEDANCH{\char"0050\char"0320} \SUPPORTEDANCH{\char"01A4\char"0320} \SUPPORTEDANCH{\char"0052\char"0320} \SUPPORTEDANCH{\char"2C64\char"0320} \SUPPORTEDANCH{\char"0054\char"0320} \SUPPORTEDANCH{\char"01AC\char"0320} \SUPPORTEDANCH{\char"01AE\char"0320} \SUPPORTEDANCH{\char"0059\char"0320} \SUPPORTEDANCH{\char"01B3\char"0320} \SUPPORTEDANCH{\char"0061\char"0320} \SUPPORTEDANCH{\char"1E9A\char"0320} \SUPPORTEDANCH{\char"0062\char"0320} \SUPPORTEDANCH{\char"0183\char"0320} \SUPPORTEDANCH{\char"0063\char"0320} \SUPPORTEDANCH{\char"0188\char"0320} \SUPPORTEDANCH{\char"0064\char"0320} \SUPPORTEDANCH{\char"018C\char"0320} \SUPPORTEDANCH{\char"0221\char"0320} \SUPPORTEDANCH{\char"0068\char"0320} \SUPPORTEDANCH{\char"2C68\char"0320} \SUPPORTEDANCH{\char"006B\char"0320} \SUPPORTEDANCH{\char"0199\char"0320} \SUPPORTEDANCH{\char"2C6A\char"0320} \SUPPORTEDANCH{\char"006C\char"0320} \SUPPORTEDANCH{\char"0234\char"0320} \SUPPORTEDANCH{\char"006E\char"0320} \SUPPORTEDANCH{\char"0235\char"0320} \SUPPORTEDANCH{\char"0070\char"0320} \SUPPORTEDANCH{\char"01A5\char"0320} \SUPPORTEDANCH{\char"0071\char"0320} \SUPPORTEDANCH{\char"02A0\char"0320} \SUPPORTEDANCH{\char"0074\char"0320} \SUPPORTEDANCH{\char"01AB\char"0320} \SUPPORTEDANCH{\char"01AD\char"0320} \SUPPORTEDANCH{\char"0236\char"0320} \SUPPORTEDANCH{\char"0075\char"0320} \SUPPORTEDANCH{\char"AB52\char"0320} \SUPPORTEDANCH{\char"0079\char"0320} \SUPPORTEDANCH{\char"01B4\char"0320} \SUPPORTEDANCH{\char"1EFF\char"0320}

\SUPPORTEDPREC{\char"0048\char"0323} \SUPPORTEDANCH{\char"2C67\char"0323} \SUPPORTEDPREC{\char"004B\char"0323} \SUPPORTEDANCH{\char"2C69\char"0323} \SUPPORTEDPREC{\char"004D\char"0323} \SUPPORTEDANCH{\char"2C6E\char"0323} \SUPPORTEDANCH{\char"0050\char"0323} \SUPPORTEDANCH{\char"01A4\char"0323} \SUPPORTEDPREC{\char"0052\char"0323} \SUPPORTEDANCH{\char"2C64\char"0323} \SUPPORTEDPREC{\char"0054\char"0323} \SUPPORTEDANCH{\char"01AC\char"0323} \SUPPORTEDANCH{\char"01AE\char"0323} \SUPPORTEDPREC{\char"0059\char"0323} \SUPPORTEDANCH{\char"01B3\char"0323} \SUPPORTEDPREC{\char"0061\char"0323} \SUPPORTEDANCH{\char"1E9A\char"0323} \SUPPORTEDPREC{\char"0062\char"0323} \SUPPORTEDANCH{\char"0183\char"0323} \SUPPORTEDANCH{\char"0063\char"0323} \SUPPORTEDANCH{\char"0188\char"0323} \SUPPORTEDPREC{\char"0064\char"0323} \SUPPORTEDANCH{\char"018C\char"0323} \SUPPORTEDANCH{\char"0221\char"0323} \SUPPORTEDPREC{\char"0068\char"0323} \SUPPORTEDANCH{\char"2C68\char"0323} \SUPPORTEDPREC{\char"006B\char"0323} \SUPPORTEDANCH{\char"0199\char"0323} \SUPPORTEDANCH{\char"2C6A\char"0323} \SUPPORTEDPREC{\char"006C\char"0323} \SUPPORTEDANCH{\char"0234\char"0323} \SUPPORTEDPREC{\char"006E\char"0323} \SUPPORTEDANCH{\char"0235\char"0323} \SUPPORTEDANCH{\char"0070\char"0323} \SUPPORTEDANCH{\char"01A5\char"0323} \SUPPORTEDANCH{\char"0071\char"0323} \SUPPORTEDANCH{\char"02A0\char"0323} \SUPPORTEDPREC{\char"0074\char"0323} \SUPPORTEDANCH{\char"01AB\char"0323} \SUPPORTEDANCH{\char"01AD\char"0323} \SUPPORTEDANCH{\char"0236\char"0323} \SUPPORTEDPREC{\char"0075\char"0323} \SUPPORTEDANCH{\char"AB52\char"0323} \SUPPORTEDPREC{\char"0079\char"0323} \SUPPORTEDANCH{\char"01B4\char"0323} \SUPPORTEDANCH{\char"1EFF\char"0323}

\SUPPORTEDANCH{\char"0048\char"0324} \SUPPORTEDANCH{\char"2C67\char"0324} \SUPPORTEDANCH{\char"004B\char"0324} \SUPPORTEDANCH{\char"2C69\char"0324} \SUPPORTEDANCH{\char"004D\char"0324} \SUPPORTEDANCH{\char"2C6E\char"0324} \SUPPORTEDANCH{\char"0050\char"0324} \SUPPORTEDANCH{\char"01A4\char"0324} \SUPPORTEDANCH{\char"0052\char"0324} \SUPPORTEDANCH{\char"2C64\char"0324} \SUPPORTEDANCH{\char"0054\char"0324} \SUPPORTEDANCH{\char"01AC\char"0324} \SUPPORTEDANCH{\char"01AE\char"0324} \SUPPORTEDANCH{\char"0059\char"0324} \SUPPORTEDANCH{\char"01B3\char"0324} \SUPPORTEDANCH{\char"0061\char"0324} \SUPPORTEDANCH{\char"1E9A\char"0324} \SUPPORTEDANCH{\char"0062\char"0324} \SUPPORTEDANCH{\char"0183\char"0324} \SUPPORTEDANCH{\char"0063\char"0324} \SUPPORTEDANCH{\char"0188\char"0324} \SUPPORTEDANCH{\char"0064\char"0324} \SUPPORTEDANCH{\char"018C\char"0324} \SUPPORTEDANCH{\char"0221\char"0324} \SUPPORTEDANCH{\char"0068\char"0324} \SUPPORTEDANCH{\char"2C68\char"0324} \SUPPORTEDANCH{\char"006B\char"0324} \SUPPORTEDANCH{\char"0199\char"0324} \SUPPORTEDANCH{\char"2C6A\char"0324} \SUPPORTEDANCH{\char"006C\char"0324} \SUPPORTEDANCH{\char"0234\char"0324} \SUPPORTEDANCH{\char"006E\char"0324} \SUPPORTEDANCH{\char"0235\char"0324} \SUPPORTEDANCH{\char"0070\char"0324} \SUPPORTEDANCH{\char"01A5\char"0324} \SUPPORTEDANCH{\char"0071\char"0324} \SUPPORTEDANCH{\char"02A0\char"0324} \SUPPORTEDANCH{\char"0074\char"0324} \SUPPORTEDANCH{\char"01AB\char"0324} \SUPPORTEDANCH{\char"01AD\char"0324} \SUPPORTEDANCH{\char"0236\char"0324} \SUPPORTEDPREC{\char"0075\char"0324} \SUPPORTEDANCH{\char"AB52\char"0324} \SUPPORTEDANCH{\char"0079\char"0324} \SUPPORTEDANCH{\char"01B4\char"0324} \SUPPORTEDANCH{\char"1EFF\char"0324}

\SUPPORTEDANCH{\char"0048\char"0325} \SUPPORTEDANCH{\char"2C67\char"0325} \SUPPORTEDANCH{\char"004B\char"0325} \SUPPORTEDANCH{\char"2C69\char"0325} \SUPPORTEDANCH{\char"004D\char"0325} \SUPPORTEDANCH{\char"2C6E\char"0325} \SUPPORTEDANCH{\char"0050\char"0325} \SUPPORTEDANCH{\char"01A4\char"0325} \SUPPORTEDANCH{\char"0052\char"0325} \SUPPORTEDANCH{\char"2C64\char"0325} \SUPPORTEDANCH{\char"0054\char"0325} \SUPPORTEDANCH{\char"01AC\char"0325} \SUPPORTEDANCH{\char"01AE\char"0325} \SUPPORTEDANCH{\char"0059\char"0325} \SUPPORTEDANCH{\char"01B3\char"0325} \SUPPORTEDPREC{\char"0061\char"0325} \SUPPORTEDANCH{\char"1E9A\char"0325} \SUPPORTEDANCH{\char"0062\char"0325} \SUPPORTEDANCH{\char"0183\char"0325} \SUPPORTEDANCH{\char"0063\char"0325} \SUPPORTEDANCH{\char"0188\char"0325} \SUPPORTEDANCH{\char"0064\char"0325} \SUPPORTEDANCH{\char"018C\char"0325} \SUPPORTEDANCH{\char"0221\char"0325} \SUPPORTEDANCH{\char"0068\char"0325} \SUPPORTEDANCH{\char"2C68\char"0325} \SUPPORTEDANCH{\char"006B\char"0325} \SUPPORTEDANCH{\char"0199\char"0325} \SUPPORTEDANCH{\char"2C6A\char"0325} \SUPPORTEDANCH{\char"006C\char"0325} \SUPPORTEDANCH{\char"0234\char"0325} \SUPPORTEDANCH{\char"006E\char"0325} \SUPPORTEDANCH{\char"0235\char"0325} \SUPPORTEDANCH{\char"0070\char"0325} \SUPPORTEDANCH{\char"01A5\char"0325} \SUPPORTEDANCH{\char"0071\char"0325} \SUPPORTEDANCH{\char"02A0\char"0325} \SUPPORTEDANCH{\char"0074\char"0325} \SUPPORTEDANCH{\char"01AB\char"0325} \SUPPORTEDANCH{\char"01AD\char"0325} \SUPPORTEDANCH{\char"0236\char"0325} \SUPPORTEDANCH{\char"0075\char"0325} \SUPPORTEDANCH{\char"AB52\char"0325} \SUPPORTEDANCH{\char"0079\char"0325} \SUPPORTEDANCH{\char"01B4\char"0325} \SUPPORTEDANCH{\char"1EFF\char"0325}

\SUPPORTEDANCH{\char"0048\char"0326} \SUPPORTEDANCH{\char"2C67\char"0326} \SUPPORTEDANCH{\char"004B\char"0326} \SUPPORTEDANCH{\char"2C69\char"0326} \SUPPORTEDANCH{\char"004D\char"0326} \SUPPORTEDANCH{\char"2C6E\char"0326} \SUPPORTEDANCH{\char"0050\char"0326} \SUPPORTEDANCH{\char"01A4\char"0326} \SUPPORTEDANCH{\char"0052\char"0326} \SUPPORTEDANCH{\char"2C64\char"0326} \SUPPORTEDPREC{\char"0054\char"0326} \SUPPORTEDANCH{\char"01AC\char"0326} \SUPPORTEDANCH{\char"01AE\char"0326} \SUPPORTEDANCH{\char"0059\char"0326} \SUPPORTEDANCH{\char"01B3\char"0326} \SUPPORTEDANCH{\char"0061\char"0326} \SUPPORTEDANCH{\char"1E9A\char"0326} \SUPPORTEDANCH{\char"0062\char"0326} \SUPPORTEDANCH{\char"0183\char"0326} \SUPPORTEDANCH{\char"0063\char"0326} \SUPPORTEDANCH{\char"0188\char"0326} \SUPPORTEDANCH{\char"0064\char"0326} \SUPPORTEDANCH{\char"018C\char"0326} \SUPPORTEDANCH{\char"0221\char"0326} \SUPPORTEDANCH{\char"0068\char"0326} \SUPPORTEDANCH{\char"2C68\char"0326} \SUPPORTEDANCH{\char"006B\char"0326} \SUPPORTEDANCH{\char"0199\char"0326} \SUPPORTEDANCH{\char"2C6A\char"0326} \SUPPORTEDANCH{\char"006C\char"0326} \SUPPORTEDANCH{\char"0234\char"0326} \SUPPORTEDANCH{\char"006E\char"0326} \SUPPORTEDANCH{\char"0235\char"0326} \SUPPORTEDANCH{\char"0070\char"0326} \SUPPORTEDANCH{\char"01A5\char"0326} \SUPPORTEDANCH{\char"0071\char"0326} \SUPPORTEDANCH{\char"02A0\char"0326} \SUPPORTEDPREC{\char"0074\char"0326} \SUPPORTEDANCH{\char"01AB\char"0326} \SUPPORTEDANCH{\char"01AD\char"0326} \SUPPORTEDANCH{\char"0236\char"0326} \SUPPORTEDANCH{\char"0075\char"0326} \SUPPORTEDANCH{\char"AB52\char"0326} \SUPPORTEDANCH{\char"0079\char"0326} \SUPPORTEDANCH{\char"01B4\char"0326} \SUPPORTEDANCH{\char"1EFF\char"0326}

\SUPPORTEDANCH{\char"0048\char"0329} \SUPPORTEDANCH{\char"2C67\char"0329} \SUPPORTEDANCH{\char"004B\char"0329} \SUPPORTEDANCH{\char"2C69\char"0329} \SUPPORTEDANCH{\char"004D\char"0329} \SUPPORTEDANCH{\char"2C6E\char"0329} \SUPPORTEDANCH{\char"0050\char"0329} \SUPPORTEDANCH{\char"01A4\char"0329} \SUPPORTEDANCH{\char"0052\char"0329} \SUPPORTEDANCH{\char"2C64\char"0329} \SUPPORTEDANCH{\char"0054\char"0329} \SUPPORTEDANCH{\char"01AC\char"0329} \SUPPORTEDANCH{\char"01AE\char"0329} \SUPPORTEDANCH{\char"0059\char"0329} \SUPPORTEDANCH{\char"01B3\char"0329} \SUPPORTEDANCH{\char"0061\char"0329} \SUPPORTEDANCH{\char"1E9A\char"0329} \SUPPORTEDANCH{\char"0062\char"0329} \SUPPORTEDANCH{\char"0183\char"0329} \SUPPORTEDANCH{\char"0063\char"0329} \SUPPORTEDANCH{\char"0188\char"0329} \SUPPORTEDANCH{\char"0064\char"0329} \SUPPORTEDANCH{\char"018C\char"0329} \SUPPORTEDANCH{\char"0221\char"0329} \SUPPORTEDANCH{\char"0068\char"0329} \SUPPORTEDANCH{\char"2C68\char"0329} \SUPPORTEDANCH{\char"006B\char"0329} \SUPPORTEDANCH{\char"0199\char"0329} \SUPPORTEDANCH{\char"2C6A\char"0329} \SUPPORTEDANCH{\char"006C\char"0329} \SUPPORTEDANCH{\char"0234\char"0329} \SUPPORTEDANCH{\char"006E\char"0329} \SUPPORTEDANCH{\char"0235\char"0329} \SUPPORTEDANCH{\char"0070\char"0329} \SUPPORTEDANCH{\char"01A5\char"0329} \SUPPORTEDANCH{\char"0071\char"0329} \SUPPORTEDANCH{\char"02A0\char"0329} \SUPPORTEDANCH{\char"0074\char"0329} \SUPPORTEDANCH{\char"01AB\char"0329} \SUPPORTEDANCH{\char"01AD\char"0329} \SUPPORTEDANCH{\char"0236\char"0329} \SUPPORTEDANCH{\char"0075\char"0329} \SUPPORTEDANCH{\char"AB52\char"0329} \SUPPORTEDANCH{\char"0079\char"0329} \SUPPORTEDANCH{\char"01B4\char"0329} \SUPPORTEDANCH{\char"1EFF\char"0329}

\SUPPORTEDANCH{\char"0048\char"032A} \SUPPORTEDANCH{\char"2C67\char"032A} \SUPPORTEDANCH{\char"004B\char"032A} \SUPPORTEDANCH{\char"2C69\char"032A} \SUPPORTEDANCH{\char"004D\char"032A} \SUPPORTEDANCH{\char"2C6E\char"032A} \SUPPORTEDANCH{\char"0050\char"032A} \SUPPORTEDANCH{\char"01A4\char"032A} \SUPPORTEDANCH{\char"0052\char"032A} \SUPPORTEDANCH{\char"2C64\char"032A} \SUPPORTEDANCH{\char"0054\char"032A} \SUPPORTEDANCH{\char"01AC\char"032A} \SUPPORTEDANCH{\char"01AE\char"032A} \SUPPORTEDANCH{\char"0059\char"032A} \SUPPORTEDANCH{\char"01B3\char"032A} \SUPPORTEDANCH{\char"0061\char"032A} \SUPPORTEDANCH{\char"1E9A\char"032A} \SUPPORTEDANCH{\char"0062\char"032A} \SUPPORTEDANCH{\char"0183\char"032A} \SUPPORTEDANCH{\char"0063\char"032A} \SUPPORTEDANCH{\char"0188\char"032A} \SUPPORTEDANCH{\char"0064\char"032A} \SUPPORTEDANCH{\char"018C\char"032A} \SUPPORTEDANCH{\char"0221\char"032A} \SUPPORTEDANCH{\char"0068\char"032A} \SUPPORTEDANCH{\char"2C68\char"032A} \SUPPORTEDANCH{\char"006B\char"032A} \SUPPORTEDANCH{\char"0199\char"032A} \SUPPORTEDANCH{\char"2C6A\char"032A} \SUPPORTEDANCH{\char"006C\char"032A} \SUPPORTEDANCH{\char"0234\char"032A} \SUPPORTEDANCH{\char"006E\char"032A} \SUPPORTEDANCH{\char"0235\char"032A} \SUPPORTEDANCH{\char"0070\char"032A} \SUPPORTEDANCH{\char"01A5\char"032A} \SUPPORTEDANCH{\char"0071\char"032A} \SUPPORTEDANCH{\char"02A0\char"032A} \SUPPORTEDANCH{\char"0074\char"032A} \SUPPORTEDANCH{\char"01AB\char"032A} \SUPPORTEDANCH{\char"01AD\char"032A} \SUPPORTEDANCH{\char"0236\char"032A} \SUPPORTEDANCH{\char"0075\char"032A} \SUPPORTEDANCH{\char"AB52\char"032A} \SUPPORTEDANCH{\char"0079\char"032A} \SUPPORTEDANCH{\char"01B4\char"032A} \SUPPORTEDANCH{\char"1EFF\char"032A}

\SUPPORTEDANCH{\char"0048\char"032B} \SUPPORTEDANCH{\char"2C67\char"032B} \SUPPORTEDANCH{\char"004B\char"032B} \SUPPORTEDANCH{\char"2C69\char"032B} \SUPPORTEDANCH{\char"004D\char"032B} \SUPPORTEDANCH{\char"2C6E\char"032B} \SUPPORTEDANCH{\char"0050\char"032B} \SUPPORTEDANCH{\char"01A4\char"032B} \SUPPORTEDANCH{\char"0052\char"032B} \SUPPORTEDANCH{\char"2C64\char"032B} \SUPPORTEDANCH{\char"0054\char"032B} \SUPPORTEDANCH{\char"01AC\char"032B} \SUPPORTEDANCH{\char"01AE\char"032B} \SUPPORTEDANCH{\char"0059\char"032B} \SUPPORTEDANCH{\char"01B3\char"032B} \SUPPORTEDANCH{\char"0061\char"032B} \SUPPORTEDANCH{\char"1E9A\char"032B} \SUPPORTEDANCH{\char"0062\char"032B} \SUPPORTEDANCH{\char"0183\char"032B} \SUPPORTEDANCH{\char"0063\char"032B} \SUPPORTEDANCH{\char"0188\char"032B} \SUPPORTEDANCH{\char"0064\char"032B} \SUPPORTEDANCH{\char"018C\char"032B} \SUPPORTEDANCH{\char"0221\char"032B} \SUPPORTEDANCH{\char"0068\char"032B} \SUPPORTEDANCH{\char"2C68\char"032B} \SUPPORTEDANCH{\char"006B\char"032B} \SUPPORTEDANCH{\char"0199\char"032B} \SUPPORTEDANCH{\char"2C6A\char"032B} \SUPPORTEDANCH{\char"006C\char"032B} \SUPPORTEDANCH{\char"0234\char"032B} \SUPPORTEDANCH{\char"006E\char"032B} \SUPPORTEDANCH{\char"0235\char"032B} \SUPPORTEDANCH{\char"0070\char"032B} \SUPPORTEDANCH{\char"01A5\char"032B} \SUPPORTEDANCH{\char"0071\char"032B} \SUPPORTEDANCH{\char"02A0\char"032B} \SUPPORTEDANCH{\char"0074\char"032B} \SUPPORTEDANCH{\char"01AB\char"032B} \SUPPORTEDANCH{\char"01AD\char"032B} \SUPPORTEDANCH{\char"0236\char"032B} \SUPPORTEDANCH{\char"0075\char"032B} \SUPPORTEDANCH{\char"AB52\char"032B} \SUPPORTEDANCH{\char"0079\char"032B} \SUPPORTEDANCH{\char"01B4\char"032B} \SUPPORTEDANCH{\char"1EFF\char"032B}

\SUPPORTEDANCH{\char"0048\char"032C} \SUPPORTEDANCH{\char"2C67\char"032C} \SUPPORTEDANCH{\char"004B\char"032C} \SUPPORTEDANCH{\char"2C69\char"032C} \SUPPORTEDANCH{\char"004D\char"032C} \SUPPORTEDANCH{\char"2C6E\char"032C} \SUPPORTEDANCH{\char"0050\char"032C} \SUPPORTEDANCH{\char"01A4\char"032C} \SUPPORTEDANCH{\char"0052\char"032C} \SUPPORTEDANCH{\char"2C64\char"032C} \SUPPORTEDANCH{\char"0054\char"032C} \SUPPORTEDANCH{\char"01AC\char"032C} \SUPPORTEDANCH{\char"01AE\char"032C} \SUPPORTEDANCH{\char"0059\char"032C} \SUPPORTEDANCH{\char"01B3\char"032C} \SUPPORTEDANCH{\char"0061\char"032C} \SUPPORTEDANCH{\char"1E9A\char"032C} \SUPPORTEDANCH{\char"0062\char"032C} \SUPPORTEDANCH{\char"0183\char"032C} \SUPPORTEDANCH{\char"0063\char"032C} \SUPPORTEDANCH{\char"0188\char"032C} \SUPPORTEDANCH{\char"0064\char"032C} \SUPPORTEDANCH{\char"018C\char"032C} \SUPPORTEDANCH{\char"0221\char"032C} \SUPPORTEDANCH{\char"0068\char"032C} \SUPPORTEDANCH{\char"2C68\char"032C} \SUPPORTEDANCH{\char"006B\char"032C} \SUPPORTEDANCH{\char"0199\char"032C} \SUPPORTEDANCH{\char"2C6A\char"032C} \SUPPORTEDANCH{\char"006C\char"032C} \SUPPORTEDANCH{\char"0234\char"032C} \SUPPORTEDANCH{\char"006E\char"032C} \SUPPORTEDANCH{\char"0235\char"032C} \SUPPORTEDANCH{\char"0070\char"032C} \SUPPORTEDANCH{\char"01A5\char"032C} \SUPPORTEDANCH{\char"0071\char"032C} \SUPPORTEDANCH{\char"02A0\char"032C} \SUPPORTEDANCH{\char"0074\char"032C} \SUPPORTEDANCH{\char"01AB\char"032C} \SUPPORTEDANCH{\char"01AD\char"032C} \SUPPORTEDANCH{\char"0236\char"032C} \SUPPORTEDANCH{\char"0075\char"032C} \SUPPORTEDANCH{\char"AB52\char"032C} \SUPPORTEDANCH{\char"0079\char"032C} \SUPPORTEDANCH{\char"01B4\char"032C} \SUPPORTEDANCH{\char"1EFF\char"032C}

\SUPPORTEDANCH{\char"0048\char"032D} \SUPPORTEDANCH{\char"2C67\char"032D} \SUPPORTEDANCH{\char"004B\char"032D} \SUPPORTEDANCH{\char"2C69\char"032D} \SUPPORTEDANCH{\char"004D\char"032D} \SUPPORTEDANCH{\char"2C6E\char"032D} \SUPPORTEDANCH{\char"0050\char"032D} \SUPPORTEDANCH{\char"01A4\char"032D} \SUPPORTEDANCH{\char"0052\char"032D} \SUPPORTEDANCH{\char"2C64\char"032D} \SUPPORTEDPREC{\char"0054\char"032D} \SUPPORTEDANCH{\char"01AC\char"032D} \SUPPORTEDANCH{\char"01AE\char"032D} \SUPPORTEDANCH{\char"0059\char"032D} \SUPPORTEDANCH{\char"01B3\char"032D} \SUPPORTEDANCH{\char"0061\char"032D} \SUPPORTEDANCH{\char"1E9A\char"032D} \SUPPORTEDANCH{\char"0062\char"032D} \SUPPORTEDANCH{\char"0183\char"032D} \SUPPORTEDANCH{\char"0063\char"032D} \SUPPORTEDANCH{\char"0188\char"032D} \SUPPORTEDPREC{\char"0064\char"032D} \SUPPORTEDANCH{\char"018C\char"032D} \SUPPORTEDANCH{\char"0221\char"032D} \SUPPORTEDANCH{\char"0068\char"032D} \SUPPORTEDANCH{\char"2C68\char"032D} \SUPPORTEDANCH{\char"006B\char"032D} \SUPPORTEDANCH{\char"0199\char"032D} \SUPPORTEDANCH{\char"2C6A\char"032D} \SUPPORTEDPREC{\char"006C\char"032D} \SUPPORTEDANCH{\char"0234\char"032D} \SUPPORTEDPREC{\char"006E\char"032D} \SUPPORTEDANCH{\char"0235\char"032D} \SUPPORTEDANCH{\char"0070\char"032D} \SUPPORTEDANCH{\char"01A5\char"032D} \SUPPORTEDANCH{\char"0071\char"032D} \SUPPORTEDANCH{\char"02A0\char"032D} \SUPPORTEDPREC{\char"0074\char"032D} \SUPPORTEDANCH{\char"01AB\char"032D} \SUPPORTEDANCH{\char"01AD\char"032D} \SUPPORTEDANCH{\char"0236\char"032D} \SUPPORTEDPREC{\char"0075\char"032D} \SUPPORTEDANCH{\char"AB52\char"032D} \SUPPORTEDANCH{\char"0079\char"032D} \SUPPORTEDANCH{\char"01B4\char"032D} \SUPPORTEDANCH{\char"1EFF\char"032D}

\SUPPORTEDPREC{\char"0048\char"032E} \SUPPORTEDANCH{\char"2C67\char"032E} \SUPPORTEDANCH{\char"004B\char"032E} \SUPPORTEDANCH{\char"2C69\char"032E} \SUPPORTEDANCH{\char"004D\char"032E} \SUPPORTEDANCH{\char"2C6E\char"032E} \SUPPORTEDANCH{\char"0050\char"032E} \SUPPORTEDANCH{\char"01A4\char"032E} \SUPPORTEDANCH{\char"0052\char"032E} \SUPPORTEDANCH{\char"2C64\char"032E} \SUPPORTEDANCH{\char"0054\char"032E} \SUPPORTEDANCH{\char"01AC\char"032E} \SUPPORTEDANCH{\char"01AE\char"032E} \SUPPORTEDANCH{\char"0059\char"032E} \SUPPORTEDANCH{\char"01B3\char"032E} \SUPPORTEDANCH{\char"0061\char"032E} \SUPPORTEDANCH{\char"1E9A\char"032E} \SUPPORTEDANCH{\char"0062\char"032E} \SUPPORTEDANCH{\char"0183\char"032E} \SUPPORTEDANCH{\char"0063\char"032E} \SUPPORTEDANCH{\char"0188\char"032E} \SUPPORTEDANCH{\char"0064\char"032E} \SUPPORTEDANCH{\char"018C\char"032E} \SUPPORTEDANCH{\char"0221\char"032E} \SUPPORTEDPREC{\char"0068\char"032E} \SUPPORTEDANCH{\char"2C68\char"032E} \SUPPORTEDANCH{\char"006B\char"032E} \SUPPORTEDANCH{\char"0199\char"032E} \SUPPORTEDANCH{\char"2C6A\char"032E} \SUPPORTEDANCH{\char"006C\char"032E} \SUPPORTEDANCH{\char"0234\char"032E} \SUPPORTEDANCH{\char"006E\char"032E} \SUPPORTEDANCH{\char"0235\char"032E} \SUPPORTEDANCH{\char"0070\char"032E} \SUPPORTEDANCH{\char"01A5\char"032E} \SUPPORTEDANCH{\char"0071\char"032E} \SUPPORTEDANCH{\char"02A0\char"032E} \SUPPORTEDANCH{\char"0074\char"032E} \SUPPORTEDANCH{\char"01AB\char"032E} \SUPPORTEDANCH{\char"01AD\char"032E} \SUPPORTEDANCH{\char"0236\char"032E} \SUPPORTEDANCH{\char"0075\char"032E} \SUPPORTEDANCH{\char"AB52\char"032E} \SUPPORTEDANCH{\char"0079\char"032E} \SUPPORTEDANCH{\char"01B4\char"032E} \SUPPORTEDANCH{\char"1EFF\char"032E}

\SUPPORTEDANCH{\char"0048\char"032F} \SUPPORTEDANCH{\char"2C67\char"032F} \SUPPORTEDANCH{\char"004B\char"032F} \SUPPORTEDANCH{\char"2C69\char"032F} \SUPPORTEDANCH{\char"004D\char"032F} \SUPPORTEDANCH{\char"2C6E\char"032F} \SUPPORTEDANCH{\char"0050\char"032F} \SUPPORTEDANCH{\char"01A4\char"032F} \SUPPORTEDANCH{\char"0052\char"032F} \SUPPORTEDANCH{\char"2C64\char"032F} \SUPPORTEDANCH{\char"0054\char"032F} \SUPPORTEDANCH{\char"01AC\char"032F} \SUPPORTEDANCH{\char"01AE\char"032F} \SUPPORTEDANCH{\char"0059\char"032F} \SUPPORTEDANCH{\char"01B3\char"032F} \SUPPORTEDANCH{\char"0061\char"032F} \SUPPORTEDANCH{\char"1E9A\char"032F} \SUPPORTEDANCH{\char"0062\char"032F} \SUPPORTEDANCH{\char"0183\char"032F} \SUPPORTEDANCH{\char"0063\char"032F} \SUPPORTEDANCH{\char"0188\char"032F} \SUPPORTEDANCH{\char"0064\char"032F} \SUPPORTEDANCH{\char"018C\char"032F} \SUPPORTEDANCH{\char"0221\char"032F} \SUPPORTEDANCH{\char"0068\char"032F} \SUPPORTEDANCH{\char"2C68\char"032F} \SUPPORTEDANCH{\char"006B\char"032F} \SUPPORTEDANCH{\char"0199\char"032F} \SUPPORTEDANCH{\char"2C6A\char"032F} \SUPPORTEDANCH{\char"006C\char"032F} \SUPPORTEDANCH{\char"0234\char"032F} \SUPPORTEDANCH{\char"006E\char"032F} \SUPPORTEDANCH{\char"0235\char"032F} \SUPPORTEDANCH{\char"0070\char"032F} \SUPPORTEDANCH{\char"01A5\char"032F} \SUPPORTEDANCH{\char"0071\char"032F} \SUPPORTEDANCH{\char"02A0\char"032F} \SUPPORTEDANCH{\char"0074\char"032F} \SUPPORTEDANCH{\char"01AB\char"032F} \SUPPORTEDANCH{\char"01AD\char"032F} \SUPPORTEDANCH{\char"0236\char"032F} \SUPPORTEDANCH{\char"0075\char"032F} \SUPPORTEDANCH{\char"AB52\char"032F} \SUPPORTEDANCH{\char"0079\char"032F} \SUPPORTEDANCH{\char"01B4\char"032F} \SUPPORTEDANCH{\char"1EFF\char"032F}

\SUPPORTEDANCH{\char"0048\char"031D} \SUPPORTEDANCH{\char"2C67\char"031D} \SUPPORTEDANCH{\char"004B\char"031D} \SUPPORTEDANCH{\char"2C69\char"031D} \SUPPORTEDANCH{\char"004D\char"031D} \SUPPORTEDANCH{\char"2C6E\char"031D} \SUPPORTEDANCH{\char"0050\char"031D} \SUPPORTEDANCH{\char"01A4\char"031D} \SUPPORTEDANCH{\char"0052\char"031D} \SUPPORTEDANCH{\char"2C64\char"031D} \SUPPORTEDANCH{\char"0054\char"031D} \SUPPORTEDANCH{\char"01AC\char"031D} \SUPPORTEDANCH{\char"01AE\char"031D} \SUPPORTEDANCH{\char"0059\char"031D} \SUPPORTEDANCH{\char"01B3\char"031D} \SUPPORTEDANCH{\char"0061\char"031D} \SUPPORTEDANCH{\char"1E9A\char"031D} \SUPPORTEDANCH{\char"0062\char"031D} \SUPPORTEDANCH{\char"0183\char"031D} \SUPPORTEDANCH{\char"0063\char"031D} \SUPPORTEDANCH{\char"0188\char"031D} \SUPPORTEDANCH{\char"0064\char"031D} \SUPPORTEDANCH{\char"018C\char"031D} \SUPPORTEDANCH{\char"0221\char"031D} \SUPPORTEDANCH{\char"0068\char"031D} \SUPPORTEDANCH{\char"2C68\char"031D} \SUPPORTEDANCH{\char"006B\char"031D} \SUPPORTEDANCH{\char"0199\char"031D} \SUPPORTEDANCH{\char"2C6A\char"031D} \SUPPORTEDANCH{\char"006C\char"031D} \SUPPORTEDANCH{\char"0234\char"031D} \SUPPORTEDANCH{\char"006E\char"031D} \SUPPORTEDANCH{\char"0235\char"031D} \SUPPORTEDANCH{\char"0070\char"031D} \SUPPORTEDANCH{\char"01A5\char"031D} \SUPPORTEDANCH{\char"0071\char"031D} \SUPPORTEDANCH{\char"02A0\char"031D} \SUPPORTEDANCH{\char"0074\char"031D} \SUPPORTEDANCH{\char"01AB\char"031D} \SUPPORTEDANCH{\char"01AD\char"031D} \SUPPORTEDANCH{\char"0236\char"031D} \SUPPORTEDANCH{\char"0075\char"031D} \SUPPORTEDANCH{\char"AB52\char"031D} \SUPPORTEDANCH{\char"0079\char"031D} \SUPPORTEDANCH{\char"01B4\char"031D} \SUPPORTEDANCH{\char"1EFF\char"031D}

\SUPPORTEDANCH{\char"0048\char"031E} \SUPPORTEDANCH{\char"2C67\char"031E} \SUPPORTEDANCH{\char"004B\char"031E} \SUPPORTEDANCH{\char"2C69\char"031E} \SUPPORTEDANCH{\char"004D\char"031E} \SUPPORTEDANCH{\char"2C6E\char"031E} \SUPPORTEDANCH{\char"0050\char"031E} \SUPPORTEDANCH{\char"01A4\char"031E} \SUPPORTEDANCH{\char"0052\char"031E} \SUPPORTEDANCH{\char"2C64\char"031E} \SUPPORTEDANCH{\char"0054\char"031E} \SUPPORTEDANCH{\char"01AC\char"031E} \SUPPORTEDANCH{\char"01AE\char"031E} \SUPPORTEDANCH{\char"0059\char"031E} \SUPPORTEDANCH{\char"01B3\char"031E} \SUPPORTEDANCH{\char"0061\char"031E} \SUPPORTEDANCH{\char"1E9A\char"031E} \SUPPORTEDANCH{\char"0062\char"031E} \SUPPORTEDANCH{\char"0183\char"031E} \SUPPORTEDANCH{\char"0063\char"031E} \SUPPORTEDANCH{\char"0188\char"031E} \SUPPORTEDANCH{\char"0064\char"031E} \SUPPORTEDANCH{\char"018C\char"031E} \SUPPORTEDANCH{\char"0221\char"031E} \SUPPORTEDANCH{\char"0068\char"031E} \SUPPORTEDANCH{\char"2C68\char"031E} \SUPPORTEDANCH{\char"006B\char"031E} \SUPPORTEDANCH{\char"0199\char"031E} \SUPPORTEDANCH{\char"2C6A\char"031E} \SUPPORTEDANCH{\char"006C\char"031E} \SUPPORTEDANCH{\char"0234\char"031E} \SUPPORTEDANCH{\char"006E\char"031E} \SUPPORTEDANCH{\char"0235\char"031E} \SUPPORTEDANCH{\char"0070\char"031E} \SUPPORTEDANCH{\char"01A5\char"031E} \SUPPORTEDANCH{\char"0071\char"031E} \SUPPORTEDANCH{\char"02A0\char"031E} \SUPPORTEDANCH{\char"0074\char"031E} \SUPPORTEDANCH{\char"01AB\char"031E} \SUPPORTEDANCH{\char"01AD\char"031E} \SUPPORTEDANCH{\char"0236\char"031E} \SUPPORTEDANCH{\char"0075\char"031E} \SUPPORTEDANCH{\char"AB52\char"031E} \SUPPORTEDANCH{\char"0079\char"031E} \SUPPORTEDANCH{\char"01B4\char"031E} \SUPPORTEDANCH{\char"1EFF\char"031E}

\SUPPORTEDANCH{\char"0048\char"031F} \SUPPORTEDANCH{\char"2C67\char"031F} \SUPPORTEDANCH{\char"004B\char"031F} \SUPPORTEDANCH{\char"2C69\char"031F} \SUPPORTEDANCH{\char"004D\char"031F} \SUPPORTEDANCH{\char"2C6E\char"031F} \SUPPORTEDANCH{\char"0050\char"031F} \SUPPORTEDANCH{\char"01A4\char"031F} \SUPPORTEDANCH{\char"0052\char"031F} \SUPPORTEDANCH{\char"2C64\char"031F} \SUPPORTEDANCH{\char"0054\char"031F} \SUPPORTEDANCH{\char"01AC\char"031F} \SUPPORTEDANCH{\char"01AE\char"031F} \SUPPORTEDANCH{\char"0059\char"031F} \SUPPORTEDANCH{\char"01B3\char"031F} \SUPPORTEDANCH{\char"0061\char"031F} \SUPPORTEDANCH{\char"1E9A\char"031F} \SUPPORTEDANCH{\char"0062\char"031F} \SUPPORTEDANCH{\char"0183\char"031F} \SUPPORTEDANCH{\char"0063\char"031F} \SUPPORTEDANCH{\char"0188\char"031F} \SUPPORTEDANCH{\char"0064\char"031F} \SUPPORTEDANCH{\char"018C\char"031F} \SUPPORTEDANCH{\char"0221\char"031F} \SUPPORTEDANCH{\char"0068\char"031F} \SUPPORTEDANCH{\char"2C68\char"031F} \SUPPORTEDANCH{\char"006B\char"031F} \SUPPORTEDANCH{\char"0199\char"031F} \SUPPORTEDANCH{\char"2C6A\char"031F} \SUPPORTEDANCH{\char"006C\char"031F} \SUPPORTEDANCH{\char"0234\char"031F} \SUPPORTEDANCH{\char"006E\char"031F} \SUPPORTEDANCH{\char"0235\char"031F} \SUPPORTEDANCH{\char"0070\char"031F} \SUPPORTEDANCH{\char"01A5\char"031F} \SUPPORTEDANCH{\char"0071\char"031F} \SUPPORTEDANCH{\char"02A0\char"031F} \SUPPORTEDANCH{\char"0074\char"031F} \SUPPORTEDANCH{\char"01AB\char"031F} \SUPPORTEDANCH{\char"01AD\char"031F} \SUPPORTEDANCH{\char"0236\char"031F} \SUPPORTEDANCH{\char"0075\char"031F} \SUPPORTEDANCH{\char"AB52\char"031F} \SUPPORTEDANCH{\char"0079\char"031F} \SUPPORTEDANCH{\char"01B4\char"031F} \SUPPORTEDANCH{\char"1EFF\char"031F}

\SUPPORTEDANCH{\char"0048\char"0330} \SUPPORTEDANCH{\char"2C67\char"0330} \SUPPORTEDANCH{\char"004B\char"0330} \SUPPORTEDANCH{\char"2C69\char"0330} \SUPPORTEDANCH{\char"004D\char"0330} \SUPPORTEDANCH{\char"2C6E\char"0330} \SUPPORTEDANCH{\char"0050\char"0330} \SUPPORTEDANCH{\char"01A4\char"0330} \SUPPORTEDANCH{\char"0052\char"0330} \SUPPORTEDANCH{\char"2C64\char"0330} \SUPPORTEDANCH{\char"0054\char"0330} \SUPPORTEDANCH{\char"01AC\char"0330} \SUPPORTEDANCH{\char"01AE\char"0330} \SUPPORTEDANCH{\char"0059\char"0330} \SUPPORTEDANCH{\char"01B3\char"0330} \SUPPORTEDANCH{\char"0061\char"0330} \SUPPORTEDANCH{\char"1E9A\char"0330} \SUPPORTEDANCH{\char"0062\char"0330} \SUPPORTEDANCH{\char"0183\char"0330} \SUPPORTEDANCH{\char"0063\char"0330} \SUPPORTEDANCH{\char"0188\char"0330} \SUPPORTEDANCH{\char"0064\char"0330} \SUPPORTEDANCH{\char"018C\char"0330} \SUPPORTEDANCH{\char"0221\char"0330} \SUPPORTEDANCH{\char"0068\char"0330} \SUPPORTEDANCH{\char"2C68\char"0330} \SUPPORTEDANCH{\char"006B\char"0330} \SUPPORTEDANCH{\char"0199\char"0330} \SUPPORTEDANCH{\char"2C6A\char"0330} \SUPPORTEDANCH{\char"006C\char"0330} \SUPPORTEDANCH{\char"0234\char"0330} \SUPPORTEDANCH{\char"006E\char"0330} \SUPPORTEDANCH{\char"0235\char"0330} \SUPPORTEDANCH{\char"0070\char"0330} \SUPPORTEDANCH{\char"01A5\char"0330} \SUPPORTEDANCH{\char"0071\char"0330} \SUPPORTEDANCH{\char"02A0\char"0330} \SUPPORTEDANCH{\char"0074\char"0330} \SUPPORTEDANCH{\char"01AB\char"0330} \SUPPORTEDANCH{\char"01AD\char"0330} \SUPPORTEDANCH{\char"0236\char"0330} \SUPPORTEDPREC{\char"0075\char"0330} \SUPPORTEDANCH{\char"AB52\char"0330} \SUPPORTEDANCH{\char"0079\char"0330} \SUPPORTEDANCH{\char"01B4\char"0330} \SUPPORTEDANCH{\char"1EFF\char"0330}

\SUPPORTEDANCH{\char"0048\char"0331} \SUPPORTEDANCH{\char"2C67\char"0331} \SUPPORTEDPREC{\char"004B\char"0331} \SUPPORTEDANCH{\char"2C69\char"0331} \SUPPORTEDANCH{\char"004D\char"0331} \SUPPORTEDANCH{\char"2C6E\char"0331} \SUPPORTEDANCH{\char"0050\char"0331} \SUPPORTEDANCH{\char"01A4\char"0331} \SUPPORTEDPREC{\char"0052\char"0331} \SUPPORTEDANCH{\char"2C64\char"0331} \SUPPORTEDPREC{\char"0054\char"0331} \SUPPORTEDANCH{\char"01AC\char"0331} \SUPPORTEDANCH{\char"01AE\char"0331} \SUPPORTEDANCH{\char"0059\char"0331} \SUPPORTEDANCH{\char"01B3\char"0331} \SUPPORTEDANCH{\char"0061\char"0331} \SUPPORTEDANCH{\char"1E9A\char"0331} \SUPPORTEDPREC{\char"0062\char"0331} \SUPPORTEDANCH{\char"0183\char"0331} \SUPPORTEDANCH{\char"0063\char"0331} \SUPPORTEDANCH{\char"0188\char"0331} \SUPPORTEDPREC{\char"0064\char"0331} \SUPPORTEDANCH{\char"018C\char"0331} \SUPPORTEDANCH{\char"0221\char"0331} \SUPPORTEDPREC{\char"0068\char"0331} \SUPPORTEDANCH{\char"2C68\char"0331} \SUPPORTEDPREC{\char"006B\char"0331} \SUPPORTEDANCH{\char"0199\char"0331} \SUPPORTEDANCH{\char"2C6A\char"0331} \SUPPORTEDPREC{\char"006C\char"0331} \SUPPORTEDANCH{\char"0234\char"0331} \SUPPORTEDPREC{\char"006E\char"0331} \SUPPORTEDANCH{\char"0235\char"0331} \SUPPORTEDANCH{\char"0070\char"0331} \SUPPORTEDANCH{\char"01A5\char"0331} \SUPPORTEDANCH{\char"0071\char"0331} \SUPPORTEDANCH{\char"02A0\char"0331} \SUPPORTEDPREC{\char"0074\char"0331} \SUPPORTEDANCH{\char"01AB\char"0331} \SUPPORTEDANCH{\char"01AD\char"0331} \SUPPORTEDANCH{\char"0236\char"0331} \SUPPORTEDANCH{\char"0075\char"0331} \SUPPORTEDANCH{\char"AB52\char"0331} \SUPPORTEDANCH{\char"0079\char"0331} \SUPPORTEDANCH{\char"01B4\char"0331} \SUPPORTEDANCH{\char"1EFF\char"0331}

\SUPPORTEDANCH{\char"0048\char"0333} \SUPPORTEDANCH{\char"2C67\char"0333} \SUPPORTEDANCH{\char"004B\char"0333} \SUPPORTEDANCH{\char"2C69\char"0333} \SUPPORTEDANCH{\char"004D\char"0333} \SUPPORTEDANCH{\char"2C6E\char"0333} \SUPPORTEDANCH{\char"0050\char"0333} \SUPPORTEDANCH{\char"01A4\char"0333} \SUPPORTEDANCH{\char"0052\char"0333} \SUPPORTEDANCH{\char"2C64\char"0333} \SUPPORTEDANCH{\char"0054\char"0333} \SUPPORTEDANCH{\char"01AC\char"0333} \SUPPORTEDANCH{\char"01AE\char"0333} \SUPPORTEDANCH{\char"0059\char"0333} \SUPPORTEDANCH{\char"01B3\char"0333} \SUPPORTEDANCH{\char"0061\char"0333} \SUPPORTEDANCH{\char"1E9A\char"0333} \SUPPORTEDANCH{\char"0062\char"0333} \SUPPORTEDANCH{\char"0183\char"0333} \SUPPORTEDANCH{\char"0063\char"0333} \SUPPORTEDANCH{\char"0188\char"0333} \SUPPORTEDANCH{\char"0064\char"0333} \SUPPORTEDANCH{\char"018C\char"0333} \SUPPORTEDANCH{\char"0221\char"0333} \SUPPORTEDANCH{\char"0068\char"0333} \SUPPORTEDANCH{\char"2C68\char"0333} \SUPPORTEDANCH{\char"006B\char"0333} \SUPPORTEDANCH{\char"0199\char"0333} \SUPPORTEDANCH{\char"2C6A\char"0333} \SUPPORTEDANCH{\char"006C\char"0333} \SUPPORTEDANCH{\char"0234\char"0333} \SUPPORTEDANCH{\char"006E\char"0333} \SUPPORTEDANCH{\char"0235\char"0333} \SUPPORTEDANCH{\char"0070\char"0333} \SUPPORTEDANCH{\char"01A5\char"0333} \SUPPORTEDANCH{\char"0071\char"0333} \SUPPORTEDANCH{\char"02A0\char"0333} \SUPPORTEDANCH{\char"0074\char"0333} \SUPPORTEDANCH{\char"01AB\char"0333} \SUPPORTEDANCH{\char"01AD\char"0333} \SUPPORTEDANCH{\char"0236\char"0333} \SUPPORTEDANCH{\char"0075\char"0333} \SUPPORTEDANCH{\char"AB52\char"0333} \SUPPORTEDANCH{\char"0079\char"0333} \SUPPORTEDANCH{\char"01B4\char"0333} \SUPPORTEDANCH{\char"1EFF\char"0333}

\SUPPORTEDANCH{\char"0048\char"0339} \SUPPORTEDANCH{\char"2C67\char"0339} \SUPPORTEDANCH{\char"004B\char"0339} \SUPPORTEDANCH{\char"2C69\char"0339} \SUPPORTEDANCH{\char"004D\char"0339} \SUPPORTEDANCH{\char"2C6E\char"0339} \SUPPORTEDANCH{\char"0050\char"0339} \SUPPORTEDANCH{\char"01A4\char"0339} \SUPPORTEDANCH{\char"0052\char"0339} \SUPPORTEDANCH{\char"2C64\char"0339} \SUPPORTEDANCH{\char"0054\char"0339} \SUPPORTEDANCH{\char"01AC\char"0339} \SUPPORTEDANCH{\char"01AE\char"0339} \SUPPORTEDANCH{\char"0059\char"0339} \SUPPORTEDANCH{\char"01B3\char"0339} \SUPPORTEDANCH{\char"0061\char"0339} \SUPPORTEDANCH{\char"1E9A\char"0339} \SUPPORTEDANCH{\char"0062\char"0339} \SUPPORTEDANCH{\char"0183\char"0339} \SUPPORTEDANCH{\char"0063\char"0339} \SUPPORTEDANCH{\char"0188\char"0339} \SUPPORTEDANCH{\char"0064\char"0339} \SUPPORTEDANCH{\char"018C\char"0339} \SUPPORTEDANCH{\char"0221\char"0339} \SUPPORTEDANCH{\char"0068\char"0339} \SUPPORTEDANCH{\char"2C68\char"0339} \SUPPORTEDANCH{\char"006B\char"0339} \SUPPORTEDANCH{\char"0199\char"0339} \SUPPORTEDANCH{\char"2C6A\char"0339} \SUPPORTEDANCH{\char"006C\char"0339} \SUPPORTEDANCH{\char"0234\char"0339} \SUPPORTEDANCH{\char"006E\char"0339} \SUPPORTEDANCH{\char"0235\char"0339} \SUPPORTEDANCH{\char"0070\char"0339} \SUPPORTEDANCH{\char"01A5\char"0339} \SUPPORTEDANCH{\char"0071\char"0339} \SUPPORTEDANCH{\char"02A0\char"0339} \SUPPORTEDANCH{\char"0074\char"0339} \SUPPORTEDANCH{\char"01AB\char"0339} \SUPPORTEDANCH{\char"01AD\char"0339} \SUPPORTEDANCH{\char"0236\char"0339} \SUPPORTEDANCH{\char"0075\char"0339} \SUPPORTEDANCH{\char"AB52\char"0339} \SUPPORTEDANCH{\char"0079\char"0339} \SUPPORTEDANCH{\char"01B4\char"0339} \SUPPORTEDANCH{\char"1EFF\char"0339}

\SUPPORTEDANCH{\char"0048\char"033A} \SUPPORTEDANCH{\char"2C67\char"033A} \SUPPORTEDANCH{\char"004B\char"033A} \SUPPORTEDANCH{\char"2C69\char"033A} \SUPPORTEDANCH{\char"004D\char"033A} \SUPPORTEDANCH{\char"2C6E\char"033A} \SUPPORTEDANCH{\char"0050\char"033A} \SUPPORTEDANCH{\char"01A4\char"033A} \SUPPORTEDANCH{\char"0052\char"033A} \SUPPORTEDANCH{\char"2C64\char"033A} \SUPPORTEDANCH{\char"0054\char"033A} \SUPPORTEDANCH{\char"01AC\char"033A} \SUPPORTEDANCH{\char"01AE\char"033A} \SUPPORTEDANCH{\char"0059\char"033A} \SUPPORTEDANCH{\char"01B3\char"033A} \SUPPORTEDANCH{\char"0061\char"033A} \SUPPORTEDANCH{\char"1E9A\char"033A} \SUPPORTEDANCH{\char"0062\char"033A} \SUPPORTEDANCH{\char"0183\char"033A} \SUPPORTEDANCH{\char"0063\char"033A} \SUPPORTEDANCH{\char"0188\char"033A} \SUPPORTEDANCH{\char"0064\char"033A} \SUPPORTEDANCH{\char"018C\char"033A} \SUPPORTEDANCH{\char"0221\char"033A} \SUPPORTEDANCH{\char"0068\char"033A} \SUPPORTEDANCH{\char"2C68\char"033A} \SUPPORTEDANCH{\char"006B\char"033A} \SUPPORTEDANCH{\char"0199\char"033A} \SUPPORTEDANCH{\char"2C6A\char"033A} \SUPPORTEDANCH{\char"006C\char"033A} \SUPPORTEDANCH{\char"0234\char"033A} \SUPPORTEDANCH{\char"006E\char"033A} \SUPPORTEDANCH{\char"0235\char"033A} \SUPPORTEDANCH{\char"0070\char"033A} \SUPPORTEDANCH{\char"01A5\char"033A} \SUPPORTEDANCH{\char"0071\char"033A} \SUPPORTEDANCH{\char"02A0\char"033A} \SUPPORTEDANCH{\char"0074\char"033A} \SUPPORTEDANCH{\char"01AB\char"033A} \SUPPORTEDANCH{\char"01AD\char"033A} \SUPPORTEDANCH{\char"0236\char"033A} \SUPPORTEDANCH{\char"0075\char"033A} \SUPPORTEDANCH{\char"AB52\char"033A} \SUPPORTEDANCH{\char"0079\char"033A} \SUPPORTEDANCH{\char"01B4\char"033A} \SUPPORTEDANCH{\char"1EFF\char"033A}

\SUPPORTEDANCH{\char"0048\char"033B} \SUPPORTEDANCH{\char"2C67\char"033B} \SUPPORTEDANCH{\char"004B\char"033B} \SUPPORTEDANCH{\char"2C69\char"033B} \SUPPORTEDANCH{\char"004D\char"033B} \SUPPORTEDANCH{\char"2C6E\char"033B} \SUPPORTEDANCH{\char"0050\char"033B} \SUPPORTEDANCH{\char"01A4\char"033B} \SUPPORTEDANCH{\char"0052\char"033B} \SUPPORTEDANCH{\char"2C64\char"033B} \SUPPORTEDANCH{\char"0054\char"033B} \SUPPORTEDANCH{\char"01AC\char"033B} \SUPPORTEDANCH{\char"01AE\char"033B} \SUPPORTEDANCH{\char"0059\char"033B} \SUPPORTEDANCH{\char"01B3\char"033B} \SUPPORTEDANCH{\char"0061\char"033B} \SUPPORTEDANCH{\char"1E9A\char"033B} \SUPPORTEDANCH{\char"0062\char"033B} \SUPPORTEDANCH{\char"0183\char"033B} \SUPPORTEDANCH{\char"0063\char"033B} \SUPPORTEDANCH{\char"0188\char"033B} \SUPPORTEDANCH{\char"0064\char"033B} \SUPPORTEDANCH{\char"018C\char"033B} \SUPPORTEDANCH{\char"0221\char"033B} \SUPPORTEDANCH{\char"0068\char"033B} \SUPPORTEDANCH{\char"2C68\char"033B} \SUPPORTEDANCH{\char"006B\char"033B} \SUPPORTEDANCH{\char"0199\char"033B} \SUPPORTEDANCH{\char"2C6A\char"033B} \SUPPORTEDANCH{\char"006C\char"033B} \SUPPORTEDANCH{\char"0234\char"033B} \SUPPORTEDANCH{\char"006E\char"033B} \SUPPORTEDANCH{\char"0235\char"033B} \SUPPORTEDANCH{\char"0070\char"033B} \SUPPORTEDANCH{\char"01A5\char"033B} \SUPPORTEDANCH{\char"0071\char"033B} \SUPPORTEDANCH{\char"02A0\char"033B} \SUPPORTEDANCH{\char"0074\char"033B} \SUPPORTEDANCH{\char"01AB\char"033B} \SUPPORTEDANCH{\char"01AD\char"033B} \SUPPORTEDANCH{\char"0236\char"033B} \SUPPORTEDANCH{\char"0075\char"033B} \SUPPORTEDANCH{\char"AB52\char"033B} \SUPPORTEDANCH{\char"0079\char"033B} \SUPPORTEDANCH{\char"01B4\char"033B} \SUPPORTEDANCH{\char"1EFF\char"033B}

\SUPPORTEDANCH{\char"0048\char"033C} \SUPPORTEDANCH{\char"2C67\char"033C} \SUPPORTEDANCH{\char"004B\char"033C} \SUPPORTEDANCH{\char"2C69\char"033C} \SUPPORTEDANCH{\char"004D\char"033C} \SUPPORTEDANCH{\char"2C6E\char"033C} \SUPPORTEDANCH{\char"0050\char"033C} \SUPPORTEDANCH{\char"01A4\char"033C} \SUPPORTEDANCH{\char"0052\char"033C} \SUPPORTEDANCH{\char"2C64\char"033C} \SUPPORTEDANCH{\char"0054\char"033C} \SUPPORTEDANCH{\char"01AC\char"033C} \SUPPORTEDANCH{\char"01AE\char"033C} \SUPPORTEDANCH{\char"0059\char"033C} \SUPPORTEDANCH{\char"01B3\char"033C} \SUPPORTEDANCH{\char"0061\char"033C} \SUPPORTEDANCH{\char"1E9A\char"033C} \SUPPORTEDANCH{\char"0062\char"033C} \SUPPORTEDANCH{\char"0183\char"033C} \SUPPORTEDANCH{\char"0063\char"033C} \SUPPORTEDANCH{\char"0188\char"033C} \SUPPORTEDANCH{\char"0064\char"033C} \SUPPORTEDANCH{\char"018C\char"033C} \SUPPORTEDANCH{\char"0221\char"033C} \SUPPORTEDANCH{\char"0068\char"033C} \SUPPORTEDANCH{\char"2C68\char"033C} \SUPPORTEDANCH{\char"006B\char"033C} \SUPPORTEDANCH{\char"0199\char"033C} \SUPPORTEDANCH{\char"2C6A\char"033C} \SUPPORTEDANCH{\char"006C\char"033C} \SUPPORTEDANCH{\char"0234\char"033C} \SUPPORTEDANCH{\char"006E\char"033C} \SUPPORTEDANCH{\char"0235\char"033C} \SUPPORTEDANCH{\char"0070\char"033C} \SUPPORTEDANCH{\char"01A5\char"033C} \SUPPORTEDANCH{\char"0071\char"033C} \SUPPORTEDANCH{\char"02A0\char"033C} \SUPPORTEDANCH{\char"0074\char"033C} \SUPPORTEDANCH{\char"01AB\char"033C} \SUPPORTEDANCH{\char"01AD\char"033C} \SUPPORTEDANCH{\char"0236\char"033C} \SUPPORTEDANCH{\char"0075\char"033C} \SUPPORTEDANCH{\char"AB52\char"033C} \SUPPORTEDANCH{\char"0079\char"033C} \SUPPORTEDANCH{\char"01B4\char"033C} \SUPPORTEDANCH{\char"1EFF\char"033C}

\SUPPORTEDANCH{\char"0048\char"0347} \SUPPORTEDANCH{\char"2C67\char"0347} \SUPPORTEDANCH{\char"004B\char"0347} \SUPPORTEDANCH{\char"2C69\char"0347} \SUPPORTEDANCH{\char"004D\char"0347} \SUPPORTEDANCH{\char"2C6E\char"0347} \SUPPORTEDANCH{\char"0050\char"0347} \SUPPORTEDANCH{\char"01A4\char"0347} \SUPPORTEDANCH{\char"0052\char"0347} \SUPPORTEDANCH{\char"2C64\char"0347} \SUPPORTEDANCH{\char"0054\char"0347} \SUPPORTEDANCH{\char"01AC\char"0347} \SUPPORTEDANCH{\char"01AE\char"0347} \SUPPORTEDANCH{\char"0059\char"0347} \SUPPORTEDANCH{\char"01B3\char"0347} \SUPPORTEDANCH{\char"0061\char"0347} \SUPPORTEDANCH{\char"1E9A\char"0347} \SUPPORTEDANCH{\char"0062\char"0347} \SUPPORTEDANCH{\char"0183\char"0347} \SUPPORTEDANCH{\char"0063\char"0347} \SUPPORTEDANCH{\char"0188\char"0347} \SUPPORTEDANCH{\char"0064\char"0347} \SUPPORTEDANCH{\char"018C\char"0347} \SUPPORTEDANCH{\char"0221\char"0347} \SUPPORTEDANCH{\char"0068\char"0347} \SUPPORTEDANCH{\char"2C68\char"0347} \SUPPORTEDANCH{\char"006B\char"0347} \SUPPORTEDANCH{\char"0199\char"0347} \SUPPORTEDANCH{\char"2C6A\char"0347} \SUPPORTEDANCH{\char"006C\char"0347} \SUPPORTEDANCH{\char"0234\char"0347} \SUPPORTEDANCH{\char"006E\char"0347} \SUPPORTEDANCH{\char"0235\char"0347} \SUPPORTEDANCH{\char"0070\char"0347} \SUPPORTEDANCH{\char"01A5\char"0347} \SUPPORTEDANCH{\char"0071\char"0347} \SUPPORTEDANCH{\char"02A0\char"0347} \SUPPORTEDANCH{\char"0074\char"0347} \SUPPORTEDANCH{\char"01AB\char"0347} \SUPPORTEDANCH{\char"01AD\char"0347} \SUPPORTEDANCH{\char"0236\char"0347} \SUPPORTEDANCH{\char"0075\char"0347} \SUPPORTEDANCH{\char"AB52\char"0347} \SUPPORTEDANCH{\char"0079\char"0347} \SUPPORTEDANCH{\char"01B4\char"0347} \SUPPORTEDANCH{\char"1EFF\char"0347}

\SUPPORTEDANCH{\char"0048\char"0348} \SUPPORTEDANCH{\char"2C67\char"0348} \SUPPORTEDANCH{\char"004B\char"0348} \SUPPORTEDANCH{\char"2C69\char"0348} \SUPPORTEDANCH{\char"004D\char"0348} \SUPPORTEDANCH{\char"2C6E\char"0348} \SUPPORTEDANCH{\char"0050\char"0348} \SUPPORTEDANCH{\char"01A4\char"0348} \SUPPORTEDANCH{\char"0052\char"0348} \SUPPORTEDANCH{\char"2C64\char"0348} \SUPPORTEDANCH{\char"0054\char"0348} \SUPPORTEDANCH{\char"01AC\char"0348} \SUPPORTEDANCH{\char"01AE\char"0348} \SUPPORTEDANCH{\char"0059\char"0348} \SUPPORTEDANCH{\char"01B3\char"0348} \SUPPORTEDANCH{\char"0061\char"0348} \SUPPORTEDANCH{\char"1E9A\char"0348} \SUPPORTEDANCH{\char"0062\char"0348} \SUPPORTEDANCH{\char"0183\char"0348} \SUPPORTEDANCH{\char"0063\char"0348} \SUPPORTEDANCH{\char"0188\char"0348} \SUPPORTEDANCH{\char"0064\char"0348} \SUPPORTEDANCH{\char"018C\char"0348} \SUPPORTEDANCH{\char"0221\char"0348} \SUPPORTEDANCH{\char"0068\char"0348} \SUPPORTEDANCH{\char"2C68\char"0348} \SUPPORTEDANCH{\char"006B\char"0348} \SUPPORTEDANCH{\char"0199\char"0348} \SUPPORTEDANCH{\char"2C6A\char"0348} \SUPPORTEDANCH{\char"006C\char"0348} \SUPPORTEDANCH{\char"0234\char"0348} \SUPPORTEDANCH{\char"006E\char"0348} \SUPPORTEDANCH{\char"0235\char"0348} \SUPPORTEDANCH{\char"0070\char"0348} \SUPPORTEDANCH{\char"01A5\char"0348} \SUPPORTEDANCH{\char"0071\char"0348} \SUPPORTEDANCH{\char"02A0\char"0348} \SUPPORTEDANCH{\char"0074\char"0348} \SUPPORTEDANCH{\char"01AB\char"0348} \SUPPORTEDANCH{\char"01AD\char"0348} \SUPPORTEDANCH{\char"0236\char"0348} \SUPPORTEDANCH{\char"0075\char"0348} \SUPPORTEDANCH{\char"AB52\char"0348} \SUPPORTEDANCH{\char"0079\char"0348} \SUPPORTEDANCH{\char"01B4\char"0348} \SUPPORTEDANCH{\char"1EFF\char"0348}

\SUPPORTEDANCH{\char"0048\char"0349} \SUPPORTEDANCH{\char"2C67\char"0349} \SUPPORTEDANCH{\char"004B\char"0349} \SUPPORTEDANCH{\char"2C69\char"0349} \SUPPORTEDANCH{\char"004D\char"0349} \SUPPORTEDANCH{\char"2C6E\char"0349} \SUPPORTEDANCH{\char"0050\char"0349} \SUPPORTEDANCH{\char"01A4\char"0349} \SUPPORTEDANCH{\char"0052\char"0349} \SUPPORTEDANCH{\char"2C64\char"0349} \SUPPORTEDANCH{\char"0054\char"0349} \SUPPORTEDANCH{\char"01AC\char"0349} \SUPPORTEDANCH{\char"01AE\char"0349} \SUPPORTEDANCH{\char"0059\char"0349} \SUPPORTEDANCH{\char"01B3\char"0349} \SUPPORTEDANCH{\char"0061\char"0349} \SUPPORTEDANCH{\char"1E9A\char"0349} \SUPPORTEDANCH{\char"0062\char"0349} \SUPPORTEDANCH{\char"0183\char"0349} \SUPPORTEDANCH{\char"0063\char"0349} \SUPPORTEDANCH{\char"0188\char"0349} \SUPPORTEDANCH{\char"0064\char"0349} \SUPPORTEDANCH{\char"018C\char"0349} \SUPPORTEDANCH{\char"0221\char"0349} \SUPPORTEDANCH{\char"0068\char"0349} \SUPPORTEDANCH{\char"2C68\char"0349} \SUPPORTEDANCH{\char"006B\char"0349} \SUPPORTEDANCH{\char"0199\char"0349} \SUPPORTEDANCH{\char"2C6A\char"0349} \SUPPORTEDANCH{\char"006C\char"0349} \SUPPORTEDANCH{\char"0234\char"0349} \SUPPORTEDANCH{\char"006E\char"0349} \SUPPORTEDANCH{\char"0235\char"0349} \SUPPORTEDANCH{\char"0070\char"0349} \SUPPORTEDANCH{\char"01A5\char"0349} \SUPPORTEDANCH{\char"0071\char"0349} \SUPPORTEDANCH{\char"02A0\char"0349} \SUPPORTEDANCH{\char"0074\char"0349} \SUPPORTEDANCH{\char"01AB\char"0349} \SUPPORTEDANCH{\char"01AD\char"0349} \SUPPORTEDANCH{\char"0236\char"0349} \SUPPORTEDANCH{\char"0075\char"0349} \SUPPORTEDANCH{\char"AB52\char"0349} \SUPPORTEDANCH{\char"0079\char"0349} \SUPPORTEDANCH{\char"01B4\char"0349} \SUPPORTEDANCH{\char"1EFF\char"0349}

\SUPPORTEDANCH{\char"0048\char"034D} \SUPPORTEDANCH{\char"2C67\char"034D} \SUPPORTEDANCH{\char"004B\char"034D} \SUPPORTEDANCH{\char"2C69\char"034D} \SUPPORTEDANCH{\char"004D\char"034D} \SUPPORTEDANCH{\char"2C6E\char"034D} \SUPPORTEDANCH{\char"0050\char"034D} \SUPPORTEDANCH{\char"01A4\char"034D} \SUPPORTEDANCH{\char"0052\char"034D} \SUPPORTEDANCH{\char"2C64\char"034D} \SUPPORTEDANCH{\char"0054\char"034D} \SUPPORTEDANCH{\char"01AC\char"034D} \SUPPORTEDANCH{\char"01AE\char"034D} \SUPPORTEDANCH{\char"0059\char"034D} \SUPPORTEDANCH{\char"01B3\char"034D} \SUPPORTEDANCH{\char"0061\char"034D} \SUPPORTEDANCH{\char"1E9A\char"034D} \SUPPORTEDANCH{\char"0062\char"034D} \SUPPORTEDANCH{\char"0183\char"034D} \SUPPORTEDANCH{\char"0063\char"034D} \SUPPORTEDANCH{\char"0188\char"034D} \SUPPORTEDANCH{\char"0064\char"034D} \SUPPORTEDANCH{\char"018C\char"034D} \SUPPORTEDANCH{\char"0221\char"034D} \SUPPORTEDANCH{\char"0068\char"034D} \SUPPORTEDANCH{\char"2C68\char"034D} \SUPPORTEDANCH{\char"006B\char"034D} \SUPPORTEDANCH{\char"0199\char"034D} \SUPPORTEDANCH{\char"2C6A\char"034D} \SUPPORTEDANCH{\char"006C\char"034D} \SUPPORTEDANCH{\char"0234\char"034D} \SUPPORTEDANCH{\char"006E\char"034D} \SUPPORTEDANCH{\char"0235\char"034D} \SUPPORTEDANCH{\char"0070\char"034D} \SUPPORTEDANCH{\char"01A5\char"034D} \SUPPORTEDANCH{\char"0071\char"034D} \SUPPORTEDANCH{\char"02A0\char"034D} \SUPPORTEDANCH{\char"0074\char"034D} \SUPPORTEDANCH{\char"01AB\char"034D} \SUPPORTEDANCH{\char"01AD\char"034D} \SUPPORTEDANCH{\char"0236\char"034D} \SUPPORTEDANCH{\char"0075\char"034D} \SUPPORTEDANCH{\char"AB52\char"034D} \SUPPORTEDANCH{\char"0079\char"034D} \SUPPORTEDANCH{\char"01B4\char"034D} \SUPPORTEDANCH{\char"1EFF\char"034D}

\SUPPORTEDANCH{\char"0048\char"034E} \SUPPORTEDANCH{\char"2C67\char"034E} \SUPPORTEDANCH{\char"004B\char"034E} \SUPPORTEDANCH{\char"2C69\char"034E} \SUPPORTEDANCH{\char"004D\char"034E} \SUPPORTEDANCH{\char"2C6E\char"034E} \SUPPORTEDANCH{\char"0050\char"034E} \SUPPORTEDANCH{\char"01A4\char"034E} \SUPPORTEDANCH{\char"0052\char"034E} \SUPPORTEDANCH{\char"2C64\char"034E} \SUPPORTEDANCH{\char"0054\char"034E} \SUPPORTEDANCH{\char"01AC\char"034E} \SUPPORTEDANCH{\char"01AE\char"034E} \SUPPORTEDANCH{\char"0059\char"034E} \SUPPORTEDANCH{\char"01B3\char"034E} \SUPPORTEDANCH{\char"0061\char"034E} \SUPPORTEDANCH{\char"1E9A\char"034E} \SUPPORTEDANCH{\char"0062\char"034E} \SUPPORTEDANCH{\char"0183\char"034E} \SUPPORTEDANCH{\char"0063\char"034E} \SUPPORTEDANCH{\char"0188\char"034E} \SUPPORTEDANCH{\char"0064\char"034E} \SUPPORTEDANCH{\char"018C\char"034E} \SUPPORTEDANCH{\char"0221\char"034E} \SUPPORTEDANCH{\char"0068\char"034E} \SUPPORTEDANCH{\char"2C68\char"034E} \SUPPORTEDANCH{\char"006B\char"034E} \SUPPORTEDANCH{\char"0199\char"034E} \SUPPORTEDANCH{\char"2C6A\char"034E} \SUPPORTEDANCH{\char"006C\char"034E} \SUPPORTEDANCH{\char"0234\char"034E} \SUPPORTEDANCH{\char"006E\char"034E} \SUPPORTEDANCH{\char"0235\char"034E} \SUPPORTEDANCH{\char"0070\char"034E} \SUPPORTEDANCH{\char"01A5\char"034E} \SUPPORTEDANCH{\char"0071\char"034E} \SUPPORTEDANCH{\char"02A0\char"034E} \SUPPORTEDANCH{\char"0074\char"034E} \SUPPORTEDANCH{\char"01AB\char"034E} \SUPPORTEDANCH{\char"01AD\char"034E} \SUPPORTEDANCH{\char"0236\char"034E} \SUPPORTEDANCH{\char"0075\char"034E} \SUPPORTEDANCH{\char"AB52\char"034E} \SUPPORTEDANCH{\char"0079\char"034E} \SUPPORTEDANCH{\char"01B4\char"034E} \SUPPORTEDANCH{\char"1EFF\char"034E}

}

\newpage
\subsection*{Italic, Group 4, below, Below}

{\itshape
% grid. columns are bases, rows are marks.
\SUPPORTEDANCH{\char"0048\char"0316} \SUPPORTEDFALL{\char"2C67\char"0316} \SUPPORTEDANCH{\char"004B\char"0316} \SUPPORTEDANCH{\char"2C69\char"0316} \SUPPORTEDANCH{\char"004D\char"0316} \MISSINGGLYPH{\char"2C6E\char"0316} \SUPPORTEDANCH{\char"0050\char"0316} \SUPPORTEDFALL{\char"01A4\char"0316} \SUPPORTEDANCH{\char"0052\char"0316} \SUPPORTEDFALL{\char"2C64\char"0316} \SUPPORTEDANCH{\char"0054\char"0316} \SUPPORTEDFALL{\char"01AC\char"0316} \SUPPORTEDFALL{\char"01AE\char"0316} \SUPPORTEDANCH{\char"0059\char"0316} \SUPPORTEDFALL{\char"01B3\char"0316} \SUPPORTEDANCH{\char"0061\char"0316} \SUPPORTEDFALL{\char"1E9A\char"0316} \SUPPORTEDANCH{\char"0062\char"0316} \SUPPORTEDANCH{\char"0183\char"0316} \SUPPORTEDANCH{\char"0063\char"0316} \SUPPORTEDFALL{\char"0188\char"0316} \SUPPORTEDANCH{\char"0064\char"0316} \SUPPORTEDANCH{\char"018C\char"0316} \SUPPORTEDFALL{\char"0221\char"0316} \SUPPORTEDANCH{\char"0068\char"0316} \SUPPORTEDANCH{\char"2C68\char"0316} \SUPPORTEDANCH{\char"006B\char"0316} \SUPPORTEDANCH{\char"0199\char"0316} \SUPPORTEDANCH{\char"2C6A\char"0316} \SUPPORTEDANCH{\char"006C\char"0316} \SUPPORTEDFALL{\char"0234\char"0316} \SUPPORTEDANCH{\char"006E\char"0316} \SUPPORTEDANCH{\char"0235\char"0316} \SUPPORTEDANCH{\char"0070\char"0316} \SUPPORTEDFALL{\char"01A5\char"0316} \SUPPORTEDANCH{\char"0071\char"0316} \SUPPORTEDFALL{\char"02A0\char"0316} \SUPPORTEDANCH{\char"0074\char"0316} \SUPPORTEDFALL{\char"01AB\char"0316} \SUPPORTEDANCH{\char"01AD\char"0316} \SUPPORTEDANCH{\char"0236\char"0316} \SUPPORTEDANCH{\char"0075\char"0316} \MISSINGGLYPH{\char"AB52\char"0316} \SUPPORTEDANCH{\char"0079\char"0316} \SUPPORTEDFALL{\char"01B4\char"0316} \MISSINGGLYPH{\char"1EFF\char"0316}

\SUPPORTEDANCH{\char"0048\char"0317} \SUPPORTEDFALL{\char"2C67\char"0317} \SUPPORTEDANCH{\char"004B\char"0317} \SUPPORTEDANCH{\char"2C69\char"0317} \SUPPORTEDANCH{\char"004D\char"0317} \MISSINGGLYPH{\char"2C6E\char"0317} \SUPPORTEDANCH{\char"0050\char"0317} \SUPPORTEDFALL{\char"01A4\char"0317} \SUPPORTEDANCH{\char"0052\char"0317} \SUPPORTEDFALL{\char"2C64\char"0317} \SUPPORTEDANCH{\char"0054\char"0317} \SUPPORTEDFALL{\char"01AC\char"0317} \SUPPORTEDFALL{\char"01AE\char"0317} \SUPPORTEDANCH{\char"0059\char"0317} \SUPPORTEDFALL{\char"01B3\char"0317} \SUPPORTEDANCH{\char"0061\char"0317} \SUPPORTEDFALL{\char"1E9A\char"0317} \SUPPORTEDANCH{\char"0062\char"0317} \SUPPORTEDANCH{\char"0183\char"0317} \SUPPORTEDANCH{\char"0063\char"0317} \SUPPORTEDFALL{\char"0188\char"0317} \SUPPORTEDANCH{\char"0064\char"0317} \SUPPORTEDANCH{\char"018C\char"0317} \SUPPORTEDFALL{\char"0221\char"0317} \SUPPORTEDANCH{\char"0068\char"0317} \SUPPORTEDANCH{\char"2C68\char"0317} \SUPPORTEDANCH{\char"006B\char"0317} \SUPPORTEDANCH{\char"0199\char"0317} \SUPPORTEDANCH{\char"2C6A\char"0317} \SUPPORTEDANCH{\char"006C\char"0317} \SUPPORTEDFALL{\char"0234\char"0317} \SUPPORTEDANCH{\char"006E\char"0317} \SUPPORTEDANCH{\char"0235\char"0317} \SUPPORTEDANCH{\char"0070\char"0317} \SUPPORTEDFALL{\char"01A5\char"0317} \SUPPORTEDANCH{\char"0071\char"0317} \SUPPORTEDFALL{\char"02A0\char"0317} \SUPPORTEDANCH{\char"0074\char"0317} \SUPPORTEDFALL{\char"01AB\char"0317} \SUPPORTEDANCH{\char"01AD\char"0317} \SUPPORTEDANCH{\char"0236\char"0317} \SUPPORTEDANCH{\char"0075\char"0317} \MISSINGGLYPH{\char"AB52\char"0317} \SUPPORTEDANCH{\char"0079\char"0317} \SUPPORTEDFALL{\char"01B4\char"0317} \MISSINGGLYPH{\char"1EFF\char"0317}

\SUPPORTEDANCH{\char"0048\char"0318} \SUPPORTEDFALL{\char"2C67\char"0318} \SUPPORTEDANCH{\char"004B\char"0318} \SUPPORTEDANCH{\char"2C69\char"0318} \SUPPORTEDANCH{\char"004D\char"0318} \MISSINGGLYPH{\char"2C6E\char"0318} \SUPPORTEDANCH{\char"0050\char"0318} \SUPPORTEDFALL{\char"01A4\char"0318} \SUPPORTEDANCH{\char"0052\char"0318} \SUPPORTEDFALL{\char"2C64\char"0318} \SUPPORTEDANCH{\char"0054\char"0318} \SUPPORTEDFALL{\char"01AC\char"0318} \SUPPORTEDFALL{\char"01AE\char"0318} \SUPPORTEDANCH{\char"0059\char"0318} \SUPPORTEDFALL{\char"01B3\char"0318} \SUPPORTEDANCH{\char"0061\char"0318} \SUPPORTEDFALL{\char"1E9A\char"0318} \SUPPORTEDANCH{\char"0062\char"0318} \SUPPORTEDANCH{\char"0183\char"0318} \SUPPORTEDANCH{\char"0063\char"0318} \SUPPORTEDFALL{\char"0188\char"0318} \SUPPORTEDANCH{\char"0064\char"0318} \SUPPORTEDANCH{\char"018C\char"0318} \SUPPORTEDFALL{\char"0221\char"0318} \SUPPORTEDANCH{\char"0068\char"0318} \SUPPORTEDANCH{\char"2C68\char"0318} \SUPPORTEDANCH{\char"006B\char"0318} \SUPPORTEDANCH{\char"0199\char"0318} \SUPPORTEDANCH{\char"2C6A\char"0318} \SUPPORTEDANCH{\char"006C\char"0318} \SUPPORTEDFALL{\char"0234\char"0318} \SUPPORTEDANCH{\char"006E\char"0318} \SUPPORTEDANCH{\char"0235\char"0318} \SUPPORTEDANCH{\char"0070\char"0318} \SUPPORTEDFALL{\char"01A5\char"0318} \SUPPORTEDANCH{\char"0071\char"0318} \SUPPORTEDFALL{\char"02A0\char"0318} \SUPPORTEDANCH{\char"0074\char"0318} \SUPPORTEDFALL{\char"01AB\char"0318} \SUPPORTEDANCH{\char"01AD\char"0318} \SUPPORTEDANCH{\char"0236\char"0318} \SUPPORTEDANCH{\char"0075\char"0318} \MISSINGGLYPH{\char"AB52\char"0318} \SUPPORTEDANCH{\char"0079\char"0318} \SUPPORTEDFALL{\char"01B4\char"0318} \MISSINGGLYPH{\char"1EFF\char"0318}

\SUPPORTEDANCH{\char"0048\char"0319} \SUPPORTEDFALL{\char"2C67\char"0319} \SUPPORTEDANCH{\char"004B\char"0319} \SUPPORTEDANCH{\char"2C69\char"0319} \SUPPORTEDANCH{\char"004D\char"0319} \MISSINGGLYPH{\char"2C6E\char"0319} \SUPPORTEDANCH{\char"0050\char"0319} \SUPPORTEDFALL{\char"01A4\char"0319} \SUPPORTEDANCH{\char"0052\char"0319} \SUPPORTEDFALL{\char"2C64\char"0319} \SUPPORTEDANCH{\char"0054\char"0319} \SUPPORTEDFALL{\char"01AC\char"0319} \SUPPORTEDFALL{\char"01AE\char"0319} \SUPPORTEDANCH{\char"0059\char"0319} \SUPPORTEDFALL{\char"01B3\char"0319} \SUPPORTEDANCH{\char"0061\char"0319} \SUPPORTEDFALL{\char"1E9A\char"0319} \SUPPORTEDANCH{\char"0062\char"0319} \SUPPORTEDANCH{\char"0183\char"0319} \SUPPORTEDANCH{\char"0063\char"0319} \SUPPORTEDFALL{\char"0188\char"0319} \SUPPORTEDANCH{\char"0064\char"0319} \SUPPORTEDANCH{\char"018C\char"0319} \SUPPORTEDFALL{\char"0221\char"0319} \SUPPORTEDANCH{\char"0068\char"0319} \SUPPORTEDANCH{\char"2C68\char"0319} \SUPPORTEDANCH{\char"006B\char"0319} \SUPPORTEDANCH{\char"0199\char"0319} \SUPPORTEDANCH{\char"2C6A\char"0319} \SUPPORTEDANCH{\char"006C\char"0319} \SUPPORTEDFALL{\char"0234\char"0319} \SUPPORTEDANCH{\char"006E\char"0319} \SUPPORTEDANCH{\char"0235\char"0319} \SUPPORTEDANCH{\char"0070\char"0319} \SUPPORTEDFALL{\char"01A5\char"0319} \SUPPORTEDANCH{\char"0071\char"0319} \SUPPORTEDFALL{\char"02A0\char"0319} \SUPPORTEDANCH{\char"0074\char"0319} \SUPPORTEDFALL{\char"01AB\char"0319} \SUPPORTEDANCH{\char"01AD\char"0319} \SUPPORTEDANCH{\char"0236\char"0319} \SUPPORTEDANCH{\char"0075\char"0319} \MISSINGGLYPH{\char"AB52\char"0319} \SUPPORTEDANCH{\char"0079\char"0319} \SUPPORTEDFALL{\char"01B4\char"0319} \MISSINGGLYPH{\char"1EFF\char"0319}

\SUPPORTEDANCH{\char"0048\char"031C} \SUPPORTEDFALL{\char"2C67\char"031C} \SUPPORTEDANCH{\char"004B\char"031C} \SUPPORTEDANCH{\char"2C69\char"031C} \SUPPORTEDANCH{\char"004D\char"031C} \MISSINGGLYPH{\char"2C6E\char"031C} \SUPPORTEDANCH{\char"0050\char"031C} \SUPPORTEDFALL{\char"01A4\char"031C} \SUPPORTEDANCH{\char"0052\char"031C} \SUPPORTEDFALL{\char"2C64\char"031C} \SUPPORTEDANCH{\char"0054\char"031C} \SUPPORTEDFALL{\char"01AC\char"031C} \SUPPORTEDFALL{\char"01AE\char"031C} \SUPPORTEDANCH{\char"0059\char"031C} \SUPPORTEDFALL{\char"01B3\char"031C} \SUPPORTEDANCH{\char"0061\char"031C} \SUPPORTEDFALL{\char"1E9A\char"031C} \SUPPORTEDANCH{\char"0062\char"031C} \SUPPORTEDANCH{\char"0183\char"031C} \SUPPORTEDANCH{\char"0063\char"031C} \SUPPORTEDFALL{\char"0188\char"031C} \SUPPORTEDANCH{\char"0064\char"031C} \SUPPORTEDANCH{\char"018C\char"031C} \SUPPORTEDFALL{\char"0221\char"031C} \SUPPORTEDANCH{\char"0068\char"031C} \SUPPORTEDANCH{\char"2C68\char"031C} \SUPPORTEDANCH{\char"006B\char"031C} \SUPPORTEDANCH{\char"0199\char"031C} \SUPPORTEDANCH{\char"2C6A\char"031C} \SUPPORTEDANCH{\char"006C\char"031C} \SUPPORTEDFALL{\char"0234\char"031C} \SUPPORTEDANCH{\char"006E\char"031C} \SUPPORTEDANCH{\char"0235\char"031C} \SUPPORTEDANCH{\char"0070\char"031C} \SUPPORTEDFALL{\char"01A5\char"031C} \SUPPORTEDANCH{\char"0071\char"031C} \SUPPORTEDFALL{\char"02A0\char"031C} \SUPPORTEDANCH{\char"0074\char"031C} \SUPPORTEDFALL{\char"01AB\char"031C} \SUPPORTEDANCH{\char"01AD\char"031C} \SUPPORTEDANCH{\char"0236\char"031C} \SUPPORTEDANCH{\char"0075\char"031C} \MISSINGGLYPH{\char"AB52\char"031C} \SUPPORTEDANCH{\char"0079\char"031C} \SUPPORTEDFALL{\char"01B4\char"031C} \MISSINGGLYPH{\char"1EFF\char"031C}

\SUPPORTEDANCH{\char"0048\char"0320} \SUPPORTEDFALL{\char"2C67\char"0320} \SUPPORTEDANCH{\char"004B\char"0320} \SUPPORTEDANCH{\char"2C69\char"0320} \SUPPORTEDANCH{\char"004D\char"0320} \MISSINGGLYPH{\char"2C6E\char"0320} \SUPPORTEDANCH{\char"0050\char"0320} \SUPPORTEDFALL{\char"01A4\char"0320} \SUPPORTEDANCH{\char"0052\char"0320} \SUPPORTEDFALL{\char"2C64\char"0320} \SUPPORTEDANCH{\char"0054\char"0320} \SUPPORTEDFALL{\char"01AC\char"0320} \SUPPORTEDFALL{\char"01AE\char"0320} \SUPPORTEDANCH{\char"0059\char"0320} \SUPPORTEDFALL{\char"01B3\char"0320} \SUPPORTEDANCH{\char"0061\char"0320} \SUPPORTEDFALL{\char"1E9A\char"0320} \SUPPORTEDANCH{\char"0062\char"0320} \SUPPORTEDANCH{\char"0183\char"0320} \SUPPORTEDANCH{\char"0063\char"0320} \SUPPORTEDFALL{\char"0188\char"0320} \SUPPORTEDANCH{\char"0064\char"0320} \SUPPORTEDANCH{\char"018C\char"0320} \SUPPORTEDFALL{\char"0221\char"0320} \SUPPORTEDANCH{\char"0068\char"0320} \SUPPORTEDANCH{\char"2C68\char"0320} \SUPPORTEDANCH{\char"006B\char"0320} \SUPPORTEDANCH{\char"0199\char"0320} \SUPPORTEDANCH{\char"2C6A\char"0320} \SUPPORTEDANCH{\char"006C\char"0320} \SUPPORTEDFALL{\char"0234\char"0320} \SUPPORTEDANCH{\char"006E\char"0320} \SUPPORTEDANCH{\char"0235\char"0320} \SUPPORTEDANCH{\char"0070\char"0320} \SUPPORTEDFALL{\char"01A5\char"0320} \SUPPORTEDANCH{\char"0071\char"0320} \SUPPORTEDFALL{\char"02A0\char"0320} \SUPPORTEDANCH{\char"0074\char"0320} \SUPPORTEDFALL{\char"01AB\char"0320} \SUPPORTEDANCH{\char"01AD\char"0320} \SUPPORTEDANCH{\char"0236\char"0320} \SUPPORTEDANCH{\char"0075\char"0320} \MISSINGGLYPH{\char"AB52\char"0320} \SUPPORTEDANCH{\char"0079\char"0320} \SUPPORTEDFALL{\char"01B4\char"0320} \MISSINGGLYPH{\char"1EFF\char"0320}

\SUPPORTEDPREC{\char"0048\char"0323} \SUPPORTEDFALL{\char"2C67\char"0323} \SUPPORTEDPREC{\char"004B\char"0323} \SUPPORTEDANCH{\char"2C69\char"0323} \SUPPORTEDPREC{\char"004D\char"0323} \MISSINGGLYPH{\char"2C6E\char"0323} \SUPPORTEDANCH{\char"0050\char"0323} \SUPPORTEDFALL{\char"01A4\char"0323} \SUPPORTEDPREC{\char"0052\char"0323} \SUPPORTEDFALL{\char"2C64\char"0323} \SUPPORTEDPREC{\char"0054\char"0323} \SUPPORTEDFALL{\char"01AC\char"0323} \SUPPORTEDFALL{\char"01AE\char"0323} \SUPPORTEDPREC{\char"0059\char"0323} \SUPPORTEDFALL{\char"01B3\char"0323} \SUPPORTEDPREC{\char"0061\char"0323} \SUPPORTEDFALL{\char"1E9A\char"0323} \SUPPORTEDPREC{\char"0062\char"0323} \SUPPORTEDANCH{\char"0183\char"0323} \SUPPORTEDANCH{\char"0063\char"0323} \SUPPORTEDFALL{\char"0188\char"0323} \SUPPORTEDPREC{\char"0064\char"0323} \SUPPORTEDANCH{\char"018C\char"0323} \SUPPORTEDFALL{\char"0221\char"0323} \SUPPORTEDPREC{\char"0068\char"0323} \SUPPORTEDANCH{\char"2C68\char"0323} \SUPPORTEDPREC{\char"006B\char"0323} \SUPPORTEDANCH{\char"0199\char"0323} \SUPPORTEDANCH{\char"2C6A\char"0323} \SUPPORTEDPREC{\char"006C\char"0323} \SUPPORTEDFALL{\char"0234\char"0323} \SUPPORTEDPREC{\char"006E\char"0323} \SUPPORTEDANCH{\char"0235\char"0323} \SUPPORTEDANCH{\char"0070\char"0323} \SUPPORTEDFALL{\char"01A5\char"0323} \SUPPORTEDANCH{\char"0071\char"0323} \SUPPORTEDFALL{\char"02A0\char"0323} \SUPPORTEDPREC{\char"0074\char"0323} \SUPPORTEDFALL{\char"01AB\char"0323} \SUPPORTEDANCH{\char"01AD\char"0323} \SUPPORTEDANCH{\char"0236\char"0323} \SUPPORTEDPREC{\char"0075\char"0323} \MISSINGGLYPH{\char"AB52\char"0323} \SUPPORTEDPREC{\char"0079\char"0323} \SUPPORTEDFALL{\char"01B4\char"0323} \MISSINGGLYPH{\char"1EFF\char"0323}

\SUPPORTEDANCH{\char"0048\char"0324} \SUPPORTEDFALL{\char"2C67\char"0324} \SUPPORTEDANCH{\char"004B\char"0324} \SUPPORTEDANCH{\char"2C69\char"0324} \SUPPORTEDANCH{\char"004D\char"0324} \MISSINGGLYPH{\char"2C6E\char"0324} \SUPPORTEDANCH{\char"0050\char"0324} \SUPPORTEDFALL{\char"01A4\char"0324} \SUPPORTEDANCH{\char"0052\char"0324} \SUPPORTEDFALL{\char"2C64\char"0324} \SUPPORTEDANCH{\char"0054\char"0324} \SUPPORTEDFALL{\char"01AC\char"0324} \SUPPORTEDFALL{\char"01AE\char"0324} \SUPPORTEDANCH{\char"0059\char"0324} \SUPPORTEDFALL{\char"01B3\char"0324} \SUPPORTEDANCH{\char"0061\char"0324} \SUPPORTEDFALL{\char"1E9A\char"0324} \SUPPORTEDANCH{\char"0062\char"0324} \SUPPORTEDANCH{\char"0183\char"0324} \SUPPORTEDANCH{\char"0063\char"0324} \SUPPORTEDFALL{\char"0188\char"0324} \SUPPORTEDANCH{\char"0064\char"0324} \SUPPORTEDANCH{\char"018C\char"0324} \SUPPORTEDFALL{\char"0221\char"0324} \SUPPORTEDANCH{\char"0068\char"0324} \SUPPORTEDANCH{\char"2C68\char"0324} \SUPPORTEDANCH{\char"006B\char"0324} \SUPPORTEDANCH{\char"0199\char"0324} \SUPPORTEDANCH{\char"2C6A\char"0324} \SUPPORTEDANCH{\char"006C\char"0324} \SUPPORTEDFALL{\char"0234\char"0324} \SUPPORTEDANCH{\char"006E\char"0324} \SUPPORTEDANCH{\char"0235\char"0324} \SUPPORTEDANCH{\char"0070\char"0324} \SUPPORTEDFALL{\char"01A5\char"0324} \SUPPORTEDANCH{\char"0071\char"0324} \SUPPORTEDFALL{\char"02A0\char"0324} \SUPPORTEDANCH{\char"0074\char"0324} \SUPPORTEDFALL{\char"01AB\char"0324} \SUPPORTEDANCH{\char"01AD\char"0324} \SUPPORTEDANCH{\char"0236\char"0324} \SUPPORTEDPREC{\char"0075\char"0324} \MISSINGGLYPH{\char"AB52\char"0324} \SUPPORTEDANCH{\char"0079\char"0324} \SUPPORTEDFALL{\char"01B4\char"0324} \MISSINGGLYPH{\char"1EFF\char"0324}

\SUPPORTEDANCH{\char"0048\char"0325} \SUPPORTEDFALL{\char"2C67\char"0325} \SUPPORTEDANCH{\char"004B\char"0325} \SUPPORTEDANCH{\char"2C69\char"0325} \SUPPORTEDANCH{\char"004D\char"0325} \MISSINGGLYPH{\char"2C6E\char"0325} \SUPPORTEDANCH{\char"0050\char"0325} \SUPPORTEDFALL{\char"01A4\char"0325} \SUPPORTEDANCH{\char"0052\char"0325} \SUPPORTEDFALL{\char"2C64\char"0325} \SUPPORTEDANCH{\char"0054\char"0325} \SUPPORTEDFALL{\char"01AC\char"0325} \SUPPORTEDFALL{\char"01AE\char"0325} \SUPPORTEDANCH{\char"0059\char"0325} \SUPPORTEDFALL{\char"01B3\char"0325} \SUPPORTEDPREC{\char"0061\char"0325} \SUPPORTEDFALL{\char"1E9A\char"0325} \SUPPORTEDANCH{\char"0062\char"0325} \SUPPORTEDANCH{\char"0183\char"0325} \SUPPORTEDANCH{\char"0063\char"0325} \SUPPORTEDFALL{\char"0188\char"0325} \SUPPORTEDANCH{\char"0064\char"0325} \SUPPORTEDANCH{\char"018C\char"0325} \SUPPORTEDFALL{\char"0221\char"0325} \SUPPORTEDANCH{\char"0068\char"0325} \SUPPORTEDANCH{\char"2C68\char"0325} \SUPPORTEDANCH{\char"006B\char"0325} \SUPPORTEDANCH{\char"0199\char"0325} \SUPPORTEDANCH{\char"2C6A\char"0325} \SUPPORTEDANCH{\char"006C\char"0325} \SUPPORTEDFALL{\char"0234\char"0325} \SUPPORTEDANCH{\char"006E\char"0325} \SUPPORTEDANCH{\char"0235\char"0325} \SUPPORTEDANCH{\char"0070\char"0325} \SUPPORTEDFALL{\char"01A5\char"0325} \SUPPORTEDANCH{\char"0071\char"0325} \SUPPORTEDFALL{\char"02A0\char"0325} \SUPPORTEDANCH{\char"0074\char"0325} \SUPPORTEDFALL{\char"01AB\char"0325} \SUPPORTEDANCH{\char"01AD\char"0325} \SUPPORTEDANCH{\char"0236\char"0325} \SUPPORTEDANCH{\char"0075\char"0325} \MISSINGGLYPH{\char"AB52\char"0325} \SUPPORTEDANCH{\char"0079\char"0325} \SUPPORTEDFALL{\char"01B4\char"0325} \MISSINGGLYPH{\char"1EFF\char"0325}

\SUPPORTEDANCH{\char"0048\char"0326} \SUPPORTEDFALL{\char"2C67\char"0326} \SUPPORTEDANCH{\char"004B\char"0326} \SUPPORTEDANCH{\char"2C69\char"0326} \SUPPORTEDANCH{\char"004D\char"0326} \MISSINGGLYPH{\char"2C6E\char"0326} \SUPPORTEDANCH{\char"0050\char"0326} \SUPPORTEDFALL{\char"01A4\char"0326} \SUPPORTEDANCH{\char"0052\char"0326} \SUPPORTEDFALL{\char"2C64\char"0326} \SUPPORTEDPREC{\char"0054\char"0326} \SUPPORTEDFALL{\char"01AC\char"0326} \SUPPORTEDFALL{\char"01AE\char"0326} \SUPPORTEDANCH{\char"0059\char"0326} \SUPPORTEDFALL{\char"01B3\char"0326} \SUPPORTEDANCH{\char"0061\char"0326} \SUPPORTEDFALL{\char"1E9A\char"0326} \SUPPORTEDANCH{\char"0062\char"0326} \SUPPORTEDANCH{\char"0183\char"0326} \SUPPORTEDANCH{\char"0063\char"0326} \SUPPORTEDFALL{\char"0188\char"0326} \SUPPORTEDANCH{\char"0064\char"0326} \SUPPORTEDANCH{\char"018C\char"0326} \SUPPORTEDFALL{\char"0221\char"0326} \SUPPORTEDANCH{\char"0068\char"0326} \SUPPORTEDANCH{\char"2C68\char"0326} \SUPPORTEDANCH{\char"006B\char"0326} \SUPPORTEDANCH{\char"0199\char"0326} \SUPPORTEDANCH{\char"2C6A\char"0326} \SUPPORTEDANCH{\char"006C\char"0326} \SUPPORTEDFALL{\char"0234\char"0326} \SUPPORTEDANCH{\char"006E\char"0326} \SUPPORTEDANCH{\char"0235\char"0326} \SUPPORTEDANCH{\char"0070\char"0326} \SUPPORTEDFALL{\char"01A5\char"0326} \SUPPORTEDANCH{\char"0071\char"0326} \SUPPORTEDFALL{\char"02A0\char"0326} \SUPPORTEDPREC{\char"0074\char"0326} \SUPPORTEDFALL{\char"01AB\char"0326} \SUPPORTEDANCH{\char"01AD\char"0326} \SUPPORTEDANCH{\char"0236\char"0326} \SUPPORTEDANCH{\char"0075\char"0326} \MISSINGGLYPH{\char"AB52\char"0326} \SUPPORTEDANCH{\char"0079\char"0326} \SUPPORTEDFALL{\char"01B4\char"0326} \MISSINGGLYPH{\char"1EFF\char"0326}

\SUPPORTEDANCH{\char"0048\char"0329} \SUPPORTEDFALL{\char"2C67\char"0329} \SUPPORTEDANCH{\char"004B\char"0329} \SUPPORTEDANCH{\char"2C69\char"0329} \SUPPORTEDANCH{\char"004D\char"0329} \MISSINGGLYPH{\char"2C6E\char"0329} \SUPPORTEDANCH{\char"0050\char"0329} \SUPPORTEDFALL{\char"01A4\char"0329} \SUPPORTEDANCH{\char"0052\char"0329} \SUPPORTEDFALL{\char"2C64\char"0329} \SUPPORTEDANCH{\char"0054\char"0329} \SUPPORTEDFALL{\char"01AC\char"0329} \SUPPORTEDFALL{\char"01AE\char"0329} \SUPPORTEDANCH{\char"0059\char"0329} \SUPPORTEDFALL{\char"01B3\char"0329} \SUPPORTEDANCH{\char"0061\char"0329} \SUPPORTEDFALL{\char"1E9A\char"0329} \SUPPORTEDANCH{\char"0062\char"0329} \SUPPORTEDANCH{\char"0183\char"0329} \SUPPORTEDANCH{\char"0063\char"0329} \SUPPORTEDFALL{\char"0188\char"0329} \SUPPORTEDANCH{\char"0064\char"0329} \SUPPORTEDANCH{\char"018C\char"0329} \SUPPORTEDFALL{\char"0221\char"0329} \SUPPORTEDANCH{\char"0068\char"0329} \SUPPORTEDANCH{\char"2C68\char"0329} \SUPPORTEDANCH{\char"006B\char"0329} \SUPPORTEDANCH{\char"0199\char"0329} \SUPPORTEDANCH{\char"2C6A\char"0329} \SUPPORTEDANCH{\char"006C\char"0329} \SUPPORTEDFALL{\char"0234\char"0329} \SUPPORTEDANCH{\char"006E\char"0329} \SUPPORTEDANCH{\char"0235\char"0329} \SUPPORTEDANCH{\char"0070\char"0329} \SUPPORTEDFALL{\char"01A5\char"0329} \SUPPORTEDANCH{\char"0071\char"0329} \SUPPORTEDFALL{\char"02A0\char"0329} \SUPPORTEDANCH{\char"0074\char"0329} \SUPPORTEDFALL{\char"01AB\char"0329} \SUPPORTEDANCH{\char"01AD\char"0329} \SUPPORTEDANCH{\char"0236\char"0329} \SUPPORTEDANCH{\char"0075\char"0329} \MISSINGGLYPH{\char"AB52\char"0329} \SUPPORTEDANCH{\char"0079\char"0329} \SUPPORTEDFALL{\char"01B4\char"0329} \MISSINGGLYPH{\char"1EFF\char"0329}

\SUPPORTEDANCH{\char"0048\char"032A} \SUPPORTEDFALL{\char"2C67\char"032A} \SUPPORTEDANCH{\char"004B\char"032A} \SUPPORTEDANCH{\char"2C69\char"032A} \SUPPORTEDANCH{\char"004D\char"032A} \MISSINGGLYPH{\char"2C6E\char"032A} \SUPPORTEDANCH{\char"0050\char"032A} \SUPPORTEDFALL{\char"01A4\char"032A} \SUPPORTEDANCH{\char"0052\char"032A} \SUPPORTEDFALL{\char"2C64\char"032A} \SUPPORTEDANCH{\char"0054\char"032A} \SUPPORTEDFALL{\char"01AC\char"032A} \SUPPORTEDFALL{\char"01AE\char"032A} \SUPPORTEDANCH{\char"0059\char"032A} \SUPPORTEDFALL{\char"01B3\char"032A} \SUPPORTEDANCH{\char"0061\char"032A} \SUPPORTEDFALL{\char"1E9A\char"032A} \SUPPORTEDANCH{\char"0062\char"032A} \SUPPORTEDANCH{\char"0183\char"032A} \SUPPORTEDANCH{\char"0063\char"032A} \SUPPORTEDFALL{\char"0188\char"032A} \SUPPORTEDANCH{\char"0064\char"032A} \SUPPORTEDANCH{\char"018C\char"032A} \SUPPORTEDFALL{\char"0221\char"032A} \SUPPORTEDANCH{\char"0068\char"032A} \SUPPORTEDANCH{\char"2C68\char"032A} \SUPPORTEDANCH{\char"006B\char"032A} \SUPPORTEDANCH{\char"0199\char"032A} \SUPPORTEDANCH{\char"2C6A\char"032A} \SUPPORTEDANCH{\char"006C\char"032A} \SUPPORTEDFALL{\char"0234\char"032A} \SUPPORTEDANCH{\char"006E\char"032A} \SUPPORTEDANCH{\char"0235\char"032A} \SUPPORTEDANCH{\char"0070\char"032A} \SUPPORTEDFALL{\char"01A5\char"032A} \SUPPORTEDANCH{\char"0071\char"032A} \SUPPORTEDFALL{\char"02A0\char"032A} \SUPPORTEDANCH{\char"0074\char"032A} \SUPPORTEDFALL{\char"01AB\char"032A} \SUPPORTEDANCH{\char"01AD\char"032A} \SUPPORTEDANCH{\char"0236\char"032A} \SUPPORTEDANCH{\char"0075\char"032A} \MISSINGGLYPH{\char"AB52\char"032A} \SUPPORTEDANCH{\char"0079\char"032A} \SUPPORTEDFALL{\char"01B4\char"032A} \MISSINGGLYPH{\char"1EFF\char"032A}

\SUPPORTEDANCH{\char"0048\char"032B} \SUPPORTEDFALL{\char"2C67\char"032B} \SUPPORTEDANCH{\char"004B\char"032B} \SUPPORTEDANCH{\char"2C69\char"032B} \SUPPORTEDANCH{\char"004D\char"032B} \MISSINGGLYPH{\char"2C6E\char"032B} \SUPPORTEDANCH{\char"0050\char"032B} \SUPPORTEDFALL{\char"01A4\char"032B} \SUPPORTEDANCH{\char"0052\char"032B} \SUPPORTEDFALL{\char"2C64\char"032B} \SUPPORTEDANCH{\char"0054\char"032B} \SUPPORTEDFALL{\char"01AC\char"032B} \SUPPORTEDFALL{\char"01AE\char"032B} \SUPPORTEDANCH{\char"0059\char"032B} \SUPPORTEDFALL{\char"01B3\char"032B} \SUPPORTEDANCH{\char"0061\char"032B} \SUPPORTEDFALL{\char"1E9A\char"032B} \SUPPORTEDANCH{\char"0062\char"032B} \SUPPORTEDANCH{\char"0183\char"032B} \SUPPORTEDANCH{\char"0063\char"032B} \SUPPORTEDFALL{\char"0188\char"032B} \SUPPORTEDANCH{\char"0064\char"032B} \SUPPORTEDANCH{\char"018C\char"032B} \SUPPORTEDFALL{\char"0221\char"032B} \SUPPORTEDANCH{\char"0068\char"032B} \SUPPORTEDANCH{\char"2C68\char"032B} \SUPPORTEDANCH{\char"006B\char"032B} \SUPPORTEDANCH{\char"0199\char"032B} \SUPPORTEDANCH{\char"2C6A\char"032B} \SUPPORTEDANCH{\char"006C\char"032B} \SUPPORTEDFALL{\char"0234\char"032B} \SUPPORTEDANCH{\char"006E\char"032B} \SUPPORTEDANCH{\char"0235\char"032B} \SUPPORTEDANCH{\char"0070\char"032B} \SUPPORTEDFALL{\char"01A5\char"032B} \SUPPORTEDANCH{\char"0071\char"032B} \SUPPORTEDFALL{\char"02A0\char"032B} \SUPPORTEDANCH{\char"0074\char"032B} \SUPPORTEDFALL{\char"01AB\char"032B} \SUPPORTEDANCH{\char"01AD\char"032B} \SUPPORTEDANCH{\char"0236\char"032B} \SUPPORTEDANCH{\char"0075\char"032B} \MISSINGGLYPH{\char"AB52\char"032B} \SUPPORTEDANCH{\char"0079\char"032B} \SUPPORTEDFALL{\char"01B4\char"032B} \MISSINGGLYPH{\char"1EFF\char"032B}

\SUPPORTEDANCH{\char"0048\char"032C} \SUPPORTEDFALL{\char"2C67\char"032C} \SUPPORTEDANCH{\char"004B\char"032C} \SUPPORTEDANCH{\char"2C69\char"032C} \SUPPORTEDANCH{\char"004D\char"032C} \MISSINGGLYPH{\char"2C6E\char"032C} \SUPPORTEDANCH{\char"0050\char"032C} \SUPPORTEDFALL{\char"01A4\char"032C} \SUPPORTEDANCH{\char"0052\char"032C} \SUPPORTEDFALL{\char"2C64\char"032C} \SUPPORTEDANCH{\char"0054\char"032C} \SUPPORTEDFALL{\char"01AC\char"032C} \SUPPORTEDFALL{\char"01AE\char"032C} \SUPPORTEDANCH{\char"0059\char"032C} \SUPPORTEDFALL{\char"01B3\char"032C} \SUPPORTEDANCH{\char"0061\char"032C} \SUPPORTEDFALL{\char"1E9A\char"032C} \SUPPORTEDANCH{\char"0062\char"032C} \SUPPORTEDANCH{\char"0183\char"032C} \SUPPORTEDANCH{\char"0063\char"032C} \SUPPORTEDFALL{\char"0188\char"032C} \SUPPORTEDANCH{\char"0064\char"032C} \SUPPORTEDANCH{\char"018C\char"032C} \SUPPORTEDFALL{\char"0221\char"032C} \SUPPORTEDANCH{\char"0068\char"032C} \SUPPORTEDANCH{\char"2C68\char"032C} \SUPPORTEDANCH{\char"006B\char"032C} \SUPPORTEDANCH{\char"0199\char"032C} \SUPPORTEDANCH{\char"2C6A\char"032C} \SUPPORTEDANCH{\char"006C\char"032C} \SUPPORTEDFALL{\char"0234\char"032C} \SUPPORTEDANCH{\char"006E\char"032C} \SUPPORTEDANCH{\char"0235\char"032C} \SUPPORTEDANCH{\char"0070\char"032C} \SUPPORTEDFALL{\char"01A5\char"032C} \SUPPORTEDANCH{\char"0071\char"032C} \SUPPORTEDFALL{\char"02A0\char"032C} \SUPPORTEDANCH{\char"0074\char"032C} \SUPPORTEDFALL{\char"01AB\char"032C} \SUPPORTEDANCH{\char"01AD\char"032C} \SUPPORTEDANCH{\char"0236\char"032C} \SUPPORTEDANCH{\char"0075\char"032C} \MISSINGGLYPH{\char"AB52\char"032C} \SUPPORTEDANCH{\char"0079\char"032C} \SUPPORTEDFALL{\char"01B4\char"032C} \MISSINGGLYPH{\char"1EFF\char"032C}

\SUPPORTEDANCH{\char"0048\char"032D} \SUPPORTEDFALL{\char"2C67\char"032D} \SUPPORTEDANCH{\char"004B\char"032D} \SUPPORTEDANCH{\char"2C69\char"032D} \SUPPORTEDANCH{\char"004D\char"032D} \MISSINGGLYPH{\char"2C6E\char"032D} \SUPPORTEDANCH{\char"0050\char"032D} \SUPPORTEDFALL{\char"01A4\char"032D} \SUPPORTEDANCH{\char"0052\char"032D} \SUPPORTEDFALL{\char"2C64\char"032D} \SUPPORTEDPREC{\char"0054\char"032D} \SUPPORTEDFALL{\char"01AC\char"032D} \SUPPORTEDFALL{\char"01AE\char"032D} \SUPPORTEDANCH{\char"0059\char"032D} \SUPPORTEDFALL{\char"01B3\char"032D} \SUPPORTEDANCH{\char"0061\char"032D} \SUPPORTEDFALL{\char"1E9A\char"032D} \SUPPORTEDANCH{\char"0062\char"032D} \SUPPORTEDANCH{\char"0183\char"032D} \SUPPORTEDANCH{\char"0063\char"032D} \SUPPORTEDFALL{\char"0188\char"032D} \SUPPORTEDPREC{\char"0064\char"032D} \SUPPORTEDANCH{\char"018C\char"032D} \SUPPORTEDFALL{\char"0221\char"032D} \SUPPORTEDANCH{\char"0068\char"032D} \SUPPORTEDANCH{\char"2C68\char"032D} \SUPPORTEDANCH{\char"006B\char"032D} \SUPPORTEDANCH{\char"0199\char"032D} \SUPPORTEDANCH{\char"2C6A\char"032D} \SUPPORTEDPREC{\char"006C\char"032D} \SUPPORTEDFALL{\char"0234\char"032D} \SUPPORTEDPREC{\char"006E\char"032D} \SUPPORTEDANCH{\char"0235\char"032D} \SUPPORTEDANCH{\char"0070\char"032D} \SUPPORTEDFALL{\char"01A5\char"032D} \SUPPORTEDANCH{\char"0071\char"032D} \SUPPORTEDFALL{\char"02A0\char"032D} \SUPPORTEDPREC{\char"0074\char"032D} \SUPPORTEDFALL{\char"01AB\char"032D} \SUPPORTEDANCH{\char"01AD\char"032D} \SUPPORTEDANCH{\char"0236\char"032D} \SUPPORTEDPREC{\char"0075\char"032D} \MISSINGGLYPH{\char"AB52\char"032D} \SUPPORTEDANCH{\char"0079\char"032D} \SUPPORTEDFALL{\char"01B4\char"032D} \MISSINGGLYPH{\char"1EFF\char"032D}

\SUPPORTEDPREC{\char"0048\char"032E} \SUPPORTEDFALL{\char"2C67\char"032E} \SUPPORTEDANCH{\char"004B\char"032E} \SUPPORTEDANCH{\char"2C69\char"032E} \SUPPORTEDANCH{\char"004D\char"032E} \MISSINGGLYPH{\char"2C6E\char"032E} \SUPPORTEDANCH{\char"0050\char"032E} \SUPPORTEDFALL{\char"01A4\char"032E} \SUPPORTEDANCH{\char"0052\char"032E} \SUPPORTEDFALL{\char"2C64\char"032E} \SUPPORTEDANCH{\char"0054\char"032E} \SUPPORTEDFALL{\char"01AC\char"032E} \SUPPORTEDFALL{\char"01AE\char"032E} \SUPPORTEDANCH{\char"0059\char"032E} \SUPPORTEDFALL{\char"01B3\char"032E} \SUPPORTEDANCH{\char"0061\char"032E} \SUPPORTEDFALL{\char"1E9A\char"032E} \SUPPORTEDANCH{\char"0062\char"032E} \SUPPORTEDANCH{\char"0183\char"032E} \SUPPORTEDANCH{\char"0063\char"032E} \SUPPORTEDFALL{\char"0188\char"032E} \SUPPORTEDANCH{\char"0064\char"032E} \SUPPORTEDANCH{\char"018C\char"032E} \SUPPORTEDFALL{\char"0221\char"032E} \SUPPORTEDPREC{\char"0068\char"032E} \SUPPORTEDANCH{\char"2C68\char"032E} \SUPPORTEDANCH{\char"006B\char"032E} \SUPPORTEDANCH{\char"0199\char"032E} \SUPPORTEDANCH{\char"2C6A\char"032E} \SUPPORTEDANCH{\char"006C\char"032E} \SUPPORTEDFALL{\char"0234\char"032E} \SUPPORTEDANCH{\char"006E\char"032E} \SUPPORTEDANCH{\char"0235\char"032E} \SUPPORTEDANCH{\char"0070\char"032E} \SUPPORTEDFALL{\char"01A5\char"032E} \SUPPORTEDANCH{\char"0071\char"032E} \SUPPORTEDFALL{\char"02A0\char"032E} \SUPPORTEDANCH{\char"0074\char"032E} \SUPPORTEDFALL{\char"01AB\char"032E} \SUPPORTEDANCH{\char"01AD\char"032E} \SUPPORTEDANCH{\char"0236\char"032E} \SUPPORTEDANCH{\char"0075\char"032E} \MISSINGGLYPH{\char"AB52\char"032E} \SUPPORTEDANCH{\char"0079\char"032E} \SUPPORTEDFALL{\char"01B4\char"032E} \MISSINGGLYPH{\char"1EFF\char"032E}

\SUPPORTEDANCH{\char"0048\char"032F} \SUPPORTEDFALL{\char"2C67\char"032F} \SUPPORTEDANCH{\char"004B\char"032F} \SUPPORTEDANCH{\char"2C69\char"032F} \SUPPORTEDANCH{\char"004D\char"032F} \MISSINGGLYPH{\char"2C6E\char"032F} \SUPPORTEDANCH{\char"0050\char"032F} \SUPPORTEDFALL{\char"01A4\char"032F} \SUPPORTEDANCH{\char"0052\char"032F} \SUPPORTEDFALL{\char"2C64\char"032F} \SUPPORTEDANCH{\char"0054\char"032F} \SUPPORTEDFALL{\char"01AC\char"032F} \SUPPORTEDFALL{\char"01AE\char"032F} \SUPPORTEDANCH{\char"0059\char"032F} \SUPPORTEDFALL{\char"01B3\char"032F} \SUPPORTEDANCH{\char"0061\char"032F} \SUPPORTEDFALL{\char"1E9A\char"032F} \SUPPORTEDANCH{\char"0062\char"032F} \SUPPORTEDANCH{\char"0183\char"032F} \SUPPORTEDANCH{\char"0063\char"032F} \SUPPORTEDFALL{\char"0188\char"032F} \SUPPORTEDANCH{\char"0064\char"032F} \SUPPORTEDANCH{\char"018C\char"032F} \SUPPORTEDFALL{\char"0221\char"032F} \SUPPORTEDANCH{\char"0068\char"032F} \SUPPORTEDANCH{\char"2C68\char"032F} \SUPPORTEDANCH{\char"006B\char"032F} \SUPPORTEDANCH{\char"0199\char"032F} \SUPPORTEDANCH{\char"2C6A\char"032F} \SUPPORTEDANCH{\char"006C\char"032F} \SUPPORTEDFALL{\char"0234\char"032F} \SUPPORTEDANCH{\char"006E\char"032F} \SUPPORTEDANCH{\char"0235\char"032F} \SUPPORTEDANCH{\char"0070\char"032F} \SUPPORTEDFALL{\char"01A5\char"032F} \SUPPORTEDANCH{\char"0071\char"032F} \SUPPORTEDFALL{\char"02A0\char"032F} \SUPPORTEDANCH{\char"0074\char"032F} \SUPPORTEDFALL{\char"01AB\char"032F} \SUPPORTEDANCH{\char"01AD\char"032F} \SUPPORTEDANCH{\char"0236\char"032F} \SUPPORTEDANCH{\char"0075\char"032F} \MISSINGGLYPH{\char"AB52\char"032F} \SUPPORTEDANCH{\char"0079\char"032F} \SUPPORTEDFALL{\char"01B4\char"032F} \MISSINGGLYPH{\char"1EFF\char"032F}

\SUPPORTEDANCH{\char"0048\char"031D} \SUPPORTEDFALL{\char"2C67\char"031D} \SUPPORTEDANCH{\char"004B\char"031D} \SUPPORTEDANCH{\char"2C69\char"031D} \SUPPORTEDANCH{\char"004D\char"031D} \MISSINGGLYPH{\char"2C6E\char"031D} \SUPPORTEDANCH{\char"0050\char"031D} \SUPPORTEDFALL{\char"01A4\char"031D} \SUPPORTEDANCH{\char"0052\char"031D} \SUPPORTEDFALL{\char"2C64\char"031D} \SUPPORTEDANCH{\char"0054\char"031D} \SUPPORTEDFALL{\char"01AC\char"031D} \SUPPORTEDFALL{\char"01AE\char"031D} \SUPPORTEDANCH{\char"0059\char"031D} \SUPPORTEDFALL{\char"01B3\char"031D} \SUPPORTEDANCH{\char"0061\char"031D} \SUPPORTEDFALL{\char"1E9A\char"031D} \SUPPORTEDANCH{\char"0062\char"031D} \SUPPORTEDANCH{\char"0183\char"031D} \SUPPORTEDANCH{\char"0063\char"031D} \SUPPORTEDFALL{\char"0188\char"031D} \SUPPORTEDANCH{\char"0064\char"031D} \SUPPORTEDANCH{\char"018C\char"031D} \SUPPORTEDFALL{\char"0221\char"031D} \SUPPORTEDANCH{\char"0068\char"031D} \SUPPORTEDANCH{\char"2C68\char"031D} \SUPPORTEDANCH{\char"006B\char"031D} \SUPPORTEDANCH{\char"0199\char"031D} \SUPPORTEDANCH{\char"2C6A\char"031D} \SUPPORTEDANCH{\char"006C\char"031D} \SUPPORTEDFALL{\char"0234\char"031D} \SUPPORTEDANCH{\char"006E\char"031D} \SUPPORTEDANCH{\char"0235\char"031D} \SUPPORTEDANCH{\char"0070\char"031D} \SUPPORTEDFALL{\char"01A5\char"031D} \SUPPORTEDANCH{\char"0071\char"031D} \SUPPORTEDFALL{\char"02A0\char"031D} \SUPPORTEDANCH{\char"0074\char"031D} \SUPPORTEDFALL{\char"01AB\char"031D} \SUPPORTEDANCH{\char"01AD\char"031D} \SUPPORTEDANCH{\char"0236\char"031D} \SUPPORTEDANCH{\char"0075\char"031D} \MISSINGGLYPH{\char"AB52\char"031D} \SUPPORTEDANCH{\char"0079\char"031D} \SUPPORTEDFALL{\char"01B4\char"031D} \MISSINGGLYPH{\char"1EFF\char"031D}

\SUPPORTEDANCH{\char"0048\char"031E} \SUPPORTEDFALL{\char"2C67\char"031E} \SUPPORTEDANCH{\char"004B\char"031E} \SUPPORTEDANCH{\char"2C69\char"031E} \SUPPORTEDANCH{\char"004D\char"031E} \MISSINGGLYPH{\char"2C6E\char"031E} \SUPPORTEDANCH{\char"0050\char"031E} \SUPPORTEDFALL{\char"01A4\char"031E} \SUPPORTEDANCH{\char"0052\char"031E} \SUPPORTEDFALL{\char"2C64\char"031E} \SUPPORTEDANCH{\char"0054\char"031E} \SUPPORTEDFALL{\char"01AC\char"031E} \SUPPORTEDFALL{\char"01AE\char"031E} \SUPPORTEDANCH{\char"0059\char"031E} \SUPPORTEDFALL{\char"01B3\char"031E} \SUPPORTEDANCH{\char"0061\char"031E} \SUPPORTEDFALL{\char"1E9A\char"031E} \SUPPORTEDANCH{\char"0062\char"031E} \SUPPORTEDANCH{\char"0183\char"031E} \SUPPORTEDANCH{\char"0063\char"031E} \SUPPORTEDFALL{\char"0188\char"031E} \SUPPORTEDANCH{\char"0064\char"031E} \SUPPORTEDANCH{\char"018C\char"031E} \SUPPORTEDFALL{\char"0221\char"031E} \SUPPORTEDANCH{\char"0068\char"031E} \SUPPORTEDANCH{\char"2C68\char"031E} \SUPPORTEDANCH{\char"006B\char"031E} \SUPPORTEDANCH{\char"0199\char"031E} \SUPPORTEDANCH{\char"2C6A\char"031E} \SUPPORTEDANCH{\char"006C\char"031E} \SUPPORTEDFALL{\char"0234\char"031E} \SUPPORTEDANCH{\char"006E\char"031E} \SUPPORTEDANCH{\char"0235\char"031E} \SUPPORTEDANCH{\char"0070\char"031E} \SUPPORTEDFALL{\char"01A5\char"031E} \SUPPORTEDANCH{\char"0071\char"031E} \SUPPORTEDFALL{\char"02A0\char"031E} \SUPPORTEDANCH{\char"0074\char"031E} \SUPPORTEDFALL{\char"01AB\char"031E} \SUPPORTEDANCH{\char"01AD\char"031E} \SUPPORTEDANCH{\char"0236\char"031E} \SUPPORTEDANCH{\char"0075\char"031E} \MISSINGGLYPH{\char"AB52\char"031E} \SUPPORTEDANCH{\char"0079\char"031E} \SUPPORTEDFALL{\char"01B4\char"031E} \MISSINGGLYPH{\char"1EFF\char"031E}

\SUPPORTEDANCH{\char"0048\char"031F} \SUPPORTEDFALL{\char"2C67\char"031F} \SUPPORTEDANCH{\char"004B\char"031F} \SUPPORTEDANCH{\char"2C69\char"031F} \SUPPORTEDANCH{\char"004D\char"031F} \MISSINGGLYPH{\char"2C6E\char"031F} \SUPPORTEDANCH{\char"0050\char"031F} \SUPPORTEDFALL{\char"01A4\char"031F} \SUPPORTEDANCH{\char"0052\char"031F} \SUPPORTEDFALL{\char"2C64\char"031F} \SUPPORTEDANCH{\char"0054\char"031F} \SUPPORTEDFALL{\char"01AC\char"031F} \SUPPORTEDFALL{\char"01AE\char"031F} \SUPPORTEDANCH{\char"0059\char"031F} \SUPPORTEDFALL{\char"01B3\char"031F} \SUPPORTEDANCH{\char"0061\char"031F} \SUPPORTEDFALL{\char"1E9A\char"031F} \SUPPORTEDANCH{\char"0062\char"031F} \SUPPORTEDANCH{\char"0183\char"031F} \SUPPORTEDANCH{\char"0063\char"031F} \SUPPORTEDFALL{\char"0188\char"031F} \SUPPORTEDANCH{\char"0064\char"031F} \SUPPORTEDANCH{\char"018C\char"031F} \SUPPORTEDFALL{\char"0221\char"031F} \SUPPORTEDANCH{\char"0068\char"031F} \SUPPORTEDANCH{\char"2C68\char"031F} \SUPPORTEDANCH{\char"006B\char"031F} \SUPPORTEDANCH{\char"0199\char"031F} \SUPPORTEDANCH{\char"2C6A\char"031F} \SUPPORTEDANCH{\char"006C\char"031F} \SUPPORTEDFALL{\char"0234\char"031F} \SUPPORTEDANCH{\char"006E\char"031F} \SUPPORTEDANCH{\char"0235\char"031F} \SUPPORTEDANCH{\char"0070\char"031F} \SUPPORTEDFALL{\char"01A5\char"031F} \SUPPORTEDANCH{\char"0071\char"031F} \SUPPORTEDFALL{\char"02A0\char"031F} \SUPPORTEDANCH{\char"0074\char"031F} \SUPPORTEDFALL{\char"01AB\char"031F} \SUPPORTEDANCH{\char"01AD\char"031F} \SUPPORTEDANCH{\char"0236\char"031F} \SUPPORTEDANCH{\char"0075\char"031F} \MISSINGGLYPH{\char"AB52\char"031F} \SUPPORTEDANCH{\char"0079\char"031F} \SUPPORTEDFALL{\char"01B4\char"031F} \MISSINGGLYPH{\char"1EFF\char"031F}

\SUPPORTEDANCH{\char"0048\char"0330} \SUPPORTEDFALL{\char"2C67\char"0330} \SUPPORTEDANCH{\char"004B\char"0330} \SUPPORTEDANCH{\char"2C69\char"0330} \SUPPORTEDANCH{\char"004D\char"0330} \MISSINGGLYPH{\char"2C6E\char"0330} \SUPPORTEDANCH{\char"0050\char"0330} \SUPPORTEDFALL{\char"01A4\char"0330} \SUPPORTEDANCH{\char"0052\char"0330} \SUPPORTEDFALL{\char"2C64\char"0330} \SUPPORTEDANCH{\char"0054\char"0330} \SUPPORTEDFALL{\char"01AC\char"0330} \SUPPORTEDFALL{\char"01AE\char"0330} \SUPPORTEDANCH{\char"0059\char"0330} \SUPPORTEDFALL{\char"01B3\char"0330} \SUPPORTEDANCH{\char"0061\char"0330} \SUPPORTEDFALL{\char"1E9A\char"0330} \SUPPORTEDANCH{\char"0062\char"0330} \SUPPORTEDANCH{\char"0183\char"0330} \SUPPORTEDANCH{\char"0063\char"0330} \SUPPORTEDFALL{\char"0188\char"0330} \SUPPORTEDANCH{\char"0064\char"0330} \SUPPORTEDANCH{\char"018C\char"0330} \SUPPORTEDFALL{\char"0221\char"0330} \SUPPORTEDANCH{\char"0068\char"0330} \SUPPORTEDANCH{\char"2C68\char"0330} \SUPPORTEDANCH{\char"006B\char"0330} \SUPPORTEDANCH{\char"0199\char"0330} \SUPPORTEDANCH{\char"2C6A\char"0330} \SUPPORTEDANCH{\char"006C\char"0330} \SUPPORTEDFALL{\char"0234\char"0330} \SUPPORTEDANCH{\char"006E\char"0330} \SUPPORTEDANCH{\char"0235\char"0330} \SUPPORTEDANCH{\char"0070\char"0330} \SUPPORTEDFALL{\char"01A5\char"0330} \SUPPORTEDANCH{\char"0071\char"0330} \SUPPORTEDFALL{\char"02A0\char"0330} \SUPPORTEDANCH{\char"0074\char"0330} \SUPPORTEDFALL{\char"01AB\char"0330} \SUPPORTEDANCH{\char"01AD\char"0330} \SUPPORTEDANCH{\char"0236\char"0330} \SUPPORTEDPREC{\char"0075\char"0330} \MISSINGGLYPH{\char"AB52\char"0330} \SUPPORTEDANCH{\char"0079\char"0330} \SUPPORTEDFALL{\char"01B4\char"0330} \MISSINGGLYPH{\char"1EFF\char"0330}

\SUPPORTEDANCH{\char"0048\char"0331} \SUPPORTEDFALL{\char"2C67\char"0331} \SUPPORTEDPREC{\char"004B\char"0331} \SUPPORTEDANCH{\char"2C69\char"0331} \SUPPORTEDANCH{\char"004D\char"0331} \MISSINGGLYPH{\char"2C6E\char"0331} \SUPPORTEDANCH{\char"0050\char"0331} \SUPPORTEDFALL{\char"01A4\char"0331} \SUPPORTEDPREC{\char"0052\char"0331} \SUPPORTEDFALL{\char"2C64\char"0331} \SUPPORTEDPREC{\char"0054\char"0331} \SUPPORTEDFALL{\char"01AC\char"0331} \SUPPORTEDFALL{\char"01AE\char"0331} \SUPPORTEDANCH{\char"0059\char"0331} \SUPPORTEDFALL{\char"01B3\char"0331} \SUPPORTEDANCH{\char"0061\char"0331} \SUPPORTEDFALL{\char"1E9A\char"0331} \SUPPORTEDPREC{\char"0062\char"0331} \SUPPORTEDANCH{\char"0183\char"0331} \SUPPORTEDANCH{\char"0063\char"0331} \SUPPORTEDFALL{\char"0188\char"0331} \SUPPORTEDPREC{\char"0064\char"0331} \SUPPORTEDANCH{\char"018C\char"0331} \SUPPORTEDFALL{\char"0221\char"0331} \SUPPORTEDPREC{\char"0068\char"0331} \SUPPORTEDANCH{\char"2C68\char"0331} \SUPPORTEDPREC{\char"006B\char"0331} \SUPPORTEDANCH{\char"0199\char"0331} \SUPPORTEDANCH{\char"2C6A\char"0331} \SUPPORTEDPREC{\char"006C\char"0331} \SUPPORTEDFALL{\char"0234\char"0331} \SUPPORTEDPREC{\char"006E\char"0331} \SUPPORTEDANCH{\char"0235\char"0331} \SUPPORTEDANCH{\char"0070\char"0331} \SUPPORTEDFALL{\char"01A5\char"0331} \SUPPORTEDANCH{\char"0071\char"0331} \SUPPORTEDFALL{\char"02A0\char"0331} \SUPPORTEDPREC{\char"0074\char"0331} \SUPPORTEDFALL{\char"01AB\char"0331} \SUPPORTEDANCH{\char"01AD\char"0331} \SUPPORTEDANCH{\char"0236\char"0331} \SUPPORTEDANCH{\char"0075\char"0331} \MISSINGGLYPH{\char"AB52\char"0331} \SUPPORTEDANCH{\char"0079\char"0331} \SUPPORTEDFALL{\char"01B4\char"0331} \MISSINGGLYPH{\char"1EFF\char"0331}

\SUPPORTEDANCH{\char"0048\char"0333} \SUPPORTEDFALL{\char"2C67\char"0333} \SUPPORTEDANCH{\char"004B\char"0333} \SUPPORTEDANCH{\char"2C69\char"0333} \SUPPORTEDANCH{\char"004D\char"0333} \MISSINGGLYPH{\char"2C6E\char"0333} \SUPPORTEDANCH{\char"0050\char"0333} \SUPPORTEDFALL{\char"01A4\char"0333} \SUPPORTEDANCH{\char"0052\char"0333} \SUPPORTEDFALL{\char"2C64\char"0333} \SUPPORTEDANCH{\char"0054\char"0333} \SUPPORTEDFALL{\char"01AC\char"0333} \SUPPORTEDFALL{\char"01AE\char"0333} \SUPPORTEDANCH{\char"0059\char"0333} \SUPPORTEDFALL{\char"01B3\char"0333} \SUPPORTEDANCH{\char"0061\char"0333} \SUPPORTEDFALL{\char"1E9A\char"0333} \SUPPORTEDANCH{\char"0062\char"0333} \SUPPORTEDANCH{\char"0183\char"0333} \SUPPORTEDANCH{\char"0063\char"0333} \SUPPORTEDFALL{\char"0188\char"0333} \SUPPORTEDANCH{\char"0064\char"0333} \SUPPORTEDANCH{\char"018C\char"0333} \SUPPORTEDFALL{\char"0221\char"0333} \SUPPORTEDANCH{\char"0068\char"0333} \SUPPORTEDANCH{\char"2C68\char"0333} \SUPPORTEDANCH{\char"006B\char"0333} \SUPPORTEDANCH{\char"0199\char"0333} \SUPPORTEDANCH{\char"2C6A\char"0333} \SUPPORTEDANCH{\char"006C\char"0333} \SUPPORTEDFALL{\char"0234\char"0333} \SUPPORTEDANCH{\char"006E\char"0333} \SUPPORTEDANCH{\char"0235\char"0333} \SUPPORTEDANCH{\char"0070\char"0333} \SUPPORTEDFALL{\char"01A5\char"0333} \SUPPORTEDANCH{\char"0071\char"0333} \SUPPORTEDFALL{\char"02A0\char"0333} \SUPPORTEDANCH{\char"0074\char"0333} \SUPPORTEDFALL{\char"01AB\char"0333} \SUPPORTEDANCH{\char"01AD\char"0333} \SUPPORTEDANCH{\char"0236\char"0333} \SUPPORTEDANCH{\char"0075\char"0333} \MISSINGGLYPH{\char"AB52\char"0333} \SUPPORTEDANCH{\char"0079\char"0333} \SUPPORTEDFALL{\char"01B4\char"0333} \MISSINGGLYPH{\char"1EFF\char"0333}

\SUPPORTEDANCH{\char"0048\char"0339} \SUPPORTEDFALL{\char"2C67\char"0339} \SUPPORTEDANCH{\char"004B\char"0339} \SUPPORTEDANCH{\char"2C69\char"0339} \SUPPORTEDANCH{\char"004D\char"0339} \MISSINGGLYPH{\char"2C6E\char"0339} \SUPPORTEDANCH{\char"0050\char"0339} \SUPPORTEDFALL{\char"01A4\char"0339} \SUPPORTEDANCH{\char"0052\char"0339} \SUPPORTEDFALL{\char"2C64\char"0339} \SUPPORTEDANCH{\char"0054\char"0339} \SUPPORTEDFALL{\char"01AC\char"0339} \SUPPORTEDFALL{\char"01AE\char"0339} \SUPPORTEDANCH{\char"0059\char"0339} \SUPPORTEDFALL{\char"01B3\char"0339} \SUPPORTEDANCH{\char"0061\char"0339} \SUPPORTEDFALL{\char"1E9A\char"0339} \SUPPORTEDANCH{\char"0062\char"0339} \SUPPORTEDANCH{\char"0183\char"0339} \SUPPORTEDANCH{\char"0063\char"0339} \SUPPORTEDFALL{\char"0188\char"0339} \SUPPORTEDANCH{\char"0064\char"0339} \SUPPORTEDANCH{\char"018C\char"0339} \SUPPORTEDFALL{\char"0221\char"0339} \SUPPORTEDANCH{\char"0068\char"0339} \SUPPORTEDANCH{\char"2C68\char"0339} \SUPPORTEDANCH{\char"006B\char"0339} \SUPPORTEDANCH{\char"0199\char"0339} \SUPPORTEDANCH{\char"2C6A\char"0339} \SUPPORTEDANCH{\char"006C\char"0339} \SUPPORTEDFALL{\char"0234\char"0339} \SUPPORTEDANCH{\char"006E\char"0339} \SUPPORTEDANCH{\char"0235\char"0339} \SUPPORTEDANCH{\char"0070\char"0339} \SUPPORTEDFALL{\char"01A5\char"0339} \SUPPORTEDANCH{\char"0071\char"0339} \SUPPORTEDFALL{\char"02A0\char"0339} \SUPPORTEDANCH{\char"0074\char"0339} \SUPPORTEDFALL{\char"01AB\char"0339} \SUPPORTEDANCH{\char"01AD\char"0339} \SUPPORTEDANCH{\char"0236\char"0339} \SUPPORTEDANCH{\char"0075\char"0339} \MISSINGGLYPH{\char"AB52\char"0339} \SUPPORTEDANCH{\char"0079\char"0339} \SUPPORTEDFALL{\char"01B4\char"0339} \MISSINGGLYPH{\char"1EFF\char"0339}

\SUPPORTEDANCH{\char"0048\char"033A} \SUPPORTEDFALL{\char"2C67\char"033A} \SUPPORTEDANCH{\char"004B\char"033A} \SUPPORTEDANCH{\char"2C69\char"033A} \SUPPORTEDANCH{\char"004D\char"033A} \MISSINGGLYPH{\char"2C6E\char"033A} \SUPPORTEDANCH{\char"0050\char"033A} \SUPPORTEDFALL{\char"01A4\char"033A} \SUPPORTEDANCH{\char"0052\char"033A} \SUPPORTEDFALL{\char"2C64\char"033A} \SUPPORTEDANCH{\char"0054\char"033A} \SUPPORTEDFALL{\char"01AC\char"033A} \SUPPORTEDFALL{\char"01AE\char"033A} \SUPPORTEDANCH{\char"0059\char"033A} \SUPPORTEDFALL{\char"01B3\char"033A} \SUPPORTEDANCH{\char"0061\char"033A} \SUPPORTEDFALL{\char"1E9A\char"033A} \SUPPORTEDANCH{\char"0062\char"033A} \SUPPORTEDANCH{\char"0183\char"033A} \SUPPORTEDANCH{\char"0063\char"033A} \SUPPORTEDFALL{\char"0188\char"033A} \SUPPORTEDANCH{\char"0064\char"033A} \SUPPORTEDANCH{\char"018C\char"033A} \SUPPORTEDFALL{\char"0221\char"033A} \SUPPORTEDANCH{\char"0068\char"033A} \SUPPORTEDANCH{\char"2C68\char"033A} \SUPPORTEDANCH{\char"006B\char"033A} \SUPPORTEDANCH{\char"0199\char"033A} \SUPPORTEDANCH{\char"2C6A\char"033A} \SUPPORTEDANCH{\char"006C\char"033A} \SUPPORTEDFALL{\char"0234\char"033A} \SUPPORTEDANCH{\char"006E\char"033A} \SUPPORTEDANCH{\char"0235\char"033A} \SUPPORTEDANCH{\char"0070\char"033A} \SUPPORTEDFALL{\char"01A5\char"033A} \SUPPORTEDANCH{\char"0071\char"033A} \SUPPORTEDFALL{\char"02A0\char"033A} \SUPPORTEDANCH{\char"0074\char"033A} \SUPPORTEDFALL{\char"01AB\char"033A} \SUPPORTEDANCH{\char"01AD\char"033A} \SUPPORTEDANCH{\char"0236\char"033A} \SUPPORTEDANCH{\char"0075\char"033A} \MISSINGGLYPH{\char"AB52\char"033A} \SUPPORTEDANCH{\char"0079\char"033A} \SUPPORTEDFALL{\char"01B4\char"033A} \MISSINGGLYPH{\char"1EFF\char"033A}

\SUPPORTEDANCH{\char"0048\char"033B} \SUPPORTEDFALL{\char"2C67\char"033B} \SUPPORTEDANCH{\char"004B\char"033B} \SUPPORTEDANCH{\char"2C69\char"033B} \SUPPORTEDANCH{\char"004D\char"033B} \MISSINGGLYPH{\char"2C6E\char"033B} \SUPPORTEDANCH{\char"0050\char"033B} \SUPPORTEDFALL{\char"01A4\char"033B} \SUPPORTEDANCH{\char"0052\char"033B} \SUPPORTEDFALL{\char"2C64\char"033B} \SUPPORTEDANCH{\char"0054\char"033B} \SUPPORTEDFALL{\char"01AC\char"033B} \SUPPORTEDFALL{\char"01AE\char"033B} \SUPPORTEDANCH{\char"0059\char"033B} \SUPPORTEDFALL{\char"01B3\char"033B} \SUPPORTEDANCH{\char"0061\char"033B} \SUPPORTEDFALL{\char"1E9A\char"033B} \SUPPORTEDANCH{\char"0062\char"033B} \SUPPORTEDANCH{\char"0183\char"033B} \SUPPORTEDANCH{\char"0063\char"033B} \SUPPORTEDFALL{\char"0188\char"033B} \SUPPORTEDANCH{\char"0064\char"033B} \SUPPORTEDANCH{\char"018C\char"033B} \SUPPORTEDFALL{\char"0221\char"033B} \SUPPORTEDANCH{\char"0068\char"033B} \SUPPORTEDANCH{\char"2C68\char"033B} \SUPPORTEDANCH{\char"006B\char"033B} \SUPPORTEDANCH{\char"0199\char"033B} \SUPPORTEDANCH{\char"2C6A\char"033B} \SUPPORTEDANCH{\char"006C\char"033B} \SUPPORTEDFALL{\char"0234\char"033B} \SUPPORTEDANCH{\char"006E\char"033B} \SUPPORTEDANCH{\char"0235\char"033B} \SUPPORTEDANCH{\char"0070\char"033B} \SUPPORTEDFALL{\char"01A5\char"033B} \SUPPORTEDANCH{\char"0071\char"033B} \SUPPORTEDFALL{\char"02A0\char"033B} \SUPPORTEDANCH{\char"0074\char"033B} \SUPPORTEDFALL{\char"01AB\char"033B} \SUPPORTEDANCH{\char"01AD\char"033B} \SUPPORTEDANCH{\char"0236\char"033B} \SUPPORTEDANCH{\char"0075\char"033B} \MISSINGGLYPH{\char"AB52\char"033B} \SUPPORTEDANCH{\char"0079\char"033B} \SUPPORTEDFALL{\char"01B4\char"033B} \MISSINGGLYPH{\char"1EFF\char"033B}

\SUPPORTEDANCH{\char"0048\char"033C} \SUPPORTEDFALL{\char"2C67\char"033C} \SUPPORTEDANCH{\char"004B\char"033C} \SUPPORTEDANCH{\char"2C69\char"033C} \SUPPORTEDANCH{\char"004D\char"033C} \MISSINGGLYPH{\char"2C6E\char"033C} \SUPPORTEDANCH{\char"0050\char"033C} \SUPPORTEDFALL{\char"01A4\char"033C} \SUPPORTEDANCH{\char"0052\char"033C} \SUPPORTEDFALL{\char"2C64\char"033C} \SUPPORTEDANCH{\char"0054\char"033C} \SUPPORTEDFALL{\char"01AC\char"033C} \SUPPORTEDFALL{\char"01AE\char"033C} \SUPPORTEDANCH{\char"0059\char"033C} \SUPPORTEDFALL{\char"01B3\char"033C} \SUPPORTEDANCH{\char"0061\char"033C} \SUPPORTEDFALL{\char"1E9A\char"033C} \SUPPORTEDANCH{\char"0062\char"033C} \SUPPORTEDANCH{\char"0183\char"033C} \SUPPORTEDANCH{\char"0063\char"033C} \SUPPORTEDFALL{\char"0188\char"033C} \SUPPORTEDANCH{\char"0064\char"033C} \SUPPORTEDANCH{\char"018C\char"033C} \SUPPORTEDFALL{\char"0221\char"033C} \SUPPORTEDANCH{\char"0068\char"033C} \SUPPORTEDANCH{\char"2C68\char"033C} \SUPPORTEDANCH{\char"006B\char"033C} \SUPPORTEDANCH{\char"0199\char"033C} \SUPPORTEDANCH{\char"2C6A\char"033C} \SUPPORTEDANCH{\char"006C\char"033C} \SUPPORTEDFALL{\char"0234\char"033C} \SUPPORTEDANCH{\char"006E\char"033C} \SUPPORTEDANCH{\char"0235\char"033C} \SUPPORTEDANCH{\char"0070\char"033C} \SUPPORTEDFALL{\char"01A5\char"033C} \SUPPORTEDANCH{\char"0071\char"033C} \SUPPORTEDFALL{\char"02A0\char"033C} \SUPPORTEDANCH{\char"0074\char"033C} \SUPPORTEDFALL{\char"01AB\char"033C} \SUPPORTEDANCH{\char"01AD\char"033C} \SUPPORTEDANCH{\char"0236\char"033C} \SUPPORTEDANCH{\char"0075\char"033C} \MISSINGGLYPH{\char"AB52\char"033C} \SUPPORTEDANCH{\char"0079\char"033C} \SUPPORTEDFALL{\char"01B4\char"033C} \MISSINGGLYPH{\char"1EFF\char"033C}

\SUPPORTEDANCH{\char"0048\char"0347} \SUPPORTEDFALL{\char"2C67\char"0347} \SUPPORTEDANCH{\char"004B\char"0347} \SUPPORTEDANCH{\char"2C69\char"0347} \SUPPORTEDANCH{\char"004D\char"0347} \MISSINGGLYPH{\char"2C6E\char"0347} \SUPPORTEDANCH{\char"0050\char"0347} \SUPPORTEDFALL{\char"01A4\char"0347} \SUPPORTEDANCH{\char"0052\char"0347} \SUPPORTEDFALL{\char"2C64\char"0347} \SUPPORTEDANCH{\char"0054\char"0347} \SUPPORTEDFALL{\char"01AC\char"0347} \SUPPORTEDFALL{\char"01AE\char"0347} \SUPPORTEDANCH{\char"0059\char"0347} \SUPPORTEDFALL{\char"01B3\char"0347} \SUPPORTEDANCH{\char"0061\char"0347} \SUPPORTEDFALL{\char"1E9A\char"0347} \SUPPORTEDANCH{\char"0062\char"0347} \SUPPORTEDANCH{\char"0183\char"0347} \SUPPORTEDANCH{\char"0063\char"0347} \SUPPORTEDFALL{\char"0188\char"0347} \SUPPORTEDANCH{\char"0064\char"0347} \SUPPORTEDANCH{\char"018C\char"0347} \SUPPORTEDFALL{\char"0221\char"0347} \SUPPORTEDANCH{\char"0068\char"0347} \SUPPORTEDANCH{\char"2C68\char"0347} \SUPPORTEDANCH{\char"006B\char"0347} \SUPPORTEDANCH{\char"0199\char"0347} \SUPPORTEDANCH{\char"2C6A\char"0347} \SUPPORTEDANCH{\char"006C\char"0347} \SUPPORTEDFALL{\char"0234\char"0347} \SUPPORTEDANCH{\char"006E\char"0347} \SUPPORTEDANCH{\char"0235\char"0347} \SUPPORTEDANCH{\char"0070\char"0347} \SUPPORTEDFALL{\char"01A5\char"0347} \SUPPORTEDANCH{\char"0071\char"0347} \SUPPORTEDFALL{\char"02A0\char"0347} \SUPPORTEDANCH{\char"0074\char"0347} \SUPPORTEDFALL{\char"01AB\char"0347} \SUPPORTEDANCH{\char"01AD\char"0347} \SUPPORTEDANCH{\char"0236\char"0347} \SUPPORTEDANCH{\char"0075\char"0347} \MISSINGGLYPH{\char"AB52\char"0347} \SUPPORTEDANCH{\char"0079\char"0347} \SUPPORTEDFALL{\char"01B4\char"0347} \MISSINGGLYPH{\char"1EFF\char"0347}

\SUPPORTEDANCH{\char"0048\char"0348} \SUPPORTEDFALL{\char"2C67\char"0348} \SUPPORTEDANCH{\char"004B\char"0348} \SUPPORTEDANCH{\char"2C69\char"0348} \SUPPORTEDANCH{\char"004D\char"0348} \MISSINGGLYPH{\char"2C6E\char"0348} \SUPPORTEDANCH{\char"0050\char"0348} \SUPPORTEDFALL{\char"01A4\char"0348} \SUPPORTEDANCH{\char"0052\char"0348} \SUPPORTEDFALL{\char"2C64\char"0348} \SUPPORTEDANCH{\char"0054\char"0348} \SUPPORTEDFALL{\char"01AC\char"0348} \SUPPORTEDFALL{\char"01AE\char"0348} \SUPPORTEDANCH{\char"0059\char"0348} \SUPPORTEDFALL{\char"01B3\char"0348} \SUPPORTEDANCH{\char"0061\char"0348} \SUPPORTEDFALL{\char"1E9A\char"0348} \SUPPORTEDANCH{\char"0062\char"0348} \SUPPORTEDANCH{\char"0183\char"0348} \SUPPORTEDANCH{\char"0063\char"0348} \SUPPORTEDFALL{\char"0188\char"0348} \SUPPORTEDANCH{\char"0064\char"0348} \SUPPORTEDANCH{\char"018C\char"0348} \SUPPORTEDFALL{\char"0221\char"0348} \SUPPORTEDANCH{\char"0068\char"0348} \SUPPORTEDANCH{\char"2C68\char"0348} \SUPPORTEDANCH{\char"006B\char"0348} \SUPPORTEDANCH{\char"0199\char"0348} \SUPPORTEDANCH{\char"2C6A\char"0348} \SUPPORTEDANCH{\char"006C\char"0348} \SUPPORTEDFALL{\char"0234\char"0348} \SUPPORTEDANCH{\char"006E\char"0348} \SUPPORTEDANCH{\char"0235\char"0348} \SUPPORTEDANCH{\char"0070\char"0348} \SUPPORTEDFALL{\char"01A5\char"0348} \SUPPORTEDANCH{\char"0071\char"0348} \SUPPORTEDFALL{\char"02A0\char"0348} \SUPPORTEDANCH{\char"0074\char"0348} \SUPPORTEDFALL{\char"01AB\char"0348} \SUPPORTEDANCH{\char"01AD\char"0348} \SUPPORTEDANCH{\char"0236\char"0348} \SUPPORTEDANCH{\char"0075\char"0348} \MISSINGGLYPH{\char"AB52\char"0348} \SUPPORTEDANCH{\char"0079\char"0348} \SUPPORTEDFALL{\char"01B4\char"0348} \MISSINGGLYPH{\char"1EFF\char"0348}

\SUPPORTEDANCH{\char"0048\char"0349} \SUPPORTEDFALL{\char"2C67\char"0349} \SUPPORTEDANCH{\char"004B\char"0349} \SUPPORTEDANCH{\char"2C69\char"0349} \SUPPORTEDANCH{\char"004D\char"0349} \MISSINGGLYPH{\char"2C6E\char"0349} \SUPPORTEDANCH{\char"0050\char"0349} \SUPPORTEDFALL{\char"01A4\char"0349} \SUPPORTEDANCH{\char"0052\char"0349} \SUPPORTEDFALL{\char"2C64\char"0349} \SUPPORTEDANCH{\char"0054\char"0349} \SUPPORTEDFALL{\char"01AC\char"0349} \SUPPORTEDFALL{\char"01AE\char"0349} \SUPPORTEDANCH{\char"0059\char"0349} \SUPPORTEDFALL{\char"01B3\char"0349} \SUPPORTEDANCH{\char"0061\char"0349} \SUPPORTEDFALL{\char"1E9A\char"0349} \SUPPORTEDANCH{\char"0062\char"0349} \SUPPORTEDANCH{\char"0183\char"0349} \SUPPORTEDANCH{\char"0063\char"0349} \SUPPORTEDFALL{\char"0188\char"0349} \SUPPORTEDANCH{\char"0064\char"0349} \SUPPORTEDANCH{\char"018C\char"0349} \SUPPORTEDFALL{\char"0221\char"0349} \SUPPORTEDANCH{\char"0068\char"0349} \SUPPORTEDANCH{\char"2C68\char"0349} \SUPPORTEDANCH{\char"006B\char"0349} \SUPPORTEDANCH{\char"0199\char"0349} \SUPPORTEDANCH{\char"2C6A\char"0349} \SUPPORTEDANCH{\char"006C\char"0349} \SUPPORTEDFALL{\char"0234\char"0349} \SUPPORTEDANCH{\char"006E\char"0349} \SUPPORTEDANCH{\char"0235\char"0349} \SUPPORTEDANCH{\char"0070\char"0349} \SUPPORTEDFALL{\char"01A5\char"0349} \SUPPORTEDANCH{\char"0071\char"0349} \SUPPORTEDFALL{\char"02A0\char"0349} \SUPPORTEDANCH{\char"0074\char"0349} \SUPPORTEDFALL{\char"01AB\char"0349} \SUPPORTEDANCH{\char"01AD\char"0349} \SUPPORTEDANCH{\char"0236\char"0349} \SUPPORTEDANCH{\char"0075\char"0349} \MISSINGGLYPH{\char"AB52\char"0349} \SUPPORTEDANCH{\char"0079\char"0349} \SUPPORTEDFALL{\char"01B4\char"0349} \MISSINGGLYPH{\char"1EFF\char"0349}

\SUPPORTEDANCH{\char"0048\char"034D} \SUPPORTEDFALL{\char"2C67\char"034D} \SUPPORTEDANCH{\char"004B\char"034D} \SUPPORTEDANCH{\char"2C69\char"034D} \SUPPORTEDANCH{\char"004D\char"034D} \MISSINGGLYPH{\char"2C6E\char"034D} \SUPPORTEDANCH{\char"0050\char"034D} \SUPPORTEDFALL{\char"01A4\char"034D} \SUPPORTEDANCH{\char"0052\char"034D} \SUPPORTEDFALL{\char"2C64\char"034D} \SUPPORTEDANCH{\char"0054\char"034D} \SUPPORTEDFALL{\char"01AC\char"034D} \SUPPORTEDFALL{\char"01AE\char"034D} \SUPPORTEDANCH{\char"0059\char"034D} \SUPPORTEDFALL{\char"01B3\char"034D} \SUPPORTEDANCH{\char"0061\char"034D} \SUPPORTEDFALL{\char"1E9A\char"034D} \SUPPORTEDANCH{\char"0062\char"034D} \SUPPORTEDANCH{\char"0183\char"034D} \SUPPORTEDANCH{\char"0063\char"034D} \SUPPORTEDFALL{\char"0188\char"034D} \SUPPORTEDANCH{\char"0064\char"034D} \SUPPORTEDANCH{\char"018C\char"034D} \SUPPORTEDFALL{\char"0221\char"034D} \SUPPORTEDANCH{\char"0068\char"034D} \SUPPORTEDANCH{\char"2C68\char"034D} \SUPPORTEDANCH{\char"006B\char"034D} \SUPPORTEDANCH{\char"0199\char"034D} \SUPPORTEDANCH{\char"2C6A\char"034D} \SUPPORTEDANCH{\char"006C\char"034D} \SUPPORTEDFALL{\char"0234\char"034D} \SUPPORTEDANCH{\char"006E\char"034D} \SUPPORTEDANCH{\char"0235\char"034D} \SUPPORTEDANCH{\char"0070\char"034D} \SUPPORTEDFALL{\char"01A5\char"034D} \SUPPORTEDANCH{\char"0071\char"034D} \SUPPORTEDFALL{\char"02A0\char"034D} \SUPPORTEDANCH{\char"0074\char"034D} \SUPPORTEDFALL{\char"01AB\char"034D} \SUPPORTEDANCH{\char"01AD\char"034D} \SUPPORTEDANCH{\char"0236\char"034D} \SUPPORTEDANCH{\char"0075\char"034D} \MISSINGGLYPH{\char"AB52\char"034D} \SUPPORTEDANCH{\char"0079\char"034D} \SUPPORTEDFALL{\char"01B4\char"034D} \MISSINGGLYPH{\char"1EFF\char"034D}

\SUPPORTEDANCH{\char"0048\char"034E} \SUPPORTEDFALL{\char"2C67\char"034E} \SUPPORTEDANCH{\char"004B\char"034E} \SUPPORTEDANCH{\char"2C69\char"034E} \SUPPORTEDANCH{\char"004D\char"034E} \MISSINGGLYPH{\char"2C6E\char"034E} \SUPPORTEDANCH{\char"0050\char"034E} \SUPPORTEDFALL{\char"01A4\char"034E} \SUPPORTEDANCH{\char"0052\char"034E} \SUPPORTEDFALL{\char"2C64\char"034E} \SUPPORTEDANCH{\char"0054\char"034E} \SUPPORTEDFALL{\char"01AC\char"034E} \SUPPORTEDFALL{\char"01AE\char"034E} \SUPPORTEDANCH{\char"0059\char"034E} \SUPPORTEDFALL{\char"01B3\char"034E} \SUPPORTEDANCH{\char"0061\char"034E} \SUPPORTEDFALL{\char"1E9A\char"034E} \SUPPORTEDANCH{\char"0062\char"034E} \SUPPORTEDANCH{\char"0183\char"034E} \SUPPORTEDANCH{\char"0063\char"034E} \SUPPORTEDFALL{\char"0188\char"034E} \SUPPORTEDANCH{\char"0064\char"034E} \SUPPORTEDANCH{\char"018C\char"034E} \SUPPORTEDFALL{\char"0221\char"034E} \SUPPORTEDANCH{\char"0068\char"034E} \SUPPORTEDANCH{\char"2C68\char"034E} \SUPPORTEDANCH{\char"006B\char"034E} \SUPPORTEDANCH{\char"0199\char"034E} \SUPPORTEDANCH{\char"2C6A\char"034E} \SUPPORTEDANCH{\char"006C\char"034E} \SUPPORTEDFALL{\char"0234\char"034E} \SUPPORTEDANCH{\char"006E\char"034E} \SUPPORTEDANCH{\char"0235\char"034E} \SUPPORTEDANCH{\char"0070\char"034E} \SUPPORTEDFALL{\char"01A5\char"034E} \SUPPORTEDANCH{\char"0071\char"034E} \SUPPORTEDFALL{\char"02A0\char"034E} \SUPPORTEDANCH{\char"0074\char"034E} \SUPPORTEDFALL{\char"01AB\char"034E} \SUPPORTEDANCH{\char"01AD\char"034E} \SUPPORTEDANCH{\char"0236\char"034E} \SUPPORTEDANCH{\char"0075\char"034E} \MISSINGGLYPH{\char"AB52\char"034E} \SUPPORTEDANCH{\char"0079\char"034E} \SUPPORTEDFALL{\char"01B4\char"034E} \MISSINGGLYPH{\char"1EFF\char"034E}

}

\newpage
\subsection*{Italic patched, Group 4, below, Below}

{\LibertinusSerifPatch\itshape
% grid. columns are bases, rows are marks.
\SUPPORTEDANCH{\char"0048\char"0316} \SUPPORTEDANCH{\char"2C67\char"0316} \SUPPORTEDANCH{\char"004B\char"0316} \SUPPORTEDANCH{\char"2C69\char"0316} \SUPPORTEDANCH{\char"004D\char"0316} \MISSINGGLYPH{\char"2C6E\char"0316} \SUPPORTEDANCH{\char"0050\char"0316} \SUPPORTEDANCH{\char"01A4\char"0316} \SUPPORTEDANCH{\char"0052\char"0316} \SUPPORTEDANCH{\char"2C64\char"0316} \SUPPORTEDANCH{\char"0054\char"0316} \SUPPORTEDANCH{\char"01AC\char"0316} \SUPPORTEDANCH{\char"01AE\char"0316} \SUPPORTEDANCH{\char"0059\char"0316} \SUPPORTEDANCH{\char"01B3\char"0316} \SUPPORTEDANCH{\char"0061\char"0316} \SUPPORTEDANCH{\char"1E9A\char"0316} \SUPPORTEDANCH{\char"0062\char"0316} \SUPPORTEDANCH{\char"0183\char"0316} \SUPPORTEDANCH{\char"0063\char"0316} \SUPPORTEDANCH{\char"0188\char"0316} \SUPPORTEDANCH{\char"0064\char"0316} \SUPPORTEDANCH{\char"018C\char"0316} \SUPPORTEDANCH{\char"0221\char"0316} \SUPPORTEDANCH{\char"0068\char"0316} \SUPPORTEDANCH{\char"2C68\char"0316} \SUPPORTEDANCH{\char"006B\char"0316} \SUPPORTEDANCH{\char"0199\char"0316} \SUPPORTEDANCH{\char"2C6A\char"0316} \SUPPORTEDANCH{\char"006C\char"0316} \SUPPORTEDANCH{\char"0234\char"0316} \SUPPORTEDANCH{\char"006E\char"0316} \SUPPORTEDANCH{\char"0235\char"0316} \SUPPORTEDANCH{\char"0070\char"0316} \SUPPORTEDANCH{\char"01A5\char"0316} \SUPPORTEDANCH{\char"0071\char"0316} \SUPPORTEDANCH{\char"02A0\char"0316} \SUPPORTEDANCH{\char"0074\char"0316} \SUPPORTEDANCH{\char"01AB\char"0316} \SUPPORTEDANCH{\char"01AD\char"0316} \SUPPORTEDANCH{\char"0236\char"0316} \SUPPORTEDANCH{\char"0075\char"0316} \MISSINGGLYPH{\char"AB52\char"0316} \SUPPORTEDANCH{\char"0079\char"0316} \SUPPORTEDANCH{\char"01B4\char"0316} \MISSINGGLYPH{\char"1EFF\char"0316}

\SUPPORTEDANCH{\char"0048\char"0317} \SUPPORTEDANCH{\char"2C67\char"0317} \SUPPORTEDANCH{\char"004B\char"0317} \SUPPORTEDANCH{\char"2C69\char"0317} \SUPPORTEDANCH{\char"004D\char"0317} \MISSINGGLYPH{\char"2C6E\char"0317} \SUPPORTEDANCH{\char"0050\char"0317} \SUPPORTEDANCH{\char"01A4\char"0317} \SUPPORTEDANCH{\char"0052\char"0317} \SUPPORTEDANCH{\char"2C64\char"0317} \SUPPORTEDANCH{\char"0054\char"0317} \SUPPORTEDANCH{\char"01AC\char"0317} \SUPPORTEDANCH{\char"01AE\char"0317} \SUPPORTEDANCH{\char"0059\char"0317} \SUPPORTEDANCH{\char"01B3\char"0317} \SUPPORTEDANCH{\char"0061\char"0317} \SUPPORTEDANCH{\char"1E9A\char"0317} \SUPPORTEDANCH{\char"0062\char"0317} \SUPPORTEDANCH{\char"0183\char"0317} \SUPPORTEDANCH{\char"0063\char"0317} \SUPPORTEDANCH{\char"0188\char"0317} \SUPPORTEDANCH{\char"0064\char"0317} \SUPPORTEDANCH{\char"018C\char"0317} \SUPPORTEDANCH{\char"0221\char"0317} \SUPPORTEDANCH{\char"0068\char"0317} \SUPPORTEDANCH{\char"2C68\char"0317} \SUPPORTEDANCH{\char"006B\char"0317} \SUPPORTEDANCH{\char"0199\char"0317} \SUPPORTEDANCH{\char"2C6A\char"0317} \SUPPORTEDANCH{\char"006C\char"0317} \SUPPORTEDANCH{\char"0234\char"0317} \SUPPORTEDANCH{\char"006E\char"0317} \SUPPORTEDANCH{\char"0235\char"0317} \SUPPORTEDANCH{\char"0070\char"0317} \SUPPORTEDANCH{\char"01A5\char"0317} \SUPPORTEDANCH{\char"0071\char"0317} \SUPPORTEDANCH{\char"02A0\char"0317} \SUPPORTEDANCH{\char"0074\char"0317} \SUPPORTEDANCH{\char"01AB\char"0317} \SUPPORTEDANCH{\char"01AD\char"0317} \SUPPORTEDANCH{\char"0236\char"0317} \SUPPORTEDANCH{\char"0075\char"0317} \MISSINGGLYPH{\char"AB52\char"0317} \SUPPORTEDANCH{\char"0079\char"0317} \SUPPORTEDANCH{\char"01B4\char"0317} \MISSINGGLYPH{\char"1EFF\char"0317}

\SUPPORTEDANCH{\char"0048\char"0318} \SUPPORTEDANCH{\char"2C67\char"0318} \SUPPORTEDANCH{\char"004B\char"0318} \SUPPORTEDANCH{\char"2C69\char"0318} \SUPPORTEDANCH{\char"004D\char"0318} \MISSINGGLYPH{\char"2C6E\char"0318} \SUPPORTEDANCH{\char"0050\char"0318} \SUPPORTEDANCH{\char"01A4\char"0318} \SUPPORTEDANCH{\char"0052\char"0318} \SUPPORTEDANCH{\char"2C64\char"0318} \SUPPORTEDANCH{\char"0054\char"0318} \SUPPORTEDANCH{\char"01AC\char"0318} \SUPPORTEDANCH{\char"01AE\char"0318} \SUPPORTEDANCH{\char"0059\char"0318} \SUPPORTEDANCH{\char"01B3\char"0318} \SUPPORTEDANCH{\char"0061\char"0318} \SUPPORTEDANCH{\char"1E9A\char"0318} \SUPPORTEDANCH{\char"0062\char"0318} \SUPPORTEDANCH{\char"0183\char"0318} \SUPPORTEDANCH{\char"0063\char"0318} \SUPPORTEDANCH{\char"0188\char"0318} \SUPPORTEDANCH{\char"0064\char"0318} \SUPPORTEDANCH{\char"018C\char"0318} \SUPPORTEDANCH{\char"0221\char"0318} \SUPPORTEDANCH{\char"0068\char"0318} \SUPPORTEDANCH{\char"2C68\char"0318} \SUPPORTEDANCH{\char"006B\char"0318} \SUPPORTEDANCH{\char"0199\char"0318} \SUPPORTEDANCH{\char"2C6A\char"0318} \SUPPORTEDANCH{\char"006C\char"0318} \SUPPORTEDANCH{\char"0234\char"0318} \SUPPORTEDANCH{\char"006E\char"0318} \SUPPORTEDANCH{\char"0235\char"0318} \SUPPORTEDANCH{\char"0070\char"0318} \SUPPORTEDANCH{\char"01A5\char"0318} \SUPPORTEDANCH{\char"0071\char"0318} \SUPPORTEDANCH{\char"02A0\char"0318} \SUPPORTEDANCH{\char"0074\char"0318} \SUPPORTEDANCH{\char"01AB\char"0318} \SUPPORTEDANCH{\char"01AD\char"0318} \SUPPORTEDANCH{\char"0236\char"0318} \SUPPORTEDANCH{\char"0075\char"0318} \MISSINGGLYPH{\char"AB52\char"0318} \SUPPORTEDANCH{\char"0079\char"0318} \SUPPORTEDANCH{\char"01B4\char"0318} \MISSINGGLYPH{\char"1EFF\char"0318}

\SUPPORTEDANCH{\char"0048\char"0319} \SUPPORTEDANCH{\char"2C67\char"0319} \SUPPORTEDANCH{\char"004B\char"0319} \SUPPORTEDANCH{\char"2C69\char"0319} \SUPPORTEDANCH{\char"004D\char"0319} \MISSINGGLYPH{\char"2C6E\char"0319} \SUPPORTEDANCH{\char"0050\char"0319} \SUPPORTEDANCH{\char"01A4\char"0319} \SUPPORTEDANCH{\char"0052\char"0319} \SUPPORTEDANCH{\char"2C64\char"0319} \SUPPORTEDANCH{\char"0054\char"0319} \SUPPORTEDANCH{\char"01AC\char"0319} \SUPPORTEDANCH{\char"01AE\char"0319} \SUPPORTEDANCH{\char"0059\char"0319} \SUPPORTEDANCH{\char"01B3\char"0319} \SUPPORTEDANCH{\char"0061\char"0319} \SUPPORTEDANCH{\char"1E9A\char"0319} \SUPPORTEDANCH{\char"0062\char"0319} \SUPPORTEDANCH{\char"0183\char"0319} \SUPPORTEDANCH{\char"0063\char"0319} \SUPPORTEDANCH{\char"0188\char"0319} \SUPPORTEDANCH{\char"0064\char"0319} \SUPPORTEDANCH{\char"018C\char"0319} \SUPPORTEDANCH{\char"0221\char"0319} \SUPPORTEDANCH{\char"0068\char"0319} \SUPPORTEDANCH{\char"2C68\char"0319} \SUPPORTEDANCH{\char"006B\char"0319} \SUPPORTEDANCH{\char"0199\char"0319} \SUPPORTEDANCH{\char"2C6A\char"0319} \SUPPORTEDANCH{\char"006C\char"0319} \SUPPORTEDANCH{\char"0234\char"0319} \SUPPORTEDANCH{\char"006E\char"0319} \SUPPORTEDANCH{\char"0235\char"0319} \SUPPORTEDANCH{\char"0070\char"0319} \SUPPORTEDANCH{\char"01A5\char"0319} \SUPPORTEDANCH{\char"0071\char"0319} \SUPPORTEDANCH{\char"02A0\char"0319} \SUPPORTEDANCH{\char"0074\char"0319} \SUPPORTEDANCH{\char"01AB\char"0319} \SUPPORTEDANCH{\char"01AD\char"0319} \SUPPORTEDANCH{\char"0236\char"0319} \SUPPORTEDANCH{\char"0075\char"0319} \MISSINGGLYPH{\char"AB52\char"0319} \SUPPORTEDANCH{\char"0079\char"0319} \SUPPORTEDANCH{\char"01B4\char"0319} \MISSINGGLYPH{\char"1EFF\char"0319}

\SUPPORTEDANCH{\char"0048\char"031C} \SUPPORTEDANCH{\char"2C67\char"031C} \SUPPORTEDANCH{\char"004B\char"031C} \SUPPORTEDANCH{\char"2C69\char"031C} \SUPPORTEDANCH{\char"004D\char"031C} \MISSINGGLYPH{\char"2C6E\char"031C} \SUPPORTEDANCH{\char"0050\char"031C} \SUPPORTEDANCH{\char"01A4\char"031C} \SUPPORTEDANCH{\char"0052\char"031C} \SUPPORTEDANCH{\char"2C64\char"031C} \SUPPORTEDANCH{\char"0054\char"031C} \SUPPORTEDANCH{\char"01AC\char"031C} \SUPPORTEDANCH{\char"01AE\char"031C} \SUPPORTEDANCH{\char"0059\char"031C} \SUPPORTEDANCH{\char"01B3\char"031C} \SUPPORTEDANCH{\char"0061\char"031C} \SUPPORTEDANCH{\char"1E9A\char"031C} \SUPPORTEDANCH{\char"0062\char"031C} \SUPPORTEDANCH{\char"0183\char"031C} \SUPPORTEDANCH{\char"0063\char"031C} \SUPPORTEDANCH{\char"0188\char"031C} \SUPPORTEDANCH{\char"0064\char"031C} \SUPPORTEDANCH{\char"018C\char"031C} \SUPPORTEDANCH{\char"0221\char"031C} \SUPPORTEDANCH{\char"0068\char"031C} \SUPPORTEDANCH{\char"2C68\char"031C} \SUPPORTEDANCH{\char"006B\char"031C} \SUPPORTEDANCH{\char"0199\char"031C} \SUPPORTEDANCH{\char"2C6A\char"031C} \SUPPORTEDANCH{\char"006C\char"031C} \SUPPORTEDANCH{\char"0234\char"031C} \SUPPORTEDANCH{\char"006E\char"031C} \SUPPORTEDANCH{\char"0235\char"031C} \SUPPORTEDANCH{\char"0070\char"031C} \SUPPORTEDANCH{\char"01A5\char"031C} \SUPPORTEDANCH{\char"0071\char"031C} \SUPPORTEDANCH{\char"02A0\char"031C} \SUPPORTEDANCH{\char"0074\char"031C} \SUPPORTEDANCH{\char"01AB\char"031C} \SUPPORTEDANCH{\char"01AD\char"031C} \SUPPORTEDANCH{\char"0236\char"031C} \SUPPORTEDANCH{\char"0075\char"031C} \MISSINGGLYPH{\char"AB52\char"031C} \SUPPORTEDANCH{\char"0079\char"031C} \SUPPORTEDANCH{\char"01B4\char"031C} \MISSINGGLYPH{\char"1EFF\char"031C}

\SUPPORTEDANCH{\char"0048\char"0320} \SUPPORTEDANCH{\char"2C67\char"0320} \SUPPORTEDANCH{\char"004B\char"0320} \SUPPORTEDANCH{\char"2C69\char"0320} \SUPPORTEDANCH{\char"004D\char"0320} \MISSINGGLYPH{\char"2C6E\char"0320} \SUPPORTEDANCH{\char"0050\char"0320} \SUPPORTEDANCH{\char"01A4\char"0320} \SUPPORTEDANCH{\char"0052\char"0320} \SUPPORTEDANCH{\char"2C64\char"0320} \SUPPORTEDANCH{\char"0054\char"0320} \SUPPORTEDANCH{\char"01AC\char"0320} \SUPPORTEDANCH{\char"01AE\char"0320} \SUPPORTEDANCH{\char"0059\char"0320} \SUPPORTEDANCH{\char"01B3\char"0320} \SUPPORTEDANCH{\char"0061\char"0320} \SUPPORTEDANCH{\char"1E9A\char"0320} \SUPPORTEDANCH{\char"0062\char"0320} \SUPPORTEDANCH{\char"0183\char"0320} \SUPPORTEDANCH{\char"0063\char"0320} \SUPPORTEDANCH{\char"0188\char"0320} \SUPPORTEDANCH{\char"0064\char"0320} \SUPPORTEDANCH{\char"018C\char"0320} \SUPPORTEDANCH{\char"0221\char"0320} \SUPPORTEDANCH{\char"0068\char"0320} \SUPPORTEDANCH{\char"2C68\char"0320} \SUPPORTEDANCH{\char"006B\char"0320} \SUPPORTEDANCH{\char"0199\char"0320} \SUPPORTEDANCH{\char"2C6A\char"0320} \SUPPORTEDANCH{\char"006C\char"0320} \SUPPORTEDANCH{\char"0234\char"0320} \SUPPORTEDANCH{\char"006E\char"0320} \SUPPORTEDANCH{\char"0235\char"0320} \SUPPORTEDANCH{\char"0070\char"0320} \SUPPORTEDANCH{\char"01A5\char"0320} \SUPPORTEDANCH{\char"0071\char"0320} \SUPPORTEDANCH{\char"02A0\char"0320} \SUPPORTEDANCH{\char"0074\char"0320} \SUPPORTEDANCH{\char"01AB\char"0320} \SUPPORTEDANCH{\char"01AD\char"0320} \SUPPORTEDANCH{\char"0236\char"0320} \SUPPORTEDANCH{\char"0075\char"0320} \MISSINGGLYPH{\char"AB52\char"0320} \SUPPORTEDANCH{\char"0079\char"0320} \SUPPORTEDANCH{\char"01B4\char"0320} \MISSINGGLYPH{\char"1EFF\char"0320}

\SUPPORTEDPREC{\char"0048\char"0323} \SUPPORTEDANCH{\char"2C67\char"0323} \SUPPORTEDPREC{\char"004B\char"0323} \SUPPORTEDANCH{\char"2C69\char"0323} \SUPPORTEDPREC{\char"004D\char"0323} \MISSINGGLYPH{\char"2C6E\char"0323} \SUPPORTEDANCH{\char"0050\char"0323} \SUPPORTEDANCH{\char"01A4\char"0323} \SUPPORTEDPREC{\char"0052\char"0323} \SUPPORTEDANCH{\char"2C64\char"0323} \SUPPORTEDPREC{\char"0054\char"0323} \SUPPORTEDANCH{\char"01AC\char"0323} \SUPPORTEDANCH{\char"01AE\char"0323} \SUPPORTEDPREC{\char"0059\char"0323} \SUPPORTEDANCH{\char"01B3\char"0323} \SUPPORTEDPREC{\char"0061\char"0323} \SUPPORTEDANCH{\char"1E9A\char"0323} \SUPPORTEDPREC{\char"0062\char"0323} \SUPPORTEDANCH{\char"0183\char"0323} \SUPPORTEDANCH{\char"0063\char"0323} \SUPPORTEDANCH{\char"0188\char"0323} \SUPPORTEDPREC{\char"0064\char"0323} \SUPPORTEDANCH{\char"018C\char"0323} \SUPPORTEDANCH{\char"0221\char"0323} \SUPPORTEDPREC{\char"0068\char"0323} \SUPPORTEDANCH{\char"2C68\char"0323} \SUPPORTEDPREC{\char"006B\char"0323} \SUPPORTEDANCH{\char"0199\char"0323} \SUPPORTEDANCH{\char"2C6A\char"0323} \SUPPORTEDPREC{\char"006C\char"0323} \SUPPORTEDANCH{\char"0234\char"0323} \SUPPORTEDPREC{\char"006E\char"0323} \SUPPORTEDANCH{\char"0235\char"0323} \SUPPORTEDANCH{\char"0070\char"0323} \SUPPORTEDANCH{\char"01A5\char"0323} \SUPPORTEDANCH{\char"0071\char"0323} \SUPPORTEDANCH{\char"02A0\char"0323} \SUPPORTEDPREC{\char"0074\char"0323} \SUPPORTEDANCH{\char"01AB\char"0323} \SUPPORTEDANCH{\char"01AD\char"0323} \SUPPORTEDANCH{\char"0236\char"0323} \SUPPORTEDPREC{\char"0075\char"0323} \MISSINGGLYPH{\char"AB52\char"0323} \SUPPORTEDPREC{\char"0079\char"0323} \SUPPORTEDANCH{\char"01B4\char"0323} \MISSINGGLYPH{\char"1EFF\char"0323}

\SUPPORTEDANCH{\char"0048\char"0324} \SUPPORTEDANCH{\char"2C67\char"0324} \SUPPORTEDANCH{\char"004B\char"0324} \SUPPORTEDANCH{\char"2C69\char"0324} \SUPPORTEDANCH{\char"004D\char"0324} \MISSINGGLYPH{\char"2C6E\char"0324} \SUPPORTEDANCH{\char"0050\char"0324} \SUPPORTEDANCH{\char"01A4\char"0324} \SUPPORTEDANCH{\char"0052\char"0324} \SUPPORTEDANCH{\char"2C64\char"0324} \SUPPORTEDANCH{\char"0054\char"0324} \SUPPORTEDANCH{\char"01AC\char"0324} \SUPPORTEDANCH{\char"01AE\char"0324} \SUPPORTEDANCH{\char"0059\char"0324} \SUPPORTEDANCH{\char"01B3\char"0324} \SUPPORTEDANCH{\char"0061\char"0324} \SUPPORTEDANCH{\char"1E9A\char"0324} \SUPPORTEDANCH{\char"0062\char"0324} \SUPPORTEDANCH{\char"0183\char"0324} \SUPPORTEDANCH{\char"0063\char"0324} \SUPPORTEDANCH{\char"0188\char"0324} \SUPPORTEDANCH{\char"0064\char"0324} \SUPPORTEDANCH{\char"018C\char"0324} \SUPPORTEDANCH{\char"0221\char"0324} \SUPPORTEDANCH{\char"0068\char"0324} \SUPPORTEDANCH{\char"2C68\char"0324} \SUPPORTEDANCH{\char"006B\char"0324} \SUPPORTEDANCH{\char"0199\char"0324} \SUPPORTEDANCH{\char"2C6A\char"0324} \SUPPORTEDANCH{\char"006C\char"0324} \SUPPORTEDANCH{\char"0234\char"0324} \SUPPORTEDANCH{\char"006E\char"0324} \SUPPORTEDANCH{\char"0235\char"0324} \SUPPORTEDANCH{\char"0070\char"0324} \SUPPORTEDANCH{\char"01A5\char"0324} \SUPPORTEDANCH{\char"0071\char"0324} \SUPPORTEDANCH{\char"02A0\char"0324} \SUPPORTEDANCH{\char"0074\char"0324} \SUPPORTEDANCH{\char"01AB\char"0324} \SUPPORTEDANCH{\char"01AD\char"0324} \SUPPORTEDANCH{\char"0236\char"0324} \SUPPORTEDPREC{\char"0075\char"0324} \MISSINGGLYPH{\char"AB52\char"0324} \SUPPORTEDANCH{\char"0079\char"0324} \SUPPORTEDANCH{\char"01B4\char"0324} \MISSINGGLYPH{\char"1EFF\char"0324}

\SUPPORTEDANCH{\char"0048\char"0325} \SUPPORTEDANCH{\char"2C67\char"0325} \SUPPORTEDANCH{\char"004B\char"0325} \SUPPORTEDANCH{\char"2C69\char"0325} \SUPPORTEDANCH{\char"004D\char"0325} \MISSINGGLYPH{\char"2C6E\char"0325} \SUPPORTEDANCH{\char"0050\char"0325} \SUPPORTEDANCH{\char"01A4\char"0325} \SUPPORTEDANCH{\char"0052\char"0325} \SUPPORTEDANCH{\char"2C64\char"0325} \SUPPORTEDANCH{\char"0054\char"0325} \SUPPORTEDANCH{\char"01AC\char"0325} \SUPPORTEDANCH{\char"01AE\char"0325} \SUPPORTEDANCH{\char"0059\char"0325} \SUPPORTEDANCH{\char"01B3\char"0325} \SUPPORTEDPREC{\char"0061\char"0325} \SUPPORTEDANCH{\char"1E9A\char"0325} \SUPPORTEDANCH{\char"0062\char"0325} \SUPPORTEDANCH{\char"0183\char"0325} \SUPPORTEDANCH{\char"0063\char"0325} \SUPPORTEDANCH{\char"0188\char"0325} \SUPPORTEDANCH{\char"0064\char"0325} \SUPPORTEDANCH{\char"018C\char"0325} \SUPPORTEDANCH{\char"0221\char"0325} \SUPPORTEDANCH{\char"0068\char"0325} \SUPPORTEDANCH{\char"2C68\char"0325} \SUPPORTEDANCH{\char"006B\char"0325} \SUPPORTEDANCH{\char"0199\char"0325} \SUPPORTEDANCH{\char"2C6A\char"0325} \SUPPORTEDANCH{\char"006C\char"0325} \SUPPORTEDANCH{\char"0234\char"0325} \SUPPORTEDANCH{\char"006E\char"0325} \SUPPORTEDANCH{\char"0235\char"0325} \SUPPORTEDANCH{\char"0070\char"0325} \SUPPORTEDANCH{\char"01A5\char"0325} \SUPPORTEDANCH{\char"0071\char"0325} \SUPPORTEDANCH{\char"02A0\char"0325} \SUPPORTEDANCH{\char"0074\char"0325} \SUPPORTEDANCH{\char"01AB\char"0325} \SUPPORTEDANCH{\char"01AD\char"0325} \SUPPORTEDANCH{\char"0236\char"0325} \SUPPORTEDANCH{\char"0075\char"0325} \MISSINGGLYPH{\char"AB52\char"0325} \SUPPORTEDANCH{\char"0079\char"0325} \SUPPORTEDANCH{\char"01B4\char"0325} \MISSINGGLYPH{\char"1EFF\char"0325}

\SUPPORTEDANCH{\char"0048\char"0326} \SUPPORTEDANCH{\char"2C67\char"0326} \SUPPORTEDANCH{\char"004B\char"0326} \SUPPORTEDANCH{\char"2C69\char"0326} \SUPPORTEDANCH{\char"004D\char"0326} \MISSINGGLYPH{\char"2C6E\char"0326} \SUPPORTEDANCH{\char"0050\char"0326} \SUPPORTEDANCH{\char"01A4\char"0326} \SUPPORTEDANCH{\char"0052\char"0326} \SUPPORTEDANCH{\char"2C64\char"0326} \SUPPORTEDPREC{\char"0054\char"0326} \SUPPORTEDANCH{\char"01AC\char"0326} \SUPPORTEDANCH{\char"01AE\char"0326} \SUPPORTEDANCH{\char"0059\char"0326} \SUPPORTEDANCH{\char"01B3\char"0326} \SUPPORTEDANCH{\char"0061\char"0326} \SUPPORTEDANCH{\char"1E9A\char"0326} \SUPPORTEDANCH{\char"0062\char"0326} \SUPPORTEDANCH{\char"0183\char"0326} \SUPPORTEDANCH{\char"0063\char"0326} \SUPPORTEDANCH{\char"0188\char"0326} \SUPPORTEDANCH{\char"0064\char"0326} \SUPPORTEDANCH{\char"018C\char"0326} \SUPPORTEDANCH{\char"0221\char"0326} \SUPPORTEDANCH{\char"0068\char"0326} \SUPPORTEDANCH{\char"2C68\char"0326} \SUPPORTEDANCH{\char"006B\char"0326} \SUPPORTEDANCH{\char"0199\char"0326} \SUPPORTEDANCH{\char"2C6A\char"0326} \SUPPORTEDANCH{\char"006C\char"0326} \SUPPORTEDANCH{\char"0234\char"0326} \SUPPORTEDANCH{\char"006E\char"0326} \SUPPORTEDANCH{\char"0235\char"0326} \SUPPORTEDANCH{\char"0070\char"0326} \SUPPORTEDANCH{\char"01A5\char"0326} \SUPPORTEDANCH{\char"0071\char"0326} \SUPPORTEDANCH{\char"02A0\char"0326} \SUPPORTEDPREC{\char"0074\char"0326} \SUPPORTEDANCH{\char"01AB\char"0326} \SUPPORTEDANCH{\char"01AD\char"0326} \SUPPORTEDANCH{\char"0236\char"0326} \SUPPORTEDANCH{\char"0075\char"0326} \MISSINGGLYPH{\char"AB52\char"0326} \SUPPORTEDANCH{\char"0079\char"0326} \SUPPORTEDANCH{\char"01B4\char"0326} \MISSINGGLYPH{\char"1EFF\char"0326}

\SUPPORTEDANCH{\char"0048\char"0329} \SUPPORTEDANCH{\char"2C67\char"0329} \SUPPORTEDANCH{\char"004B\char"0329} \SUPPORTEDANCH{\char"2C69\char"0329} \SUPPORTEDANCH{\char"004D\char"0329} \MISSINGGLYPH{\char"2C6E\char"0329} \SUPPORTEDANCH{\char"0050\char"0329} \SUPPORTEDANCH{\char"01A4\char"0329} \SUPPORTEDANCH{\char"0052\char"0329} \SUPPORTEDANCH{\char"2C64\char"0329} \SUPPORTEDANCH{\char"0054\char"0329} \SUPPORTEDANCH{\char"01AC\char"0329} \SUPPORTEDANCH{\char"01AE\char"0329} \SUPPORTEDANCH{\char"0059\char"0329} \SUPPORTEDANCH{\char"01B3\char"0329} \SUPPORTEDANCH{\char"0061\char"0329} \SUPPORTEDANCH{\char"1E9A\char"0329} \SUPPORTEDANCH{\char"0062\char"0329} \SUPPORTEDANCH{\char"0183\char"0329} \SUPPORTEDANCH{\char"0063\char"0329} \SUPPORTEDANCH{\char"0188\char"0329} \SUPPORTEDANCH{\char"0064\char"0329} \SUPPORTEDANCH{\char"018C\char"0329} \SUPPORTEDANCH{\char"0221\char"0329} \SUPPORTEDANCH{\char"0068\char"0329} \SUPPORTEDANCH{\char"2C68\char"0329} \SUPPORTEDANCH{\char"006B\char"0329} \SUPPORTEDANCH{\char"0199\char"0329} \SUPPORTEDANCH{\char"2C6A\char"0329} \SUPPORTEDANCH{\char"006C\char"0329} \SUPPORTEDANCH{\char"0234\char"0329} \SUPPORTEDANCH{\char"006E\char"0329} \SUPPORTEDANCH{\char"0235\char"0329} \SUPPORTEDANCH{\char"0070\char"0329} \SUPPORTEDANCH{\char"01A5\char"0329} \SUPPORTEDANCH{\char"0071\char"0329} \SUPPORTEDANCH{\char"02A0\char"0329} \SUPPORTEDANCH{\char"0074\char"0329} \SUPPORTEDANCH{\char"01AB\char"0329} \SUPPORTEDANCH{\char"01AD\char"0329} \SUPPORTEDANCH{\char"0236\char"0329} \SUPPORTEDANCH{\char"0075\char"0329} \MISSINGGLYPH{\char"AB52\char"0329} \SUPPORTEDANCH{\char"0079\char"0329} \SUPPORTEDANCH{\char"01B4\char"0329} \MISSINGGLYPH{\char"1EFF\char"0329}

\SUPPORTEDANCH{\char"0048\char"032A} \SUPPORTEDANCH{\char"2C67\char"032A} \SUPPORTEDANCH{\char"004B\char"032A} \SUPPORTEDANCH{\char"2C69\char"032A} \SUPPORTEDANCH{\char"004D\char"032A} \MISSINGGLYPH{\char"2C6E\char"032A} \SUPPORTEDANCH{\char"0050\char"032A} \SUPPORTEDANCH{\char"01A4\char"032A} \SUPPORTEDANCH{\char"0052\char"032A} \SUPPORTEDANCH{\char"2C64\char"032A} \SUPPORTEDANCH{\char"0054\char"032A} \SUPPORTEDANCH{\char"01AC\char"032A} \SUPPORTEDANCH{\char"01AE\char"032A} \SUPPORTEDANCH{\char"0059\char"032A} \SUPPORTEDANCH{\char"01B3\char"032A} \SUPPORTEDANCH{\char"0061\char"032A} \SUPPORTEDANCH{\char"1E9A\char"032A} \SUPPORTEDANCH{\char"0062\char"032A} \SUPPORTEDANCH{\char"0183\char"032A} \SUPPORTEDANCH{\char"0063\char"032A} \SUPPORTEDANCH{\char"0188\char"032A} \SUPPORTEDANCH{\char"0064\char"032A} \SUPPORTEDANCH{\char"018C\char"032A} \SUPPORTEDANCH{\char"0221\char"032A} \SUPPORTEDANCH{\char"0068\char"032A} \SUPPORTEDANCH{\char"2C68\char"032A} \SUPPORTEDANCH{\char"006B\char"032A} \SUPPORTEDANCH{\char"0199\char"032A} \SUPPORTEDANCH{\char"2C6A\char"032A} \SUPPORTEDANCH{\char"006C\char"032A} \SUPPORTEDANCH{\char"0234\char"032A} \SUPPORTEDANCH{\char"006E\char"032A} \SUPPORTEDANCH{\char"0235\char"032A} \SUPPORTEDANCH{\char"0070\char"032A} \SUPPORTEDANCH{\char"01A5\char"032A} \SUPPORTEDANCH{\char"0071\char"032A} \SUPPORTEDANCH{\char"02A0\char"032A} \SUPPORTEDANCH{\char"0074\char"032A} \SUPPORTEDANCH{\char"01AB\char"032A} \SUPPORTEDANCH{\char"01AD\char"032A} \SUPPORTEDANCH{\char"0236\char"032A} \SUPPORTEDANCH{\char"0075\char"032A} \MISSINGGLYPH{\char"AB52\char"032A} \SUPPORTEDANCH{\char"0079\char"032A} \SUPPORTEDANCH{\char"01B4\char"032A} \MISSINGGLYPH{\char"1EFF\char"032A}

\SUPPORTEDANCH{\char"0048\char"032B} \SUPPORTEDANCH{\char"2C67\char"032B} \SUPPORTEDANCH{\char"004B\char"032B} \SUPPORTEDANCH{\char"2C69\char"032B} \SUPPORTEDANCH{\char"004D\char"032B} \MISSINGGLYPH{\char"2C6E\char"032B} \SUPPORTEDANCH{\char"0050\char"032B} \SUPPORTEDANCH{\char"01A4\char"032B} \SUPPORTEDANCH{\char"0052\char"032B} \SUPPORTEDANCH{\char"2C64\char"032B} \SUPPORTEDANCH{\char"0054\char"032B} \SUPPORTEDANCH{\char"01AC\char"032B} \SUPPORTEDANCH{\char"01AE\char"032B} \SUPPORTEDANCH{\char"0059\char"032B} \SUPPORTEDANCH{\char"01B3\char"032B} \SUPPORTEDANCH{\char"0061\char"032B} \SUPPORTEDANCH{\char"1E9A\char"032B} \SUPPORTEDANCH{\char"0062\char"032B} \SUPPORTEDANCH{\char"0183\char"032B} \SUPPORTEDANCH{\char"0063\char"032B} \SUPPORTEDANCH{\char"0188\char"032B} \SUPPORTEDANCH{\char"0064\char"032B} \SUPPORTEDANCH{\char"018C\char"032B} \SUPPORTEDANCH{\char"0221\char"032B} \SUPPORTEDANCH{\char"0068\char"032B} \SUPPORTEDANCH{\char"2C68\char"032B} \SUPPORTEDANCH{\char"006B\char"032B} \SUPPORTEDANCH{\char"0199\char"032B} \SUPPORTEDANCH{\char"2C6A\char"032B} \SUPPORTEDANCH{\char"006C\char"032B} \SUPPORTEDANCH{\char"0234\char"032B} \SUPPORTEDANCH{\char"006E\char"032B} \SUPPORTEDANCH{\char"0235\char"032B} \SUPPORTEDANCH{\char"0070\char"032B} \SUPPORTEDANCH{\char"01A5\char"032B} \SUPPORTEDANCH{\char"0071\char"032B} \SUPPORTEDANCH{\char"02A0\char"032B} \SUPPORTEDANCH{\char"0074\char"032B} \SUPPORTEDANCH{\char"01AB\char"032B} \SUPPORTEDANCH{\char"01AD\char"032B} \SUPPORTEDANCH{\char"0236\char"032B} \SUPPORTEDANCH{\char"0075\char"032B} \MISSINGGLYPH{\char"AB52\char"032B} \SUPPORTEDANCH{\char"0079\char"032B} \SUPPORTEDANCH{\char"01B4\char"032B} \MISSINGGLYPH{\char"1EFF\char"032B}

\SUPPORTEDANCH{\char"0048\char"032C} \SUPPORTEDANCH{\char"2C67\char"032C} \SUPPORTEDANCH{\char"004B\char"032C} \SUPPORTEDANCH{\char"2C69\char"032C} \SUPPORTEDANCH{\char"004D\char"032C} \MISSINGGLYPH{\char"2C6E\char"032C} \SUPPORTEDANCH{\char"0050\char"032C} \SUPPORTEDANCH{\char"01A4\char"032C} \SUPPORTEDANCH{\char"0052\char"032C} \SUPPORTEDANCH{\char"2C64\char"032C} \SUPPORTEDANCH{\char"0054\char"032C} \SUPPORTEDANCH{\char"01AC\char"032C} \SUPPORTEDANCH{\char"01AE\char"032C} \SUPPORTEDANCH{\char"0059\char"032C} \SUPPORTEDANCH{\char"01B3\char"032C} \SUPPORTEDANCH{\char"0061\char"032C} \SUPPORTEDANCH{\char"1E9A\char"032C} \SUPPORTEDANCH{\char"0062\char"032C} \SUPPORTEDANCH{\char"0183\char"032C} \SUPPORTEDANCH{\char"0063\char"032C} \SUPPORTEDANCH{\char"0188\char"032C} \SUPPORTEDANCH{\char"0064\char"032C} \SUPPORTEDANCH{\char"018C\char"032C} \SUPPORTEDANCH{\char"0221\char"032C} \SUPPORTEDANCH{\char"0068\char"032C} \SUPPORTEDANCH{\char"2C68\char"032C} \SUPPORTEDANCH{\char"006B\char"032C} \SUPPORTEDANCH{\char"0199\char"032C} \SUPPORTEDANCH{\char"2C6A\char"032C} \SUPPORTEDANCH{\char"006C\char"032C} \SUPPORTEDANCH{\char"0234\char"032C} \SUPPORTEDANCH{\char"006E\char"032C} \SUPPORTEDANCH{\char"0235\char"032C} \SUPPORTEDANCH{\char"0070\char"032C} \SUPPORTEDANCH{\char"01A5\char"032C} \SUPPORTEDANCH{\char"0071\char"032C} \SUPPORTEDANCH{\char"02A0\char"032C} \SUPPORTEDANCH{\char"0074\char"032C} \SUPPORTEDANCH{\char"01AB\char"032C} \SUPPORTEDANCH{\char"01AD\char"032C} \SUPPORTEDANCH{\char"0236\char"032C} \SUPPORTEDANCH{\char"0075\char"032C} \MISSINGGLYPH{\char"AB52\char"032C} \SUPPORTEDANCH{\char"0079\char"032C} \SUPPORTEDANCH{\char"01B4\char"032C} \MISSINGGLYPH{\char"1EFF\char"032C}

\SUPPORTEDANCH{\char"0048\char"032D} \SUPPORTEDANCH{\char"2C67\char"032D} \SUPPORTEDANCH{\char"004B\char"032D} \SUPPORTEDANCH{\char"2C69\char"032D} \SUPPORTEDANCH{\char"004D\char"032D} \MISSINGGLYPH{\char"2C6E\char"032D} \SUPPORTEDANCH{\char"0050\char"032D} \SUPPORTEDANCH{\char"01A4\char"032D} \SUPPORTEDANCH{\char"0052\char"032D} \SUPPORTEDANCH{\char"2C64\char"032D} \SUPPORTEDPREC{\char"0054\char"032D} \SUPPORTEDANCH{\char"01AC\char"032D} \SUPPORTEDANCH{\char"01AE\char"032D} \SUPPORTEDANCH{\char"0059\char"032D} \SUPPORTEDANCH{\char"01B3\char"032D} \SUPPORTEDANCH{\char"0061\char"032D} \SUPPORTEDANCH{\char"1E9A\char"032D} \SUPPORTEDANCH{\char"0062\char"032D} \SUPPORTEDANCH{\char"0183\char"032D} \SUPPORTEDANCH{\char"0063\char"032D} \SUPPORTEDANCH{\char"0188\char"032D} \SUPPORTEDPREC{\char"0064\char"032D} \SUPPORTEDANCH{\char"018C\char"032D} \SUPPORTEDANCH{\char"0221\char"032D} \SUPPORTEDANCH{\char"0068\char"032D} \SUPPORTEDANCH{\char"2C68\char"032D} \SUPPORTEDANCH{\char"006B\char"032D} \SUPPORTEDANCH{\char"0199\char"032D} \SUPPORTEDANCH{\char"2C6A\char"032D} \SUPPORTEDPREC{\char"006C\char"032D} \SUPPORTEDANCH{\char"0234\char"032D} \SUPPORTEDPREC{\char"006E\char"032D} \SUPPORTEDANCH{\char"0235\char"032D} \SUPPORTEDANCH{\char"0070\char"032D} \SUPPORTEDANCH{\char"01A5\char"032D} \SUPPORTEDANCH{\char"0071\char"032D} \SUPPORTEDANCH{\char"02A0\char"032D} \SUPPORTEDPREC{\char"0074\char"032D} \SUPPORTEDANCH{\char"01AB\char"032D} \SUPPORTEDANCH{\char"01AD\char"032D} \SUPPORTEDANCH{\char"0236\char"032D} \SUPPORTEDPREC{\char"0075\char"032D} \MISSINGGLYPH{\char"AB52\char"032D} \SUPPORTEDANCH{\char"0079\char"032D} \SUPPORTEDANCH{\char"01B4\char"032D} \MISSINGGLYPH{\char"1EFF\char"032D}

\SUPPORTEDPREC{\char"0048\char"032E} \SUPPORTEDANCH{\char"2C67\char"032E} \SUPPORTEDANCH{\char"004B\char"032E} \SUPPORTEDANCH{\char"2C69\char"032E} \SUPPORTEDANCH{\char"004D\char"032E} \MISSINGGLYPH{\char"2C6E\char"032E} \SUPPORTEDANCH{\char"0050\char"032E} \SUPPORTEDANCH{\char"01A4\char"032E} \SUPPORTEDANCH{\char"0052\char"032E} \SUPPORTEDANCH{\char"2C64\char"032E} \SUPPORTEDANCH{\char"0054\char"032E} \SUPPORTEDANCH{\char"01AC\char"032E} \SUPPORTEDANCH{\char"01AE\char"032E} \SUPPORTEDANCH{\char"0059\char"032E} \SUPPORTEDANCH{\char"01B3\char"032E} \SUPPORTEDANCH{\char"0061\char"032E} \SUPPORTEDANCH{\char"1E9A\char"032E} \SUPPORTEDANCH{\char"0062\char"032E} \SUPPORTEDANCH{\char"0183\char"032E} \SUPPORTEDANCH{\char"0063\char"032E} \SUPPORTEDANCH{\char"0188\char"032E} \SUPPORTEDANCH{\char"0064\char"032E} \SUPPORTEDANCH{\char"018C\char"032E} \SUPPORTEDANCH{\char"0221\char"032E} \SUPPORTEDPREC{\char"0068\char"032E} \SUPPORTEDANCH{\char"2C68\char"032E} \SUPPORTEDANCH{\char"006B\char"032E} \SUPPORTEDANCH{\char"0199\char"032E} \SUPPORTEDANCH{\char"2C6A\char"032E} \SUPPORTEDANCH{\char"006C\char"032E} \SUPPORTEDANCH{\char"0234\char"032E} \SUPPORTEDANCH{\char"006E\char"032E} \SUPPORTEDANCH{\char"0235\char"032E} \SUPPORTEDANCH{\char"0070\char"032E} \SUPPORTEDANCH{\char"01A5\char"032E} \SUPPORTEDANCH{\char"0071\char"032E} \SUPPORTEDANCH{\char"02A0\char"032E} \SUPPORTEDANCH{\char"0074\char"032E} \SUPPORTEDANCH{\char"01AB\char"032E} \SUPPORTEDANCH{\char"01AD\char"032E} \SUPPORTEDANCH{\char"0236\char"032E} \SUPPORTEDANCH{\char"0075\char"032E} \MISSINGGLYPH{\char"AB52\char"032E} \SUPPORTEDANCH{\char"0079\char"032E} \SUPPORTEDANCH{\char"01B4\char"032E} \MISSINGGLYPH{\char"1EFF\char"032E}

\SUPPORTEDANCH{\char"0048\char"032F} \SUPPORTEDANCH{\char"2C67\char"032F} \SUPPORTEDANCH{\char"004B\char"032F} \SUPPORTEDANCH{\char"2C69\char"032F} \SUPPORTEDANCH{\char"004D\char"032F} \MISSINGGLYPH{\char"2C6E\char"032F} \SUPPORTEDANCH{\char"0050\char"032F} \SUPPORTEDANCH{\char"01A4\char"032F} \SUPPORTEDANCH{\char"0052\char"032F} \SUPPORTEDANCH{\char"2C64\char"032F} \SUPPORTEDANCH{\char"0054\char"032F} \SUPPORTEDANCH{\char"01AC\char"032F} \SUPPORTEDANCH{\char"01AE\char"032F} \SUPPORTEDANCH{\char"0059\char"032F} \SUPPORTEDANCH{\char"01B3\char"032F} \SUPPORTEDANCH{\char"0061\char"032F} \SUPPORTEDANCH{\char"1E9A\char"032F} \SUPPORTEDANCH{\char"0062\char"032F} \SUPPORTEDANCH{\char"0183\char"032F} \SUPPORTEDANCH{\char"0063\char"032F} \SUPPORTEDANCH{\char"0188\char"032F} \SUPPORTEDANCH{\char"0064\char"032F} \SUPPORTEDANCH{\char"018C\char"032F} \SUPPORTEDANCH{\char"0221\char"032F} \SUPPORTEDANCH{\char"0068\char"032F} \SUPPORTEDANCH{\char"2C68\char"032F} \SUPPORTEDANCH{\char"006B\char"032F} \SUPPORTEDANCH{\char"0199\char"032F} \SUPPORTEDANCH{\char"2C6A\char"032F} \SUPPORTEDANCH{\char"006C\char"032F} \SUPPORTEDANCH{\char"0234\char"032F} \SUPPORTEDANCH{\char"006E\char"032F} \SUPPORTEDANCH{\char"0235\char"032F} \SUPPORTEDANCH{\char"0070\char"032F} \SUPPORTEDANCH{\char"01A5\char"032F} \SUPPORTEDANCH{\char"0071\char"032F} \SUPPORTEDANCH{\char"02A0\char"032F} \SUPPORTEDANCH{\char"0074\char"032F} \SUPPORTEDANCH{\char"01AB\char"032F} \SUPPORTEDANCH{\char"01AD\char"032F} \SUPPORTEDANCH{\char"0236\char"032F} \SUPPORTEDANCH{\char"0075\char"032F} \MISSINGGLYPH{\char"AB52\char"032F} \SUPPORTEDANCH{\char"0079\char"032F} \SUPPORTEDANCH{\char"01B4\char"032F} \MISSINGGLYPH{\char"1EFF\char"032F}

\SUPPORTEDANCH{\char"0048\char"031D} \SUPPORTEDANCH{\char"2C67\char"031D} \SUPPORTEDANCH{\char"004B\char"031D} \SUPPORTEDANCH{\char"2C69\char"031D} \SUPPORTEDANCH{\char"004D\char"031D} \MISSINGGLYPH{\char"2C6E\char"031D} \SUPPORTEDANCH{\char"0050\char"031D} \SUPPORTEDANCH{\char"01A4\char"031D} \SUPPORTEDANCH{\char"0052\char"031D} \SUPPORTEDANCH{\char"2C64\char"031D} \SUPPORTEDANCH{\char"0054\char"031D} \SUPPORTEDANCH{\char"01AC\char"031D} \SUPPORTEDANCH{\char"01AE\char"031D} \SUPPORTEDANCH{\char"0059\char"031D} \SUPPORTEDANCH{\char"01B3\char"031D} \SUPPORTEDANCH{\char"0061\char"031D} \SUPPORTEDANCH{\char"1E9A\char"031D} \SUPPORTEDANCH{\char"0062\char"031D} \SUPPORTEDANCH{\char"0183\char"031D} \SUPPORTEDANCH{\char"0063\char"031D} \SUPPORTEDANCH{\char"0188\char"031D} \SUPPORTEDANCH{\char"0064\char"031D} \SUPPORTEDANCH{\char"018C\char"031D} \SUPPORTEDANCH{\char"0221\char"031D} \SUPPORTEDANCH{\char"0068\char"031D} \SUPPORTEDANCH{\char"2C68\char"031D} \SUPPORTEDANCH{\char"006B\char"031D} \SUPPORTEDANCH{\char"0199\char"031D} \SUPPORTEDANCH{\char"2C6A\char"031D} \SUPPORTEDANCH{\char"006C\char"031D} \SUPPORTEDANCH{\char"0234\char"031D} \SUPPORTEDANCH{\char"006E\char"031D} \SUPPORTEDANCH{\char"0235\char"031D} \SUPPORTEDANCH{\char"0070\char"031D} \SUPPORTEDANCH{\char"01A5\char"031D} \SUPPORTEDANCH{\char"0071\char"031D} \SUPPORTEDANCH{\char"02A0\char"031D} \SUPPORTEDANCH{\char"0074\char"031D} \SUPPORTEDANCH{\char"01AB\char"031D} \SUPPORTEDANCH{\char"01AD\char"031D} \SUPPORTEDANCH{\char"0236\char"031D} \SUPPORTEDANCH{\char"0075\char"031D} \MISSINGGLYPH{\char"AB52\char"031D} \SUPPORTEDANCH{\char"0079\char"031D} \SUPPORTEDANCH{\char"01B4\char"031D} \MISSINGGLYPH{\char"1EFF\char"031D}

\SUPPORTEDANCH{\char"0048\char"031E} \SUPPORTEDANCH{\char"2C67\char"031E} \SUPPORTEDANCH{\char"004B\char"031E} \SUPPORTEDANCH{\char"2C69\char"031E} \SUPPORTEDANCH{\char"004D\char"031E} \MISSINGGLYPH{\char"2C6E\char"031E} \SUPPORTEDANCH{\char"0050\char"031E} \SUPPORTEDANCH{\char"01A4\char"031E} \SUPPORTEDANCH{\char"0052\char"031E} \SUPPORTEDANCH{\char"2C64\char"031E} \SUPPORTEDANCH{\char"0054\char"031E} \SUPPORTEDANCH{\char"01AC\char"031E} \SUPPORTEDANCH{\char"01AE\char"031E} \SUPPORTEDANCH{\char"0059\char"031E} \SUPPORTEDANCH{\char"01B3\char"031E} \SUPPORTEDANCH{\char"0061\char"031E} \SUPPORTEDANCH{\char"1E9A\char"031E} \SUPPORTEDANCH{\char"0062\char"031E} \SUPPORTEDANCH{\char"0183\char"031E} \SUPPORTEDANCH{\char"0063\char"031E} \SUPPORTEDANCH{\char"0188\char"031E} \SUPPORTEDANCH{\char"0064\char"031E} \SUPPORTEDANCH{\char"018C\char"031E} \SUPPORTEDANCH{\char"0221\char"031E} \SUPPORTEDANCH{\char"0068\char"031E} \SUPPORTEDANCH{\char"2C68\char"031E} \SUPPORTEDANCH{\char"006B\char"031E} \SUPPORTEDANCH{\char"0199\char"031E} \SUPPORTEDANCH{\char"2C6A\char"031E} \SUPPORTEDANCH{\char"006C\char"031E} \SUPPORTEDANCH{\char"0234\char"031E} \SUPPORTEDANCH{\char"006E\char"031E} \SUPPORTEDANCH{\char"0235\char"031E} \SUPPORTEDANCH{\char"0070\char"031E} \SUPPORTEDANCH{\char"01A5\char"031E} \SUPPORTEDANCH{\char"0071\char"031E} \SUPPORTEDANCH{\char"02A0\char"031E} \SUPPORTEDANCH{\char"0074\char"031E} \SUPPORTEDANCH{\char"01AB\char"031E} \SUPPORTEDANCH{\char"01AD\char"031E} \SUPPORTEDANCH{\char"0236\char"031E} \SUPPORTEDANCH{\char"0075\char"031E} \MISSINGGLYPH{\char"AB52\char"031E} \SUPPORTEDANCH{\char"0079\char"031E} \SUPPORTEDANCH{\char"01B4\char"031E} \MISSINGGLYPH{\char"1EFF\char"031E}

\SUPPORTEDANCH{\char"0048\char"031F} \SUPPORTEDANCH{\char"2C67\char"031F} \SUPPORTEDANCH{\char"004B\char"031F} \SUPPORTEDANCH{\char"2C69\char"031F} \SUPPORTEDANCH{\char"004D\char"031F} \MISSINGGLYPH{\char"2C6E\char"031F} \SUPPORTEDANCH{\char"0050\char"031F} \SUPPORTEDANCH{\char"01A4\char"031F} \SUPPORTEDANCH{\char"0052\char"031F} \SUPPORTEDANCH{\char"2C64\char"031F} \SUPPORTEDANCH{\char"0054\char"031F} \SUPPORTEDANCH{\char"01AC\char"031F} \SUPPORTEDANCH{\char"01AE\char"031F} \SUPPORTEDANCH{\char"0059\char"031F} \SUPPORTEDANCH{\char"01B3\char"031F} \SUPPORTEDANCH{\char"0061\char"031F} \SUPPORTEDANCH{\char"1E9A\char"031F} \SUPPORTEDANCH{\char"0062\char"031F} \SUPPORTEDANCH{\char"0183\char"031F} \SUPPORTEDANCH{\char"0063\char"031F} \SUPPORTEDANCH{\char"0188\char"031F} \SUPPORTEDANCH{\char"0064\char"031F} \SUPPORTEDANCH{\char"018C\char"031F} \SUPPORTEDANCH{\char"0221\char"031F} \SUPPORTEDANCH{\char"0068\char"031F} \SUPPORTEDANCH{\char"2C68\char"031F} \SUPPORTEDANCH{\char"006B\char"031F} \SUPPORTEDANCH{\char"0199\char"031F} \SUPPORTEDANCH{\char"2C6A\char"031F} \SUPPORTEDANCH{\char"006C\char"031F} \SUPPORTEDANCH{\char"0234\char"031F} \SUPPORTEDANCH{\char"006E\char"031F} \SUPPORTEDANCH{\char"0235\char"031F} \SUPPORTEDANCH{\char"0070\char"031F} \SUPPORTEDANCH{\char"01A5\char"031F} \SUPPORTEDANCH{\char"0071\char"031F} \SUPPORTEDANCH{\char"02A0\char"031F} \SUPPORTEDANCH{\char"0074\char"031F} \SUPPORTEDANCH{\char"01AB\char"031F} \SUPPORTEDANCH{\char"01AD\char"031F} \SUPPORTEDANCH{\char"0236\char"031F} \SUPPORTEDANCH{\char"0075\char"031F} \MISSINGGLYPH{\char"AB52\char"031F} \SUPPORTEDANCH{\char"0079\char"031F} \SUPPORTEDANCH{\char"01B4\char"031F} \MISSINGGLYPH{\char"1EFF\char"031F}

\SUPPORTEDANCH{\char"0048\char"0330} \SUPPORTEDANCH{\char"2C67\char"0330} \SUPPORTEDANCH{\char"004B\char"0330} \SUPPORTEDANCH{\char"2C69\char"0330} \SUPPORTEDANCH{\char"004D\char"0330} \MISSINGGLYPH{\char"2C6E\char"0330} \SUPPORTEDANCH{\char"0050\char"0330} \SUPPORTEDANCH{\char"01A4\char"0330} \SUPPORTEDANCH{\char"0052\char"0330} \SUPPORTEDANCH{\char"2C64\char"0330} \SUPPORTEDANCH{\char"0054\char"0330} \SUPPORTEDANCH{\char"01AC\char"0330} \SUPPORTEDANCH{\char"01AE\char"0330} \SUPPORTEDANCH{\char"0059\char"0330} \SUPPORTEDANCH{\char"01B3\char"0330} \SUPPORTEDANCH{\char"0061\char"0330} \SUPPORTEDANCH{\char"1E9A\char"0330} \SUPPORTEDANCH{\char"0062\char"0330} \SUPPORTEDANCH{\char"0183\char"0330} \SUPPORTEDANCH{\char"0063\char"0330} \SUPPORTEDANCH{\char"0188\char"0330} \SUPPORTEDANCH{\char"0064\char"0330} \SUPPORTEDANCH{\char"018C\char"0330} \SUPPORTEDANCH{\char"0221\char"0330} \SUPPORTEDANCH{\char"0068\char"0330} \SUPPORTEDANCH{\char"2C68\char"0330} \SUPPORTEDANCH{\char"006B\char"0330} \SUPPORTEDANCH{\char"0199\char"0330} \SUPPORTEDANCH{\char"2C6A\char"0330} \SUPPORTEDANCH{\char"006C\char"0330} \SUPPORTEDANCH{\char"0234\char"0330} \SUPPORTEDANCH{\char"006E\char"0330} \SUPPORTEDANCH{\char"0235\char"0330} \SUPPORTEDANCH{\char"0070\char"0330} \SUPPORTEDANCH{\char"01A5\char"0330} \SUPPORTEDANCH{\char"0071\char"0330} \SUPPORTEDANCH{\char"02A0\char"0330} \SUPPORTEDANCH{\char"0074\char"0330} \SUPPORTEDANCH{\char"01AB\char"0330} \SUPPORTEDANCH{\char"01AD\char"0330} \SUPPORTEDANCH{\char"0236\char"0330} \SUPPORTEDPREC{\char"0075\char"0330} \MISSINGGLYPH{\char"AB52\char"0330} \SUPPORTEDANCH{\char"0079\char"0330} \SUPPORTEDANCH{\char"01B4\char"0330} \MISSINGGLYPH{\char"1EFF\char"0330}

\SUPPORTEDANCH{\char"0048\char"0331} \SUPPORTEDANCH{\char"2C67\char"0331} \SUPPORTEDPREC{\char"004B\char"0331} \SUPPORTEDANCH{\char"2C69\char"0331} \SUPPORTEDANCH{\char"004D\char"0331} \MISSINGGLYPH{\char"2C6E\char"0331} \SUPPORTEDANCH{\char"0050\char"0331} \SUPPORTEDANCH{\char"01A4\char"0331} \SUPPORTEDPREC{\char"0052\char"0331} \SUPPORTEDANCH{\char"2C64\char"0331} \SUPPORTEDPREC{\char"0054\char"0331} \SUPPORTEDANCH{\char"01AC\char"0331} \SUPPORTEDANCH{\char"01AE\char"0331} \SUPPORTEDANCH{\char"0059\char"0331} \SUPPORTEDANCH{\char"01B3\char"0331} \SUPPORTEDANCH{\char"0061\char"0331} \SUPPORTEDANCH{\char"1E9A\char"0331} \SUPPORTEDPREC{\char"0062\char"0331} \SUPPORTEDANCH{\char"0183\char"0331} \SUPPORTEDANCH{\char"0063\char"0331} \SUPPORTEDANCH{\char"0188\char"0331} \SUPPORTEDPREC{\char"0064\char"0331} \SUPPORTEDANCH{\char"018C\char"0331} \SUPPORTEDANCH{\char"0221\char"0331} \SUPPORTEDPREC{\char"0068\char"0331} \SUPPORTEDANCH{\char"2C68\char"0331} \SUPPORTEDPREC{\char"006B\char"0331} \SUPPORTEDANCH{\char"0199\char"0331} \SUPPORTEDANCH{\char"2C6A\char"0331} \SUPPORTEDPREC{\char"006C\char"0331} \SUPPORTEDANCH{\char"0234\char"0331} \SUPPORTEDPREC{\char"006E\char"0331} \SUPPORTEDANCH{\char"0235\char"0331} \SUPPORTEDANCH{\char"0070\char"0331} \SUPPORTEDANCH{\char"01A5\char"0331} \SUPPORTEDANCH{\char"0071\char"0331} \SUPPORTEDANCH{\char"02A0\char"0331} \SUPPORTEDPREC{\char"0074\char"0331} \SUPPORTEDANCH{\char"01AB\char"0331} \SUPPORTEDANCH{\char"01AD\char"0331} \SUPPORTEDANCH{\char"0236\char"0331} \SUPPORTEDANCH{\char"0075\char"0331} \MISSINGGLYPH{\char"AB52\char"0331} \SUPPORTEDANCH{\char"0079\char"0331} \SUPPORTEDANCH{\char"01B4\char"0331} \MISSINGGLYPH{\char"1EFF\char"0331}

\SUPPORTEDANCH{\char"0048\char"0333} \SUPPORTEDANCH{\char"2C67\char"0333} \SUPPORTEDANCH{\char"004B\char"0333} \SUPPORTEDANCH{\char"2C69\char"0333} \SUPPORTEDANCH{\char"004D\char"0333} \MISSINGGLYPH{\char"2C6E\char"0333} \SUPPORTEDANCH{\char"0050\char"0333} \SUPPORTEDANCH{\char"01A4\char"0333} \SUPPORTEDANCH{\char"0052\char"0333} \SUPPORTEDANCH{\char"2C64\char"0333} \SUPPORTEDANCH{\char"0054\char"0333} \SUPPORTEDANCH{\char"01AC\char"0333} \SUPPORTEDANCH{\char"01AE\char"0333} \SUPPORTEDANCH{\char"0059\char"0333} \SUPPORTEDANCH{\char"01B3\char"0333} \SUPPORTEDANCH{\char"0061\char"0333} \SUPPORTEDANCH{\char"1E9A\char"0333} \SUPPORTEDANCH{\char"0062\char"0333} \SUPPORTEDANCH{\char"0183\char"0333} \SUPPORTEDANCH{\char"0063\char"0333} \SUPPORTEDANCH{\char"0188\char"0333} \SUPPORTEDANCH{\char"0064\char"0333} \SUPPORTEDANCH{\char"018C\char"0333} \SUPPORTEDANCH{\char"0221\char"0333} \SUPPORTEDANCH{\char"0068\char"0333} \SUPPORTEDANCH{\char"2C68\char"0333} \SUPPORTEDANCH{\char"006B\char"0333} \SUPPORTEDANCH{\char"0199\char"0333} \SUPPORTEDANCH{\char"2C6A\char"0333} \SUPPORTEDANCH{\char"006C\char"0333} \SUPPORTEDANCH{\char"0234\char"0333} \SUPPORTEDANCH{\char"006E\char"0333} \SUPPORTEDANCH{\char"0235\char"0333} \SUPPORTEDANCH{\char"0070\char"0333} \SUPPORTEDANCH{\char"01A5\char"0333} \SUPPORTEDANCH{\char"0071\char"0333} \SUPPORTEDANCH{\char"02A0\char"0333} \SUPPORTEDANCH{\char"0074\char"0333} \SUPPORTEDANCH{\char"01AB\char"0333} \SUPPORTEDANCH{\char"01AD\char"0333} \SUPPORTEDANCH{\char"0236\char"0333} \SUPPORTEDANCH{\char"0075\char"0333} \MISSINGGLYPH{\char"AB52\char"0333} \SUPPORTEDANCH{\char"0079\char"0333} \SUPPORTEDANCH{\char"01B4\char"0333} \MISSINGGLYPH{\char"1EFF\char"0333}

\SUPPORTEDANCH{\char"0048\char"0339} \SUPPORTEDANCH{\char"2C67\char"0339} \SUPPORTEDANCH{\char"004B\char"0339} \SUPPORTEDANCH{\char"2C69\char"0339} \SUPPORTEDANCH{\char"004D\char"0339} \MISSINGGLYPH{\char"2C6E\char"0339} \SUPPORTEDANCH{\char"0050\char"0339} \SUPPORTEDANCH{\char"01A4\char"0339} \SUPPORTEDANCH{\char"0052\char"0339} \SUPPORTEDANCH{\char"2C64\char"0339} \SUPPORTEDANCH{\char"0054\char"0339} \SUPPORTEDANCH{\char"01AC\char"0339} \SUPPORTEDANCH{\char"01AE\char"0339} \SUPPORTEDANCH{\char"0059\char"0339} \SUPPORTEDANCH{\char"01B3\char"0339} \SUPPORTEDANCH{\char"0061\char"0339} \SUPPORTEDANCH{\char"1E9A\char"0339} \SUPPORTEDANCH{\char"0062\char"0339} \SUPPORTEDANCH{\char"0183\char"0339} \SUPPORTEDANCH{\char"0063\char"0339} \SUPPORTEDANCH{\char"0188\char"0339} \SUPPORTEDANCH{\char"0064\char"0339} \SUPPORTEDANCH{\char"018C\char"0339} \SUPPORTEDANCH{\char"0221\char"0339} \SUPPORTEDANCH{\char"0068\char"0339} \SUPPORTEDANCH{\char"2C68\char"0339} \SUPPORTEDANCH{\char"006B\char"0339} \SUPPORTEDANCH{\char"0199\char"0339} \SUPPORTEDANCH{\char"2C6A\char"0339} \SUPPORTEDANCH{\char"006C\char"0339} \SUPPORTEDANCH{\char"0234\char"0339} \SUPPORTEDANCH{\char"006E\char"0339} \SUPPORTEDANCH{\char"0235\char"0339} \SUPPORTEDANCH{\char"0070\char"0339} \SUPPORTEDANCH{\char"01A5\char"0339} \SUPPORTEDANCH{\char"0071\char"0339} \SUPPORTEDANCH{\char"02A0\char"0339} \SUPPORTEDANCH{\char"0074\char"0339} \SUPPORTEDANCH{\char"01AB\char"0339} \SUPPORTEDANCH{\char"01AD\char"0339} \SUPPORTEDANCH{\char"0236\char"0339} \SUPPORTEDANCH{\char"0075\char"0339} \MISSINGGLYPH{\char"AB52\char"0339} \SUPPORTEDANCH{\char"0079\char"0339} \SUPPORTEDANCH{\char"01B4\char"0339} \MISSINGGLYPH{\char"1EFF\char"0339}

\SUPPORTEDANCH{\char"0048\char"033A} \SUPPORTEDANCH{\char"2C67\char"033A} \SUPPORTEDANCH{\char"004B\char"033A} \SUPPORTEDANCH{\char"2C69\char"033A} \SUPPORTEDANCH{\char"004D\char"033A} \MISSINGGLYPH{\char"2C6E\char"033A} \SUPPORTEDANCH{\char"0050\char"033A} \SUPPORTEDANCH{\char"01A4\char"033A} \SUPPORTEDANCH{\char"0052\char"033A} \SUPPORTEDANCH{\char"2C64\char"033A} \SUPPORTEDANCH{\char"0054\char"033A} \SUPPORTEDANCH{\char"01AC\char"033A} \SUPPORTEDANCH{\char"01AE\char"033A} \SUPPORTEDANCH{\char"0059\char"033A} \SUPPORTEDANCH{\char"01B3\char"033A} \SUPPORTEDANCH{\char"0061\char"033A} \SUPPORTEDANCH{\char"1E9A\char"033A} \SUPPORTEDANCH{\char"0062\char"033A} \SUPPORTEDANCH{\char"0183\char"033A} \SUPPORTEDANCH{\char"0063\char"033A} \SUPPORTEDANCH{\char"0188\char"033A} \SUPPORTEDANCH{\char"0064\char"033A} \SUPPORTEDANCH{\char"018C\char"033A} \SUPPORTEDANCH{\char"0221\char"033A} \SUPPORTEDANCH{\char"0068\char"033A} \SUPPORTEDANCH{\char"2C68\char"033A} \SUPPORTEDANCH{\char"006B\char"033A} \SUPPORTEDANCH{\char"0199\char"033A} \SUPPORTEDANCH{\char"2C6A\char"033A} \SUPPORTEDANCH{\char"006C\char"033A} \SUPPORTEDANCH{\char"0234\char"033A} \SUPPORTEDANCH{\char"006E\char"033A} \SUPPORTEDANCH{\char"0235\char"033A} \SUPPORTEDANCH{\char"0070\char"033A} \SUPPORTEDANCH{\char"01A5\char"033A} \SUPPORTEDANCH{\char"0071\char"033A} \SUPPORTEDANCH{\char"02A0\char"033A} \SUPPORTEDANCH{\char"0074\char"033A} \SUPPORTEDANCH{\char"01AB\char"033A} \SUPPORTEDANCH{\char"01AD\char"033A} \SUPPORTEDANCH{\char"0236\char"033A} \SUPPORTEDANCH{\char"0075\char"033A} \MISSINGGLYPH{\char"AB52\char"033A} \SUPPORTEDANCH{\char"0079\char"033A} \SUPPORTEDANCH{\char"01B4\char"033A} \MISSINGGLYPH{\char"1EFF\char"033A}

\SUPPORTEDANCH{\char"0048\char"033B} \SUPPORTEDANCH{\char"2C67\char"033B} \SUPPORTEDANCH{\char"004B\char"033B} \SUPPORTEDANCH{\char"2C69\char"033B} \SUPPORTEDANCH{\char"004D\char"033B} \MISSINGGLYPH{\char"2C6E\char"033B} \SUPPORTEDANCH{\char"0050\char"033B} \SUPPORTEDANCH{\char"01A4\char"033B} \SUPPORTEDANCH{\char"0052\char"033B} \SUPPORTEDANCH{\char"2C64\char"033B} \SUPPORTEDANCH{\char"0054\char"033B} \SUPPORTEDANCH{\char"01AC\char"033B} \SUPPORTEDANCH{\char"01AE\char"033B} \SUPPORTEDANCH{\char"0059\char"033B} \SUPPORTEDANCH{\char"01B3\char"033B} \SUPPORTEDANCH{\char"0061\char"033B} \SUPPORTEDANCH{\char"1E9A\char"033B} \SUPPORTEDANCH{\char"0062\char"033B} \SUPPORTEDANCH{\char"0183\char"033B} \SUPPORTEDANCH{\char"0063\char"033B} \SUPPORTEDANCH{\char"0188\char"033B} \SUPPORTEDANCH{\char"0064\char"033B} \SUPPORTEDANCH{\char"018C\char"033B} \SUPPORTEDANCH{\char"0221\char"033B} \SUPPORTEDANCH{\char"0068\char"033B} \SUPPORTEDANCH{\char"2C68\char"033B} \SUPPORTEDANCH{\char"006B\char"033B} \SUPPORTEDANCH{\char"0199\char"033B} \SUPPORTEDANCH{\char"2C6A\char"033B} \SUPPORTEDANCH{\char"006C\char"033B} \SUPPORTEDANCH{\char"0234\char"033B} \SUPPORTEDANCH{\char"006E\char"033B} \SUPPORTEDANCH{\char"0235\char"033B} \SUPPORTEDANCH{\char"0070\char"033B} \SUPPORTEDANCH{\char"01A5\char"033B} \SUPPORTEDANCH{\char"0071\char"033B} \SUPPORTEDANCH{\char"02A0\char"033B} \SUPPORTEDANCH{\char"0074\char"033B} \SUPPORTEDANCH{\char"01AB\char"033B} \SUPPORTEDANCH{\char"01AD\char"033B} \SUPPORTEDANCH{\char"0236\char"033B} \SUPPORTEDANCH{\char"0075\char"033B} \MISSINGGLYPH{\char"AB52\char"033B} \SUPPORTEDANCH{\char"0079\char"033B} \SUPPORTEDANCH{\char"01B4\char"033B} \MISSINGGLYPH{\char"1EFF\char"033B}

\SUPPORTEDANCH{\char"0048\char"033C} \SUPPORTEDANCH{\char"2C67\char"033C} \SUPPORTEDANCH{\char"004B\char"033C} \SUPPORTEDANCH{\char"2C69\char"033C} \SUPPORTEDANCH{\char"004D\char"033C} \MISSINGGLYPH{\char"2C6E\char"033C} \SUPPORTEDANCH{\char"0050\char"033C} \SUPPORTEDANCH{\char"01A4\char"033C} \SUPPORTEDANCH{\char"0052\char"033C} \SUPPORTEDANCH{\char"2C64\char"033C} \SUPPORTEDANCH{\char"0054\char"033C} \SUPPORTEDANCH{\char"01AC\char"033C} \SUPPORTEDANCH{\char"01AE\char"033C} \SUPPORTEDANCH{\char"0059\char"033C} \SUPPORTEDANCH{\char"01B3\char"033C} \SUPPORTEDANCH{\char"0061\char"033C} \SUPPORTEDANCH{\char"1E9A\char"033C} \SUPPORTEDANCH{\char"0062\char"033C} \SUPPORTEDANCH{\char"0183\char"033C} \SUPPORTEDANCH{\char"0063\char"033C} \SUPPORTEDANCH{\char"0188\char"033C} \SUPPORTEDANCH{\char"0064\char"033C} \SUPPORTEDANCH{\char"018C\char"033C} \SUPPORTEDANCH{\char"0221\char"033C} \SUPPORTEDANCH{\char"0068\char"033C} \SUPPORTEDANCH{\char"2C68\char"033C} \SUPPORTEDANCH{\char"006B\char"033C} \SUPPORTEDANCH{\char"0199\char"033C} \SUPPORTEDANCH{\char"2C6A\char"033C} \SUPPORTEDANCH{\char"006C\char"033C} \SUPPORTEDANCH{\char"0234\char"033C} \SUPPORTEDANCH{\char"006E\char"033C} \SUPPORTEDANCH{\char"0235\char"033C} \SUPPORTEDANCH{\char"0070\char"033C} \SUPPORTEDANCH{\char"01A5\char"033C} \SUPPORTEDANCH{\char"0071\char"033C} \SUPPORTEDANCH{\char"02A0\char"033C} \SUPPORTEDANCH{\char"0074\char"033C} \SUPPORTEDANCH{\char"01AB\char"033C} \SUPPORTEDANCH{\char"01AD\char"033C} \SUPPORTEDANCH{\char"0236\char"033C} \SUPPORTEDANCH{\char"0075\char"033C} \MISSINGGLYPH{\char"AB52\char"033C} \SUPPORTEDANCH{\char"0079\char"033C} \SUPPORTEDANCH{\char"01B4\char"033C} \MISSINGGLYPH{\char"1EFF\char"033C}

\SUPPORTEDANCH{\char"0048\char"0347} \SUPPORTEDANCH{\char"2C67\char"0347} \SUPPORTEDANCH{\char"004B\char"0347} \SUPPORTEDANCH{\char"2C69\char"0347} \SUPPORTEDANCH{\char"004D\char"0347} \MISSINGGLYPH{\char"2C6E\char"0347} \SUPPORTEDANCH{\char"0050\char"0347} \SUPPORTEDANCH{\char"01A4\char"0347} \SUPPORTEDANCH{\char"0052\char"0347} \SUPPORTEDANCH{\char"2C64\char"0347} \SUPPORTEDANCH{\char"0054\char"0347} \SUPPORTEDANCH{\char"01AC\char"0347} \SUPPORTEDANCH{\char"01AE\char"0347} \SUPPORTEDANCH{\char"0059\char"0347} \SUPPORTEDANCH{\char"01B3\char"0347} \SUPPORTEDANCH{\char"0061\char"0347} \SUPPORTEDANCH{\char"1E9A\char"0347} \SUPPORTEDANCH{\char"0062\char"0347} \SUPPORTEDANCH{\char"0183\char"0347} \SUPPORTEDANCH{\char"0063\char"0347} \SUPPORTEDANCH{\char"0188\char"0347} \SUPPORTEDANCH{\char"0064\char"0347} \SUPPORTEDANCH{\char"018C\char"0347} \SUPPORTEDANCH{\char"0221\char"0347} \SUPPORTEDANCH{\char"0068\char"0347} \SUPPORTEDANCH{\char"2C68\char"0347} \SUPPORTEDANCH{\char"006B\char"0347} \SUPPORTEDANCH{\char"0199\char"0347} \SUPPORTEDANCH{\char"2C6A\char"0347} \SUPPORTEDANCH{\char"006C\char"0347} \SUPPORTEDANCH{\char"0234\char"0347} \SUPPORTEDANCH{\char"006E\char"0347} \SUPPORTEDANCH{\char"0235\char"0347} \SUPPORTEDANCH{\char"0070\char"0347} \SUPPORTEDANCH{\char"01A5\char"0347} \SUPPORTEDANCH{\char"0071\char"0347} \SUPPORTEDANCH{\char"02A0\char"0347} \SUPPORTEDANCH{\char"0074\char"0347} \SUPPORTEDANCH{\char"01AB\char"0347} \SUPPORTEDANCH{\char"01AD\char"0347} \SUPPORTEDANCH{\char"0236\char"0347} \SUPPORTEDANCH{\char"0075\char"0347} \MISSINGGLYPH{\char"AB52\char"0347} \SUPPORTEDANCH{\char"0079\char"0347} \SUPPORTEDANCH{\char"01B4\char"0347} \MISSINGGLYPH{\char"1EFF\char"0347}

\SUPPORTEDANCH{\char"0048\char"0348} \SUPPORTEDANCH{\char"2C67\char"0348} \SUPPORTEDANCH{\char"004B\char"0348} \SUPPORTEDANCH{\char"2C69\char"0348} \SUPPORTEDANCH{\char"004D\char"0348} \MISSINGGLYPH{\char"2C6E\char"0348} \SUPPORTEDANCH{\char"0050\char"0348} \SUPPORTEDANCH{\char"01A4\char"0348} \SUPPORTEDANCH{\char"0052\char"0348} \SUPPORTEDANCH{\char"2C64\char"0348} \SUPPORTEDANCH{\char"0054\char"0348} \SUPPORTEDANCH{\char"01AC\char"0348} \SUPPORTEDANCH{\char"01AE\char"0348} \SUPPORTEDANCH{\char"0059\char"0348} \SUPPORTEDANCH{\char"01B3\char"0348} \SUPPORTEDANCH{\char"0061\char"0348} \SUPPORTEDANCH{\char"1E9A\char"0348} \SUPPORTEDANCH{\char"0062\char"0348} \SUPPORTEDANCH{\char"0183\char"0348} \SUPPORTEDANCH{\char"0063\char"0348} \SUPPORTEDANCH{\char"0188\char"0348} \SUPPORTEDANCH{\char"0064\char"0348} \SUPPORTEDANCH{\char"018C\char"0348} \SUPPORTEDANCH{\char"0221\char"0348} \SUPPORTEDANCH{\char"0068\char"0348} \SUPPORTEDANCH{\char"2C68\char"0348} \SUPPORTEDANCH{\char"006B\char"0348} \SUPPORTEDANCH{\char"0199\char"0348} \SUPPORTEDANCH{\char"2C6A\char"0348} \SUPPORTEDANCH{\char"006C\char"0348} \SUPPORTEDANCH{\char"0234\char"0348} \SUPPORTEDANCH{\char"006E\char"0348} \SUPPORTEDANCH{\char"0235\char"0348} \SUPPORTEDANCH{\char"0070\char"0348} \SUPPORTEDANCH{\char"01A5\char"0348} \SUPPORTEDANCH{\char"0071\char"0348} \SUPPORTEDANCH{\char"02A0\char"0348} \SUPPORTEDANCH{\char"0074\char"0348} \SUPPORTEDANCH{\char"01AB\char"0348} \SUPPORTEDANCH{\char"01AD\char"0348} \SUPPORTEDANCH{\char"0236\char"0348} \SUPPORTEDANCH{\char"0075\char"0348} \MISSINGGLYPH{\char"AB52\char"0348} \SUPPORTEDANCH{\char"0079\char"0348} \SUPPORTEDANCH{\char"01B4\char"0348} \MISSINGGLYPH{\char"1EFF\char"0348}

\SUPPORTEDANCH{\char"0048\char"0349} \SUPPORTEDANCH{\char"2C67\char"0349} \SUPPORTEDANCH{\char"004B\char"0349} \SUPPORTEDANCH{\char"2C69\char"0349} \SUPPORTEDANCH{\char"004D\char"0349} \MISSINGGLYPH{\char"2C6E\char"0349} \SUPPORTEDANCH{\char"0050\char"0349} \SUPPORTEDANCH{\char"01A4\char"0349} \SUPPORTEDANCH{\char"0052\char"0349} \SUPPORTEDANCH{\char"2C64\char"0349} \SUPPORTEDANCH{\char"0054\char"0349} \SUPPORTEDANCH{\char"01AC\char"0349} \SUPPORTEDANCH{\char"01AE\char"0349} \SUPPORTEDANCH{\char"0059\char"0349} \SUPPORTEDANCH{\char"01B3\char"0349} \SUPPORTEDANCH{\char"0061\char"0349} \SUPPORTEDANCH{\char"1E9A\char"0349} \SUPPORTEDANCH{\char"0062\char"0349} \SUPPORTEDANCH{\char"0183\char"0349} \SUPPORTEDANCH{\char"0063\char"0349} \SUPPORTEDANCH{\char"0188\char"0349} \SUPPORTEDANCH{\char"0064\char"0349} \SUPPORTEDANCH{\char"018C\char"0349} \SUPPORTEDANCH{\char"0221\char"0349} \SUPPORTEDANCH{\char"0068\char"0349} \SUPPORTEDANCH{\char"2C68\char"0349} \SUPPORTEDANCH{\char"006B\char"0349} \SUPPORTEDANCH{\char"0199\char"0349} \SUPPORTEDANCH{\char"2C6A\char"0349} \SUPPORTEDANCH{\char"006C\char"0349} \SUPPORTEDANCH{\char"0234\char"0349} \SUPPORTEDANCH{\char"006E\char"0349} \SUPPORTEDANCH{\char"0235\char"0349} \SUPPORTEDANCH{\char"0070\char"0349} \SUPPORTEDANCH{\char"01A5\char"0349} \SUPPORTEDANCH{\char"0071\char"0349} \SUPPORTEDANCH{\char"02A0\char"0349} \SUPPORTEDANCH{\char"0074\char"0349} \SUPPORTEDANCH{\char"01AB\char"0349} \SUPPORTEDANCH{\char"01AD\char"0349} \SUPPORTEDANCH{\char"0236\char"0349} \SUPPORTEDANCH{\char"0075\char"0349} \MISSINGGLYPH{\char"AB52\char"0349} \SUPPORTEDANCH{\char"0079\char"0349} \SUPPORTEDANCH{\char"01B4\char"0349} \MISSINGGLYPH{\char"1EFF\char"0349}

\SUPPORTEDANCH{\char"0048\char"034D} \SUPPORTEDANCH{\char"2C67\char"034D} \SUPPORTEDANCH{\char"004B\char"034D} \SUPPORTEDANCH{\char"2C69\char"034D} \SUPPORTEDANCH{\char"004D\char"034D} \MISSINGGLYPH{\char"2C6E\char"034D} \SUPPORTEDANCH{\char"0050\char"034D} \SUPPORTEDANCH{\char"01A4\char"034D} \SUPPORTEDANCH{\char"0052\char"034D} \SUPPORTEDANCH{\char"2C64\char"034D} \SUPPORTEDANCH{\char"0054\char"034D} \SUPPORTEDANCH{\char"01AC\char"034D} \SUPPORTEDANCH{\char"01AE\char"034D} \SUPPORTEDANCH{\char"0059\char"034D} \SUPPORTEDANCH{\char"01B3\char"034D} \SUPPORTEDANCH{\char"0061\char"034D} \SUPPORTEDANCH{\char"1E9A\char"034D} \SUPPORTEDANCH{\char"0062\char"034D} \SUPPORTEDANCH{\char"0183\char"034D} \SUPPORTEDANCH{\char"0063\char"034D} \SUPPORTEDANCH{\char"0188\char"034D} \SUPPORTEDANCH{\char"0064\char"034D} \SUPPORTEDANCH{\char"018C\char"034D} \SUPPORTEDANCH{\char"0221\char"034D} \SUPPORTEDANCH{\char"0068\char"034D} \SUPPORTEDANCH{\char"2C68\char"034D} \SUPPORTEDANCH{\char"006B\char"034D} \SUPPORTEDANCH{\char"0199\char"034D} \SUPPORTEDANCH{\char"2C6A\char"034D} \SUPPORTEDANCH{\char"006C\char"034D} \SUPPORTEDANCH{\char"0234\char"034D} \SUPPORTEDANCH{\char"006E\char"034D} \SUPPORTEDANCH{\char"0235\char"034D} \SUPPORTEDANCH{\char"0070\char"034D} \SUPPORTEDANCH{\char"01A5\char"034D} \SUPPORTEDANCH{\char"0071\char"034D} \SUPPORTEDANCH{\char"02A0\char"034D} \SUPPORTEDANCH{\char"0074\char"034D} \SUPPORTEDANCH{\char"01AB\char"034D} \SUPPORTEDANCH{\char"01AD\char"034D} \SUPPORTEDANCH{\char"0236\char"034D} \SUPPORTEDANCH{\char"0075\char"034D} \MISSINGGLYPH{\char"AB52\char"034D} \SUPPORTEDANCH{\char"0079\char"034D} \SUPPORTEDANCH{\char"01B4\char"034D} \MISSINGGLYPH{\char"1EFF\char"034D}

\SUPPORTEDANCH{\char"0048\char"034E} \SUPPORTEDANCH{\char"2C67\char"034E} \SUPPORTEDANCH{\char"004B\char"034E} \SUPPORTEDANCH{\char"2C69\char"034E} \SUPPORTEDANCH{\char"004D\char"034E} \MISSINGGLYPH{\char"2C6E\char"034E} \SUPPORTEDANCH{\char"0050\char"034E} \SUPPORTEDANCH{\char"01A4\char"034E} \SUPPORTEDANCH{\char"0052\char"034E} \SUPPORTEDANCH{\char"2C64\char"034E} \SUPPORTEDANCH{\char"0054\char"034E} \SUPPORTEDANCH{\char"01AC\char"034E} \SUPPORTEDANCH{\char"01AE\char"034E} \SUPPORTEDANCH{\char"0059\char"034E} \SUPPORTEDANCH{\char"01B3\char"034E} \SUPPORTEDANCH{\char"0061\char"034E} \SUPPORTEDANCH{\char"1E9A\char"034E} \SUPPORTEDANCH{\char"0062\char"034E} \SUPPORTEDANCH{\char"0183\char"034E} \SUPPORTEDANCH{\char"0063\char"034E} \SUPPORTEDANCH{\char"0188\char"034E} \SUPPORTEDANCH{\char"0064\char"034E} \SUPPORTEDANCH{\char"018C\char"034E} \SUPPORTEDANCH{\char"0221\char"034E} \SUPPORTEDANCH{\char"0068\char"034E} \SUPPORTEDANCH{\char"2C68\char"034E} \SUPPORTEDANCH{\char"006B\char"034E} \SUPPORTEDANCH{\char"0199\char"034E} \SUPPORTEDANCH{\char"2C6A\char"034E} \SUPPORTEDANCH{\char"006C\char"034E} \SUPPORTEDANCH{\char"0234\char"034E} \SUPPORTEDANCH{\char"006E\char"034E} \SUPPORTEDANCH{\char"0235\char"034E} \SUPPORTEDANCH{\char"0070\char"034E} \SUPPORTEDANCH{\char"01A5\char"034E} \SUPPORTEDANCH{\char"0071\char"034E} \SUPPORTEDANCH{\char"02A0\char"034E} \SUPPORTEDANCH{\char"0074\char"034E} \SUPPORTEDANCH{\char"01AB\char"034E} \SUPPORTEDANCH{\char"01AD\char"034E} \SUPPORTEDANCH{\char"0236\char"034E} \SUPPORTEDANCH{\char"0075\char"034E} \MISSINGGLYPH{\char"AB52\char"034E} \SUPPORTEDANCH{\char"0079\char"034E} \SUPPORTEDANCH{\char"01B4\char"034E} \MISSINGGLYPH{\char"1EFF\char"034E}

}

\newpage
\subsection*{Semibold, Group 4, below, Below}

{\bfseries
% grid. columns are bases, rows are marks.
\SUPPORTEDANCH{\char"0048\char"0316} \SUPPORTEDFALL{\char"2C67\char"0316} \SUPPORTEDFALL{\char"004B\char"0316} \SUPPORTEDFALL{\char"2C69\char"0316} \SUPPORTEDFALL{\char"004D\char"0316} \MISSINGGLYPH{\char"2C6E\char"0316} \SUPPORTEDFALL{\char"0050\char"0316} \SUPPORTEDFALL{\char"01A4\char"0316} \SUPPORTEDANCH{\char"0052\char"0316} \SUPPORTEDFALL{\char"2C64\char"0316} \SUPPORTEDANCH{\char"0054\char"0316} \SUPPORTEDFALL{\char"01AC\char"0316} \SUPPORTEDFALL{\char"01AE\char"0316} \SUPPORTEDANCH{\char"0059\char"0316} \SUPPORTEDFALL{\char"01B3\char"0316} \SUPPORTEDFALL{\char"0061\char"0316} \SUPPORTEDFALL{\char"1E9A\char"0316} \SUPPORTEDFALL{\char"0062\char"0316} \SUPPORTEDFALL{\char"0183\char"0316} \SUPPORTEDFALL{\char"0063\char"0316} \SUPPORTEDFALL{\char"0188\char"0316} \SUPPORTEDFALL{\char"0064\char"0316} \SUPPORTEDFALL{\char"018C\char"0316} \SUPPORTEDFALL{\char"0221\char"0316} \SUPPORTEDANCH{\char"0068\char"0316} \SUPPORTEDFALL{\char"2C68\char"0316} \SUPPORTEDFALL{\char"006B\char"0316} \SUPPORTEDFALL{\char"0199\char"0316} \SUPPORTEDFALL{\char"2C6A\char"0316} \SUPPORTEDANCH{\char"006C\char"0316} \SUPPORTEDFALL{\char"0234\char"0316} \SUPPORTEDANCH{\char"006E\char"0316} \SUPPORTEDFALL{\char"0235\char"0316} \SUPPORTEDFALL{\char"0070\char"0316} \SUPPORTEDFALL{\char"01A5\char"0316} \SUPPORTEDFALL{\char"0071\char"0316} \SUPPORTEDFALL{\char"02A0\char"0316} \SUPPORTEDANCH{\char"0074\char"0316} \SUPPORTEDFALL{\char"01AB\char"0316} \SUPPORTEDFALL{\char"01AD\char"0316} \SUPPORTEDFALL{\char"0236\char"0316} \SUPPORTEDANCH{\char"0075\char"0316} \MISSINGGLYPH{\char"AB52\char"0316} \SUPPORTEDFALL{\char"0079\char"0316} \SUPPORTEDFALL{\char"01B4\char"0316} \MISSINGGLYPH{\char"1EFF\char"0316}

\SUPPORTEDANCH{\char"0048\char"0317} \SUPPORTEDFALL{\char"2C67\char"0317} \SUPPORTEDFALL{\char"004B\char"0317} \SUPPORTEDFALL{\char"2C69\char"0317} \SUPPORTEDFALL{\char"004D\char"0317} \MISSINGGLYPH{\char"2C6E\char"0317} \SUPPORTEDFALL{\char"0050\char"0317} \SUPPORTEDFALL{\char"01A4\char"0317} \SUPPORTEDANCH{\char"0052\char"0317} \SUPPORTEDFALL{\char"2C64\char"0317} \SUPPORTEDANCH{\char"0054\char"0317} \SUPPORTEDFALL{\char"01AC\char"0317} \SUPPORTEDFALL{\char"01AE\char"0317} \SUPPORTEDANCH{\char"0059\char"0317} \SUPPORTEDFALL{\char"01B3\char"0317} \SUPPORTEDFALL{\char"0061\char"0317} \SUPPORTEDFALL{\char"1E9A\char"0317} \SUPPORTEDFALL{\char"0062\char"0317} \SUPPORTEDFALL{\char"0183\char"0317} \SUPPORTEDFALL{\char"0063\char"0317} \SUPPORTEDFALL{\char"0188\char"0317} \SUPPORTEDFALL{\char"0064\char"0317} \SUPPORTEDFALL{\char"018C\char"0317} \SUPPORTEDFALL{\char"0221\char"0317} \SUPPORTEDANCH{\char"0068\char"0317} \SUPPORTEDFALL{\char"2C68\char"0317} \SUPPORTEDFALL{\char"006B\char"0317} \SUPPORTEDFALL{\char"0199\char"0317} \SUPPORTEDFALL{\char"2C6A\char"0317} \SUPPORTEDANCH{\char"006C\char"0317} \SUPPORTEDFALL{\char"0234\char"0317} \SUPPORTEDANCH{\char"006E\char"0317} \SUPPORTEDFALL{\char"0235\char"0317} \SUPPORTEDFALL{\char"0070\char"0317} \SUPPORTEDFALL{\char"01A5\char"0317} \SUPPORTEDFALL{\char"0071\char"0317} \SUPPORTEDFALL{\char"02A0\char"0317} \SUPPORTEDANCH{\char"0074\char"0317} \SUPPORTEDFALL{\char"01AB\char"0317} \SUPPORTEDFALL{\char"01AD\char"0317} \SUPPORTEDFALL{\char"0236\char"0317} \SUPPORTEDANCH{\char"0075\char"0317} \MISSINGGLYPH{\char"AB52\char"0317} \SUPPORTEDFALL{\char"0079\char"0317} \SUPPORTEDFALL{\char"01B4\char"0317} \MISSINGGLYPH{\char"1EFF\char"0317}

\SUPPORTEDANCH{\char"0048\char"0318} \SUPPORTEDFALL{\char"2C67\char"0318} \SUPPORTEDFALL{\char"004B\char"0318} \SUPPORTEDFALL{\char"2C69\char"0318} \SUPPORTEDFALL{\char"004D\char"0318} \MISSINGGLYPH{\char"2C6E\char"0318} \SUPPORTEDFALL{\char"0050\char"0318} \SUPPORTEDFALL{\char"01A4\char"0318} \SUPPORTEDANCH{\char"0052\char"0318} \SUPPORTEDFALL{\char"2C64\char"0318} \SUPPORTEDANCH{\char"0054\char"0318} \SUPPORTEDFALL{\char"01AC\char"0318} \SUPPORTEDFALL{\char"01AE\char"0318} \SUPPORTEDANCH{\char"0059\char"0318} \SUPPORTEDFALL{\char"01B3\char"0318} \SUPPORTEDFALL{\char"0061\char"0318} \SUPPORTEDFALL{\char"1E9A\char"0318} \SUPPORTEDFALL{\char"0062\char"0318} \SUPPORTEDFALL{\char"0183\char"0318} \SUPPORTEDFALL{\char"0063\char"0318} \SUPPORTEDFALL{\char"0188\char"0318} \SUPPORTEDFALL{\char"0064\char"0318} \SUPPORTEDFALL{\char"018C\char"0318} \SUPPORTEDFALL{\char"0221\char"0318} \SUPPORTEDANCH{\char"0068\char"0318} \SUPPORTEDFALL{\char"2C68\char"0318} \SUPPORTEDFALL{\char"006B\char"0318} \SUPPORTEDFALL{\char"0199\char"0318} \SUPPORTEDFALL{\char"2C6A\char"0318} \SUPPORTEDANCH{\char"006C\char"0318} \SUPPORTEDFALL{\char"0234\char"0318} \SUPPORTEDANCH{\char"006E\char"0318} \SUPPORTEDFALL{\char"0235\char"0318} \SUPPORTEDFALL{\char"0070\char"0318} \SUPPORTEDFALL{\char"01A5\char"0318} \SUPPORTEDFALL{\char"0071\char"0318} \SUPPORTEDFALL{\char"02A0\char"0318} \SUPPORTEDANCH{\char"0074\char"0318} \SUPPORTEDFALL{\char"01AB\char"0318} \SUPPORTEDFALL{\char"01AD\char"0318} \SUPPORTEDFALL{\char"0236\char"0318} \SUPPORTEDANCH{\char"0075\char"0318} \MISSINGGLYPH{\char"AB52\char"0318} \SUPPORTEDFALL{\char"0079\char"0318} \SUPPORTEDFALL{\char"01B4\char"0318} \MISSINGGLYPH{\char"1EFF\char"0318}

\SUPPORTEDANCH{\char"0048\char"0319} \SUPPORTEDFALL{\char"2C67\char"0319} \SUPPORTEDFALL{\char"004B\char"0319} \SUPPORTEDFALL{\char"2C69\char"0319} \SUPPORTEDFALL{\char"004D\char"0319} \MISSINGGLYPH{\char"2C6E\char"0319} \SUPPORTEDFALL{\char"0050\char"0319} \SUPPORTEDFALL{\char"01A4\char"0319} \SUPPORTEDANCH{\char"0052\char"0319} \SUPPORTEDFALL{\char"2C64\char"0319} \SUPPORTEDANCH{\char"0054\char"0319} \SUPPORTEDFALL{\char"01AC\char"0319} \SUPPORTEDFALL{\char"01AE\char"0319} \SUPPORTEDANCH{\char"0059\char"0319} \SUPPORTEDFALL{\char"01B3\char"0319} \SUPPORTEDFALL{\char"0061\char"0319} \SUPPORTEDFALL{\char"1E9A\char"0319} \SUPPORTEDFALL{\char"0062\char"0319} \SUPPORTEDFALL{\char"0183\char"0319} \SUPPORTEDFALL{\char"0063\char"0319} \SUPPORTEDFALL{\char"0188\char"0319} \SUPPORTEDFALL{\char"0064\char"0319} \SUPPORTEDFALL{\char"018C\char"0319} \SUPPORTEDFALL{\char"0221\char"0319} \SUPPORTEDANCH{\char"0068\char"0319} \SUPPORTEDFALL{\char"2C68\char"0319} \SUPPORTEDFALL{\char"006B\char"0319} \SUPPORTEDFALL{\char"0199\char"0319} \SUPPORTEDFALL{\char"2C6A\char"0319} \SUPPORTEDANCH{\char"006C\char"0319} \SUPPORTEDFALL{\char"0234\char"0319} \SUPPORTEDANCH{\char"006E\char"0319} \SUPPORTEDFALL{\char"0235\char"0319} \SUPPORTEDFALL{\char"0070\char"0319} \SUPPORTEDFALL{\char"01A5\char"0319} \SUPPORTEDFALL{\char"0071\char"0319} \SUPPORTEDFALL{\char"02A0\char"0319} \SUPPORTEDANCH{\char"0074\char"0319} \SUPPORTEDFALL{\char"01AB\char"0319} \SUPPORTEDFALL{\char"01AD\char"0319} \SUPPORTEDFALL{\char"0236\char"0319} \SUPPORTEDANCH{\char"0075\char"0319} \MISSINGGLYPH{\char"AB52\char"0319} \SUPPORTEDFALL{\char"0079\char"0319} \SUPPORTEDFALL{\char"01B4\char"0319} \MISSINGGLYPH{\char"1EFF\char"0319}

\SUPPORTEDANCH{\char"0048\char"031C} \SUPPORTEDFALL{\char"2C67\char"031C} \SUPPORTEDFALL{\char"004B\char"031C} \SUPPORTEDFALL{\char"2C69\char"031C} \SUPPORTEDFALL{\char"004D\char"031C} \MISSINGGLYPH{\char"2C6E\char"031C} \SUPPORTEDFALL{\char"0050\char"031C} \SUPPORTEDFALL{\char"01A4\char"031C} \SUPPORTEDANCH{\char"0052\char"031C} \SUPPORTEDFALL{\char"2C64\char"031C} \SUPPORTEDANCH{\char"0054\char"031C} \SUPPORTEDFALL{\char"01AC\char"031C} \SUPPORTEDFALL{\char"01AE\char"031C} \SUPPORTEDANCH{\char"0059\char"031C} \SUPPORTEDFALL{\char"01B3\char"031C} \SUPPORTEDFALL{\char"0061\char"031C} \SUPPORTEDFALL{\char"1E9A\char"031C} \SUPPORTEDFALL{\char"0062\char"031C} \SUPPORTEDFALL{\char"0183\char"031C} \SUPPORTEDFALL{\char"0063\char"031C} \SUPPORTEDFALL{\char"0188\char"031C} \SUPPORTEDFALL{\char"0064\char"031C} \SUPPORTEDFALL{\char"018C\char"031C} \SUPPORTEDFALL{\char"0221\char"031C} \SUPPORTEDANCH{\char"0068\char"031C} \SUPPORTEDFALL{\char"2C68\char"031C} \SUPPORTEDFALL{\char"006B\char"031C} \SUPPORTEDFALL{\char"0199\char"031C} \SUPPORTEDFALL{\char"2C6A\char"031C} \SUPPORTEDANCH{\char"006C\char"031C} \SUPPORTEDFALL{\char"0234\char"031C} \SUPPORTEDANCH{\char"006E\char"031C} \SUPPORTEDFALL{\char"0235\char"031C} \SUPPORTEDFALL{\char"0070\char"031C} \SUPPORTEDFALL{\char"01A5\char"031C} \SUPPORTEDFALL{\char"0071\char"031C} \SUPPORTEDFALL{\char"02A0\char"031C} \SUPPORTEDANCH{\char"0074\char"031C} \SUPPORTEDFALL{\char"01AB\char"031C} \SUPPORTEDFALL{\char"01AD\char"031C} \SUPPORTEDFALL{\char"0236\char"031C} \SUPPORTEDANCH{\char"0075\char"031C} \MISSINGGLYPH{\char"AB52\char"031C} \SUPPORTEDFALL{\char"0079\char"031C} \SUPPORTEDFALL{\char"01B4\char"031C} \MISSINGGLYPH{\char"1EFF\char"031C}

\SUPPORTEDANCH{\char"0048\char"0320} \SUPPORTEDFALL{\char"2C67\char"0320} \SUPPORTEDFALL{\char"004B\char"0320} \SUPPORTEDFALL{\char"2C69\char"0320} \SUPPORTEDFALL{\char"004D\char"0320} \MISSINGGLYPH{\char"2C6E\char"0320} \SUPPORTEDFALL{\char"0050\char"0320} \SUPPORTEDFALL{\char"01A4\char"0320} \SUPPORTEDANCH{\char"0052\char"0320} \SUPPORTEDFALL{\char"2C64\char"0320} \SUPPORTEDANCH{\char"0054\char"0320} \SUPPORTEDFALL{\char"01AC\char"0320} \SUPPORTEDFALL{\char"01AE\char"0320} \SUPPORTEDANCH{\char"0059\char"0320} \SUPPORTEDFALL{\char"01B3\char"0320} \SUPPORTEDFALL{\char"0061\char"0320} \SUPPORTEDFALL{\char"1E9A\char"0320} \SUPPORTEDFALL{\char"0062\char"0320} \SUPPORTEDFALL{\char"0183\char"0320} \SUPPORTEDFALL{\char"0063\char"0320} \SUPPORTEDFALL{\char"0188\char"0320} \SUPPORTEDFALL{\char"0064\char"0320} \SUPPORTEDFALL{\char"018C\char"0320} \SUPPORTEDFALL{\char"0221\char"0320} \SUPPORTEDANCH{\char"0068\char"0320} \SUPPORTEDFALL{\char"2C68\char"0320} \SUPPORTEDFALL{\char"006B\char"0320} \SUPPORTEDFALL{\char"0199\char"0320} \SUPPORTEDFALL{\char"2C6A\char"0320} \SUPPORTEDANCH{\char"006C\char"0320} \SUPPORTEDFALL{\char"0234\char"0320} \SUPPORTEDANCH{\char"006E\char"0320} \SUPPORTEDFALL{\char"0235\char"0320} \SUPPORTEDFALL{\char"0070\char"0320} \SUPPORTEDFALL{\char"01A5\char"0320} \SUPPORTEDFALL{\char"0071\char"0320} \SUPPORTEDFALL{\char"02A0\char"0320} \SUPPORTEDANCH{\char"0074\char"0320} \SUPPORTEDFALL{\char"01AB\char"0320} \SUPPORTEDFALL{\char"01AD\char"0320} \SUPPORTEDFALL{\char"0236\char"0320} \SUPPORTEDANCH{\char"0075\char"0320} \MISSINGGLYPH{\char"AB52\char"0320} \SUPPORTEDFALL{\char"0079\char"0320} \SUPPORTEDFALL{\char"01B4\char"0320} \MISSINGGLYPH{\char"1EFF\char"0320}

\SUPPORTEDPREC{\char"0048\char"0323} \SUPPORTEDFALL{\char"2C67\char"0323} \SUPPORTEDPREC{\char"004B\char"0323} \SUPPORTEDFALL{\char"2C69\char"0323} \SUPPORTEDPREC{\char"004D\char"0323} \MISSINGGLYPH{\char"2C6E\char"0323} \SUPPORTEDFALL{\char"0050\char"0323} \SUPPORTEDFALL{\char"01A4\char"0323} \SUPPORTEDPREC{\char"0052\char"0323} \SUPPORTEDFALL{\char"2C64\char"0323} \SUPPORTEDPREC{\char"0054\char"0323} \SUPPORTEDFALL{\char"01AC\char"0323} \SUPPORTEDFALL{\char"01AE\char"0323} \SUPPORTEDPREC{\char"0059\char"0323} \SUPPORTEDFALL{\char"01B3\char"0323} \SUPPORTEDPREC{\char"0061\char"0323} \SUPPORTEDFALL{\char"1E9A\char"0323} \SUPPORTEDPREC{\char"0062\char"0323} \SUPPORTEDFALL{\char"0183\char"0323} \SUPPORTEDFALL{\char"0063\char"0323} \SUPPORTEDFALL{\char"0188\char"0323} \SUPPORTEDPREC{\char"0064\char"0323} \SUPPORTEDFALL{\char"018C\char"0323} \SUPPORTEDFALL{\char"0221\char"0323} \SUPPORTEDPREC{\char"0068\char"0323} \SUPPORTEDFALL{\char"2C68\char"0323} \SUPPORTEDPREC{\char"006B\char"0323} \SUPPORTEDFALL{\char"0199\char"0323} \SUPPORTEDFALL{\char"2C6A\char"0323} \SUPPORTEDPREC{\char"006C\char"0323} \SUPPORTEDFALL{\char"0234\char"0323} \SUPPORTEDPREC{\char"006E\char"0323} \SUPPORTEDFALL{\char"0235\char"0323} \SUPPORTEDFALL{\char"0070\char"0323} \SUPPORTEDFALL{\char"01A5\char"0323} \SUPPORTEDFALL{\char"0071\char"0323} \SUPPORTEDFALL{\char"02A0\char"0323} \SUPPORTEDPREC{\char"0074\char"0323} \SUPPORTEDFALL{\char"01AB\char"0323} \SUPPORTEDFALL{\char"01AD\char"0323} \SUPPORTEDFALL{\char"0236\char"0323} \SUPPORTEDPREC{\char"0075\char"0323} \MISSINGGLYPH{\char"AB52\char"0323} \SUPPORTEDPREC{\char"0079\char"0323} \SUPPORTEDFALL{\char"01B4\char"0323} \MISSINGGLYPH{\char"1EFF\char"0323}

\SUPPORTEDANCH{\char"0048\char"0324} \SUPPORTEDFALL{\char"2C67\char"0324} \SUPPORTEDFALL{\char"004B\char"0324} \SUPPORTEDFALL{\char"2C69\char"0324} \SUPPORTEDFALL{\char"004D\char"0324} \MISSINGGLYPH{\char"2C6E\char"0324} \SUPPORTEDFALL{\char"0050\char"0324} \SUPPORTEDFALL{\char"01A4\char"0324} \SUPPORTEDANCH{\char"0052\char"0324} \SUPPORTEDFALL{\char"2C64\char"0324} \SUPPORTEDANCH{\char"0054\char"0324} \SUPPORTEDFALL{\char"01AC\char"0324} \SUPPORTEDFALL{\char"01AE\char"0324} \SUPPORTEDANCH{\char"0059\char"0324} \SUPPORTEDFALL{\char"01B3\char"0324} \SUPPORTEDFALL{\char"0061\char"0324} \SUPPORTEDFALL{\char"1E9A\char"0324} \SUPPORTEDFALL{\char"0062\char"0324} \SUPPORTEDFALL{\char"0183\char"0324} \SUPPORTEDFALL{\char"0063\char"0324} \SUPPORTEDFALL{\char"0188\char"0324} \SUPPORTEDFALL{\char"0064\char"0324} \SUPPORTEDFALL{\char"018C\char"0324} \SUPPORTEDFALL{\char"0221\char"0324} \SUPPORTEDANCH{\char"0068\char"0324} \SUPPORTEDFALL{\char"2C68\char"0324} \SUPPORTEDFALL{\char"006B\char"0324} \SUPPORTEDFALL{\char"0199\char"0324} \SUPPORTEDFALL{\char"2C6A\char"0324} \SUPPORTEDANCH{\char"006C\char"0324} \SUPPORTEDFALL{\char"0234\char"0324} \SUPPORTEDANCH{\char"006E\char"0324} \SUPPORTEDFALL{\char"0235\char"0324} \SUPPORTEDFALL{\char"0070\char"0324} \SUPPORTEDFALL{\char"01A5\char"0324} \SUPPORTEDFALL{\char"0071\char"0324} \SUPPORTEDFALL{\char"02A0\char"0324} \SUPPORTEDANCH{\char"0074\char"0324} \SUPPORTEDFALL{\char"01AB\char"0324} \SUPPORTEDFALL{\char"01AD\char"0324} \SUPPORTEDFALL{\char"0236\char"0324} \SUPPORTEDPREC{\char"0075\char"0324} \MISSINGGLYPH{\char"AB52\char"0324} \SUPPORTEDFALL{\char"0079\char"0324} \SUPPORTEDFALL{\char"01B4\char"0324} \MISSINGGLYPH{\char"1EFF\char"0324}

\SUPPORTEDANCH{\char"0048\char"0325} \SUPPORTEDFALL{\char"2C67\char"0325} \SUPPORTEDFALL{\char"004B\char"0325} \SUPPORTEDFALL{\char"2C69\char"0325} \SUPPORTEDFALL{\char"004D\char"0325} \MISSINGGLYPH{\char"2C6E\char"0325} \SUPPORTEDFALL{\char"0050\char"0325} \SUPPORTEDFALL{\char"01A4\char"0325} \SUPPORTEDANCH{\char"0052\char"0325} \SUPPORTEDFALL{\char"2C64\char"0325} \SUPPORTEDANCH{\char"0054\char"0325} \SUPPORTEDFALL{\char"01AC\char"0325} \SUPPORTEDFALL{\char"01AE\char"0325} \SUPPORTEDANCH{\char"0059\char"0325} \SUPPORTEDFALL{\char"01B3\char"0325} \SUPPORTEDPREC{\char"0061\char"0325} \SUPPORTEDFALL{\char"1E9A\char"0325} \SUPPORTEDFALL{\char"0062\char"0325} \SUPPORTEDFALL{\char"0183\char"0325} \SUPPORTEDFALL{\char"0063\char"0325} \SUPPORTEDFALL{\char"0188\char"0325} \SUPPORTEDFALL{\char"0064\char"0325} \SUPPORTEDFALL{\char"018C\char"0325} \SUPPORTEDFALL{\char"0221\char"0325} \SUPPORTEDANCH{\char"0068\char"0325} \SUPPORTEDFALL{\char"2C68\char"0325} \SUPPORTEDFALL{\char"006B\char"0325} \SUPPORTEDFALL{\char"0199\char"0325} \SUPPORTEDFALL{\char"2C6A\char"0325} \SUPPORTEDANCH{\char"006C\char"0325} \SUPPORTEDFALL{\char"0234\char"0325} \SUPPORTEDANCH{\char"006E\char"0325} \SUPPORTEDFALL{\char"0235\char"0325} \SUPPORTEDFALL{\char"0070\char"0325} \SUPPORTEDFALL{\char"01A5\char"0325} \SUPPORTEDFALL{\char"0071\char"0325} \SUPPORTEDFALL{\char"02A0\char"0325} \SUPPORTEDANCH{\char"0074\char"0325} \SUPPORTEDFALL{\char"01AB\char"0325} \SUPPORTEDFALL{\char"01AD\char"0325} \SUPPORTEDFALL{\char"0236\char"0325} \SUPPORTEDANCH{\char"0075\char"0325} \MISSINGGLYPH{\char"AB52\char"0325} \SUPPORTEDFALL{\char"0079\char"0325} \SUPPORTEDFALL{\char"01B4\char"0325} \MISSINGGLYPH{\char"1EFF\char"0325}

\SUPPORTEDANCH{\char"0048\char"0326} \SUPPORTEDFALL{\char"2C67\char"0326} \SUPPORTEDFALL{\char"004B\char"0326} \SUPPORTEDFALL{\char"2C69\char"0326} \SUPPORTEDFALL{\char"004D\char"0326} \MISSINGGLYPH{\char"2C6E\char"0326} \SUPPORTEDFALL{\char"0050\char"0326} \SUPPORTEDFALL{\char"01A4\char"0326} \SUPPORTEDANCH{\char"0052\char"0326} \SUPPORTEDFALL{\char"2C64\char"0326} \SUPPORTEDPREC{\char"0054\char"0326} \SUPPORTEDFALL{\char"01AC\char"0326} \SUPPORTEDFALL{\char"01AE\char"0326} \SUPPORTEDANCH{\char"0059\char"0326} \SUPPORTEDFALL{\char"01B3\char"0326} \SUPPORTEDFALL{\char"0061\char"0326} \SUPPORTEDFALL{\char"1E9A\char"0326} \SUPPORTEDFALL{\char"0062\char"0326} \SUPPORTEDFALL{\char"0183\char"0326} \SUPPORTEDFALL{\char"0063\char"0326} \SUPPORTEDFALL{\char"0188\char"0326} \SUPPORTEDFALL{\char"0064\char"0326} \SUPPORTEDFALL{\char"018C\char"0326} \SUPPORTEDFALL{\char"0221\char"0326} \SUPPORTEDANCH{\char"0068\char"0326} \SUPPORTEDFALL{\char"2C68\char"0326} \SUPPORTEDFALL{\char"006B\char"0326} \SUPPORTEDFALL{\char"0199\char"0326} \SUPPORTEDFALL{\char"2C6A\char"0326} \SUPPORTEDANCH{\char"006C\char"0326} \SUPPORTEDFALL{\char"0234\char"0326} \SUPPORTEDANCH{\char"006E\char"0326} \SUPPORTEDFALL{\char"0235\char"0326} \SUPPORTEDFALL{\char"0070\char"0326} \SUPPORTEDFALL{\char"01A5\char"0326} \SUPPORTEDFALL{\char"0071\char"0326} \SUPPORTEDFALL{\char"02A0\char"0326} \SUPPORTEDPREC{\char"0074\char"0326} \SUPPORTEDFALL{\char"01AB\char"0326} \SUPPORTEDFALL{\char"01AD\char"0326} \SUPPORTEDFALL{\char"0236\char"0326} \SUPPORTEDANCH{\char"0075\char"0326} \MISSINGGLYPH{\char"AB52\char"0326} \SUPPORTEDFALL{\char"0079\char"0326} \SUPPORTEDFALL{\char"01B4\char"0326} \MISSINGGLYPH{\char"1EFF\char"0326}

\SUPPORTEDANCH{\char"0048\char"0329} \SUPPORTEDFALL{\char"2C67\char"0329} \SUPPORTEDFALL{\char"004B\char"0329} \SUPPORTEDFALL{\char"2C69\char"0329} \SUPPORTEDFALL{\char"004D\char"0329} \MISSINGGLYPH{\char"2C6E\char"0329} \SUPPORTEDFALL{\char"0050\char"0329} \SUPPORTEDFALL{\char"01A4\char"0329} \SUPPORTEDANCH{\char"0052\char"0329} \SUPPORTEDFALL{\char"2C64\char"0329} \SUPPORTEDANCH{\char"0054\char"0329} \SUPPORTEDFALL{\char"01AC\char"0329} \SUPPORTEDFALL{\char"01AE\char"0329} \SUPPORTEDANCH{\char"0059\char"0329} \SUPPORTEDFALL{\char"01B3\char"0329} \SUPPORTEDFALL{\char"0061\char"0329} \SUPPORTEDFALL{\char"1E9A\char"0329} \SUPPORTEDFALL{\char"0062\char"0329} \SUPPORTEDFALL{\char"0183\char"0329} \SUPPORTEDFALL{\char"0063\char"0329} \SUPPORTEDFALL{\char"0188\char"0329} \SUPPORTEDFALL{\char"0064\char"0329} \SUPPORTEDFALL{\char"018C\char"0329} \SUPPORTEDFALL{\char"0221\char"0329} \SUPPORTEDANCH{\char"0068\char"0329} \SUPPORTEDFALL{\char"2C68\char"0329} \SUPPORTEDFALL{\char"006B\char"0329} \SUPPORTEDFALL{\char"0199\char"0329} \SUPPORTEDFALL{\char"2C6A\char"0329} \SUPPORTEDANCH{\char"006C\char"0329} \SUPPORTEDFALL{\char"0234\char"0329} \SUPPORTEDANCH{\char"006E\char"0329} \SUPPORTEDFALL{\char"0235\char"0329} \SUPPORTEDFALL{\char"0070\char"0329} \SUPPORTEDFALL{\char"01A5\char"0329} \SUPPORTEDFALL{\char"0071\char"0329} \SUPPORTEDFALL{\char"02A0\char"0329} \SUPPORTEDANCH{\char"0074\char"0329} \SUPPORTEDFALL{\char"01AB\char"0329} \SUPPORTEDFALL{\char"01AD\char"0329} \SUPPORTEDFALL{\char"0236\char"0329} \SUPPORTEDANCH{\char"0075\char"0329} \MISSINGGLYPH{\char"AB52\char"0329} \SUPPORTEDFALL{\char"0079\char"0329} \SUPPORTEDFALL{\char"01B4\char"0329} \MISSINGGLYPH{\char"1EFF\char"0329}

\SUPPORTEDANCH{\char"0048\char"032A} \SUPPORTEDFALL{\char"2C67\char"032A} \SUPPORTEDFALL{\char"004B\char"032A} \SUPPORTEDFALL{\char"2C69\char"032A} \SUPPORTEDFALL{\char"004D\char"032A} \MISSINGGLYPH{\char"2C6E\char"032A} \SUPPORTEDFALL{\char"0050\char"032A} \SUPPORTEDFALL{\char"01A4\char"032A} \SUPPORTEDANCH{\char"0052\char"032A} \SUPPORTEDFALL{\char"2C64\char"032A} \SUPPORTEDANCH{\char"0054\char"032A} \SUPPORTEDFALL{\char"01AC\char"032A} \SUPPORTEDFALL{\char"01AE\char"032A} \SUPPORTEDANCH{\char"0059\char"032A} \SUPPORTEDFALL{\char"01B3\char"032A} \SUPPORTEDFALL{\char"0061\char"032A} \SUPPORTEDFALL{\char"1E9A\char"032A} \SUPPORTEDFALL{\char"0062\char"032A} \SUPPORTEDFALL{\char"0183\char"032A} \SUPPORTEDFALL{\char"0063\char"032A} \SUPPORTEDFALL{\char"0188\char"032A} \SUPPORTEDFALL{\char"0064\char"032A} \SUPPORTEDFALL{\char"018C\char"032A} \SUPPORTEDFALL{\char"0221\char"032A} \SUPPORTEDANCH{\char"0068\char"032A} \SUPPORTEDFALL{\char"2C68\char"032A} \SUPPORTEDFALL{\char"006B\char"032A} \SUPPORTEDFALL{\char"0199\char"032A} \SUPPORTEDFALL{\char"2C6A\char"032A} \SUPPORTEDANCH{\char"006C\char"032A} \SUPPORTEDFALL{\char"0234\char"032A} \SUPPORTEDANCH{\char"006E\char"032A} \SUPPORTEDFALL{\char"0235\char"032A} \SUPPORTEDFALL{\char"0070\char"032A} \SUPPORTEDFALL{\char"01A5\char"032A} \SUPPORTEDFALL{\char"0071\char"032A} \SUPPORTEDFALL{\char"02A0\char"032A} \SUPPORTEDANCH{\char"0074\char"032A} \SUPPORTEDFALL{\char"01AB\char"032A} \SUPPORTEDFALL{\char"01AD\char"032A} \SUPPORTEDFALL{\char"0236\char"032A} \SUPPORTEDANCH{\char"0075\char"032A} \MISSINGGLYPH{\char"AB52\char"032A} \SUPPORTEDFALL{\char"0079\char"032A} \SUPPORTEDFALL{\char"01B4\char"032A} \MISSINGGLYPH{\char"1EFF\char"032A}

\SUPPORTEDANCH{\char"0048\char"032B} \SUPPORTEDFALL{\char"2C67\char"032B} \SUPPORTEDFALL{\char"004B\char"032B} \SUPPORTEDFALL{\char"2C69\char"032B} \SUPPORTEDFALL{\char"004D\char"032B} \MISSINGGLYPH{\char"2C6E\char"032B} \SUPPORTEDFALL{\char"0050\char"032B} \SUPPORTEDFALL{\char"01A4\char"032B} \SUPPORTEDANCH{\char"0052\char"032B} \SUPPORTEDFALL{\char"2C64\char"032B} \SUPPORTEDANCH{\char"0054\char"032B} \SUPPORTEDFALL{\char"01AC\char"032B} \SUPPORTEDFALL{\char"01AE\char"032B} \SUPPORTEDANCH{\char"0059\char"032B} \SUPPORTEDFALL{\char"01B3\char"032B} \SUPPORTEDFALL{\char"0061\char"032B} \SUPPORTEDFALL{\char"1E9A\char"032B} \SUPPORTEDFALL{\char"0062\char"032B} \SUPPORTEDFALL{\char"0183\char"032B} \SUPPORTEDFALL{\char"0063\char"032B} \SUPPORTEDFALL{\char"0188\char"032B} \SUPPORTEDFALL{\char"0064\char"032B} \SUPPORTEDFALL{\char"018C\char"032B} \SUPPORTEDFALL{\char"0221\char"032B} \SUPPORTEDANCH{\char"0068\char"032B} \SUPPORTEDFALL{\char"2C68\char"032B} \SUPPORTEDFALL{\char"006B\char"032B} \SUPPORTEDFALL{\char"0199\char"032B} \SUPPORTEDFALL{\char"2C6A\char"032B} \SUPPORTEDANCH{\char"006C\char"032B} \SUPPORTEDFALL{\char"0234\char"032B} \SUPPORTEDANCH{\char"006E\char"032B} \SUPPORTEDFALL{\char"0235\char"032B} \SUPPORTEDFALL{\char"0070\char"032B} \SUPPORTEDFALL{\char"01A5\char"032B} \SUPPORTEDFALL{\char"0071\char"032B} \SUPPORTEDFALL{\char"02A0\char"032B} \SUPPORTEDANCH{\char"0074\char"032B} \SUPPORTEDFALL{\char"01AB\char"032B} \SUPPORTEDFALL{\char"01AD\char"032B} \SUPPORTEDFALL{\char"0236\char"032B} \SUPPORTEDANCH{\char"0075\char"032B} \MISSINGGLYPH{\char"AB52\char"032B} \SUPPORTEDFALL{\char"0079\char"032B} \SUPPORTEDFALL{\char"01B4\char"032B} \MISSINGGLYPH{\char"1EFF\char"032B}

\SUPPORTEDANCH{\char"0048\char"032C} \SUPPORTEDFALL{\char"2C67\char"032C} \SUPPORTEDFALL{\char"004B\char"032C} \SUPPORTEDFALL{\char"2C69\char"032C} \SUPPORTEDFALL{\char"004D\char"032C} \MISSINGGLYPH{\char"2C6E\char"032C} \SUPPORTEDFALL{\char"0050\char"032C} \SUPPORTEDFALL{\char"01A4\char"032C} \SUPPORTEDANCH{\char"0052\char"032C} \SUPPORTEDFALL{\char"2C64\char"032C} \SUPPORTEDANCH{\char"0054\char"032C} \SUPPORTEDFALL{\char"01AC\char"032C} \SUPPORTEDFALL{\char"01AE\char"032C} \SUPPORTEDANCH{\char"0059\char"032C} \SUPPORTEDFALL{\char"01B3\char"032C} \SUPPORTEDFALL{\char"0061\char"032C} \SUPPORTEDFALL{\char"1E9A\char"032C} \SUPPORTEDFALL{\char"0062\char"032C} \SUPPORTEDFALL{\char"0183\char"032C} \SUPPORTEDFALL{\char"0063\char"032C} \SUPPORTEDFALL{\char"0188\char"032C} \SUPPORTEDFALL{\char"0064\char"032C} \SUPPORTEDFALL{\char"018C\char"032C} \SUPPORTEDFALL{\char"0221\char"032C} \SUPPORTEDANCH{\char"0068\char"032C} \SUPPORTEDFALL{\char"2C68\char"032C} \SUPPORTEDFALL{\char"006B\char"032C} \SUPPORTEDFALL{\char"0199\char"032C} \SUPPORTEDFALL{\char"2C6A\char"032C} \SUPPORTEDANCH{\char"006C\char"032C} \SUPPORTEDFALL{\char"0234\char"032C} \SUPPORTEDANCH{\char"006E\char"032C} \SUPPORTEDFALL{\char"0235\char"032C} \SUPPORTEDFALL{\char"0070\char"032C} \SUPPORTEDFALL{\char"01A5\char"032C} \SUPPORTEDFALL{\char"0071\char"032C} \SUPPORTEDFALL{\char"02A0\char"032C} \SUPPORTEDANCH{\char"0074\char"032C} \SUPPORTEDFALL{\char"01AB\char"032C} \SUPPORTEDFALL{\char"01AD\char"032C} \SUPPORTEDFALL{\char"0236\char"032C} \SUPPORTEDANCH{\char"0075\char"032C} \MISSINGGLYPH{\char"AB52\char"032C} \SUPPORTEDFALL{\char"0079\char"032C} \SUPPORTEDFALL{\char"01B4\char"032C} \MISSINGGLYPH{\char"1EFF\char"032C}

\SUPPORTEDANCH{\char"0048\char"032D} \SUPPORTEDFALL{\char"2C67\char"032D} \SUPPORTEDFALL{\char"004B\char"032D} \SUPPORTEDFALL{\char"2C69\char"032D} \SUPPORTEDFALL{\char"004D\char"032D} \MISSINGGLYPH{\char"2C6E\char"032D} \SUPPORTEDFALL{\char"0050\char"032D} \SUPPORTEDFALL{\char"01A4\char"032D} \SUPPORTEDANCH{\char"0052\char"032D} \SUPPORTEDFALL{\char"2C64\char"032D} \SUPPORTEDPREC{\char"0054\char"032D} \SUPPORTEDFALL{\char"01AC\char"032D} \SUPPORTEDFALL{\char"01AE\char"032D} \SUPPORTEDANCH{\char"0059\char"032D} \SUPPORTEDFALL{\char"01B3\char"032D} \SUPPORTEDFALL{\char"0061\char"032D} \SUPPORTEDFALL{\char"1E9A\char"032D} \SUPPORTEDFALL{\char"0062\char"032D} \SUPPORTEDFALL{\char"0183\char"032D} \SUPPORTEDFALL{\char"0063\char"032D} \SUPPORTEDFALL{\char"0188\char"032D} \SUPPORTEDPREC{\char"0064\char"032D} \SUPPORTEDFALL{\char"018C\char"032D} \SUPPORTEDFALL{\char"0221\char"032D} \SUPPORTEDANCH{\char"0068\char"032D} \SUPPORTEDFALL{\char"2C68\char"032D} \SUPPORTEDFALL{\char"006B\char"032D} \SUPPORTEDFALL{\char"0199\char"032D} \SUPPORTEDFALL{\char"2C6A\char"032D} \SUPPORTEDPREC{\char"006C\char"032D} \SUPPORTEDFALL{\char"0234\char"032D} \SUPPORTEDPREC{\char"006E\char"032D} \SUPPORTEDFALL{\char"0235\char"032D} \SUPPORTEDFALL{\char"0070\char"032D} \SUPPORTEDFALL{\char"01A5\char"032D} \SUPPORTEDFALL{\char"0071\char"032D} \SUPPORTEDFALL{\char"02A0\char"032D} \SUPPORTEDPREC{\char"0074\char"032D} \SUPPORTEDFALL{\char"01AB\char"032D} \SUPPORTEDFALL{\char"01AD\char"032D} \SUPPORTEDFALL{\char"0236\char"032D} \SUPPORTEDPREC{\char"0075\char"032D} \MISSINGGLYPH{\char"AB52\char"032D} \SUPPORTEDFALL{\char"0079\char"032D} \SUPPORTEDFALL{\char"01B4\char"032D} \MISSINGGLYPH{\char"1EFF\char"032D}

\SUPPORTEDPREC{\char"0048\char"032E} \SUPPORTEDFALL{\char"2C67\char"032E} \SUPPORTEDFALL{\char"004B\char"032E} \SUPPORTEDFALL{\char"2C69\char"032E} \SUPPORTEDFALL{\char"004D\char"032E} \MISSINGGLYPH{\char"2C6E\char"032E} \SUPPORTEDFALL{\char"0050\char"032E} \SUPPORTEDFALL{\char"01A4\char"032E} \SUPPORTEDANCH{\char"0052\char"032E} \SUPPORTEDFALL{\char"2C64\char"032E} \SUPPORTEDANCH{\char"0054\char"032E} \SUPPORTEDFALL{\char"01AC\char"032E} \SUPPORTEDFALL{\char"01AE\char"032E} \SUPPORTEDANCH{\char"0059\char"032E} \SUPPORTEDFALL{\char"01B3\char"032E} \SUPPORTEDFALL{\char"0061\char"032E} \SUPPORTEDFALL{\char"1E9A\char"032E} \SUPPORTEDFALL{\char"0062\char"032E} \SUPPORTEDFALL{\char"0183\char"032E} \SUPPORTEDFALL{\char"0063\char"032E} \SUPPORTEDFALL{\char"0188\char"032E} \SUPPORTEDFALL{\char"0064\char"032E} \SUPPORTEDFALL{\char"018C\char"032E} \SUPPORTEDFALL{\char"0221\char"032E} \SUPPORTEDPREC{\char"0068\char"032E} \SUPPORTEDFALL{\char"2C68\char"032E} \SUPPORTEDFALL{\char"006B\char"032E} \SUPPORTEDFALL{\char"0199\char"032E} \SUPPORTEDFALL{\char"2C6A\char"032E} \SUPPORTEDANCH{\char"006C\char"032E} \SUPPORTEDFALL{\char"0234\char"032E} \SUPPORTEDANCH{\char"006E\char"032E} \SUPPORTEDFALL{\char"0235\char"032E} \SUPPORTEDFALL{\char"0070\char"032E} \SUPPORTEDFALL{\char"01A5\char"032E} \SUPPORTEDFALL{\char"0071\char"032E} \SUPPORTEDFALL{\char"02A0\char"032E} \SUPPORTEDANCH{\char"0074\char"032E} \SUPPORTEDFALL{\char"01AB\char"032E} \SUPPORTEDFALL{\char"01AD\char"032E} \SUPPORTEDFALL{\char"0236\char"032E} \SUPPORTEDANCH{\char"0075\char"032E} \MISSINGGLYPH{\char"AB52\char"032E} \SUPPORTEDFALL{\char"0079\char"032E} \SUPPORTEDFALL{\char"01B4\char"032E} \MISSINGGLYPH{\char"1EFF\char"032E}

\SUPPORTEDANCH{\char"0048\char"032F} \SUPPORTEDFALL{\char"2C67\char"032F} \SUPPORTEDFALL{\char"004B\char"032F} \SUPPORTEDFALL{\char"2C69\char"032F} \SUPPORTEDFALL{\char"004D\char"032F} \MISSINGGLYPH{\char"2C6E\char"032F} \SUPPORTEDFALL{\char"0050\char"032F} \SUPPORTEDFALL{\char"01A4\char"032F} \SUPPORTEDANCH{\char"0052\char"032F} \SUPPORTEDFALL{\char"2C64\char"032F} \SUPPORTEDANCH{\char"0054\char"032F} \SUPPORTEDFALL{\char"01AC\char"032F} \SUPPORTEDFALL{\char"01AE\char"032F} \SUPPORTEDANCH{\char"0059\char"032F} \SUPPORTEDFALL{\char"01B3\char"032F} \SUPPORTEDFALL{\char"0061\char"032F} \SUPPORTEDFALL{\char"1E9A\char"032F} \SUPPORTEDFALL{\char"0062\char"032F} \SUPPORTEDFALL{\char"0183\char"032F} \SUPPORTEDFALL{\char"0063\char"032F} \SUPPORTEDFALL{\char"0188\char"032F} \SUPPORTEDFALL{\char"0064\char"032F} \SUPPORTEDFALL{\char"018C\char"032F} \SUPPORTEDFALL{\char"0221\char"032F} \SUPPORTEDANCH{\char"0068\char"032F} \SUPPORTEDFALL{\char"2C68\char"032F} \SUPPORTEDFALL{\char"006B\char"032F} \SUPPORTEDFALL{\char"0199\char"032F} \SUPPORTEDFALL{\char"2C6A\char"032F} \SUPPORTEDANCH{\char"006C\char"032F} \SUPPORTEDFALL{\char"0234\char"032F} \SUPPORTEDANCH{\char"006E\char"032F} \SUPPORTEDFALL{\char"0235\char"032F} \SUPPORTEDFALL{\char"0070\char"032F} \SUPPORTEDFALL{\char"01A5\char"032F} \SUPPORTEDFALL{\char"0071\char"032F} \SUPPORTEDFALL{\char"02A0\char"032F} \SUPPORTEDANCH{\char"0074\char"032F} \SUPPORTEDFALL{\char"01AB\char"032F} \SUPPORTEDFALL{\char"01AD\char"032F} \SUPPORTEDFALL{\char"0236\char"032F} \SUPPORTEDANCH{\char"0075\char"032F} \MISSINGGLYPH{\char"AB52\char"032F} \SUPPORTEDFALL{\char"0079\char"032F} \SUPPORTEDFALL{\char"01B4\char"032F} \MISSINGGLYPH{\char"1EFF\char"032F}

\SUPPORTEDANCH{\char"0048\char"031D} \SUPPORTEDFALL{\char"2C67\char"031D} \SUPPORTEDFALL{\char"004B\char"031D} \SUPPORTEDFALL{\char"2C69\char"031D} \SUPPORTEDFALL{\char"004D\char"031D} \MISSINGGLYPH{\char"2C6E\char"031D} \SUPPORTEDFALL{\char"0050\char"031D} \SUPPORTEDFALL{\char"01A4\char"031D} \SUPPORTEDANCH{\char"0052\char"031D} \SUPPORTEDFALL{\char"2C64\char"031D} \SUPPORTEDANCH{\char"0054\char"031D} \SUPPORTEDFALL{\char"01AC\char"031D} \SUPPORTEDFALL{\char"01AE\char"031D} \SUPPORTEDANCH{\char"0059\char"031D} \SUPPORTEDFALL{\char"01B3\char"031D} \SUPPORTEDFALL{\char"0061\char"031D} \SUPPORTEDFALL{\char"1E9A\char"031D} \SUPPORTEDFALL{\char"0062\char"031D} \SUPPORTEDFALL{\char"0183\char"031D} \SUPPORTEDFALL{\char"0063\char"031D} \SUPPORTEDFALL{\char"0188\char"031D} \SUPPORTEDFALL{\char"0064\char"031D} \SUPPORTEDFALL{\char"018C\char"031D} \SUPPORTEDFALL{\char"0221\char"031D} \SUPPORTEDANCH{\char"0068\char"031D} \SUPPORTEDFALL{\char"2C68\char"031D} \SUPPORTEDFALL{\char"006B\char"031D} \SUPPORTEDFALL{\char"0199\char"031D} \SUPPORTEDFALL{\char"2C6A\char"031D} \SUPPORTEDANCH{\char"006C\char"031D} \SUPPORTEDFALL{\char"0234\char"031D} \SUPPORTEDANCH{\char"006E\char"031D} \SUPPORTEDFALL{\char"0235\char"031D} \SUPPORTEDFALL{\char"0070\char"031D} \SUPPORTEDFALL{\char"01A5\char"031D} \SUPPORTEDFALL{\char"0071\char"031D} \SUPPORTEDFALL{\char"02A0\char"031D} \SUPPORTEDANCH{\char"0074\char"031D} \SUPPORTEDFALL{\char"01AB\char"031D} \SUPPORTEDFALL{\char"01AD\char"031D} \SUPPORTEDFALL{\char"0236\char"031D} \SUPPORTEDANCH{\char"0075\char"031D} \MISSINGGLYPH{\char"AB52\char"031D} \SUPPORTEDFALL{\char"0079\char"031D} \SUPPORTEDFALL{\char"01B4\char"031D} \MISSINGGLYPH{\char"1EFF\char"031D}

\SUPPORTEDANCH{\char"0048\char"031E} \SUPPORTEDFALL{\char"2C67\char"031E} \SUPPORTEDFALL{\char"004B\char"031E} \SUPPORTEDFALL{\char"2C69\char"031E} \SUPPORTEDFALL{\char"004D\char"031E} \MISSINGGLYPH{\char"2C6E\char"031E} \SUPPORTEDFALL{\char"0050\char"031E} \SUPPORTEDFALL{\char"01A4\char"031E} \SUPPORTEDANCH{\char"0052\char"031E} \SUPPORTEDFALL{\char"2C64\char"031E} \SUPPORTEDANCH{\char"0054\char"031E} \SUPPORTEDFALL{\char"01AC\char"031E} \SUPPORTEDFALL{\char"01AE\char"031E} \SUPPORTEDANCH{\char"0059\char"031E} \SUPPORTEDFALL{\char"01B3\char"031E} \SUPPORTEDFALL{\char"0061\char"031E} \SUPPORTEDFALL{\char"1E9A\char"031E} \SUPPORTEDFALL{\char"0062\char"031E} \SUPPORTEDFALL{\char"0183\char"031E} \SUPPORTEDFALL{\char"0063\char"031E} \SUPPORTEDFALL{\char"0188\char"031E} \SUPPORTEDFALL{\char"0064\char"031E} \SUPPORTEDFALL{\char"018C\char"031E} \SUPPORTEDFALL{\char"0221\char"031E} \SUPPORTEDANCH{\char"0068\char"031E} \SUPPORTEDFALL{\char"2C68\char"031E} \SUPPORTEDFALL{\char"006B\char"031E} \SUPPORTEDFALL{\char"0199\char"031E} \SUPPORTEDFALL{\char"2C6A\char"031E} \SUPPORTEDANCH{\char"006C\char"031E} \SUPPORTEDFALL{\char"0234\char"031E} \SUPPORTEDANCH{\char"006E\char"031E} \SUPPORTEDFALL{\char"0235\char"031E} \SUPPORTEDFALL{\char"0070\char"031E} \SUPPORTEDFALL{\char"01A5\char"031E} \SUPPORTEDFALL{\char"0071\char"031E} \SUPPORTEDFALL{\char"02A0\char"031E} \SUPPORTEDANCH{\char"0074\char"031E} \SUPPORTEDFALL{\char"01AB\char"031E} \SUPPORTEDFALL{\char"01AD\char"031E} \SUPPORTEDFALL{\char"0236\char"031E} \SUPPORTEDANCH{\char"0075\char"031E} \MISSINGGLYPH{\char"AB52\char"031E} \SUPPORTEDFALL{\char"0079\char"031E} \SUPPORTEDFALL{\char"01B4\char"031E} \MISSINGGLYPH{\char"1EFF\char"031E}

\SUPPORTEDANCH{\char"0048\char"031F} \SUPPORTEDFALL{\char"2C67\char"031F} \SUPPORTEDFALL{\char"004B\char"031F} \SUPPORTEDFALL{\char"2C69\char"031F} \SUPPORTEDFALL{\char"004D\char"031F} \MISSINGGLYPH{\char"2C6E\char"031F} \SUPPORTEDFALL{\char"0050\char"031F} \SUPPORTEDFALL{\char"01A4\char"031F} \SUPPORTEDANCH{\char"0052\char"031F} \SUPPORTEDFALL{\char"2C64\char"031F} \SUPPORTEDANCH{\char"0054\char"031F} \SUPPORTEDFALL{\char"01AC\char"031F} \SUPPORTEDFALL{\char"01AE\char"031F} \SUPPORTEDANCH{\char"0059\char"031F} \SUPPORTEDFALL{\char"01B3\char"031F} \SUPPORTEDFALL{\char"0061\char"031F} \SUPPORTEDFALL{\char"1E9A\char"031F} \SUPPORTEDFALL{\char"0062\char"031F} \SUPPORTEDFALL{\char"0183\char"031F} \SUPPORTEDFALL{\char"0063\char"031F} \SUPPORTEDFALL{\char"0188\char"031F} \SUPPORTEDFALL{\char"0064\char"031F} \SUPPORTEDFALL{\char"018C\char"031F} \SUPPORTEDFALL{\char"0221\char"031F} \SUPPORTEDANCH{\char"0068\char"031F} \SUPPORTEDFALL{\char"2C68\char"031F} \SUPPORTEDFALL{\char"006B\char"031F} \SUPPORTEDFALL{\char"0199\char"031F} \SUPPORTEDFALL{\char"2C6A\char"031F} \SUPPORTEDANCH{\char"006C\char"031F} \SUPPORTEDFALL{\char"0234\char"031F} \SUPPORTEDANCH{\char"006E\char"031F} \SUPPORTEDFALL{\char"0235\char"031F} \SUPPORTEDFALL{\char"0070\char"031F} \SUPPORTEDFALL{\char"01A5\char"031F} \SUPPORTEDFALL{\char"0071\char"031F} \SUPPORTEDFALL{\char"02A0\char"031F} \SUPPORTEDANCH{\char"0074\char"031F} \SUPPORTEDFALL{\char"01AB\char"031F} \SUPPORTEDFALL{\char"01AD\char"031F} \SUPPORTEDFALL{\char"0236\char"031F} \SUPPORTEDANCH{\char"0075\char"031F} \MISSINGGLYPH{\char"AB52\char"031F} \SUPPORTEDFALL{\char"0079\char"031F} \SUPPORTEDFALL{\char"01B4\char"031F} \MISSINGGLYPH{\char"1EFF\char"031F}

\SUPPORTEDANCH{\char"0048\char"0330} \SUPPORTEDFALL{\char"2C67\char"0330} \SUPPORTEDFALL{\char"004B\char"0330} \SUPPORTEDFALL{\char"2C69\char"0330} \SUPPORTEDFALL{\char"004D\char"0330} \MISSINGGLYPH{\char"2C6E\char"0330} \SUPPORTEDFALL{\char"0050\char"0330} \SUPPORTEDFALL{\char"01A4\char"0330} \SUPPORTEDANCH{\char"0052\char"0330} \SUPPORTEDFALL{\char"2C64\char"0330} \SUPPORTEDANCH{\char"0054\char"0330} \SUPPORTEDFALL{\char"01AC\char"0330} \SUPPORTEDFALL{\char"01AE\char"0330} \SUPPORTEDANCH{\char"0059\char"0330} \SUPPORTEDFALL{\char"01B3\char"0330} \SUPPORTEDFALL{\char"0061\char"0330} \SUPPORTEDFALL{\char"1E9A\char"0330} \SUPPORTEDFALL{\char"0062\char"0330} \SUPPORTEDFALL{\char"0183\char"0330} \SUPPORTEDFALL{\char"0063\char"0330} \SUPPORTEDFALL{\char"0188\char"0330} \SUPPORTEDFALL{\char"0064\char"0330} \SUPPORTEDFALL{\char"018C\char"0330} \SUPPORTEDFALL{\char"0221\char"0330} \SUPPORTEDANCH{\char"0068\char"0330} \SUPPORTEDFALL{\char"2C68\char"0330} \SUPPORTEDFALL{\char"006B\char"0330} \SUPPORTEDFALL{\char"0199\char"0330} \SUPPORTEDFALL{\char"2C6A\char"0330} \SUPPORTEDANCH{\char"006C\char"0330} \SUPPORTEDFALL{\char"0234\char"0330} \SUPPORTEDANCH{\char"006E\char"0330} \SUPPORTEDFALL{\char"0235\char"0330} \SUPPORTEDFALL{\char"0070\char"0330} \SUPPORTEDFALL{\char"01A5\char"0330} \SUPPORTEDFALL{\char"0071\char"0330} \SUPPORTEDFALL{\char"02A0\char"0330} \SUPPORTEDANCH{\char"0074\char"0330} \SUPPORTEDFALL{\char"01AB\char"0330} \SUPPORTEDFALL{\char"01AD\char"0330} \SUPPORTEDFALL{\char"0236\char"0330} \SUPPORTEDPREC{\char"0075\char"0330} \MISSINGGLYPH{\char"AB52\char"0330} \SUPPORTEDFALL{\char"0079\char"0330} \SUPPORTEDFALL{\char"01B4\char"0330} \MISSINGGLYPH{\char"1EFF\char"0330}

\SUPPORTEDANCH{\char"0048\char"0331} \SUPPORTEDFALL{\char"2C67\char"0331} \SUPPORTEDPREC{\char"004B\char"0331} \SUPPORTEDFALL{\char"2C69\char"0331} \SUPPORTEDFALL{\char"004D\char"0331} \MISSINGGLYPH{\char"2C6E\char"0331} \SUPPORTEDFALL{\char"0050\char"0331} \SUPPORTEDFALL{\char"01A4\char"0331} \SUPPORTEDPREC{\char"0052\char"0331} \SUPPORTEDFALL{\char"2C64\char"0331} \SUPPORTEDPREC{\char"0054\char"0331} \SUPPORTEDFALL{\char"01AC\char"0331} \SUPPORTEDFALL{\char"01AE\char"0331} \SUPPORTEDANCH{\char"0059\char"0331} \SUPPORTEDFALL{\char"01B3\char"0331} \SUPPORTEDFALL{\char"0061\char"0331} \SUPPORTEDFALL{\char"1E9A\char"0331} \SUPPORTEDPREC{\char"0062\char"0331} \SUPPORTEDFALL{\char"0183\char"0331} \SUPPORTEDFALL{\char"0063\char"0331} \SUPPORTEDFALL{\char"0188\char"0331} \SUPPORTEDPREC{\char"0064\char"0331} \SUPPORTEDFALL{\char"018C\char"0331} \SUPPORTEDFALL{\char"0221\char"0331} \SUPPORTEDPREC{\char"0068\char"0331} \SUPPORTEDFALL{\char"2C68\char"0331} \SUPPORTEDPREC{\char"006B\char"0331} \SUPPORTEDFALL{\char"0199\char"0331} \SUPPORTEDFALL{\char"2C6A\char"0331} \SUPPORTEDPREC{\char"006C\char"0331} \SUPPORTEDFALL{\char"0234\char"0331} \SUPPORTEDPREC{\char"006E\char"0331} \SUPPORTEDFALL{\char"0235\char"0331} \SUPPORTEDFALL{\char"0070\char"0331} \SUPPORTEDFALL{\char"01A5\char"0331} \SUPPORTEDFALL{\char"0071\char"0331} \SUPPORTEDFALL{\char"02A0\char"0331} \SUPPORTEDPREC{\char"0074\char"0331} \SUPPORTEDFALL{\char"01AB\char"0331} \SUPPORTEDFALL{\char"01AD\char"0331} \SUPPORTEDFALL{\char"0236\char"0331} \SUPPORTEDANCH{\char"0075\char"0331} \MISSINGGLYPH{\char"AB52\char"0331} \SUPPORTEDFALL{\char"0079\char"0331} \SUPPORTEDFALL{\char"01B4\char"0331} \MISSINGGLYPH{\char"1EFF\char"0331}

\MISSINGGLYPH{\char"0048} \MISSINGGLYPH{\char"2C67} \MISSINGGLYPH{\char"004B} \MISSINGGLYPH{\char"2C69} \MISSINGGLYPH{\char"004D} \MISSINGGLYPH{\char"2C6E} \MISSINGGLYPH{\char"0050} \MISSINGGLYPH{\char"01A4} \MISSINGGLYPH{\char"0052} \MISSINGGLYPH{\char"2C64} \MISSINGGLYPH{\char"0054} \MISSINGGLYPH{\char"01AC} \MISSINGGLYPH{\char"01AE} \MISSINGGLYPH{\char"0059} \MISSINGGLYPH{\char"01B3} \MISSINGGLYPH{\char"0061} \MISSINGGLYPH{\char"1E9A} \MISSINGGLYPH{\char"0062} \MISSINGGLYPH{\char"0183} \MISSINGGLYPH{\char"0063} \MISSINGGLYPH{\char"0188} \MISSINGGLYPH{\char"0064} \MISSINGGLYPH{\char"018C} \MISSINGGLYPH{\char"0221} \MISSINGGLYPH{\char"0068} \MISSINGGLYPH{\char"2C68} \MISSINGGLYPH{\char"006B} \MISSINGGLYPH{\char"0199} \MISSINGGLYPH{\char"2C6A} \MISSINGGLYPH{\char"006C} \MISSINGGLYPH{\char"0234} \MISSINGGLYPH{\char"006E} \MISSINGGLYPH{\char"0235} \MISSINGGLYPH{\char"0070} \MISSINGGLYPH{\char"01A5} \MISSINGGLYPH{\char"0071} \MISSINGGLYPH{\char"02A0} \MISSINGGLYPH{\char"0074} \MISSINGGLYPH{\char"01AB} \MISSINGGLYPH{\char"01AD} \MISSINGGLYPH{\char"0236} \MISSINGGLYPH{\char"0075} \MISSINGGLYPH{\char"AB52} \MISSINGGLYPH{\char"0079} \MISSINGGLYPH{\char"01B4} \MISSINGGLYPH{\char"1EFF}

\MISSINGGLYPH{\char"0048} \MISSINGGLYPH{\char"2C67} \MISSINGGLYPH{\char"004B} \MISSINGGLYPH{\char"2C69} \MISSINGGLYPH{\char"004D} \MISSINGGLYPH{\char"2C6E} \MISSINGGLYPH{\char"0050} \MISSINGGLYPH{\char"01A4} \MISSINGGLYPH{\char"0052} \MISSINGGLYPH{\char"2C64} \MISSINGGLYPH{\char"0054} \MISSINGGLYPH{\char"01AC} \MISSINGGLYPH{\char"01AE} \MISSINGGLYPH{\char"0059} \MISSINGGLYPH{\char"01B3} \MISSINGGLYPH{\char"0061} \MISSINGGLYPH{\char"1E9A} \MISSINGGLYPH{\char"0062} \MISSINGGLYPH{\char"0183} \MISSINGGLYPH{\char"0063} \MISSINGGLYPH{\char"0188} \MISSINGGLYPH{\char"0064} \MISSINGGLYPH{\char"018C} \MISSINGGLYPH{\char"0221} \MISSINGGLYPH{\char"0068} \MISSINGGLYPH{\char"2C68} \MISSINGGLYPH{\char"006B} \MISSINGGLYPH{\char"0199} \MISSINGGLYPH{\char"2C6A} \MISSINGGLYPH{\char"006C} \MISSINGGLYPH{\char"0234} \MISSINGGLYPH{\char"006E} \MISSINGGLYPH{\char"0235} \MISSINGGLYPH{\char"0070} \MISSINGGLYPH{\char"01A5} \MISSINGGLYPH{\char"0071} \MISSINGGLYPH{\char"02A0} \MISSINGGLYPH{\char"0074} \MISSINGGLYPH{\char"01AB} \MISSINGGLYPH{\char"01AD} \MISSINGGLYPH{\char"0236} \MISSINGGLYPH{\char"0075} \MISSINGGLYPH{\char"AB52} \MISSINGGLYPH{\char"0079} \MISSINGGLYPH{\char"01B4} \MISSINGGLYPH{\char"1EFF}

\MISSINGGLYPH{\char"0048} \MISSINGGLYPH{\char"2C67} \MISSINGGLYPH{\char"004B} \MISSINGGLYPH{\char"2C69} \MISSINGGLYPH{\char"004D} \MISSINGGLYPH{\char"2C6E} \MISSINGGLYPH{\char"0050} \MISSINGGLYPH{\char"01A4} \MISSINGGLYPH{\char"0052} \MISSINGGLYPH{\char"2C64} \MISSINGGLYPH{\char"0054} \MISSINGGLYPH{\char"01AC} \MISSINGGLYPH{\char"01AE} \MISSINGGLYPH{\char"0059} \MISSINGGLYPH{\char"01B3} \MISSINGGLYPH{\char"0061} \MISSINGGLYPH{\char"1E9A} \MISSINGGLYPH{\char"0062} \MISSINGGLYPH{\char"0183} \MISSINGGLYPH{\char"0063} \MISSINGGLYPH{\char"0188} \MISSINGGLYPH{\char"0064} \MISSINGGLYPH{\char"018C} \MISSINGGLYPH{\char"0221} \MISSINGGLYPH{\char"0068} \MISSINGGLYPH{\char"2C68} \MISSINGGLYPH{\char"006B} \MISSINGGLYPH{\char"0199} \MISSINGGLYPH{\char"2C6A} \MISSINGGLYPH{\char"006C} \MISSINGGLYPH{\char"0234} \MISSINGGLYPH{\char"006E} \MISSINGGLYPH{\char"0235} \MISSINGGLYPH{\char"0070} \MISSINGGLYPH{\char"01A5} \MISSINGGLYPH{\char"0071} \MISSINGGLYPH{\char"02A0} \MISSINGGLYPH{\char"0074} \MISSINGGLYPH{\char"01AB} \MISSINGGLYPH{\char"01AD} \MISSINGGLYPH{\char"0236} \MISSINGGLYPH{\char"0075} \MISSINGGLYPH{\char"AB52} \MISSINGGLYPH{\char"0079} \MISSINGGLYPH{\char"01B4} \MISSINGGLYPH{\char"1EFF}

\MISSINGGLYPH{\char"0048} \MISSINGGLYPH{\char"2C67} \MISSINGGLYPH{\char"004B} \MISSINGGLYPH{\char"2C69} \MISSINGGLYPH{\char"004D} \MISSINGGLYPH{\char"2C6E} \MISSINGGLYPH{\char"0050} \MISSINGGLYPH{\char"01A4} \MISSINGGLYPH{\char"0052} \MISSINGGLYPH{\char"2C64} \MISSINGGLYPH{\char"0054} \MISSINGGLYPH{\char"01AC} \MISSINGGLYPH{\char"01AE} \MISSINGGLYPH{\char"0059} \MISSINGGLYPH{\char"01B3} \MISSINGGLYPH{\char"0061} \MISSINGGLYPH{\char"1E9A} \MISSINGGLYPH{\char"0062} \MISSINGGLYPH{\char"0183} \MISSINGGLYPH{\char"0063} \MISSINGGLYPH{\char"0188} \MISSINGGLYPH{\char"0064} \MISSINGGLYPH{\char"018C} \MISSINGGLYPH{\char"0221} \MISSINGGLYPH{\char"0068} \MISSINGGLYPH{\char"2C68} \MISSINGGLYPH{\char"006B} \MISSINGGLYPH{\char"0199} \MISSINGGLYPH{\char"2C6A} \MISSINGGLYPH{\char"006C} \MISSINGGLYPH{\char"0234} \MISSINGGLYPH{\char"006E} \MISSINGGLYPH{\char"0235} \MISSINGGLYPH{\char"0070} \MISSINGGLYPH{\char"01A5} \MISSINGGLYPH{\char"0071} \MISSINGGLYPH{\char"02A0} \MISSINGGLYPH{\char"0074} \MISSINGGLYPH{\char"01AB} \MISSINGGLYPH{\char"01AD} \MISSINGGLYPH{\char"0236} \MISSINGGLYPH{\char"0075} \MISSINGGLYPH{\char"AB52} \MISSINGGLYPH{\char"0079} \MISSINGGLYPH{\char"01B4} \MISSINGGLYPH{\char"1EFF}

\MISSINGGLYPH{\char"0048} \MISSINGGLYPH{\char"2C67} \MISSINGGLYPH{\char"004B} \MISSINGGLYPH{\char"2C69} \MISSINGGLYPH{\char"004D} \MISSINGGLYPH{\char"2C6E} \MISSINGGLYPH{\char"0050} \MISSINGGLYPH{\char"01A4} \MISSINGGLYPH{\char"0052} \MISSINGGLYPH{\char"2C64} \MISSINGGLYPH{\char"0054} \MISSINGGLYPH{\char"01AC} \MISSINGGLYPH{\char"01AE} \MISSINGGLYPH{\char"0059} \MISSINGGLYPH{\char"01B3} \MISSINGGLYPH{\char"0061} \MISSINGGLYPH{\char"1E9A} \MISSINGGLYPH{\char"0062} \MISSINGGLYPH{\char"0183} \MISSINGGLYPH{\char"0063} \MISSINGGLYPH{\char"0188} \MISSINGGLYPH{\char"0064} \MISSINGGLYPH{\char"018C} \MISSINGGLYPH{\char"0221} \MISSINGGLYPH{\char"0068} \MISSINGGLYPH{\char"2C68} \MISSINGGLYPH{\char"006B} \MISSINGGLYPH{\char"0199} \MISSINGGLYPH{\char"2C6A} \MISSINGGLYPH{\char"006C} \MISSINGGLYPH{\char"0234} \MISSINGGLYPH{\char"006E} \MISSINGGLYPH{\char"0235} \MISSINGGLYPH{\char"0070} \MISSINGGLYPH{\char"01A5} \MISSINGGLYPH{\char"0071} \MISSINGGLYPH{\char"02A0} \MISSINGGLYPH{\char"0074} \MISSINGGLYPH{\char"01AB} \MISSINGGLYPH{\char"01AD} \MISSINGGLYPH{\char"0236} \MISSINGGLYPH{\char"0075} \MISSINGGLYPH{\char"AB52} \MISSINGGLYPH{\char"0079} \MISSINGGLYPH{\char"01B4} \MISSINGGLYPH{\char"1EFF}

\MISSINGGLYPH{\char"0048} \MISSINGGLYPH{\char"2C67} \MISSINGGLYPH{\char"004B} \MISSINGGLYPH{\char"2C69} \MISSINGGLYPH{\char"004D} \MISSINGGLYPH{\char"2C6E} \MISSINGGLYPH{\char"0050} \MISSINGGLYPH{\char"01A4} \MISSINGGLYPH{\char"0052} \MISSINGGLYPH{\char"2C64} \MISSINGGLYPH{\char"0054} \MISSINGGLYPH{\char"01AC} \MISSINGGLYPH{\char"01AE} \MISSINGGLYPH{\char"0059} \MISSINGGLYPH{\char"01B3} \MISSINGGLYPH{\char"0061} \MISSINGGLYPH{\char"1E9A} \MISSINGGLYPH{\char"0062} \MISSINGGLYPH{\char"0183} \MISSINGGLYPH{\char"0063} \MISSINGGLYPH{\char"0188} \MISSINGGLYPH{\char"0064} \MISSINGGLYPH{\char"018C} \MISSINGGLYPH{\char"0221} \MISSINGGLYPH{\char"0068} \MISSINGGLYPH{\char"2C68} \MISSINGGLYPH{\char"006B} \MISSINGGLYPH{\char"0199} \MISSINGGLYPH{\char"2C6A} \MISSINGGLYPH{\char"006C} \MISSINGGLYPH{\char"0234} \MISSINGGLYPH{\char"006E} \MISSINGGLYPH{\char"0235} \MISSINGGLYPH{\char"0070} \MISSINGGLYPH{\char"01A5} \MISSINGGLYPH{\char"0071} \MISSINGGLYPH{\char"02A0} \MISSINGGLYPH{\char"0074} \MISSINGGLYPH{\char"01AB} \MISSINGGLYPH{\char"01AD} \MISSINGGLYPH{\char"0236} \MISSINGGLYPH{\char"0075} \MISSINGGLYPH{\char"AB52} \MISSINGGLYPH{\char"0079} \MISSINGGLYPH{\char"01B4} \MISSINGGLYPH{\char"1EFF}

\MISSINGGLYPH{\char"0048} \MISSINGGLYPH{\char"2C67} \MISSINGGLYPH{\char"004B} \MISSINGGLYPH{\char"2C69} \MISSINGGLYPH{\char"004D} \MISSINGGLYPH{\char"2C6E} \MISSINGGLYPH{\char"0050} \MISSINGGLYPH{\char"01A4} \MISSINGGLYPH{\char"0052} \MISSINGGLYPH{\char"2C64} \MISSINGGLYPH{\char"0054} \MISSINGGLYPH{\char"01AC} \MISSINGGLYPH{\char"01AE} \MISSINGGLYPH{\char"0059} \MISSINGGLYPH{\char"01B3} \MISSINGGLYPH{\char"0061} \MISSINGGLYPH{\char"1E9A} \MISSINGGLYPH{\char"0062} \MISSINGGLYPH{\char"0183} \MISSINGGLYPH{\char"0063} \MISSINGGLYPH{\char"0188} \MISSINGGLYPH{\char"0064} \MISSINGGLYPH{\char"018C} \MISSINGGLYPH{\char"0221} \MISSINGGLYPH{\char"0068} \MISSINGGLYPH{\char"2C68} \MISSINGGLYPH{\char"006B} \MISSINGGLYPH{\char"0199} \MISSINGGLYPH{\char"2C6A} \MISSINGGLYPH{\char"006C} \MISSINGGLYPH{\char"0234} \MISSINGGLYPH{\char"006E} \MISSINGGLYPH{\char"0235} \MISSINGGLYPH{\char"0070} \MISSINGGLYPH{\char"01A5} \MISSINGGLYPH{\char"0071} \MISSINGGLYPH{\char"02A0} \MISSINGGLYPH{\char"0074} \MISSINGGLYPH{\char"01AB} \MISSINGGLYPH{\char"01AD} \MISSINGGLYPH{\char"0236} \MISSINGGLYPH{\char"0075} \MISSINGGLYPH{\char"AB52} \MISSINGGLYPH{\char"0079} \MISSINGGLYPH{\char"01B4} \MISSINGGLYPH{\char"1EFF}

\MISSINGGLYPH{\char"0048} \MISSINGGLYPH{\char"2C67} \MISSINGGLYPH{\char"004B} \MISSINGGLYPH{\char"2C69} \MISSINGGLYPH{\char"004D} \MISSINGGLYPH{\char"2C6E} \MISSINGGLYPH{\char"0050} \MISSINGGLYPH{\char"01A4} \MISSINGGLYPH{\char"0052} \MISSINGGLYPH{\char"2C64} \MISSINGGLYPH{\char"0054} \MISSINGGLYPH{\char"01AC} \MISSINGGLYPH{\char"01AE} \MISSINGGLYPH{\char"0059} \MISSINGGLYPH{\char"01B3} \MISSINGGLYPH{\char"0061} \MISSINGGLYPH{\char"1E9A} \MISSINGGLYPH{\char"0062} \MISSINGGLYPH{\char"0183} \MISSINGGLYPH{\char"0063} \MISSINGGLYPH{\char"0188} \MISSINGGLYPH{\char"0064} \MISSINGGLYPH{\char"018C} \MISSINGGLYPH{\char"0221} \MISSINGGLYPH{\char"0068} \MISSINGGLYPH{\char"2C68} \MISSINGGLYPH{\char"006B} \MISSINGGLYPH{\char"0199} \MISSINGGLYPH{\char"2C6A} \MISSINGGLYPH{\char"006C} \MISSINGGLYPH{\char"0234} \MISSINGGLYPH{\char"006E} \MISSINGGLYPH{\char"0235} \MISSINGGLYPH{\char"0070} \MISSINGGLYPH{\char"01A5} \MISSINGGLYPH{\char"0071} \MISSINGGLYPH{\char"02A0} \MISSINGGLYPH{\char"0074} \MISSINGGLYPH{\char"01AB} \MISSINGGLYPH{\char"01AD} \MISSINGGLYPH{\char"0236} \MISSINGGLYPH{\char"0075} \MISSINGGLYPH{\char"AB52} \MISSINGGLYPH{\char"0079} \MISSINGGLYPH{\char"01B4} \MISSINGGLYPH{\char"1EFF}

\MISSINGGLYPH{\char"0048} \MISSINGGLYPH{\char"2C67} \MISSINGGLYPH{\char"004B} \MISSINGGLYPH{\char"2C69} \MISSINGGLYPH{\char"004D} \MISSINGGLYPH{\char"2C6E} \MISSINGGLYPH{\char"0050} \MISSINGGLYPH{\char"01A4} \MISSINGGLYPH{\char"0052} \MISSINGGLYPH{\char"2C64} \MISSINGGLYPH{\char"0054} \MISSINGGLYPH{\char"01AC} \MISSINGGLYPH{\char"01AE} \MISSINGGLYPH{\char"0059} \MISSINGGLYPH{\char"01B3} \MISSINGGLYPH{\char"0061} \MISSINGGLYPH{\char"1E9A} \MISSINGGLYPH{\char"0062} \MISSINGGLYPH{\char"0183} \MISSINGGLYPH{\char"0063} \MISSINGGLYPH{\char"0188} \MISSINGGLYPH{\char"0064} \MISSINGGLYPH{\char"018C} \MISSINGGLYPH{\char"0221} \MISSINGGLYPH{\char"0068} \MISSINGGLYPH{\char"2C68} \MISSINGGLYPH{\char"006B} \MISSINGGLYPH{\char"0199} \MISSINGGLYPH{\char"2C6A} \MISSINGGLYPH{\char"006C} \MISSINGGLYPH{\char"0234} \MISSINGGLYPH{\char"006E} \MISSINGGLYPH{\char"0235} \MISSINGGLYPH{\char"0070} \MISSINGGLYPH{\char"01A5} \MISSINGGLYPH{\char"0071} \MISSINGGLYPH{\char"02A0} \MISSINGGLYPH{\char"0074} \MISSINGGLYPH{\char"01AB} \MISSINGGLYPH{\char"01AD} \MISSINGGLYPH{\char"0236} \MISSINGGLYPH{\char"0075} \MISSINGGLYPH{\char"AB52} \MISSINGGLYPH{\char"0079} \MISSINGGLYPH{\char"01B4} \MISSINGGLYPH{\char"1EFF}

\MISSINGGLYPH{\char"0048} \MISSINGGLYPH{\char"2C67} \MISSINGGLYPH{\char"004B} \MISSINGGLYPH{\char"2C69} \MISSINGGLYPH{\char"004D} \MISSINGGLYPH{\char"2C6E} \MISSINGGLYPH{\char"0050} \MISSINGGLYPH{\char"01A4} \MISSINGGLYPH{\char"0052} \MISSINGGLYPH{\char"2C64} \MISSINGGLYPH{\char"0054} \MISSINGGLYPH{\char"01AC} \MISSINGGLYPH{\char"01AE} \MISSINGGLYPH{\char"0059} \MISSINGGLYPH{\char"01B3} \MISSINGGLYPH{\char"0061} \MISSINGGLYPH{\char"1E9A} \MISSINGGLYPH{\char"0062} \MISSINGGLYPH{\char"0183} \MISSINGGLYPH{\char"0063} \MISSINGGLYPH{\char"0188} \MISSINGGLYPH{\char"0064} \MISSINGGLYPH{\char"018C} \MISSINGGLYPH{\char"0221} \MISSINGGLYPH{\char"0068} \MISSINGGLYPH{\char"2C68} \MISSINGGLYPH{\char"006B} \MISSINGGLYPH{\char"0199} \MISSINGGLYPH{\char"2C6A} \MISSINGGLYPH{\char"006C} \MISSINGGLYPH{\char"0234} \MISSINGGLYPH{\char"006E} \MISSINGGLYPH{\char"0235} \MISSINGGLYPH{\char"0070} \MISSINGGLYPH{\char"01A5} \MISSINGGLYPH{\char"0071} \MISSINGGLYPH{\char"02A0} \MISSINGGLYPH{\char"0074} \MISSINGGLYPH{\char"01AB} \MISSINGGLYPH{\char"01AD} \MISSINGGLYPH{\char"0236} \MISSINGGLYPH{\char"0075} \MISSINGGLYPH{\char"AB52} \MISSINGGLYPH{\char"0079} \MISSINGGLYPH{\char"01B4} \MISSINGGLYPH{\char"1EFF}

}

\newpage
\subsection*{Semibold patched, Group 4, below, Below}

{\LibertinusSerifPatch\bfseries
% grid. columns are bases, rows are marks.
\SUPPORTEDANCH{\char"0048\char"0316} \SUPPORTEDANCH{\char"2C67\char"0316} \SUPPORTEDANCH{\char"004B\char"0316} \SUPPORTEDANCH{\char"2C69\char"0316} \SUPPORTEDANCH{\char"004D\char"0316} \MISSINGGLYPH{\char"2C6E\char"0316} \SUPPORTEDANCH{\char"0050\char"0316} \SUPPORTEDANCH{\char"01A4\char"0316} \SUPPORTEDANCH{\char"0052\char"0316} \SUPPORTEDANCH{\char"2C64\char"0316} \SUPPORTEDANCH{\char"0054\char"0316} \SUPPORTEDANCH{\char"01AC\char"0316} \SUPPORTEDANCH{\char"01AE\char"0316} \SUPPORTEDANCH{\char"0059\char"0316} \SUPPORTEDANCH{\char"01B3\char"0316} \SUPPORTEDANCH{\char"0061\char"0316} \SUPPORTEDANCH{\char"1E9A\char"0316} \SUPPORTEDANCH{\char"0062\char"0316} \SUPPORTEDANCH{\char"0183\char"0316} \SUPPORTEDANCH{\char"0063\char"0316} \SUPPORTEDANCH{\char"0188\char"0316} \SUPPORTEDANCH{\char"0064\char"0316} \SUPPORTEDANCH{\char"018C\char"0316} \SUPPORTEDANCH{\char"0221\char"0316} \SUPPORTEDANCH{\char"0068\char"0316} \SUPPORTEDANCH{\char"2C68\char"0316} \SUPPORTEDANCH{\char"006B\char"0316} \SUPPORTEDANCH{\char"0199\char"0316} \SUPPORTEDANCH{\char"2C6A\char"0316} \SUPPORTEDANCH{\char"006C\char"0316} \SUPPORTEDANCH{\char"0234\char"0316} \SUPPORTEDANCH{\char"006E\char"0316} \SUPPORTEDANCH{\char"0235\char"0316} \SUPPORTEDANCH{\char"0070\char"0316} \SUPPORTEDANCH{\char"01A5\char"0316} \SUPPORTEDANCH{\char"0071\char"0316} \SUPPORTEDANCH{\char"02A0\char"0316} \SUPPORTEDANCH{\char"0074\char"0316} \SUPPORTEDANCH{\char"01AB\char"0316} \SUPPORTEDANCH{\char"01AD\char"0316} \SUPPORTEDANCH{\char"0236\char"0316} \SUPPORTEDANCH{\char"0075\char"0316} \MISSINGGLYPH{\char"AB52\char"0316} \SUPPORTEDANCH{\char"0079\char"0316} \SUPPORTEDANCH{\char"01B4\char"0316} \MISSINGGLYPH{\char"1EFF\char"0316}

\SUPPORTEDANCH{\char"0048\char"0317} \SUPPORTEDANCH{\char"2C67\char"0317} \SUPPORTEDANCH{\char"004B\char"0317} \SUPPORTEDANCH{\char"2C69\char"0317} \SUPPORTEDANCH{\char"004D\char"0317} \MISSINGGLYPH{\char"2C6E\char"0317} \SUPPORTEDANCH{\char"0050\char"0317} \SUPPORTEDANCH{\char"01A4\char"0317} \SUPPORTEDANCH{\char"0052\char"0317} \SUPPORTEDANCH{\char"2C64\char"0317} \SUPPORTEDANCH{\char"0054\char"0317} \SUPPORTEDANCH{\char"01AC\char"0317} \SUPPORTEDANCH{\char"01AE\char"0317} \SUPPORTEDANCH{\char"0059\char"0317} \SUPPORTEDANCH{\char"01B3\char"0317} \SUPPORTEDANCH{\char"0061\char"0317} \SUPPORTEDANCH{\char"1E9A\char"0317} \SUPPORTEDANCH{\char"0062\char"0317} \SUPPORTEDANCH{\char"0183\char"0317} \SUPPORTEDANCH{\char"0063\char"0317} \SUPPORTEDANCH{\char"0188\char"0317} \SUPPORTEDANCH{\char"0064\char"0317} \SUPPORTEDANCH{\char"018C\char"0317} \SUPPORTEDANCH{\char"0221\char"0317} \SUPPORTEDANCH{\char"0068\char"0317} \SUPPORTEDANCH{\char"2C68\char"0317} \SUPPORTEDANCH{\char"006B\char"0317} \SUPPORTEDANCH{\char"0199\char"0317} \SUPPORTEDANCH{\char"2C6A\char"0317} \SUPPORTEDANCH{\char"006C\char"0317} \SUPPORTEDANCH{\char"0234\char"0317} \SUPPORTEDANCH{\char"006E\char"0317} \SUPPORTEDANCH{\char"0235\char"0317} \SUPPORTEDANCH{\char"0070\char"0317} \SUPPORTEDANCH{\char"01A5\char"0317} \SUPPORTEDANCH{\char"0071\char"0317} \SUPPORTEDANCH{\char"02A0\char"0317} \SUPPORTEDANCH{\char"0074\char"0317} \SUPPORTEDANCH{\char"01AB\char"0317} \SUPPORTEDANCH{\char"01AD\char"0317} \SUPPORTEDANCH{\char"0236\char"0317} \SUPPORTEDANCH{\char"0075\char"0317} \MISSINGGLYPH{\char"AB52\char"0317} \SUPPORTEDANCH{\char"0079\char"0317} \SUPPORTEDANCH{\char"01B4\char"0317} \MISSINGGLYPH{\char"1EFF\char"0317}

\SUPPORTEDANCH{\char"0048\char"0318} \SUPPORTEDANCH{\char"2C67\char"0318} \SUPPORTEDANCH{\char"004B\char"0318} \SUPPORTEDANCH{\char"2C69\char"0318} \SUPPORTEDANCH{\char"004D\char"0318} \MISSINGGLYPH{\char"2C6E\char"0318} \SUPPORTEDANCH{\char"0050\char"0318} \SUPPORTEDANCH{\char"01A4\char"0318} \SUPPORTEDANCH{\char"0052\char"0318} \SUPPORTEDANCH{\char"2C64\char"0318} \SUPPORTEDANCH{\char"0054\char"0318} \SUPPORTEDANCH{\char"01AC\char"0318} \SUPPORTEDANCH{\char"01AE\char"0318} \SUPPORTEDANCH{\char"0059\char"0318} \SUPPORTEDANCH{\char"01B3\char"0318} \SUPPORTEDANCH{\char"0061\char"0318} \SUPPORTEDANCH{\char"1E9A\char"0318} \SUPPORTEDANCH{\char"0062\char"0318} \SUPPORTEDANCH{\char"0183\char"0318} \SUPPORTEDANCH{\char"0063\char"0318} \SUPPORTEDANCH{\char"0188\char"0318} \SUPPORTEDANCH{\char"0064\char"0318} \SUPPORTEDANCH{\char"018C\char"0318} \SUPPORTEDANCH{\char"0221\char"0318} \SUPPORTEDANCH{\char"0068\char"0318} \SUPPORTEDANCH{\char"2C68\char"0318} \SUPPORTEDANCH{\char"006B\char"0318} \SUPPORTEDANCH{\char"0199\char"0318} \SUPPORTEDANCH{\char"2C6A\char"0318} \SUPPORTEDANCH{\char"006C\char"0318} \SUPPORTEDANCH{\char"0234\char"0318} \SUPPORTEDANCH{\char"006E\char"0318} \SUPPORTEDANCH{\char"0235\char"0318} \SUPPORTEDANCH{\char"0070\char"0318} \SUPPORTEDANCH{\char"01A5\char"0318} \SUPPORTEDANCH{\char"0071\char"0318} \SUPPORTEDANCH{\char"02A0\char"0318} \SUPPORTEDANCH{\char"0074\char"0318} \SUPPORTEDANCH{\char"01AB\char"0318} \SUPPORTEDANCH{\char"01AD\char"0318} \SUPPORTEDANCH{\char"0236\char"0318} \SUPPORTEDANCH{\char"0075\char"0318} \MISSINGGLYPH{\char"AB52\char"0318} \SUPPORTEDANCH{\char"0079\char"0318} \SUPPORTEDANCH{\char"01B4\char"0318} \MISSINGGLYPH{\char"1EFF\char"0318}

\SUPPORTEDANCH{\char"0048\char"0319} \SUPPORTEDANCH{\char"2C67\char"0319} \SUPPORTEDANCH{\char"004B\char"0319} \SUPPORTEDANCH{\char"2C69\char"0319} \SUPPORTEDANCH{\char"004D\char"0319} \MISSINGGLYPH{\char"2C6E\char"0319} \SUPPORTEDANCH{\char"0050\char"0319} \SUPPORTEDANCH{\char"01A4\char"0319} \SUPPORTEDANCH{\char"0052\char"0319} \SUPPORTEDANCH{\char"2C64\char"0319} \SUPPORTEDANCH{\char"0054\char"0319} \SUPPORTEDANCH{\char"01AC\char"0319} \SUPPORTEDANCH{\char"01AE\char"0319} \SUPPORTEDANCH{\char"0059\char"0319} \SUPPORTEDANCH{\char"01B3\char"0319} \SUPPORTEDANCH{\char"0061\char"0319} \SUPPORTEDANCH{\char"1E9A\char"0319} \SUPPORTEDANCH{\char"0062\char"0319} \SUPPORTEDANCH{\char"0183\char"0319} \SUPPORTEDANCH{\char"0063\char"0319} \SUPPORTEDANCH{\char"0188\char"0319} \SUPPORTEDANCH{\char"0064\char"0319} \SUPPORTEDANCH{\char"018C\char"0319} \SUPPORTEDANCH{\char"0221\char"0319} \SUPPORTEDANCH{\char"0068\char"0319} \SUPPORTEDANCH{\char"2C68\char"0319} \SUPPORTEDANCH{\char"006B\char"0319} \SUPPORTEDANCH{\char"0199\char"0319} \SUPPORTEDANCH{\char"2C6A\char"0319} \SUPPORTEDANCH{\char"006C\char"0319} \SUPPORTEDANCH{\char"0234\char"0319} \SUPPORTEDANCH{\char"006E\char"0319} \SUPPORTEDANCH{\char"0235\char"0319} \SUPPORTEDANCH{\char"0070\char"0319} \SUPPORTEDANCH{\char"01A5\char"0319} \SUPPORTEDANCH{\char"0071\char"0319} \SUPPORTEDANCH{\char"02A0\char"0319} \SUPPORTEDANCH{\char"0074\char"0319} \SUPPORTEDANCH{\char"01AB\char"0319} \SUPPORTEDANCH{\char"01AD\char"0319} \SUPPORTEDANCH{\char"0236\char"0319} \SUPPORTEDANCH{\char"0075\char"0319} \MISSINGGLYPH{\char"AB52\char"0319} \SUPPORTEDANCH{\char"0079\char"0319} \SUPPORTEDANCH{\char"01B4\char"0319} \MISSINGGLYPH{\char"1EFF\char"0319}

\SUPPORTEDANCH{\char"0048\char"031C} \SUPPORTEDANCH{\char"2C67\char"031C} \SUPPORTEDANCH{\char"004B\char"031C} \SUPPORTEDANCH{\char"2C69\char"031C} \SUPPORTEDANCH{\char"004D\char"031C} \MISSINGGLYPH{\char"2C6E\char"031C} \SUPPORTEDANCH{\char"0050\char"031C} \SUPPORTEDANCH{\char"01A4\char"031C} \SUPPORTEDANCH{\char"0052\char"031C} \SUPPORTEDANCH{\char"2C64\char"031C} \SUPPORTEDANCH{\char"0054\char"031C} \SUPPORTEDANCH{\char"01AC\char"031C} \SUPPORTEDANCH{\char"01AE\char"031C} \SUPPORTEDANCH{\char"0059\char"031C} \SUPPORTEDANCH{\char"01B3\char"031C} \SUPPORTEDANCH{\char"0061\char"031C} \SUPPORTEDANCH{\char"1E9A\char"031C} \SUPPORTEDANCH{\char"0062\char"031C} \SUPPORTEDANCH{\char"0183\char"031C} \SUPPORTEDANCH{\char"0063\char"031C} \SUPPORTEDANCH{\char"0188\char"031C} \SUPPORTEDANCH{\char"0064\char"031C} \SUPPORTEDANCH{\char"018C\char"031C} \SUPPORTEDANCH{\char"0221\char"031C} \SUPPORTEDANCH{\char"0068\char"031C} \SUPPORTEDANCH{\char"2C68\char"031C} \SUPPORTEDANCH{\char"006B\char"031C} \SUPPORTEDANCH{\char"0199\char"031C} \SUPPORTEDANCH{\char"2C6A\char"031C} \SUPPORTEDANCH{\char"006C\char"031C} \SUPPORTEDANCH{\char"0234\char"031C} \SUPPORTEDANCH{\char"006E\char"031C} \SUPPORTEDANCH{\char"0235\char"031C} \SUPPORTEDANCH{\char"0070\char"031C} \SUPPORTEDANCH{\char"01A5\char"031C} \SUPPORTEDANCH{\char"0071\char"031C} \SUPPORTEDANCH{\char"02A0\char"031C} \SUPPORTEDANCH{\char"0074\char"031C} \SUPPORTEDANCH{\char"01AB\char"031C} \SUPPORTEDANCH{\char"01AD\char"031C} \SUPPORTEDANCH{\char"0236\char"031C} \SUPPORTEDANCH{\char"0075\char"031C} \MISSINGGLYPH{\char"AB52\char"031C} \SUPPORTEDANCH{\char"0079\char"031C} \SUPPORTEDANCH{\char"01B4\char"031C} \MISSINGGLYPH{\char"1EFF\char"031C}

\SUPPORTEDANCH{\char"0048\char"0320} \SUPPORTEDANCH{\char"2C67\char"0320} \SUPPORTEDANCH{\char"004B\char"0320} \SUPPORTEDANCH{\char"2C69\char"0320} \SUPPORTEDANCH{\char"004D\char"0320} \MISSINGGLYPH{\char"2C6E\char"0320} \SUPPORTEDANCH{\char"0050\char"0320} \SUPPORTEDANCH{\char"01A4\char"0320} \SUPPORTEDANCH{\char"0052\char"0320} \SUPPORTEDANCH{\char"2C64\char"0320} \SUPPORTEDANCH{\char"0054\char"0320} \SUPPORTEDANCH{\char"01AC\char"0320} \SUPPORTEDANCH{\char"01AE\char"0320} \SUPPORTEDANCH{\char"0059\char"0320} \SUPPORTEDANCH{\char"01B3\char"0320} \SUPPORTEDANCH{\char"0061\char"0320} \SUPPORTEDANCH{\char"1E9A\char"0320} \SUPPORTEDANCH{\char"0062\char"0320} \SUPPORTEDANCH{\char"0183\char"0320} \SUPPORTEDANCH{\char"0063\char"0320} \SUPPORTEDANCH{\char"0188\char"0320} \SUPPORTEDANCH{\char"0064\char"0320} \SUPPORTEDANCH{\char"018C\char"0320} \SUPPORTEDANCH{\char"0221\char"0320} \SUPPORTEDANCH{\char"0068\char"0320} \SUPPORTEDANCH{\char"2C68\char"0320} \SUPPORTEDANCH{\char"006B\char"0320} \SUPPORTEDANCH{\char"0199\char"0320} \SUPPORTEDANCH{\char"2C6A\char"0320} \SUPPORTEDANCH{\char"006C\char"0320} \SUPPORTEDANCH{\char"0234\char"0320} \SUPPORTEDANCH{\char"006E\char"0320} \SUPPORTEDANCH{\char"0235\char"0320} \SUPPORTEDANCH{\char"0070\char"0320} \SUPPORTEDANCH{\char"01A5\char"0320} \SUPPORTEDANCH{\char"0071\char"0320} \SUPPORTEDANCH{\char"02A0\char"0320} \SUPPORTEDANCH{\char"0074\char"0320} \SUPPORTEDANCH{\char"01AB\char"0320} \SUPPORTEDANCH{\char"01AD\char"0320} \SUPPORTEDANCH{\char"0236\char"0320} \SUPPORTEDANCH{\char"0075\char"0320} \MISSINGGLYPH{\char"AB52\char"0320} \SUPPORTEDANCH{\char"0079\char"0320} \SUPPORTEDANCH{\char"01B4\char"0320} \MISSINGGLYPH{\char"1EFF\char"0320}

\SUPPORTEDPREC{\char"0048\char"0323} \SUPPORTEDANCH{\char"2C67\char"0323} \SUPPORTEDPREC{\char"004B\char"0323} \SUPPORTEDANCH{\char"2C69\char"0323} \SUPPORTEDPREC{\char"004D\char"0323} \MISSINGGLYPH{\char"2C6E\char"0323} \SUPPORTEDANCH{\char"0050\char"0323} \SUPPORTEDANCH{\char"01A4\char"0323} \SUPPORTEDPREC{\char"0052\char"0323} \SUPPORTEDANCH{\char"2C64\char"0323} \SUPPORTEDPREC{\char"0054\char"0323} \SUPPORTEDANCH{\char"01AC\char"0323} \SUPPORTEDANCH{\char"01AE\char"0323} \SUPPORTEDPREC{\char"0059\char"0323} \SUPPORTEDANCH{\char"01B3\char"0323} \SUPPORTEDPREC{\char"0061\char"0323} \SUPPORTEDANCH{\char"1E9A\char"0323} \SUPPORTEDPREC{\char"0062\char"0323} \SUPPORTEDANCH{\char"0183\char"0323} \SUPPORTEDANCH{\char"0063\char"0323} \SUPPORTEDANCH{\char"0188\char"0323} \SUPPORTEDPREC{\char"0064\char"0323} \SUPPORTEDANCH{\char"018C\char"0323} \SUPPORTEDANCH{\char"0221\char"0323} \SUPPORTEDPREC{\char"0068\char"0323} \SUPPORTEDANCH{\char"2C68\char"0323} \SUPPORTEDPREC{\char"006B\char"0323} \SUPPORTEDANCH{\char"0199\char"0323} \SUPPORTEDANCH{\char"2C6A\char"0323} \SUPPORTEDPREC{\char"006C\char"0323} \SUPPORTEDANCH{\char"0234\char"0323} \SUPPORTEDPREC{\char"006E\char"0323} \SUPPORTEDANCH{\char"0235\char"0323} \SUPPORTEDANCH{\char"0070\char"0323} \SUPPORTEDANCH{\char"01A5\char"0323} \SUPPORTEDANCH{\char"0071\char"0323} \SUPPORTEDANCH{\char"02A0\char"0323} \SUPPORTEDPREC{\char"0074\char"0323} \SUPPORTEDANCH{\char"01AB\char"0323} \SUPPORTEDANCH{\char"01AD\char"0323} \SUPPORTEDANCH{\char"0236\char"0323} \SUPPORTEDPREC{\char"0075\char"0323} \MISSINGGLYPH{\char"AB52\char"0323} \SUPPORTEDPREC{\char"0079\char"0323} \SUPPORTEDANCH{\char"01B4\char"0323} \MISSINGGLYPH{\char"1EFF\char"0323}

\SUPPORTEDANCH{\char"0048\char"0324} \SUPPORTEDANCH{\char"2C67\char"0324} \SUPPORTEDANCH{\char"004B\char"0324} \SUPPORTEDANCH{\char"2C69\char"0324} \SUPPORTEDANCH{\char"004D\char"0324} \MISSINGGLYPH{\char"2C6E\char"0324} \SUPPORTEDANCH{\char"0050\char"0324} \SUPPORTEDANCH{\char"01A4\char"0324} \SUPPORTEDANCH{\char"0052\char"0324} \SUPPORTEDANCH{\char"2C64\char"0324} \SUPPORTEDANCH{\char"0054\char"0324} \SUPPORTEDANCH{\char"01AC\char"0324} \SUPPORTEDANCH{\char"01AE\char"0324} \SUPPORTEDANCH{\char"0059\char"0324} \SUPPORTEDANCH{\char"01B3\char"0324} \SUPPORTEDANCH{\char"0061\char"0324} \SUPPORTEDANCH{\char"1E9A\char"0324} \SUPPORTEDANCH{\char"0062\char"0324} \SUPPORTEDANCH{\char"0183\char"0324} \SUPPORTEDANCH{\char"0063\char"0324} \SUPPORTEDANCH{\char"0188\char"0324} \SUPPORTEDANCH{\char"0064\char"0324} \SUPPORTEDANCH{\char"018C\char"0324} \SUPPORTEDANCH{\char"0221\char"0324} \SUPPORTEDANCH{\char"0068\char"0324} \SUPPORTEDANCH{\char"2C68\char"0324} \SUPPORTEDANCH{\char"006B\char"0324} \SUPPORTEDANCH{\char"0199\char"0324} \SUPPORTEDANCH{\char"2C6A\char"0324} \SUPPORTEDANCH{\char"006C\char"0324} \SUPPORTEDANCH{\char"0234\char"0324} \SUPPORTEDANCH{\char"006E\char"0324} \SUPPORTEDANCH{\char"0235\char"0324} \SUPPORTEDANCH{\char"0070\char"0324} \SUPPORTEDANCH{\char"01A5\char"0324} \SUPPORTEDANCH{\char"0071\char"0324} \SUPPORTEDANCH{\char"02A0\char"0324} \SUPPORTEDANCH{\char"0074\char"0324} \SUPPORTEDANCH{\char"01AB\char"0324} \SUPPORTEDANCH{\char"01AD\char"0324} \SUPPORTEDANCH{\char"0236\char"0324} \SUPPORTEDPREC{\char"0075\char"0324} \MISSINGGLYPH{\char"AB52\char"0324} \SUPPORTEDANCH{\char"0079\char"0324} \SUPPORTEDANCH{\char"01B4\char"0324} \MISSINGGLYPH{\char"1EFF\char"0324}

\SUPPORTEDANCH{\char"0048\char"0325} \SUPPORTEDANCH{\char"2C67\char"0325} \SUPPORTEDANCH{\char"004B\char"0325} \SUPPORTEDANCH{\char"2C69\char"0325} \SUPPORTEDANCH{\char"004D\char"0325} \MISSINGGLYPH{\char"2C6E\char"0325} \SUPPORTEDANCH{\char"0050\char"0325} \SUPPORTEDANCH{\char"01A4\char"0325} \SUPPORTEDANCH{\char"0052\char"0325} \SUPPORTEDANCH{\char"2C64\char"0325} \SUPPORTEDANCH{\char"0054\char"0325} \SUPPORTEDANCH{\char"01AC\char"0325} \SUPPORTEDANCH{\char"01AE\char"0325} \SUPPORTEDANCH{\char"0059\char"0325} \SUPPORTEDANCH{\char"01B3\char"0325} \SUPPORTEDPREC{\char"0061\char"0325} \SUPPORTEDANCH{\char"1E9A\char"0325} \SUPPORTEDANCH{\char"0062\char"0325} \SUPPORTEDANCH{\char"0183\char"0325} \SUPPORTEDANCH{\char"0063\char"0325} \SUPPORTEDANCH{\char"0188\char"0325} \SUPPORTEDANCH{\char"0064\char"0325} \SUPPORTEDANCH{\char"018C\char"0325} \SUPPORTEDANCH{\char"0221\char"0325} \SUPPORTEDANCH{\char"0068\char"0325} \SUPPORTEDANCH{\char"2C68\char"0325} \SUPPORTEDANCH{\char"006B\char"0325} \SUPPORTEDANCH{\char"0199\char"0325} \SUPPORTEDANCH{\char"2C6A\char"0325} \SUPPORTEDANCH{\char"006C\char"0325} \SUPPORTEDANCH{\char"0234\char"0325} \SUPPORTEDANCH{\char"006E\char"0325} \SUPPORTEDANCH{\char"0235\char"0325} \SUPPORTEDANCH{\char"0070\char"0325} \SUPPORTEDANCH{\char"01A5\char"0325} \SUPPORTEDANCH{\char"0071\char"0325} \SUPPORTEDANCH{\char"02A0\char"0325} \SUPPORTEDANCH{\char"0074\char"0325} \SUPPORTEDANCH{\char"01AB\char"0325} \SUPPORTEDANCH{\char"01AD\char"0325} \SUPPORTEDANCH{\char"0236\char"0325} \SUPPORTEDANCH{\char"0075\char"0325} \MISSINGGLYPH{\char"AB52\char"0325} \SUPPORTEDANCH{\char"0079\char"0325} \SUPPORTEDANCH{\char"01B4\char"0325} \MISSINGGLYPH{\char"1EFF\char"0325}

\SUPPORTEDANCH{\char"0048\char"0326} \SUPPORTEDANCH{\char"2C67\char"0326} \SUPPORTEDANCH{\char"004B\char"0326} \SUPPORTEDANCH{\char"2C69\char"0326} \SUPPORTEDANCH{\char"004D\char"0326} \MISSINGGLYPH{\char"2C6E\char"0326} \SUPPORTEDANCH{\char"0050\char"0326} \SUPPORTEDANCH{\char"01A4\char"0326} \SUPPORTEDANCH{\char"0052\char"0326} \SUPPORTEDANCH{\char"2C64\char"0326} \SUPPORTEDPREC{\char"0054\char"0326} \SUPPORTEDANCH{\char"01AC\char"0326} \SUPPORTEDANCH{\char"01AE\char"0326} \SUPPORTEDANCH{\char"0059\char"0326} \SUPPORTEDANCH{\char"01B3\char"0326} \SUPPORTEDANCH{\char"0061\char"0326} \SUPPORTEDANCH{\char"1E9A\char"0326} \SUPPORTEDANCH{\char"0062\char"0326} \SUPPORTEDANCH{\char"0183\char"0326} \SUPPORTEDANCH{\char"0063\char"0326} \SUPPORTEDANCH{\char"0188\char"0326} \SUPPORTEDANCH{\char"0064\char"0326} \SUPPORTEDANCH{\char"018C\char"0326} \SUPPORTEDANCH{\char"0221\char"0326} \SUPPORTEDANCH{\char"0068\char"0326} \SUPPORTEDANCH{\char"2C68\char"0326} \SUPPORTEDANCH{\char"006B\char"0326} \SUPPORTEDANCH{\char"0199\char"0326} \SUPPORTEDANCH{\char"2C6A\char"0326} \SUPPORTEDANCH{\char"006C\char"0326} \SUPPORTEDANCH{\char"0234\char"0326} \SUPPORTEDANCH{\char"006E\char"0326} \SUPPORTEDANCH{\char"0235\char"0326} \SUPPORTEDANCH{\char"0070\char"0326} \SUPPORTEDANCH{\char"01A5\char"0326} \SUPPORTEDANCH{\char"0071\char"0326} \SUPPORTEDANCH{\char"02A0\char"0326} \SUPPORTEDPREC{\char"0074\char"0326} \SUPPORTEDANCH{\char"01AB\char"0326} \SUPPORTEDANCH{\char"01AD\char"0326} \SUPPORTEDANCH{\char"0236\char"0326} \SUPPORTEDANCH{\char"0075\char"0326} \MISSINGGLYPH{\char"AB52\char"0326} \SUPPORTEDANCH{\char"0079\char"0326} \SUPPORTEDANCH{\char"01B4\char"0326} \MISSINGGLYPH{\char"1EFF\char"0326}

\SUPPORTEDANCH{\char"0048\char"0329} \SUPPORTEDANCH{\char"2C67\char"0329} \SUPPORTEDANCH{\char"004B\char"0329} \SUPPORTEDANCH{\char"2C69\char"0329} \SUPPORTEDANCH{\char"004D\char"0329} \MISSINGGLYPH{\char"2C6E\char"0329} \SUPPORTEDANCH{\char"0050\char"0329} \SUPPORTEDANCH{\char"01A4\char"0329} \SUPPORTEDANCH{\char"0052\char"0329} \SUPPORTEDANCH{\char"2C64\char"0329} \SUPPORTEDANCH{\char"0054\char"0329} \SUPPORTEDANCH{\char"01AC\char"0329} \SUPPORTEDANCH{\char"01AE\char"0329} \SUPPORTEDANCH{\char"0059\char"0329} \SUPPORTEDANCH{\char"01B3\char"0329} \SUPPORTEDANCH{\char"0061\char"0329} \SUPPORTEDANCH{\char"1E9A\char"0329} \SUPPORTEDANCH{\char"0062\char"0329} \SUPPORTEDANCH{\char"0183\char"0329} \SUPPORTEDANCH{\char"0063\char"0329} \SUPPORTEDANCH{\char"0188\char"0329} \SUPPORTEDANCH{\char"0064\char"0329} \SUPPORTEDANCH{\char"018C\char"0329} \SUPPORTEDANCH{\char"0221\char"0329} \SUPPORTEDANCH{\char"0068\char"0329} \SUPPORTEDANCH{\char"2C68\char"0329} \SUPPORTEDANCH{\char"006B\char"0329} \SUPPORTEDANCH{\char"0199\char"0329} \SUPPORTEDANCH{\char"2C6A\char"0329} \SUPPORTEDANCH{\char"006C\char"0329} \SUPPORTEDANCH{\char"0234\char"0329} \SUPPORTEDANCH{\char"006E\char"0329} \SUPPORTEDANCH{\char"0235\char"0329} \SUPPORTEDANCH{\char"0070\char"0329} \SUPPORTEDANCH{\char"01A5\char"0329} \SUPPORTEDANCH{\char"0071\char"0329} \SUPPORTEDANCH{\char"02A0\char"0329} \SUPPORTEDANCH{\char"0074\char"0329} \SUPPORTEDANCH{\char"01AB\char"0329} \SUPPORTEDANCH{\char"01AD\char"0329} \SUPPORTEDANCH{\char"0236\char"0329} \SUPPORTEDANCH{\char"0075\char"0329} \MISSINGGLYPH{\char"AB52\char"0329} \SUPPORTEDANCH{\char"0079\char"0329} \SUPPORTEDANCH{\char"01B4\char"0329} \MISSINGGLYPH{\char"1EFF\char"0329}

\SUPPORTEDANCH{\char"0048\char"032A} \SUPPORTEDANCH{\char"2C67\char"032A} \SUPPORTEDANCH{\char"004B\char"032A} \SUPPORTEDANCH{\char"2C69\char"032A} \SUPPORTEDANCH{\char"004D\char"032A} \MISSINGGLYPH{\char"2C6E\char"032A} \SUPPORTEDANCH{\char"0050\char"032A} \SUPPORTEDANCH{\char"01A4\char"032A} \SUPPORTEDANCH{\char"0052\char"032A} \SUPPORTEDANCH{\char"2C64\char"032A} \SUPPORTEDANCH{\char"0054\char"032A} \SUPPORTEDANCH{\char"01AC\char"032A} \SUPPORTEDANCH{\char"01AE\char"032A} \SUPPORTEDANCH{\char"0059\char"032A} \SUPPORTEDANCH{\char"01B3\char"032A} \SUPPORTEDANCH{\char"0061\char"032A} \SUPPORTEDANCH{\char"1E9A\char"032A} \SUPPORTEDANCH{\char"0062\char"032A} \SUPPORTEDANCH{\char"0183\char"032A} \SUPPORTEDANCH{\char"0063\char"032A} \SUPPORTEDANCH{\char"0188\char"032A} \SUPPORTEDANCH{\char"0064\char"032A} \SUPPORTEDANCH{\char"018C\char"032A} \SUPPORTEDANCH{\char"0221\char"032A} \SUPPORTEDANCH{\char"0068\char"032A} \SUPPORTEDANCH{\char"2C68\char"032A} \SUPPORTEDANCH{\char"006B\char"032A} \SUPPORTEDANCH{\char"0199\char"032A} \SUPPORTEDANCH{\char"2C6A\char"032A} \SUPPORTEDANCH{\char"006C\char"032A} \SUPPORTEDANCH{\char"0234\char"032A} \SUPPORTEDANCH{\char"006E\char"032A} \SUPPORTEDANCH{\char"0235\char"032A} \SUPPORTEDANCH{\char"0070\char"032A} \SUPPORTEDANCH{\char"01A5\char"032A} \SUPPORTEDANCH{\char"0071\char"032A} \SUPPORTEDANCH{\char"02A0\char"032A} \SUPPORTEDANCH{\char"0074\char"032A} \SUPPORTEDANCH{\char"01AB\char"032A} \SUPPORTEDANCH{\char"01AD\char"032A} \SUPPORTEDANCH{\char"0236\char"032A} \SUPPORTEDANCH{\char"0075\char"032A} \MISSINGGLYPH{\char"AB52\char"032A} \SUPPORTEDANCH{\char"0079\char"032A} \SUPPORTEDANCH{\char"01B4\char"032A} \MISSINGGLYPH{\char"1EFF\char"032A}

\SUPPORTEDANCH{\char"0048\char"032B} \SUPPORTEDANCH{\char"2C67\char"032B} \SUPPORTEDANCH{\char"004B\char"032B} \SUPPORTEDANCH{\char"2C69\char"032B} \SUPPORTEDANCH{\char"004D\char"032B} \MISSINGGLYPH{\char"2C6E\char"032B} \SUPPORTEDANCH{\char"0050\char"032B} \SUPPORTEDANCH{\char"01A4\char"032B} \SUPPORTEDANCH{\char"0052\char"032B} \SUPPORTEDANCH{\char"2C64\char"032B} \SUPPORTEDANCH{\char"0054\char"032B} \SUPPORTEDANCH{\char"01AC\char"032B} \SUPPORTEDANCH{\char"01AE\char"032B} \SUPPORTEDANCH{\char"0059\char"032B} \SUPPORTEDANCH{\char"01B3\char"032B} \SUPPORTEDANCH{\char"0061\char"032B} \SUPPORTEDANCH{\char"1E9A\char"032B} \SUPPORTEDANCH{\char"0062\char"032B} \SUPPORTEDANCH{\char"0183\char"032B} \SUPPORTEDANCH{\char"0063\char"032B} \SUPPORTEDANCH{\char"0188\char"032B} \SUPPORTEDANCH{\char"0064\char"032B} \SUPPORTEDANCH{\char"018C\char"032B} \SUPPORTEDANCH{\char"0221\char"032B} \SUPPORTEDANCH{\char"0068\char"032B} \SUPPORTEDANCH{\char"2C68\char"032B} \SUPPORTEDANCH{\char"006B\char"032B} \SUPPORTEDANCH{\char"0199\char"032B} \SUPPORTEDANCH{\char"2C6A\char"032B} \SUPPORTEDANCH{\char"006C\char"032B} \SUPPORTEDANCH{\char"0234\char"032B} \SUPPORTEDANCH{\char"006E\char"032B} \SUPPORTEDANCH{\char"0235\char"032B} \SUPPORTEDANCH{\char"0070\char"032B} \SUPPORTEDANCH{\char"01A5\char"032B} \SUPPORTEDANCH{\char"0071\char"032B} \SUPPORTEDANCH{\char"02A0\char"032B} \SUPPORTEDANCH{\char"0074\char"032B} \SUPPORTEDANCH{\char"01AB\char"032B} \SUPPORTEDANCH{\char"01AD\char"032B} \SUPPORTEDANCH{\char"0236\char"032B} \SUPPORTEDANCH{\char"0075\char"032B} \MISSINGGLYPH{\char"AB52\char"032B} \SUPPORTEDANCH{\char"0079\char"032B} \SUPPORTEDANCH{\char"01B4\char"032B} \MISSINGGLYPH{\char"1EFF\char"032B}

\SUPPORTEDANCH{\char"0048\char"032C} \SUPPORTEDANCH{\char"2C67\char"032C} \SUPPORTEDANCH{\char"004B\char"032C} \SUPPORTEDANCH{\char"2C69\char"032C} \SUPPORTEDANCH{\char"004D\char"032C} \MISSINGGLYPH{\char"2C6E\char"032C} \SUPPORTEDANCH{\char"0050\char"032C} \SUPPORTEDANCH{\char"01A4\char"032C} \SUPPORTEDANCH{\char"0052\char"032C} \SUPPORTEDANCH{\char"2C64\char"032C} \SUPPORTEDANCH{\char"0054\char"032C} \SUPPORTEDANCH{\char"01AC\char"032C} \SUPPORTEDANCH{\char"01AE\char"032C} \SUPPORTEDANCH{\char"0059\char"032C} \SUPPORTEDANCH{\char"01B3\char"032C} \SUPPORTEDANCH{\char"0061\char"032C} \SUPPORTEDANCH{\char"1E9A\char"032C} \SUPPORTEDANCH{\char"0062\char"032C} \SUPPORTEDANCH{\char"0183\char"032C} \SUPPORTEDANCH{\char"0063\char"032C} \SUPPORTEDANCH{\char"0188\char"032C} \SUPPORTEDANCH{\char"0064\char"032C} \SUPPORTEDANCH{\char"018C\char"032C} \SUPPORTEDANCH{\char"0221\char"032C} \SUPPORTEDANCH{\char"0068\char"032C} \SUPPORTEDANCH{\char"2C68\char"032C} \SUPPORTEDANCH{\char"006B\char"032C} \SUPPORTEDANCH{\char"0199\char"032C} \SUPPORTEDANCH{\char"2C6A\char"032C} \SUPPORTEDANCH{\char"006C\char"032C} \SUPPORTEDANCH{\char"0234\char"032C} \SUPPORTEDANCH{\char"006E\char"032C} \SUPPORTEDANCH{\char"0235\char"032C} \SUPPORTEDANCH{\char"0070\char"032C} \SUPPORTEDANCH{\char"01A5\char"032C} \SUPPORTEDANCH{\char"0071\char"032C} \SUPPORTEDANCH{\char"02A0\char"032C} \SUPPORTEDANCH{\char"0074\char"032C} \SUPPORTEDANCH{\char"01AB\char"032C} \SUPPORTEDANCH{\char"01AD\char"032C} \SUPPORTEDANCH{\char"0236\char"032C} \SUPPORTEDANCH{\char"0075\char"032C} \MISSINGGLYPH{\char"AB52\char"032C} \SUPPORTEDANCH{\char"0079\char"032C} \SUPPORTEDANCH{\char"01B4\char"032C} \MISSINGGLYPH{\char"1EFF\char"032C}

\SUPPORTEDANCH{\char"0048\char"032D} \SUPPORTEDANCH{\char"2C67\char"032D} \SUPPORTEDANCH{\char"004B\char"032D} \SUPPORTEDANCH{\char"2C69\char"032D} \SUPPORTEDANCH{\char"004D\char"032D} \MISSINGGLYPH{\char"2C6E\char"032D} \SUPPORTEDANCH{\char"0050\char"032D} \SUPPORTEDANCH{\char"01A4\char"032D} \SUPPORTEDANCH{\char"0052\char"032D} \SUPPORTEDANCH{\char"2C64\char"032D} \SUPPORTEDPREC{\char"0054\char"032D} \SUPPORTEDANCH{\char"01AC\char"032D} \SUPPORTEDANCH{\char"01AE\char"032D} \SUPPORTEDANCH{\char"0059\char"032D} \SUPPORTEDANCH{\char"01B3\char"032D} \SUPPORTEDANCH{\char"0061\char"032D} \SUPPORTEDANCH{\char"1E9A\char"032D} \SUPPORTEDANCH{\char"0062\char"032D} \SUPPORTEDANCH{\char"0183\char"032D} \SUPPORTEDANCH{\char"0063\char"032D} \SUPPORTEDANCH{\char"0188\char"032D} \SUPPORTEDPREC{\char"0064\char"032D} \SUPPORTEDANCH{\char"018C\char"032D} \SUPPORTEDANCH{\char"0221\char"032D} \SUPPORTEDANCH{\char"0068\char"032D} \SUPPORTEDANCH{\char"2C68\char"032D} \SUPPORTEDANCH{\char"006B\char"032D} \SUPPORTEDANCH{\char"0199\char"032D} \SUPPORTEDANCH{\char"2C6A\char"032D} \SUPPORTEDPREC{\char"006C\char"032D} \SUPPORTEDANCH{\char"0234\char"032D} \SUPPORTEDPREC{\char"006E\char"032D} \SUPPORTEDANCH{\char"0235\char"032D} \SUPPORTEDANCH{\char"0070\char"032D} \SUPPORTEDANCH{\char"01A5\char"032D} \SUPPORTEDANCH{\char"0071\char"032D} \SUPPORTEDANCH{\char"02A0\char"032D} \SUPPORTEDPREC{\char"0074\char"032D} \SUPPORTEDANCH{\char"01AB\char"032D} \SUPPORTEDANCH{\char"01AD\char"032D} \SUPPORTEDANCH{\char"0236\char"032D} \SUPPORTEDPREC{\char"0075\char"032D} \MISSINGGLYPH{\char"AB52\char"032D} \SUPPORTEDANCH{\char"0079\char"032D} \SUPPORTEDANCH{\char"01B4\char"032D} \MISSINGGLYPH{\char"1EFF\char"032D}

\SUPPORTEDPREC{\char"0048\char"032E} \SUPPORTEDANCH{\char"2C67\char"032E} \SUPPORTEDANCH{\char"004B\char"032E} \SUPPORTEDANCH{\char"2C69\char"032E} \SUPPORTEDANCH{\char"004D\char"032E} \MISSINGGLYPH{\char"2C6E\char"032E} \SUPPORTEDANCH{\char"0050\char"032E} \SUPPORTEDANCH{\char"01A4\char"032E} \SUPPORTEDANCH{\char"0052\char"032E} \SUPPORTEDANCH{\char"2C64\char"032E} \SUPPORTEDANCH{\char"0054\char"032E} \SUPPORTEDANCH{\char"01AC\char"032E} \SUPPORTEDANCH{\char"01AE\char"032E} \SUPPORTEDANCH{\char"0059\char"032E} \SUPPORTEDANCH{\char"01B3\char"032E} \SUPPORTEDANCH{\char"0061\char"032E} \SUPPORTEDANCH{\char"1E9A\char"032E} \SUPPORTEDANCH{\char"0062\char"032E} \SUPPORTEDANCH{\char"0183\char"032E} \SUPPORTEDANCH{\char"0063\char"032E} \SUPPORTEDANCH{\char"0188\char"032E} \SUPPORTEDANCH{\char"0064\char"032E} \SUPPORTEDANCH{\char"018C\char"032E} \SUPPORTEDANCH{\char"0221\char"032E} \SUPPORTEDPREC{\char"0068\char"032E} \SUPPORTEDANCH{\char"2C68\char"032E} \SUPPORTEDANCH{\char"006B\char"032E} \SUPPORTEDANCH{\char"0199\char"032E} \SUPPORTEDANCH{\char"2C6A\char"032E} \SUPPORTEDANCH{\char"006C\char"032E} \SUPPORTEDANCH{\char"0234\char"032E} \SUPPORTEDANCH{\char"006E\char"032E} \SUPPORTEDANCH{\char"0235\char"032E} \SUPPORTEDANCH{\char"0070\char"032E} \SUPPORTEDANCH{\char"01A5\char"032E} \SUPPORTEDANCH{\char"0071\char"032E} \SUPPORTEDANCH{\char"02A0\char"032E} \SUPPORTEDANCH{\char"0074\char"032E} \SUPPORTEDANCH{\char"01AB\char"032E} \SUPPORTEDANCH{\char"01AD\char"032E} \SUPPORTEDANCH{\char"0236\char"032E} \SUPPORTEDANCH{\char"0075\char"032E} \MISSINGGLYPH{\char"AB52\char"032E} \SUPPORTEDANCH{\char"0079\char"032E} \SUPPORTEDANCH{\char"01B4\char"032E} \MISSINGGLYPH{\char"1EFF\char"032E}

\SUPPORTEDANCH{\char"0048\char"032F} \SUPPORTEDANCH{\char"2C67\char"032F} \SUPPORTEDANCH{\char"004B\char"032F} \SUPPORTEDANCH{\char"2C69\char"032F} \SUPPORTEDANCH{\char"004D\char"032F} \MISSINGGLYPH{\char"2C6E\char"032F} \SUPPORTEDANCH{\char"0050\char"032F} \SUPPORTEDANCH{\char"01A4\char"032F} \SUPPORTEDANCH{\char"0052\char"032F} \SUPPORTEDANCH{\char"2C64\char"032F} \SUPPORTEDANCH{\char"0054\char"032F} \SUPPORTEDANCH{\char"01AC\char"032F} \SUPPORTEDANCH{\char"01AE\char"032F} \SUPPORTEDANCH{\char"0059\char"032F} \SUPPORTEDANCH{\char"01B3\char"032F} \SUPPORTEDANCH{\char"0061\char"032F} \SUPPORTEDANCH{\char"1E9A\char"032F} \SUPPORTEDANCH{\char"0062\char"032F} \SUPPORTEDANCH{\char"0183\char"032F} \SUPPORTEDANCH{\char"0063\char"032F} \SUPPORTEDANCH{\char"0188\char"032F} \SUPPORTEDANCH{\char"0064\char"032F} \SUPPORTEDANCH{\char"018C\char"032F} \SUPPORTEDANCH{\char"0221\char"032F} \SUPPORTEDANCH{\char"0068\char"032F} \SUPPORTEDANCH{\char"2C68\char"032F} \SUPPORTEDANCH{\char"006B\char"032F} \SUPPORTEDANCH{\char"0199\char"032F} \SUPPORTEDANCH{\char"2C6A\char"032F} \SUPPORTEDANCH{\char"006C\char"032F} \SUPPORTEDANCH{\char"0234\char"032F} \SUPPORTEDANCH{\char"006E\char"032F} \SUPPORTEDANCH{\char"0235\char"032F} \SUPPORTEDANCH{\char"0070\char"032F} \SUPPORTEDANCH{\char"01A5\char"032F} \SUPPORTEDANCH{\char"0071\char"032F} \SUPPORTEDANCH{\char"02A0\char"032F} \SUPPORTEDANCH{\char"0074\char"032F} \SUPPORTEDANCH{\char"01AB\char"032F} \SUPPORTEDANCH{\char"01AD\char"032F} \SUPPORTEDANCH{\char"0236\char"032F} \SUPPORTEDANCH{\char"0075\char"032F} \MISSINGGLYPH{\char"AB52\char"032F} \SUPPORTEDANCH{\char"0079\char"032F} \SUPPORTEDANCH{\char"01B4\char"032F} \MISSINGGLYPH{\char"1EFF\char"032F}

\SUPPORTEDANCH{\char"0048\char"031D} \SUPPORTEDANCH{\char"2C67\char"031D} \SUPPORTEDANCH{\char"004B\char"031D} \SUPPORTEDANCH{\char"2C69\char"031D} \SUPPORTEDANCH{\char"004D\char"031D} \MISSINGGLYPH{\char"2C6E\char"031D} \SUPPORTEDANCH{\char"0050\char"031D} \SUPPORTEDANCH{\char"01A4\char"031D} \SUPPORTEDANCH{\char"0052\char"031D} \SUPPORTEDANCH{\char"2C64\char"031D} \SUPPORTEDANCH{\char"0054\char"031D} \SUPPORTEDANCH{\char"01AC\char"031D} \SUPPORTEDANCH{\char"01AE\char"031D} \SUPPORTEDANCH{\char"0059\char"031D} \SUPPORTEDANCH{\char"01B3\char"031D} \SUPPORTEDANCH{\char"0061\char"031D} \SUPPORTEDANCH{\char"1E9A\char"031D} \SUPPORTEDANCH{\char"0062\char"031D} \SUPPORTEDANCH{\char"0183\char"031D} \SUPPORTEDANCH{\char"0063\char"031D} \SUPPORTEDANCH{\char"0188\char"031D} \SUPPORTEDANCH{\char"0064\char"031D} \SUPPORTEDANCH{\char"018C\char"031D} \SUPPORTEDANCH{\char"0221\char"031D} \SUPPORTEDANCH{\char"0068\char"031D} \SUPPORTEDANCH{\char"2C68\char"031D} \SUPPORTEDANCH{\char"006B\char"031D} \SUPPORTEDANCH{\char"0199\char"031D} \SUPPORTEDANCH{\char"2C6A\char"031D} \SUPPORTEDANCH{\char"006C\char"031D} \SUPPORTEDANCH{\char"0234\char"031D} \SUPPORTEDANCH{\char"006E\char"031D} \SUPPORTEDANCH{\char"0235\char"031D} \SUPPORTEDANCH{\char"0070\char"031D} \SUPPORTEDANCH{\char"01A5\char"031D} \SUPPORTEDANCH{\char"0071\char"031D} \SUPPORTEDANCH{\char"02A0\char"031D} \SUPPORTEDANCH{\char"0074\char"031D} \SUPPORTEDANCH{\char"01AB\char"031D} \SUPPORTEDANCH{\char"01AD\char"031D} \SUPPORTEDANCH{\char"0236\char"031D} \SUPPORTEDANCH{\char"0075\char"031D} \MISSINGGLYPH{\char"AB52\char"031D} \SUPPORTEDANCH{\char"0079\char"031D} \SUPPORTEDANCH{\char"01B4\char"031D} \MISSINGGLYPH{\char"1EFF\char"031D}

\SUPPORTEDANCH{\char"0048\char"031E} \SUPPORTEDANCH{\char"2C67\char"031E} \SUPPORTEDANCH{\char"004B\char"031E} \SUPPORTEDANCH{\char"2C69\char"031E} \SUPPORTEDANCH{\char"004D\char"031E} \MISSINGGLYPH{\char"2C6E\char"031E} \SUPPORTEDANCH{\char"0050\char"031E} \SUPPORTEDANCH{\char"01A4\char"031E} \SUPPORTEDANCH{\char"0052\char"031E} \SUPPORTEDANCH{\char"2C64\char"031E} \SUPPORTEDANCH{\char"0054\char"031E} \SUPPORTEDANCH{\char"01AC\char"031E} \SUPPORTEDANCH{\char"01AE\char"031E} \SUPPORTEDANCH{\char"0059\char"031E} \SUPPORTEDANCH{\char"01B3\char"031E} \SUPPORTEDANCH{\char"0061\char"031E} \SUPPORTEDANCH{\char"1E9A\char"031E} \SUPPORTEDANCH{\char"0062\char"031E} \SUPPORTEDANCH{\char"0183\char"031E} \SUPPORTEDANCH{\char"0063\char"031E} \SUPPORTEDANCH{\char"0188\char"031E} \SUPPORTEDANCH{\char"0064\char"031E} \SUPPORTEDANCH{\char"018C\char"031E} \SUPPORTEDANCH{\char"0221\char"031E} \SUPPORTEDANCH{\char"0068\char"031E} \SUPPORTEDANCH{\char"2C68\char"031E} \SUPPORTEDANCH{\char"006B\char"031E} \SUPPORTEDANCH{\char"0199\char"031E} \SUPPORTEDANCH{\char"2C6A\char"031E} \SUPPORTEDANCH{\char"006C\char"031E} \SUPPORTEDANCH{\char"0234\char"031E} \SUPPORTEDANCH{\char"006E\char"031E} \SUPPORTEDANCH{\char"0235\char"031E} \SUPPORTEDANCH{\char"0070\char"031E} \SUPPORTEDANCH{\char"01A5\char"031E} \SUPPORTEDANCH{\char"0071\char"031E} \SUPPORTEDANCH{\char"02A0\char"031E} \SUPPORTEDANCH{\char"0074\char"031E} \SUPPORTEDANCH{\char"01AB\char"031E} \SUPPORTEDANCH{\char"01AD\char"031E} \SUPPORTEDANCH{\char"0236\char"031E} \SUPPORTEDANCH{\char"0075\char"031E} \MISSINGGLYPH{\char"AB52\char"031E} \SUPPORTEDANCH{\char"0079\char"031E} \SUPPORTEDANCH{\char"01B4\char"031E} \MISSINGGLYPH{\char"1EFF\char"031E}

\SUPPORTEDANCH{\char"0048\char"031F} \SUPPORTEDANCH{\char"2C67\char"031F} \SUPPORTEDANCH{\char"004B\char"031F} \SUPPORTEDANCH{\char"2C69\char"031F} \SUPPORTEDANCH{\char"004D\char"031F} \MISSINGGLYPH{\char"2C6E\char"031F} \SUPPORTEDANCH{\char"0050\char"031F} \SUPPORTEDANCH{\char"01A4\char"031F} \SUPPORTEDANCH{\char"0052\char"031F} \SUPPORTEDANCH{\char"2C64\char"031F} \SUPPORTEDANCH{\char"0054\char"031F} \SUPPORTEDANCH{\char"01AC\char"031F} \SUPPORTEDANCH{\char"01AE\char"031F} \SUPPORTEDANCH{\char"0059\char"031F} \SUPPORTEDANCH{\char"01B3\char"031F} \SUPPORTEDANCH{\char"0061\char"031F} \SUPPORTEDANCH{\char"1E9A\char"031F} \SUPPORTEDANCH{\char"0062\char"031F} \SUPPORTEDANCH{\char"0183\char"031F} \SUPPORTEDANCH{\char"0063\char"031F} \SUPPORTEDANCH{\char"0188\char"031F} \SUPPORTEDANCH{\char"0064\char"031F} \SUPPORTEDANCH{\char"018C\char"031F} \SUPPORTEDANCH{\char"0221\char"031F} \SUPPORTEDANCH{\char"0068\char"031F} \SUPPORTEDANCH{\char"2C68\char"031F} \SUPPORTEDANCH{\char"006B\char"031F} \SUPPORTEDANCH{\char"0199\char"031F} \SUPPORTEDANCH{\char"2C6A\char"031F} \SUPPORTEDANCH{\char"006C\char"031F} \SUPPORTEDANCH{\char"0234\char"031F} \SUPPORTEDANCH{\char"006E\char"031F} \SUPPORTEDANCH{\char"0235\char"031F} \SUPPORTEDANCH{\char"0070\char"031F} \SUPPORTEDANCH{\char"01A5\char"031F} \SUPPORTEDANCH{\char"0071\char"031F} \SUPPORTEDANCH{\char"02A0\char"031F} \SUPPORTEDANCH{\char"0074\char"031F} \SUPPORTEDANCH{\char"01AB\char"031F} \SUPPORTEDANCH{\char"01AD\char"031F} \SUPPORTEDANCH{\char"0236\char"031F} \SUPPORTEDANCH{\char"0075\char"031F} \MISSINGGLYPH{\char"AB52\char"031F} \SUPPORTEDANCH{\char"0079\char"031F} \SUPPORTEDANCH{\char"01B4\char"031F} \MISSINGGLYPH{\char"1EFF\char"031F}

\SUPPORTEDANCH{\char"0048\char"0330} \SUPPORTEDANCH{\char"2C67\char"0330} \SUPPORTEDANCH{\char"004B\char"0330} \SUPPORTEDANCH{\char"2C69\char"0330} \SUPPORTEDANCH{\char"004D\char"0330} \MISSINGGLYPH{\char"2C6E\char"0330} \SUPPORTEDANCH{\char"0050\char"0330} \SUPPORTEDANCH{\char"01A4\char"0330} \SUPPORTEDANCH{\char"0052\char"0330} \SUPPORTEDANCH{\char"2C64\char"0330} \SUPPORTEDANCH{\char"0054\char"0330} \SUPPORTEDANCH{\char"01AC\char"0330} \SUPPORTEDANCH{\char"01AE\char"0330} \SUPPORTEDANCH{\char"0059\char"0330} \SUPPORTEDANCH{\char"01B3\char"0330} \SUPPORTEDANCH{\char"0061\char"0330} \SUPPORTEDANCH{\char"1E9A\char"0330} \SUPPORTEDANCH{\char"0062\char"0330} \SUPPORTEDANCH{\char"0183\char"0330} \SUPPORTEDANCH{\char"0063\char"0330} \SUPPORTEDANCH{\char"0188\char"0330} \SUPPORTEDANCH{\char"0064\char"0330} \SUPPORTEDANCH{\char"018C\char"0330} \SUPPORTEDANCH{\char"0221\char"0330} \SUPPORTEDANCH{\char"0068\char"0330} \SUPPORTEDANCH{\char"2C68\char"0330} \SUPPORTEDANCH{\char"006B\char"0330} \SUPPORTEDANCH{\char"0199\char"0330} \SUPPORTEDANCH{\char"2C6A\char"0330} \SUPPORTEDANCH{\char"006C\char"0330} \SUPPORTEDANCH{\char"0234\char"0330} \SUPPORTEDANCH{\char"006E\char"0330} \SUPPORTEDANCH{\char"0235\char"0330} \SUPPORTEDANCH{\char"0070\char"0330} \SUPPORTEDANCH{\char"01A5\char"0330} \SUPPORTEDANCH{\char"0071\char"0330} \SUPPORTEDANCH{\char"02A0\char"0330} \SUPPORTEDANCH{\char"0074\char"0330} \SUPPORTEDANCH{\char"01AB\char"0330} \SUPPORTEDANCH{\char"01AD\char"0330} \SUPPORTEDANCH{\char"0236\char"0330} \SUPPORTEDPREC{\char"0075\char"0330} \MISSINGGLYPH{\char"AB52\char"0330} \SUPPORTEDANCH{\char"0079\char"0330} \SUPPORTEDANCH{\char"01B4\char"0330} \MISSINGGLYPH{\char"1EFF\char"0330}

\SUPPORTEDANCH{\char"0048\char"0331} \SUPPORTEDANCH{\char"2C67\char"0331} \SUPPORTEDPREC{\char"004B\char"0331} \SUPPORTEDANCH{\char"2C69\char"0331} \SUPPORTEDANCH{\char"004D\char"0331} \MISSINGGLYPH{\char"2C6E\char"0331} \SUPPORTEDANCH{\char"0050\char"0331} \SUPPORTEDANCH{\char"01A4\char"0331} \SUPPORTEDPREC{\char"0052\char"0331} \SUPPORTEDANCH{\char"2C64\char"0331} \SUPPORTEDPREC{\char"0054\char"0331} \SUPPORTEDANCH{\char"01AC\char"0331} \SUPPORTEDANCH{\char"01AE\char"0331} \SUPPORTEDANCH{\char"0059\char"0331} \SUPPORTEDANCH{\char"01B3\char"0331} \SUPPORTEDANCH{\char"0061\char"0331} \SUPPORTEDANCH{\char"1E9A\char"0331} \SUPPORTEDPREC{\char"0062\char"0331} \SUPPORTEDANCH{\char"0183\char"0331} \SUPPORTEDANCH{\char"0063\char"0331} \SUPPORTEDANCH{\char"0188\char"0331} \SUPPORTEDPREC{\char"0064\char"0331} \SUPPORTEDANCH{\char"018C\char"0331} \SUPPORTEDANCH{\char"0221\char"0331} \SUPPORTEDPREC{\char"0068\char"0331} \SUPPORTEDANCH{\char"2C68\char"0331} \SUPPORTEDPREC{\char"006B\char"0331} \SUPPORTEDANCH{\char"0199\char"0331} \SUPPORTEDANCH{\char"2C6A\char"0331} \SUPPORTEDPREC{\char"006C\char"0331} \SUPPORTEDANCH{\char"0234\char"0331} \SUPPORTEDPREC{\char"006E\char"0331} \SUPPORTEDANCH{\char"0235\char"0331} \SUPPORTEDANCH{\char"0070\char"0331} \SUPPORTEDANCH{\char"01A5\char"0331} \SUPPORTEDANCH{\char"0071\char"0331} \SUPPORTEDANCH{\char"02A0\char"0331} \SUPPORTEDPREC{\char"0074\char"0331} \SUPPORTEDANCH{\char"01AB\char"0331} \SUPPORTEDANCH{\char"01AD\char"0331} \SUPPORTEDANCH{\char"0236\char"0331} \SUPPORTEDANCH{\char"0075\char"0331} \MISSINGGLYPH{\char"AB52\char"0331} \SUPPORTEDANCH{\char"0079\char"0331} \SUPPORTEDANCH{\char"01B4\char"0331} \MISSINGGLYPH{\char"1EFF\char"0331}

\SUPPORTEDANCH{\char"0048\char"0333} \SUPPORTEDANCH{\char"2C67\char"0333} \SUPPORTEDANCH{\char"004B\char"0333} \SUPPORTEDANCH{\char"2C69\char"0333} \SUPPORTEDANCH{\char"004D\char"0333} \MISSINGGLYPH{\char"2C6E\char"0333} \SUPPORTEDANCH{\char"0050\char"0333} \SUPPORTEDANCH{\char"01A4\char"0333} \SUPPORTEDANCH{\char"0052\char"0333} \SUPPORTEDANCH{\char"2C64\char"0333} \SUPPORTEDANCH{\char"0054\char"0333} \SUPPORTEDANCH{\char"01AC\char"0333} \SUPPORTEDANCH{\char"01AE\char"0333} \SUPPORTEDANCH{\char"0059\char"0333} \SUPPORTEDANCH{\char"01B3\char"0333} \SUPPORTEDANCH{\char"0061\char"0333} \SUPPORTEDANCH{\char"1E9A\char"0333} \SUPPORTEDANCH{\char"0062\char"0333} \SUPPORTEDANCH{\char"0183\char"0333} \SUPPORTEDANCH{\char"0063\char"0333} \SUPPORTEDANCH{\char"0188\char"0333} \SUPPORTEDANCH{\char"0064\char"0333} \SUPPORTEDANCH{\char"018C\char"0333} \SUPPORTEDANCH{\char"0221\char"0333} \SUPPORTEDANCH{\char"0068\char"0333} \SUPPORTEDANCH{\char"2C68\char"0333} \SUPPORTEDANCH{\char"006B\char"0333} \SUPPORTEDANCH{\char"0199\char"0333} \SUPPORTEDANCH{\char"2C6A\char"0333} \SUPPORTEDANCH{\char"006C\char"0333} \SUPPORTEDANCH{\char"0234\char"0333} \SUPPORTEDANCH{\char"006E\char"0333} \SUPPORTEDANCH{\char"0235\char"0333} \SUPPORTEDANCH{\char"0070\char"0333} \SUPPORTEDANCH{\char"01A5\char"0333} \SUPPORTEDANCH{\char"0071\char"0333} \SUPPORTEDANCH{\char"02A0\char"0333} \SUPPORTEDANCH{\char"0074\char"0333} \SUPPORTEDANCH{\char"01AB\char"0333} \SUPPORTEDANCH{\char"01AD\char"0333} \SUPPORTEDANCH{\char"0236\char"0333} \SUPPORTEDANCH{\char"0075\char"0333} \MISSINGGLYPH{\char"AB52\char"0333} \SUPPORTEDANCH{\char"0079\char"0333} \SUPPORTEDANCH{\char"01B4\char"0333} \MISSINGGLYPH{\char"1EFF\char"0333}

\SUPPORTEDANCH{\char"0048\char"0339} \SUPPORTEDANCH{\char"2C67\char"0339} \SUPPORTEDANCH{\char"004B\char"0339} \SUPPORTEDANCH{\char"2C69\char"0339} \SUPPORTEDANCH{\char"004D\char"0339} \MISSINGGLYPH{\char"2C6E\char"0339} \SUPPORTEDANCH{\char"0050\char"0339} \SUPPORTEDANCH{\char"01A4\char"0339} \SUPPORTEDANCH{\char"0052\char"0339} \SUPPORTEDANCH{\char"2C64\char"0339} \SUPPORTEDANCH{\char"0054\char"0339} \SUPPORTEDANCH{\char"01AC\char"0339} \SUPPORTEDANCH{\char"01AE\char"0339} \SUPPORTEDANCH{\char"0059\char"0339} \SUPPORTEDANCH{\char"01B3\char"0339} \SUPPORTEDANCH{\char"0061\char"0339} \SUPPORTEDANCH{\char"1E9A\char"0339} \SUPPORTEDANCH{\char"0062\char"0339} \SUPPORTEDANCH{\char"0183\char"0339} \SUPPORTEDANCH{\char"0063\char"0339} \SUPPORTEDANCH{\char"0188\char"0339} \SUPPORTEDANCH{\char"0064\char"0339} \SUPPORTEDANCH{\char"018C\char"0339} \SUPPORTEDANCH{\char"0221\char"0339} \SUPPORTEDANCH{\char"0068\char"0339} \SUPPORTEDANCH{\char"2C68\char"0339} \SUPPORTEDANCH{\char"006B\char"0339} \SUPPORTEDANCH{\char"0199\char"0339} \SUPPORTEDANCH{\char"2C6A\char"0339} \SUPPORTEDANCH{\char"006C\char"0339} \SUPPORTEDANCH{\char"0234\char"0339} \SUPPORTEDANCH{\char"006E\char"0339} \SUPPORTEDANCH{\char"0235\char"0339} \SUPPORTEDANCH{\char"0070\char"0339} \SUPPORTEDANCH{\char"01A5\char"0339} \SUPPORTEDANCH{\char"0071\char"0339} \SUPPORTEDANCH{\char"02A0\char"0339} \SUPPORTEDANCH{\char"0074\char"0339} \SUPPORTEDANCH{\char"01AB\char"0339} \SUPPORTEDANCH{\char"01AD\char"0339} \SUPPORTEDANCH{\char"0236\char"0339} \SUPPORTEDANCH{\char"0075\char"0339} \MISSINGGLYPH{\char"AB52\char"0339} \SUPPORTEDANCH{\char"0079\char"0339} \SUPPORTEDANCH{\char"01B4\char"0339} \MISSINGGLYPH{\char"1EFF\char"0339}

\SUPPORTEDANCH{\char"0048\char"033A} \SUPPORTEDANCH{\char"2C67\char"033A} \SUPPORTEDANCH{\char"004B\char"033A} \SUPPORTEDANCH{\char"2C69\char"033A} \SUPPORTEDANCH{\char"004D\char"033A} \MISSINGGLYPH{\char"2C6E\char"033A} \SUPPORTEDANCH{\char"0050\char"033A} \SUPPORTEDANCH{\char"01A4\char"033A} \SUPPORTEDANCH{\char"0052\char"033A} \SUPPORTEDANCH{\char"2C64\char"033A} \SUPPORTEDANCH{\char"0054\char"033A} \SUPPORTEDANCH{\char"01AC\char"033A} \SUPPORTEDANCH{\char"01AE\char"033A} \SUPPORTEDANCH{\char"0059\char"033A} \SUPPORTEDANCH{\char"01B3\char"033A} \SUPPORTEDANCH{\char"0061\char"033A} \SUPPORTEDANCH{\char"1E9A\char"033A} \SUPPORTEDANCH{\char"0062\char"033A} \SUPPORTEDANCH{\char"0183\char"033A} \SUPPORTEDANCH{\char"0063\char"033A} \SUPPORTEDANCH{\char"0188\char"033A} \SUPPORTEDANCH{\char"0064\char"033A} \SUPPORTEDANCH{\char"018C\char"033A} \SUPPORTEDANCH{\char"0221\char"033A} \SUPPORTEDANCH{\char"0068\char"033A} \SUPPORTEDANCH{\char"2C68\char"033A} \SUPPORTEDANCH{\char"006B\char"033A} \SUPPORTEDANCH{\char"0199\char"033A} \SUPPORTEDANCH{\char"2C6A\char"033A} \SUPPORTEDANCH{\char"006C\char"033A} \SUPPORTEDANCH{\char"0234\char"033A} \SUPPORTEDANCH{\char"006E\char"033A} \SUPPORTEDANCH{\char"0235\char"033A} \SUPPORTEDANCH{\char"0070\char"033A} \SUPPORTEDANCH{\char"01A5\char"033A} \SUPPORTEDANCH{\char"0071\char"033A} \SUPPORTEDANCH{\char"02A0\char"033A} \SUPPORTEDANCH{\char"0074\char"033A} \SUPPORTEDANCH{\char"01AB\char"033A} \SUPPORTEDANCH{\char"01AD\char"033A} \SUPPORTEDANCH{\char"0236\char"033A} \SUPPORTEDANCH{\char"0075\char"033A} \MISSINGGLYPH{\char"AB52\char"033A} \SUPPORTEDANCH{\char"0079\char"033A} \SUPPORTEDANCH{\char"01B4\char"033A} \MISSINGGLYPH{\char"1EFF\char"033A}

\SUPPORTEDANCH{\char"0048\char"033B} \SUPPORTEDANCH{\char"2C67\char"033B} \SUPPORTEDANCH{\char"004B\char"033B} \SUPPORTEDANCH{\char"2C69\char"033B} \SUPPORTEDANCH{\char"004D\char"033B} \MISSINGGLYPH{\char"2C6E\char"033B} \SUPPORTEDANCH{\char"0050\char"033B} \SUPPORTEDANCH{\char"01A4\char"033B} \SUPPORTEDANCH{\char"0052\char"033B} \SUPPORTEDANCH{\char"2C64\char"033B} \SUPPORTEDANCH{\char"0054\char"033B} \SUPPORTEDANCH{\char"01AC\char"033B} \SUPPORTEDANCH{\char"01AE\char"033B} \SUPPORTEDANCH{\char"0059\char"033B} \SUPPORTEDANCH{\char"01B3\char"033B} \SUPPORTEDANCH{\char"0061\char"033B} \SUPPORTEDANCH{\char"1E9A\char"033B} \SUPPORTEDANCH{\char"0062\char"033B} \SUPPORTEDANCH{\char"0183\char"033B} \SUPPORTEDANCH{\char"0063\char"033B} \SUPPORTEDANCH{\char"0188\char"033B} \SUPPORTEDANCH{\char"0064\char"033B} \SUPPORTEDANCH{\char"018C\char"033B} \SUPPORTEDANCH{\char"0221\char"033B} \SUPPORTEDANCH{\char"0068\char"033B} \SUPPORTEDANCH{\char"2C68\char"033B} \SUPPORTEDANCH{\char"006B\char"033B} \SUPPORTEDANCH{\char"0199\char"033B} \SUPPORTEDANCH{\char"2C6A\char"033B} \SUPPORTEDANCH{\char"006C\char"033B} \SUPPORTEDANCH{\char"0234\char"033B} \SUPPORTEDANCH{\char"006E\char"033B} \SUPPORTEDANCH{\char"0235\char"033B} \SUPPORTEDANCH{\char"0070\char"033B} \SUPPORTEDANCH{\char"01A5\char"033B} \SUPPORTEDANCH{\char"0071\char"033B} \SUPPORTEDANCH{\char"02A0\char"033B} \SUPPORTEDANCH{\char"0074\char"033B} \SUPPORTEDANCH{\char"01AB\char"033B} \SUPPORTEDANCH{\char"01AD\char"033B} \SUPPORTEDANCH{\char"0236\char"033B} \SUPPORTEDANCH{\char"0075\char"033B} \MISSINGGLYPH{\char"AB52\char"033B} \SUPPORTEDANCH{\char"0079\char"033B} \SUPPORTEDANCH{\char"01B4\char"033B} \MISSINGGLYPH{\char"1EFF\char"033B}

\SUPPORTEDANCH{\char"0048\char"033C} \SUPPORTEDANCH{\char"2C67\char"033C} \SUPPORTEDANCH{\char"004B\char"033C} \SUPPORTEDANCH{\char"2C69\char"033C} \SUPPORTEDANCH{\char"004D\char"033C} \MISSINGGLYPH{\char"2C6E\char"033C} \SUPPORTEDANCH{\char"0050\char"033C} \SUPPORTEDANCH{\char"01A4\char"033C} \SUPPORTEDANCH{\char"0052\char"033C} \SUPPORTEDANCH{\char"2C64\char"033C} \SUPPORTEDANCH{\char"0054\char"033C} \SUPPORTEDANCH{\char"01AC\char"033C} \SUPPORTEDANCH{\char"01AE\char"033C} \SUPPORTEDANCH{\char"0059\char"033C} \SUPPORTEDANCH{\char"01B3\char"033C} \SUPPORTEDANCH{\char"0061\char"033C} \SUPPORTEDANCH{\char"1E9A\char"033C} \SUPPORTEDANCH{\char"0062\char"033C} \SUPPORTEDANCH{\char"0183\char"033C} \SUPPORTEDANCH{\char"0063\char"033C} \SUPPORTEDANCH{\char"0188\char"033C} \SUPPORTEDANCH{\char"0064\char"033C} \SUPPORTEDANCH{\char"018C\char"033C} \SUPPORTEDANCH{\char"0221\char"033C} \SUPPORTEDANCH{\char"0068\char"033C} \SUPPORTEDANCH{\char"2C68\char"033C} \SUPPORTEDANCH{\char"006B\char"033C} \SUPPORTEDANCH{\char"0199\char"033C} \SUPPORTEDANCH{\char"2C6A\char"033C} \SUPPORTEDANCH{\char"006C\char"033C} \SUPPORTEDANCH{\char"0234\char"033C} \SUPPORTEDANCH{\char"006E\char"033C} \SUPPORTEDANCH{\char"0235\char"033C} \SUPPORTEDANCH{\char"0070\char"033C} \SUPPORTEDANCH{\char"01A5\char"033C} \SUPPORTEDANCH{\char"0071\char"033C} \SUPPORTEDANCH{\char"02A0\char"033C} \SUPPORTEDANCH{\char"0074\char"033C} \SUPPORTEDANCH{\char"01AB\char"033C} \SUPPORTEDANCH{\char"01AD\char"033C} \SUPPORTEDANCH{\char"0236\char"033C} \SUPPORTEDANCH{\char"0075\char"033C} \MISSINGGLYPH{\char"AB52\char"033C} \SUPPORTEDANCH{\char"0079\char"033C} \SUPPORTEDANCH{\char"01B4\char"033C} \MISSINGGLYPH{\char"1EFF\char"033C}

\SUPPORTEDANCH{\char"0048\char"0347} \SUPPORTEDANCH{\char"2C67\char"0347} \SUPPORTEDANCH{\char"004B\char"0347} \SUPPORTEDANCH{\char"2C69\char"0347} \SUPPORTEDANCH{\char"004D\char"0347} \MISSINGGLYPH{\char"2C6E\char"0347} \SUPPORTEDANCH{\char"0050\char"0347} \SUPPORTEDANCH{\char"01A4\char"0347} \SUPPORTEDANCH{\char"0052\char"0347} \SUPPORTEDANCH{\char"2C64\char"0347} \SUPPORTEDANCH{\char"0054\char"0347} \SUPPORTEDANCH{\char"01AC\char"0347} \SUPPORTEDANCH{\char"01AE\char"0347} \SUPPORTEDANCH{\char"0059\char"0347} \SUPPORTEDANCH{\char"01B3\char"0347} \SUPPORTEDANCH{\char"0061\char"0347} \SUPPORTEDANCH{\char"1E9A\char"0347} \SUPPORTEDANCH{\char"0062\char"0347} \SUPPORTEDANCH{\char"0183\char"0347} \SUPPORTEDANCH{\char"0063\char"0347} \SUPPORTEDANCH{\char"0188\char"0347} \SUPPORTEDANCH{\char"0064\char"0347} \SUPPORTEDANCH{\char"018C\char"0347} \SUPPORTEDANCH{\char"0221\char"0347} \SUPPORTEDANCH{\char"0068\char"0347} \SUPPORTEDANCH{\char"2C68\char"0347} \SUPPORTEDANCH{\char"006B\char"0347} \SUPPORTEDANCH{\char"0199\char"0347} \SUPPORTEDANCH{\char"2C6A\char"0347} \SUPPORTEDANCH{\char"006C\char"0347} \SUPPORTEDANCH{\char"0234\char"0347} \SUPPORTEDANCH{\char"006E\char"0347} \SUPPORTEDANCH{\char"0235\char"0347} \SUPPORTEDANCH{\char"0070\char"0347} \SUPPORTEDANCH{\char"01A5\char"0347} \SUPPORTEDANCH{\char"0071\char"0347} \SUPPORTEDANCH{\char"02A0\char"0347} \SUPPORTEDANCH{\char"0074\char"0347} \SUPPORTEDANCH{\char"01AB\char"0347} \SUPPORTEDANCH{\char"01AD\char"0347} \SUPPORTEDANCH{\char"0236\char"0347} \SUPPORTEDANCH{\char"0075\char"0347} \MISSINGGLYPH{\char"AB52\char"0347} \SUPPORTEDANCH{\char"0079\char"0347} \SUPPORTEDANCH{\char"01B4\char"0347} \MISSINGGLYPH{\char"1EFF\char"0347}

\SUPPORTEDANCH{\char"0048\char"0348} \SUPPORTEDANCH{\char"2C67\char"0348} \SUPPORTEDANCH{\char"004B\char"0348} \SUPPORTEDANCH{\char"2C69\char"0348} \SUPPORTEDANCH{\char"004D\char"0348} \MISSINGGLYPH{\char"2C6E\char"0348} \SUPPORTEDANCH{\char"0050\char"0348} \SUPPORTEDANCH{\char"01A4\char"0348} \SUPPORTEDANCH{\char"0052\char"0348} \SUPPORTEDANCH{\char"2C64\char"0348} \SUPPORTEDANCH{\char"0054\char"0348} \SUPPORTEDANCH{\char"01AC\char"0348} \SUPPORTEDANCH{\char"01AE\char"0348} \SUPPORTEDANCH{\char"0059\char"0348} \SUPPORTEDANCH{\char"01B3\char"0348} \SUPPORTEDANCH{\char"0061\char"0348} \SUPPORTEDANCH{\char"1E9A\char"0348} \SUPPORTEDANCH{\char"0062\char"0348} \SUPPORTEDANCH{\char"0183\char"0348} \SUPPORTEDANCH{\char"0063\char"0348} \SUPPORTEDANCH{\char"0188\char"0348} \SUPPORTEDANCH{\char"0064\char"0348} \SUPPORTEDANCH{\char"018C\char"0348} \SUPPORTEDANCH{\char"0221\char"0348} \SUPPORTEDANCH{\char"0068\char"0348} \SUPPORTEDANCH{\char"2C68\char"0348} \SUPPORTEDANCH{\char"006B\char"0348} \SUPPORTEDANCH{\char"0199\char"0348} \SUPPORTEDANCH{\char"2C6A\char"0348} \SUPPORTEDANCH{\char"006C\char"0348} \SUPPORTEDANCH{\char"0234\char"0348} \SUPPORTEDANCH{\char"006E\char"0348} \SUPPORTEDANCH{\char"0235\char"0348} \SUPPORTEDANCH{\char"0070\char"0348} \SUPPORTEDANCH{\char"01A5\char"0348} \SUPPORTEDANCH{\char"0071\char"0348} \SUPPORTEDANCH{\char"02A0\char"0348} \SUPPORTEDANCH{\char"0074\char"0348} \SUPPORTEDANCH{\char"01AB\char"0348} \SUPPORTEDANCH{\char"01AD\char"0348} \SUPPORTEDANCH{\char"0236\char"0348} \SUPPORTEDANCH{\char"0075\char"0348} \MISSINGGLYPH{\char"AB52\char"0348} \SUPPORTEDANCH{\char"0079\char"0348} \SUPPORTEDANCH{\char"01B4\char"0348} \MISSINGGLYPH{\char"1EFF\char"0348}

\SUPPORTEDANCH{\char"0048\char"0349} \SUPPORTEDANCH{\char"2C67\char"0349} \SUPPORTEDANCH{\char"004B\char"0349} \SUPPORTEDANCH{\char"2C69\char"0349} \SUPPORTEDANCH{\char"004D\char"0349} \MISSINGGLYPH{\char"2C6E\char"0349} \SUPPORTEDANCH{\char"0050\char"0349} \SUPPORTEDANCH{\char"01A4\char"0349} \SUPPORTEDANCH{\char"0052\char"0349} \SUPPORTEDANCH{\char"2C64\char"0349} \SUPPORTEDANCH{\char"0054\char"0349} \SUPPORTEDANCH{\char"01AC\char"0349} \SUPPORTEDANCH{\char"01AE\char"0349} \SUPPORTEDANCH{\char"0059\char"0349} \SUPPORTEDANCH{\char"01B3\char"0349} \SUPPORTEDANCH{\char"0061\char"0349} \SUPPORTEDANCH{\char"1E9A\char"0349} \SUPPORTEDANCH{\char"0062\char"0349} \SUPPORTEDANCH{\char"0183\char"0349} \SUPPORTEDANCH{\char"0063\char"0349} \SUPPORTEDANCH{\char"0188\char"0349} \SUPPORTEDANCH{\char"0064\char"0349} \SUPPORTEDANCH{\char"018C\char"0349} \SUPPORTEDANCH{\char"0221\char"0349} \SUPPORTEDANCH{\char"0068\char"0349} \SUPPORTEDANCH{\char"2C68\char"0349} \SUPPORTEDANCH{\char"006B\char"0349} \SUPPORTEDANCH{\char"0199\char"0349} \SUPPORTEDANCH{\char"2C6A\char"0349} \SUPPORTEDANCH{\char"006C\char"0349} \SUPPORTEDANCH{\char"0234\char"0349} \SUPPORTEDANCH{\char"006E\char"0349} \SUPPORTEDANCH{\char"0235\char"0349} \SUPPORTEDANCH{\char"0070\char"0349} \SUPPORTEDANCH{\char"01A5\char"0349} \SUPPORTEDANCH{\char"0071\char"0349} \SUPPORTEDANCH{\char"02A0\char"0349} \SUPPORTEDANCH{\char"0074\char"0349} \SUPPORTEDANCH{\char"01AB\char"0349} \SUPPORTEDANCH{\char"01AD\char"0349} \SUPPORTEDANCH{\char"0236\char"0349} \SUPPORTEDANCH{\char"0075\char"0349} \MISSINGGLYPH{\char"AB52\char"0349} \SUPPORTEDANCH{\char"0079\char"0349} \SUPPORTEDANCH{\char"01B4\char"0349} \MISSINGGLYPH{\char"1EFF\char"0349}

\SUPPORTEDANCH{\char"0048\char"034D} \SUPPORTEDANCH{\char"2C67\char"034D} \SUPPORTEDANCH{\char"004B\char"034D} \SUPPORTEDANCH{\char"2C69\char"034D} \SUPPORTEDANCH{\char"004D\char"034D} \MISSINGGLYPH{\char"2C6E\char"034D} \SUPPORTEDANCH{\char"0050\char"034D} \SUPPORTEDANCH{\char"01A4\char"034D} \SUPPORTEDANCH{\char"0052\char"034D} \SUPPORTEDANCH{\char"2C64\char"034D} \SUPPORTEDANCH{\char"0054\char"034D} \SUPPORTEDANCH{\char"01AC\char"034D} \SUPPORTEDANCH{\char"01AE\char"034D} \SUPPORTEDANCH{\char"0059\char"034D} \SUPPORTEDANCH{\char"01B3\char"034D} \SUPPORTEDANCH{\char"0061\char"034D} \SUPPORTEDANCH{\char"1E9A\char"034D} \SUPPORTEDANCH{\char"0062\char"034D} \SUPPORTEDANCH{\char"0183\char"034D} \SUPPORTEDANCH{\char"0063\char"034D} \SUPPORTEDANCH{\char"0188\char"034D} \SUPPORTEDANCH{\char"0064\char"034D} \SUPPORTEDANCH{\char"018C\char"034D} \SUPPORTEDANCH{\char"0221\char"034D} \SUPPORTEDANCH{\char"0068\char"034D} \SUPPORTEDANCH{\char"2C68\char"034D} \SUPPORTEDANCH{\char"006B\char"034D} \SUPPORTEDANCH{\char"0199\char"034D} \SUPPORTEDANCH{\char"2C6A\char"034D} \SUPPORTEDANCH{\char"006C\char"034D} \SUPPORTEDANCH{\char"0234\char"034D} \SUPPORTEDANCH{\char"006E\char"034D} \SUPPORTEDANCH{\char"0235\char"034D} \SUPPORTEDANCH{\char"0070\char"034D} \SUPPORTEDANCH{\char"01A5\char"034D} \SUPPORTEDANCH{\char"0071\char"034D} \SUPPORTEDANCH{\char"02A0\char"034D} \SUPPORTEDANCH{\char"0074\char"034D} \SUPPORTEDANCH{\char"01AB\char"034D} \SUPPORTEDANCH{\char"01AD\char"034D} \SUPPORTEDANCH{\char"0236\char"034D} \SUPPORTEDANCH{\char"0075\char"034D} \MISSINGGLYPH{\char"AB52\char"034D} \SUPPORTEDANCH{\char"0079\char"034D} \SUPPORTEDANCH{\char"01B4\char"034D} \MISSINGGLYPH{\char"1EFF\char"034D}

\SUPPORTEDANCH{\char"0048\char"034E} \SUPPORTEDANCH{\char"2C67\char"034E} \SUPPORTEDANCH{\char"004B\char"034E} \SUPPORTEDANCH{\char"2C69\char"034E} \SUPPORTEDANCH{\char"004D\char"034E} \MISSINGGLYPH{\char"2C6E\char"034E} \SUPPORTEDANCH{\char"0050\char"034E} \SUPPORTEDANCH{\char"01A4\char"034E} \SUPPORTEDANCH{\char"0052\char"034E} \SUPPORTEDANCH{\char"2C64\char"034E} \SUPPORTEDANCH{\char"0054\char"034E} \SUPPORTEDANCH{\char"01AC\char"034E} \SUPPORTEDANCH{\char"01AE\char"034E} \SUPPORTEDANCH{\char"0059\char"034E} \SUPPORTEDANCH{\char"01B3\char"034E} \SUPPORTEDANCH{\char"0061\char"034E} \SUPPORTEDANCH{\char"1E9A\char"034E} \SUPPORTEDANCH{\char"0062\char"034E} \SUPPORTEDANCH{\char"0183\char"034E} \SUPPORTEDANCH{\char"0063\char"034E} \SUPPORTEDANCH{\char"0188\char"034E} \SUPPORTEDANCH{\char"0064\char"034E} \SUPPORTEDANCH{\char"018C\char"034E} \SUPPORTEDANCH{\char"0221\char"034E} \SUPPORTEDANCH{\char"0068\char"034E} \SUPPORTEDANCH{\char"2C68\char"034E} \SUPPORTEDANCH{\char"006B\char"034E} \SUPPORTEDANCH{\char"0199\char"034E} \SUPPORTEDANCH{\char"2C6A\char"034E} \SUPPORTEDANCH{\char"006C\char"034E} \SUPPORTEDANCH{\char"0234\char"034E} \SUPPORTEDANCH{\char"006E\char"034E} \SUPPORTEDANCH{\char"0235\char"034E} \SUPPORTEDANCH{\char"0070\char"034E} \SUPPORTEDANCH{\char"01A5\char"034E} \SUPPORTEDANCH{\char"0071\char"034E} \SUPPORTEDANCH{\char"02A0\char"034E} \SUPPORTEDANCH{\char"0074\char"034E} \SUPPORTEDANCH{\char"01AB\char"034E} \SUPPORTEDANCH{\char"01AD\char"034E} \SUPPORTEDANCH{\char"0236\char"034E} \SUPPORTEDANCH{\char"0075\char"034E} \MISSINGGLYPH{\char"AB52\char"034E} \SUPPORTEDANCH{\char"0079\char"034E} \SUPPORTEDANCH{\char"01B4\char"034E} \MISSINGGLYPH{\char"1EFF\char"034E}

}

\newpage
\subsection*{Semibold italic, Group 4, below, Below}

{\bfseries\itshape
% grid. columns are bases, rows are marks.
\SUPPORTEDANCH{\char"0048\char"0316} \SUPPORTEDFALL{\char"2C67\char"0316} \SUPPORTEDANCH{\char"004B\char"0316} \SUPPORTEDANCH{\char"2C69\char"0316} \SUPPORTEDANCH{\char"004D\char"0316} \MISSINGGLYPH{\char"2C6E\char"0316} \SUPPORTEDANCH{\char"0050\char"0316} \SUPPORTEDANCH{\char"01A4\char"0316} \SUPPORTEDANCH{\char"0052\char"0316} \SUPPORTEDFALL{\char"2C64\char"0316} \SUPPORTEDANCH{\char"0054\char"0316} \SUPPORTEDFALL{\char"01AC\char"0316} \SUPPORTEDFALL{\char"01AE\char"0316} \SUPPORTEDANCH{\char"0059\char"0316} \SUPPORTEDFALL{\char"01B3\char"0316} \SUPPORTEDANCH{\char"0061\char"0316} \SUPPORTEDFALL{\char"1E9A\char"0316} \SUPPORTEDANCH{\char"0062\char"0316} \SUPPORTEDANCH{\char"0183\char"0316} \SUPPORTEDANCH{\char"0063\char"0316} \SUPPORTEDFALL{\char"0188\char"0316} \SUPPORTEDANCH{\char"0064\char"0316} \SUPPORTEDANCH{\char"018C\char"0316} \SUPPORTEDFALL{\char"0221\char"0316} \SUPPORTEDANCH{\char"0068\char"0316} \SUPPORTEDANCH{\char"2C68\char"0316} \SUPPORTEDANCH{\char"006B\char"0316} \SUPPORTEDANCH{\char"0199\char"0316} \SUPPORTEDANCH{\char"2C6A\char"0316} \SUPPORTEDANCH{\char"006C\char"0316} \SUPPORTEDFALL{\char"0234\char"0316} \SUPPORTEDANCH{\char"006E\char"0316} \SUPPORTEDANCH{\char"0235\char"0316} \SUPPORTEDANCH{\char"0070\char"0316} \SUPPORTEDFALL{\char"01A5\char"0316} \SUPPORTEDANCH{\char"0071\char"0316} \SUPPORTEDFALL{\char"02A0\char"0316} \SUPPORTEDANCH{\char"0074\char"0316} \SUPPORTEDFALL{\char"01AB\char"0316} \SUPPORTEDANCH{\char"01AD\char"0316} \SUPPORTEDANCH{\char"0236\char"0316} \SUPPORTEDANCH{\char"0075\char"0316} \MISSINGGLYPH{\char"AB52\char"0316} \SUPPORTEDANCH{\char"0079\char"0316} \SUPPORTEDFALL{\char"01B4\char"0316} \MISSINGGLYPH{\char"1EFF\char"0316}

\SUPPORTEDANCH{\char"0048\char"0317} \SUPPORTEDFALL{\char"2C67\char"0317} \SUPPORTEDANCH{\char"004B\char"0317} \SUPPORTEDANCH{\char"2C69\char"0317} \SUPPORTEDANCH{\char"004D\char"0317} \MISSINGGLYPH{\char"2C6E\char"0317} \SUPPORTEDANCH{\char"0050\char"0317} \SUPPORTEDANCH{\char"01A4\char"0317} \SUPPORTEDANCH{\char"0052\char"0317} \SUPPORTEDFALL{\char"2C64\char"0317} \SUPPORTEDANCH{\char"0054\char"0317} \SUPPORTEDFALL{\char"01AC\char"0317} \SUPPORTEDFALL{\char"01AE\char"0317} \SUPPORTEDANCH{\char"0059\char"0317} \SUPPORTEDFALL{\char"01B3\char"0317} \SUPPORTEDANCH{\char"0061\char"0317} \SUPPORTEDFALL{\char"1E9A\char"0317} \SUPPORTEDANCH{\char"0062\char"0317} \SUPPORTEDANCH{\char"0183\char"0317} \SUPPORTEDANCH{\char"0063\char"0317} \SUPPORTEDFALL{\char"0188\char"0317} \SUPPORTEDANCH{\char"0064\char"0317} \SUPPORTEDANCH{\char"018C\char"0317} \SUPPORTEDFALL{\char"0221\char"0317} \SUPPORTEDANCH{\char"0068\char"0317} \SUPPORTEDANCH{\char"2C68\char"0317} \SUPPORTEDANCH{\char"006B\char"0317} \SUPPORTEDANCH{\char"0199\char"0317} \SUPPORTEDANCH{\char"2C6A\char"0317} \SUPPORTEDANCH{\char"006C\char"0317} \SUPPORTEDFALL{\char"0234\char"0317} \SUPPORTEDANCH{\char"006E\char"0317} \SUPPORTEDANCH{\char"0235\char"0317} \SUPPORTEDANCH{\char"0070\char"0317} \SUPPORTEDFALL{\char"01A5\char"0317} \SUPPORTEDANCH{\char"0071\char"0317} \SUPPORTEDFALL{\char"02A0\char"0317} \SUPPORTEDANCH{\char"0074\char"0317} \SUPPORTEDFALL{\char"01AB\char"0317} \SUPPORTEDANCH{\char"01AD\char"0317} \SUPPORTEDANCH{\char"0236\char"0317} \SUPPORTEDANCH{\char"0075\char"0317} \MISSINGGLYPH{\char"AB52\char"0317} \SUPPORTEDANCH{\char"0079\char"0317} \SUPPORTEDFALL{\char"01B4\char"0317} \MISSINGGLYPH{\char"1EFF\char"0317}

\SUPPORTEDANCH{\char"0048\char"0318} \SUPPORTEDFALL{\char"2C67\char"0318} \SUPPORTEDANCH{\char"004B\char"0318} \SUPPORTEDANCH{\char"2C69\char"0318} \SUPPORTEDANCH{\char"004D\char"0318} \MISSINGGLYPH{\char"2C6E\char"0318} \SUPPORTEDANCH{\char"0050\char"0318} \SUPPORTEDANCH{\char"01A4\char"0318} \SUPPORTEDANCH{\char"0052\char"0318} \SUPPORTEDFALL{\char"2C64\char"0318} \SUPPORTEDANCH{\char"0054\char"0318} \SUPPORTEDFALL{\char"01AC\char"0318} \SUPPORTEDFALL{\char"01AE\char"0318} \SUPPORTEDANCH{\char"0059\char"0318} \SUPPORTEDFALL{\char"01B3\char"0318} \SUPPORTEDANCH{\char"0061\char"0318} \SUPPORTEDFALL{\char"1E9A\char"0318} \SUPPORTEDANCH{\char"0062\char"0318} \SUPPORTEDANCH{\char"0183\char"0318} \SUPPORTEDANCH{\char"0063\char"0318} \SUPPORTEDFALL{\char"0188\char"0318} \SUPPORTEDANCH{\char"0064\char"0318} \SUPPORTEDANCH{\char"018C\char"0318} \SUPPORTEDFALL{\char"0221\char"0318} \SUPPORTEDANCH{\char"0068\char"0318} \SUPPORTEDANCH{\char"2C68\char"0318} \SUPPORTEDANCH{\char"006B\char"0318} \SUPPORTEDANCH{\char"0199\char"0318} \SUPPORTEDANCH{\char"2C6A\char"0318} \SUPPORTEDANCH{\char"006C\char"0318} \SUPPORTEDFALL{\char"0234\char"0318} \SUPPORTEDANCH{\char"006E\char"0318} \SUPPORTEDANCH{\char"0235\char"0318} \SUPPORTEDANCH{\char"0070\char"0318} \SUPPORTEDFALL{\char"01A5\char"0318} \SUPPORTEDANCH{\char"0071\char"0318} \SUPPORTEDFALL{\char"02A0\char"0318} \SUPPORTEDANCH{\char"0074\char"0318} \SUPPORTEDFALL{\char"01AB\char"0318} \SUPPORTEDANCH{\char"01AD\char"0318} \SUPPORTEDANCH{\char"0236\char"0318} \SUPPORTEDANCH{\char"0075\char"0318} \MISSINGGLYPH{\char"AB52\char"0318} \SUPPORTEDANCH{\char"0079\char"0318} \SUPPORTEDFALL{\char"01B4\char"0318} \MISSINGGLYPH{\char"1EFF\char"0318}

\SUPPORTEDANCH{\char"0048\char"0319} \SUPPORTEDFALL{\char"2C67\char"0319} \SUPPORTEDANCH{\char"004B\char"0319} \SUPPORTEDANCH{\char"2C69\char"0319} \SUPPORTEDANCH{\char"004D\char"0319} \MISSINGGLYPH{\char"2C6E\char"0319} \SUPPORTEDANCH{\char"0050\char"0319} \SUPPORTEDANCH{\char"01A4\char"0319} \SUPPORTEDANCH{\char"0052\char"0319} \SUPPORTEDFALL{\char"2C64\char"0319} \SUPPORTEDANCH{\char"0054\char"0319} \SUPPORTEDFALL{\char"01AC\char"0319} \SUPPORTEDFALL{\char"01AE\char"0319} \SUPPORTEDANCH{\char"0059\char"0319} \SUPPORTEDFALL{\char"01B3\char"0319} \SUPPORTEDANCH{\char"0061\char"0319} \SUPPORTEDFALL{\char"1E9A\char"0319} \SUPPORTEDANCH{\char"0062\char"0319} \SUPPORTEDANCH{\char"0183\char"0319} \SUPPORTEDANCH{\char"0063\char"0319} \SUPPORTEDFALL{\char"0188\char"0319} \SUPPORTEDANCH{\char"0064\char"0319} \SUPPORTEDANCH{\char"018C\char"0319} \SUPPORTEDFALL{\char"0221\char"0319} \SUPPORTEDANCH{\char"0068\char"0319} \SUPPORTEDANCH{\char"2C68\char"0319} \SUPPORTEDANCH{\char"006B\char"0319} \SUPPORTEDANCH{\char"0199\char"0319} \SUPPORTEDANCH{\char"2C6A\char"0319} \SUPPORTEDANCH{\char"006C\char"0319} \SUPPORTEDFALL{\char"0234\char"0319} \SUPPORTEDANCH{\char"006E\char"0319} \SUPPORTEDANCH{\char"0235\char"0319} \SUPPORTEDANCH{\char"0070\char"0319} \SUPPORTEDFALL{\char"01A5\char"0319} \SUPPORTEDANCH{\char"0071\char"0319} \SUPPORTEDFALL{\char"02A0\char"0319} \SUPPORTEDANCH{\char"0074\char"0319} \SUPPORTEDFALL{\char"01AB\char"0319} \SUPPORTEDANCH{\char"01AD\char"0319} \SUPPORTEDANCH{\char"0236\char"0319} \SUPPORTEDANCH{\char"0075\char"0319} \MISSINGGLYPH{\char"AB52\char"0319} \SUPPORTEDANCH{\char"0079\char"0319} \SUPPORTEDFALL{\char"01B4\char"0319} \MISSINGGLYPH{\char"1EFF\char"0319}

\SUPPORTEDANCH{\char"0048\char"031C} \SUPPORTEDFALL{\char"2C67\char"031C} \SUPPORTEDANCH{\char"004B\char"031C} \SUPPORTEDANCH{\char"2C69\char"031C} \SUPPORTEDANCH{\char"004D\char"031C} \MISSINGGLYPH{\char"2C6E\char"031C} \SUPPORTEDANCH{\char"0050\char"031C} \SUPPORTEDANCH{\char"01A4\char"031C} \SUPPORTEDANCH{\char"0052\char"031C} \SUPPORTEDFALL{\char"2C64\char"031C} \SUPPORTEDANCH{\char"0054\char"031C} \SUPPORTEDFALL{\char"01AC\char"031C} \SUPPORTEDFALL{\char"01AE\char"031C} \SUPPORTEDANCH{\char"0059\char"031C} \SUPPORTEDFALL{\char"01B3\char"031C} \SUPPORTEDANCH{\char"0061\char"031C} \SUPPORTEDFALL{\char"1E9A\char"031C} \SUPPORTEDANCH{\char"0062\char"031C} \SUPPORTEDANCH{\char"0183\char"031C} \SUPPORTEDANCH{\char"0063\char"031C} \SUPPORTEDFALL{\char"0188\char"031C} \SUPPORTEDANCH{\char"0064\char"031C} \SUPPORTEDANCH{\char"018C\char"031C} \SUPPORTEDFALL{\char"0221\char"031C} \SUPPORTEDANCH{\char"0068\char"031C} \SUPPORTEDANCH{\char"2C68\char"031C} \SUPPORTEDANCH{\char"006B\char"031C} \SUPPORTEDANCH{\char"0199\char"031C} \SUPPORTEDANCH{\char"2C6A\char"031C} \SUPPORTEDANCH{\char"006C\char"031C} \SUPPORTEDFALL{\char"0234\char"031C} \SUPPORTEDANCH{\char"006E\char"031C} \SUPPORTEDANCH{\char"0235\char"031C} \SUPPORTEDANCH{\char"0070\char"031C} \SUPPORTEDFALL{\char"01A5\char"031C} \SUPPORTEDANCH{\char"0071\char"031C} \SUPPORTEDFALL{\char"02A0\char"031C} \SUPPORTEDANCH{\char"0074\char"031C} \SUPPORTEDFALL{\char"01AB\char"031C} \SUPPORTEDANCH{\char"01AD\char"031C} \SUPPORTEDANCH{\char"0236\char"031C} \SUPPORTEDANCH{\char"0075\char"031C} \MISSINGGLYPH{\char"AB52\char"031C} \SUPPORTEDANCH{\char"0079\char"031C} \SUPPORTEDFALL{\char"01B4\char"031C} \MISSINGGLYPH{\char"1EFF\char"031C}

\SUPPORTEDANCH{\char"0048\char"0320} \SUPPORTEDFALL{\char"2C67\char"0320} \SUPPORTEDANCH{\char"004B\char"0320} \SUPPORTEDANCH{\char"2C69\char"0320} \SUPPORTEDANCH{\char"004D\char"0320} \MISSINGGLYPH{\char"2C6E\char"0320} \SUPPORTEDANCH{\char"0050\char"0320} \SUPPORTEDANCH{\char"01A4\char"0320} \SUPPORTEDANCH{\char"0052\char"0320} \SUPPORTEDFALL{\char"2C64\char"0320} \SUPPORTEDANCH{\char"0054\char"0320} \SUPPORTEDFALL{\char"01AC\char"0320} \SUPPORTEDFALL{\char"01AE\char"0320} \SUPPORTEDANCH{\char"0059\char"0320} \SUPPORTEDFALL{\char"01B3\char"0320} \SUPPORTEDANCH{\char"0061\char"0320} \SUPPORTEDFALL{\char"1E9A\char"0320} \SUPPORTEDANCH{\char"0062\char"0320} \SUPPORTEDANCH{\char"0183\char"0320} \SUPPORTEDANCH{\char"0063\char"0320} \SUPPORTEDFALL{\char"0188\char"0320} \SUPPORTEDANCH{\char"0064\char"0320} \SUPPORTEDANCH{\char"018C\char"0320} \SUPPORTEDFALL{\char"0221\char"0320} \SUPPORTEDANCH{\char"0068\char"0320} \SUPPORTEDANCH{\char"2C68\char"0320} \SUPPORTEDANCH{\char"006B\char"0320} \SUPPORTEDANCH{\char"0199\char"0320} \SUPPORTEDANCH{\char"2C6A\char"0320} \SUPPORTEDANCH{\char"006C\char"0320} \SUPPORTEDFALL{\char"0234\char"0320} \SUPPORTEDANCH{\char"006E\char"0320} \SUPPORTEDANCH{\char"0235\char"0320} \SUPPORTEDANCH{\char"0070\char"0320} \SUPPORTEDFALL{\char"01A5\char"0320} \SUPPORTEDANCH{\char"0071\char"0320} \SUPPORTEDFALL{\char"02A0\char"0320} \SUPPORTEDANCH{\char"0074\char"0320} \SUPPORTEDFALL{\char"01AB\char"0320} \SUPPORTEDANCH{\char"01AD\char"0320} \SUPPORTEDANCH{\char"0236\char"0320} \SUPPORTEDANCH{\char"0075\char"0320} \MISSINGGLYPH{\char"AB52\char"0320} \SUPPORTEDANCH{\char"0079\char"0320} \SUPPORTEDFALL{\char"01B4\char"0320} \MISSINGGLYPH{\char"1EFF\char"0320}

\SUPPORTEDPREC{\char"0048\char"0323} \SUPPORTEDFALL{\char"2C67\char"0323} \SUPPORTEDPREC{\char"004B\char"0323} \SUPPORTEDANCH{\char"2C69\char"0323} \SUPPORTEDPREC{\char"004D\char"0323} \MISSINGGLYPH{\char"2C6E\char"0323} \SUPPORTEDANCH{\char"0050\char"0323} \SUPPORTEDANCH{\char"01A4\char"0323} \SUPPORTEDPREC{\char"0052\char"0323} \SUPPORTEDFALL{\char"2C64\char"0323} \SUPPORTEDPREC{\char"0054\char"0323} \SUPPORTEDFALL{\char"01AC\char"0323} \SUPPORTEDFALL{\char"01AE\char"0323} \SUPPORTEDPREC{\char"0059\char"0323} \SUPPORTEDFALL{\char"01B3\char"0323} \SUPPORTEDPREC{\char"0061\char"0323} \SUPPORTEDFALL{\char"1E9A\char"0323} \SUPPORTEDPREC{\char"0062\char"0323} \SUPPORTEDANCH{\char"0183\char"0323} \SUPPORTEDANCH{\char"0063\char"0323} \SUPPORTEDFALL{\char"0188\char"0323} \SUPPORTEDPREC{\char"0064\char"0323} \SUPPORTEDANCH{\char"018C\char"0323} \SUPPORTEDFALL{\char"0221\char"0323} \SUPPORTEDPREC{\char"0068\char"0323} \SUPPORTEDANCH{\char"2C68\char"0323} \SUPPORTEDPREC{\char"006B\char"0323} \SUPPORTEDANCH{\char"0199\char"0323} \SUPPORTEDANCH{\char"2C6A\char"0323} \SUPPORTEDPREC{\char"006C\char"0323} \SUPPORTEDFALL{\char"0234\char"0323} \SUPPORTEDPREC{\char"006E\char"0323} \SUPPORTEDANCH{\char"0235\char"0323} \SUPPORTEDANCH{\char"0070\char"0323} \SUPPORTEDFALL{\char"01A5\char"0323} \SUPPORTEDANCH{\char"0071\char"0323} \SUPPORTEDFALL{\char"02A0\char"0323} \SUPPORTEDPREC{\char"0074\char"0323} \SUPPORTEDFALL{\char"01AB\char"0323} \SUPPORTEDANCH{\char"01AD\char"0323} \SUPPORTEDANCH{\char"0236\char"0323} \SUPPORTEDPREC{\char"0075\char"0323} \MISSINGGLYPH{\char"AB52\char"0323} \SUPPORTEDPREC{\char"0079\char"0323} \SUPPORTEDFALL{\char"01B4\char"0323} \MISSINGGLYPH{\char"1EFF\char"0323}

\SUPPORTEDANCH{\char"0048\char"0324} \SUPPORTEDFALL{\char"2C67\char"0324} \SUPPORTEDANCH{\char"004B\char"0324} \SUPPORTEDANCH{\char"2C69\char"0324} \SUPPORTEDANCH{\char"004D\char"0324} \MISSINGGLYPH{\char"2C6E\char"0324} \SUPPORTEDANCH{\char"0050\char"0324} \SUPPORTEDANCH{\char"01A4\char"0324} \SUPPORTEDANCH{\char"0052\char"0324} \SUPPORTEDFALL{\char"2C64\char"0324} \SUPPORTEDANCH{\char"0054\char"0324} \SUPPORTEDFALL{\char"01AC\char"0324} \SUPPORTEDFALL{\char"01AE\char"0324} \SUPPORTEDANCH{\char"0059\char"0324} \SUPPORTEDFALL{\char"01B3\char"0324} \SUPPORTEDANCH{\char"0061\char"0324} \SUPPORTEDFALL{\char"1E9A\char"0324} \SUPPORTEDANCH{\char"0062\char"0324} \SUPPORTEDANCH{\char"0183\char"0324} \SUPPORTEDANCH{\char"0063\char"0324} \SUPPORTEDFALL{\char"0188\char"0324} \SUPPORTEDANCH{\char"0064\char"0324} \SUPPORTEDANCH{\char"018C\char"0324} \SUPPORTEDFALL{\char"0221\char"0324} \SUPPORTEDANCH{\char"0068\char"0324} \SUPPORTEDANCH{\char"2C68\char"0324} \SUPPORTEDANCH{\char"006B\char"0324} \SUPPORTEDANCH{\char"0199\char"0324} \SUPPORTEDANCH{\char"2C6A\char"0324} \SUPPORTEDANCH{\char"006C\char"0324} \SUPPORTEDFALL{\char"0234\char"0324} \SUPPORTEDANCH{\char"006E\char"0324} \SUPPORTEDANCH{\char"0235\char"0324} \SUPPORTEDANCH{\char"0070\char"0324} \SUPPORTEDFALL{\char"01A5\char"0324} \SUPPORTEDANCH{\char"0071\char"0324} \SUPPORTEDFALL{\char"02A0\char"0324} \SUPPORTEDANCH{\char"0074\char"0324} \SUPPORTEDFALL{\char"01AB\char"0324} \SUPPORTEDANCH{\char"01AD\char"0324} \SUPPORTEDANCH{\char"0236\char"0324} \SUPPORTEDPREC{\char"0075\char"0324} \MISSINGGLYPH{\char"AB52\char"0324} \SUPPORTEDANCH{\char"0079\char"0324} \SUPPORTEDFALL{\char"01B4\char"0324} \MISSINGGLYPH{\char"1EFF\char"0324}

\SUPPORTEDANCH{\char"0048\char"0325} \SUPPORTEDFALL{\char"2C67\char"0325} \SUPPORTEDANCH{\char"004B\char"0325} \SUPPORTEDANCH{\char"2C69\char"0325} \SUPPORTEDANCH{\char"004D\char"0325} \MISSINGGLYPH{\char"2C6E\char"0325} \SUPPORTEDANCH{\char"0050\char"0325} \SUPPORTEDANCH{\char"01A4\char"0325} \SUPPORTEDANCH{\char"0052\char"0325} \SUPPORTEDFALL{\char"2C64\char"0325} \SUPPORTEDANCH{\char"0054\char"0325} \SUPPORTEDFALL{\char"01AC\char"0325} \SUPPORTEDFALL{\char"01AE\char"0325} \SUPPORTEDANCH{\char"0059\char"0325} \SUPPORTEDFALL{\char"01B3\char"0325} \SUPPORTEDPREC{\char"0061\char"0325} \SUPPORTEDFALL{\char"1E9A\char"0325} \SUPPORTEDANCH{\char"0062\char"0325} \SUPPORTEDANCH{\char"0183\char"0325} \SUPPORTEDANCH{\char"0063\char"0325} \SUPPORTEDFALL{\char"0188\char"0325} \SUPPORTEDANCH{\char"0064\char"0325} \SUPPORTEDANCH{\char"018C\char"0325} \SUPPORTEDFALL{\char"0221\char"0325} \SUPPORTEDANCH{\char"0068\char"0325} \SUPPORTEDANCH{\char"2C68\char"0325} \SUPPORTEDANCH{\char"006B\char"0325} \SUPPORTEDANCH{\char"0199\char"0325} \SUPPORTEDANCH{\char"2C6A\char"0325} \SUPPORTEDANCH{\char"006C\char"0325} \SUPPORTEDFALL{\char"0234\char"0325} \SUPPORTEDANCH{\char"006E\char"0325} \SUPPORTEDANCH{\char"0235\char"0325} \SUPPORTEDANCH{\char"0070\char"0325} \SUPPORTEDFALL{\char"01A5\char"0325} \SUPPORTEDANCH{\char"0071\char"0325} \SUPPORTEDFALL{\char"02A0\char"0325} \SUPPORTEDANCH{\char"0074\char"0325} \SUPPORTEDFALL{\char"01AB\char"0325} \SUPPORTEDANCH{\char"01AD\char"0325} \SUPPORTEDANCH{\char"0236\char"0325} \SUPPORTEDANCH{\char"0075\char"0325} \MISSINGGLYPH{\char"AB52\char"0325} \SUPPORTEDANCH{\char"0079\char"0325} \SUPPORTEDFALL{\char"01B4\char"0325} \MISSINGGLYPH{\char"1EFF\char"0325}

\SUPPORTEDANCH{\char"0048\char"0326} \SUPPORTEDFALL{\char"2C67\char"0326} \SUPPORTEDANCH{\char"004B\char"0326} \SUPPORTEDANCH{\char"2C69\char"0326} \SUPPORTEDANCH{\char"004D\char"0326} \MISSINGGLYPH{\char"2C6E\char"0326} \SUPPORTEDANCH{\char"0050\char"0326} \SUPPORTEDANCH{\char"01A4\char"0326} \SUPPORTEDANCH{\char"0052\char"0326} \SUPPORTEDFALL{\char"2C64\char"0326} \SUPPORTEDPREC{\char"0054\char"0326} \SUPPORTEDFALL{\char"01AC\char"0326} \SUPPORTEDFALL{\char"01AE\char"0326} \SUPPORTEDANCH{\char"0059\char"0326} \SUPPORTEDFALL{\char"01B3\char"0326} \SUPPORTEDANCH{\char"0061\char"0326} \SUPPORTEDFALL{\char"1E9A\char"0326} \SUPPORTEDANCH{\char"0062\char"0326} \SUPPORTEDANCH{\char"0183\char"0326} \SUPPORTEDANCH{\char"0063\char"0326} \SUPPORTEDFALL{\char"0188\char"0326} \SUPPORTEDANCH{\char"0064\char"0326} \SUPPORTEDANCH{\char"018C\char"0326} \SUPPORTEDFALL{\char"0221\char"0326} \SUPPORTEDANCH{\char"0068\char"0326} \SUPPORTEDANCH{\char"2C68\char"0326} \SUPPORTEDANCH{\char"006B\char"0326} \SUPPORTEDANCH{\char"0199\char"0326} \SUPPORTEDANCH{\char"2C6A\char"0326} \SUPPORTEDANCH{\char"006C\char"0326} \SUPPORTEDFALL{\char"0234\char"0326} \SUPPORTEDANCH{\char"006E\char"0326} \SUPPORTEDANCH{\char"0235\char"0326} \SUPPORTEDANCH{\char"0070\char"0326} \SUPPORTEDFALL{\char"01A5\char"0326} \SUPPORTEDANCH{\char"0071\char"0326} \SUPPORTEDFALL{\char"02A0\char"0326} \SUPPORTEDPREC{\char"0074\char"0326} \SUPPORTEDFALL{\char"01AB\char"0326} \SUPPORTEDANCH{\char"01AD\char"0326} \SUPPORTEDANCH{\char"0236\char"0326} \SUPPORTEDANCH{\char"0075\char"0326} \MISSINGGLYPH{\char"AB52\char"0326} \SUPPORTEDANCH{\char"0079\char"0326} \SUPPORTEDFALL{\char"01B4\char"0326} \MISSINGGLYPH{\char"1EFF\char"0326}

\SUPPORTEDANCH{\char"0048\char"0329} \SUPPORTEDFALL{\char"2C67\char"0329} \SUPPORTEDANCH{\char"004B\char"0329} \SUPPORTEDANCH{\char"2C69\char"0329} \SUPPORTEDANCH{\char"004D\char"0329} \MISSINGGLYPH{\char"2C6E\char"0329} \SUPPORTEDANCH{\char"0050\char"0329} \SUPPORTEDANCH{\char"01A4\char"0329} \SUPPORTEDANCH{\char"0052\char"0329} \SUPPORTEDFALL{\char"2C64\char"0329} \SUPPORTEDANCH{\char"0054\char"0329} \SUPPORTEDFALL{\char"01AC\char"0329} \SUPPORTEDFALL{\char"01AE\char"0329} \SUPPORTEDANCH{\char"0059\char"0329} \SUPPORTEDFALL{\char"01B3\char"0329} \SUPPORTEDANCH{\char"0061\char"0329} \SUPPORTEDFALL{\char"1E9A\char"0329} \SUPPORTEDANCH{\char"0062\char"0329} \SUPPORTEDANCH{\char"0183\char"0329} \SUPPORTEDANCH{\char"0063\char"0329} \SUPPORTEDFALL{\char"0188\char"0329} \SUPPORTEDANCH{\char"0064\char"0329} \SUPPORTEDANCH{\char"018C\char"0329} \SUPPORTEDFALL{\char"0221\char"0329} \SUPPORTEDANCH{\char"0068\char"0329} \SUPPORTEDANCH{\char"2C68\char"0329} \SUPPORTEDANCH{\char"006B\char"0329} \SUPPORTEDANCH{\char"0199\char"0329} \SUPPORTEDANCH{\char"2C6A\char"0329} \SUPPORTEDANCH{\char"006C\char"0329} \SUPPORTEDFALL{\char"0234\char"0329} \SUPPORTEDANCH{\char"006E\char"0329} \SUPPORTEDANCH{\char"0235\char"0329} \SUPPORTEDANCH{\char"0070\char"0329} \SUPPORTEDFALL{\char"01A5\char"0329} \SUPPORTEDANCH{\char"0071\char"0329} \SUPPORTEDFALL{\char"02A0\char"0329} \SUPPORTEDANCH{\char"0074\char"0329} \SUPPORTEDFALL{\char"01AB\char"0329} \SUPPORTEDANCH{\char"01AD\char"0329} \SUPPORTEDANCH{\char"0236\char"0329} \SUPPORTEDANCH{\char"0075\char"0329} \MISSINGGLYPH{\char"AB52\char"0329} \SUPPORTEDANCH{\char"0079\char"0329} \SUPPORTEDFALL{\char"01B4\char"0329} \MISSINGGLYPH{\char"1EFF\char"0329}

\SUPPORTEDANCH{\char"0048\char"032A} \SUPPORTEDFALL{\char"2C67\char"032A} \SUPPORTEDANCH{\char"004B\char"032A} \SUPPORTEDANCH{\char"2C69\char"032A} \SUPPORTEDANCH{\char"004D\char"032A} \MISSINGGLYPH{\char"2C6E\char"032A} \SUPPORTEDANCH{\char"0050\char"032A} \SUPPORTEDANCH{\char"01A4\char"032A} \SUPPORTEDANCH{\char"0052\char"032A} \SUPPORTEDFALL{\char"2C64\char"032A} \SUPPORTEDANCH{\char"0054\char"032A} \SUPPORTEDFALL{\char"01AC\char"032A} \SUPPORTEDFALL{\char"01AE\char"032A} \SUPPORTEDANCH{\char"0059\char"032A} \SUPPORTEDFALL{\char"01B3\char"032A} \SUPPORTEDANCH{\char"0061\char"032A} \SUPPORTEDFALL{\char"1E9A\char"032A} \SUPPORTEDANCH{\char"0062\char"032A} \SUPPORTEDANCH{\char"0183\char"032A} \SUPPORTEDANCH{\char"0063\char"032A} \SUPPORTEDFALL{\char"0188\char"032A} \SUPPORTEDANCH{\char"0064\char"032A} \SUPPORTEDANCH{\char"018C\char"032A} \SUPPORTEDFALL{\char"0221\char"032A} \SUPPORTEDANCH{\char"0068\char"032A} \SUPPORTEDANCH{\char"2C68\char"032A} \SUPPORTEDANCH{\char"006B\char"032A} \SUPPORTEDANCH{\char"0199\char"032A} \SUPPORTEDANCH{\char"2C6A\char"032A} \SUPPORTEDANCH{\char"006C\char"032A} \SUPPORTEDFALL{\char"0234\char"032A} \SUPPORTEDANCH{\char"006E\char"032A} \SUPPORTEDANCH{\char"0235\char"032A} \SUPPORTEDANCH{\char"0070\char"032A} \SUPPORTEDFALL{\char"01A5\char"032A} \SUPPORTEDANCH{\char"0071\char"032A} \SUPPORTEDFALL{\char"02A0\char"032A} \SUPPORTEDANCH{\char"0074\char"032A} \SUPPORTEDFALL{\char"01AB\char"032A} \SUPPORTEDANCH{\char"01AD\char"032A} \SUPPORTEDANCH{\char"0236\char"032A} \SUPPORTEDANCH{\char"0075\char"032A} \MISSINGGLYPH{\char"AB52\char"032A} \SUPPORTEDANCH{\char"0079\char"032A} \SUPPORTEDFALL{\char"01B4\char"032A} \MISSINGGLYPH{\char"1EFF\char"032A}

\SUPPORTEDANCH{\char"0048\char"032B} \SUPPORTEDFALL{\char"2C67\char"032B} \SUPPORTEDANCH{\char"004B\char"032B} \SUPPORTEDANCH{\char"2C69\char"032B} \SUPPORTEDANCH{\char"004D\char"032B} \MISSINGGLYPH{\char"2C6E\char"032B} \SUPPORTEDANCH{\char"0050\char"032B} \SUPPORTEDANCH{\char"01A4\char"032B} \SUPPORTEDANCH{\char"0052\char"032B} \SUPPORTEDFALL{\char"2C64\char"032B} \SUPPORTEDANCH{\char"0054\char"032B} \SUPPORTEDFALL{\char"01AC\char"032B} \SUPPORTEDFALL{\char"01AE\char"032B} \SUPPORTEDANCH{\char"0059\char"032B} \SUPPORTEDFALL{\char"01B3\char"032B} \SUPPORTEDANCH{\char"0061\char"032B} \SUPPORTEDFALL{\char"1E9A\char"032B} \SUPPORTEDANCH{\char"0062\char"032B} \SUPPORTEDANCH{\char"0183\char"032B} \SUPPORTEDANCH{\char"0063\char"032B} \SUPPORTEDFALL{\char"0188\char"032B} \SUPPORTEDANCH{\char"0064\char"032B} \SUPPORTEDANCH{\char"018C\char"032B} \SUPPORTEDFALL{\char"0221\char"032B} \SUPPORTEDANCH{\char"0068\char"032B} \SUPPORTEDANCH{\char"2C68\char"032B} \SUPPORTEDANCH{\char"006B\char"032B} \SUPPORTEDANCH{\char"0199\char"032B} \SUPPORTEDANCH{\char"2C6A\char"032B} \SUPPORTEDANCH{\char"006C\char"032B} \SUPPORTEDFALL{\char"0234\char"032B} \SUPPORTEDANCH{\char"006E\char"032B} \SUPPORTEDANCH{\char"0235\char"032B} \SUPPORTEDANCH{\char"0070\char"032B} \SUPPORTEDFALL{\char"01A5\char"032B} \SUPPORTEDANCH{\char"0071\char"032B} \SUPPORTEDFALL{\char"02A0\char"032B} \SUPPORTEDANCH{\char"0074\char"032B} \SUPPORTEDFALL{\char"01AB\char"032B} \SUPPORTEDANCH{\char"01AD\char"032B} \SUPPORTEDANCH{\char"0236\char"032B} \SUPPORTEDANCH{\char"0075\char"032B} \MISSINGGLYPH{\char"AB52\char"032B} \SUPPORTEDANCH{\char"0079\char"032B} \SUPPORTEDFALL{\char"01B4\char"032B} \MISSINGGLYPH{\char"1EFF\char"032B}

\SUPPORTEDANCH{\char"0048\char"032C} \SUPPORTEDFALL{\char"2C67\char"032C} \SUPPORTEDANCH{\char"004B\char"032C} \SUPPORTEDANCH{\char"2C69\char"032C} \SUPPORTEDANCH{\char"004D\char"032C} \MISSINGGLYPH{\char"2C6E\char"032C} \SUPPORTEDANCH{\char"0050\char"032C} \SUPPORTEDANCH{\char"01A4\char"032C} \SUPPORTEDANCH{\char"0052\char"032C} \SUPPORTEDFALL{\char"2C64\char"032C} \SUPPORTEDANCH{\char"0054\char"032C} \SUPPORTEDFALL{\char"01AC\char"032C} \SUPPORTEDFALL{\char"01AE\char"032C} \SUPPORTEDANCH{\char"0059\char"032C} \SUPPORTEDFALL{\char"01B3\char"032C} \SUPPORTEDANCH{\char"0061\char"032C} \SUPPORTEDFALL{\char"1E9A\char"032C} \SUPPORTEDANCH{\char"0062\char"032C} \SUPPORTEDANCH{\char"0183\char"032C} \SUPPORTEDANCH{\char"0063\char"032C} \SUPPORTEDFALL{\char"0188\char"032C} \SUPPORTEDANCH{\char"0064\char"032C} \SUPPORTEDANCH{\char"018C\char"032C} \SUPPORTEDFALL{\char"0221\char"032C} \SUPPORTEDANCH{\char"0068\char"032C} \SUPPORTEDANCH{\char"2C68\char"032C} \SUPPORTEDANCH{\char"006B\char"032C} \SUPPORTEDANCH{\char"0199\char"032C} \SUPPORTEDANCH{\char"2C6A\char"032C} \SUPPORTEDANCH{\char"006C\char"032C} \SUPPORTEDFALL{\char"0234\char"032C} \SUPPORTEDANCH{\char"006E\char"032C} \SUPPORTEDANCH{\char"0235\char"032C} \SUPPORTEDANCH{\char"0070\char"032C} \SUPPORTEDFALL{\char"01A5\char"032C} \SUPPORTEDANCH{\char"0071\char"032C} \SUPPORTEDFALL{\char"02A0\char"032C} \SUPPORTEDANCH{\char"0074\char"032C} \SUPPORTEDFALL{\char"01AB\char"032C} \SUPPORTEDANCH{\char"01AD\char"032C} \SUPPORTEDANCH{\char"0236\char"032C} \SUPPORTEDANCH{\char"0075\char"032C} \MISSINGGLYPH{\char"AB52\char"032C} \SUPPORTEDANCH{\char"0079\char"032C} \SUPPORTEDFALL{\char"01B4\char"032C} \MISSINGGLYPH{\char"1EFF\char"032C}

\SUPPORTEDANCH{\char"0048\char"032D} \SUPPORTEDFALL{\char"2C67\char"032D} \SUPPORTEDANCH{\char"004B\char"032D} \SUPPORTEDANCH{\char"2C69\char"032D} \SUPPORTEDANCH{\char"004D\char"032D} \MISSINGGLYPH{\char"2C6E\char"032D} \SUPPORTEDANCH{\char"0050\char"032D} \SUPPORTEDANCH{\char"01A4\char"032D} \SUPPORTEDANCH{\char"0052\char"032D} \SUPPORTEDFALL{\char"2C64\char"032D} \SUPPORTEDPREC{\char"0054\char"032D} \SUPPORTEDFALL{\char"01AC\char"032D} \SUPPORTEDFALL{\char"01AE\char"032D} \SUPPORTEDANCH{\char"0059\char"032D} \SUPPORTEDFALL{\char"01B3\char"032D} \SUPPORTEDANCH{\char"0061\char"032D} \SUPPORTEDFALL{\char"1E9A\char"032D} \SUPPORTEDANCH{\char"0062\char"032D} \SUPPORTEDANCH{\char"0183\char"032D} \SUPPORTEDANCH{\char"0063\char"032D} \SUPPORTEDFALL{\char"0188\char"032D} \SUPPORTEDPREC{\char"0064\char"032D} \SUPPORTEDANCH{\char"018C\char"032D} \SUPPORTEDFALL{\char"0221\char"032D} \SUPPORTEDANCH{\char"0068\char"032D} \SUPPORTEDANCH{\char"2C68\char"032D} \SUPPORTEDANCH{\char"006B\char"032D} \SUPPORTEDANCH{\char"0199\char"032D} \SUPPORTEDANCH{\char"2C6A\char"032D} \SUPPORTEDPREC{\char"006C\char"032D} \SUPPORTEDFALL{\char"0234\char"032D} \SUPPORTEDPREC{\char"006E\char"032D} \SUPPORTEDANCH{\char"0235\char"032D} \SUPPORTEDANCH{\char"0070\char"032D} \SUPPORTEDFALL{\char"01A5\char"032D} \SUPPORTEDANCH{\char"0071\char"032D} \SUPPORTEDFALL{\char"02A0\char"032D} \SUPPORTEDPREC{\char"0074\char"032D} \SUPPORTEDFALL{\char"01AB\char"032D} \SUPPORTEDANCH{\char"01AD\char"032D} \SUPPORTEDANCH{\char"0236\char"032D} \SUPPORTEDPREC{\char"0075\char"032D} \MISSINGGLYPH{\char"AB52\char"032D} \SUPPORTEDANCH{\char"0079\char"032D} \SUPPORTEDFALL{\char"01B4\char"032D} \MISSINGGLYPH{\char"1EFF\char"032D}

\SUPPORTEDPREC{\char"0048\char"032E} \SUPPORTEDFALL{\char"2C67\char"032E} \SUPPORTEDANCH{\char"004B\char"032E} \SUPPORTEDANCH{\char"2C69\char"032E} \SUPPORTEDANCH{\char"004D\char"032E} \MISSINGGLYPH{\char"2C6E\char"032E} \SUPPORTEDANCH{\char"0050\char"032E} \SUPPORTEDANCH{\char"01A4\char"032E} \SUPPORTEDANCH{\char"0052\char"032E} \SUPPORTEDFALL{\char"2C64\char"032E} \SUPPORTEDANCH{\char"0054\char"032E} \SUPPORTEDFALL{\char"01AC\char"032E} \SUPPORTEDFALL{\char"01AE\char"032E} \SUPPORTEDANCH{\char"0059\char"032E} \SUPPORTEDFALL{\char"01B3\char"032E} \SUPPORTEDANCH{\char"0061\char"032E} \SUPPORTEDFALL{\char"1E9A\char"032E} \SUPPORTEDANCH{\char"0062\char"032E} \SUPPORTEDANCH{\char"0183\char"032E} \SUPPORTEDANCH{\char"0063\char"032E} \SUPPORTEDFALL{\char"0188\char"032E} \SUPPORTEDANCH{\char"0064\char"032E} \SUPPORTEDANCH{\char"018C\char"032E} \SUPPORTEDFALL{\char"0221\char"032E} \SUPPORTEDPREC{\char"0068\char"032E} \SUPPORTEDANCH{\char"2C68\char"032E} \SUPPORTEDANCH{\char"006B\char"032E} \SUPPORTEDANCH{\char"0199\char"032E} \SUPPORTEDANCH{\char"2C6A\char"032E} \SUPPORTEDANCH{\char"006C\char"032E} \SUPPORTEDFALL{\char"0234\char"032E} \SUPPORTEDANCH{\char"006E\char"032E} \SUPPORTEDANCH{\char"0235\char"032E} \SUPPORTEDANCH{\char"0070\char"032E} \SUPPORTEDFALL{\char"01A5\char"032E} \SUPPORTEDANCH{\char"0071\char"032E} \SUPPORTEDFALL{\char"02A0\char"032E} \SUPPORTEDANCH{\char"0074\char"032E} \SUPPORTEDFALL{\char"01AB\char"032E} \SUPPORTEDANCH{\char"01AD\char"032E} \SUPPORTEDANCH{\char"0236\char"032E} \SUPPORTEDANCH{\char"0075\char"032E} \MISSINGGLYPH{\char"AB52\char"032E} \SUPPORTEDANCH{\char"0079\char"032E} \SUPPORTEDFALL{\char"01B4\char"032E} \MISSINGGLYPH{\char"1EFF\char"032E}

\SUPPORTEDANCH{\char"0048\char"032F} \SUPPORTEDFALL{\char"2C67\char"032F} \SUPPORTEDANCH{\char"004B\char"032F} \SUPPORTEDANCH{\char"2C69\char"032F} \SUPPORTEDANCH{\char"004D\char"032F} \MISSINGGLYPH{\char"2C6E\char"032F} \SUPPORTEDANCH{\char"0050\char"032F} \SUPPORTEDANCH{\char"01A4\char"032F} \SUPPORTEDANCH{\char"0052\char"032F} \SUPPORTEDFALL{\char"2C64\char"032F} \SUPPORTEDANCH{\char"0054\char"032F} \SUPPORTEDFALL{\char"01AC\char"032F} \SUPPORTEDFALL{\char"01AE\char"032F} \SUPPORTEDANCH{\char"0059\char"032F} \SUPPORTEDFALL{\char"01B3\char"032F} \SUPPORTEDANCH{\char"0061\char"032F} \SUPPORTEDFALL{\char"1E9A\char"032F} \SUPPORTEDANCH{\char"0062\char"032F} \SUPPORTEDANCH{\char"0183\char"032F} \SUPPORTEDANCH{\char"0063\char"032F} \SUPPORTEDFALL{\char"0188\char"032F} \SUPPORTEDANCH{\char"0064\char"032F} \SUPPORTEDANCH{\char"018C\char"032F} \SUPPORTEDFALL{\char"0221\char"032F} \SUPPORTEDANCH{\char"0068\char"032F} \SUPPORTEDANCH{\char"2C68\char"032F} \SUPPORTEDANCH{\char"006B\char"032F} \SUPPORTEDANCH{\char"0199\char"032F} \SUPPORTEDANCH{\char"2C6A\char"032F} \SUPPORTEDANCH{\char"006C\char"032F} \SUPPORTEDFALL{\char"0234\char"032F} \SUPPORTEDANCH{\char"006E\char"032F} \SUPPORTEDANCH{\char"0235\char"032F} \SUPPORTEDANCH{\char"0070\char"032F} \SUPPORTEDFALL{\char"01A5\char"032F} \SUPPORTEDANCH{\char"0071\char"032F} \SUPPORTEDFALL{\char"02A0\char"032F} \SUPPORTEDANCH{\char"0074\char"032F} \SUPPORTEDFALL{\char"01AB\char"032F} \SUPPORTEDANCH{\char"01AD\char"032F} \SUPPORTEDANCH{\char"0236\char"032F} \SUPPORTEDANCH{\char"0075\char"032F} \MISSINGGLYPH{\char"AB52\char"032F} \SUPPORTEDANCH{\char"0079\char"032F} \SUPPORTEDFALL{\char"01B4\char"032F} \MISSINGGLYPH{\char"1EFF\char"032F}

\SUPPORTEDANCH{\char"0048\char"031D} \SUPPORTEDFALL{\char"2C67\char"031D} \SUPPORTEDANCH{\char"004B\char"031D} \SUPPORTEDANCH{\char"2C69\char"031D} \SUPPORTEDANCH{\char"004D\char"031D} \MISSINGGLYPH{\char"2C6E\char"031D} \SUPPORTEDANCH{\char"0050\char"031D} \SUPPORTEDANCH{\char"01A4\char"031D} \SUPPORTEDANCH{\char"0052\char"031D} \SUPPORTEDFALL{\char"2C64\char"031D} \SUPPORTEDANCH{\char"0054\char"031D} \SUPPORTEDFALL{\char"01AC\char"031D} \SUPPORTEDFALL{\char"01AE\char"031D} \SUPPORTEDANCH{\char"0059\char"031D} \SUPPORTEDFALL{\char"01B3\char"031D} \SUPPORTEDANCH{\char"0061\char"031D} \SUPPORTEDFALL{\char"1E9A\char"031D} \SUPPORTEDANCH{\char"0062\char"031D} \SUPPORTEDANCH{\char"0183\char"031D} \SUPPORTEDANCH{\char"0063\char"031D} \SUPPORTEDFALL{\char"0188\char"031D} \SUPPORTEDANCH{\char"0064\char"031D} \SUPPORTEDANCH{\char"018C\char"031D} \SUPPORTEDFALL{\char"0221\char"031D} \SUPPORTEDANCH{\char"0068\char"031D} \SUPPORTEDANCH{\char"2C68\char"031D} \SUPPORTEDANCH{\char"006B\char"031D} \SUPPORTEDANCH{\char"0199\char"031D} \SUPPORTEDANCH{\char"2C6A\char"031D} \SUPPORTEDANCH{\char"006C\char"031D} \SUPPORTEDFALL{\char"0234\char"031D} \SUPPORTEDANCH{\char"006E\char"031D} \SUPPORTEDANCH{\char"0235\char"031D} \SUPPORTEDANCH{\char"0070\char"031D} \SUPPORTEDFALL{\char"01A5\char"031D} \SUPPORTEDANCH{\char"0071\char"031D} \SUPPORTEDFALL{\char"02A0\char"031D} \SUPPORTEDANCH{\char"0074\char"031D} \SUPPORTEDFALL{\char"01AB\char"031D} \SUPPORTEDANCH{\char"01AD\char"031D} \SUPPORTEDANCH{\char"0236\char"031D} \SUPPORTEDANCH{\char"0075\char"031D} \MISSINGGLYPH{\char"AB52\char"031D} \SUPPORTEDANCH{\char"0079\char"031D} \SUPPORTEDFALL{\char"01B4\char"031D} \MISSINGGLYPH{\char"1EFF\char"031D}

\SUPPORTEDANCH{\char"0048\char"031E} \SUPPORTEDFALL{\char"2C67\char"031E} \SUPPORTEDANCH{\char"004B\char"031E} \SUPPORTEDANCH{\char"2C69\char"031E} \SUPPORTEDANCH{\char"004D\char"031E} \MISSINGGLYPH{\char"2C6E\char"031E} \SUPPORTEDANCH{\char"0050\char"031E} \SUPPORTEDANCH{\char"01A4\char"031E} \SUPPORTEDANCH{\char"0052\char"031E} \SUPPORTEDFALL{\char"2C64\char"031E} \SUPPORTEDANCH{\char"0054\char"031E} \SUPPORTEDFALL{\char"01AC\char"031E} \SUPPORTEDFALL{\char"01AE\char"031E} \SUPPORTEDANCH{\char"0059\char"031E} \SUPPORTEDFALL{\char"01B3\char"031E} \SUPPORTEDANCH{\char"0061\char"031E} \SUPPORTEDFALL{\char"1E9A\char"031E} \SUPPORTEDANCH{\char"0062\char"031E} \SUPPORTEDANCH{\char"0183\char"031E} \SUPPORTEDANCH{\char"0063\char"031E} \SUPPORTEDFALL{\char"0188\char"031E} \SUPPORTEDANCH{\char"0064\char"031E} \SUPPORTEDANCH{\char"018C\char"031E} \SUPPORTEDFALL{\char"0221\char"031E} \SUPPORTEDANCH{\char"0068\char"031E} \SUPPORTEDANCH{\char"2C68\char"031E} \SUPPORTEDANCH{\char"006B\char"031E} \SUPPORTEDANCH{\char"0199\char"031E} \SUPPORTEDANCH{\char"2C6A\char"031E} \SUPPORTEDANCH{\char"006C\char"031E} \SUPPORTEDFALL{\char"0234\char"031E} \SUPPORTEDANCH{\char"006E\char"031E} \SUPPORTEDANCH{\char"0235\char"031E} \SUPPORTEDANCH{\char"0070\char"031E} \SUPPORTEDFALL{\char"01A5\char"031E} \SUPPORTEDANCH{\char"0071\char"031E} \SUPPORTEDFALL{\char"02A0\char"031E} \SUPPORTEDANCH{\char"0074\char"031E} \SUPPORTEDFALL{\char"01AB\char"031E} \SUPPORTEDANCH{\char"01AD\char"031E} \SUPPORTEDANCH{\char"0236\char"031E} \SUPPORTEDANCH{\char"0075\char"031E} \MISSINGGLYPH{\char"AB52\char"031E} \SUPPORTEDANCH{\char"0079\char"031E} \SUPPORTEDFALL{\char"01B4\char"031E} \MISSINGGLYPH{\char"1EFF\char"031E}

\SUPPORTEDANCH{\char"0048\char"031F} \SUPPORTEDFALL{\char"2C67\char"031F} \SUPPORTEDANCH{\char"004B\char"031F} \SUPPORTEDANCH{\char"2C69\char"031F} \SUPPORTEDANCH{\char"004D\char"031F} \MISSINGGLYPH{\char"2C6E\char"031F} \SUPPORTEDANCH{\char"0050\char"031F} \SUPPORTEDANCH{\char"01A4\char"031F} \SUPPORTEDANCH{\char"0052\char"031F} \SUPPORTEDFALL{\char"2C64\char"031F} \SUPPORTEDANCH{\char"0054\char"031F} \SUPPORTEDFALL{\char"01AC\char"031F} \SUPPORTEDFALL{\char"01AE\char"031F} \SUPPORTEDANCH{\char"0059\char"031F} \SUPPORTEDFALL{\char"01B3\char"031F} \SUPPORTEDANCH{\char"0061\char"031F} \SUPPORTEDFALL{\char"1E9A\char"031F} \SUPPORTEDANCH{\char"0062\char"031F} \SUPPORTEDANCH{\char"0183\char"031F} \SUPPORTEDANCH{\char"0063\char"031F} \SUPPORTEDFALL{\char"0188\char"031F} \SUPPORTEDANCH{\char"0064\char"031F} \SUPPORTEDANCH{\char"018C\char"031F} \SUPPORTEDFALL{\char"0221\char"031F} \SUPPORTEDANCH{\char"0068\char"031F} \SUPPORTEDANCH{\char"2C68\char"031F} \SUPPORTEDANCH{\char"006B\char"031F} \SUPPORTEDANCH{\char"0199\char"031F} \SUPPORTEDANCH{\char"2C6A\char"031F} \SUPPORTEDANCH{\char"006C\char"031F} \SUPPORTEDFALL{\char"0234\char"031F} \SUPPORTEDANCH{\char"006E\char"031F} \SUPPORTEDANCH{\char"0235\char"031F} \SUPPORTEDANCH{\char"0070\char"031F} \SUPPORTEDFALL{\char"01A5\char"031F} \SUPPORTEDANCH{\char"0071\char"031F} \SUPPORTEDFALL{\char"02A0\char"031F} \SUPPORTEDANCH{\char"0074\char"031F} \SUPPORTEDFALL{\char"01AB\char"031F} \SUPPORTEDANCH{\char"01AD\char"031F} \SUPPORTEDANCH{\char"0236\char"031F} \SUPPORTEDANCH{\char"0075\char"031F} \MISSINGGLYPH{\char"AB52\char"031F} \SUPPORTEDANCH{\char"0079\char"031F} \SUPPORTEDFALL{\char"01B4\char"031F} \MISSINGGLYPH{\char"1EFF\char"031F}

\SUPPORTEDANCH{\char"0048\char"0330} \SUPPORTEDFALL{\char"2C67\char"0330} \SUPPORTEDANCH{\char"004B\char"0330} \SUPPORTEDANCH{\char"2C69\char"0330} \SUPPORTEDANCH{\char"004D\char"0330} \MISSINGGLYPH{\char"2C6E\char"0330} \SUPPORTEDANCH{\char"0050\char"0330} \SUPPORTEDANCH{\char"01A4\char"0330} \SUPPORTEDANCH{\char"0052\char"0330} \SUPPORTEDFALL{\char"2C64\char"0330} \SUPPORTEDANCH{\char"0054\char"0330} \SUPPORTEDFALL{\char"01AC\char"0330} \SUPPORTEDFALL{\char"01AE\char"0330} \SUPPORTEDANCH{\char"0059\char"0330} \SUPPORTEDFALL{\char"01B3\char"0330} \SUPPORTEDANCH{\char"0061\char"0330} \SUPPORTEDFALL{\char"1E9A\char"0330} \SUPPORTEDANCH{\char"0062\char"0330} \SUPPORTEDANCH{\char"0183\char"0330} \SUPPORTEDANCH{\char"0063\char"0330} \SUPPORTEDFALL{\char"0188\char"0330} \SUPPORTEDANCH{\char"0064\char"0330} \SUPPORTEDANCH{\char"018C\char"0330} \SUPPORTEDFALL{\char"0221\char"0330} \SUPPORTEDANCH{\char"0068\char"0330} \SUPPORTEDANCH{\char"2C68\char"0330} \SUPPORTEDANCH{\char"006B\char"0330} \SUPPORTEDANCH{\char"0199\char"0330} \SUPPORTEDANCH{\char"2C6A\char"0330} \SUPPORTEDANCH{\char"006C\char"0330} \SUPPORTEDFALL{\char"0234\char"0330} \SUPPORTEDANCH{\char"006E\char"0330} \SUPPORTEDANCH{\char"0235\char"0330} \SUPPORTEDANCH{\char"0070\char"0330} \SUPPORTEDFALL{\char"01A5\char"0330} \SUPPORTEDANCH{\char"0071\char"0330} \SUPPORTEDFALL{\char"02A0\char"0330} \SUPPORTEDANCH{\char"0074\char"0330} \SUPPORTEDFALL{\char"01AB\char"0330} \SUPPORTEDANCH{\char"01AD\char"0330} \SUPPORTEDANCH{\char"0236\char"0330} \SUPPORTEDPREC{\char"0075\char"0330} \MISSINGGLYPH{\char"AB52\char"0330} \SUPPORTEDANCH{\char"0079\char"0330} \SUPPORTEDFALL{\char"01B4\char"0330} \MISSINGGLYPH{\char"1EFF\char"0330}

\SUPPORTEDANCH{\char"0048\char"0331} \SUPPORTEDFALL{\char"2C67\char"0331} \SUPPORTEDPREC{\char"004B\char"0331} \SUPPORTEDANCH{\char"2C69\char"0331} \SUPPORTEDANCH{\char"004D\char"0331} \MISSINGGLYPH{\char"2C6E\char"0331} \SUPPORTEDANCH{\char"0050\char"0331} \SUPPORTEDANCH{\char"01A4\char"0331} \SUPPORTEDPREC{\char"0052\char"0331} \SUPPORTEDFALL{\char"2C64\char"0331} \SUPPORTEDPREC{\char"0054\char"0331} \SUPPORTEDFALL{\char"01AC\char"0331} \SUPPORTEDFALL{\char"01AE\char"0331} \SUPPORTEDANCH{\char"0059\char"0331} \SUPPORTEDFALL{\char"01B3\char"0331} \SUPPORTEDANCH{\char"0061\char"0331} \SUPPORTEDFALL{\char"1E9A\char"0331} \SUPPORTEDPREC{\char"0062\char"0331} \SUPPORTEDANCH{\char"0183\char"0331} \SUPPORTEDANCH{\char"0063\char"0331} \SUPPORTEDFALL{\char"0188\char"0331} \SUPPORTEDPREC{\char"0064\char"0331} \SUPPORTEDANCH{\char"018C\char"0331} \SUPPORTEDFALL{\char"0221\char"0331} \SUPPORTEDPREC{\char"0068\char"0331} \SUPPORTEDANCH{\char"2C68\char"0331} \SUPPORTEDPREC{\char"006B\char"0331} \SUPPORTEDANCH{\char"0199\char"0331} \SUPPORTEDANCH{\char"2C6A\char"0331} \SUPPORTEDPREC{\char"006C\char"0331} \SUPPORTEDFALL{\char"0234\char"0331} \SUPPORTEDPREC{\char"006E\char"0331} \SUPPORTEDANCH{\char"0235\char"0331} \SUPPORTEDANCH{\char"0070\char"0331} \SUPPORTEDFALL{\char"01A5\char"0331} \SUPPORTEDANCH{\char"0071\char"0331} \SUPPORTEDFALL{\char"02A0\char"0331} \SUPPORTEDPREC{\char"0074\char"0331} \SUPPORTEDFALL{\char"01AB\char"0331} \SUPPORTEDANCH{\char"01AD\char"0331} \SUPPORTEDANCH{\char"0236\char"0331} \SUPPORTEDANCH{\char"0075\char"0331} \MISSINGGLYPH{\char"AB52\char"0331} \SUPPORTEDANCH{\char"0079\char"0331} \SUPPORTEDFALL{\char"01B4\char"0331} \MISSINGGLYPH{\char"1EFF\char"0331}

\MISSINGGLYPH{\char"0048} \MISSINGGLYPH{\char"2C67} \MISSINGGLYPH{\char"004B} \MISSINGGLYPH{\char"2C69} \MISSINGGLYPH{\char"004D} \MISSINGGLYPH{\char"2C6E} \MISSINGGLYPH{\char"0050} \MISSINGGLYPH{\char"01A4} \MISSINGGLYPH{\char"0052} \MISSINGGLYPH{\char"2C64} \MISSINGGLYPH{\char"0054} \MISSINGGLYPH{\char"01AC} \MISSINGGLYPH{\char"01AE} \MISSINGGLYPH{\char"0059} \MISSINGGLYPH{\char"01B3} \MISSINGGLYPH{\char"0061} \MISSINGGLYPH{\char"1E9A} \MISSINGGLYPH{\char"0062} \MISSINGGLYPH{\char"0183} \MISSINGGLYPH{\char"0063} \MISSINGGLYPH{\char"0188} \MISSINGGLYPH{\char"0064} \MISSINGGLYPH{\char"018C} \MISSINGGLYPH{\char"0221} \MISSINGGLYPH{\char"0068} \MISSINGGLYPH{\char"2C68} \MISSINGGLYPH{\char"006B} \MISSINGGLYPH{\char"0199} \MISSINGGLYPH{\char"2C6A} \MISSINGGLYPH{\char"006C} \MISSINGGLYPH{\char"0234} \MISSINGGLYPH{\char"006E} \MISSINGGLYPH{\char"0235} \MISSINGGLYPH{\char"0070} \MISSINGGLYPH{\char"01A5} \MISSINGGLYPH{\char"0071} \MISSINGGLYPH{\char"02A0} \MISSINGGLYPH{\char"0074} \MISSINGGLYPH{\char"01AB} \MISSINGGLYPH{\char"01AD} \MISSINGGLYPH{\char"0236} \MISSINGGLYPH{\char"0075} \MISSINGGLYPH{\char"AB52} \MISSINGGLYPH{\char"0079} \MISSINGGLYPH{\char"01B4} \MISSINGGLYPH{\char"1EFF}

\MISSINGGLYPH{\char"0048} \MISSINGGLYPH{\char"2C67} \MISSINGGLYPH{\char"004B} \MISSINGGLYPH{\char"2C69} \MISSINGGLYPH{\char"004D} \MISSINGGLYPH{\char"2C6E} \MISSINGGLYPH{\char"0050} \MISSINGGLYPH{\char"01A4} \MISSINGGLYPH{\char"0052} \MISSINGGLYPH{\char"2C64} \MISSINGGLYPH{\char"0054} \MISSINGGLYPH{\char"01AC} \MISSINGGLYPH{\char"01AE} \MISSINGGLYPH{\char"0059} \MISSINGGLYPH{\char"01B3} \MISSINGGLYPH{\char"0061} \MISSINGGLYPH{\char"1E9A} \MISSINGGLYPH{\char"0062} \MISSINGGLYPH{\char"0183} \MISSINGGLYPH{\char"0063} \MISSINGGLYPH{\char"0188} \MISSINGGLYPH{\char"0064} \MISSINGGLYPH{\char"018C} \MISSINGGLYPH{\char"0221} \MISSINGGLYPH{\char"0068} \MISSINGGLYPH{\char"2C68} \MISSINGGLYPH{\char"006B} \MISSINGGLYPH{\char"0199} \MISSINGGLYPH{\char"2C6A} \MISSINGGLYPH{\char"006C} \MISSINGGLYPH{\char"0234} \MISSINGGLYPH{\char"006E} \MISSINGGLYPH{\char"0235} \MISSINGGLYPH{\char"0070} \MISSINGGLYPH{\char"01A5} \MISSINGGLYPH{\char"0071} \MISSINGGLYPH{\char"02A0} \MISSINGGLYPH{\char"0074} \MISSINGGLYPH{\char"01AB} \MISSINGGLYPH{\char"01AD} \MISSINGGLYPH{\char"0236} \MISSINGGLYPH{\char"0075} \MISSINGGLYPH{\char"AB52} \MISSINGGLYPH{\char"0079} \MISSINGGLYPH{\char"01B4} \MISSINGGLYPH{\char"1EFF}

\MISSINGGLYPH{\char"0048} \MISSINGGLYPH{\char"2C67} \MISSINGGLYPH{\char"004B} \MISSINGGLYPH{\char"2C69} \MISSINGGLYPH{\char"004D} \MISSINGGLYPH{\char"2C6E} \MISSINGGLYPH{\char"0050} \MISSINGGLYPH{\char"01A4} \MISSINGGLYPH{\char"0052} \MISSINGGLYPH{\char"2C64} \MISSINGGLYPH{\char"0054} \MISSINGGLYPH{\char"01AC} \MISSINGGLYPH{\char"01AE} \MISSINGGLYPH{\char"0059} \MISSINGGLYPH{\char"01B3} \MISSINGGLYPH{\char"0061} \MISSINGGLYPH{\char"1E9A} \MISSINGGLYPH{\char"0062} \MISSINGGLYPH{\char"0183} \MISSINGGLYPH{\char"0063} \MISSINGGLYPH{\char"0188} \MISSINGGLYPH{\char"0064} \MISSINGGLYPH{\char"018C} \MISSINGGLYPH{\char"0221} \MISSINGGLYPH{\char"0068} \MISSINGGLYPH{\char"2C68} \MISSINGGLYPH{\char"006B} \MISSINGGLYPH{\char"0199} \MISSINGGLYPH{\char"2C6A} \MISSINGGLYPH{\char"006C} \MISSINGGLYPH{\char"0234} \MISSINGGLYPH{\char"006E} \MISSINGGLYPH{\char"0235} \MISSINGGLYPH{\char"0070} \MISSINGGLYPH{\char"01A5} \MISSINGGLYPH{\char"0071} \MISSINGGLYPH{\char"02A0} \MISSINGGLYPH{\char"0074} \MISSINGGLYPH{\char"01AB} \MISSINGGLYPH{\char"01AD} \MISSINGGLYPH{\char"0236} \MISSINGGLYPH{\char"0075} \MISSINGGLYPH{\char"AB52} \MISSINGGLYPH{\char"0079} \MISSINGGLYPH{\char"01B4} \MISSINGGLYPH{\char"1EFF}

\MISSINGGLYPH{\char"0048} \MISSINGGLYPH{\char"2C67} \MISSINGGLYPH{\char"004B} \MISSINGGLYPH{\char"2C69} \MISSINGGLYPH{\char"004D} \MISSINGGLYPH{\char"2C6E} \MISSINGGLYPH{\char"0050} \MISSINGGLYPH{\char"01A4} \MISSINGGLYPH{\char"0052} \MISSINGGLYPH{\char"2C64} \MISSINGGLYPH{\char"0054} \MISSINGGLYPH{\char"01AC} \MISSINGGLYPH{\char"01AE} \MISSINGGLYPH{\char"0059} \MISSINGGLYPH{\char"01B3} \MISSINGGLYPH{\char"0061} \MISSINGGLYPH{\char"1E9A} \MISSINGGLYPH{\char"0062} \MISSINGGLYPH{\char"0183} \MISSINGGLYPH{\char"0063} \MISSINGGLYPH{\char"0188} \MISSINGGLYPH{\char"0064} \MISSINGGLYPH{\char"018C} \MISSINGGLYPH{\char"0221} \MISSINGGLYPH{\char"0068} \MISSINGGLYPH{\char"2C68} \MISSINGGLYPH{\char"006B} \MISSINGGLYPH{\char"0199} \MISSINGGLYPH{\char"2C6A} \MISSINGGLYPH{\char"006C} \MISSINGGLYPH{\char"0234} \MISSINGGLYPH{\char"006E} \MISSINGGLYPH{\char"0235} \MISSINGGLYPH{\char"0070} \MISSINGGLYPH{\char"01A5} \MISSINGGLYPH{\char"0071} \MISSINGGLYPH{\char"02A0} \MISSINGGLYPH{\char"0074} \MISSINGGLYPH{\char"01AB} \MISSINGGLYPH{\char"01AD} \MISSINGGLYPH{\char"0236} \MISSINGGLYPH{\char"0075} \MISSINGGLYPH{\char"AB52} \MISSINGGLYPH{\char"0079} \MISSINGGLYPH{\char"01B4} \MISSINGGLYPH{\char"1EFF}

\MISSINGGLYPH{\char"0048} \MISSINGGLYPH{\char"2C67} \MISSINGGLYPH{\char"004B} \MISSINGGLYPH{\char"2C69} \MISSINGGLYPH{\char"004D} \MISSINGGLYPH{\char"2C6E} \MISSINGGLYPH{\char"0050} \MISSINGGLYPH{\char"01A4} \MISSINGGLYPH{\char"0052} \MISSINGGLYPH{\char"2C64} \MISSINGGLYPH{\char"0054} \MISSINGGLYPH{\char"01AC} \MISSINGGLYPH{\char"01AE} \MISSINGGLYPH{\char"0059} \MISSINGGLYPH{\char"01B3} \MISSINGGLYPH{\char"0061} \MISSINGGLYPH{\char"1E9A} \MISSINGGLYPH{\char"0062} \MISSINGGLYPH{\char"0183} \MISSINGGLYPH{\char"0063} \MISSINGGLYPH{\char"0188} \MISSINGGLYPH{\char"0064} \MISSINGGLYPH{\char"018C} \MISSINGGLYPH{\char"0221} \MISSINGGLYPH{\char"0068} \MISSINGGLYPH{\char"2C68} \MISSINGGLYPH{\char"006B} \MISSINGGLYPH{\char"0199} \MISSINGGLYPH{\char"2C6A} \MISSINGGLYPH{\char"006C} \MISSINGGLYPH{\char"0234} \MISSINGGLYPH{\char"006E} \MISSINGGLYPH{\char"0235} \MISSINGGLYPH{\char"0070} \MISSINGGLYPH{\char"01A5} \MISSINGGLYPH{\char"0071} \MISSINGGLYPH{\char"02A0} \MISSINGGLYPH{\char"0074} \MISSINGGLYPH{\char"01AB} \MISSINGGLYPH{\char"01AD} \MISSINGGLYPH{\char"0236} \MISSINGGLYPH{\char"0075} \MISSINGGLYPH{\char"AB52} \MISSINGGLYPH{\char"0079} \MISSINGGLYPH{\char"01B4} \MISSINGGLYPH{\char"1EFF}

\MISSINGGLYPH{\char"0048} \MISSINGGLYPH{\char"2C67} \MISSINGGLYPH{\char"004B} \MISSINGGLYPH{\char"2C69} \MISSINGGLYPH{\char"004D} \MISSINGGLYPH{\char"2C6E} \MISSINGGLYPH{\char"0050} \MISSINGGLYPH{\char"01A4} \MISSINGGLYPH{\char"0052} \MISSINGGLYPH{\char"2C64} \MISSINGGLYPH{\char"0054} \MISSINGGLYPH{\char"01AC} \MISSINGGLYPH{\char"01AE} \MISSINGGLYPH{\char"0059} \MISSINGGLYPH{\char"01B3} \MISSINGGLYPH{\char"0061} \MISSINGGLYPH{\char"1E9A} \MISSINGGLYPH{\char"0062} \MISSINGGLYPH{\char"0183} \MISSINGGLYPH{\char"0063} \MISSINGGLYPH{\char"0188} \MISSINGGLYPH{\char"0064} \MISSINGGLYPH{\char"018C} \MISSINGGLYPH{\char"0221} \MISSINGGLYPH{\char"0068} \MISSINGGLYPH{\char"2C68} \MISSINGGLYPH{\char"006B} \MISSINGGLYPH{\char"0199} \MISSINGGLYPH{\char"2C6A} \MISSINGGLYPH{\char"006C} \MISSINGGLYPH{\char"0234} \MISSINGGLYPH{\char"006E} \MISSINGGLYPH{\char"0235} \MISSINGGLYPH{\char"0070} \MISSINGGLYPH{\char"01A5} \MISSINGGLYPH{\char"0071} \MISSINGGLYPH{\char"02A0} \MISSINGGLYPH{\char"0074} \MISSINGGLYPH{\char"01AB} \MISSINGGLYPH{\char"01AD} \MISSINGGLYPH{\char"0236} \MISSINGGLYPH{\char"0075} \MISSINGGLYPH{\char"AB52} \MISSINGGLYPH{\char"0079} \MISSINGGLYPH{\char"01B4} \MISSINGGLYPH{\char"1EFF}

\MISSINGGLYPH{\char"0048} \MISSINGGLYPH{\char"2C67} \MISSINGGLYPH{\char"004B} \MISSINGGLYPH{\char"2C69} \MISSINGGLYPH{\char"004D} \MISSINGGLYPH{\char"2C6E} \MISSINGGLYPH{\char"0050} \MISSINGGLYPH{\char"01A4} \MISSINGGLYPH{\char"0052} \MISSINGGLYPH{\char"2C64} \MISSINGGLYPH{\char"0054} \MISSINGGLYPH{\char"01AC} \MISSINGGLYPH{\char"01AE} \MISSINGGLYPH{\char"0059} \MISSINGGLYPH{\char"01B3} \MISSINGGLYPH{\char"0061} \MISSINGGLYPH{\char"1E9A} \MISSINGGLYPH{\char"0062} \MISSINGGLYPH{\char"0183} \MISSINGGLYPH{\char"0063} \MISSINGGLYPH{\char"0188} \MISSINGGLYPH{\char"0064} \MISSINGGLYPH{\char"018C} \MISSINGGLYPH{\char"0221} \MISSINGGLYPH{\char"0068} \MISSINGGLYPH{\char"2C68} \MISSINGGLYPH{\char"006B} \MISSINGGLYPH{\char"0199} \MISSINGGLYPH{\char"2C6A} \MISSINGGLYPH{\char"006C} \MISSINGGLYPH{\char"0234} \MISSINGGLYPH{\char"006E} \MISSINGGLYPH{\char"0235} \MISSINGGLYPH{\char"0070} \MISSINGGLYPH{\char"01A5} \MISSINGGLYPH{\char"0071} \MISSINGGLYPH{\char"02A0} \MISSINGGLYPH{\char"0074} \MISSINGGLYPH{\char"01AB} \MISSINGGLYPH{\char"01AD} \MISSINGGLYPH{\char"0236} \MISSINGGLYPH{\char"0075} \MISSINGGLYPH{\char"AB52} \MISSINGGLYPH{\char"0079} \MISSINGGLYPH{\char"01B4} \MISSINGGLYPH{\char"1EFF}

\MISSINGGLYPH{\char"0048} \MISSINGGLYPH{\char"2C67} \MISSINGGLYPH{\char"004B} \MISSINGGLYPH{\char"2C69} \MISSINGGLYPH{\char"004D} \MISSINGGLYPH{\char"2C6E} \MISSINGGLYPH{\char"0050} \MISSINGGLYPH{\char"01A4} \MISSINGGLYPH{\char"0052} \MISSINGGLYPH{\char"2C64} \MISSINGGLYPH{\char"0054} \MISSINGGLYPH{\char"01AC} \MISSINGGLYPH{\char"01AE} \MISSINGGLYPH{\char"0059} \MISSINGGLYPH{\char"01B3} \MISSINGGLYPH{\char"0061} \MISSINGGLYPH{\char"1E9A} \MISSINGGLYPH{\char"0062} \MISSINGGLYPH{\char"0183} \MISSINGGLYPH{\char"0063} \MISSINGGLYPH{\char"0188} \MISSINGGLYPH{\char"0064} \MISSINGGLYPH{\char"018C} \MISSINGGLYPH{\char"0221} \MISSINGGLYPH{\char"0068} \MISSINGGLYPH{\char"2C68} \MISSINGGLYPH{\char"006B} \MISSINGGLYPH{\char"0199} \MISSINGGLYPH{\char"2C6A} \MISSINGGLYPH{\char"006C} \MISSINGGLYPH{\char"0234} \MISSINGGLYPH{\char"006E} \MISSINGGLYPH{\char"0235} \MISSINGGLYPH{\char"0070} \MISSINGGLYPH{\char"01A5} \MISSINGGLYPH{\char"0071} \MISSINGGLYPH{\char"02A0} \MISSINGGLYPH{\char"0074} \MISSINGGLYPH{\char"01AB} \MISSINGGLYPH{\char"01AD} \MISSINGGLYPH{\char"0236} \MISSINGGLYPH{\char"0075} \MISSINGGLYPH{\char"AB52} \MISSINGGLYPH{\char"0079} \MISSINGGLYPH{\char"01B4} \MISSINGGLYPH{\char"1EFF}

\MISSINGGLYPH{\char"0048} \MISSINGGLYPH{\char"2C67} \MISSINGGLYPH{\char"004B} \MISSINGGLYPH{\char"2C69} \MISSINGGLYPH{\char"004D} \MISSINGGLYPH{\char"2C6E} \MISSINGGLYPH{\char"0050} \MISSINGGLYPH{\char"01A4} \MISSINGGLYPH{\char"0052} \MISSINGGLYPH{\char"2C64} \MISSINGGLYPH{\char"0054} \MISSINGGLYPH{\char"01AC} \MISSINGGLYPH{\char"01AE} \MISSINGGLYPH{\char"0059} \MISSINGGLYPH{\char"01B3} \MISSINGGLYPH{\char"0061} \MISSINGGLYPH{\char"1E9A} \MISSINGGLYPH{\char"0062} \MISSINGGLYPH{\char"0183} \MISSINGGLYPH{\char"0063} \MISSINGGLYPH{\char"0188} \MISSINGGLYPH{\char"0064} \MISSINGGLYPH{\char"018C} \MISSINGGLYPH{\char"0221} \MISSINGGLYPH{\char"0068} \MISSINGGLYPH{\char"2C68} \MISSINGGLYPH{\char"006B} \MISSINGGLYPH{\char"0199} \MISSINGGLYPH{\char"2C6A} \MISSINGGLYPH{\char"006C} \MISSINGGLYPH{\char"0234} \MISSINGGLYPH{\char"006E} \MISSINGGLYPH{\char"0235} \MISSINGGLYPH{\char"0070} \MISSINGGLYPH{\char"01A5} \MISSINGGLYPH{\char"0071} \MISSINGGLYPH{\char"02A0} \MISSINGGLYPH{\char"0074} \MISSINGGLYPH{\char"01AB} \MISSINGGLYPH{\char"01AD} \MISSINGGLYPH{\char"0236} \MISSINGGLYPH{\char"0075} \MISSINGGLYPH{\char"AB52} \MISSINGGLYPH{\char"0079} \MISSINGGLYPH{\char"01B4} \MISSINGGLYPH{\char"1EFF}

\MISSINGGLYPH{\char"0048} \MISSINGGLYPH{\char"2C67} \MISSINGGLYPH{\char"004B} \MISSINGGLYPH{\char"2C69} \MISSINGGLYPH{\char"004D} \MISSINGGLYPH{\char"2C6E} \MISSINGGLYPH{\char"0050} \MISSINGGLYPH{\char"01A4} \MISSINGGLYPH{\char"0052} \MISSINGGLYPH{\char"2C64} \MISSINGGLYPH{\char"0054} \MISSINGGLYPH{\char"01AC} \MISSINGGLYPH{\char"01AE} \MISSINGGLYPH{\char"0059} \MISSINGGLYPH{\char"01B3} \MISSINGGLYPH{\char"0061} \MISSINGGLYPH{\char"1E9A} \MISSINGGLYPH{\char"0062} \MISSINGGLYPH{\char"0183} \MISSINGGLYPH{\char"0063} \MISSINGGLYPH{\char"0188} \MISSINGGLYPH{\char"0064} \MISSINGGLYPH{\char"018C} \MISSINGGLYPH{\char"0221} \MISSINGGLYPH{\char"0068} \MISSINGGLYPH{\char"2C68} \MISSINGGLYPH{\char"006B} \MISSINGGLYPH{\char"0199} \MISSINGGLYPH{\char"2C6A} \MISSINGGLYPH{\char"006C} \MISSINGGLYPH{\char"0234} \MISSINGGLYPH{\char"006E} \MISSINGGLYPH{\char"0235} \MISSINGGLYPH{\char"0070} \MISSINGGLYPH{\char"01A5} \MISSINGGLYPH{\char"0071} \MISSINGGLYPH{\char"02A0} \MISSINGGLYPH{\char"0074} \MISSINGGLYPH{\char"01AB} \MISSINGGLYPH{\char"01AD} \MISSINGGLYPH{\char"0236} \MISSINGGLYPH{\char"0075} \MISSINGGLYPH{\char"AB52} \MISSINGGLYPH{\char"0079} \MISSINGGLYPH{\char"01B4} \MISSINGGLYPH{\char"1EFF}

}

\newpage
\subsection*{Semibold italic patched, Group 4, below, Below}

{\LibertinusSerifPatch\bfseries\itshape
% grid. columns are bases, rows are marks.
\SUPPORTEDANCH{\char"0048\char"0316} \SUPPORTEDANCH{\char"2C67\char"0316} \SUPPORTEDANCH{\char"004B\char"0316} \SUPPORTEDANCH{\char"2C69\char"0316} \SUPPORTEDANCH{\char"004D\char"0316} \MISSINGGLYPH{\char"2C6E\char"0316} \SUPPORTEDANCH{\char"0050\char"0316} \SUPPORTEDANCH{\char"01A4\char"0316} \SUPPORTEDANCH{\char"0052\char"0316} \SUPPORTEDANCH{\char"2C64\char"0316} \SUPPORTEDANCH{\char"0054\char"0316} \SUPPORTEDANCH{\char"01AC\char"0316} \SUPPORTEDANCH{\char"01AE\char"0316} \SUPPORTEDANCH{\char"0059\char"0316} \SUPPORTEDANCH{\char"01B3\char"0316} \SUPPORTEDANCH{\char"0061\char"0316} \SUPPORTEDANCH{\char"1E9A\char"0316} \SUPPORTEDANCH{\char"0062\char"0316} \SUPPORTEDANCH{\char"0183\char"0316} \SUPPORTEDANCH{\char"0063\char"0316} \SUPPORTEDANCH{\char"0188\char"0316} \SUPPORTEDANCH{\char"0064\char"0316} \SUPPORTEDANCH{\char"018C\char"0316} \SUPPORTEDANCH{\char"0221\char"0316} \SUPPORTEDANCH{\char"0068\char"0316} \SUPPORTEDANCH{\char"2C68\char"0316} \SUPPORTEDANCH{\char"006B\char"0316} \SUPPORTEDANCH{\char"0199\char"0316} \SUPPORTEDANCH{\char"2C6A\char"0316} \SUPPORTEDANCH{\char"006C\char"0316} \SUPPORTEDANCH{\char"0234\char"0316} \SUPPORTEDANCH{\char"006E\char"0316} \SUPPORTEDANCH{\char"0235\char"0316} \SUPPORTEDANCH{\char"0070\char"0316} \SUPPORTEDANCH{\char"01A5\char"0316} \SUPPORTEDANCH{\char"0071\char"0316} \SUPPORTEDANCH{\char"02A0\char"0316} \SUPPORTEDANCH{\char"0074\char"0316} \SUPPORTEDANCH{\char"01AB\char"0316} \SUPPORTEDANCH{\char"01AD\char"0316} \SUPPORTEDANCH{\char"0236\char"0316} \SUPPORTEDANCH{\char"0075\char"0316} \MISSINGGLYPH{\char"AB52\char"0316} \SUPPORTEDANCH{\char"0079\char"0316} \SUPPORTEDANCH{\char"01B4\char"0316} \MISSINGGLYPH{\char"1EFF\char"0316}

\SUPPORTEDANCH{\char"0048\char"0317} \SUPPORTEDANCH{\char"2C67\char"0317} \SUPPORTEDANCH{\char"004B\char"0317} \SUPPORTEDANCH{\char"2C69\char"0317} \SUPPORTEDANCH{\char"004D\char"0317} \MISSINGGLYPH{\char"2C6E\char"0317} \SUPPORTEDANCH{\char"0050\char"0317} \SUPPORTEDANCH{\char"01A4\char"0317} \SUPPORTEDANCH{\char"0052\char"0317} \SUPPORTEDANCH{\char"2C64\char"0317} \SUPPORTEDANCH{\char"0054\char"0317} \SUPPORTEDANCH{\char"01AC\char"0317} \SUPPORTEDANCH{\char"01AE\char"0317} \SUPPORTEDANCH{\char"0059\char"0317} \SUPPORTEDANCH{\char"01B3\char"0317} \SUPPORTEDANCH{\char"0061\char"0317} \SUPPORTEDANCH{\char"1E9A\char"0317} \SUPPORTEDANCH{\char"0062\char"0317} \SUPPORTEDANCH{\char"0183\char"0317} \SUPPORTEDANCH{\char"0063\char"0317} \SUPPORTEDANCH{\char"0188\char"0317} \SUPPORTEDANCH{\char"0064\char"0317} \SUPPORTEDANCH{\char"018C\char"0317} \SUPPORTEDANCH{\char"0221\char"0317} \SUPPORTEDANCH{\char"0068\char"0317} \SUPPORTEDANCH{\char"2C68\char"0317} \SUPPORTEDANCH{\char"006B\char"0317} \SUPPORTEDANCH{\char"0199\char"0317} \SUPPORTEDANCH{\char"2C6A\char"0317} \SUPPORTEDANCH{\char"006C\char"0317} \SUPPORTEDANCH{\char"0234\char"0317} \SUPPORTEDANCH{\char"006E\char"0317} \SUPPORTEDANCH{\char"0235\char"0317} \SUPPORTEDANCH{\char"0070\char"0317} \SUPPORTEDANCH{\char"01A5\char"0317} \SUPPORTEDANCH{\char"0071\char"0317} \SUPPORTEDANCH{\char"02A0\char"0317} \SUPPORTEDANCH{\char"0074\char"0317} \SUPPORTEDANCH{\char"01AB\char"0317} \SUPPORTEDANCH{\char"01AD\char"0317} \SUPPORTEDANCH{\char"0236\char"0317} \SUPPORTEDANCH{\char"0075\char"0317} \MISSINGGLYPH{\char"AB52\char"0317} \SUPPORTEDANCH{\char"0079\char"0317} \SUPPORTEDANCH{\char"01B4\char"0317} \MISSINGGLYPH{\char"1EFF\char"0317}

\SUPPORTEDANCH{\char"0048\char"0318} \SUPPORTEDANCH{\char"2C67\char"0318} \SUPPORTEDANCH{\char"004B\char"0318} \SUPPORTEDANCH{\char"2C69\char"0318} \SUPPORTEDANCH{\char"004D\char"0318} \MISSINGGLYPH{\char"2C6E\char"0318} \SUPPORTEDANCH{\char"0050\char"0318} \SUPPORTEDANCH{\char"01A4\char"0318} \SUPPORTEDANCH{\char"0052\char"0318} \SUPPORTEDANCH{\char"2C64\char"0318} \SUPPORTEDANCH{\char"0054\char"0318} \SUPPORTEDANCH{\char"01AC\char"0318} \SUPPORTEDANCH{\char"01AE\char"0318} \SUPPORTEDANCH{\char"0059\char"0318} \SUPPORTEDANCH{\char"01B3\char"0318} \SUPPORTEDANCH{\char"0061\char"0318} \SUPPORTEDANCH{\char"1E9A\char"0318} \SUPPORTEDANCH{\char"0062\char"0318} \SUPPORTEDANCH{\char"0183\char"0318} \SUPPORTEDANCH{\char"0063\char"0318} \SUPPORTEDANCH{\char"0188\char"0318} \SUPPORTEDANCH{\char"0064\char"0318} \SUPPORTEDANCH{\char"018C\char"0318} \SUPPORTEDANCH{\char"0221\char"0318} \SUPPORTEDANCH{\char"0068\char"0318} \SUPPORTEDANCH{\char"2C68\char"0318} \SUPPORTEDANCH{\char"006B\char"0318} \SUPPORTEDANCH{\char"0199\char"0318} \SUPPORTEDANCH{\char"2C6A\char"0318} \SUPPORTEDANCH{\char"006C\char"0318} \SUPPORTEDANCH{\char"0234\char"0318} \SUPPORTEDANCH{\char"006E\char"0318} \SUPPORTEDANCH{\char"0235\char"0318} \SUPPORTEDANCH{\char"0070\char"0318} \SUPPORTEDANCH{\char"01A5\char"0318} \SUPPORTEDANCH{\char"0071\char"0318} \SUPPORTEDANCH{\char"02A0\char"0318} \SUPPORTEDANCH{\char"0074\char"0318} \SUPPORTEDANCH{\char"01AB\char"0318} \SUPPORTEDANCH{\char"01AD\char"0318} \SUPPORTEDANCH{\char"0236\char"0318} \SUPPORTEDANCH{\char"0075\char"0318} \MISSINGGLYPH{\char"AB52\char"0318} \SUPPORTEDANCH{\char"0079\char"0318} \SUPPORTEDANCH{\char"01B4\char"0318} \MISSINGGLYPH{\char"1EFF\char"0318}

\SUPPORTEDANCH{\char"0048\char"0319} \SUPPORTEDANCH{\char"2C67\char"0319} \SUPPORTEDANCH{\char"004B\char"0319} \SUPPORTEDANCH{\char"2C69\char"0319} \SUPPORTEDANCH{\char"004D\char"0319} \MISSINGGLYPH{\char"2C6E\char"0319} \SUPPORTEDANCH{\char"0050\char"0319} \SUPPORTEDANCH{\char"01A4\char"0319} \SUPPORTEDANCH{\char"0052\char"0319} \SUPPORTEDANCH{\char"2C64\char"0319} \SUPPORTEDANCH{\char"0054\char"0319} \SUPPORTEDANCH{\char"01AC\char"0319} \SUPPORTEDANCH{\char"01AE\char"0319} \SUPPORTEDANCH{\char"0059\char"0319} \SUPPORTEDANCH{\char"01B3\char"0319} \SUPPORTEDANCH{\char"0061\char"0319} \SUPPORTEDANCH{\char"1E9A\char"0319} \SUPPORTEDANCH{\char"0062\char"0319} \SUPPORTEDANCH{\char"0183\char"0319} \SUPPORTEDANCH{\char"0063\char"0319} \SUPPORTEDANCH{\char"0188\char"0319} \SUPPORTEDANCH{\char"0064\char"0319} \SUPPORTEDANCH{\char"018C\char"0319} \SUPPORTEDANCH{\char"0221\char"0319} \SUPPORTEDANCH{\char"0068\char"0319} \SUPPORTEDANCH{\char"2C68\char"0319} \SUPPORTEDANCH{\char"006B\char"0319} \SUPPORTEDANCH{\char"0199\char"0319} \SUPPORTEDANCH{\char"2C6A\char"0319} \SUPPORTEDANCH{\char"006C\char"0319} \SUPPORTEDANCH{\char"0234\char"0319} \SUPPORTEDANCH{\char"006E\char"0319} \SUPPORTEDANCH{\char"0235\char"0319} \SUPPORTEDANCH{\char"0070\char"0319} \SUPPORTEDANCH{\char"01A5\char"0319} \SUPPORTEDANCH{\char"0071\char"0319} \SUPPORTEDANCH{\char"02A0\char"0319} \SUPPORTEDANCH{\char"0074\char"0319} \SUPPORTEDANCH{\char"01AB\char"0319} \SUPPORTEDANCH{\char"01AD\char"0319} \SUPPORTEDANCH{\char"0236\char"0319} \SUPPORTEDANCH{\char"0075\char"0319} \MISSINGGLYPH{\char"AB52\char"0319} \SUPPORTEDANCH{\char"0079\char"0319} \SUPPORTEDANCH{\char"01B4\char"0319} \MISSINGGLYPH{\char"1EFF\char"0319}

\SUPPORTEDANCH{\char"0048\char"031C} \SUPPORTEDANCH{\char"2C67\char"031C} \SUPPORTEDANCH{\char"004B\char"031C} \SUPPORTEDANCH{\char"2C69\char"031C} \SUPPORTEDANCH{\char"004D\char"031C} \MISSINGGLYPH{\char"2C6E\char"031C} \SUPPORTEDANCH{\char"0050\char"031C} \SUPPORTEDANCH{\char"01A4\char"031C} \SUPPORTEDANCH{\char"0052\char"031C} \SUPPORTEDANCH{\char"2C64\char"031C} \SUPPORTEDANCH{\char"0054\char"031C} \SUPPORTEDANCH{\char"01AC\char"031C} \SUPPORTEDANCH{\char"01AE\char"031C} \SUPPORTEDANCH{\char"0059\char"031C} \SUPPORTEDANCH{\char"01B3\char"031C} \SUPPORTEDANCH{\char"0061\char"031C} \SUPPORTEDANCH{\char"1E9A\char"031C} \SUPPORTEDANCH{\char"0062\char"031C} \SUPPORTEDANCH{\char"0183\char"031C} \SUPPORTEDANCH{\char"0063\char"031C} \SUPPORTEDANCH{\char"0188\char"031C} \SUPPORTEDANCH{\char"0064\char"031C} \SUPPORTEDANCH{\char"018C\char"031C} \SUPPORTEDANCH{\char"0221\char"031C} \SUPPORTEDANCH{\char"0068\char"031C} \SUPPORTEDANCH{\char"2C68\char"031C} \SUPPORTEDANCH{\char"006B\char"031C} \SUPPORTEDANCH{\char"0199\char"031C} \SUPPORTEDANCH{\char"2C6A\char"031C} \SUPPORTEDANCH{\char"006C\char"031C} \SUPPORTEDANCH{\char"0234\char"031C} \SUPPORTEDANCH{\char"006E\char"031C} \SUPPORTEDANCH{\char"0235\char"031C} \SUPPORTEDANCH{\char"0070\char"031C} \SUPPORTEDANCH{\char"01A5\char"031C} \SUPPORTEDANCH{\char"0071\char"031C} \SUPPORTEDANCH{\char"02A0\char"031C} \SUPPORTEDANCH{\char"0074\char"031C} \SUPPORTEDANCH{\char"01AB\char"031C} \SUPPORTEDANCH{\char"01AD\char"031C} \SUPPORTEDANCH{\char"0236\char"031C} \SUPPORTEDANCH{\char"0075\char"031C} \MISSINGGLYPH{\char"AB52\char"031C} \SUPPORTEDANCH{\char"0079\char"031C} \SUPPORTEDANCH{\char"01B4\char"031C} \MISSINGGLYPH{\char"1EFF\char"031C}

\SUPPORTEDANCH{\char"0048\char"0320} \SUPPORTEDANCH{\char"2C67\char"0320} \SUPPORTEDANCH{\char"004B\char"0320} \SUPPORTEDANCH{\char"2C69\char"0320} \SUPPORTEDANCH{\char"004D\char"0320} \MISSINGGLYPH{\char"2C6E\char"0320} \SUPPORTEDANCH{\char"0050\char"0320} \SUPPORTEDANCH{\char"01A4\char"0320} \SUPPORTEDANCH{\char"0052\char"0320} \SUPPORTEDANCH{\char"2C64\char"0320} \SUPPORTEDANCH{\char"0054\char"0320} \SUPPORTEDANCH{\char"01AC\char"0320} \SUPPORTEDANCH{\char"01AE\char"0320} \SUPPORTEDANCH{\char"0059\char"0320} \SUPPORTEDANCH{\char"01B3\char"0320} \SUPPORTEDANCH{\char"0061\char"0320} \SUPPORTEDANCH{\char"1E9A\char"0320} \SUPPORTEDANCH{\char"0062\char"0320} \SUPPORTEDANCH{\char"0183\char"0320} \SUPPORTEDANCH{\char"0063\char"0320} \SUPPORTEDANCH{\char"0188\char"0320} \SUPPORTEDANCH{\char"0064\char"0320} \SUPPORTEDANCH{\char"018C\char"0320} \SUPPORTEDANCH{\char"0221\char"0320} \SUPPORTEDANCH{\char"0068\char"0320} \SUPPORTEDANCH{\char"2C68\char"0320} \SUPPORTEDANCH{\char"006B\char"0320} \SUPPORTEDANCH{\char"0199\char"0320} \SUPPORTEDANCH{\char"2C6A\char"0320} \SUPPORTEDANCH{\char"006C\char"0320} \SUPPORTEDANCH{\char"0234\char"0320} \SUPPORTEDANCH{\char"006E\char"0320} \SUPPORTEDANCH{\char"0235\char"0320} \SUPPORTEDANCH{\char"0070\char"0320} \SUPPORTEDANCH{\char"01A5\char"0320} \SUPPORTEDANCH{\char"0071\char"0320} \SUPPORTEDANCH{\char"02A0\char"0320} \SUPPORTEDANCH{\char"0074\char"0320} \SUPPORTEDANCH{\char"01AB\char"0320} \SUPPORTEDANCH{\char"01AD\char"0320} \SUPPORTEDANCH{\char"0236\char"0320} \SUPPORTEDANCH{\char"0075\char"0320} \MISSINGGLYPH{\char"AB52\char"0320} \SUPPORTEDANCH{\char"0079\char"0320} \SUPPORTEDANCH{\char"01B4\char"0320} \MISSINGGLYPH{\char"1EFF\char"0320}

\SUPPORTEDPREC{\char"0048\char"0323} \SUPPORTEDANCH{\char"2C67\char"0323} \SUPPORTEDPREC{\char"004B\char"0323} \SUPPORTEDANCH{\char"2C69\char"0323} \SUPPORTEDPREC{\char"004D\char"0323} \MISSINGGLYPH{\char"2C6E\char"0323} \SUPPORTEDANCH{\char"0050\char"0323} \SUPPORTEDANCH{\char"01A4\char"0323} \SUPPORTEDPREC{\char"0052\char"0323} \SUPPORTEDANCH{\char"2C64\char"0323} \SUPPORTEDPREC{\char"0054\char"0323} \SUPPORTEDANCH{\char"01AC\char"0323} \SUPPORTEDANCH{\char"01AE\char"0323} \SUPPORTEDPREC{\char"0059\char"0323} \SUPPORTEDANCH{\char"01B3\char"0323} \SUPPORTEDPREC{\char"0061\char"0323} \SUPPORTEDANCH{\char"1E9A\char"0323} \SUPPORTEDPREC{\char"0062\char"0323} \SUPPORTEDANCH{\char"0183\char"0323} \SUPPORTEDANCH{\char"0063\char"0323} \SUPPORTEDANCH{\char"0188\char"0323} \SUPPORTEDPREC{\char"0064\char"0323} \SUPPORTEDANCH{\char"018C\char"0323} \SUPPORTEDANCH{\char"0221\char"0323} \SUPPORTEDPREC{\char"0068\char"0323} \SUPPORTEDANCH{\char"2C68\char"0323} \SUPPORTEDPREC{\char"006B\char"0323} \SUPPORTEDANCH{\char"0199\char"0323} \SUPPORTEDANCH{\char"2C6A\char"0323} \SUPPORTEDPREC{\char"006C\char"0323} \SUPPORTEDANCH{\char"0234\char"0323} \SUPPORTEDPREC{\char"006E\char"0323} \SUPPORTEDANCH{\char"0235\char"0323} \SUPPORTEDANCH{\char"0070\char"0323} \SUPPORTEDANCH{\char"01A5\char"0323} \SUPPORTEDANCH{\char"0071\char"0323} \SUPPORTEDANCH{\char"02A0\char"0323} \SUPPORTEDPREC{\char"0074\char"0323} \SUPPORTEDANCH{\char"01AB\char"0323} \SUPPORTEDANCH{\char"01AD\char"0323} \SUPPORTEDANCH{\char"0236\char"0323} \SUPPORTEDPREC{\char"0075\char"0323} \MISSINGGLYPH{\char"AB52\char"0323} \SUPPORTEDPREC{\char"0079\char"0323} \SUPPORTEDANCH{\char"01B4\char"0323} \MISSINGGLYPH{\char"1EFF\char"0323}

\SUPPORTEDANCH{\char"0048\char"0324} \SUPPORTEDANCH{\char"2C67\char"0324} \SUPPORTEDANCH{\char"004B\char"0324} \SUPPORTEDANCH{\char"2C69\char"0324} \SUPPORTEDANCH{\char"004D\char"0324} \MISSINGGLYPH{\char"2C6E\char"0324} \SUPPORTEDANCH{\char"0050\char"0324} \SUPPORTEDANCH{\char"01A4\char"0324} \SUPPORTEDANCH{\char"0052\char"0324} \SUPPORTEDANCH{\char"2C64\char"0324} \SUPPORTEDANCH{\char"0054\char"0324} \SUPPORTEDANCH{\char"01AC\char"0324} \SUPPORTEDANCH{\char"01AE\char"0324} \SUPPORTEDANCH{\char"0059\char"0324} \SUPPORTEDANCH{\char"01B3\char"0324} \SUPPORTEDANCH{\char"0061\char"0324} \SUPPORTEDANCH{\char"1E9A\char"0324} \SUPPORTEDANCH{\char"0062\char"0324} \SUPPORTEDANCH{\char"0183\char"0324} \SUPPORTEDANCH{\char"0063\char"0324} \SUPPORTEDANCH{\char"0188\char"0324} \SUPPORTEDANCH{\char"0064\char"0324} \SUPPORTEDANCH{\char"018C\char"0324} \SUPPORTEDANCH{\char"0221\char"0324} \SUPPORTEDANCH{\char"0068\char"0324} \SUPPORTEDANCH{\char"2C68\char"0324} \SUPPORTEDANCH{\char"006B\char"0324} \SUPPORTEDANCH{\char"0199\char"0324} \SUPPORTEDANCH{\char"2C6A\char"0324} \SUPPORTEDANCH{\char"006C\char"0324} \SUPPORTEDANCH{\char"0234\char"0324} \SUPPORTEDANCH{\char"006E\char"0324} \SUPPORTEDANCH{\char"0235\char"0324} \SUPPORTEDANCH{\char"0070\char"0324} \SUPPORTEDANCH{\char"01A5\char"0324} \SUPPORTEDANCH{\char"0071\char"0324} \SUPPORTEDANCH{\char"02A0\char"0324} \SUPPORTEDANCH{\char"0074\char"0324} \SUPPORTEDANCH{\char"01AB\char"0324} \SUPPORTEDANCH{\char"01AD\char"0324} \SUPPORTEDANCH{\char"0236\char"0324} \SUPPORTEDPREC{\char"0075\char"0324} \MISSINGGLYPH{\char"AB52\char"0324} \SUPPORTEDANCH{\char"0079\char"0324} \SUPPORTEDANCH{\char"01B4\char"0324} \MISSINGGLYPH{\char"1EFF\char"0324}

\SUPPORTEDANCH{\char"0048\char"0325} \SUPPORTEDANCH{\char"2C67\char"0325} \SUPPORTEDANCH{\char"004B\char"0325} \SUPPORTEDANCH{\char"2C69\char"0325} \SUPPORTEDANCH{\char"004D\char"0325} \MISSINGGLYPH{\char"2C6E\char"0325} \SUPPORTEDANCH{\char"0050\char"0325} \SUPPORTEDANCH{\char"01A4\char"0325} \SUPPORTEDANCH{\char"0052\char"0325} \SUPPORTEDANCH{\char"2C64\char"0325} \SUPPORTEDANCH{\char"0054\char"0325} \SUPPORTEDANCH{\char"01AC\char"0325} \SUPPORTEDANCH{\char"01AE\char"0325} \SUPPORTEDANCH{\char"0059\char"0325} \SUPPORTEDANCH{\char"01B3\char"0325} \SUPPORTEDPREC{\char"0061\char"0325} \SUPPORTEDANCH{\char"1E9A\char"0325} \SUPPORTEDANCH{\char"0062\char"0325} \SUPPORTEDANCH{\char"0183\char"0325} \SUPPORTEDANCH{\char"0063\char"0325} \SUPPORTEDANCH{\char"0188\char"0325} \SUPPORTEDANCH{\char"0064\char"0325} \SUPPORTEDANCH{\char"018C\char"0325} \SUPPORTEDANCH{\char"0221\char"0325} \SUPPORTEDANCH{\char"0068\char"0325} \SUPPORTEDANCH{\char"2C68\char"0325} \SUPPORTEDANCH{\char"006B\char"0325} \SUPPORTEDANCH{\char"0199\char"0325} \SUPPORTEDANCH{\char"2C6A\char"0325} \SUPPORTEDANCH{\char"006C\char"0325} \SUPPORTEDANCH{\char"0234\char"0325} \SUPPORTEDANCH{\char"006E\char"0325} \SUPPORTEDANCH{\char"0235\char"0325} \SUPPORTEDANCH{\char"0070\char"0325} \SUPPORTEDANCH{\char"01A5\char"0325} \SUPPORTEDANCH{\char"0071\char"0325} \SUPPORTEDANCH{\char"02A0\char"0325} \SUPPORTEDANCH{\char"0074\char"0325} \SUPPORTEDANCH{\char"01AB\char"0325} \SUPPORTEDANCH{\char"01AD\char"0325} \SUPPORTEDANCH{\char"0236\char"0325} \SUPPORTEDANCH{\char"0075\char"0325} \MISSINGGLYPH{\char"AB52\char"0325} \SUPPORTEDANCH{\char"0079\char"0325} \SUPPORTEDANCH{\char"01B4\char"0325} \MISSINGGLYPH{\char"1EFF\char"0325}

\SUPPORTEDANCH{\char"0048\char"0326} \SUPPORTEDANCH{\char"2C67\char"0326} \SUPPORTEDANCH{\char"004B\char"0326} \SUPPORTEDANCH{\char"2C69\char"0326} \SUPPORTEDANCH{\char"004D\char"0326} \MISSINGGLYPH{\char"2C6E\char"0326} \SUPPORTEDANCH{\char"0050\char"0326} \SUPPORTEDANCH{\char"01A4\char"0326} \SUPPORTEDANCH{\char"0052\char"0326} \SUPPORTEDANCH{\char"2C64\char"0326} \SUPPORTEDPREC{\char"0054\char"0326} \SUPPORTEDANCH{\char"01AC\char"0326} \SUPPORTEDANCH{\char"01AE\char"0326} \SUPPORTEDANCH{\char"0059\char"0326} \SUPPORTEDANCH{\char"01B3\char"0326} \SUPPORTEDANCH{\char"0061\char"0326} \SUPPORTEDANCH{\char"1E9A\char"0326} \SUPPORTEDANCH{\char"0062\char"0326} \SUPPORTEDANCH{\char"0183\char"0326} \SUPPORTEDANCH{\char"0063\char"0326} \SUPPORTEDANCH{\char"0188\char"0326} \SUPPORTEDANCH{\char"0064\char"0326} \SUPPORTEDANCH{\char"018C\char"0326} \SUPPORTEDANCH{\char"0221\char"0326} \SUPPORTEDANCH{\char"0068\char"0326} \SUPPORTEDANCH{\char"2C68\char"0326} \SUPPORTEDANCH{\char"006B\char"0326} \SUPPORTEDANCH{\char"0199\char"0326} \SUPPORTEDANCH{\char"2C6A\char"0326} \SUPPORTEDANCH{\char"006C\char"0326} \SUPPORTEDANCH{\char"0234\char"0326} \SUPPORTEDANCH{\char"006E\char"0326} \SUPPORTEDANCH{\char"0235\char"0326} \SUPPORTEDANCH{\char"0070\char"0326} \SUPPORTEDANCH{\char"01A5\char"0326} \SUPPORTEDANCH{\char"0071\char"0326} \SUPPORTEDANCH{\char"02A0\char"0326} \SUPPORTEDPREC{\char"0074\char"0326} \SUPPORTEDANCH{\char"01AB\char"0326} \SUPPORTEDANCH{\char"01AD\char"0326} \SUPPORTEDANCH{\char"0236\char"0326} \SUPPORTEDANCH{\char"0075\char"0326} \MISSINGGLYPH{\char"AB52\char"0326} \SUPPORTEDANCH{\char"0079\char"0326} \SUPPORTEDANCH{\char"01B4\char"0326} \MISSINGGLYPH{\char"1EFF\char"0326}

\SUPPORTEDANCH{\char"0048\char"0329} \SUPPORTEDANCH{\char"2C67\char"0329} \SUPPORTEDANCH{\char"004B\char"0329} \SUPPORTEDANCH{\char"2C69\char"0329} \SUPPORTEDANCH{\char"004D\char"0329} \MISSINGGLYPH{\char"2C6E\char"0329} \SUPPORTEDANCH{\char"0050\char"0329} \SUPPORTEDANCH{\char"01A4\char"0329} \SUPPORTEDANCH{\char"0052\char"0329} \SUPPORTEDANCH{\char"2C64\char"0329} \SUPPORTEDANCH{\char"0054\char"0329} \SUPPORTEDANCH{\char"01AC\char"0329} \SUPPORTEDANCH{\char"01AE\char"0329} \SUPPORTEDANCH{\char"0059\char"0329} \SUPPORTEDANCH{\char"01B3\char"0329} \SUPPORTEDANCH{\char"0061\char"0329} \SUPPORTEDANCH{\char"1E9A\char"0329} \SUPPORTEDANCH{\char"0062\char"0329} \SUPPORTEDANCH{\char"0183\char"0329} \SUPPORTEDANCH{\char"0063\char"0329} \SUPPORTEDANCH{\char"0188\char"0329} \SUPPORTEDANCH{\char"0064\char"0329} \SUPPORTEDANCH{\char"018C\char"0329} \SUPPORTEDANCH{\char"0221\char"0329} \SUPPORTEDANCH{\char"0068\char"0329} \SUPPORTEDANCH{\char"2C68\char"0329} \SUPPORTEDANCH{\char"006B\char"0329} \SUPPORTEDANCH{\char"0199\char"0329} \SUPPORTEDANCH{\char"2C6A\char"0329} \SUPPORTEDANCH{\char"006C\char"0329} \SUPPORTEDANCH{\char"0234\char"0329} \SUPPORTEDANCH{\char"006E\char"0329} \SUPPORTEDANCH{\char"0235\char"0329} \SUPPORTEDANCH{\char"0070\char"0329} \SUPPORTEDANCH{\char"01A5\char"0329} \SUPPORTEDANCH{\char"0071\char"0329} \SUPPORTEDANCH{\char"02A0\char"0329} \SUPPORTEDANCH{\char"0074\char"0329} \SUPPORTEDANCH{\char"01AB\char"0329} \SUPPORTEDANCH{\char"01AD\char"0329} \SUPPORTEDANCH{\char"0236\char"0329} \SUPPORTEDANCH{\char"0075\char"0329} \MISSINGGLYPH{\char"AB52\char"0329} \SUPPORTEDANCH{\char"0079\char"0329} \SUPPORTEDANCH{\char"01B4\char"0329} \MISSINGGLYPH{\char"1EFF\char"0329}

\SUPPORTEDANCH{\char"0048\char"032A} \SUPPORTEDANCH{\char"2C67\char"032A} \SUPPORTEDANCH{\char"004B\char"032A} \SUPPORTEDANCH{\char"2C69\char"032A} \SUPPORTEDANCH{\char"004D\char"032A} \MISSINGGLYPH{\char"2C6E\char"032A} \SUPPORTEDANCH{\char"0050\char"032A} \SUPPORTEDANCH{\char"01A4\char"032A} \SUPPORTEDANCH{\char"0052\char"032A} \SUPPORTEDANCH{\char"2C64\char"032A} \SUPPORTEDANCH{\char"0054\char"032A} \SUPPORTEDANCH{\char"01AC\char"032A} \SUPPORTEDANCH{\char"01AE\char"032A} \SUPPORTEDANCH{\char"0059\char"032A} \SUPPORTEDANCH{\char"01B3\char"032A} \SUPPORTEDANCH{\char"0061\char"032A} \SUPPORTEDANCH{\char"1E9A\char"032A} \SUPPORTEDANCH{\char"0062\char"032A} \SUPPORTEDANCH{\char"0183\char"032A} \SUPPORTEDANCH{\char"0063\char"032A} \SUPPORTEDANCH{\char"0188\char"032A} \SUPPORTEDANCH{\char"0064\char"032A} \SUPPORTEDANCH{\char"018C\char"032A} \SUPPORTEDANCH{\char"0221\char"032A} \SUPPORTEDANCH{\char"0068\char"032A} \SUPPORTEDANCH{\char"2C68\char"032A} \SUPPORTEDANCH{\char"006B\char"032A} \SUPPORTEDANCH{\char"0199\char"032A} \SUPPORTEDANCH{\char"2C6A\char"032A} \SUPPORTEDANCH{\char"006C\char"032A} \SUPPORTEDANCH{\char"0234\char"032A} \SUPPORTEDANCH{\char"006E\char"032A} \SUPPORTEDANCH{\char"0235\char"032A} \SUPPORTEDANCH{\char"0070\char"032A} \SUPPORTEDANCH{\char"01A5\char"032A} \SUPPORTEDANCH{\char"0071\char"032A} \SUPPORTEDANCH{\char"02A0\char"032A} \SUPPORTEDANCH{\char"0074\char"032A} \SUPPORTEDANCH{\char"01AB\char"032A} \SUPPORTEDANCH{\char"01AD\char"032A} \SUPPORTEDANCH{\char"0236\char"032A} \SUPPORTEDANCH{\char"0075\char"032A} \MISSINGGLYPH{\char"AB52\char"032A} \SUPPORTEDANCH{\char"0079\char"032A} \SUPPORTEDANCH{\char"01B4\char"032A} \MISSINGGLYPH{\char"1EFF\char"032A}

\SUPPORTEDANCH{\char"0048\char"032B} \SUPPORTEDANCH{\char"2C67\char"032B} \SUPPORTEDANCH{\char"004B\char"032B} \SUPPORTEDANCH{\char"2C69\char"032B} \SUPPORTEDANCH{\char"004D\char"032B} \MISSINGGLYPH{\char"2C6E\char"032B} \SUPPORTEDANCH{\char"0050\char"032B} \SUPPORTEDANCH{\char"01A4\char"032B} \SUPPORTEDANCH{\char"0052\char"032B} \SUPPORTEDANCH{\char"2C64\char"032B} \SUPPORTEDANCH{\char"0054\char"032B} \SUPPORTEDANCH{\char"01AC\char"032B} \SUPPORTEDANCH{\char"01AE\char"032B} \SUPPORTEDANCH{\char"0059\char"032B} \SUPPORTEDANCH{\char"01B3\char"032B} \SUPPORTEDANCH{\char"0061\char"032B} \SUPPORTEDANCH{\char"1E9A\char"032B} \SUPPORTEDANCH{\char"0062\char"032B} \SUPPORTEDANCH{\char"0183\char"032B} \SUPPORTEDANCH{\char"0063\char"032B} \SUPPORTEDANCH{\char"0188\char"032B} \SUPPORTEDANCH{\char"0064\char"032B} \SUPPORTEDANCH{\char"018C\char"032B} \SUPPORTEDANCH{\char"0221\char"032B} \SUPPORTEDANCH{\char"0068\char"032B} \SUPPORTEDANCH{\char"2C68\char"032B} \SUPPORTEDANCH{\char"006B\char"032B} \SUPPORTEDANCH{\char"0199\char"032B} \SUPPORTEDANCH{\char"2C6A\char"032B} \SUPPORTEDANCH{\char"006C\char"032B} \SUPPORTEDANCH{\char"0234\char"032B} \SUPPORTEDANCH{\char"006E\char"032B} \SUPPORTEDANCH{\char"0235\char"032B} \SUPPORTEDANCH{\char"0070\char"032B} \SUPPORTEDANCH{\char"01A5\char"032B} \SUPPORTEDANCH{\char"0071\char"032B} \SUPPORTEDANCH{\char"02A0\char"032B} \SUPPORTEDANCH{\char"0074\char"032B} \SUPPORTEDANCH{\char"01AB\char"032B} \SUPPORTEDANCH{\char"01AD\char"032B} \SUPPORTEDANCH{\char"0236\char"032B} \SUPPORTEDANCH{\char"0075\char"032B} \MISSINGGLYPH{\char"AB52\char"032B} \SUPPORTEDANCH{\char"0079\char"032B} \SUPPORTEDANCH{\char"01B4\char"032B} \MISSINGGLYPH{\char"1EFF\char"032B}

\SUPPORTEDANCH{\char"0048\char"032C} \SUPPORTEDANCH{\char"2C67\char"032C} \SUPPORTEDANCH{\char"004B\char"032C} \SUPPORTEDANCH{\char"2C69\char"032C} \SUPPORTEDANCH{\char"004D\char"032C} \MISSINGGLYPH{\char"2C6E\char"032C} \SUPPORTEDANCH{\char"0050\char"032C} \SUPPORTEDANCH{\char"01A4\char"032C} \SUPPORTEDANCH{\char"0052\char"032C} \SUPPORTEDANCH{\char"2C64\char"032C} \SUPPORTEDANCH{\char"0054\char"032C} \SUPPORTEDANCH{\char"01AC\char"032C} \SUPPORTEDANCH{\char"01AE\char"032C} \SUPPORTEDANCH{\char"0059\char"032C} \SUPPORTEDANCH{\char"01B3\char"032C} \SUPPORTEDANCH{\char"0061\char"032C} \SUPPORTEDANCH{\char"1E9A\char"032C} \SUPPORTEDANCH{\char"0062\char"032C} \SUPPORTEDANCH{\char"0183\char"032C} \SUPPORTEDANCH{\char"0063\char"032C} \SUPPORTEDANCH{\char"0188\char"032C} \SUPPORTEDANCH{\char"0064\char"032C} \SUPPORTEDANCH{\char"018C\char"032C} \SUPPORTEDANCH{\char"0221\char"032C} \SUPPORTEDANCH{\char"0068\char"032C} \SUPPORTEDANCH{\char"2C68\char"032C} \SUPPORTEDANCH{\char"006B\char"032C} \SUPPORTEDANCH{\char"0199\char"032C} \SUPPORTEDANCH{\char"2C6A\char"032C} \SUPPORTEDANCH{\char"006C\char"032C} \SUPPORTEDANCH{\char"0234\char"032C} \SUPPORTEDANCH{\char"006E\char"032C} \SUPPORTEDANCH{\char"0235\char"032C} \SUPPORTEDANCH{\char"0070\char"032C} \SUPPORTEDANCH{\char"01A5\char"032C} \SUPPORTEDANCH{\char"0071\char"032C} \SUPPORTEDANCH{\char"02A0\char"032C} \SUPPORTEDANCH{\char"0074\char"032C} \SUPPORTEDANCH{\char"01AB\char"032C} \SUPPORTEDANCH{\char"01AD\char"032C} \SUPPORTEDANCH{\char"0236\char"032C} \SUPPORTEDANCH{\char"0075\char"032C} \MISSINGGLYPH{\char"AB52\char"032C} \SUPPORTEDANCH{\char"0079\char"032C} \SUPPORTEDANCH{\char"01B4\char"032C} \MISSINGGLYPH{\char"1EFF\char"032C}

\SUPPORTEDANCH{\char"0048\char"032D} \SUPPORTEDANCH{\char"2C67\char"032D} \SUPPORTEDANCH{\char"004B\char"032D} \SUPPORTEDANCH{\char"2C69\char"032D} \SUPPORTEDANCH{\char"004D\char"032D} \MISSINGGLYPH{\char"2C6E\char"032D} \SUPPORTEDANCH{\char"0050\char"032D} \SUPPORTEDANCH{\char"01A4\char"032D} \SUPPORTEDANCH{\char"0052\char"032D} \SUPPORTEDANCH{\char"2C64\char"032D} \SUPPORTEDPREC{\char"0054\char"032D} \SUPPORTEDANCH{\char"01AC\char"032D} \SUPPORTEDANCH{\char"01AE\char"032D} \SUPPORTEDANCH{\char"0059\char"032D} \SUPPORTEDANCH{\char"01B3\char"032D} \SUPPORTEDANCH{\char"0061\char"032D} \SUPPORTEDANCH{\char"1E9A\char"032D} \SUPPORTEDANCH{\char"0062\char"032D} \SUPPORTEDANCH{\char"0183\char"032D} \SUPPORTEDANCH{\char"0063\char"032D} \SUPPORTEDANCH{\char"0188\char"032D} \SUPPORTEDPREC{\char"0064\char"032D} \SUPPORTEDANCH{\char"018C\char"032D} \SUPPORTEDANCH{\char"0221\char"032D} \SUPPORTEDANCH{\char"0068\char"032D} \SUPPORTEDANCH{\char"2C68\char"032D} \SUPPORTEDANCH{\char"006B\char"032D} \SUPPORTEDANCH{\char"0199\char"032D} \SUPPORTEDANCH{\char"2C6A\char"032D} \SUPPORTEDPREC{\char"006C\char"032D} \SUPPORTEDANCH{\char"0234\char"032D} \SUPPORTEDPREC{\char"006E\char"032D} \SUPPORTEDANCH{\char"0235\char"032D} \SUPPORTEDANCH{\char"0070\char"032D} \SUPPORTEDANCH{\char"01A5\char"032D} \SUPPORTEDANCH{\char"0071\char"032D} \SUPPORTEDANCH{\char"02A0\char"032D} \SUPPORTEDPREC{\char"0074\char"032D} \SUPPORTEDANCH{\char"01AB\char"032D} \SUPPORTEDANCH{\char"01AD\char"032D} \SUPPORTEDANCH{\char"0236\char"032D} \SUPPORTEDPREC{\char"0075\char"032D} \MISSINGGLYPH{\char"AB52\char"032D} \SUPPORTEDANCH{\char"0079\char"032D} \SUPPORTEDANCH{\char"01B4\char"032D} \MISSINGGLYPH{\char"1EFF\char"032D}

\SUPPORTEDPREC{\char"0048\char"032E} \SUPPORTEDANCH{\char"2C67\char"032E} \SUPPORTEDANCH{\char"004B\char"032E} \SUPPORTEDANCH{\char"2C69\char"032E} \SUPPORTEDANCH{\char"004D\char"032E} \MISSINGGLYPH{\char"2C6E\char"032E} \SUPPORTEDANCH{\char"0050\char"032E} \SUPPORTEDANCH{\char"01A4\char"032E} \SUPPORTEDANCH{\char"0052\char"032E} \SUPPORTEDANCH{\char"2C64\char"032E} \SUPPORTEDANCH{\char"0054\char"032E} \SUPPORTEDANCH{\char"01AC\char"032E} \SUPPORTEDANCH{\char"01AE\char"032E} \SUPPORTEDANCH{\char"0059\char"032E} \SUPPORTEDANCH{\char"01B3\char"032E} \SUPPORTEDANCH{\char"0061\char"032E} \SUPPORTEDANCH{\char"1E9A\char"032E} \SUPPORTEDANCH{\char"0062\char"032E} \SUPPORTEDANCH{\char"0183\char"032E} \SUPPORTEDANCH{\char"0063\char"032E} \SUPPORTEDANCH{\char"0188\char"032E} \SUPPORTEDANCH{\char"0064\char"032E} \SUPPORTEDANCH{\char"018C\char"032E} \SUPPORTEDANCH{\char"0221\char"032E} \SUPPORTEDPREC{\char"0068\char"032E} \SUPPORTEDANCH{\char"2C68\char"032E} \SUPPORTEDANCH{\char"006B\char"032E} \SUPPORTEDANCH{\char"0199\char"032E} \SUPPORTEDANCH{\char"2C6A\char"032E} \SUPPORTEDANCH{\char"006C\char"032E} \SUPPORTEDANCH{\char"0234\char"032E} \SUPPORTEDANCH{\char"006E\char"032E} \SUPPORTEDANCH{\char"0235\char"032E} \SUPPORTEDANCH{\char"0070\char"032E} \SUPPORTEDANCH{\char"01A5\char"032E} \SUPPORTEDANCH{\char"0071\char"032E} \SUPPORTEDANCH{\char"02A0\char"032E} \SUPPORTEDANCH{\char"0074\char"032E} \SUPPORTEDANCH{\char"01AB\char"032E} \SUPPORTEDANCH{\char"01AD\char"032E} \SUPPORTEDANCH{\char"0236\char"032E} \SUPPORTEDANCH{\char"0075\char"032E} \MISSINGGLYPH{\char"AB52\char"032E} \SUPPORTEDANCH{\char"0079\char"032E} \SUPPORTEDANCH{\char"01B4\char"032E} \MISSINGGLYPH{\char"1EFF\char"032E}

\SUPPORTEDANCH{\char"0048\char"032F} \SUPPORTEDANCH{\char"2C67\char"032F} \SUPPORTEDANCH{\char"004B\char"032F} \SUPPORTEDANCH{\char"2C69\char"032F} \SUPPORTEDANCH{\char"004D\char"032F} \MISSINGGLYPH{\char"2C6E\char"032F} \SUPPORTEDANCH{\char"0050\char"032F} \SUPPORTEDANCH{\char"01A4\char"032F} \SUPPORTEDANCH{\char"0052\char"032F} \SUPPORTEDANCH{\char"2C64\char"032F} \SUPPORTEDANCH{\char"0054\char"032F} \SUPPORTEDANCH{\char"01AC\char"032F} \SUPPORTEDANCH{\char"01AE\char"032F} \SUPPORTEDANCH{\char"0059\char"032F} \SUPPORTEDANCH{\char"01B3\char"032F} \SUPPORTEDANCH{\char"0061\char"032F} \SUPPORTEDANCH{\char"1E9A\char"032F} \SUPPORTEDANCH{\char"0062\char"032F} \SUPPORTEDANCH{\char"0183\char"032F} \SUPPORTEDANCH{\char"0063\char"032F} \SUPPORTEDANCH{\char"0188\char"032F} \SUPPORTEDANCH{\char"0064\char"032F} \SUPPORTEDANCH{\char"018C\char"032F} \SUPPORTEDANCH{\char"0221\char"032F} \SUPPORTEDANCH{\char"0068\char"032F} \SUPPORTEDANCH{\char"2C68\char"032F} \SUPPORTEDANCH{\char"006B\char"032F} \SUPPORTEDANCH{\char"0199\char"032F} \SUPPORTEDANCH{\char"2C6A\char"032F} \SUPPORTEDANCH{\char"006C\char"032F} \SUPPORTEDANCH{\char"0234\char"032F} \SUPPORTEDANCH{\char"006E\char"032F} \SUPPORTEDANCH{\char"0235\char"032F} \SUPPORTEDANCH{\char"0070\char"032F} \SUPPORTEDANCH{\char"01A5\char"032F} \SUPPORTEDANCH{\char"0071\char"032F} \SUPPORTEDANCH{\char"02A0\char"032F} \SUPPORTEDANCH{\char"0074\char"032F} \SUPPORTEDANCH{\char"01AB\char"032F} \SUPPORTEDANCH{\char"01AD\char"032F} \SUPPORTEDANCH{\char"0236\char"032F} \SUPPORTEDANCH{\char"0075\char"032F} \MISSINGGLYPH{\char"AB52\char"032F} \SUPPORTEDANCH{\char"0079\char"032F} \SUPPORTEDANCH{\char"01B4\char"032F} \MISSINGGLYPH{\char"1EFF\char"032F}

\SUPPORTEDANCH{\char"0048\char"031D} \SUPPORTEDANCH{\char"2C67\char"031D} \SUPPORTEDANCH{\char"004B\char"031D} \SUPPORTEDANCH{\char"2C69\char"031D} \SUPPORTEDANCH{\char"004D\char"031D} \MISSINGGLYPH{\char"2C6E\char"031D} \SUPPORTEDANCH{\char"0050\char"031D} \SUPPORTEDANCH{\char"01A4\char"031D} \SUPPORTEDANCH{\char"0052\char"031D} \SUPPORTEDANCH{\char"2C64\char"031D} \SUPPORTEDANCH{\char"0054\char"031D} \SUPPORTEDANCH{\char"01AC\char"031D} \SUPPORTEDANCH{\char"01AE\char"031D} \SUPPORTEDANCH{\char"0059\char"031D} \SUPPORTEDANCH{\char"01B3\char"031D} \SUPPORTEDANCH{\char"0061\char"031D} \SUPPORTEDANCH{\char"1E9A\char"031D} \SUPPORTEDANCH{\char"0062\char"031D} \SUPPORTEDANCH{\char"0183\char"031D} \SUPPORTEDANCH{\char"0063\char"031D} \SUPPORTEDANCH{\char"0188\char"031D} \SUPPORTEDANCH{\char"0064\char"031D} \SUPPORTEDANCH{\char"018C\char"031D} \SUPPORTEDANCH{\char"0221\char"031D} \SUPPORTEDANCH{\char"0068\char"031D} \SUPPORTEDANCH{\char"2C68\char"031D} \SUPPORTEDANCH{\char"006B\char"031D} \SUPPORTEDANCH{\char"0199\char"031D} \SUPPORTEDANCH{\char"2C6A\char"031D} \SUPPORTEDANCH{\char"006C\char"031D} \SUPPORTEDANCH{\char"0234\char"031D} \SUPPORTEDANCH{\char"006E\char"031D} \SUPPORTEDANCH{\char"0235\char"031D} \SUPPORTEDANCH{\char"0070\char"031D} \SUPPORTEDANCH{\char"01A5\char"031D} \SUPPORTEDANCH{\char"0071\char"031D} \SUPPORTEDANCH{\char"02A0\char"031D} \SUPPORTEDANCH{\char"0074\char"031D} \SUPPORTEDANCH{\char"01AB\char"031D} \SUPPORTEDANCH{\char"01AD\char"031D} \SUPPORTEDANCH{\char"0236\char"031D} \SUPPORTEDANCH{\char"0075\char"031D} \MISSINGGLYPH{\char"AB52\char"031D} \SUPPORTEDANCH{\char"0079\char"031D} \SUPPORTEDANCH{\char"01B4\char"031D} \MISSINGGLYPH{\char"1EFF\char"031D}

\SUPPORTEDANCH{\char"0048\char"031E} \SUPPORTEDANCH{\char"2C67\char"031E} \SUPPORTEDANCH{\char"004B\char"031E} \SUPPORTEDANCH{\char"2C69\char"031E} \SUPPORTEDANCH{\char"004D\char"031E} \MISSINGGLYPH{\char"2C6E\char"031E} \SUPPORTEDANCH{\char"0050\char"031E} \SUPPORTEDANCH{\char"01A4\char"031E} \SUPPORTEDANCH{\char"0052\char"031E} \SUPPORTEDANCH{\char"2C64\char"031E} \SUPPORTEDANCH{\char"0054\char"031E} \SUPPORTEDANCH{\char"01AC\char"031E} \SUPPORTEDANCH{\char"01AE\char"031E} \SUPPORTEDANCH{\char"0059\char"031E} \SUPPORTEDANCH{\char"01B3\char"031E} \SUPPORTEDANCH{\char"0061\char"031E} \SUPPORTEDANCH{\char"1E9A\char"031E} \SUPPORTEDANCH{\char"0062\char"031E} \SUPPORTEDANCH{\char"0183\char"031E} \SUPPORTEDANCH{\char"0063\char"031E} \SUPPORTEDANCH{\char"0188\char"031E} \SUPPORTEDANCH{\char"0064\char"031E} \SUPPORTEDANCH{\char"018C\char"031E} \SUPPORTEDANCH{\char"0221\char"031E} \SUPPORTEDANCH{\char"0068\char"031E} \SUPPORTEDANCH{\char"2C68\char"031E} \SUPPORTEDANCH{\char"006B\char"031E} \SUPPORTEDANCH{\char"0199\char"031E} \SUPPORTEDANCH{\char"2C6A\char"031E} \SUPPORTEDANCH{\char"006C\char"031E} \SUPPORTEDANCH{\char"0234\char"031E} \SUPPORTEDANCH{\char"006E\char"031E} \SUPPORTEDANCH{\char"0235\char"031E} \SUPPORTEDANCH{\char"0070\char"031E} \SUPPORTEDANCH{\char"01A5\char"031E} \SUPPORTEDANCH{\char"0071\char"031E} \SUPPORTEDANCH{\char"02A0\char"031E} \SUPPORTEDANCH{\char"0074\char"031E} \SUPPORTEDANCH{\char"01AB\char"031E} \SUPPORTEDANCH{\char"01AD\char"031E} \SUPPORTEDANCH{\char"0236\char"031E} \SUPPORTEDANCH{\char"0075\char"031E} \MISSINGGLYPH{\char"AB52\char"031E} \SUPPORTEDANCH{\char"0079\char"031E} \SUPPORTEDANCH{\char"01B4\char"031E} \MISSINGGLYPH{\char"1EFF\char"031E}

\SUPPORTEDANCH{\char"0048\char"031F} \SUPPORTEDANCH{\char"2C67\char"031F} \SUPPORTEDANCH{\char"004B\char"031F} \SUPPORTEDANCH{\char"2C69\char"031F} \SUPPORTEDANCH{\char"004D\char"031F} \MISSINGGLYPH{\char"2C6E\char"031F} \SUPPORTEDANCH{\char"0050\char"031F} \SUPPORTEDANCH{\char"01A4\char"031F} \SUPPORTEDANCH{\char"0052\char"031F} \SUPPORTEDANCH{\char"2C64\char"031F} \SUPPORTEDANCH{\char"0054\char"031F} \SUPPORTEDANCH{\char"01AC\char"031F} \SUPPORTEDANCH{\char"01AE\char"031F} \SUPPORTEDANCH{\char"0059\char"031F} \SUPPORTEDANCH{\char"01B3\char"031F} \SUPPORTEDANCH{\char"0061\char"031F} \SUPPORTEDANCH{\char"1E9A\char"031F} \SUPPORTEDANCH{\char"0062\char"031F} \SUPPORTEDANCH{\char"0183\char"031F} \SUPPORTEDANCH{\char"0063\char"031F} \SUPPORTEDANCH{\char"0188\char"031F} \SUPPORTEDANCH{\char"0064\char"031F} \SUPPORTEDANCH{\char"018C\char"031F} \SUPPORTEDANCH{\char"0221\char"031F} \SUPPORTEDANCH{\char"0068\char"031F} \SUPPORTEDANCH{\char"2C68\char"031F} \SUPPORTEDANCH{\char"006B\char"031F} \SUPPORTEDANCH{\char"0199\char"031F} \SUPPORTEDANCH{\char"2C6A\char"031F} \SUPPORTEDANCH{\char"006C\char"031F} \SUPPORTEDANCH{\char"0234\char"031F} \SUPPORTEDANCH{\char"006E\char"031F} \SUPPORTEDANCH{\char"0235\char"031F} \SUPPORTEDANCH{\char"0070\char"031F} \SUPPORTEDANCH{\char"01A5\char"031F} \SUPPORTEDANCH{\char"0071\char"031F} \SUPPORTEDANCH{\char"02A0\char"031F} \SUPPORTEDANCH{\char"0074\char"031F} \SUPPORTEDANCH{\char"01AB\char"031F} \SUPPORTEDANCH{\char"01AD\char"031F} \SUPPORTEDANCH{\char"0236\char"031F} \SUPPORTEDANCH{\char"0075\char"031F} \MISSINGGLYPH{\char"AB52\char"031F} \SUPPORTEDANCH{\char"0079\char"031F} \SUPPORTEDANCH{\char"01B4\char"031F} \MISSINGGLYPH{\char"1EFF\char"031F}

\SUPPORTEDANCH{\char"0048\char"0330} \SUPPORTEDANCH{\char"2C67\char"0330} \SUPPORTEDANCH{\char"004B\char"0330} \SUPPORTEDANCH{\char"2C69\char"0330} \SUPPORTEDANCH{\char"004D\char"0330} \MISSINGGLYPH{\char"2C6E\char"0330} \SUPPORTEDANCH{\char"0050\char"0330} \SUPPORTEDANCH{\char"01A4\char"0330} \SUPPORTEDANCH{\char"0052\char"0330} \SUPPORTEDANCH{\char"2C64\char"0330} \SUPPORTEDANCH{\char"0054\char"0330} \SUPPORTEDANCH{\char"01AC\char"0330} \SUPPORTEDANCH{\char"01AE\char"0330} \SUPPORTEDANCH{\char"0059\char"0330} \SUPPORTEDANCH{\char"01B3\char"0330} \SUPPORTEDANCH{\char"0061\char"0330} \SUPPORTEDANCH{\char"1E9A\char"0330} \SUPPORTEDANCH{\char"0062\char"0330} \SUPPORTEDANCH{\char"0183\char"0330} \SUPPORTEDANCH{\char"0063\char"0330} \SUPPORTEDANCH{\char"0188\char"0330} \SUPPORTEDANCH{\char"0064\char"0330} \SUPPORTEDANCH{\char"018C\char"0330} \SUPPORTEDANCH{\char"0221\char"0330} \SUPPORTEDANCH{\char"0068\char"0330} \SUPPORTEDANCH{\char"2C68\char"0330} \SUPPORTEDANCH{\char"006B\char"0330} \SUPPORTEDANCH{\char"0199\char"0330} \SUPPORTEDANCH{\char"2C6A\char"0330} \SUPPORTEDANCH{\char"006C\char"0330} \SUPPORTEDANCH{\char"0234\char"0330} \SUPPORTEDANCH{\char"006E\char"0330} \SUPPORTEDANCH{\char"0235\char"0330} \SUPPORTEDANCH{\char"0070\char"0330} \SUPPORTEDANCH{\char"01A5\char"0330} \SUPPORTEDANCH{\char"0071\char"0330} \SUPPORTEDANCH{\char"02A0\char"0330} \SUPPORTEDANCH{\char"0074\char"0330} \SUPPORTEDANCH{\char"01AB\char"0330} \SUPPORTEDANCH{\char"01AD\char"0330} \SUPPORTEDANCH{\char"0236\char"0330} \SUPPORTEDPREC{\char"0075\char"0330} \MISSINGGLYPH{\char"AB52\char"0330} \SUPPORTEDANCH{\char"0079\char"0330} \SUPPORTEDANCH{\char"01B4\char"0330} \MISSINGGLYPH{\char"1EFF\char"0330}

\SUPPORTEDANCH{\char"0048\char"0331} \SUPPORTEDANCH{\char"2C67\char"0331} \SUPPORTEDPREC{\char"004B\char"0331} \SUPPORTEDANCH{\char"2C69\char"0331} \SUPPORTEDANCH{\char"004D\char"0331} \MISSINGGLYPH{\char"2C6E\char"0331} \SUPPORTEDANCH{\char"0050\char"0331} \SUPPORTEDANCH{\char"01A4\char"0331} \SUPPORTEDPREC{\char"0052\char"0331} \SUPPORTEDANCH{\char"2C64\char"0331} \SUPPORTEDPREC{\char"0054\char"0331} \SUPPORTEDANCH{\char"01AC\char"0331} \SUPPORTEDANCH{\char"01AE\char"0331} \SUPPORTEDANCH{\char"0059\char"0331} \SUPPORTEDANCH{\char"01B3\char"0331} \SUPPORTEDANCH{\char"0061\char"0331} \SUPPORTEDANCH{\char"1E9A\char"0331} \SUPPORTEDPREC{\char"0062\char"0331} \SUPPORTEDANCH{\char"0183\char"0331} \SUPPORTEDANCH{\char"0063\char"0331} \SUPPORTEDANCH{\char"0188\char"0331} \SUPPORTEDPREC{\char"0064\char"0331} \SUPPORTEDANCH{\char"018C\char"0331} \SUPPORTEDANCH{\char"0221\char"0331} \SUPPORTEDPREC{\char"0068\char"0331} \SUPPORTEDANCH{\char"2C68\char"0331} \SUPPORTEDPREC{\char"006B\char"0331} \SUPPORTEDANCH{\char"0199\char"0331} \SUPPORTEDANCH{\char"2C6A\char"0331} \SUPPORTEDPREC{\char"006C\char"0331} \SUPPORTEDANCH{\char"0234\char"0331} \SUPPORTEDPREC{\char"006E\char"0331} \SUPPORTEDANCH{\char"0235\char"0331} \SUPPORTEDANCH{\char"0070\char"0331} \SUPPORTEDANCH{\char"01A5\char"0331} \SUPPORTEDANCH{\char"0071\char"0331} \SUPPORTEDANCH{\char"02A0\char"0331} \SUPPORTEDPREC{\char"0074\char"0331} \SUPPORTEDANCH{\char"01AB\char"0331} \SUPPORTEDANCH{\char"01AD\char"0331} \SUPPORTEDANCH{\char"0236\char"0331} \SUPPORTEDANCH{\char"0075\char"0331} \MISSINGGLYPH{\char"AB52\char"0331} \SUPPORTEDANCH{\char"0079\char"0331} \SUPPORTEDANCH{\char"01B4\char"0331} \MISSINGGLYPH{\char"1EFF\char"0331}

\SUPPORTEDANCH{\char"0048\char"0333} \SUPPORTEDANCH{\char"2C67\char"0333} \SUPPORTEDANCH{\char"004B\char"0333} \SUPPORTEDANCH{\char"2C69\char"0333} \SUPPORTEDANCH{\char"004D\char"0333} \MISSINGGLYPH{\char"2C6E\char"0333} \SUPPORTEDANCH{\char"0050\char"0333} \SUPPORTEDANCH{\char"01A4\char"0333} \SUPPORTEDANCH{\char"0052\char"0333} \SUPPORTEDANCH{\char"2C64\char"0333} \SUPPORTEDANCH{\char"0054\char"0333} \SUPPORTEDANCH{\char"01AC\char"0333} \SUPPORTEDANCH{\char"01AE\char"0333} \SUPPORTEDANCH{\char"0059\char"0333} \SUPPORTEDANCH{\char"01B3\char"0333} \SUPPORTEDANCH{\char"0061\char"0333} \SUPPORTEDANCH{\char"1E9A\char"0333} \SUPPORTEDANCH{\char"0062\char"0333} \SUPPORTEDANCH{\char"0183\char"0333} \SUPPORTEDANCH{\char"0063\char"0333} \SUPPORTEDANCH{\char"0188\char"0333} \SUPPORTEDANCH{\char"0064\char"0333} \SUPPORTEDANCH{\char"018C\char"0333} \SUPPORTEDANCH{\char"0221\char"0333} \SUPPORTEDANCH{\char"0068\char"0333} \SUPPORTEDANCH{\char"2C68\char"0333} \SUPPORTEDANCH{\char"006B\char"0333} \SUPPORTEDANCH{\char"0199\char"0333} \SUPPORTEDANCH{\char"2C6A\char"0333} \SUPPORTEDANCH{\char"006C\char"0333} \SUPPORTEDANCH{\char"0234\char"0333} \SUPPORTEDANCH{\char"006E\char"0333} \SUPPORTEDANCH{\char"0235\char"0333} \SUPPORTEDANCH{\char"0070\char"0333} \SUPPORTEDANCH{\char"01A5\char"0333} \SUPPORTEDANCH{\char"0071\char"0333} \SUPPORTEDANCH{\char"02A0\char"0333} \SUPPORTEDANCH{\char"0074\char"0333} \SUPPORTEDANCH{\char"01AB\char"0333} \SUPPORTEDANCH{\char"01AD\char"0333} \SUPPORTEDANCH{\char"0236\char"0333} \SUPPORTEDANCH{\char"0075\char"0333} \MISSINGGLYPH{\char"AB52\char"0333} \SUPPORTEDANCH{\char"0079\char"0333} \SUPPORTEDANCH{\char"01B4\char"0333} \MISSINGGLYPH{\char"1EFF\char"0333}

\SUPPORTEDANCH{\char"0048\char"0339} \SUPPORTEDANCH{\char"2C67\char"0339} \SUPPORTEDANCH{\char"004B\char"0339} \SUPPORTEDANCH{\char"2C69\char"0339} \SUPPORTEDANCH{\char"004D\char"0339} \MISSINGGLYPH{\char"2C6E\char"0339} \SUPPORTEDANCH{\char"0050\char"0339} \SUPPORTEDANCH{\char"01A4\char"0339} \SUPPORTEDANCH{\char"0052\char"0339} \SUPPORTEDANCH{\char"2C64\char"0339} \SUPPORTEDANCH{\char"0054\char"0339} \SUPPORTEDANCH{\char"01AC\char"0339} \SUPPORTEDANCH{\char"01AE\char"0339} \SUPPORTEDANCH{\char"0059\char"0339} \SUPPORTEDANCH{\char"01B3\char"0339} \SUPPORTEDANCH{\char"0061\char"0339} \SUPPORTEDANCH{\char"1E9A\char"0339} \SUPPORTEDANCH{\char"0062\char"0339} \SUPPORTEDANCH{\char"0183\char"0339} \SUPPORTEDANCH{\char"0063\char"0339} \SUPPORTEDANCH{\char"0188\char"0339} \SUPPORTEDANCH{\char"0064\char"0339} \SUPPORTEDANCH{\char"018C\char"0339} \SUPPORTEDANCH{\char"0221\char"0339} \SUPPORTEDANCH{\char"0068\char"0339} \SUPPORTEDANCH{\char"2C68\char"0339} \SUPPORTEDANCH{\char"006B\char"0339} \SUPPORTEDANCH{\char"0199\char"0339} \SUPPORTEDANCH{\char"2C6A\char"0339} \SUPPORTEDANCH{\char"006C\char"0339} \SUPPORTEDANCH{\char"0234\char"0339} \SUPPORTEDANCH{\char"006E\char"0339} \SUPPORTEDANCH{\char"0235\char"0339} \SUPPORTEDANCH{\char"0070\char"0339} \SUPPORTEDANCH{\char"01A5\char"0339} \SUPPORTEDANCH{\char"0071\char"0339} \SUPPORTEDANCH{\char"02A0\char"0339} \SUPPORTEDANCH{\char"0074\char"0339} \SUPPORTEDANCH{\char"01AB\char"0339} \SUPPORTEDANCH{\char"01AD\char"0339} \SUPPORTEDANCH{\char"0236\char"0339} \SUPPORTEDANCH{\char"0075\char"0339} \MISSINGGLYPH{\char"AB52\char"0339} \SUPPORTEDANCH{\char"0079\char"0339} \SUPPORTEDANCH{\char"01B4\char"0339} \MISSINGGLYPH{\char"1EFF\char"0339}

\SUPPORTEDANCH{\char"0048\char"033A} \SUPPORTEDANCH{\char"2C67\char"033A} \SUPPORTEDANCH{\char"004B\char"033A} \SUPPORTEDANCH{\char"2C69\char"033A} \SUPPORTEDANCH{\char"004D\char"033A} \MISSINGGLYPH{\char"2C6E\char"033A} \SUPPORTEDANCH{\char"0050\char"033A} \SUPPORTEDANCH{\char"01A4\char"033A} \SUPPORTEDANCH{\char"0052\char"033A} \SUPPORTEDANCH{\char"2C64\char"033A} \SUPPORTEDANCH{\char"0054\char"033A} \SUPPORTEDANCH{\char"01AC\char"033A} \SUPPORTEDANCH{\char"01AE\char"033A} \SUPPORTEDANCH{\char"0059\char"033A} \SUPPORTEDANCH{\char"01B3\char"033A} \SUPPORTEDANCH{\char"0061\char"033A} \SUPPORTEDANCH{\char"1E9A\char"033A} \SUPPORTEDANCH{\char"0062\char"033A} \SUPPORTEDANCH{\char"0183\char"033A} \SUPPORTEDANCH{\char"0063\char"033A} \SUPPORTEDANCH{\char"0188\char"033A} \SUPPORTEDANCH{\char"0064\char"033A} \SUPPORTEDANCH{\char"018C\char"033A} \SUPPORTEDANCH{\char"0221\char"033A} \SUPPORTEDANCH{\char"0068\char"033A} \SUPPORTEDANCH{\char"2C68\char"033A} \SUPPORTEDANCH{\char"006B\char"033A} \SUPPORTEDANCH{\char"0199\char"033A} \SUPPORTEDANCH{\char"2C6A\char"033A} \SUPPORTEDANCH{\char"006C\char"033A} \SUPPORTEDANCH{\char"0234\char"033A} \SUPPORTEDANCH{\char"006E\char"033A} \SUPPORTEDANCH{\char"0235\char"033A} \SUPPORTEDANCH{\char"0070\char"033A} \SUPPORTEDANCH{\char"01A5\char"033A} \SUPPORTEDANCH{\char"0071\char"033A} \SUPPORTEDANCH{\char"02A0\char"033A} \SUPPORTEDANCH{\char"0074\char"033A} \SUPPORTEDANCH{\char"01AB\char"033A} \SUPPORTEDANCH{\char"01AD\char"033A} \SUPPORTEDANCH{\char"0236\char"033A} \SUPPORTEDANCH{\char"0075\char"033A} \MISSINGGLYPH{\char"AB52\char"033A} \SUPPORTEDANCH{\char"0079\char"033A} \SUPPORTEDANCH{\char"01B4\char"033A} \MISSINGGLYPH{\char"1EFF\char"033A}

\SUPPORTEDANCH{\char"0048\char"033B} \SUPPORTEDANCH{\char"2C67\char"033B} \SUPPORTEDANCH{\char"004B\char"033B} \SUPPORTEDANCH{\char"2C69\char"033B} \SUPPORTEDANCH{\char"004D\char"033B} \MISSINGGLYPH{\char"2C6E\char"033B} \SUPPORTEDANCH{\char"0050\char"033B} \SUPPORTEDANCH{\char"01A4\char"033B} \SUPPORTEDANCH{\char"0052\char"033B} \SUPPORTEDANCH{\char"2C64\char"033B} \SUPPORTEDANCH{\char"0054\char"033B} \SUPPORTEDANCH{\char"01AC\char"033B} \SUPPORTEDANCH{\char"01AE\char"033B} \SUPPORTEDANCH{\char"0059\char"033B} \SUPPORTEDANCH{\char"01B3\char"033B} \SUPPORTEDANCH{\char"0061\char"033B} \SUPPORTEDANCH{\char"1E9A\char"033B} \SUPPORTEDANCH{\char"0062\char"033B} \SUPPORTEDANCH{\char"0183\char"033B} \SUPPORTEDANCH{\char"0063\char"033B} \SUPPORTEDANCH{\char"0188\char"033B} \SUPPORTEDANCH{\char"0064\char"033B} \SUPPORTEDANCH{\char"018C\char"033B} \SUPPORTEDANCH{\char"0221\char"033B} \SUPPORTEDANCH{\char"0068\char"033B} \SUPPORTEDANCH{\char"2C68\char"033B} \SUPPORTEDANCH{\char"006B\char"033B} \SUPPORTEDANCH{\char"0199\char"033B} \SUPPORTEDANCH{\char"2C6A\char"033B} \SUPPORTEDANCH{\char"006C\char"033B} \SUPPORTEDANCH{\char"0234\char"033B} \SUPPORTEDANCH{\char"006E\char"033B} \SUPPORTEDANCH{\char"0235\char"033B} \SUPPORTEDANCH{\char"0070\char"033B} \SUPPORTEDANCH{\char"01A5\char"033B} \SUPPORTEDANCH{\char"0071\char"033B} \SUPPORTEDANCH{\char"02A0\char"033B} \SUPPORTEDANCH{\char"0074\char"033B} \SUPPORTEDANCH{\char"01AB\char"033B} \SUPPORTEDANCH{\char"01AD\char"033B} \SUPPORTEDANCH{\char"0236\char"033B} \SUPPORTEDANCH{\char"0075\char"033B} \MISSINGGLYPH{\char"AB52\char"033B} \SUPPORTEDANCH{\char"0079\char"033B} \SUPPORTEDANCH{\char"01B4\char"033B} \MISSINGGLYPH{\char"1EFF\char"033B}

\SUPPORTEDANCH{\char"0048\char"033C} \SUPPORTEDANCH{\char"2C67\char"033C} \SUPPORTEDANCH{\char"004B\char"033C} \SUPPORTEDANCH{\char"2C69\char"033C} \SUPPORTEDANCH{\char"004D\char"033C} \MISSINGGLYPH{\char"2C6E\char"033C} \SUPPORTEDANCH{\char"0050\char"033C} \SUPPORTEDANCH{\char"01A4\char"033C} \SUPPORTEDANCH{\char"0052\char"033C} \SUPPORTEDANCH{\char"2C64\char"033C} \SUPPORTEDANCH{\char"0054\char"033C} \SUPPORTEDANCH{\char"01AC\char"033C} \SUPPORTEDANCH{\char"01AE\char"033C} \SUPPORTEDANCH{\char"0059\char"033C} \SUPPORTEDANCH{\char"01B3\char"033C} \SUPPORTEDANCH{\char"0061\char"033C} \SUPPORTEDANCH{\char"1E9A\char"033C} \SUPPORTEDANCH{\char"0062\char"033C} \SUPPORTEDANCH{\char"0183\char"033C} \SUPPORTEDANCH{\char"0063\char"033C} \SUPPORTEDANCH{\char"0188\char"033C} \SUPPORTEDANCH{\char"0064\char"033C} \SUPPORTEDANCH{\char"018C\char"033C} \SUPPORTEDANCH{\char"0221\char"033C} \SUPPORTEDANCH{\char"0068\char"033C} \SUPPORTEDANCH{\char"2C68\char"033C} \SUPPORTEDANCH{\char"006B\char"033C} \SUPPORTEDANCH{\char"0199\char"033C} \SUPPORTEDANCH{\char"2C6A\char"033C} \SUPPORTEDANCH{\char"006C\char"033C} \SUPPORTEDANCH{\char"0234\char"033C} \SUPPORTEDANCH{\char"006E\char"033C} \SUPPORTEDANCH{\char"0235\char"033C} \SUPPORTEDANCH{\char"0070\char"033C} \SUPPORTEDANCH{\char"01A5\char"033C} \SUPPORTEDANCH{\char"0071\char"033C} \SUPPORTEDANCH{\char"02A0\char"033C} \SUPPORTEDANCH{\char"0074\char"033C} \SUPPORTEDANCH{\char"01AB\char"033C} \SUPPORTEDANCH{\char"01AD\char"033C} \SUPPORTEDANCH{\char"0236\char"033C} \SUPPORTEDANCH{\char"0075\char"033C} \MISSINGGLYPH{\char"AB52\char"033C} \SUPPORTEDANCH{\char"0079\char"033C} \SUPPORTEDANCH{\char"01B4\char"033C} \MISSINGGLYPH{\char"1EFF\char"033C}

\SUPPORTEDANCH{\char"0048\char"0347} \SUPPORTEDANCH{\char"2C67\char"0347} \SUPPORTEDANCH{\char"004B\char"0347} \SUPPORTEDANCH{\char"2C69\char"0347} \SUPPORTEDANCH{\char"004D\char"0347} \MISSINGGLYPH{\char"2C6E\char"0347} \SUPPORTEDANCH{\char"0050\char"0347} \SUPPORTEDANCH{\char"01A4\char"0347} \SUPPORTEDANCH{\char"0052\char"0347} \SUPPORTEDANCH{\char"2C64\char"0347} \SUPPORTEDANCH{\char"0054\char"0347} \SUPPORTEDANCH{\char"01AC\char"0347} \SUPPORTEDANCH{\char"01AE\char"0347} \SUPPORTEDANCH{\char"0059\char"0347} \SUPPORTEDANCH{\char"01B3\char"0347} \SUPPORTEDANCH{\char"0061\char"0347} \SUPPORTEDANCH{\char"1E9A\char"0347} \SUPPORTEDANCH{\char"0062\char"0347} \SUPPORTEDANCH{\char"0183\char"0347} \SUPPORTEDANCH{\char"0063\char"0347} \SUPPORTEDANCH{\char"0188\char"0347} \SUPPORTEDANCH{\char"0064\char"0347} \SUPPORTEDANCH{\char"018C\char"0347} \SUPPORTEDANCH{\char"0221\char"0347} \SUPPORTEDANCH{\char"0068\char"0347} \SUPPORTEDANCH{\char"2C68\char"0347} \SUPPORTEDANCH{\char"006B\char"0347} \SUPPORTEDANCH{\char"0199\char"0347} \SUPPORTEDANCH{\char"2C6A\char"0347} \SUPPORTEDANCH{\char"006C\char"0347} \SUPPORTEDANCH{\char"0234\char"0347} \SUPPORTEDANCH{\char"006E\char"0347} \SUPPORTEDANCH{\char"0235\char"0347} \SUPPORTEDANCH{\char"0070\char"0347} \SUPPORTEDANCH{\char"01A5\char"0347} \SUPPORTEDANCH{\char"0071\char"0347} \SUPPORTEDANCH{\char"02A0\char"0347} \SUPPORTEDANCH{\char"0074\char"0347} \SUPPORTEDANCH{\char"01AB\char"0347} \SUPPORTEDANCH{\char"01AD\char"0347} \SUPPORTEDANCH{\char"0236\char"0347} \SUPPORTEDANCH{\char"0075\char"0347} \MISSINGGLYPH{\char"AB52\char"0347} \SUPPORTEDANCH{\char"0079\char"0347} \SUPPORTEDANCH{\char"01B4\char"0347} \MISSINGGLYPH{\char"1EFF\char"0347}

\SUPPORTEDANCH{\char"0048\char"0348} \SUPPORTEDANCH{\char"2C67\char"0348} \SUPPORTEDANCH{\char"004B\char"0348} \SUPPORTEDANCH{\char"2C69\char"0348} \SUPPORTEDANCH{\char"004D\char"0348} \MISSINGGLYPH{\char"2C6E\char"0348} \SUPPORTEDANCH{\char"0050\char"0348} \SUPPORTEDANCH{\char"01A4\char"0348} \SUPPORTEDANCH{\char"0052\char"0348} \SUPPORTEDANCH{\char"2C64\char"0348} \SUPPORTEDANCH{\char"0054\char"0348} \SUPPORTEDANCH{\char"01AC\char"0348} \SUPPORTEDANCH{\char"01AE\char"0348} \SUPPORTEDANCH{\char"0059\char"0348} \SUPPORTEDANCH{\char"01B3\char"0348} \SUPPORTEDANCH{\char"0061\char"0348} \SUPPORTEDANCH{\char"1E9A\char"0348} \SUPPORTEDANCH{\char"0062\char"0348} \SUPPORTEDANCH{\char"0183\char"0348} \SUPPORTEDANCH{\char"0063\char"0348} \SUPPORTEDANCH{\char"0188\char"0348} \SUPPORTEDANCH{\char"0064\char"0348} \SUPPORTEDANCH{\char"018C\char"0348} \SUPPORTEDANCH{\char"0221\char"0348} \SUPPORTEDANCH{\char"0068\char"0348} \SUPPORTEDANCH{\char"2C68\char"0348} \SUPPORTEDANCH{\char"006B\char"0348} \SUPPORTEDANCH{\char"0199\char"0348} \SUPPORTEDANCH{\char"2C6A\char"0348} \SUPPORTEDANCH{\char"006C\char"0348} \SUPPORTEDANCH{\char"0234\char"0348} \SUPPORTEDANCH{\char"006E\char"0348} \SUPPORTEDANCH{\char"0235\char"0348} \SUPPORTEDANCH{\char"0070\char"0348} \SUPPORTEDANCH{\char"01A5\char"0348} \SUPPORTEDANCH{\char"0071\char"0348} \SUPPORTEDANCH{\char"02A0\char"0348} \SUPPORTEDANCH{\char"0074\char"0348} \SUPPORTEDANCH{\char"01AB\char"0348} \SUPPORTEDANCH{\char"01AD\char"0348} \SUPPORTEDANCH{\char"0236\char"0348} \SUPPORTEDANCH{\char"0075\char"0348} \MISSINGGLYPH{\char"AB52\char"0348} \SUPPORTEDANCH{\char"0079\char"0348} \SUPPORTEDANCH{\char"01B4\char"0348} \MISSINGGLYPH{\char"1EFF\char"0348}

\SUPPORTEDANCH{\char"0048\char"0349} \SUPPORTEDANCH{\char"2C67\char"0349} \SUPPORTEDANCH{\char"004B\char"0349} \SUPPORTEDANCH{\char"2C69\char"0349} \SUPPORTEDANCH{\char"004D\char"0349} \MISSINGGLYPH{\char"2C6E\char"0349} \SUPPORTEDANCH{\char"0050\char"0349} \SUPPORTEDANCH{\char"01A4\char"0349} \SUPPORTEDANCH{\char"0052\char"0349} \SUPPORTEDANCH{\char"2C64\char"0349} \SUPPORTEDANCH{\char"0054\char"0349} \SUPPORTEDANCH{\char"01AC\char"0349} \SUPPORTEDANCH{\char"01AE\char"0349} \SUPPORTEDANCH{\char"0059\char"0349} \SUPPORTEDANCH{\char"01B3\char"0349} \SUPPORTEDANCH{\char"0061\char"0349} \SUPPORTEDANCH{\char"1E9A\char"0349} \SUPPORTEDANCH{\char"0062\char"0349} \SUPPORTEDANCH{\char"0183\char"0349} \SUPPORTEDANCH{\char"0063\char"0349} \SUPPORTEDANCH{\char"0188\char"0349} \SUPPORTEDANCH{\char"0064\char"0349} \SUPPORTEDANCH{\char"018C\char"0349} \SUPPORTEDANCH{\char"0221\char"0349} \SUPPORTEDANCH{\char"0068\char"0349} \SUPPORTEDANCH{\char"2C68\char"0349} \SUPPORTEDANCH{\char"006B\char"0349} \SUPPORTEDANCH{\char"0199\char"0349} \SUPPORTEDANCH{\char"2C6A\char"0349} \SUPPORTEDANCH{\char"006C\char"0349} \SUPPORTEDANCH{\char"0234\char"0349} \SUPPORTEDANCH{\char"006E\char"0349} \SUPPORTEDANCH{\char"0235\char"0349} \SUPPORTEDANCH{\char"0070\char"0349} \SUPPORTEDANCH{\char"01A5\char"0349} \SUPPORTEDANCH{\char"0071\char"0349} \SUPPORTEDANCH{\char"02A0\char"0349} \SUPPORTEDANCH{\char"0074\char"0349} \SUPPORTEDANCH{\char"01AB\char"0349} \SUPPORTEDANCH{\char"01AD\char"0349} \SUPPORTEDANCH{\char"0236\char"0349} \SUPPORTEDANCH{\char"0075\char"0349} \MISSINGGLYPH{\char"AB52\char"0349} \SUPPORTEDANCH{\char"0079\char"0349} \SUPPORTEDANCH{\char"01B4\char"0349} \MISSINGGLYPH{\char"1EFF\char"0349}

\SUPPORTEDANCH{\char"0048\char"034D} \SUPPORTEDANCH{\char"2C67\char"034D} \SUPPORTEDANCH{\char"004B\char"034D} \SUPPORTEDANCH{\char"2C69\char"034D} \SUPPORTEDANCH{\char"004D\char"034D} \MISSINGGLYPH{\char"2C6E\char"034D} \SUPPORTEDANCH{\char"0050\char"034D} \SUPPORTEDANCH{\char"01A4\char"034D} \SUPPORTEDANCH{\char"0052\char"034D} \SUPPORTEDANCH{\char"2C64\char"034D} \SUPPORTEDANCH{\char"0054\char"034D} \SUPPORTEDANCH{\char"01AC\char"034D} \SUPPORTEDANCH{\char"01AE\char"034D} \SUPPORTEDANCH{\char"0059\char"034D} \SUPPORTEDANCH{\char"01B3\char"034D} \SUPPORTEDANCH{\char"0061\char"034D} \SUPPORTEDANCH{\char"1E9A\char"034D} \SUPPORTEDANCH{\char"0062\char"034D} \SUPPORTEDANCH{\char"0183\char"034D} \SUPPORTEDANCH{\char"0063\char"034D} \SUPPORTEDANCH{\char"0188\char"034D} \SUPPORTEDANCH{\char"0064\char"034D} \SUPPORTEDANCH{\char"018C\char"034D} \SUPPORTEDANCH{\char"0221\char"034D} \SUPPORTEDANCH{\char"0068\char"034D} \SUPPORTEDANCH{\char"2C68\char"034D} \SUPPORTEDANCH{\char"006B\char"034D} \SUPPORTEDANCH{\char"0199\char"034D} \SUPPORTEDANCH{\char"2C6A\char"034D} \SUPPORTEDANCH{\char"006C\char"034D} \SUPPORTEDANCH{\char"0234\char"034D} \SUPPORTEDANCH{\char"006E\char"034D} \SUPPORTEDANCH{\char"0235\char"034D} \SUPPORTEDANCH{\char"0070\char"034D} \SUPPORTEDANCH{\char"01A5\char"034D} \SUPPORTEDANCH{\char"0071\char"034D} \SUPPORTEDANCH{\char"02A0\char"034D} \SUPPORTEDANCH{\char"0074\char"034D} \SUPPORTEDANCH{\char"01AB\char"034D} \SUPPORTEDANCH{\char"01AD\char"034D} \SUPPORTEDANCH{\char"0236\char"034D} \SUPPORTEDANCH{\char"0075\char"034D} \MISSINGGLYPH{\char"AB52\char"034D} \SUPPORTEDANCH{\char"0079\char"034D} \SUPPORTEDANCH{\char"01B4\char"034D} \MISSINGGLYPH{\char"1EFF\char"034D}

\SUPPORTEDANCH{\char"0048\char"034E} \SUPPORTEDANCH{\char"2C67\char"034E} \SUPPORTEDANCH{\char"004B\char"034E} \SUPPORTEDANCH{\char"2C69\char"034E} \SUPPORTEDANCH{\char"004D\char"034E} \MISSINGGLYPH{\char"2C6E\char"034E} \SUPPORTEDANCH{\char"0050\char"034E} \SUPPORTEDANCH{\char"01A4\char"034E} \SUPPORTEDANCH{\char"0052\char"034E} \SUPPORTEDANCH{\char"2C64\char"034E} \SUPPORTEDANCH{\char"0054\char"034E} \SUPPORTEDANCH{\char"01AC\char"034E} \SUPPORTEDANCH{\char"01AE\char"034E} \SUPPORTEDANCH{\char"0059\char"034E} \SUPPORTEDANCH{\char"01B3\char"034E} \SUPPORTEDANCH{\char"0061\char"034E} \SUPPORTEDANCH{\char"1E9A\char"034E} \SUPPORTEDANCH{\char"0062\char"034E} \SUPPORTEDANCH{\char"0183\char"034E} \SUPPORTEDANCH{\char"0063\char"034E} \SUPPORTEDANCH{\char"0188\char"034E} \SUPPORTEDANCH{\char"0064\char"034E} \SUPPORTEDANCH{\char"018C\char"034E} \SUPPORTEDANCH{\char"0221\char"034E} \SUPPORTEDANCH{\char"0068\char"034E} \SUPPORTEDANCH{\char"2C68\char"034E} \SUPPORTEDANCH{\char"006B\char"034E} \SUPPORTEDANCH{\char"0199\char"034E} \SUPPORTEDANCH{\char"2C6A\char"034E} \SUPPORTEDANCH{\char"006C\char"034E} \SUPPORTEDANCH{\char"0234\char"034E} \SUPPORTEDANCH{\char"006E\char"034E} \SUPPORTEDANCH{\char"0235\char"034E} \SUPPORTEDANCH{\char"0070\char"034E} \SUPPORTEDANCH{\char"01A5\char"034E} \SUPPORTEDANCH{\char"0071\char"034E} \SUPPORTEDANCH{\char"02A0\char"034E} \SUPPORTEDANCH{\char"0074\char"034E} \SUPPORTEDANCH{\char"01AB\char"034E} \SUPPORTEDANCH{\char"01AD\char"034E} \SUPPORTEDANCH{\char"0236\char"034E} \SUPPORTEDANCH{\char"0075\char"034E} \MISSINGGLYPH{\char"AB52\char"034E} \SUPPORTEDANCH{\char"0079\char"034E} \SUPPORTEDANCH{\char"01B4\char"034E} \MISSINGGLYPH{\char"1EFF\char"034E}

}

\end{document}
